\xchapter{Procatechesis, or Prologue to the Catechetical Lectures of our Holy Father, Cyril, Archbishop of Jerusalem.}

1. \textsc{Already} there is an odour of blessedness upon you, O ye who are soon to be enlightened\footnote{The “blessedness” is the grace of Baptism, the hope of which is as a fragrant odour already borne towards the Candidates. These were called no longer Catechumens, but \textgreek{φωτιζόμενοι}, as already on the way “to be enlightened.” Compare xvi. 26, the last sentence, and see Index, “enlighten.”}: already ye are gathering the spiritual\footnote{\textgreek{νοητά}. The word is much used by Plato to distinguish things which can be discerned only by the mind from the objects of sight and sense. Here “the spiritual (or, mental) flowers” are the Divine truths in which “the fragrance of the Holy Spirit” breathes.} flowers, to weave heavenly crowns: already the fragrance of the Holy Spirit has breathed upon you: already ye have gathered round the vestibule of the King’s palace\footnote{By “the vestibule” is meant “the outer hall of the Baptistery” (xix. 2), and by “the King’s Palace” the Baptistery itself, which Cyril calls “the inner chamber” (xx. 1) and “the bride-chamber” (iii. 2; xxii. 2). See Index, “Baptistery.” Here the local terms have also an allegorical sense, Baptism being regarded as the marriage of the Soul to Christ.}; may ye be led in also by the King! For blossoms now have appeared upon the trees\footnote{Another allegory, from the season of Spring, when the Lectures were delivered.}; may the fruit also be found perfect! Thus far there has been an inscription of your names\footnote{\textgreek{ὀνοματογραφία}. See Index.}, and a call to service, and torches\footnote{That the Candidates on their first admission carried torches or lighted tapers in procession is a conjecture founded on this passage and Lect. I. 1: “Ye who have just lighted the torches of faith, preserve them in your hands unquenched.” But see Index, “Lights.”} of the bridal train, and a longing for heavenly citizenship, and a good purpose, and hope attendant thereon. For he lieth not who said, \emph{that to them that love God all things work together for good}. God is lavish in beneficence, yet He waits for each man’s genuine will: therefore the Apostle added and said, \emph{to them that are called according to a purpose}\footnote{Rom. viii. 28. In S. Paul’s argument the “purpose” is God’s eternal purpose of salvation through Christ (Eph. i. 11; iii. 11): but Cyril applies it here to sincerity of purpose in coming to Baptism.}. The honesty of purpose makes thee called: for if thy body be here but not thy mind, it profiteth thee nothing.

2. Even Simon Magus once came to the Laver\footnote{Acts viii. 13.}: he was baptized, but was not enlightened; and though he dipped his body in water, he enlightened not his heart with the Spirit: his body went down and came up, but his soul was not buried with Christ, nor raised with Him\footnote{Rom. vi. 4; Col. ii. 12.}. Now I mention the statements\footnote{Greek, \textgreek{ὑπογραφή}, meaning either an “indictment,” or a descriptive “sketch.” For the former meaning, see Plato, \emph{Theaet}. 172, E. \textgreek{ὑπογραφὴν …ἣν ἀντωμοσίαν καλοῦσιν}.} of (men’s) falls, that thou mayest not fall: for these things happened to them by way of example, \emph{and they are written for the admonition}\footnote{1 Cor. x. 11.} of those who to this day draw near. Let none of you be found tempting His grace, \emph{lest any root of bitterness spring up and trouble you}\footnote{Heb. xii. 15.}. Let none of you enter saying, Let us see what the faithful\footnote{“The faithful” are those who have been already baptized, and instructed in those mysteries of the Christian Faith which were reserved for the initiated. See Index, “Faithful.”} are doing: let me go in and see, that I may learn what is being done. Dost thou expect to see, and not expect to be seen? And thinkest thou, that whilst thou art searching out what is going on, God is not searching thy heart?

3. A certain man in the Gospels once pried into the marriage feast\footnote{Matt. xxii. 12. The same passage is applied to Baptism in Cat. iii. 2.}, and took an unbecoming garment, and came in, sat down, and ate: for the bridegroom permitted it. But when he saw them all clad in white\footnote{See Cat. xxii. 8 and Index, “White.”}, he ought to have assumed a garment of the same kind himself: whereas he partook of the like food, but was unlike them in fashion and in purpose. The bridegroom, however, though bountiful, was not undiscerning: and in going round to each of the guests and observing them (for his care was not for their eating, but for their seemly behaviour), he saw a stranger \emph{not having on a wedding garment}, and said to him, \emph{Friend, how camest thou in hither?} In what a colour\footnote{The Greed word (\textgreek{χρῶμα}) is used by Ignatius in the beginning of his \emph{Epistle to the Romans} of a discolouring stain.}! With what a conscience! What though the door-keeper forbade thee not, because of the bountifulness of the entertainer? what though thou wert ignorant in what fashion thou shouldest come in to the banquet?—thou didst come in, and didst see the glittering fashions of the guests: shouldest thou not have been taught even by what was before thine eyes? Shouldest thou not have retired in good season, that thou mightest enter in good season again? But now thou hast come in unseasonably, to be unseasonably cast out. So he commands the servants, \emph{Bind his feet}, which daringly intruded: \emph{bind his hands}, which knew not how to put a bright garment around him: \emph{and cast him into the outer darkness;} for he is unworthy of the wedding torches\footnote{Compare § 1, note 6.}. Thou seest what happened to that man: make thine own condition safe.

4. For we, the ministers of Christ, have admitted every one, and occupying, as it were, the place of door-keepers we left the door open: and possibly thou didst enter with thy soul bemired with sins, and with a will defiled. Enter thou didst, and wast allowed: thy name was inscribed. Tell me, dost thou behold this venerable constitution of the Church? Dost thou view her order and discipline\footnote{The Greek word (\textgreek{ἐπιστήμη}) which commonly means “knowledge” or “understanding,” is applied here and in vi. 35 to the intelligence and skill displayed in the arrangement of the public services of the Church. Compare \emph{Apostolic Constitutions}, ii. 57, where the Bishop is exhorted to have the assemblies arranged \textgreek{μετὰ πάσης ἐπιστήμης}.}, the reading of Scriptures\footnote{In the same passage of the Apostolic Constitutions precise directions are given for reading a Lesson from the Old Testament, singing the Psalms, and reading the Epistle and Gospel.}, the presence of the ordained\footnote{By “the ordained” (\textgreek{κανονικῶν}) are meant all whose names were registered as bearing office in the Church, Priests, Deacons, Deaconesses, Monks, Virgins, Widows, all having their appointed placed and proper duties. \emph{Apost. Canon}. 70, \textgreek{εἴ τις ἐπίσκοπος, ἢ πρεσβύτερος, ἢ διάκονος, ἢ ὅλως τοῦ καταλόγου τῶν κληρικῶν, κ.τ.λ}.}, the course of instruction\footnote{Compare \emph{Apost. Const.} as above: “Let the Presbyters one by one, not all together, exhort the people; and the Bishop last, as being the commander.”}? Be abashed at the place, and be taught by what thou seest\footnote{S. Aug. \emph{de Civit. Dei}., ii. 28: “Though some come to mock at such admonitions, all their insolence is either humbled by a sudden conversation (immutatio) or suppressed by fear or shame.”}. Go out opportunely now, and enter most opportunely to-morrow.

If the fashion of thy soul is avarice, put on another fashion and come in. Put off thy former fashion, cloke it not up. Put off, I pray thee, fornication and uncleanness, and put on the brightest robe of chastity. This charge I give thee, before Jesus the Bridegroom of souls come in and see their fashions. A long notice\footnote{Greek, \textgreek{προθεσμία}. Compare Gal. iv. 2: “the time appointed of the father.” At Athens it meant a “limitation,” or fixed period within which a debt must be claimed or paid, or an action commenced.} is allowed thee; thou hast forty\footnote{Index, “Lent.”} days for repentance: thou hast full opportunity both to put off, and wash, and to put on and enter. But if thou persist in an evil purpose, the speaker is blameless, but thou must not look for the grace: for the water will receive, but the Spirit will not accept thee\footnote{Compare xvii. 36.}. If any one is conscious of his wound, let him take the salve; if any has fallen, let him arise. Let there be no Simon among you, no hypocrisy, no idle curiosity about the matter.

5. Possibly too thou art come on another pretext. It is possible that a man is wishing to pay court to a woman, and came hither on that account\footnote{S. Ambrose on the 119th Psalm, \emph{Serm}. xx. § 48, speaks of some who pretended to be Christians in order to marry one whose parents would not give her in marriage to a heathen.}. The remark applies in like manner to women also in their turn. A slave also perhaps wishes to please his master, and a friend his friend. I accept this bait for the hook, and welcome thee, though thou camest with an evil purpose, yet as one to be saved by a good hope. Perhaps thou knewest not whither thou wert coming, nor in what kind of net thou art taken. Thou art come within the Church’s nets\footnote{Matt. xiii. 47.}: be taken alive, flee not: for Jesus is angling for thee, not in order to kill, but by killing to make alive: for thou must die and rise again. For thou hast heard the Apostle say, \emph{Dead indeed unto sin, but living unto righteousness\footnote{Rom. vi. 11, 14.}}. Die to thy sins, and live to righteousness, live from this very day.

6. See, I pray thee, how great a dignity Jesus bestows on thee. Thou wert called a Catechumen, while the word echoed\footnote{S. Cyril plays upon the word “Catechumen,” which has the same root as “echo.”} round thee from without; hearing of hope, and knowing it not; hearing mysteries, and not understanding them; hearing Scriptures, and not knowing their depth. The echo is no longer around thee, but within thee; for \emph{the indwelling Spirit\footnote{Rom. viii. 9, 11.}} henceforth makes thy mind a house of God. When thou shalt have heard what is written concerning the mysteries, then wilt thou understand things which thou knewest not. And think not that thou receivest a small thing: though a miserable man, thou receivest one of God’s titles. Hear St. Paul saying, \emph{God is faithful\footnote{1 Cor. i. 9.}}. Hear another Scripture saying, \emph{God is faithful and just\footnote{1 John i. 9.}}. Foreseeing this, the Psalmist, because men are to receive a title of God, spoke thus in the person of God: I said, \emph{Ye are Gods, and are all sons of the Most High}\footnote{Ps. lxxxi. 6.}. But beware lest thou have the title of “\emph{faithful},” but the will of the faithless. Thou hast entered into a contest, toil on through the race: another such opportunity thou canst not have\footnote{Compare xvii. 36.}. Were it thy wedding-day before thee, wouldest thou not have disregarded all else, and set about the preparation for the feast? And on the eve of consecrating thy soul to the heavenly Bridegroom, wilt thou not cease from carnal things, that thou mayest win spiritual?

7. We may not receive Baptism twice or thrice; else it might be said, Though I have failed once, I shall set it right a second time: whereas if thou fail once, the thing cannot be set right; for there is \emph{one Lord, and one faith, and one baptism}\footnote{Eph. iv. 5.}: for only the heretics are re-baptized\footnote{This sentence is omitted in one \textsc{ms.} (Paris, 1824), but probably only through the repetition of the word “baptism.” On the laws of the Church against the repetition of Baptism, and concerning the re-baptism of heretics, see Tertull. \emph{de Baptismo}, c. xv: \emph{Apost. Const.} xv.: Bingham, xii. 5: Hefele, \emph{Councils}, Lib. I. c. 2: Dictionary Christian Antiq. I. p. 167 a.}, because the former was no baptism.

8. For God seeks nothing else from us, save a good purpose. Say not, How are my sins blotted out? I tell thee, By willing, by believing\footnote{Rufinus, in the \emph{Exposition of the Creed}, on the \emph{Remission of sins}: “The Pagans are wont to say in derision of us, that we deceive ourselves in thinking that crimes which have been committed in deed can be washed out by words.”}. What can be shorter than this? But if, while thy lips declare thee willing, thy heart be silent, He knoweth the heart, who judgeth thee. Cease from this day from every evil deed. Let not thy tongue speak unseemly words, let thine eye abstain from sin, and from roving\footnote{The reading in the Benedictine Edition, \textgreek{μηδὲ ὁ νοῦς σου ῥεμβέσθω}, has little authority, and is quite unsuitable. See below, \textgreek{τὸ βλέμμα ῥεμβόμενον}.} after things unprofitable.

9. Let thy feet hasten to the catechisings; receive with earnestness the exorcisms\footnote{Index, “Exorcism.”}: whether thou be breathed upon or exorcised, the act is to thee salvation. Suppose thou hast gold unwrought and alloyed, mixed with various substances, copper, and tin, and iron, and lead: we seek to have the gold alone; can gold be purified from the foreign substances without fire? Even so without exorcisms the soul cannot be purified; and these exorcisms are divine, having been collected out of the divine Scriptures. Thy face has been veiled\footnote{Index, “Veiling.”}, that thy mind may henceforward be free, lest the eye by roving make the heart rove also. But when thine eyes are veiled, thine ears are not hindered from receiving the means of salvation. For in like manner as those who are skilled in the goldsmith’s craft throw in their breath upon the fire through certain delicate instruments, and blowing up the gold which is hidden in the crucible stir the flame which surrounds it, and so find what they are seeking; even so when the exorcists inspire terror by the Spirit of God, and set the soul, as it were, on fire in the crucible of the body, the hostile demon flees away, and there abide salvation and the hope of eternal life, and the soul henceforth is cleansed from its sins and hath salvation. Let us then, brethren, abide in hope, and surrender ourselves, and hope, in order that the God of all may see our purpose, and cleanse us from our sins, and impart to us good hopes of our estate, and grant us repentance that bringeth salvation. God hath called, and His call is to thee.

10. Attend closely to the catechisings, and though we should prolong our discourse, let not thy mind be wearied out. For thou art receiving armour against the adverse power, armour against heresies, against Jews, and Samaritans\footnote{The Samaritans are frequently mentioned by Epiphanius and other writers of the 4th century among the chief adversaries of Christianity. “In their humble synagogue, at the foot of the mountain (Gerizim), the Samaritans still worship, the oldest and the smallest sect in the world.” (Stanley, \emph{Sinai and Palestine}, p. 240.)}, and Gentiles. Thou hast many enemies; take to thee many darts, for thou hast many to hurl them at: and thou hast need to learn how to strike down the Greek, how to contend against heretic, against Jew and Samaritan. And the armour is ready, and most ready \emph{the sword of the Spirit}\footnote{Eph. vi. 17.}: but thou also must stretch forth thy right hand with good resolution, that thou mayest war the Lord’s warfare, and overcome adverse powers, and become invincible against every heretical attempt.

11. Let me give thee this charge also. Study our teachings and keep them for ever. Think not that they are the ordinary homilies\footnote{See above, § 4, note 3.}; for though they also are good and trustworthy, yet if we should neglect them to-day we may study them to-morrow. But if the teaching concerning the laver of regeneration delivered in a consecutive course be neglected to-day, when shall it be made right? Suppose it is the season for planting trees: if we do not dig, and dig deep, when else can that be planted rightly which has once been planted ill? Suppose, pray, that the Catechising is a kind of building: if we do not bind the house together by regular bonds in the building, lest some gap be found, and the building become unsound, even our former labour is of no use. But stone must follow stone by course, and corner match with corner, and by our smoothing off inequalities the building must thus rise evenly. In like manner we are bringing to thee stones, as it were, of knowledge. Thou must hear concerning the living God, thou must hear of Judgment, must hear of Christ, and of the Resurrection. And many things there are to be discussed in succession, which though now dropped one by one are afterwards to be presented in harmonious connexion. But unless thou fit them together in the one whole, and remember what is first, and what is second, the builder may build, but thou wilt find the building unsound.

12. When, therefore, the Lecture is delivered, if a Catechumen ask thee what the teachers have said, tell nothing to him that is without\footnote{On the Disciplina Arcani, or rule against publishing the Christian Creed and Mysteries to Catechumens and Gentiles, see Index, “\emph{Mysteries}.”}. For we deliver to thee a mystery, and a hope of the life to come. Guard the mystery for Him who gives the reward. Let none ever say to thee, What harm to thee, if I also know it? So too the sick ask for wine; but if it be given at a wrong time it causes delirium, and two evils arise; the sick man dies, and the physician is blamed. Thus is it also with the Catechumen, if he hear anything from the believer: both the Catechumen becomes delirious (for he understands not what he has heard, and finds fault with the thing, and scoffs at what is said), and the believer is condemned as a traitor. But thou art now standing on the border: take heed, pray, to tell nothing out; not that the things spoken are not worthy to be told, but because his ear is unworthy to receive. Thou wast once thyself a Catechumen, and I described not what lay before thee. When by experience thou hast learned how high are the matters of our teaching, then thou wilt know that the Catechumens are not worthy to hear them.

13. Ye who have been enrolled are become sons and daughters of one Mother. When ye have come in before the hour of the exorcisms, let each one of you speak things tending to godliness: and if any of your number be not present, seek for him. If thou wert called to a banquet, wouldest thou not wait for thy fellow guest? If thou hadst a brother, wouldest thou not seek thy brother’s good?

Afterwards busy not thyself about unprofitable matters: neither, what the city has done, nor the village, nor the King\footnote{The title “King” (\textgreek{Βασιλεύς}) is used in the Greek Liturgies and Fathers of the Roman Emperor, as in the Clementine Liturgy: \textgreek{ὑπὲρ τοῦ βασιλέως, καὶ τῶν ἐν ὑπεροχῇ}, where it is taken from 1 Tim. ii. 2. Compare Cat. xiv. 14, and 22: \textgreek{Κωνσταντίνου τοῦ βασιλέως}.}, nor the Bishop, nor the Presbyter. Look upward; that is what thy present hour needeth. \emph{Be still}\footnote{Ps. xlvi. 10. Sept. \textgreek{σχολάσατε}, “give attention freely.”}, \emph{and know that I am God}. If thou seest the believers ministering, and shewing no care, they enjoy security, they know what they have received, they are in possession of grace. But thou standest just now in the turn of the scale, to be received or not: copy not those who have freedom from anxiety, but cherish fear.

14. And when the Exorcism has been done, until the others who are being exorcised have come\footnote{From S. Augustine, \emph{de Symbolo}, i. 1 (Migne T. vi. p. 930), we learn that the Candidates were brought in before the Congregation one by one for exorcism; and so, as Cyril here shews, they had to wait outside till the others returned.}, let men be with men, and women with women. For now I need the example of Noah’s ark: in which were Noah and his sons, and his wife and his sons’ wives. For though the ark was one, and the door was shut, yet had things been suitably arranged. If the Church is shut, and you are all inside, yet let there be a separation, men with men, and women with women\footnote{Chrys. \emph{in Matt. Hom}. lxxiv. § 3: “You ought to have within you the wall that separates you from the women: but since ye will not, our fathers have thought it necessary to separate you at least by these boards; for I have heard from my elders that there were not these walls in old times.” These barriers had not yet been introduced at Jerusalem, or Cyril’s admonition would have been needless. Compare \emph{Apostolic Constitutions}, II. 57.}: lest the pretext of salvation become an occasion of destruction. Even if there be a fair pretext for sitting near each other, let passions be put away. Further, let the men when sitting have a useful book; and let one read, and another listen: and if there be no book, let one pray, and another speak something useful. And again let the party of young women sit together in like manner, either singing or reading quietly, so that their lips speak, but others’ ears catch not the sound: \emph{for I suffer not a woman to speak in the Church}\footnote{1 Cor. xiv. 34; 1 Tim. ii. 12.}. And let the married woman also follow the same example, and pray; and let her lips move, but her voice be unheard, that a Samuel\footnote{1 Sam. i. 13, 20. On the various interpretations of the name Samuel, see \emph{Dict. Bib.} “Samuel,” and Driver on the passage. Cyril adopts the meaning “heard of God.”} may come, and thy barren soul give birth to the salvation of “God who hath heard thy prayer;” for this is the interpretation of the name Samuel.

15. I shall observe each man’s earnestness, each woman’s reverence. Let your mind be refined as by fire unto reverence; let your soul be forged as metal: let the stubbornness of unbelief be hammered out: let the superfluous scales of the iron drop off, and what is pure remain; let the rust of the iron be rubbed off, and the true metal remain. May God sometime shew you that night, the darkness which shines like the day, concerning which it is said, \emph{The darkness shall not be hidden from thee, and the night shall shine as the day}\footnote{Ps. cxxxix. 12. On Easter Eve the Church was full of lights which were kept burning all night, and the newly-baptized carried torches. Gregory of Nyssa, preaching on the Resurrection (\emph{Orat}. iv.) describes the scene: “This brilliant night, by mingling the flames of torches with the morning rays of the sun, has made one continuous day, not divided by the interposition of darkness.”}. Then may the gate of Paradise be opened to every man and every woman among you. Then may you enjoy the Christ-bearing waters in their fragrance\footnote{Or, as the Benedictine Editor conjectures, “the waters which have a Christ-bearing (\textgreek{χριστοφόρον}) fragrance.” On the epithet \textgreek{χριστοφόρος}, see Bishop Lightfoot’s note on Ignat. \emph{ad Eph}. § 1 and § 9. Its meaning, as well as that of \textgreek{Θεοφόρος} is defined in the answer of Ignatius to Trajan, \textgreek{῾Ο Χριστὸν ἔχων ἐν στέρνοις} (\emph{Martyr. Ign. Ant.} § 2).}. Then may you receive the name of Christ\footnote{Cat. xxi. 1: “made partakers therefore of Christ, ye are rightly called Christs.”}, and the power of things divine. Even now, I beseech you, lift up the eye of the mind: even now imagine the choirs of Angels, and God the Lord of all there sitting, and His Only-begotten Son sitting with Him on His right hand, and the Spirit present with them; and Thrones and Dominions doing service, and every man of you and every woman receiving salvation. Even now let your ears ring, as it were, with that glorious sound, when over your salvation the angels shall chant, \emph{Blessed are they whose iniquities are forgiven, and whose sins are covered}\footnote{Ps. xxxii. 1, which verse is still chanted in the Greek Church as soon as the Baptism is completed.}: when like stars of the Church you shall enter in, bright in the body and radiant in the soul.

16. Great is the Baptism that lies before you\footnote{S. Basil has a passage in praise of Baptism almost the same, word for word, with this. It is more likely to have been borrowed from Cyril by Basil and other Fathers, than to be a later interpolation here.}: a ransom to captives; a remission of offences; a death of sin; a new-birth of the soul; a garment of light; a holy indissoluble seal; a chariot to heaven; the delight of Paradise; a welcome into the kingdom; the gift of adoption! But there is a serpent by the wayside watching those who pass by: beware lest he bite thee with unbelief. He sees so many receiving salvation, and is \emph{seeking whom he may devour}\footnote{1 Pet. v. 8.}. Thou art coming in unto the Father of Spirits, but thou art going past that serpent. How then mayest thou pass him? Have \emph{thy feet shod with the preparation of the gospel of peace}\footnote{Eph. vi. 15.}; that even if he bite, he may not hurt thee. Have faith in-dwelling, stedfast hope, a strong sandal, that thou mayest pass the enemy, and enter the presence of thy Lord. Prepare thine own heart for reception of doctrine, for fellowship in holy mysteries. Pray more frequently, that God may make thee worthy of the heavenly and immortal mysteries. Cease not day nor night: but when sleep is banished from thine eyes, then let thy mind be free for prayer. And if thou find any shameful thought rise up in thy mind, turn to meditation upon Judgment to remind thee of Salvation. Give thy mind wholly to study, that it may forget base things. If thou find any one saying to thee, Art thou then going in, to descend into the water? Has the city just now no baths? take notice that it is \emph{the dragon of the sea}\footnote{Is. xxvii. 1.} who is laying these plots against thee. Attend not to the lips of the talker, but to God who worketh in thee. Guard thine own soul, that thou be not ensnared, to the end that abiding in hope thou mayest become an heir of everlasting salvation.

17. We for our part as men charge and teach you thus: but make not ye our building \emph{hay and stubble} and chaff, lest we \emph{suffer loss}, from our \emph{work being burnt up:} but make ye our work gold, \emph{and silver, and precious stones}\footnote{1 Cor. iii. 12, 15.}! For it lies in me to speak, but in thee to set thy mind\footnote{Greek \textgreek{προσθέσθαι}, Sept. Deut. xiii. 4, “cleave unto Him.” Compare Josh. xxiii. 12; Ps. lxii. 10, “Set not your heart upon them.”} upon it, and in God to make perfect. Let us nerve our minds, and brace up our souls, and prepare our hearts. The race is for our soul: our hope is of things eternal: and God, who knoweth your hearts, and observeth who is sincere, and who a hypocrite, is able both to guard the sincere, and to give faith to the hypocrite: for even to the unbeliever, if only he give his heart, God is able to give faith. So may He \emph{blot out the handwriting that is against you}\footnote{Col. ii. 14.}, and grant you forgiveness of your former trespasses; may He plant you into His Church, and enlist you in His own service, and put on you \emph{the armour of righteousness}\footnote{2 Cor. vi. 7; Rom. vi. 13.}: may He fill you with the heavenly things of the New Covenant, and give you the seal of the Holy Spirit indelible throughout all ages, in Christ Jesus Our Lord: to whom be the glory for ever and ever! Amen.

(\emph{To the Reader}\footnote{It is doubtful whether this caution proceeded from Cyril himself when issuing a written copy of his Lectures, or from some later editor. Eusebius (\emph{E.H.} v. 20) has preserved an adjuration by Irenæus at the end of his treatise, \emph{On the Ogdoad}: I adjure thee, who mayest transcribe this book, by Our Lord Jesus Christ, and by His glorious advent, when He cometh to judge the quick and the dead, to compare what thou hast written and correct it carefully by this copy, from which thou hast transcribed it; this adjuration also thou shalt write in like manner, and set it in the copy.}.)

These Catechetical Lectures for those who are to be enlightened thou mayest lend to candidates for Baptism, and to believers who are already baptized, to read, but give not at all\footnote{Gr. \textgreek{τὸ σύνολον}. Plat. Leg. 654 B; Soph. 220 B.}, neither to Catechumens, nor to any others who are not Christians, as thou shalt answer to the Lord. And if thou make a copy, write this in the beginning, as in the sight of the Lord.

\xchapter{Lecture I. To those who are to be Enlightened, delivered extempore at Jerusalem, as an Introductory Lecture to those who had come forward for Baptism.}

To those who are to be Enlightened, delivered extempore at Jerusalem, as an Introductory Lecture to those who had come forward for Baptism\footnote{The title prefixed to this Lecture is given in full. In the following Lectures the form will be abbreviated. See Index, \textgreek{ἀνάγνωσις} and \textgreek{σχεδιασθεῖσα}.}:

\textsc{With a reading from} Isaiah

\emph{Wash you, make you clean; put away your iniquities from your souls, from before mine eyes}, and the rest\footnote{Is. i. 16.}.

1. \textsc{Disciples} of the New Testament and partakers of the mysteries of Christ, as yet by calling only, but ere long by grace also, \emph{make you a new heart and a new spirit}\footnote{Ezek. xviii. 31.}, that there may be gladness among the inhabitants of heaven: for \emph{if over one sinner that repenteth there is joy}, according to the Gospel\footnote{Luke xv. 7.}, how much more shall the salvation of so many souls move the inhabitants of heaven to gladness. As ye have entered upon a good and most glorious path, run with reverence the race of godliness. For the Only-begotten Son of God is present here most ready to redeem you, saying, \emph{Come unto Me all that labour and are heavy laden, and I will give you rest}\footnote{Matt. xi. 28.}. Ye that are clothed with the rough garment\footnote{Compare xv. 25.} of your offences, who are \emph{holden with the cords of your own sins}, hear the voice of the Prophet saying, \emph{Wash you, make you clean, put away your iniquities from before Mine eyes}\footnote{Is. i. 16.}: that the choir of Angels may chant over you, \emph{Blessed are they whose iniquities are forgiven, and whose sins are covered}\footnote{Ps. xxxii. 1. See Procat. 15.}. Ye who have just lighted the torches of faith\footnote{Procat. 1, note 6.}, guard them carefully in your hands unquenched; that He, who erewhile on this all-holy Golgotha opened Paradise to the robber on account of his faith, may grant to you to sing the bridal song.

2. If any here is a slave of sin, let him promptly prepare himself through faith for the new birth into freedom and adoption; and having put off the miserable bondage of his sins, and taken on him the most blessed bondage of the Lord, so may he be counted worthy to inherit the kingdom of heaven. \emph{Put off}, by confession\footnote{See Index, “Confession.”}, \emph{the old man, which waxeth corrupt after the lusts of deceit}, that ye may \emph{put on the new man, which is renewed according to knowledge of Him that created him}\footnote{Eph. iv. 22; Col. iii. 10.}. Get you \emph{the earnest of the Holy Spirit}\footnote{2 Cor. i. 22.} through faith, that ye may be able to be received \emph{into the everlasting habitations}\footnote{Luke xvi. 9.}. Come for the mystical Seal, that ye may be easily recognised by the Master; be ye numbered among the holy and spiritual flock of Christ, to be set apart on His right hand, and inherit the life prepared for you. For they to whom the rough garment\footnote{Compare xv. 25.} of their sins still clings are found on the left hand, because they came not to the grace of God which is given through Christ at the new birth of Baptism: new birth I mean not of bodies, but the spiritual new birth of the soul. For our bodies are begotten by parents who are seen, but our souls are begotten anew through faith: \emph{for the Spirit bloweth where it listeth}\footnote{John iii. 8.}: and then, if thou be found worthy, thou mayest hear, \emph{Well done, good and faithful servant}\footnote{Matt. xxv. 21.}, when thou art found to have no defilement of hypocrisy in thy conscience.

3. For if any of those who are present should think to tempt God’s grace, he deceives himself, and knows not its power. Keep thy soul free from hypocrisy, O man, because of Him \emph{who searcheth hearts and reins}\footnote{Ps. vii. 10.}. For as those who are going to make a levy for war examine the ages and the bodies of those who are taking service, so also the Lord in enlisting souls examines their purpose: and if any has a secret hypocrisy, He rejects the man as unfit for His true service; but if He finds one worthy, to him He readily gives His grace. He gives not holy things to the dogs\footnote{Matt. vii. 6.}; but where He discerns the good conscience, there He gives the Seal of salvation, that wondrous Seal, which devils tremble at, and Angels recognise; that the one may be driven to flight, and the others may watch around it as kindred to themselves. Those therefore who receive this spiritual and saving Seal, have need also of the disposition akin to it. For as a writing-reed or a dart has need of one to use it, so grace also has need of believing minds.

4. Thou art receiving not a perishable but a spiritual shield. Henceforth thou art planted in the invisible\footnote{Gr. \textgreek{νοητόν}, i.e. the true Paradise, to be seen by the mind, not by the eye. Apoc. xii. 7, 17.} Paradise. Thou receivest a new name, which thou hadst not before. Heretofore thou wast a Catechumen, but now thou wilt be called a Believer. Thou art transplanted henceforth among the spiritual\footnote{See preceding note.} olive-trees, being grafted from the wild into the good olive-tree\footnote{Rom. xi. 24.}, from sins into righteousness, from pollutions into purity. Thou art made partaker of the Holy Vine\footnote{John xv. 1, 4, 5.}. Well then, if thou abide in the Vine, thou growest as a fruitful branch; but if thou abide not, thou wilt be consumed by the fire. Let us therefore bear fruit worthily. God forbid that in us should be done what befell that barren fig-tree\footnote{Matt. xxi. 19.}, that Jesus come not even now and curse us for our barrenness. But may all be able to use that other saying, \emph{But I am like a fruitful olive-tree in the house of God: I have trusted in the mercy of God for ever}\footnote{Ps. lii. 10.},—an olive-tree not to be perceived by sense, but by the mind\footnote{\textgreek{νοητή}, see note 1, above.}, and full of light. As then it is His part to plant and to water\footnote{1 Cor. iii. 6. When Paul plants and Apollos waters, it is God Himself who works through His ministers.}, so it is thine to bear fruit: it is God’s to grant grace, but thine to receive and guard it. Despise not the grace because it is freely given, but receive and treasure it devoutly.

5. The present is the season of confession: confess what thou hast done in word or in deed, by night or by day; \emph{confess in an acceptable time, and in the day of salvation}\footnote{2 Cor. vi. 2.} receive the heavenly treasure. Devote thy time to the Exorcisms: be assiduous at the Catechisings, and remember the things that shall be spoken, for they are spoken not for thine ears only, but that by faith thou mayest seal them up in the memory. Blot out from thy mind all earthly\footnote{Literally “human.”} care: for thou art running for thy soul. Thou art utterly forsaking the things of the world: little are the things which thou art forsaking, great what the Lord is giving. Forsake things present, and put thy trust in things to come. Hast thou run so many circles of the years busied in vain about the world, and hast thou not forty days to be free (for prayer\footnote{Some \textsc{mss.} omit \textgreek{τῇ προσευχῇ} after \textgreek{σχολάζεις}.}), for thine own soul’s sake? \emph{Be still}\footnote{Ps. xlvi. 10: \textgreek{σχολάσατε}. Compare Procat. 13.}, \emph{and know that I am God}, saith the Scripture. Excuse thyself from talking many idle words: neither backbite, nor lend a willing ear to backbiters; but rather be prompt to prayer. Shew in ascetic exercise that thy heart is nerved\footnote{Compare Procat. 17: xviii. 1.}. Cleanse thy vessel, that thou mayest receive grace more abundantly. For though remission of sins is given equally to all, the communion of the Holy Ghost is bestowed in proportion to each man’s faith. If thou hast laboured little, thou receivest little; but if thou hast wrought much, the reward is great. Thou art running for thyself, see to thine own interest.

6. If thou hast aught against any man, forgive it: thou comest here to receive forgiveness of sins, and thou also must forgive him that hath sinned against thee. Else with what face wilt thou say to the Lord, Forgive me my many sins, if thou hast not thyself forgiven thy fellow-servant even his little sins. Attend diligently the Church assemblies\footnote{See Index, \textgreek{σύναξις}.}; not only now when diligent attendance is required of thee by the Clergy, but also after thou hast received the grace. For if, before thou hast received it, the practice is good, is it not also good after the bestowal? If before thou be grafted in, it is a safe course to be watered and tended, is it not far better after the planting? Wrestle for thine own soul, especially in such days as these. Nourish thy soul with sacred readings; for the Lord hath prepared for thee a spiritual table; therefore say thou also after the Psalmist, \emph{The Lord is my shepherd, and I shall lack nothing: in a place of grass, there hath He made me rest; He hath fed me beside the waters of comfort, He hath converted my soul}\footnote{Ps. xxiii. 1–3.}:—that Angels also may share your joy, and Christ Himself the great High Priest, having accepted your resolve, may present you all to the Father, saying, \emph{Behold, I and the children whom God hath given Me}\footnote{Is. viii. 18; Heb. ii. 13.}. May He keep you all well-pleasing in His sight! To whom be the glory, and the power unto the endless ages of eternity. Amen.

\xchapter{Lecture II. On Repentance and Remission of Sins, and Concerning the Adversary.}

\textsc{Ezekiel xviii. 20–23}

The righteousness of the righteous shall be upon him, and the wickedness of the wicked shall be upon him. But if the wicked will turn from all his sins, \&c.

1. \textsc{A fearful} thing is sin, and the sorest disease of the soul is transgression, secretly cutting its sinews, and becoming also the cause of eternal fire; an evil of a man’s own choosing, an offspring of the will.\footnote{For references to Cyril’s doctrine of Free-will, see Index, “Soul.”} For that we sin of our own free will the Prophet says plainly in a certain place: \emph{Yet I planted thee a fruitful vine, wholly true: how art thou turned to bitterness, (and become) the strange vine}\footnote{Jer. ii. 21.}? The planting was good, the fruit coming from the will is evil; and therefore the planter is blameless, but the vine shall be burnt with fire since it was planted for good, and bore fruit unto evil of its own will. \emph{For God}, according to the Preacher, \emph{made man upright, and they have themselves sought out many inventions}\footnote{Eccles. vii. 29.}. \emph{For we are His workmanship}, says the Apostle, \emph{created unto good works, which God afore prepared, that we should walk in them}\footnote{Eph. ii. 10.}. So then the Creator, being good, created for good works; but the creature turned of its own free will to wickedness. Sin then is, as we have said, a fearful evil, but not incurable; fearful for him who clings to it, but easy of cure for him who by repentance puts it from him. For suppose that a man is holding fire in his hand; as long as he holds fast the live coal he is sure to be burned, but should he put away the coal, he would have cast away the flame also with it. If however any one thinks that he is not being burned when sinning, to him the Scripture saith, \emph{Shall a man wrap up fire in his bosom, and not burn his clothes}\footnote{Prov. vi. 27.}? For sin burns the sinews of the soul, [and breaks the spiritual bones of the mind, and darkens the light of the heart\footnote{Milles and the Benedictine Editor omit these clauses, but the more recent editions of Reischl and Alexandrides insert them on the authority of Munich, Jerusalem, and other good \textsc{mss.}}].

2. But some one will say, What can sin be? Is it a living thing? Is it an angel? Is it a demon? What is this which works within us? It is not an enemy, O man, that assails thee from without, but an evil shoot growing up out of thyself. \emph{Look right on with thine eyes}\footnote{Prov. iv. 25.}, and there is no lust. [Keep thine own, and\footnote{Omitted by recent editors with the best \textsc{mss.}}] seize not the things of others, and robbery has ceased\footnote{Gr. \textgreek{κεκοίμηται} “has fallen asleep.”}. Remember the Judgment, and neither fornication, nor adultery, nor murder, nor any transgression of the law shall prevail with thee. But whenever thou forgettest God, forthwith thou beginnest to devise wickedness and to commit iniquity.

3. Yet thou art not the sole author of the evil, but there is also another most wicked prompter, the devil. He indeed suggests, but does not get the mastery by force over those who do not consent. Therefore saith the Preacher, \emph{If the spirit of him that hath power rise up against thee, quit not thy place}\footnote{Eccles. x. 4. Compare Eph. iv. 27: “Neither give place to the devil.”}. Shut thy door, and put him far from thee, and he shall not hurt thee. But if thou indifferently admit the thought of lust, it strikes root in thee by its suggestions, and enthrals thy mind, and drags thee down into a pit of evils.

But perhaps thou sayest, I am a believer, and lust does not gain the ascendant over me, even if I think upon it frequently. Knowest thou not that a root breaks even a rock by long persistence? Admit not the seed, since it will rend thy faith asunder: tear out the evil by the root before it blossom, lest from being careless at the beginning thou have afterwards to seek for axes and fire. When thine eyes begin to be diseased, get them cured in good time, lest thou become blind, and then have to seek the physician.

4. The devil then is the first author of sin, and the father of the wicked: and this is the Lord’s saying, not mine, \emph{that the devil sinneth} \emph{from the beginning}\footnote{1 John iii. 8; John viii. 44.}: none sinned before him. But he sinned, not as having received necessarily from nature the propensity to sin, since then the cause of sin is traced back again to Him that made him so; but having been created good, he has of his own free will become a devil, and received that name from his action. For being an Archangel\footnote{On Cyril’s doctrine of the Angels, see Index, “Angels.”} he was afterwards called a devil from his slandering: from being a good servant of God he has become rightly named Satan; for “Satan” is interpreted \emph{the adversary}\footnote{1 Kings v. 4, \&c.}. And this is not my teaching, but that of the inspired prophet Ezekiel: for he takes up a lamentation over him and says, \emph{Thou wast a seal of likeness, and a crown of beauty; in the Paradise of God wast thou born}\footnote{Ezek. xxviii. 12–17, an obscure passage, addressed to the Prince of Tyre, and meaning that he was “the perfect pattern” of earthly glory, set in a condition like that of Adam in Paradise, and, seemingly, blameless as Adam before his fall. Cyril seems to regard the Prince of Tyre as an embodiment of Satan, because he was deified as the object of national worship: v. 1, “Thou hast said, I am a God, I sit in the seat of God.”}: and soon after, \emph{Thou wast born blameless in thy days, from the day in which thou wast created, until thine iniquities were found in thee.} Very rightly hath he said, \emph{were found in thee}; for they were not brought in from without, but thou didst thyself beget the evil. The cause also he mentions forthwith: \emph{Thine heart was lifted up because of thy beauty: for the multitude of thy sins wast thou wounded, and I did cast thee to the ground}. In agreement with this the Lord says again in the Gospels: \emph{I beheld Satan as lightning fall from heaven}\footnote{Luke x. 18.}. Thou seest the harmony of the Old Testament with the New. He when cast out drew many away with him. It is he that puts lusts into them that listen to him: from him come adultery, fornication, and every kind of evil. Through him our forefather Adam was cast out for disobedience, and exchanged a Paradise bringing forth wondrous fruits of its own accord for the ground which bringeth forth thorns.

5. What then? some one will say. We have been beguiled and are lost. Is there then no salvation left? We have fallen: Is it not possible to rise again? We have been blinded: May we not recover our sight? We have become crippled: Can we never walk upright? In a word, we are dead: May we not rise again? He that woke Lazarus who was four days dead and already stank, shall He not, O man, much more easily raise thee who art alive? He who shed His precious blood for us, shall Himself deliver us from sin. Let us not despair of ourselves, brethren; let us not abandon ourselves to a hopeless condition. For it is a fearful thing not to believe in a hope of repentance. For he that looks not for salvation spares not to add evil to evil: but to him that hopes for cure, it is henceforth easy to be careful over himself. The robber who looks not for pardon grows desperate; but, if he hopes for forgiveness, often comes to repentance. What then, does the serpent cast its slough\footnote{Literally, “its old age” (\textgreek{τὸ γῆρας}). Compare iii. 7, and Dict. Chr. Biogr., \emph{Macarius}, p. 770 a.}, and shall not we cast off our sin? Thorny ground also, if cultivated well, is turned into fruitful; and is salvation to us irrecoverable? Nay rather, our nature admits of salvation, but the will also is required.

6. God is loving to man, and loving in no small measure. For say not, I have committed fornication and adultery: I have done dreadful things, and not once only, but often: will He forgive? Will He grant pardon? Hear what the Psalmist says: \emph{How great is the multitude of Thy goodness, O Lord}\footnote{Ps. xxxi. 20.}! Thine accumulated offences surpass not the multitude of God’s mercies: thy wounds surpass not the great Physician’s skill. Only give thyself up in faith: tell the Physician thine ailment: say thou also, like David: \emph{I said, I will confess me my sin unto the Lord:} and the same shall be done in thy case, which he says forthwith: \emph{And thou forgavest the wickedness of my heart}\footnote{Ps. xxxii. 5.}.

7. Wouldest thou see the loving-kindness of God, O thou that art lately come to the catechising? Wouldest thou see the loving-kindness of God, and the abundance of His long-suffering? Hear about Adam. Adam, God’s first-formed man, transgressed: could He not at once have brought death upon him? But see what the Lord does, in His great love towards man. He casts him out from Paradise, for because of sin he was unworthy to live there; but He \emph{puts him to dwell over against Paradise}\footnote{This is the reading of the Septuagint instead of—“He placed at the east of the garden of Eden.”}: that seeing whence he had fallen, and from what and into what a state he was brought down, he might afterwards be saved by repentance. Cain the first-born man became his brother’s murderer, the inventor of evils, the first author of murders, and the first envious man. Yet after slaying his brother to what is he condemned? \emph{Groaning and trembling shalt thou be upon the earth}\footnote{Gen. iv. 12: “A fugitive and a vagabond shalt thou be upon the earth.”}. How great the offence, the sentence how light!

8. Even this then was truly loving-kindness in God, but little as yet in comparison with what follows. For consider what happened in the days of Noe. The giants sinned, and much wickedness was then spread over the earth, and because of this the flood was to come upon them: and in the five hundredth year God utters His threatening; but in the six hundredth He brought the flood upon the earth. Seest thou the breadth of God’s loving-kindness extending to a hundred years? Could He not have done immediately what He did then after the hundred years? But He extended (the time) on purpose, granting a respite for repentance. Seest thou God’s goodness? And if the men of that time had repented, they would not have missed the loving-kindness of God.

9. Come with me now to the other class, those who were saved by repentance. But perhaps even among women some one will say, I have committed fornication, and adultery, I have defiled my body by excesses of all kinds: is there salvation for me? Turn thine eyes, O woman, upon Rahab, and look thou also for salvation; for if she who had been openly and publicly a harlot was saved by repentance, is not she who on some one occasion before receiving grace committed fornication to be saved by repentance and fasting? For inquire how she was saved: this only she said: \emph{For your God is God in heaven and upon earth}\footnote{Josh. ii. 11.}. \emph{Your God;} for her own she did not dare to say, because of her wanton life. And if you wish to receive Scriptural testimony of her having been saved, you have it written in the Psalms: \emph{I will make mention of Rahab and Babylon among them that know me\footnote{Ps. lxxxvii. 4. “Rahab” is there a poetical name of Egypt, and the passage has nothing to do with Rahab the harlot. The Benedictine Editor rightly disregards S. Jerome’s suggestion, that Rahab is, like Egypt, a type of the Gentile Church.}}. O the greatness of God’s loving-kindness, making mention even of harlots in the Scriptures: nay, not simply \emph{I will make mention of Rahab and Babylon}, but with the addition, \emph{among them that know me}. There is then in the case both of men and of women alike the salvation which is ushered in by repentance.

10. Nay more, if a whole people sin, this surpasses not the loving-kindness of God. The people made a calf, yet God ceased not from His loving-kindness. Men denied God, but God denied not Himself\footnote{2 Tim. ii. 13.}. \emph{These be thy gods, O Israel}\footnote{Ex. xxxii. 4.}, they said: yet again, as He was wont, the God of Israel became their Saviour. And not only the people sinned, but also Aaron the High Priest. For it is Moses that says: \emph{And the anger of the Lord came upon Aaron: and I prayed for him}, saith he, \emph{and God forgave him}\footnote{Deut. ix. 20.}. What then, did Moses praying for a High Priest that sinned prevail with God, and shall not Jesus, His Only-begotten, prevail with God when He prays for us? And if He did not hinder Aaron, because of his offence, from entering upon the High Priesthood, will He hinder thee, who art come out from the Gentiles, from entering into salvation? Only, O man, repent thou also in like manner, and grace is not forbidden thee. Render thy way of life henceforth unblameable; for God is truly loving unto man, nor can all time\footnote{For “all time,” the reading of the best \textsc{mss.}, the Benedictine text has “all mankind.”} worthily tell out His loving kindness; nay, not if all the tongues of men unite together will they be able even so to declare any considerable part of His loving-kindness. For we tell some part of what is written concerning His loving-kindness to men, but how much He forgave the Angels we know not: for them also He forgives, since One alone is without sin, even Jesus who purgeth our sins. And of them we have said enough.

11. But if concerning us men thou wilt have other examples also set before thee\footnote{The Benedictine has, “But if thou wilt I will set before thee other examples also of our state? Come on to the blessed David.”}, come on to the blessed David, and take him for an example of repentance. Great as he was, he fell: after his sleep, walking in the eventide on the housetop, he cast a careless look, and felt a human passion. His sin was completed, but there died not with it his candour concerning the confession of his fault. Nathan the Prophet came, a swift accuser, and a healer of the wound. \emph{The Lord is wroth}, he says, \emph{and thou hast sinned}\footnote{2 Sam. xii.}. So spake the subject to the reigning king. But David the king\footnote{Bened. “The king, the wearer of the purple.”} was not indignant, for he regarded not the speaker, but God who had sent him. He was not puffed up\footnote{Bened. “blinded.”} by the array of soldiers standing round: for he had seen in thought the angel-host of the Lord, and he trembled \emph{as seeing Him who is invisible}\footnote{Heb. xi. 27.}; and to the messenger, or rather by him in answer to God who sent him, he said, \emph{I have sinned against the Lord}\footnote{2 Sam. xii. 13.}. Seest thou the humility of the king? Seest thou his confession? For had he been convicted by any one? Were many privy to the matter? The deed was quickly done, and straightway the Prophet appeared as accuser, and the offender confesses the fault. And because he candidly confessed, he received a most speedy cure. For Nathan the Prophet who had uttered the threat, said immediately, \emph{The Lord also hath put away thy sin}. Thou seest the swift relenting of a merciful God. He says, however, \emph{Thou hast greatly provoked the enemies of the} \emph{Lord}. Though thou hadst many enemies because of thy righteousness, thy self-control protected thee; but now that thou hast surrendered thy strongest armour, thine enemies are risen up, and stand ready against thee.

12. Thus then did the Prophet comfort him, but the blessed David, for all he heard it said, \emph{The \textsc{Lord} hath put away thy sin}, did not cease from repentance, king though he was, but put on sackcloth instead of purple, and instead of a golden throne, he sat, a king, in ashes on the ground; nay, not only sat in ashes, but also had ashes for his food, even as he saith himself, \emph{I have eaten ashes as it were bread}\footnote{Ps. cii. 10.}. His lustful eye he wasted away with tears saying, \emph{Every night will I wash my couch, and water my bed with my tears}\footnote{Ib. vii. 7.}. When his officers besought him to eat bread he would not listen. He prolonged his fast unto seven whole days. If a king thus made confession oughtest not thou, a private person, to confess? Again, after Absalom’s insurrection, though there were many roads for him to escape, he chose to flee by the Mount of Olives, in thought, as it were, invoking the Redeemer who was to go up thence into the heavens\footnote{2 Sam. xvi. 10, 11.}. And when Shimei cursed him bitterly, he said, \emph{Let him alone}, for he knew that “to him that forgiveth it shall be forgiven\footnote{Resch. (\emph{Agrapha}, p. 137) quotes various forms of this saying from early writers, and regards it as a fragment of an extracanonical Gospel. But see Lightfoot, \emph{Clem. Rom}. c. xiii.}.”

13. Thou seest that it is good to make confession. Thou seest that there is salvation for them that repent. Solomon also fell but what saith he? \emph{Afterwards I repented}\footnote{Prov. xxiv. 32, Sept. Heb. “Set my heart.” The passage has no reference to repentance: it means, “I considered the field of the slothful.” Hilary, Ps. lii.; Ambrose, \emph{Apolog}. 1, \emph{Prophetæ David}, c. iii. and other Fathers affirm the repentance of Solomon. Augustine (\emph{c. Faustum}, Lib. xxii. c. 88) maintains that Scripture says nothing of his repentance or forgiveness. See Dante, \emph{Paradiso, Canto} x. 109.}. Ahab, too, the King of Samaria, became a most wicked idolater, an outrageous man, the murderer of the Prophets\footnote{1 Kings xviii. 4.}, a stranger to godliness, a coveter of other men’s fields and vineyards. Yet when by Jezebel’s means he had slain Naboth, and the Prophet Elias came and merely threatened him, he rent his garments, and put on sackcloth. And what saith the merciful God to Elias? \emph{Hast than seen how Ahab is pricked in the heart before Me}\footnote{Ib. xxi. 29.}? as if almost He would persuade the fiery zeal of the Prophet to condescend to the penitent. For He saith, \emph{I will not bring the evil in his days}. And though after this forgiveness he was sure not to depart from his wickedness, nevertheless the forgiving God forgave him, not as being ignorant of the future, but as granting a forgiveness corresponding to his present season of repentance. For it is the part of a righteous judge to give sentence according to each case that has occurred.

14. Again, Jeroboam was standing at the altar sacrificing to the idols: his hand became withered, because he commanded the Prophet who reproved him to be seized: but having by experience learned the power of the man before him, he says, \emph{Entreat the face of the Lord thy God}\footnote{1 Kings xiii. 6.}; and because of this saying his hand was restored again. If the Prophet healed Jeroboam, is Christ not able to heal and deliver thee from thy sins? Manasses also was utterly wicked, who sawed Isaiah asunder\footnote{Justin Martyr, \emph{Dialogue with Trypho}, § 120 charges the Jews with having cut out a passage referring to the death of Isaiah. Theophylact commenting on Heb. xi. 37, says: “They were sawn asunder, as Isaiah by Manasses: and they say that he was sawn with a wooden saw, that his punishment might be the more painful to him from being prolonged.” Jerome on Is. i. 10, says that he was slain because of his calling the Jews “princes of Sodom and people of Gomorra,” and because he said, “I saw the Lord sitting upon a throne, high and lifted up.”}, and was defiled with all kinds of idolatries, and \emph{filled Jerusalem with innocent blood}\footnote{2 Chron. xxxiii. 12, 13.}; but having been led captive to Babylon he used his experience of misfortune for a healing course of repentance: for the Scripture saith that Manasses \emph{humbled himself before the Lord, and prayed, and the Lord heard him, and brought him back to his kingdom}. If He who sawed the Prophet asunder was saved by repentance, shall not thou then, having done no such great wickedness, be saved?

15. Take heed lest without reason thou mistrust the power of repentance. Wouldst thou know what power repentance has? Wouldst thou know the strong weapon of salvation, and learn what the force of confession is? Hezekiah by means of confession routed a hundred and fourscore and five thousand of his enemies. A great thing verily was this, but still small in comparison with what remains to be told: the same king by repentance obtained the recall of a divine sentence which had already gone forth. For when he had fallen sick, Esaias said to him, \emph{Set thine house in order; for thou shalt die, and not live}\footnote{2 Kings xx. 1.}. What expectation remained, what hope of recovery, when the Prophet said, \emph{for thou shalt die}? Yet Hezekiah did not desist from repentance; but remembering what is written, \emph{When thou shalt turn and lament, then shalt thou be saved}\footnote{Is. xxx. 15.}, he turned to the wall, and from his bed lifting his mind to heaven (for thickness of walls is no hindrance to prayers sent up with devotion), he said, “Remember me, O Lord, for it is sufficient for my healing that Thou remember me. Thou art not subject to times, but art Thyself the giver of the law of life. For our life depends not on a nativity, nor on a conjunction of stars, as some idly talk; but both of life and its duration. Then art Thyself the Lawgiver according to Thy Will.” And he, who could not hope to live because of the prophetic sentence, had fifteen years added to his life, and for the sign the sun ran backward in his course. Well then, for Ezekias’ sake the sun turned back but for Christ the sun was eclipsed, not retracing his steps, but suffering eclipse\footnote{Isaiah xxxviii. 8.}, and therefore shewing the difference between them, I mean between Ezekias and Jesus. The former prevailed to the cancelling of God’s decree, and cannot Jesus grant remission of sins? Turn and bewail thyself, shut thy door, and pray to be forgiven, pray that He may remove from thee the burning flames. For confession has power to quench even fire, power to tame even lions\footnote{From this point the \textsc{mss.} differ so widely that the Benedictine Editor gives two complete recensions of the whole Lecture. The Codd. Coislin, Ottob. 2, and Grodec, with the editions of Prevot and Milles, forming as it were one family of \textsc{mss.}, constitute the received text. On the other hand the older Munich Codex, with Codd. Roe and Casaubon, exhibit a recension of the Lecture differing from the editions. Reischl wishing to retain the received text unaltered, though preferring the other in particular passages, intended to append the other recension complete, but having left his work half finished, failed to do so. The chief variations are given in the following notes.}.

16. But if thou disbelieve, consider what befel Ananias and his companions. What streams did they pour out\footnote{Roe and Casaubon (R.C.) add: “into the furnace of fire.”}? How many vessels\footnote{R.C. “What measure.”} of water could quench the flame that rose up forty-nine cubits high\footnote{Song of the Three Children, v. 24.}? Nay, but where the flame mounted up a little\footnote{R.C. “Much.”} too high, faith was there poured out as a river, and there spake they the spell against all ills\footnote{R.C. “A great stream of repentance was poured forth, when they said, For Thou art righteous,” \&c.}: \emph{Righteous art Thou, O Lord, in all the things that Thou hast done to us: for we have sinned, and transgressed Thy law}\footnote{Song of the Three Children, v. 4.}. And their repentance quelled the flames\footnote{R.C. “Did then repentance quench the flames of the furnace, and dost thou disbelieve that it is able also to quench the fire of hell?”}. If thou believest not that repentance is able to quench the fire of hell, learn it from what happened in regard to Ananias\footnote{The Gospel only says, “There was darkness over all the land.” An eclipse of the sun was impossible at the time of the Paschal full moon.}. But some keen hearer will say, Those men God rescued justly in that case: because they refused to commit idolatry, God gave them that power. And since this thought has occurred, I come next to a different example of penitence\footnote{R.C. “That the narrative is not appropriate to those who are here present. For it was because Ananias and his companions refused to worship the idol, that God gave them that marvellous power. Adapting myself, therefore, to such a hearer, and looking to the profusion of instances, I come next to a different example of repentance.”}.

17. What thinkest thou of Nabuchodonosor? Hast thou not heard out of the Scriptures that he was bloodthirsty, fierce\footnote{R.C. “most impious, and most fierce in temper.”}, lion-like in disposition? Hast thou not heard that he brought out the bones of the kings from their graves into the light\footnote{Jer. viii. 1; Baruch ii. 25.}? Hast thou not heard\footnote{“Knowest thou not…”} that he carried the people away captive? Hast thou not heard that he put out the eyes of the king, after he had already seen his children slain\footnote{2 Kings xxv. 7.}? Hast thou not heard that he brake in pieces\footnote{R.C. “carried off.”} the Cherubim? I do not mean the invisible\footnote{\textgreek{νοητά}. R.C. add “and heavenly.”} beings;—away with such a thought, O man\footnote{Omitted by R.C.},—but the sculptured images, and the mercy-seat, in the midst of which God spake with His voice\footnote{R.C. “But those which had been constructed in the Temple, which were over the mercy-seat of the Ark.” Besides the two Cherubim of solid gold which Moses placed on the two ends of the Mercy-seat (Ex. xxxvii. 7 ff.), Solomon set “within the oracle” two Cherubim of olive wood overlaid with gold, ten feet high with outstretched wings overshadowing the Ark (1 Kings vi. 23–26; viii. 6, 7). All these were either carried off or destroyed, when Nebuchadnezzar took away “all the treasures of the house of the Lord” and “cut in pieces all the vessels of gold which Solomon, King of Israel, had made in the Temple of the Lord” (2 Kings xxiv. 13; 1 Esdras i. 54; 2 Esdras x. 22). The Benedictine editor is concerned because Cyril has paid no attention to the strange fiction in 2 Maccabees ii. 4 that Jeremy the Prophet “commanded the Tabernacle and the Ark to go with him” to Mount Horeb, and there hid them, with the Altar of Incense, in a hollow cave, to remain “unknown until the time that God gathers His people again together.”}. The veil of the Sanctuary\footnote{The Greek word rendered “Sanctuary” is \textgreek{ἡ ἁγιωσύνη}, literally “the holiness.”} he trampled under foot: the altar of incense he took and carried away to an idol-temple\footnote{2 Chron. xxxvi. 7.}: all the offerings he took away: the Temple he burned from the foundations\footnote{R.C. “The veil of the Sanctuary he tore down, he overturned the altar, and took all the vessels and carried them away to an idol temple. The Temple itself he burned.”}. How great punishments did he deserve, for slaying kings, for setting fire to the Sanctuary, for taking the people captive, for setting the sacred vessels in the house of idols? Did he not deserve ten thousand deaths?

18. Thou hast seen the greatness of his evil deeds: come now to God’s loving-kindness. He was turned into a wild beast\footnote{R.C. Afterwards he was turned into a wild beast: “he who was like a wild beast and most cruel in disposition; but he was turned into a wild beast, not that he might perish, but that by repentance he might be saved.”}, he abode in the wilderness, he was scourged, that he might be saved. He had claws as a lion\footnote{R.C. “of birds.” See Dan. iv. 33.}; for he was a ravager of the Sanctuary. He had a lion’s mane: for he was a ravening and a roaring lion. He ate grass like an ox: for a brute beast he was, not knowing Him who had given him the kingdom. His body was wet from the dew; because after seeing the fire quenched by the dew he believed not\footnote{R.C. “after the midst of the furnace had become to Ananias and his companions as the tinkling breath of rain, he saw and believed not.”}. And what happened\footnote{R.C. “But afterwards he came to his senses and repented, as he says himself.”}? \emph{After this}, saith he, \emph{I, Nabuchodonosor, lifted up} \emph{mine eyes unto heaven, and I blessed the Most High, and to Him that liveth for ever I gave praise and glory}\footnote{Dan. iv. 34.}. When, therefore, he recognised the Most High\footnote{R.C. “And after he had been scourged many years, he gave praise to Him that liveth for ever, and acknowledged Him that had given him the kingdom, and recognised the King of kings. And though he had often sinned in deeds, on making confession only in words, he received the benefit of God’s unspeakable loving kindness. He who was of all men most wicked, by the Divine judgment and loving-kindness of God who chastised him, crowned himself again with the royal diadem, and recovered his imperial throne.”}, and sent up these words of thankfulness to God, and repented himself for what he had done, and recognised his own weakness, then God gave back to him the honour of the kingdom.

19. What then\footnote{R.C. “If then there is present among you any from among the Heathen who has ever spoken evil against Christians, or in times of persecution plotted against the Holy Churches, let him take Nabuchodonsor as an example of salvation: let him confess in like manner, that he may also find the like forgiveness. If any has been defiled by lust and passions, let him take up the repentance of the blessed David: if any has denied like Peter, let him die like him for the sake of the Lord Jesus. For He who to his tears begrudged not the Apostleship, will not refuse thee the gospel mysteries. And for women let Rahab be a pattern unto salvation, and for men the manifold examples mentioned of the men of old times.}? When Nabuchodonosor, after having done such deeds, had made confession, did God give him pardon and the kingdom, and when thou repentest shall He not give thee the remission of sins, and the kingdom of heaven, if thou live a worthy life? The \textsc{Lord} is loving unto man, and swift to pardon, but slow to punish. Let no man therefore despair of his own salvation. Peter, the chiefest and foremost of the Apostles, denied the Lord thrice before a little maid: but he repented himself, and wept bitterly. Now weeping shews the repentance of the heart: and therefore he not only received forgiveness for his denial, but also held his Apostolic dignity unforfeited.

20. Having therefore, brethren, many examples of those who have sinned and repented and been saved, do ye also heartily make confession unto the Lord, that ye may both receive the forgiveness of your former sins, and be counted worthy of the heavenly gift, and inherit the heavenly kingdom with all the saints in Christ Jesus; to Whom is the glory for ever and ever. Amen\footnote{R.C. “And be ye all of good hope, having regard to the lovingkindness of God; not that we may fall back into the same sins, but that having had the benefit of redemption, and lived in a manner worthy of His grace, we may be able to blot out the handwriting that is against us by good works; in the power of the Only-begotten, the Son of God, and our Lord Jesus Christ, with whom be glory to the Father, with the Holy Ghost, both now and ever, and unto all the ages of eternity. Amen.”}.

\xchapter{Lecture III. On Baptism.}

\textsc{Romans vi. 3, 4}

Or know ye not that all we who were baptized into Christ Jesus were baptized into His death? were buried therefore with Him by our baptism into death, \&c.

1. \emph{Rejoice, ye heavens, and let the earth be glad}\footnote{Ps. xcvi. 11.}, for those who are to be sprinkled with hyssop, and cleansed with the spiritual\footnote{The invisible or spiritual (\textgreek{νοητός}) hyssop is the cleansing power of the Holy Ghost in Baptism. Compare Ps. li. 7.} hyssop, the power of Him to whom at His Passion drink was offered on hyssop and a reed\footnote{S. Cyril here, and still more emphatically in xiii. 39, distinguishes the hyssop (John xix. 29) from the reed (Matt. xxvii. 48), implying that the sponge filled with vinegar was bound round with hyssop, and then fixed on a reed. Another opinion is that the reed itself was that of hyssop. See Dictionary of the Bible, “Hyssop.”}. And while the Heavenly Powers rejoice, let the souls that are to be united to the spiritual Bridegroom make themselves ready. For \emph{the voice} is heard \emph{of one crying in the wilderness, Prepare ye the way of the Lord}\footnote{Is. xl. 3.}. For this is no light matter, no ordinary and indiscriminate union according to the flesh\footnote{\textgreek{σωμάτων}.}, but the All-searching Spirit’s election according to faith. For the inter-marriages and contracts of the world are not made altogether with judgment: but wherever there is wealth or beauty, there the bridegroom speedily approves: but here it is not beauty of person, but the soul’s clear conscience; not the condemned Mammon, but the wealth of the soul in godliness.

2. Listen then, O ye children of righteousness, to John’s exhortation when he says, \emph{Make straight the way of the Lord}. Take away all obstacles and stumbling-blocks, that ye may walk straight onward to eternal life. Make ready the vessels\footnote{So in § 15, the soul is regarded as a vessel for receiving grace.} of the soul, cleansed by unfeigned faith, for reception of the Holy Ghost. Begin at once to wash your robes in repentance, that when called to the bride-chamber ye may be found clean. For the Bridegroom invites all without distinction, because His grace is bounteous; and the cry of loud-voiced heralds assembles them all: but the same Bridegroom afterwards separates those who have come in to the figurative marriage. O may none of those whose names have now been enrolled hear the words, \emph{Friend, how camest thou in hither, not having a wedding garment}\footnote{Matt. xxii. 12.}? But may you all hear, \emph{Well done, good and faithful servant; thou wast faithful over a few things, I will set thee over many things: enter thou into the joy of thy lord}\footnote{Matt. xxv. 21.}.

For now meanwhile thou standest outside the door: but God grant that you all may say, \emph{The King hath brought me into His chamber}\footnote{Cant. i. 4.}. \emph{Let my soul rejoice in the Lord: for He hath clothed me with a garment of salvation, and a robe of gladness: He hath crowned me with a garland as a bridegroom}\footnote{Is. lxi. 10. Compare Cant. iii. 11: \emph{Go forth, O ye daughters of Zion, and behold King Solomon, with the crown wherewith his mother hath crowned him in the day of his espousals.} In the passage of Isaiah the bridegroom’s crown is likened to the priestly mitre.}, \emph{and decked me with ornaments as a bride}: that the soul of every one of you may be found \emph{not having spot or wrinkle or any such thing}\footnote{Eph. v. 27.}; I do not mean before you have received the grace, for how could that be? since it is for remission of sins that ye have been called; but that, when the grace is to be given, your conscience being found uncondemned may concur with the grace.

3. This is in truth a serious matter, brethren, and you must approach it with good heed. Each one of you is about to be presented to God before tens of thousands of the Angelic Hosts: the Holy Ghost is about to seal\footnote{See Index, “Seal.”} your souls: ye are to be enrolled in the army of the Great King. Therefore make you ready, and equip yourselves, by putting on I mean, not bright apparel\footnote{Index, “White.”}, but piety of soul with a good conscience. Regard not the Laver as simple water, but rather regard the spiritual grace that is given with the water. For just as the offerings brought to the heathen altars\footnote{\textgreek{βωμοῖς} used of heathen altars only, in Septuagint and N.T.}, though simple in their nature, become defiled by the invocation of the idols\footnote{Both here and in xix. 7, Cyril speaks of things offered to idols just as S. Paul in 1 Cor. x. 20. The Benediction of the water of Baptism is found in the \emph{Apostolic Constitutions}, vii. 43: “Look down from heaven, and sanctify this water, and give it grace and power, that so he that is to be baptized according to the command of Thy Christ, may be crucified with Him, and may die with Him, and be buried with Him, and may rise with him to the adoption which is in him, that he may be dead to sin and live to righteousness.”}, so contrariwise the simple water having received the invocation of the Holy Ghost, and of Christ, and of the Father, acquires a new power of holiness.

4. For since man is of twofold nature, soul and body, the purification also is twofold, the one incorporeal for the incorporeal part, and the other bodily for the body: the water cleanses the body, and the Spirit seals the soul; that we may draw near unto God, \emph{having our heart sprinkled} by the Spirit, \emph{and our body washed with pure water}\footnote{Heb. x. 22.}. When going down, therefore, into the water, think not of the bare element, but look for salvation by the power of the Holy Ghost: for without both thou canst not possibly be made perfect\footnote{See the note on “the twofold grace perfected by water and the Spirit,” at the end of this Lecture.}. It is not I that say this, but the Lord Jesus Christ, who has the power in this matter: for He saith, \emph{Except a man be born anew} (and He adds the words) \emph{of water and of the Spirit, he cannot enter into the kingdom of God}\footnote{John iii. 3.}. Neither doth he that is baptized with water, but not found worthy of the Spirit, receive the grace in perfection; nor if a man be virtuous in his deeds, but receive not the seal by water, shall he enter into the kingdom of heaven. A bold saying, but not mine, for it is Jesus who hath declared it: and here is the proof of the statement from Holy Scripture. Cornelius was a just man, who was honoured with a vision of Angels, and had set up his prayers and alms-deeds as a good memorial\footnote{\textgreek{στηλή}, Sept. A pillar of stone, bearing an inscription, was a common form of memorial among the Israelites and other ancient nations. See Dictionary of the Bible, “Pillar.”} before God in heaven. Peter came, and the Spirit was poured out upon them that believed, and they spake with other tongues, and prophesied: and after the grace of the Spirit the Scripture saith that Peter \emph{commanded them to be baptized in the name of Jesus Christ}\footnote{Acts x. 48.}; in order that, the soul having been born again by faith\footnote{S. Cyril considers that Cornelius and his friends were regenerated, as the Apostles were, apart from Baptism; as August. \emph{Serm.} 269, \emph{n.} 2, and Chrysost. \emph{in Act. Apost. Hom.} 25, seem to do. R.W.C.}, the body also might by the water partake of the grace.

5. But if any one wishes to know why the grace is given by water and not by a different element, let him take up the Divine Scriptures and he shall learn. For water is a grand thing, and the noblest of the four visible elements of the world. Heaven is the dwelling-place of Angels, but the heavens are from the waters\footnote{Compare ix. 5.}: the earth is the place of men, but the earth is from the waters: and before the whole six days’ formation of the things that were made, \emph{the Spirit of God moved upon the face of the water}\footnote{Gen. i. 2.}. The water was the beginning of the world, and Jordan the beginning of the Gospel tidings: for Israel deliverance from Pharaoh was through the sea, and for the world deliverance from sins \emph{by the washing of water with the word}\footnote{Ephes. v. 26.} of God. Where a covenant is made with any, there is water also. After the flood, a covenant was made with Noah: a covenant for Israel from Mount Sinai, but \emph{with water, and scarlet wool, and hyssop}\footnote{Heb. ix. 19.}. Elias is taken up, but not apart from water: for first he crosses the Jordan, then in a chariot mounts the heaven. The high-priest is first washed, then offers incense; for Aaron first washed, then was made high-priest: for how could one who had not yet been purified by water pray for the rest? Also as a symbol of Baptism there was a laver set apart within the Tabernacle.

6. Baptism is the end of the Old Testament, and beginning of the New. For its author was John, than whom was \emph{none greater among them that are born of women}. The end he was of the Prophets: \emph{for all the Prophets and the law were until John}\footnote{Matt. xi. 13.}: but of the Gospel history he was the first-fruit. For it saith, \emph{The beginning of the Gospel of Jesus Christ, \&c}.: \emph{John came baptising in the wilderness}\footnote{Mark i. 1, 4.}. You may mention Elias the Tishbite who was taken up into heaven, yet he is not greater than John: Enoch was translated, but he is not greater than John: Moses was a very great lawgiver, and all the Prophets were admirable, but not greater than John. It is not I that dare to compare Prophets with Prophets: but their Master and ours, the Lord Jesus, declared it: \emph{Among them that are born of women there hath not risen a greater than John}\footnote{Matt. xi. 11.}: He saith not “among them that are born of virgins,” but \emph{of women}\footnote{From the Clementine Recognitions, I. 54 and 60, we learn that there were some who asserted that John was the Christ, and not Jesus, inasmuch as Jesus Himself declared that John was greater than all men, and all Prophets. The answer is there given, that John was greater than all who are born of women, yet not greater than the Son of Man.}. The comparison is between the great servant and his fellow-servants: but the pre-eminence and the grace of the Son is beyond comparison with servants. Seest thou how great a man God chose as the first minister of this grace?—a man possessing nothing, and a lover of the desert, yet no hater of mankind: who ate locusts, and winged his soul for heaven\footnote{The locust being winged suggest the idea of growing wings for the soul. Is. xl. 31: \textgreek{πτεροφυησουσιν ὡς ἀετοί}.}: feeding upon honey, and speaking things both sweeter and more salutary than honey: clothed with a garment of camel’s hair, and shewing in himself the pattern of the ascetic life; who also was sanctified by the Holy Ghost while yet he was carried in his mother’s womb. Jeremiah was sanctified, but did not prophesy, in the womb\footnote{Jer. i. 5.}: John alone while carried in the womb leaped for joy\footnote{Luke i. 44.}, and though he saw not with the eyes of flesh, knew his Master by the Spirit: for since the grace of Baptism was great, it required greatness in its founder also.

7. This man was baptizing in Jordan, and \emph{there went out unto him all Jerusalem}\footnote{Matt. iii. 5.}, to enjoy the first-fruits of baptisms: for in Jerusalem is the prerogative of all things good. But learn, O ye inhabitants of Jerusalem, how they that came out were baptized by him: \emph{confessing their sins}, it is said\footnote{Matt. iii. 6.}. First they shewed their wounds, then he applied the remedies, and to them that believed gave redemption from eternal fire. And if thou wilt be convinced of this very point, that the baptism of John is a redemption from the threat of the fire, hear how he says, \emph{O generation of vipers, who hath warned you to flee from the wrath to come}\footnote{Ib. iii. 7.}\emph{? Be not then henceforth a viper, but as thou hast been formerly a viper’s brood, put off, saith he, the slough}\footnote{The Greek word (\textgreek{ὑπόστασις}) is used by Polybius (xxxiv. 9) for the deposit of silver from crushed ore, and by Hippocrates for any sediment or deposit. Here it means, as the context clearly shews, the old skin cast by a snake. Compare ii. 5.} \emph{of thy former sinful life. For every serpent creeps into a hole and casts its old slough, and having rubbed off the old skin, grows young again in body. In like manner enter thou also through the strait and narrow gate}\footnote{Matt. vii. 13, 14.}: rub off thy former self by fasting, and drive out that which is destroying thee. \emph{Put off the old man with his doings}\footnote{Col. iii. 9.}, and quote that saying in the Canticles, \emph{I have put off my coat, how shall I put it on}\footnote{Cant. v. 3. In the Song, this saying is an excuse for not rising from bed. S. Cyril applies it in a different way.}?

But there is perhaps among you some hypocrite, a man-pleaser, and one who makes a pretence of piety, but believes not from the heart; having the hypocrisy of Simon Magus; one who has come hither not in order to receive of the grace, but to spy out what is given: let him also learn from John: \emph{And now also the axe is laid unto the root of the trees, Every tree therefore that bringeth not forth good fruit is hewn down, and cast into the fire}\footnote{Matt. iii. 10.}. The Judge is inexorable; put away thine hypocrisy.

8. What then must you do? And what are the fruits of repentance? \emph{Let him that hath two coats give to him that hath none}\footnote{Luke iii. 11.}: the teacher was worthy of credit, since he was also the first to practise what he taught: he was not ashamed to speak, for conscience hindered not his tongue: \emph{and he that hath meat, let him do likewise}. Wouldst thou enjoy the grace of the Holy Spirit, yet judgest the poor not worthy of bodily food? Seekest thou the great gifts, and impartest not of the small? Though thou be a publican, or a fornicator, have hope of salvation: \emph{the publicans and the harlots go into the kingdom of God before you}\footnote{Matt. xxi. 31.}. Paul also is witness, saying, \emph{Neither fornicators, nor adulterers, nor} the rest, \emph{shall inherit the kingdom of God. And such were some of you: but ye were washed, but ye were sanctified}\footnote{1 Cor. vi. 9, 10.}. He said not, \emph{such are some of you}, but \emph{such were some of you}. Sin committed in the state of ignorance is pardoned, but persistent wickedness is condemned.

9. Thou hast as the glory of Baptism the Son Himself, the Only-begotten of God. For why should I speak any more of man? John was great, but what is he to the Lord? His was a loud-sounding voice, but what in comparison with the Word? Very noble was the herald, but what in comparison with the King? Noble was he that baptized with water, but what to Him that baptizeth \emph{with the Holy Ghost and with fire}\footnote{Matt. iii. 11.}? The Saviour baptized the Apostles with the Holy Ghost and with fire, when \emph{suddenly there came a sound from heaven as of the rushing of a mighty wind, and it filled all the house where they were sitting. And there appeared unto them cloven tongues like as of fire: and it sat upon each one of them, and they were all filled with the Holy Ghost}\footnote{Acts ii. 2.}.

10. If any man receive not Baptism, he hath not salvation; except only Martyrs, who even without the water receive the kingdom. For when the Saviour, in redeeming the world by His Cross, was pierced in the side, He shed forth blood and water; that men, living in times of peace, might be baptized in water, and, in times of persecution, in their own blood. For martyrdom also the Saviour is wont to call a baptism, saying, \emph{Can ye drink the cup which I drink, and be baptized with the baptism that I am baptized with}\footnote{Mark x. 38.}? And the Martyrs confess, by \emph{being made a spectacle unto the world, and to Angels, and to men}\footnote{1 Cor. iv. 9.}; and thou wilt soon confess:—but it is not yet the time for thee to hear of this.

11. Jesus sanctified Baptism by being Himself baptized. If the Son of God was baptized, what godly man is he that despiseth Baptism? But He was baptized not that He might receive remission of sins, for He was sinless; but being sinless, He was baptized, that He might give to them that are baptized a divine and excellent grace. For \emph{since the children are partakers of flesh and blood, He also Himself likewise partook of the same}\footnote{Heb. ii. 14.}, that having been made partakers of His presence in the flesh we might be made partakers also of His Divine grace: thus Jesus was baptized, that thereby we again by our participation might receive both salvation and honour. According to Job, there was in the waters the dragon that \emph{draweth up Jordan into his mouth}\footnote{Job xl. 23.}. Since, therefore, it was necessary to break the heads of the dragon in pieces\footnote{Ps. lxxiv. 14.}, He went down and bound the strong one in the waters, that we might receive power to \emph{tread upon serpents and scorpions}\footnote{Luke x. 19.}. The beast was great and terrible. \emph{No fishing-vessel was able to carry one scale of his tail}\footnote{Job xl. 26, in the Sept. in place of xli. 7: Canst thou fill his skin with barbed irons, or his head with fish spears? (\textsc{A.V.} and \textsc{R.V.})}: destruction ran before him\footnote{Job xli. 13, Sept. but in R.V. xli. 22: And terror danceth before him.}, ravaging all that met him. The Life encountered him, that the mouth of Death might henceforth be stopped, and all we that are saved might say, \emph{O death, where is thy sting? O grave, where is thy victory}\footnote{1 Cor. xv. 55.}? The sting of death is drawn by Baptism.

12. For thou goest down into the water, bearing thy sins, but the invocation of grace\footnote{Compare III. 3, and see Index, “Baptism.”}, having sealed thy soul, suffereth thee not afterwards to be swallowed up by the terrible dragon. Having gone down dead in sins, thou comest up quickened in righteousness. For if thou hast been \emph{united with the likeness of the Saviour’s death}\footnote{Rom. vi. 5.}, thou shalt also be deemed worthy of His Resurrection. For as Jesus took upon Him the sins of the world, and died, that by putting sin to death He might rise again in righteousness; so thou by going down into the water, and being in a manner buried in the waters, as He was in the rock, art raised again \emph{walking in newness of life}\footnote{Rom. vi. 4. Instead of “might rise again” (Roe, Casaub. Mon.), the older Editions have “might raise thee up,” which is less appropriate in this part of the sentence.}.

13. Moreover, when thou hast been deemed worthy of the grace, He then giveth thee strength to wrestle against the adverse powers. For as after His Baptism He was tempted forty days (not that He was unable to gain the victory before, but because He wished to do all things in due order and succession), so thou likewise, though not daring before thy baptism to wrestle with the adversaries, yet after thou hast received the grace and art henceforth confident in \emph{the armour of righteousness}\footnote{2 Cor. vi. 7.}, must then do battle, and preach the Gospel, if thou wilt.

14. Jesus Christ was the Son of God, yet He preached not the Gospel before His Baptism. If the Master Himself followed the right time in due order, ought we, His servants, to venture out of order? \emph{From that time Jesus began to preach\footnote{Matt. iv. 17.}}, when \emph{the Holy Spirit had descended upon Him in a bodily shape, like a dove}\footnote{Luke iii. 22.}; not that Jesus might see Him first, for He knew Him even before He came in a bodily shape, but that John, who was baptizing Him, might behold Him. For \emph{I}, saith he, \emph{knew Him not: but He that sent me to baptize with water, He said unto me, Upon whomsoever thou shalt see the Spirit descending and abiding on Him, that is He}\footnote{John i. 33.}. If thou too hast unfeigned piety, the Holy Ghost cometh down on thee also, and a Father’s voice sounds over thee from on high—not, “\emph{This is My Son},” but, “This has now been made My son;” for the “\emph{is}” belongs to Him alone, because \emph{In the beginning was the Word, and the Word was with God, and the Word was God}\footnote{Ib. i. 1.}. To Him belongs the “\emph{is},” since He is always the Son of God: but to thee “has now been made:” since thou hast not the sonship by nature, but receivest it by adoption. He eternally “\emph{is};” but thou receivest the grace by advancement.

15. Make ready then the vessel of thy soul, that thou mayest become a son of God, and \emph{an heir of God, and joint-heir with Christ}\footnote{Rom. viii. 17.}; if, indeed, thou art preparing thyself that thou mayest receive; if thou art drawing nigh in faith that thou mayest be made faithful; if of set purpose thou art putting off the old man. For all things whatsoever thou hast done shall be forgiven thee, whether it be fornication, or adultery, or any other such form of licentiousness. What can be a greater sin than to crucify Christ? Yet even of this Baptism can purify. For so spake Peter to the three thousand who came to him, to those who had crucified the Lord, when they asked him, saying, \emph{Men and brethren, what shall we do}\footnote{Acts ii. 37.}? For the wound is great. Thou hast made us think of our fall, O Peter, by saying, \emph{Ye killed the Prince of Life}\footnote{Ib. iii. 15.}. What salve is there for so great a wound? What cleansing for such foulness? What is the salvation for such perdition? \emph{Repent}, saith he, \emph{and be baptized every one of you in the name of Jesus Christ our Lord, for the remission of sins, and ye shall receive the gift of the Holy Ghost}\footnote{Ib. ii. 58.}. O unspeakable loving-kindness of God! They have no hope of being saved, and yet they are thought worthy of the Holy Ghost. Thou seest the power of Baptism! If any of you has crucified the Christ by blasphemous words; if any of you in ignorance has denied Him before men; if any by wicked works has caused the doctrine to be blasphemed; let him repent and be of good hope, for the same grace is present even now.

16. \emph{Be of good courage, O Jerusalem; the Lord will take away all thine iniquities}\footnote{Zeph. iii. 14, 15.}. \emph{The Lord will wash away the filth of His sons and of His daughters by the Spirit of judgment, and by the Spirit of burning}\footnote{Is. iv. 4.}.\emph{ He will sprinkle clean water upon you, and ye shall be cleansed from all your sin}\footnote{Ezek. xxxvi. 25.}. Angels shall dance around you, and say, Who is this that cometh up in white array, leaning upon her beloved\footnote{Cant. viii. 4, Gr. \textgreek{ἀδελφιδόν}, “brother,” “kinsman.”}? For the soul that was formerly a slave has now adopted her Master Himself as her kinsman: and He accepting the unfeigned purpose will answer: \emph{Behold, thou art fair, my love; behold, thou art fair: thy teeth are like flocks of sheep new shorn,} (because of the confession of a good conscience: and further) \emph{which have all of them twins}\footnote{Ib. iv. 1, 2.}; because of the twofold grace, I mean that which is perfected of water and of the Spirit\footnote{The Fathers sometimes speak as if Baptism was primarily the Sacrament of remission of sins, and \emph{upon} that came the gift of the Spirit, which notwithstanding was but begun in Baptism and completed in Confirmation. Vid. Tertullian. \emph{de Bapt}. 7, 8, \emph{supr}. i. 5 \emph{fin}. Hence, as in the text, Baptism may be said to be made up of \emph{two} gifts, Water, which is Christ’s blood, and the Spirit. There is no real difference between this and the ordinary way of speaking on the subject;—Water, which \emph{conveys} both gifts, is considered as a \emph{type} of one especially,—\emph{conveys} both remission of sins through Christ’s blood and the grace of the Spirit, but is the \emph{type} of one, \emph{viz}. the blood of Christ, as the Oil in Confirmation is of the other. And again, remission of sins is a complete gift given at once, sanctification an increasing one. (\textsc{R.W.C.}) See Index, “Baptism.”}, or that which is announced by the Old and by the New Testament. And God grant that all of you when you have finished the course of the fast, may remember what I say, and bringing forth fruit in good works, may stand blameless beside the Spiritual Bridegroom, and obtain the remission of your sins from God; to whom with the Son and Holy Spirit be the glory for ever. Amen.

\xchapter{Lecture IV. On the Ten Points of Doctrine.}

On the Ten\footnote{The number “ten” is confirmed by Theodoret, who quotes the article on Christ’s “Birth of the Virgin” as from Cyril’s fourth Catechetical Lecture “On the ten Doctrines.” The \textsc{mss.} vary between “ten” and “eleven,” and differ also in the special titles and numeration of the separate Articles.} Points of Doctrine.

\textsc{Colossians ii. 8}.

Beware lest any man spoil you through philosophy and vain deceit, after the tradition of men, after the rudiments of the world, \&c.

1. \textsc{Vice} mimics virtue, and the tares strive to be thought wheat, growing like the wheat in appearance, but being detected by good judges from the taste. \emph{The devil also transfigures himself into an angel of light}\footnote{2 Cor. xi. 14.}; not that he may reascend to where he was, for having made \emph{his heart hard as an anvil}\footnote{Job xli. 24, Sept.; xli. 15: \textgreek{ἡ καρδία αὐτοῦ…ἕστηκεν ὥσπερ ἄκμων ἀνήλατος}. These statements concerning the Devil seem to be directed against Origen’s opinion (\emph{De Principiis} I. 2), that the Angels “who have been removed from their primal state of blessedness have not been removed irrecoverably.” The question is discussed, and the opinions of several Fathers quoted, by Huet, \emph{Origeniana}, II. c. 25.}, he has henceforth a will that cannot repent; but in order that he may envelope those who are living an Angelic life in a mist of blindness, and a pestilent condition of unbelief. Many wolves are going about \emph{in sheeps’ clothing}\footnote{Matt. vii. 15. The same text is applied to Heretics by Ignatius, \emph{Philadelph.} ii. and by Irenæus, L. I. c. i. § 2.}, their clothing being that of sheep, not so their claws and teeth: but clad in their soft skin, and deceiving the innocent by their appearance, they shed upon them from their fangs the destructive poison of ungodliness. We have need therefore of divine grace, and of a sober mind, and of eyes that see, lest from eating tares as wheat we suffer harm from ignorance, and lest from taking the wolf to be a sheep we become his prey, and from supposing the destroying Devil to be a beneficent Angel we be devoured: for, as the Scripture saith, \emph{he goeth about as a roaring lion, seeking whom he may devour}\footnote{1 Pet. v. 8.}. This is the cause of the Church’s admonitions, the cause of the present instructions, and of the lessons which are read.

2. For the method of godliness consists of these two things, pious doctrines, and virtuous practice: and neither are the doctrines acceptable to God apart from good works, nor does God accept the works which are not perfected with pious doctrines. For what profit is it, to know well the doctrines concerning God, and yet to be a vile fornicator? And again, what profit is it, to be nobly temperate, and an impious blasphemer? A most precious possession therefore is the knowledge of doctrines: also there is need of a wakeful soul, since there are many \emph{that make spoil through philosophy and vain deceit}\footnote{Col. ii. 8.}. The Greeks on the one hand draw men away by their smooth tongue, \emph{for honey droppeth from a harlot’s lips}\footnote{Prov. v. 3.}: whereas they of the Circumcision deceive those who come to them by means of the Divine Scriptures, which they miserably misinterpret though \emph{studying them from childhood to old age}\footnote{Is. xlvi. 3. Sept. \textgreek{παιδευόμενοι ἐκ παιδίου ἕως γήρως}.}, and growing old in ignorance. But the children of heretics, \emph{by their good words and smooth tongue, deceive the hearts of the innocent}\footnote{Rom. xvi. 17. Cyril has \textgreek{εὐγλωττίας} in place of \textgreek{εὐλογίας}.}, disguising with the name of Christ as it were with honey the poisoned arrows\footnote{Compare Ignatius, \emph{Trall}. vi.} of their impious doctrines: concerning all of whom together the Lord saith, Take heed lest any man mislead you\footnote{Matt. xxiv. 4.}. This is the reason for the teaching of the Creed and for expositions upon it.

3. But before delivering you over to the Creed\footnote{Compare Rom. vi. 17: “\emph{that form of teaching whereunto ye were delivered.}” The instruction of Catechumens in the Articles of the Faith was commonly called the “Traditio Symboli,” or “Delivery of the Creed.”}, I think it is well to make use at present of a short summary of necessary doctrines; that the multitude of things to be spoken, and the long interval of the days of all this holy Lent, may not cause forgetfulness in the mind of the more simple among you; but that, having strewn some seeds now in a summary way, we may not forget the same when afterwards more widely tilled. But let those here present whose habit of mind is mature, and who \emph{have their senses already exercised to discern good and evil}\footnote{Heb. v. 14.}, endure patiently to listen to things fitted rather for children, and to an introductory course, as it were, of milk: that at the same time both those who have need of the instruction may be benefited, and those who have the knowledge may rekindle the remembrance of things which they already know.

I. Of God.

4. First then let there be laid as a foundation in your soul the doctrine concerning God; that God is One, alone unbegotten, without beginning, change, or variation\footnote{Compare Hermas, \emph{Mandat}. I. Athan. \emph{Epist. de Decretis Nic. Syn.} xxii.: \textgreek{οὕτω καὶ τὸ ἄτρεπτον καὶ ἀναλλοίωτον αὐτὸν εἶναι σωθήσεται}. So Aristotle (\emph{Metaphys}. XI. c. iv. 13) describes the First Cause as \textgreek{ἀπαθὲς καὶ ἀναλλοίωτον}.}; neither begotten of another, nor having another to succeed Him in His life; who neither began to live in time, nor endeth ever: and that He is both good and just; that if ever thou hear a heretic say, that there is one God who is just, and another who is good\footnote{Irenæus, I. c. xxvii. says that Cerdo taught that the God of the Law and the Prophets was not the Father of our Lord Jesus Christ: for that He is known, but the other unknown, and the one is just, but the other good. Also III. c. 25, § 3: “Marcion himself, therefore, by dividing God into two, and calling the one good, and the other judicial, on both sides puts an end to Deity.” Compare Tertullian, \emph{c. Marcion}. I. 2, and 6; Origen, \emph{c. Cels}. iv. 54.}, thou mayest immediately remember, and discern the poisoned arrow of heresy. For some have impiously dared to divide the One God in their teaching: and some have said that one is the Creator and Lord of the soul, and another of the body\footnote{This tenet was held by the Manichæans and other heretics, and is traced back to the Apostolic age by Bishop Pearson (\emph{Exposition of the Creed}, Art. i. p. 79, note c). Compare Athanasius \emph{c.} \emph{Apollinarium}, I. 21; II. 8; \emph{c.} \emph{Gentes}, § 6; \emph{de Incarnatione}, § 2, in this series, and Augustine (\emph{c.} \emph{Faustum}, xx. 15, 21, and xxi. 4).}; a doctrine at once absurd and impious. For how can a man become the one servant of two masters, when our Lord says in the Gospels, \emph{No man can serve two masters}\footnote{Matt. vi. 24; Luke xvi. 13.}? There is then One Only God, the Maker both of souls and bodies: One the Creator of heaven and earth, the Maker of Angels and Archangels: of many the Creator, but of One only the Father before all ages,—of One only, His Only-begotten Son, our Lord Jesus Christ, by Whom He made \emph{all things visible and invisible}\footnote{John i. 3; Col. i. 16.}.

5. This Father of our Lord Jesus Christ is not circumscribed in any place\footnote{S. Aug. \emph{in Ps}. lxxv. 6: Si in aliquo loco esset, non esset Deus. \emph{Sermo} 342: Deus habitando continet non continetur. Origen, \emph{c.} \emph{Cels}. vii. 34: “God is of too excellent a nature for any place: He holds all things in His power, and is Himself not confined by anything whatever.” Compare the quotation from Sir Isaac Newton’s \emph{Principia}, in the note on Cat. vi. 8.}, nor is He less than the heaven; \emph{but the heavens are the works of His fingers}\footnote{Ps. viii. 3.}, and \emph{the whole earth is held in His grasp}\footnote{Is. xl. 12.}: He is in all things and around all. Think not that the sun is brighter than He\footnote{See Cat. xv. 3, and note there.}, or equal to Him: for He who at first formed the sun must needs be incomparably greater and brighter. He foreknoweth the things that shall be, and is mightier than all, knowing all things and doing as He will; not being subject to any necessary sequence of events, nor to nativity, nor chance, nor fate; in all things perfect, and equally possessing every absolute form\footnote{\textgreek{ἰδέαν}. Cyril uses the word in the Platonic sense, as in the next sentence he adopts the formula, which Plato commonly uses in describing the “idea:” \textgreek{ἀεὶ κατὰ τὰ αὐτὰ καὶ ὡσαύτως ἔχειν}. Phaed. 78 c.} of virtue, neither diminishing nor increasing, but in mode and conditions ever the same; who hath prepared punishment for sinners, and a crown for the righteous.

6. Seeing then that many have gone astray in divers ways from the One God, some having deified the sun, that when the sun sets they may abide in the night season without God; others the moon, to have no God by day\footnote{Job xxxi. 26, 27. The worship of Sun and Moon under various names was almost universal.}; others the other parts of the world\footnote{Gaea or Tellus, the earth; Zeus or Jupiter, the sky; rivers, fountains, \&c.}; others the arts\footnote{Music, Medicine, Hunting, War, Agriculture, Metallurgy, \&c., represented by Apollo, Æsculapius, Diana, Mars, Ceres, Vulcan.}; others their various kinds of food\footnote{Herodotus, Book II., describes the Egyptian worship of various birds, fishes, and quadrupeds. Leeks and onions also were held sacred: Porrum et caepe nefas violare, Juv. \emph{Sat}. xv. 9. Compare Clement of Alexandria, \emph{Protrept}. c. ii. § 39, Klotz.}; others their pleasures\footnote{Eros, Dionysus.}; while some, mad after women, have set up on high an image of a naked woman, and called it Aphrodite\footnote{Clement of Alexandria (\emph{Protrept}. c. iv. § 53, Klotz) states that the courtesan Phryne was taken as a model for Aphrodite. “Praxiteles when fashioning the statue of Aphrodite of Cnidus made it like the form of Cratine his paramour.” \emph{Ibid}.}, and worshipped their own lust in a visible form; and others dazzled by the brightness of gold have deified it\footnote{Plutus.} and the other kinds of matter;—whereas if one lay as a first foundation in his heart the doctrine of the unity\footnote{\textgreek{τῆς μοναρχίας τοῦ θεοῦ}. See note on the title of Cat. VI. Praxeas made use of the term “Monarchy” to exclude the Son (and the Spirit) from the Godhead. Tertullian in his treatise against Praxeas maintains the true doctrine that the Son is no obstacle to the “Monarchy,” because He is of the substance of the Father, does nothing without the Father’s will, and has received all power from the Father, to Whom He will in the end deliver up the kingdom. In this sense Dionysius, Bishop of Rome, speaks of the Divine Monarchy as “that most sacred doctrine of the Church of God.” Compare Athanas. \emph{de Decretis, Nic. Syn.} c. vi. § 3 and Dr. Newman’s note. In \emph{Orat.} iv. \emph{c.} \emph{Arian}. p 606 (617), Athanasius derives the term from \textgreek{ἀρχή}, in the sense of “beginning:” \textgreek{οὕτως μία ἀρχὴ θεότητος καὶ οὐ δύο ἀρχαί, ὅθεν κυρίως καὶ μοναρχία ἐστίν}. See the full discussion of Monarchianism in \emph{Athanasius}, p. xxiii. ff. in this series, and Newman’s Introduction to \emph{Athan. Or}. iv.} of God, and trust to Him, he roots out at once the whole crop\footnote{For \textgreek{φοράν} (Bened.) many \textsc{mss.} read \textgreek{φθοράν}, “corruption.”} of the evils of idolatry, and of the error of the heretics: lay thou, therefore, this first doctrine of religion as a foundation in thy soul by faith.

Of Christ.

7. Believe also in the Son of God, One and Only, our Lord Jesus Christ, Who was begotten God of God, begotten Life of Life, begotten Light of Light\footnote{Compare xi. 4, 9, 18.}, Who is in all things like\footnote{\textgreek{Τὸν ὅμοιον κατὰ πάντα τῷ γεννησαντι}. On the meaning and history of this phrase, proposed by the Semi-Arians at the Council of Ariminum as a substitute for \textgreek{ὁμοούσιον}, see Athan. \emph{de Syn}. § 8, \emph{sqq.}} to Him that begat, Who received not His being in time, but was before all ages eternally and incomprehensibly begotten of the Father: The Wisdom and the Power of God, and His Righteousness personally subsisting\footnote{\textgreek{ἐνυπόστατος}. Cf. xi. 10; Athan. \emph{c. Apollinar}. I. 20, 21.}: Who sitteth on the right hand of the Father before all ages.

For the throne at God’s right hand He received not, as some have thought, because of His patient endurance, being crowned as it were by God after His Passion; but throughout His being,—a being by eternal generation\footnote{The \textsc{mss.} vary much, but I have followed the Benedictine text.},—He holds His royal dignity, and shares the Father’s seat, being God and Wisdom and Power, as hath been said; reigning together with the Father, and creating all things for the Father, yet lacking nothing in the dignity of Godhead, and knowing Him that hath begotten Him, even as He is known of Him that hath begotten; and to speak briefly, remember thou what is written in the Gospels, that \emph{none knoweth the Son but the Father, neither knoweth any the Father save the Son}\footnote{Matt. xi. 27; John x. 15; xvii. 25.}.

8. Further, do thou neither separate\footnote{This was a point earnestly maintained by the orthodox Bishops at Nicæa, that the Son begotten of the substance of the Father is ever inseparably in the Father. Athan. \emph{de Decretis Syn.} c. 20 ; Tertullian \emph{c. Marc}. IV. c. 6. Cf. Ignat. \emph{ad Trall}. vi. (Long Recension): \textgreek{τὸν μὲν γὰρ Χριστὸν ἀλλοτριουσι τοῦ Πατρός}.} the Son from the Father, nor by making a confusion believe in a Son-Fatherhood\footnote{\textgreek{υἱοπατορία}. A term of derision applied to the doctrine of Sabellius. Compare Athanas. \emph{Expositio Fidei}, c. 2: “neither do we imagine a Son-Father, as the Sabellians.” See Index, \textgreek{Υιοπάτωρ}.}; but believe that of One God there is One Only-begotten Son, who is before all ages God the Word; not the uttered\footnote{\textgreek{Λόγος προφορικός}, the term used by Paul of Samosata, implied that the Word was impersonal, being conceived as a particular activity of God. See Dorner, \emph{Person of Christ}, Div. I. vol. ii. p. 436 (English Tr.): and compare Athanasius, \emph{Expositio Fidei,} c. 1; \textgreek{υἱὸν ἐκ τοῦ Πατρὸς ἀνάρχως καὶ ἀϊδίως γεγεννημένον, λόγον δὲ οὐ προφορικόν, οὐκ ἐνδιάθετον}. Cardinal Newman (\emph{Athan. c. Arianos}, I. 7, note) observes that some Christian writers of the 2nd Century “seem to speak of the Divine generation as taking place immediately before the creation of the world, that is, as if not eternal, though at the same time they teach that our Lord existed before that generation. In other words they seem to teach that He was the Word from eternity, and became the Son at the beginning of all things; some of them expressly considering Him, first as the \textgreek{λόγος ἐνδιάθετος}, or Reason, in the Father, or (as may be speciously represented) a mere attribute; next, as the \textgreek{λόγος προφορικός}, or Word.” \par The terms \textgreek{λόγος ἐνδιάθετος}, or ‘word conceived in the mind,’ and \textgreek{λόγος προφορικός}, or ‘word expressed’ (\emph{emissum}, or \emph{prolalivum}), were in use among the Gnostics (\emph{Iren}. II. c. 12, § 5). As applied to the Son both terms, though sometimes used in a right sense, were condemned as inadequate. Compare xi. 10.} word diffused into the air, nor to be likened to impersonal words\footnote{\textgreek{ἀνυποστάτοις λόγοις}. Athan. \emph{c. Arianos Orat}. iv. c. 8: \textgreek{πάλιν οἱ λέγοντες μόνον ὄνομα εἶναι υἱοῦ, ἀνούσιον δὲ καὶ ἀνυπόστατον εἶναι τὸν υἱὸν τοῦ Θεοῦ, κ.τ.λ.}}; but the Word the Son, Maker of all who partake of reason, the Word who heareth the Father, and Himself speaketh. And on these points, should God permit, we will speak more at large in due season; for we do not forget our present purpose to give a summary introduction to the Faith.

Concerning His Birth of the Virgin.

9. Believe then that this Only-begotten Son of God for our sins came down from heaven upon earth, and took upon Him this human nature of like passions\footnote{\textgreek{ὁμοιοπαθῆ}. Compare Acts xiv. 15; Jas. v. 17.} with us, and was begotten of the Holy Virgin and of the Holy Ghost, and was made Man, not in seeming and mere show\footnote{On the origin of the Docetic heresy, see vi. 14.}, but in truth; nor yet by passing through the Virgin as through a channel\footnote{Valentinus the Gnostic taught that God produced a Son of an animal nature who “passed through Mary just as water through a tube, and that on him the Saviour descended at his Baptism.” Irenæus, I. vii. 2.}; but was of her made truly flesh, [and truly nourished with milk\footnote{The words which the Benedictine Editor introduces in the brackets are found in Theodoret, and adopted by recent Editors, with Codd. M.A.}], and did truly eat as we do, and truly drink as we do. For if the Incarnation was a phantom, salvation is a phantom also. The Christ was of two natures, Man in what was seen, but God in what was not seen; as Man truly eating like us, for He had the like feeling of the flesh with us; but as God feeding the five thousand from five loaves; as Man truly dying, but as God raising him that had been dead four days; truly sleeping in the ship as Man, and walking upon the waters as God.

Of the Cross.

10. He was truly crucified for our sins. For if thou wouldest deny it, the place refutes thee visibly, this blessed Golgotha\footnote{Eusebius, \emph{Life of Constantine}, iii. 28.}, in which we are now assembled for the sake of Him who was here crucified; and the whole world has since been filled with pieces of the wood of the Cross\footnote{The discovery of the “True Cross” is related with many marvellous particulars by Socrates, \emph{Eccles. Hist.} i. 17; and Sozomen, \emph{E. H.} ii. 1. A portion was said to have been left by Helena at Jerusalem, enclosed in a silver case; and another portion sent to Constantinople, where Constantine privately enclosed it in his own statue, to be a safeguard to the city. Eusebius, \emph{Life of Constantine}, iii. 25–30 , gives a long account of the discovery of the Holy Sepulchre, but makes no mention of the Cross. Cyril seems to have been the first to record it, 25 years after. Cf. Greg. Nyss. \emph{Bapt. Christi} (p. 519).}. But He was crucified not for sins of His own, but that we might be delivered from \emph{our} sins. And though as Man He was at that time \emph{despised of men}, and was buffeted, yet He was acknowledged by the Creation as God: for when the sun saw his Lord dishonoured, he grew dim and trembled, not enduring the sight.

Of His Burial.

11. He was truly laid as Man in a tomb of rock; but rocks were rent asunder by terror because of Him. He went down into the regions beneath the earth, that thence also He might redeem the righteous\footnote{Compare xiv. 18, 19, on the Descent into Hades.}. For, tell me, couldst thou wish the living only to enjoy His grace, and that, though most of them are unholy; and not wish those who from Adam had for a long while been imprisoned to have now gained their liberty? Esaias the Prophet proclaimed with loud voice so many things concerning Him; wouldst thou not wish that the King should go down and redeem His herald? David was there, and Samuel, and all the Prophets\footnote{The same Old Testament saints are named in xiv. 19, as redeemed by Christ in Hades.}, John himself also, who by his messengers said, \emph{Art thou He that should come, or look we for another}\footnote{Matt. xi. 3.}? Wouldst thou not wish that He should descend and redeem such as these?

Of the Resurrection.

12. But He who descended into the regions beneath the earth came up again; and Jesus, who was buried, truly rose again the third day. And if the Jews ever worry thee, meet them at once by asking thus: Did Jonah come forth from the whale on the third day, and hath not Christ then risen from the earth on the third day? Is a dead man raised to life on touching the bones of Elisha, and is it not much easier for the Maker of mankind to be raised by the power of the Father? Well then, He truly rose, and after He had risen was seen again of the disciples: and twelve disciples were witnesses of His Resurrection, who bare witness not in pleasing words, but contended even unto torture and death for the truth of the Resurrection. What then, \emph{shall every word be established at the mouth of two of three witnesses}\footnote{Deut. xix. 15.}, according to the Scripture, and, though twelve bear witness to the Resurrection of Christ, art thou still incredulous in regard to His Resurrection?

Concerning the Ascension.

13. But when Jesus had finished His course of patient endurance, and had redeemed mankind from their sins, He ascended again into the heavens, a cloud receiving Him up: and as He went up Angels were beside Him, and Apostles were beholding. But if any man disbelieves the words which I speak, let him believe the actual power of the things now seen. All kings when they die have their power extinguished with their life: but Christ crucified is worshipped by the whole world. We proclaim The Crucified, and the devils tremble now. Many have been crucified at various times; but of what other who was crucified did the invocation ever drive the devils away?

14. Let us, therefore, not be ashamed of the Cross of Christ; but though another hide it, do thou openly seal it upon thy forehead, that the devils may behold the royal sign and flee trembling far away\footnote{Justin M. \emph{Dialogue with Trypho}, 247 C: We call Him Helper and Redeemer, the power of whose Name even demons do fear; and at this day, when exorcised in the name of Jesus Christ, crucified under Pontius Pilate, Governor of Judæa, they are overcome.}. Make then this sign at eating and drinking, at sitting, at lying down, at rising up, at speaking, at walking: in a word, at every act\footnote{Tertullian, \emph{de Coronâ}, 3: At every forward step and movement, at every going in and out, when we put on our clothes and shoes, when we bathe, when we sit at table, when we light the lamps, on couch, on seat, in all the ordinary actions of daily life, we trace upon the forehead the Sign. If for these, and other such rules, you insist upon having positive Scripture injunction, you will find none. Tradition will be held forth to you as the originator of them, custom as their strengthener, and faith as their observer.}. For He who was here crucified is in heaven above. If after being crucified and buried He had remained in the tomb, we should have had cause to be ashamed; but, in fact, He who was crucified on Golgotha here, has ascended into heaven from the Mount of Olives on the East. For after having gone down hence into Hades, and come up again to us, He ascended again from us into heaven, His Father addressing Him, and saying, \emph{Sit Thou on My right hand, until I make Thine enemies Thy footstool}\footnote{Ps. cx. 1.}.

Of Judgment to Come.

15. This Jesus Christ who is gone up shall come again, not from earth but from heaven: and I say, “not from earth,” because there are many Antichrists to come at this time from earth. For already, as thou hast seen, many have begun to say, \emph{I am the Christ}\footnote{Matt. xxiv. 5.}: and \emph{the abomination of desolation}\footnote{Matt. xxiv. 15. Compare Cat. xv. 9, 15.} is yet to come, assuming to himself the false title of Christ. But look thou for the true Christ, the Only-begotten Son of God, coming henceforth no more from earth, but from heaven, appearing to all more bright than any lightning and brilliancy of light, with angel guards attended, that He may judge both quick and dead, and reign in a heavenly, eternal kingdom, which shall have no end. For on this point also, I pray thee, make thyself sure, since there are many who say that Christ’s Kingdom hath an end\footnote{Compare xv. 27, where the followers of Marcellus of Ancyra are indicated as holding this opinion.}.

Of the Holy Ghost.

16. Believe thou also in the Holy Ghost, and hold the same opinion concerning Him, which thou hast \emph{received to hold} concerning the Father and the Son, and follow not those who teach blasphemous things of Him\footnote{In xvi. 6–10, Cyril gives a long list of heresies concerning the Holy Ghost.}. But learn thou that this Holy Spirit is One, indivisible, of manifold power; having many operations, yet not Himself divided; Who knoweth the mysteries, Who \emph{searcheth all things, even the deep things of God}\footnote{1 Cor. ii. 10.}: Who descended upon the Lord Jesus Christ in form of a dove; Who wrought in the Law and in the Prophets; Who now also at the season of Baptism sealeth thy soul; of Whose holiness also every intellectual nature hath need: against Whom \emph{if any dare to blaspheme, he hath no forgiveness, neither in this world, nor in that which is to come}\footnote{Matt. xii. 32.}: “Who with the Father and the Son together\footnote{This clause is not in the Creed of Nicæa, but is added in the Creed of Constantinople, \textsc{a.d.} 381.}” is honoured with the glory of the Godhead: of Whom also \emph{thrones, and dominions, principalities, and powers} have need\footnote{Col. i. 16.}. For there is One God, the Father of Christ; and One Lord Jesus Christ, the Only-begotten Son of the Only God; and One Holy Ghost, the sanctifier and deifier of all\footnote{\textgreek{θεοποιόν} is omitted in Codd. Roe, Casaubon, and A.}, Who spake in the Law and in the Prophets, in the Old and in the New Testament.

17. Have thou ever in thy mind this seal\footnote{The Benedictine Editor argues from Cat. i. 5, “that thou mayest by faith seal up the things that are spoken;” and xxiii. 18: “sealing up the Prayer by the Amen,” that Cyril means by “this seal” the firm belief of Christian doctrine. Compare John iii. 33. But Milles understands by the “seal” the Creed itself, which agrees better with the following context.}, which for the present has been lightly touched in my discourse, by way of summary, but shall be stated, should the Lord permit, to the best of my power with the proof from the Scriptures. For concerning the divine and holy mysteries of the Faith, not even a casual statement must be delivered without the Holy Scriptures; nor must we be drawn aside by mere plausibility and artifices of speech. Even to me, who tell thee these things, give not absolute credence, unless thou receive the proof of the things which I announce from the Divine Scriptures. For this salvation which we believe depends not on ingenious reasoning\footnote{\textgreek{ἡ σωτηρία γὰρ αὕτη τῆς πίστεως ἡμῶν}, which might be rendered, “this our salvation by faith,” or, with Milles, “this safety of our Faith.” For the rendering in the text compare Heb. iii. 1: \textgreek{ἀρχιερέα τῆς ὁμολογίας ἡμῶν}. On \textgreek{εὑρεσιλογία}, see Polybius xviii. 29, § 3: \textgreek{διὰ τῆς προς ἀλλήλους εὑρεσιλογίας}.}, but on demonstration of the Holy Scriptures.

Of the Soul.

18. Next to the knowledge of this venerable and glorious and all-holy Faith, learn further what thou thyself art: that as man thou art of a two-fold nature, consisting of soul and body; and that, as was said a short time ago, the same God is the Creator both of soul and body\footnote{iv. 4.}. Know also that thou hast a soul self-governed, the noblest work of God, made after the image of its Creator\footnote{In the Clementine Homily xvi. 16, the soul having come forth from God, clothed with His breath, is said to be of the same substance, and yet not God. In Tertull. \emph{c. Marcion} II. c. 9, the soul is the \emph{affatus} (\textgreek{πνοή} not \textgreek{πνεῦμα}) of God, i.e. the image of the Spirit, and inferior to it, though possessing the true lineaments of divinity, immortality, freedom, its own mastery over itself.}: immortal because of God that gives it immortality; a living being, rational, imperishable, because of Him that bestowed these gifts: having free power to do what it willeth\footnote{Tertull. \emph{c. Marc}. II. 6: It was proper that he who is the image and likeness of God should be formed with a free will, and a mastery of himself, so that this very thing, namely freedom of will and self-command, might be reckoned as the image and likeness of God in him.}. For it is not according to thy nativity that thou sinnest, nor is it by the power of chance that thou committest fornication, nor, as some idly talk, do the conjunctions of the stars compel thee to give thyself to wantonness\footnote{Compare Aug. \emph{de Civ. Dei}. v. 1, where he says that the astrologers (Mathematici) say, not merely such or such a position of Mars signifies that a man will be a murderer, but makes him a murderer. See Dict. of Christian Antiq., “Astrology.”}. Why dost thou shrink from confessing thine own evil deeds, and ascribe the blame to the innocent stars? Give no more heed, pray, to astrologers; for of these the divine Scripture saith, \emph{Let the stargazers of the heaven stand up and save thee}, and what follows: \emph{Behold, they all shall be consumed as stubble on the fire, and shall not deliver their soul from the flame}\footnote{Is. xlvii. 13.}.

19. And learn this also, that the soul, before it came into this world, had committed no sin\footnote{“The Orphic poets were under the impression that the soul is suffering the punishment of sin, and that the body is an enclosure or prison in which the soul is incarcerated and kept (\textgreek{σώζεται}) as the name \textgreek{σῶμα} implies, until the penalty is paid.” Plato, \emph{Cratyl}. 400. Clement of Alexandria (\emph{Strom}. III. iii. 17), after referring to this passage of Plato, quotes Philolaus the Pythagorean, as saying: “The ancient theologians and soothsayers also testify that the soul has been chained to the body for a kind of punishment, and is buried in it as in a tomb.“}, but having come in sinless, we now sin of our free-will. Listen not, I pray thee, to any one perversely interpreting the words, \emph{But if I do that which I would not}\footnote{Rom. vii. 16.}: but remember Him who saith, \emph{If ye be willing, and hearken unto Me, ye shall eat the good things of the land: but if ye be not willing, neither hearken unto Me, the sword shall devour you, \&c.}\footnote{Is. i. 19, 20.}: and again, \emph{As ye presented your members as servants to uncleanness and to iniquity unto iniquity, even so now present your members as servants to righteousness unto sanctification}\footnote{Rom. vi. 19.}. Remember also the Scripture, which saith, \emph{Even as they did not like to retain God in their knowledge}\footnote{Rom. i. 28.}: and, \emph{That which may be known of God is mani}\emph{festin them}\footnote{Rom. i. 19.}; and again, \emph{their eyes they have closed}\footnote{Matt. xiii. 15.}. Also remember how God again accuseth them, and saith, \emph{Yet I planted thee a fruitful vine, wholly true: how art thou turned to bitterness, thou the strange vine}\footnote{Jer. ii. 21.}?

20. The soul is immortal, and all souls are alike both of men and women; for only the members of the body are distinguished\footnote{Apelles, the heretic, attributed the difference of sex to the soul, which existing before the body impressed its sex upon it. Tertull. \emph{On the Soul}, c. xxxvi.}. There is not a class of souls sinning by nature, and a class of souls practising righteousness by nature\footnote{Irenæus I. vii. 5: “They (the Valentinians) conceive of three kinds of men, spiritual, material, and animal.…These three natures are no longer found in one person, but constitute various kinds of men.…And again subdividing the animal souls themselves, they say that some are by nature good, and others by nature evil.” Origen \emph{on Romans}, Lib. VIII. § 10: “I know not how those who come from the School of Valentinus and Basilides…suppose that there are souls of one nature which are always safe and never perish, and others which always perish, and are never saved.”}: but both act from choice, the substance of their souls being of one kind only, and alike in all. I know, however, that I am talking much, and that the time is already long: but what is more precious than salvation? Art thou not willing to take trouble in getting provisions for the way against the heretics? And wilt thou not learn the bye-paths of the road, lest from ignorance thou fall down a precipice? If thy teachers think it no small gain for thee to learn these things, shouldest not thou the learner gladly receive the multitude of things told thee?

21. The soul is self-governed: and though the devil can suggest, he has not the power to compel against the will. He pictures to thee the thought of fornication: if thou wilt, thou acceptest it; if thou wilt not, thou rejectest. For if thou wert a fornicator by necessity, then for what cause did God prepare hell? If thou were a doer of righteousness by nature and not by will, wherefore did God prepare crowns of ineffable glory? The sheep is gentle, but never was it crowned for its gentleness: since its gentle quality belongs to it not from choice but by nature.

Of the Body.

22. Thou hast learned, beloved, the nature of the soul, as far as there is time at present: now do thy best to receive the doctrine of the body also. Suffer none of those who say that this body is no work of God\footnote{See iv. 18.}: for they who believe that the body is independent of God, and that the soul dwells in it as in a strange vessel, readily abuse it to fornication\footnote{On the impure practices of the Manichees, see vi. 33, 34.}. And yet what fault have they found in this wonderful body? For what is lacking in comeliness? And what in its structure is not full of skill? Ought they not to have observed the luminous construction of the eyes? And how the ears being set obliquely receive the sound unhindered? And how the smell is able to distinguish scents, and to perceive exhalations? And how the tongue ministers to two purposes, the sense of taste, and the power of speech? How the lungs placed out of sight are unceasing in their respiration of the air? Who imparted the incessant pulsation of the heart? Who made the distribution into so many veins and arteries? Who skilfully knitted together the bones with the sinews? Who assigned a part of the food to our substance, and separated a part for decent secretion, and hid away the unseemly members in more seemly places? Who when the human race must have died out, rendered it by a simple intercourse perpetual?

23. Tell me not that the body is a cause of sin\footnote{Fortunatus, the Manichee, in August. \emph{Disput}. ii. 20, \emph{contra Fortunat}. is represented as saying, What we assert is this, that the soul is compelled to sin by a substance of contrary nature.}. For if the body is a cause of sin, why does not a dead body sin? Put a sword in the right hand of one just dead, and no murder takes place. Let beauties of every kind pass before a youth just dead, and no impure desire arises. Why? Because the body sins not of itself, but the soul through the body. The body is an instrument, and, as it were, a garment and robe of the soul: and if by this latter it be given over to fornication, it becomes defiled: but if it dwell with a holy soul, it becomes a temple of the Holy Ghost. It is not I that say this, but the Apostle Paul hath said, \emph{Know ye not, that your bodies are the temple of the Holy Ghost which is in you}\footnote{1 Cor. vi. 19.}? Be tender, therefore, of thy body as being a temple of the Holy Ghost. Pollute not thy flesh in fornication: defile not this thy fairest robe: and if ever thou hast defiled it, now cleanse it by repentance: get thyself washed, while time permits.

24. And to the doctrine of chastity let the first to give heed be the order of Solitaries\footnote{\textgreek{μονάζοντες}. Compare xii. 33; xvi. 22. The origin of Monasticism is usually traced to the time of the Decian persecution, the middle of the third century. Previously “there were no monks, but only ascetics in the Church; from that time to the reign of Constantine, Monachism was confined to the anchorets living in private cells in the wilderness: but when Pachomius had erected monasteries in Egypt, other countries presently followed the example.…Hilarion, who was scholar to Antonius, was the first monk that ever lived in Palestine or Syria.” Bingham, VII. i. 4.} and of Virgins, who maintain the angelic life in the world; and let the rest of the Church’s people follow them. For you, brethren, a great crown is laid up: barter not away a great dignity for a petty pleasure: listen to the Apostle speaking: \emph{Lest there be any fornicator or profane person, as Esau, who for one mess of} \emph{meat sold his own birthright}\footnote{Heb. xii. 16.}. Enrolled henceforth in the Angelic books for thy profession of chastity, see that thou be not blotted out again for thy practice of fornication.

25. Nor again, on the other hand, in maintaining thy chastity be thou puffed up against those who walk in the humbler path of matrimony. For as the Apostle saith, \emph{Let marriage be had in honour among all, and let the bed be undefiled}\footnote{Heb. xiii. 4.}. Thou too who retainest thy chastity, wast thou not begotten of those who had married? Because thou hast a possession of gold, do not on that account reprobate the silver. But let those also be of good cheer, who being married use marriage lawfully; who make a marriage according to God’s ordinance, and not of wantonness for the sake of unbounded license; who recognise seasons of abstinence, \emph{that they may give themselves unto prayer}\footnote{1 Cor. vii. 5.}; who in our assemblies bring clean bodies as well as clean garments into the Church; who have entered upon matrimony for the procreation of children, but not for indulgence.

26. Let those also who marry but once not reprobate those who have consented to a second marriage\footnote{The condemnation of a second marriage, which the Benedictine Editor and others import into this passage, is not to be found in it. \textgreek{τοὺς δευτέρῳ γάμῳ συμπεριενεχθέντας} neither means “qui ad secundas nuptias ultro se dejecere,” nor even “who have \emph{involved} themselves” (R.W.C.), but simply “who have consented to,”—or, “consented together in—a second marriage,” without any intimation of censure. See V. 9; VI. 13: Ecclus. xxv. 1; \textgreek{γυνὴ καὶ ἁνὴρ ἑαυτοῖς συμπεριφερόμενοι}; 2 Macc. ix. 27; Euseb. \emph{H. E.} ix. 9, 7: \textgreek{ἀνεξικάκως καὶ συμμέτρως συμπεριφέροιντο αὐτοῖς}; Zeno, \emph{ap. Diog. Laert}. vii. 18; \textgreek{τὸ συμπεριφερεσθαι τοῖς φίλοις}. \emph{Diog. Laert}. vii. 13: \textgreek{εὐσυμπερίφορος}. Polyb. IV. 35, § 7, and II. 17, § 12. The gentleness with which Cyril here speaks of second marriages is in striking contrast with the passionate vehemence of Tertullian in the treatise \emph{de Monogamia}, and elsewhere. Aug. \emph{de Hæresibus}, cc. 26, 38, reckons the condemnation of second marriage among the heretical doctrines of the Montanists and Cathari. In the treatise \emph{de Bono Viduitatis}, c. 6, he argues that a second marriage is not to be condemned, but is less honourable than widowhood, and severely rebukes the heretical teaching on this point of Tertullian, the Montanists, and the Novatians. \emph{De Bono Conjugali}, c. 21: Sacramentum nuptiarum temporis nostri sic ad unum virum et unam uxorem redactum est, ut Ecclesiæ dispensatorem non liceat ordinare nisi unius uxoris virum. On the practice of the Church at various times see Bingham, IV. v. 1–4; Suicer, \emph{Thesaur}. \textgreek{Διγαμία}.}: for though continence is a noble and admirable thing, yet it is also permissible to enter upon a second marriage, that the weak may not fall into fornication. For \emph{it is good for them}, saith the Apostle, \emph{if they abide even as I. But if they have not continency, let them marry: for it is better to marry than to burn}\footnote{1 Cor. vii. 8, 9.}. But let all the other practices be banished afar, fornication, adultery, and every kind of licentiousness: and let the body be kept pure for the Lord, that the Lord also may have respect unto the body. And let the body be nourished with food, that it may live, and serve without hindrance; not, however, that it may be given up to luxuries.

Concerning Meats.

27. And concerning food let these be your ordinances, since in regard to meats also many stumble. For some deal indifferently with things offered to idols\footnote{The Nicolaitans (\emph{Apocal}. ii. 14, 20); and the Valentinians, of whom Irenæus (II. xiv. 5), says that they derived their opinion as to the indifference of meats from the Cynics. See also Irenæus I. vi. 3; and xxvi. 3.}, while others discipline themselves, but condemn those that eat: and in different ways men’s souls are defiled in the matter of meats, from ignorance of the useful reasons for eating and not eating. For we fast by abstaining from wine and flesh, not because we abhor them as abominations, but because we look for our reward; that having scorned things sensible, we may enjoy a spiritual and intellectual feast; and that \emph{having now sown in tears we may reap in joy}\footnote{Ps. cxxvi. 5.} in the world to come. Despise not therefore them that eat, and because of the weakness of their bodies partake of food: nor yet blame these who \emph{use a little wine for their stomach’s sake and their often infirmities}\footnote{1 Tim. v. 23.}: and neither condemn the men as sinners, nor abhor the flesh as strange food; for the Apostle knows some of this sort, when he says: \emph{forbidding to marry, and commanding to abstain from meats, which God created to be received with thanksgiving by them that believe}\footnote{1 Tim. iv. 3.}. In abstaining then from these things, abstain not as from things abominable\footnote{The various sects of Gnostics, and the Manichees, considered certain meats and drinks, as flesh and wine, to be polluting. Vid. Iren. \emph{Hær}. i. 28. Clem. \emph{Pæd}. ii. 2. p. 186. Epiph. \emph{Hær}. xlvi. 2, xlvii. 1, \&c., \&c. August. \emph{Hær}. 46, vid. Canon. \emph{Apost}. 43. “If any Bishop, \&c., abstain from marriage, flesh, and wine, not for discipline (\textgreek{δι᾽ ἄσκησιν}) but as abhorring them, forgetting that they are all very good, \&c., and speaking blasphemy against the creation, let him amend or be deposed,” \&c. \textsc{R.W.C.}}, else thou hast no reward: but as being good things disregard them for the sake of the better spiritual things set before thee.

28. Guard thy soul safely, lest at any time thou eat of things offered to idols: for concerning meats of this kind, not only I at this time, but ere now Apostles also, and James the bishop of this Church, have had earnest care: and the Apostles and Elders write a Catholic epistle to all the Gentiles, that they should \emph{abstain} first \emph{from things offered to idols, and} then \emph{from blood} also \emph{and from things strangled}\footnote{Acts xv. 20, 29. The prohibition of blood and things strangled has continued to the present day in the Eastern Church, though already disregarded by the Latins in the time of S. Augustine (\emph{c. Faustum}. xxxii. 13).}. For many men being of savage nature, and living like dogs, both lap up blood\footnote{Tertullian (\emph{Apologeticus}, c. 9) speaks of those “who at the gladiator shows, for the cure of epilepsy, quaff with greedy thirst the blood of criminals slain in the arena,” and of others “who make meals on the flesh of wild beasts at the place of combat:” and contrasts the habits of Christians, who abstain from things strangled, to avoid pollution by the blood.}, in imitation of the manner of the fiercest beasts, and greedily devour things strangled. But do thou, the servant of Christ, in eating observe to eat with reverence. And so enough concerning meats.

Of Apparel.

29. But let thine apparel be plain, not for adornment, but for necessary covering: not to minister to thy vanity, but to keep thee warm in winter, and to hide the unseemliness of the body: lest under pretence of hiding the unseemliness, thou fall into another kind of unseemliness by thy extravagant dress.

Of the Resurrection.

30. Be tender, I beseech thee, of this body, and understand that thou wilt be raised from the dead, to be judged with this body. But if there steal into thy mind any thought of unbelief, as though the thing were impossible, judge of the things unseen by what happens to thyself. For tell me; a hundred years ago or more, think where wast thou thyself: and from what a most minute and mean substance thou art come to so great a stature, and so much dignity of beauty\footnote{XVIII. 9.}. What then? Cannot He who brought the non-existent into being, raise up again that which already exists and has decayed\footnote{Compare xviii. 6, 9; Athenagoras, \emph{On the Resurrection of the Dead}, c. 3.}? He who raises the corn, which is sown for our sakes, as year by year it dies,—will He have difficulty in raising us up, for whose sakes that corn also has been raised\footnote{XVIII. 6. John xii. 24; 1 Cor. xv. 36.}? Seest thou how the trees stand now for many months without either fruit or leaves: but when the winter is past they spring up whole into life again as if from the dead\footnote{XVIII. 7.}: shall not we much rather and more easily return to life? The rod of Moses was transformed by the will of God into the unfamiliar nature of a serpent: and cannot a man, who has fallen into death, be restored to himself again?

31. Heed not those who say that this body is not raised; for it is raised: and Esaias is witness, when he says: \emph{The dead shall arise, and they that are in the tombs shall awake}\footnote{Is. xxvi. 19.}: and according to Daniel, \emph{Many of them that sleep in the dust of the earth shall arise, some to everlasting life, and some to everlasting shame}\footnote{Dan. xii. 2.}. But though to rise again is common to all men, yet the resurrection is not alike to all: for the bodies received by us all are eternal, but not like bodies by all: for the just receive them, that through eternity they may join the Choirs of Angels; but the sinners, that they may endure for ever the torment of their sins.

Of the Laver.

32. For this cause the Lord, preventing us according to His loving-kindness, has granted repentance at Baptism\footnote{Gr. \textgreek{λουτροῦ μετάνοιαν}. Other readings are \textgreek{λύτρον μετανοίας}, “redemption by repentance,” and \textgreek{λουτρὸν μετανοίας} “a laver (baptism) of repentance.”}, in order that we may cast off the chief—nay rather the whole burden of our sins, and having received the seal by the Holy Ghost, may be made heirs of eternal life. But as we have spoken sufficiently concerning the Laver the day before yesterday, let us now return to the remaining subjects of our introductory teaching.

Of the Divine Scriptures.

33. Now these the divinely-inspired Scriptures of both the Old and the New Testament teach us. For the God of the two Testaments is One, Who in the Old Testament foretold the Christ Who appeared in the New; Who by the Law and the Prophets led us to Christ’s school.\emph{ For before faith came, we were kept in ward under the law}, and, \emph{the law hath been our tutor to bring us unto Christ}\footnote{Gal. iii. 24. The \textgreek{Παιδαγωγός} is described by Clement of Alexandria (\emph{Paedag}. i. 7) as one who both conducts a boy to school, and helps to teach him,—an usher: “under-master” (Wicliff).}. And if ever thou hear any of the heretics speaking evil of the Law or the Prophets, answer in the sound of the Saviour’s voice, saying, Jesus \emph{came not to destroy the Law, but to fulfil it}\footnote{Matt. v. 17.}. Learn also diligently, and from the Church, what are the books of the Old Testament, and what those of the New. And, pray, read none of the apocryphal writings\footnote{\textgreek{τῶν ἀποκρύφων}. The sense in which Cyril uses this term may be learned from Rufinus (\emph{Expositio Symboli}, § 38), who distinguishes three classes of books: (1) The Canonical Books of the Old and New Testaments, which alone are to be used in proof of doctrine; (2) Ecclesiastical, which may be read in Churches, including Wisdom, Ecclesiasticus, Tobit, Judith, and the Books of the Maccabees, in the Old Testament, and \emph{The Shepherd} of Hermas, and \emph{The Two Ways} in the New Testament; (3) The other writings they called “Apocryphal,” which they would not have read in Churches. The distinction is useful, though the second class is not complete.}: for why dost thou, who knowest not those which are acknowledged among all, trouble thyself in vain about those which are disputed? Read the Divine Scriptures, the twenty-two books of the Old Testament, these that have been translated by the Seventy-two Interpreters\footnote{The original source of this account of the Septuagint version is a letter purporting to have been written by Aristeas, or Aristæus, a confidential minister of Ptolemy Philadelphus, to his brother Philocrates. Though the letter is not regarded as genuine its statements are in part admitted to be true, being confirmed by a fragment, preserved by Eusebius (\emph{Præparatio Evangelica}, ix. 6.), of a work of Aristobulus, a Jewish philosopher who wrote in the reign of Ptolemy Philometor, 181–146, \textsc{b.c.} Upon these testimonies it is generally admitted that “the whole Law,” i.e. the Pentateuch was translated into Greek at Alexandria in the reign either of Ptolemy Soter (323–285, \textsc{b.c.}), or of his son Ptolemy Philadelphus (285–247, \textsc{b.c.}), under the direction of Demetrius Phalereus, curator of the King’s library.}.

34. For after the death of Alexander, the king of the Macedonians, and the division of his kingdom into four principalities, into Babylonia, and Macedonia, and Asia, and Egypt, one of those who reigned over Egypt, Ptolemy Philadelphus, being a king very fond of learning, while collecting the books that were in every place, heard from Demetrius Phalereus, the curator of his library, of the Divine Scriptures of the Law and the Prophets, and judged it much nobler, not to get the books from the possessors by force against their will, but rather to propitiate them by gifts and friendship; and knowing that what is extorted is often adulterated, being given unwillingly, while that which is willingly supplied is freely given with all sincerity, he sent to Eleazar, who was then High Priest, a great many gifts for the Temple here at Jerusalem, and caused him to send him six interpreters from each of the twelve tribes of Israel for the translation\footnote{Up to this point Cyril’s account is based upon the statements of the Pseudo-Aristeas. The fabulous incidents which follow, concerning the separate cells, the completion of the whole version by each translator, the miraculous agreement in the very words, proving a Divine inspiration, are found in Philo Judæus, \emph{Life of Moses}, II. 7. Josephus, \emph{Antiquities}, XII. \emph{c}. ii. 3–14, following the letter of Aristeas, gives long descriptions of the magnificent presents sent by Philadelphus to Jerusalem, and of his splendid hospitality to the translators, but makes no allusion to the separate cells or miraculous agreement. On the contrary he represents the 72 interpreters as meeting together for consultation, agreeing on the text to be adopted, and completing their joint labours in 72 days. The slightest comparison of the Version with the original Hebrew must convince any reasonable person that the idea of divine inspiration or supernatural assistance, borrowed by Justin Martyr, Irenæus, and other Fathers, apparently from Philo, is a mere invention of the imagination, disproved by the facts. Compare the article “Septuagint” in Murray’s Dictionary of the Bible.}. Then, further, to make experiment whether the books were Divine or not, he took precaution that those who had been sent should not combine among themselves, by assigning to each of the interpreters who had come his separate chamber in the island called Pharos, which lies over against Alexandria, and committed to each the whole Scriptures to translate. And when they had fulfilled the task in seventy-two days, he brought together all their translations, which they had made in different chambers without sending them one to another, and found that they agreed not only in the sense but even in words. For the process was no word-craft, nor contrivance of human devices: but the translation of the Divine Scriptures, spoken by the Holy Ghost, was of the Holy Ghost accomplished.

35. Of these read the two and twenty books, but have nothing to do with the apocryphal writings. Study earnestly these only which we read openly in the Church. Far wiser and more pious than thyself were the Apostles, and the bishops of old time, the presidents of the Church who handed down these books. Being therefore a child of the Church, trench\footnote{The rendering “trench not” (\textsc{R.W.C.}) agrees well with the etymology of the verb (\textgreek{παραχαράσσω}). Its more usual signification seems to be “counterfeit,” “forge.” The sense required here, apart from any metaphor, is “transgress” (Heurtley).} thou not upon its statutes. And of the Old Testament, as we have said, study the two and twenty books, which, if thou art desirous of learning, strive to remember by name, as I recite them. For of the Law the books of Moses are the first five, Genesis, Exodus, Leviticus, Numbers, Deuteronomy. And next, Joshua the son of Nave\footnote{The name “Nun” is represented by “Nave” in the Septuagint, which Cyril used.}, and the book of Judges, including Ruth, counted as seventh. And of the other historical books, the first and second books of the Kings\footnote{The two books of Samuel.} are among the Hebrews one book; also the third and fourth one book. And in like manner, the first and second of Chronicles are with them one book; and the first and second of Esdras are counted one. Esther is the twelfth book; and these are the Historical writings. But those which are written in verses are five, Job, and the book of Psalms, and Proverbs, and Ecclesiastes, and the Song of Songs, which is the seventeenth book. And after these come the five Prophetic books: of the Twelve Prophets one book, of Isaiah one, of Jeremiah one, including Baruch and Lamentations and the Epistle\footnote{The Epistle of Jeremy, which now appears in the Apocrypha as the last chapter of Baruch. On the number and arrangement of the Books of the Old and New Testaments the student should consult an interesting Essay by Professor Sanday (\emph{Studia Biblica}, vol. iii.), who traces the introduction of a fixed order to the time when papyrus \emph{rolls} were superseded by \emph{codices}, in which the sheets of skin were folded and bound together, as in printed books. This change had commenced before the Diocletian persecution, \textsc{a.d.} 303, when among the sacred books taken from the Christians \emph{codices} were much more numerous than \emph{rolls}. On the contents of the Jewish Canon, see Dictionary of the Bible, “Canon.” B.F.W. “Josephus enumerates 20 books ‘which are justly believed to be divine.’” One of the earliest attempts by a Christian to ascertain correctly the number and order of the Books of the O.T. was made by Melito, Bishop of Sardis, who travelled for this purpose to Palestine, in the latter part of the 2nd Century. His list is as follows:—“Of Moses five (books); Genesis, Exodus, Numbers, Leviticus, Deuteronomy, Jesus son of Nave, Judges, Ruth, four Books of Kings, two of Chronicles, Psalms of David, Solomon’s Proverbs, which is also called Wisdom, Ecclesiastes, Song of Songs, Job, Prophets, Isaiah, Jeremiah, the Twelve in one Book, Daniel, Ezekiel, Esdras.” (Eusebius, \emph{H.E.} III. cap. 10, note I, in this series.) Cyril’s List agrees with that of Athanasius (\emph{Festal Epistle}, 373 \textsc{a.d.}), except that Job is placed by Ath. after Canticles instead of before Psalms.}; then Ezekiel, and the Book of Daniel, the twenty-second of the Old Testament.

36. Then of the New Testament there are the four Gospels only, for the rest have false titles\footnote{Gr. \textgreek{ψευδεπίγραφα}. For an account of the many Apocryphal Gospels, see the article by Lipsius in the “\emph{Dictionary of Christian Biography},” Smith and Wace, and the English translations in Clark’s Ante-Nicene Library.} and are mischievous. The Manichæans also wrote a Gospel according to Thomas, which being tinctured with the fragrance of the evangelic title corrupts the souls of the simple sort. Receive also the Acts of the Twelve Apostles; and in addition to these the seven Catholic Epistles of James, Peter, John, and Jude; and as a seal upon them all, and the last work of the disciples, the fourteen Epistles of Paul\footnote{Cyril includes in this list all the books which we receive, except the Apocalypse. See Bishop Westcott’s Article “Canon,” in the \emph{Dictionary of the Bible}, and Origen’s Catalogue in Euseb. \emph{Hist.} vi. 25 (Nicene and Post-Nicene Fathers, vol. i.).}. But let all the rest be put aside in a secondary rank. And whatever books are not read in Churches, these read not even by thyself, as thou hast heard me say. Thus much of these subjects.

37. But shun thou every diabolical operation, and believe not the apostate Serpent, whose transformation from a good nature was of his own free choice: who can over-persuade the willing, but can compel no one. Also give heed neither to observations of the stars nor auguries, nor omens, nor to the fabulous divinations of the Greeks\footnote{Compare xix. 8. where all such acts of divination are said to be service of the devil.}. Witchcraft, and enchantment, and the wicked practices of necromancy, admit not even to a hearing. From every kind of intemperance stand aloof, giving thyself neither to gluttony nor licentiousness, rising superior to all covetousness and usury. Neither venture thyself at heathen assemblies for public spectacles, nor ever use amulets in sicknesses; shun also all the vulgarity of tavern-haunting. Fall not away either into the sect of the Samaritans, or into Judaism: for Jesus Christ henceforth hath ransomed thee. Stand aloof from all observance of Sabbaths\footnote{Compare Gal. iv. 10, “Ye observe days.”}, and from calling any indifferent meats \emph{common or unclean}. But especially abhor all the assemblies of wicked heretics; and in every way make thine own soul safe, by fastings, prayers, almsgivings, and reading the oracles of God; that having lived the rest of thy life in the flesh in soberness and godly doctrine, thou mayest enjoy the one salvation which flows from Baptism; and thus enrolled in the armies of heaven by God and the Father, mayest also be deemed worthy of the heavenly crowns, in Christ Jesus our Lord, to Whom be the glory for ever and ever. Amen.

\xchapter{Lecture V. Of Faith.}

\textsc{Hebrews xi. 1, 2}

Now faith is the substance of things hoped for, the evidence of things not seen. For by it the elders obtained a good report.

1. \textsc{How} great a dignity the Lord bestows on you in transferring you from the order of Catechumens to that of the Faithful, the Apostle Paul shews, when he affirms, \emph{God is faithful, by Whom ye were called into the fellowship of His Son Jesus Christ}\footnote{1 Cor. i. 9.}. For since God is called Faithful, thou also in receiving this title receivest a great dignity. For as God is called Good, and Just, and Almighty, and Maker of the Universe, so is He also called Faithful. Consider therefore to what a dignity thou art rising, seeing thou art to become partaker of a title of God\footnote{See Procatechesis 6, and Index, \emph{Faithful}.}.

2. Here then it is further required, that each of you be found faithful in his conscience: \emph{for a faithful man it is hard to find}\footnote{Prov. xx. 6.}: not that thou shouldest shew thy conscience to me, for thou art not to \emph{be judged of man’s judgment}\footnote{1 Cor. iv. 3. See Index, \emph{Confession.}}; but that thou shew the sincerity of thy faith to God, \emph{who trieth the reins and hearts}\footnote{Ps. vii. 9.}, and \emph{knoweth the thoughts of men}\footnote{Ps. xciv. 11.}. A great thing is a faithful man, being richest of all rich men. For \emph{to the faithful man belongs the whole world of wealth}\footnote{This sentence is a spurious addition to the text of the Septuagint, variously placed after Prov. xvii. 4, and xvii. 6. The thought is there completed by the antithesis, \emph{but to the faithless not even an obol}. The origin of the interpolation is unknown.}, in that he disdains and tramples on it. For they who in appearance are rich, and have many possessions, are poor in soul: since the more they gather, the more they pine with longing for what is still lacking. But the faithful man, most strange paradox, in poverty is rich: for knowing that we need only to have \emph{food and raiment}, and being \emph{therewith content}\footnote{1 Tim. vi. 8.}, he has trodden riches under foot.

3. Nor is it only among us, who bear the name of Christ, that the dignity of faith is great\footnote{It was a common objection of Pagan philosophers that the Christian religion was not founded upon reason but only on faith. \par Cyril’s answer that faith is necessary in the ordinary affairs of life is the same which Origen had employed against Celsus (I. 11): “Why should it not be more reasonable, since all human affairs are dependent upon faith, to believe God rather than men? For who takes a voyage, or marries, or begets children, or casts seeds into the ground, without believing that better things will result, although the contrary might and sometimes does happen?” See also Arnobius, \emph{adversus Gentes}, II. 8; and Hooker’s allusion to the scornful reproach of Julian the Apostate, “The highest point of your wisdom is \emph{believe}” (\emph{Eccles. Pol.} V. lxiii. 1.).}: but likewise all things that are accomplished in the world, even by those who are aliens\footnote{By “aliens from the Church,” and “those who are without,” S. Cyril here means Pagans: so Tertullian, \emph{de Idololatriâ}, c. xiv. But the latter term is applied to a Catechumen in Procatechesis. c. 12, and was also a common description of heretics: see Tertullian, \emph{de Baptismo}, c. xv.} from the Church, are accomplished by faith.

By faith the laws of marriage yoke together those who have lived as strangers: and because of the faith in marriage contracts a stranger is made partner of a stranger’s person and possessions. By faith husbandry also is sustained, for he who believes not that he shall receive a harvest endures not the toils. By faith sea-faring men, trusting to the thinnest plank, exchange that most solid element, the land, for the restless motion of the waves, committing themselves to uncertain hopes, and carrying with them a faith more sure than any anchor. By faith therefore most of men’s affairs are held together: and not among us only has there been this belief, but also, as I have said, among those who are without\footnote{By “aliens from the Church,” and “those who are without,” S. Cyril here means Pagans: so Tertullian, \emph{de Idololatriâ}, c. xiv. But the latter term is applied to a Catechumen in Procatechesis. c. 12, and was also a common description of heretics: see Tertullian, \emph{de Baptismo}, c. xv.}. For if they receive not the Scriptures, but bring forward certain doctrines of their own, even these they accept by faith.

4. The lesson also which was read to-day invites you to the true faith, by setting before you the way in which you also must please God: for it affirms that \emph{without faith it is impossible to please Him}\footnote{Heb. xi. 6.}. For when will a man resolve to serve God, unless he believes that \emph{He is a giver of reward?} When will a young woman choose a virgin life, or a young man live soberly, if they believe not that for chastity there is \emph{a crown that fadeth not away}\footnote{1 Pet. v. 4.}? Faith is an eye that enlightens every conscience, and imparts understanding; for the Prophet saith, \emph{And if ye believe not, ye shall not understand}\footnote{Is. vii. 9, according to the Septuagint. But A.V. and R.V. both render: \emph{If ye will not believe, surely ye shall not be established}.}.

Faith \emph{stoppeth the mouths of lions}\footnote{Heb. xi. 34.}, as in Daniel’s case: for the Scripture saith concerning him, that \emph{Daniel was brought up out of the den, and no manner of hurt was found upon him, because he believed in his God}\footnote{Dan. vi. 23.}. Is there anything more fearful than the devil? Yet even against him we have no other shield than faith\footnote{1 Pet. v. 9: \emph{Whom resist, stedfast in the faith.}}, an impalpable buckler against an unseen foe. For he sends forth divers arrows, and \emph{shoots down in the dark night}\footnote{Ps. xi. 2, \emph{that they may shoot in darkness at the upright in heart} (R.V.). The Hebrew word \texthebrew{ל\texthebrew{פֶא}}, signifying deep darkness (Job iii. 6; x. 22) is vigorously rendered by the Seventy \textgreek{σκοτομήνη}, which is explained by the Scholiast on Homer (Od. xiv. 457: \textgreek{Νὺξ δ᾽ ἄρ᾽ ἐπῆλθε κακὴ σκοτομήνιος}) to be the deep darkness of the night preceding the new moon.} those that watch not; but, since the enemy is unseen, we have faith as our strong armour, according to the saying of the Apostle, \emph{In all things taking the shield of faith, wherewith ye shall be able to quench all the fiery darts of the wicked one}\footnote{Eph. vi. 16.}. A fiery dart of desire of base indulgence is often cast forth from the devil: but faith, suggesting a picture of the judgment, cools down the mind, and quenches the dart.

5. There is much to tell of faith, and the whole day would not be time sufficient for us to describe it fully. At present let us be content with Abraham only, as one of the examples from the Old Testament, seeing that we have been made his sons through faith. He was justified not only by works, but also by faith\footnote{James ii. 21. Casaubon omitted \textgreek{μόνον}, which is found in every \textsc{ms.,} thus making the meaning to be, “He was justified not by works but by faith,” which directly contradicts the statement of S. James, and is inconsistent with the following context in S. Cyril.}: for though he did many things well, yet he was never called the friend of God\footnote{James ii. 23; 2 Chron. xx. 7; Is. xli. 8; Gen. xv. 6.}, except when he believed. Moreover, his every work was performed in faith. Through faith he left his parents; left country, and place, and home through faith\footnote{Heb. xi. 8–10.}. In like manner, therefore, as he was justified be thou justified also. In his body he was already dead in regard to offspring, and Sarah his wife was now old, and there was no hope left of having children. God promises the old man a child, and Abraham \emph{without being weakened in faith, though he considered his own body now as good as dead}\footnote{Rom. iv. 19.}, heeded not the weakness of his body, but the power of Him who promised, because \emph{he counted Him faithful who had promised}\footnote{Heb. xi. 11, 12.}, and so beyond all expectation gained the child from bodies as it were already dead. And when, after he had gained his son, he was commanded to offer him up, although he had heard the word, \emph{In Isaac shall thy seed be called}\footnote{Gen. xxi. 12; xxii. 2.}, he proceeded to offer up his son, his only son, to God, believing \emph{that God is able to raise up even from the dead}\footnote{Heb. xi. 19.}. And having bound his son, and laid him on the wood, he did in purpose offer him, but by the goodness of God in delivering to him a lamb instead of his child, he received his son alive. Being faithful in these things, he was sealed for righteousness, \emph{and received circumcision as a seal of the faith which he had while he was in uncircumcision}\footnote{Rom. iv. 11.}, having received a promise \emph{that he should be the father of many nations}\footnote{Gen. xvii. 5.}.

6. Let us see, then, how Abraham is the father of many nations\footnote{Rom. iv. 17, 18.}. Of Jews he is confessedly the father, through succession according to the flesh. But if we hold to the succession according to the flesh, we shall be compelled to say that the oracle was false. For according to the flesh he is no longer father of us all: but the example of his faith makes us all sons of Abraham. How? and in what manner? With men it is incredible that one should rise from the dead; as in like manner it is incredible also that there should be offspring from aged persons as good as dead. But when Christ is preached as having been crucified on the tree, and as having died and risen again, we believe it. By the likeness therefore of our faith we are adopted into the sonship of Abraham. And then, following upon our faith, we receive like him the spiritual seal, being circumcised by the Holy Spirit through Baptism, not in the foreskin of the body, but in the heart, according to Jeremiah, saying, \emph{And ye shall be circumcised unto God in the foreskin of your heart}\footnote{Jer. iv. 4: \emph{Circumcise yourselves to the Lord, and take away the foreskins of your heart.} The Septuagint agrees closely with the Hebrew, but Cyril quotes freely from memory.}: and according to the Apostle, \emph{in the circumcision of Christ, having been buried with Him in baptism}, and the rest\footnote{Col. ii. 11, 12.}.

7. This faith if we keep we shall be free from condemnation, and shall be adorned with all kinds of virtues. For so great is the strength of faith, as even to buoy men up in walking on the sea. Peter was a man like ourselves, made up of flesh and blood, and living upon like food. But when Jesus said, \emph{Come}\footnote{Matt. xiv. 29.}, he believed, and walked upon the waters, and found his faith safer upon the waters than any ground; and his heavy body was upheld by the buoyancy of his faith. But though he had safe footing over the water as long as he believed, yet when he doubted, at once he began to sink: for as his faith gradually relaxed, his body also was drawn down with it. And when He saw his distress, Jesus who remedies the distresses of our souls, said, \emph{O thou of little faith, wherefore didst thou doubt}\footnote{Mark xiv. 31.}? And being nerved again by Him who grasped his right hand, he had no sooner recovered his faith, than, led by the hand of the Master, he resumed the same walking upon the waters: for this the Gospel indirectly mentioned, saying, \emph{when they were gone up into the ship}\footnote{Ib. 32.}. For it says not that Peter swam across and went up, but gives us to understand that, after returning the same distance that he went to meet Jesus, he went up again into the ship.

8. Yea, so much power hath faith, that not the believer only is saved, but some have been saved by others believing. The paralytic in Capernaum was not a believer, but they believed who brought him, and let him down through the tiles\footnote{Mark ii. 4.}: for the sick man’s soul shared the sickness of his body. And think not that I accuse him without cause: the Gospel itself says, \emph{when Jesus saw}, not his faith, but \emph{their faith}, He saith to the sick of the palsy, Arise\footnote{Matt. ix. 2, 6.}! The bearers believed, and the sick of the palsy enjoyed the blessing of the cure.

9. Wouldest thou see yet more surely that some are saved by others’ faith? Lazarus died\footnote{John xi. 14–44.}: one day had passed, and a second, and a third: his sinews\footnote{\textgreek{νεῦρα}. “Sinews” is the original meaning, the application to “nerves,” as distinct organs of sensation, being later.} were decayed, and corruption was preying already upon his body. How could one four days dead believe, and entreat the Redeemer on his own behalf? But what the dead man lacked was supplied by his true sisters. For when the Lord was come, the sister fell down before Him, and when He said, \emph{Where have ye laid him?} and she had made answer, \emph{Lord, by this time he stinketh; for he hath been four days dead}, the Lord said, \emph{If thou believe, thou shalt see the glory of God}; as much as saying, Supply thou the dead man’s lack of faith: and the sisters’ faith had so much power, that it recalled the dead from the gates of hell. Have then men by believing, the one on behalf of the other, been able to raise\footnote{For \textgreek{ἀναστῆναι}, retained by the Benedictine Editor and Reischl, read \textgreek{ἀναστῆσαι}, with Roe, Casaubon, and Alexandrides.} the dead, and shalt not thou, if thou believe sincerely on thine own behalf, be much rather profited? Nay, even if thou be faithless, or of little faith, the Lord is loving unto man; He condescends to thee on thy repentance: only on thy part say with honest mind, \emph{Lord, I believe, help thou mine unbelief}\footnote{Mark ix. 24.}. But if thou thinkest that thou really art faithful, but hast not yet the fulness of faith, thou too hast need to say like the Apostles, \emph{Lord, increase our faith}\footnote{Luke xvii. 5.}: for some part thou hast of thyself, but the greater part thou receivest from Him.

10. For the name of Faith is in the form of speech\footnote{\textgreek{κατὰ τὴν προσηγορίαν}. Compare Aristotle, \emph{Categories}, V. 30: \textgreek{τῷ σχήματι τῆς προσηγορίας}. Cyril’s description of faith as twofold, and of dogmatic faith as an assent (\textgreek{συγκατάθεσις}) of the soul to something as credible, seems to be derived from Clement of Alexandria, Strom. II. c. 12. Compare by all means Pearson on the Creed, Art. I. and his Notes a, b, c.} one, but has two distinct senses. For there is one kind of faith, the dogmatic, involving an assent of the soul on some particular point: and it is profitable to the soul, as the Lord saith: \emph{He that heareth My words, and believeth Him that sent Me, hath everlasting life, and cometh not into judgment}\footnote{John v. 24.}: and again, \emph{He that believeth in the Son is not judged, but hath passed from death unto life}\footnote{Ib. iii. 18; v. 24.}. Oh the great loving-kindness of God! For the righteous were many years in pleasing Him: but what they succeeded in gaining by many years of well-pleasing\footnote{\textgreek{εὐαρεστήσεως} , Bened. and Reischl, with best \textsc{mss.} Milles and the earlier editions have \textgreek{ἐρευνήσεως}, “searching.”}, this Jesus now bestows on thee in a single hour. For if thou shalt believe that Jesus Christ is Lord, and that God raised Him from the dead, thou shalt be saved, and shalt be transported into Paradise by Him who brought in thither the robber. And doubt not whether it is possible; for He who on this sacred Golgotha saved the robber after one single hour of belief, the same shall save thee also on thy believing\footnote{Luke xxiii. 43; the argument is used again in Cat. xiii. 31.}.

11. But there is a second kind of faith, which is bestowed by Christ as a gift of grace. \emph{For to one is given through the Spirit the word of wisdom, and to another the word of knowledge according to the same Spirit: to another faith, by the same Spirit, and to another gifts of healing}\footnote{1 Cor. xii. 8, 9.}. This faith then which is given of grace from the Spirit is not merely doctrinal, but also worketh things above man’s power. For whosoever hath this faith, \emph{shall say to this mountain, Remove hence to yonder place, and it shall remove}\footnote{Mark xi. 23.}. For whenever any one shall say this in faith, \emph{believing that it cometh to pass, and shall not doubt in his heart, then receiveth he the grace}.

And of this faith it is said, \emph{If ye have faith as a grain of mustard seed}\footnote{Matt. xvii. 20.}. For just as the grain of mustard seed is small in size, but fiery in its operation, and though sown in a small space has a circle of great branches, and when grown up is able even to shelter the fowls\footnote{Matt. xiii. 32.}; so, likewise, faith in the swiftest moment works the greatest effects in the soul. For, when enlightened by faith, the soul hath visions of God, and as far as is possible beholds God, and ranges round the bounds of the universe, and before the end of this world already beholds the Judgment, and the payment of the promised rewards. Have thou therefore that faith in Him which cometh from thine own self, that thou mayest also receive from Him that faith which worketh things above man\footnote{S. Chrysostom (Hom. xxix. in 1 Cor. xii. 9, 10) in like manner distinguishes dogmatic faith from the faith which is “the mother of miracles.” The former S. Cyril calls our own, not meaning that God’s help is not needed for it, but because, as he has shewn in § 10, it consists in the mind’s assent, and voluntary approval of the doctrines set before it: but the latter is a pure gift of grace working in man without his own help. Compare \emph{Apostolic Constitutions}, VIII. c. 1.}.

12. But in learning the Faith and in professing it, acquire and keep that only, which is now delivered\footnote{This Lecture was to be immediately followed by a first recitation of the Creed. See Index, \emph{Creed}.} to thee by the Church, and which has been built up strongly out of all the Scriptures. For since all cannot read the Scriptures, some being hindered as to the knowledge of them by want of learning, and others by a want of leisure, in order that the soul may not perish from ignorance, we comprise the whole doctrine of the Faith in a few lines. This summary I wish you both to commit to memory when I recite it\footnote{\textgreek{ἐπ᾽ αὐτῆς τῆς λέξεως}. “in ipsâ lectione” (Milles): “ipsis verbis” (Bened.): “in the very phrase” (\textsc{R.W.C.}). See below, note 4.}, and to rehearse it with all diligence among yourselves, not writing it out on paper\footnote{Compare S. August. Serm. ccxii., “At the delivery of the Creed,” and Index, \emph{Creed}.}, but engraving it by the memory upon your heart\footnote{Compare Æschylus, \emph{Prometheus} V. 789: \textgreek{ἣν ἐγγράφου σὺ μνήμοσιν δέλτοις φρενῶν}.}, taking care while you rehearse it that no Catechumen chance to overhear the things which have been delivered to you. I wish you also to keep this as a provision\footnote{\textgreek{ἐφόδιον}, \emph{Viaticum}, i.e. provision for a journey, and here for the journey through this life. It is applied metaphorically by other Fathers (a) in this general sense, to the reading of Holy Scripture, Prayer, and Baptism, and (b) in a special sense to the Holy Eucharist when administered to the sick and dying, as a preparation for departure to the life after death. Council of Nicæa (\textsc{a.d.} 325), Canon xiii. “With respect to the dying, the old rule of the Church should continue to be observed, which forbids that any one who is on the point of death should be deprived of the last and most necessary \emph{viaticum} (\textgreek{ἐφόδιον}).”} through the whole course of your life, and beside this to receive no other, neither if we ourselves should change and contradict our present teaching, nor if an adverse angel, \emph{transformed into an angel of light}\footnote{2 Cor. xi. 14.} should wish to lead you astray. \emph{For though we or an angel from heaven preach to you any other gospel than that ye have received, let him be to you anathema}\footnote{Gal. i. 8, 9.}. So for the present listen while I simply say the Creed\footnote{\textgreek{ἐπ᾽ αὐτῆς τῆς λέξεως}. (Bened. Reischl. with best \textsc{mss.}). \textgreek{ταύτης τῆς λέξεως}, “this my recitation,” (Milles).}, and commit it to memory; but at the proper season expect the confirmation out of Holy Scripture of each part of the contents. For the articles of the Faith were not composed as seemed good to men; but the most important points collected out of all the Scripture make up one complete teaching of the Faith. And just as the mustard seed in one small grain contains many branches, so also this Faith has embraced in few words all the knowledge of godliness in the Old and New Testaments. Take heed then, brethren, and \emph{hold fast the traditions}\footnote{2 Thess. ii. 15. Compare Cat. xxiii. 23.} which ye now receive, and \emph{write them an the table of your heart}\footnote{Prov. vii. 3. Note 9, above.}.

13. Guard them with reverence, lest per chance the enemy despoil any who have grown slack; or lest some heretic pervert any of the truths delivered to you. For faith is like putting money into the bank\footnote{Matt. xxv. 27; Luke xix. 23. See note on Catech. vi. 36: “Be thou a good banker.”}, even as we have now done; but from you God requires the accounts of the deposit. \emph{I charge you}, as the Apostle saith, \emph{before God, who quickeneth all things, and Christ Jesus, who before Pontius Pilate witnessed the good confession, that ye keep} this faith which is committed to you, \emph{without spot, until the appearing of our Lord Jesus Christ}\footnote{1 Tim. v. 21; vi. 13, 14.}. A treasure of life has now been committed to thee, and the Master demandeth the deposit at His appearing, \emph{which in His own times He shall shew, Who is the blessed and only Potentate, the King of kings, and Lord of lords; Who only hath immortality, dwelling in light which no man can approach unto; Whom no man hath seen nor can see. To Whom be glory, honour, and power}\footnote{1 Tim. vi. 15, 16.} for ever and ever. Amen.

\xchapter{Lecture VI. Concerning the Unity of God. On the Article, I Believe in One God. Also Concerning Heresies.}

Concerning the Unity of God\footnote{\textgreek{Περὶ Θεοῦ Μοναρχίας}. The word \textgreek{μοναρχία}, as used by Plato (\emph{Polit}. 291 C), Aristotle (\emph{Polit}. III. xiv. 11. \textgreek{εἶδος μοναρχίας βασιλικῆς}), Philo Judæus (\emph{de Circumcisione}, § 2; \emph{de Monarchia}, Titul.), means “sole government.” Compare Tertullian (\emph{adv. Praxean}. c. iii.): “If I have gained any knowledge of either language, I am sure that \textgreek{Μοναρχία} has no other meaning than ‘single and individual rule.’” Athanasius (\emph{de Decretis Nicænæ Synodi}, § 26) has preserved part of an Epistle of Dionysius, Bishop of Rome (259–269, \textsc{a.d}.), against the Sabellians: “It will be natural for me now to speak against those who divide, and cut into pieces, and destroy that most sacred doctrine of the Church of God, the Monarchia, making it, as it were, three powers and divided hypostases, and three Godheads;” (\emph{ibid}.): “It is the doctrine of the presumptuous Marcion to sever and divide the Monarchia into three origins (\textgreek{ἀρχάς}).” We see here the sense which \textgreek{Μοναρχία} had acquired in Christian Theology: it meant the “Unity of God,” as the one principle and origin of all things. “By the Monarchy is meant the doctrine that the Second and Third Persons in the Ever-blessed Trinity are ever to be referred in our thoughts to the First, as the Fountain of Godhead” (Newman, Athanas. \emph{de Decretis Nic}. \emph{Syn}. § 26, note h). Justin Martyr (Euseb. \emph{H.E}. IV. 18), and Irenæus (\emph{ibid}. V. 20), had each written a treatise \textgreek{περὶ Μοναρχίας}. On the history of Monarchianism see, in this Series, Athanasius, \emph{Prolegomena}, p. xxiii. \emph{sqq}.}. On the Article, I Believe in One God. Also Concerning Heresies.

\textsc{Isaiah xlv. 16, 17}\textsc{. (Sept.)}

Sanctify yourselves unto Me, O islands. Israel is saved by the Lord with an everlasting salvation; they shall not be ashamed, neither shall they be confounded for ever, \&c.

1. \emph{Blessed be the God and Father of our Lord Jesus Christ}\footnote{2 Cor. i. 3.}. Blessed also be His Only-begotten Son\footnote{This clause is omitted in some \textsc{mss.} Various forms of the Doxology were adopted in Cyril’s time by various parties in the Church. Thus Theodoret (\emph{Hist. Eccles}. II. c. 19) relates that Leontius, Bishop of Antioch, \textsc{a.d.} 348–357, observing that the Clergy and the Congregation were divided into two parties, the one using the form “and to the Son, and to the Holy Ghost,” the other “through the Son, in the Holy Ghost,” used to repeat the Doxology silently, so that those who were near could hear only “world without end.” \par The form which was regarded as the most orthodox, and adopted in the Liturgies ran thus: “Glory to the Father, and to the Son, and to the Holy Ghost, now and ever, and to the ages of the ages.” See Suicer’s Thesaurus, \textgreek{Δοξολογία}.}. For with the thought of \emph{God} let the thought of \emph{Father} at once be joined, that the ascription of glory to the Father and the Son may be made indivisible. For the Father hath not one glory, and the Son another, but one and the same, since He is the Father’s Only-begotten Son; and when the Father is glorified, the Son also shares the glory with Him, because the glory of the Son flows from His Father’s honour: and again, when the Son is glorified, the Father of so great a blessing is highly honoured.

2. Now though the mind is most rapid in its thoughts, yet the tongue needs words, and a long recital of intermediary speech. For the eye embraces at once a multitude of the ‘starry quire;’ but when any one wishes to describe them one by one, which is the Morning-star, and which, the Evening-star, and which each one of them, he has need of many words. In like manner again the mind in the briefest moment compasses earth and sea and all the bounds of the universe; but what it conceives in an instant, it uses many words to describe\footnote{Irenæus II. xxviii. 4: “But since God is all mind, all reason, all active Spirit, all light, and always exists as one and the same, such conditions and divisions (of operation) cannot fittingly be ascribed to Him. For our tongue, as being made of flesh, is not able to minister to the rapidity of man’s sense, because that is of a spiritual nature; for which reason our speech is restrained (\emph{suffocatur}) within us, and is not at once expressed as it has been conceived in the mind but is uttered by successive efforts, just as the tongue is able to serve it.”}. Yet forcible as is the example I have mentioned, still it is after all weak and inadequate. For of God we speak not all we ought (for that is known to Him only), but so much as the capacity of human nature has received, and so much as our weakness can bear. For we explain not what God is but candidly confess that we have not exact knowledge concerning Him. For in what concerns God to confess our ignorance is the best knowledge\footnote{Tertullian, \emph{Apologeticus}, § 17: “That which is infinite is known only to itself. This it is which gives some notion of God, while yet beyond all our conceptions—our very incapacity of fully grasping Him affords us the idea of what He really is. He is presented to our minds in His transcendent greatness, as at once known and unknown.” Cf. Phil. Jud. \emph{de Monarch}. i. 4: Hooker, \emph{Eccles. Pol}. I. ii. 3: “Whom although to know be life, and joy to make mention of His name; yet our soundest knowledge is to know that we know Him not as He is, neither can know Him.”}. Therefore magnify the Lord with me, and let us exalt His Name together\footnote{Ps. xxxiv. 3.},—all of us in common, for one alone is powerless; nay rather, even if we be all united together, we shall yet not do it as we ought. I mean not you only who are here present, but even if all the nurslings of the whole Church throughout the world, both that which now is and that which shall be, should meet together, they would not be able worthily to sing the praises of their Shepherd.

3. A great and honourable man was Abraham, but only great in comparison with men; and when he came before God, then speaking the truth candidly he saith, \emph{I am earth and ashes}\footnote{Gen. xviii. 27.}. He did not say ‘earth,’ and then cease, lest he should call himself by the name of that great element; but he added \emph{‘and ashes,’} that he might represent his perishable and frail nature. Is there anything, he saith, smaller or lighter than ashes? For take, saith he, the comparison of ashes to a house, of a house to a city, a city to a province, a province to the Roman Empire, and the Roman Empire to the whole earth and all its bounds, and the whole earth to the heaven in which it is embosomed;—the earth, which bears the same proportion to the heaven as the centre to the whole circumference of a wheel, for the earth is no more than this in comparison with the heaven\footnote{The opinion of Aristarchus of Samos, as stated by Archimedes (\emph{Arenarius}, p. 320, Oxon), was that the sphere of the fixed stars was so large, that it bore to the earth’s orbit the same proportion as a sphere to its centre, or more correctly (as Archimedes explains) the same proportion as the earth’s orbit round the sun to the earth itself. Compare Cat. xv. 24.}: consider then that this first heaven which is seen is less than the second, and the second than the third, for so far Scripture has named them, not that they are only so many, but because it was expedient for us to know so many only. And when in thought thou hast surveyed all the heavens, not yet will even the heavens be able to praise God as He is, nay, not if they should resound with a voice louder than thunder. But if these great vaults of the heavens cannot worthily sing God’s praise, when shall \emph{‘earth and ashes,’} the smallest and least of things existing, be able to send up a worthy hymn of praise to God, or worthily to speak of God, \emph{that sitteth upon the circle of the earth, and holdeth the inhabitants thereof as grasshoppers}\footnote{Is. xl. 22.}.

4. If any man attempt to speak of God, let him first describe the bounds of the earth. Thou dwellest on the earth, and the limit of this earth which is thy dwelling thou knowest not: how then shalt thou be able to form a worthy thought of its Creator? Thou beholdest the stars, but their Maker thou beholdest not: count these which are visible, and then describe Him who is invisible, \emph{Who telleth the number of the stars, and calleth them all by their names}\footnote{Ps. cxlvii. 4.}. Violent rains lately came pouring down upon us, and nearly destroyed us: number the drops in this city alone: nay, I say not in the city, but number the drops on thine own house for one single hour, if thou canst: but thou canst not. Learn then thine own weakness; learn from this instance the mightiness of God: for \emph{He hath numbered the drops of rain}\footnote{Job xxxvi. 27: \textgreek{ἀριθμηταὶ δὲ αὐτῷ σταγόνες ὑετοῦ}. \textsc{R.V.} \emph{For He draweth up the drops of water}.}, which have been poured down on all the earth, not only now but in all time. The sun is a work of God, which, great though it be, is but a spot in comparison with the whole heaven; first gaze stedfastly upon the sun, and then curiously scan the Lord of the sun. \emph{Seek not the things that are too deep for thee, neither search out the things that are above thy strength: what is commanded thee, think thereupon}\footnote{Ecclus. iii. 21, 22.}.

5. But some one will say, If the Divine substance is incomprehensible, why then dost thou discourse of these things? So then, because I cannot drink up all the river, am I not even to take in moderation what is expedient for me? Because with eyes so constituted as mine I cannot take in all the sun, am I not even to look upon him enough to satisfy my wants? Or again, because I have entered into a great garden, and cannot eat all the supply of fruits, wouldst thou have me go away altogether hungry? I praise and glorify Him that made us; for it is a divine command which saith, \emph{Let every breath praise the Lord}\footnote{Ps. cl. 6.}. I am attempting now to glorify the Lord, but not to describe Him, knowing nevertheless that I shall fall short of glorifying Him worthily, yet deeming it a work of piety even to attempt it at all. For the Lord Jesus encourageth my weakness, by saying, \emph{No man hath seen God at any time}\footnote{John i. 18. They are the Evangelist’s own words.}.

6. What then, some man will say, is it not written, \emph{The little ones’ Angels do always behold the face of My Father which is in heaven}\footnote{Matt. xviii. 10.}? Yes, but the Angels see God not as He is, but as far as they themselves are capable. For it is Jesus Himself who saith, \emph{Not that any man hath seen the Father, save He which is of God, He hath seen the Father}\footnote{John vi. 46.}. The Angels therefore behold as much as they can bear, and Archangels as much as they are able; and Thrones and Dominions more than the former, but yet less than His worthiness: for with the Son the Holy Ghost alone can rightly behold Him: for \emph{He searcheth all things, and knoweth even the deep things of God}\footnote{1 Cor. ii. 10.}: as indeed the Only-begotten Son also, with the Holy Ghost, knoweth the Father fully: For \emph{neither}, saith He, \emph{knoweth any man the Father, save the Son, and he to whom the Son will reveal Him}\footnote{Matt. xi. 27.}. For He fully beholdeth, and, according as each can bear, revealeth God through the Spirit: since the Only-begotten Son together with the Holy Ghost is a partaker of the Father’s Godhead. He, who\footnote{The Benedictine and earlier printed texts read \textgreek{ὁ γεννηθεὶς [ἀπαθῶς πρὸ τῶν χρόνων αἰωνίων]}: but the words in brackets are not found in the best \textsc{mss.} The false grammar betrays a spurious insertion, which also interrupts the sense. On the meaning of the phrase \textgreek{ὁ γεννηθεὶς ἀπαθῶς}, see note on vii. 5: \textgreek{οὐ πάθει πατὴρ γενόμενος}.} was begotten knoweth Him who begat; and He Who begat knoweth Him who is begotten. Since Angels then are ignorant (for to each according to his own capacity doth the Only-begotten reveal Him through the Holy Ghost, as we have said), let no man be ashamed to confess his ignorance. I am speaking now, as all do on occasion: but how we speak, we cannot tell: how then can I declare Him who hath given us speech? I who have a soul, and cannot tell its distinctive properties, how shall I be able to describe its Giver?

7. For devotion it suffices us simply to know that we have a God; a God who is One, a living\footnote{Gr. \textgreek{ὄντα, ἀεὶ ὄντα}.}, an ever-living God; always like unto Himself\footnote{Iren. II. xiii. 3: “He is altogether like and equal to Himself; since He is all sense, and all spirit, and all feeling, and all thought, and all reason, and all hearing, and all ear, and all eye, and all light, and all a fount of every good,—even as the religious and pious are wont to speak of God.”}; who has no Father, none mightier than Himself, no successor to thrust Him out from His kingdom: Who in name is manifold, in power infinite, in substance uniform\footnote{\textgreek{μονοειδῆ}. A Platonic word. \emph{Phædo}, 80 B: \textgreek{τῷ μὲν θείω καὶ ἀθανάτῳ καὶ νοητῷ καὶ μονοειδεῖ καὶ ἀδιαλύτῳ καὶ ἀεὶ ὡσαύτως κατὰ τὰ αὐτὰ ἔχοντι ἑαυτῷ ὁμοιότατον εἶναι ψυχήν}. See Index, “Hypostasis.”}. For though He is called Good, and Just, and Almighty and Sabaoth\footnote{Iren. II. xxxv. 3: “If any object that in the Hebrew language different expressions occur, such as Sabaoth, Elöe, Adonai, and all other such terms, striving to prove from these that there are different powers and Gods, let them learn that all expressions of this kind are titles and announcements of one and the same Being.”}, He is not on that account diverse and various; but being one and the same, He sends forth countless operations of His Godhead, not exceeding here and deficient there, but being in all things like unto Himself. Not great in loving-kindness only, and little in wisdom, but with wisdom and loving-kindness in equal power: not seeing in part, and in part devoid of sight; but being all eye, and all ear, and all mind\footnote{See the passages of Irenæus quoted above, § 2 note 4, and § 7 note 3.}: not like us perceiving in part and in part not knowing; for such a statement were blasphemous, and unworthy of the Divine substance. He foreknoweth the things that be; He is Holy, and Almighty, and excelleth all in goodness, and majesty, and wisdom: of Whom we can declare neither beginning, nor form, nor shape. For \emph{ye have neither heard His voice at any time, nor seen His shape}\footnote{John v. 37.}, saith Holy Scripture. Wherefore Moses saith also to the Israelites: \emph{And take ye good heed to your own souls, for ye saw no similitude}\footnote{Deut. iv. 15.}. For if it is wholly impossible to imagine His likeness, how shall thought come near His substance?

8. There have been many imaginations by many persons, and all have failed. Some have thought that God is fire; others that He is, as it were, a man with wings, because of a true text ill understood, \emph{Thou shalt hide me under the shadow of Thy wings}\footnote{Ps. xvii. 8.}. They forgot that our Lord Jesus Christ, the Only-begotten, speaks in like manner concerning Himself to Jerusalem, \emph{How often would I have gathered thy children together even as a hen doth gather her chickens under her wings, and ye would not}\footnote{Matt. xxiii. 37.}. For whereas God’s protecting power was conceived as wings, they failing to understand this sank down to the level of things human, and supposed that the Unsearchable exists in the likeness of man. Some again dared to say that He has seven eyes, because it is written, \emph{seven eyes of the Lord looking upon the whole earth}\footnote{Zech. iv. 10.}. For if He has but seven eyes surrounding Him in part, His seeing is therefore partial and not perfect: but to say this of God is blasphemous; for we must believe that God is in all things perfect, according to our Saviour’s word, which saith, \emph{Your Father in heaven is perfect}\footnote{Matt. v. 48.}: perfect in sight, perfect in power, perfect in greatness, perfect in foreknowledge, perfect in goodness, perfect in justice, perfect in loving-kindness: not circumscribed in any space, but the Creator of all space, existing in all, and circumscribed by none\footnote{Philo Judæus (\emph{Leg. Alleg}. I. 14. p. 52). \textgreek{Θεοῦ γὰρ οὐδὲ ὁ σύμπας κόσμος ἀξίον ἂν εἴη χωρίον καὶ ἐνδιαίτημα, ἐπεὶ αὐτὸς ἑαυτῷ τήπος}. So Sir Isaac Newton, at the end of the Principia, asserts that God by His eternal and infinite existence constitutes Time and Space: “Non est duratio vel spatium, sed durat et adest, et existendo semper et ubique spatium et durationem constituit.”}. \emph{Heaven is His throne}, but higher is He that sitteth thereon: \emph{and earth is His footstool}\footnote{Is. lxvi. 1.}, but His power reacheth unto things under the earth.

9. One He is, everywhere present, beholding all things, perceiving all things, creating all things through Christ: \emph{For all things were made by Him, and without Him was not anything made}\footnote{John i. 3.}. A fountain of every good, abundant and unfailing, a river of blessings, an eternal light of never-failing splendour, an insuperable power condescending to our infirmities: whose very Name we dare not hear\footnote{The sacred name (\texthebrew{ה\texthebrew{והי}}) was not pronounced, but Adonai was substituted.}. \emph{Wilt thou find a footstep of the Lord?} saith Job, \emph{or hast thou attained unto the least things which the Almighty hath made}\footnote{Job xi. 7 (R.V.): \emph{Canst thou by searching find out God? Canst thou find out the Almighty unto perfection?} Cyril seems to have understood \textgreek{τὰ ἔσχατα} as “the least,” not as “the utmost.”}? If the least of His works are incomprehensible, shall He be comprehended who made them all? \emph{Eye hath not seen, and ear hath not heard, neither have entered into the heart of man, the things which God hath prepared for them that love Him}\footnote{1 Cor. ii. 9.}. If the things which God hath prepared are incomprehensible to our thoughts, how can we comprehend with our mind Himself who hath prepared them? \emph{O the depth of the riches, and wisdom, and knowledge of God! How unsearchable are His judgments, and His ways past finding out}\footnote{Rom. xi. 33.}! saith the Apostle. If His judgments and His ways are incomprehensible, can He Himself be comprehended?

10. God then being thus great, and yet greater, (for even were I to change my whole substance into tongue, I could not speak His excellence: nay more, not even if all Angels should assemble, could they ever speak His worth), God being therefore so great in goodness and majesty, man hath yet dared to say to a stone that he hath graven, \emph{Thou art my God}\footnote{Is. xliv. 17.}! O monstrous blindness, that from majesty so great came down so low! The tree which was planted by God, and nourished by the rain, and afterwards burnt and turned into ashes by the fire,—this is addressed as God, and the true God is despised. But the wickedness of idolatry grew yet more prodigal, and cat, and dog, and wolf\footnote{The cat was sacred to the goddess Pasht, called by the Greeks Bubastis, and identified by Herodotus (ii. 137) with Artemis or Diana. Cats were embalmed after death, and their mummies are found at various places, but especially at Bubastis (\emph{Herod}. ii. 67). \par “The Dogs are interred in the cities to which they belong, in sacred burial-places” (\emph{Herod}. ii. 67), but chiefly at Cynopolis (“City of Dogs”) where the dog-headed deity Anubis was worshipped. \par Mummies of wolves are found in chambers excavated in the rocks at Lycopolis, where Osiris was worshipped under the symbol of a wolf.} were worshipped instead of God: the man-eating lion\footnote{The lion was held sacred at Leontopolis (Strabo, xvii. p. 812).} also was worshipped instead of God, the most loving friend of man. The snake and the serpent\footnote{“In the neighbourhood of Thebes there are sacred serpents perfectly harmless to man. These they bury in the temple of Zeus, the god to whom they are sacred.” (\emph{Herod}. ii. 74.) \par At Epidaurus in Argolis the serpent was held sacred as the symbol of Æsculapius. Clement of Alexandria (\emph{Exhort}. c. ii.) gives a fuller list of animals worshipped by various nations. Compare also \emph{Clement. Recogn}. V. 20.}, counterfeit of him who thrust us out of Paradise, were worshipped, and He who planted Paradise was despised. And I am ashamed to say, and yet do say it, even onions\footnote{Juvenal \emph{Sat}. xv. 7. \par Illic aeluros, hic piscem fluminis, illic \par Oppida tota canem venerantur, nemo Dianam. \par Possum et caepe nefas violare et frangere morsu.} were worshipped among some. Wine was given \emph{to make glad the heart of man}\footnote{Ps. civ. 15.}: and Dionysus (Bacchus) was worshipped instead of God. God made corn by saying, \emph{Let the earth bring forth grass, yielding seed after his kind and after his likeness}\footnote{Gen. i. 11.}, \emph{that bread may strengthen man’s heart}\footnote{Ps. civ. 15.}: why then was Demeter (Ceres) worshipped? Fire cometh forth from striking stones together even to this day: how then was Hephæstus (Vulcan) the creator of fire?

11. Whence came the polytheistic error of the Greeks\footnote{The early Creeds of the Eastern Churches, like that which Eusebius of Cæsarea proposed at Nicæa, expressly declare the unity of God, in opposition both to the heathen Polytheism, and to the various heresies which introduced two or more Gods. See below in this Lecture, §§ 12–18; and compare Athan. (\emph{contra Gentes}, § 6, \emph{sqq}.)}? God has no body: whence then the adulteries alleged among those who are by them called gods? I say nothing of the transformations of Zeus into a swan: I am ashamed to speak of his transformations into a bull: for bellowings are unworthy of a god. The god of the Greeks has been found an adulterer, yet are they not ashamed: for if he is an adulterer let him not be called a god. They tell also of deaths\footnote{Clement of Alexandria (\emph{Exhort}. cap. ii. § 37), quotes a passage from a hymn of Callimachus, implying the death of Zeus: \par “For even thy tomb, O king, \par The Cretans fashioned.” \par Adonis, or “Thammuz yearly wounded,” was said to live and die in alternate years.}, and falls\footnote{By the word “falls” (\textgreek{ἀποπτώσεις}) Cyril evidently refers to the story of Hephæstus, or Vulcan, to which Milton alludes (\emph{Paradise Lost}, I. 740):— \par “Men call’d him Mulciber, and how he fell \par From heaven they fabled, thrown by angry Jove \par Sheer o’er the crystal battlements: from morn \par To noon he fell, from noon to dewy eve, \par A summer’s day.”}, and thunder-strokes\footnote{The “thunder-strokes” refer to “Titan heaven’s first-born, With his enormous brood” (\emph{Par. Lost}, I. 510). Cf. Virgil, \emph{Æn}. vi. 580:— \par “Hic genus antiquum Terræ, Titania pubes, \par Fulmine dejecti fundo volvuntur in imo.” \par Ibid. \emph{v}. 585:— \par “Vidi et crudeles dantem Salmonea pœnas, \par Dum flammas Jovis et sonitus imitatur Olympi.” \par Clem. Alex. (\emph{Exhort}. II. § 37):—“Æsculapius lies struck with lightning in the regions of Cynosuris.” Cf. Virg. \emph{Æn}. vii. 770 \emph{ss}.} of their gods. Seest thou from how great a height and how low they have fallen? Was it without reason then that the Son of God came down from heaven? or was it that He might heal so great a wound? Was it without reason that the Son came? or was it in order that the Father might be acknowledged? Thou hast learned what moved the Only-begotten to come down from the throne at God’s right hand. The Father was despised, the Son must needs correct the error: for He \textsc{Through} \textsc{Whom} \textsc{All} \textsc{Things} \textsc{Were} \textsc{Made} must bring them all as offerings to the Lord of all. The wound must be healed: for what could be worse than this disease, that a stone should be worshipped instead of God?

Of Heresies.

12. And not among the heathen only did the devil make these assaults; for many of those who are falsely called Christians, and wrongfully addressed by the sweet name of Christ, have ere now impiously dared to banish God from His own creation. I mean the brood of heretics, those most ungodly men of evil name, pretending to be friends of Christ but utterly hating Him. For he who blasphemes the Father of the Christ is an enemy of the Son. These men have dared to speak of two Godheads, one good and one evil\footnote{The theory of two Gods, one good and the other evil, was held by Cerdo, and Marcion (Hippolytus, \emph{Refut. omnium Hær.} VII. cap. 17: Irenæus, III. xxv. 3, quoted in note on Cat. iv. 4). The Manichees also held that the Creator of the world was distinct from the Supreme God (Alexander Lycop. \emph{de Manichæorum Sententiis}, cap. iii.).}! O monstrous blindness! If a Godhead, then assuredly good. But if not good, why called a Godhead? For if goodness is an attribute of God; if loving-kindness, beneficence, almighty power, are proper to God, then of two things one, either in calling Him God let the name and operation be united; or if they would rob Him of His operations, let them not give Him the bare name.

13. Heretics have dared to say that there are two Gods, and of good and evil two sources, and these unbegotten. If both are unbegotten it is certain that they are also equal, and both mighty. How then doth the light destroy the darkness? And do they ever exist together, or are they separated? Together they cannot be; \emph{for what fellowship hath light with darkness?} saith the Apostle\footnote{2 Cor. vi. 14. Cyril’s description applies especially to the heresy of Manes. See § 36, note 3, at the end of this Lecture; also Cat. xi. 21. and Cat. xv. 3.}. But if they are far from each other, it is certain that they hold also each his own place; and if they hold their own separate places, we are certainly in the realm of one God, and certainly worship one God. For thus we must conclude, even if we assent to their folly, that we must worship one God. Let us examine also what they say of the good God. Hath He power or no power? If He hath power, how did evil arise against His will? And how doth the evil substance intrude, if He be not willing? For if He knows but cannot hinder it, they charge Him with want of power; but if He has the power, yet hinders not, they accuse Him of treachery. Mark too their want of sense. At one time they say that the Evil One hath no communion with the good God in the creation of the world; but at another time they say that he hath the fourth part only. Also they say that the good God is the Father of Christ; but Christ they call this sun. If, therefore according to them, the world was made by the Evil One, and the sun is in the world, how is the Son of the Good an unwilling slave in the kingdom of the Evil? We bemire ourselves in speaking of these things, but we do it lest any of those present should from ignorance fall into the mire of the heretics. I know that I have defiled my own mouth and the ears of my listeners: yet it is expedient. For it is much better to hear absurdities charged against others, than to fall into them from ignorance: far better that thou know the mire and hate it, than unawares fall into it. For the godless system of the heresies is a road with many branches, and whenever a man has strayed from the one straight way, then he falls down precipices again and again.

14. The inventor of all heresy was Simon Magus\footnote{So Irenæus (I. xxiii. 2) says that “from this Simon of Samaria all kinds of heresies derive their origin.”}: that Simon, who in the Acts of the Apostles thought to purchase with money the unsaleable grace of the Spirit, and heard the words, \emph{Thou hast neither part nor lot in this matter}\footnote{Acts viii. 18–21.}, and the rest: concerning whom also it is written, \emph{They went out from us, but they were not of us; for if they had been of us, they would have remained with us}\footnote{1 John ii. 19.}. This man, after he had been cast out by the Apostles, came to Rome, and gaining over one Helena a harlot\footnote{Irenæus (I. xxiii. 2): “Having purchased from Tyre, a city of Phœnicia, a certain harlot named Helena, he used to carry her about with him, declaring that this woman was the first conception of his mind, the mother of all, by whom in the beginning he conceived in his mind the creation of Angels and Archangels.”}, was the first that dared with blasphemous mouth to say that it was himself who appeared on Mount Sinai as the Father, and afterwards appeared among the Jews, not in real flesh but in seeming\footnote{Cf. Epiphan. (\emph{Hæres}. p. 55, B): “He said that he was the Son and had not really suffered, but only in appearance (\textgreek{δοκήσει}).”}, as Christ Jesus, and afterwards as the Holy Spirit whom Christ promised to send as the Paraclete\footnote{Irenæus (I. xxiii. 1): “He taught that it was himself who appeared among the Jews as the Son, and descended in Samaria as the Father, but came to other nations as the Holy Spirit.” \par Cyril here departs from his authority by substituting Mount Sinai for Samaria, and thereby falls into error. Simon had first appeared in Samaria, being a native of Gitton: moreover in claiming to be the Father he meant to set himself far above the inferior Deity who had given the Law on Sinai, saying that he was “the highest of all Powers, that is the Father who is over all.”}. And he so deceived the City of Rome that Claudius set up his statue, and wrote beneath it, in the language of the Romans, “Simoni Deo Sancto,” which being interpreted signifies, “To Simon the Holy God\footnote{“Justin Martyr in his first Apology, addressed to Antoninus Pius, writes thus (c. 26): ‘There was one Simon a Samaritan, of the village called Gitton, who in the reign of Claudius Cæsar, and in your royal city of Rome, did mighty feats of magic by the art of dæmons working in him. He was considered a god, and as a god was honoured among you with a statue, which statue was set up in the river Tiber between the two bridges, and bears this inscription in Latin: \par \emph{Simoni Deo Sancto}; \par which is, \par To Simon the holy God. \par “The substance of this story is repeated by Irenæus (\emph{adv. Hær}. I. xxiii. 1), and by Tertullian (\emph{Apol}. c. 13), who reproaches the Romans for installing Simon Magus in their Pantheon, and giving him a statue and the title ‘Holy God.’ \par “In \textsc{a.d.} 1574, a stone, which had formed the base of a statue, was dug up on the site described by Justin, the Island in the Tiber, bearing an inscription—‘Semoni Sanco Deo Fidio Sacrum, \&c.’ Hence it has been supposed that Justin mistook a statue of the Sabine God, ‘Semo Sancus,’ for one of Simon Magus. See the notes in Otto’s Justin Martyr, and Stieren’s Irenæus. \par “On the other hand Tillemont (\emph{Memoires}, t. ii. p. 482) maintains that Justin in an Apology addressed to the emperor and written in Rome itself cannot reasonably be supposed to have fallen into so manifest an error. Whichever view we take of Justin’s accuracy concerning the inscription and the statue, there is nothing improbable in his statement that Simon Magus was at Rome in the reign of Claudius.” (Extracted by permission from the Speaker’s Commentary, \emph{Introduction to the Epistle to the Romans}, p. 4.)}.”

15. As the delusion was extending, Peter and Paul, a noble pair, chief rulers of the Church, arrived and set the error right\footnote{“Justin says not one word about St. Peter’s alleged visit to Rome, and his encounter with Simon Magus.” But “Eusebius in his \emph{Ecclesiastical History} (c. \textsc{a.d.} 325), quotes Justin Martyr’s story about Simon Magus (\emph{E. H.} ii. c. 13), and then, without referring to any authority, goes on to assert (c. 14) that ‘immediately in the same reign of Claudius divine Providence led Peter the great Apostle to Rome to encounter this great destroyer of life,’ and that he thus brought the light of the Gospel from the East to the West’ (\emph{ibidem}). \par Eusebius probably borrowed this story “from the strange fictions of the \emph{Clementine Recognitions} and \emph{Homilies}, and \emph{Apostolic Constitutions}.” See \emph{Recogn}. III. 63–65; \emph{Hom}. I. 15, III. 58; \emph{Apost. Constit}. VI. 7, 8, 9. Cyril’s account of Simon’s death is taken from the same untrustworthy sources.}; and when the supposed god Simon wished to shew himself off, they straightway shewed him as a corpse. For Simon promised to rise aloft to heaven, and came riding in a dæmons’ chariot on the air; but the servants of God fell on their knees, and having shewn that agreement of which Jesus spake, that \emph{If two of you shall agree concerning anything that they shall ask, it shall be done unto them}\footnote{Matt. xviii. 19.}, they launched the weapon of their concord in prayer against Magus, and struck him down to the earth. And marvellous though it was, yet no marvel. For Peter was there, who carrieth the keys of heaven\footnote{Ib. xvi. 19.}: and nothing wonderful, for Paul was there\footnote{It is certain that S. Paul was not at Rome at this time. This story of Simon Magus and his ‘fiery car’ is told, with variations, by Arnobius (\emph{adv. Gentes}, II. 12), and in \emph{Apost. Constit.} VI. 9.}, who was \emph{caught up to the third heaven}, and \emph{into Paradise, and heard unspeakable words, which it is not lawful far a man to utter}\footnote{2 Cor. xii. 2, 4.}. These brought the supposed God down from the sky to earth, thence to be taken down to the regions below the earth. In this man first the serpent of wickedness appeared; but when one head had been cut off, the root of wickedness was found again with many heads.

16. For Cerinthus\footnote{Cerinthus taught that the world was not made by the supreme God, but by a separate Power ignorant of Him. See Irenæus, \emph{Hær}. I. xxvi., Euseb. \emph{E.H.} iii. 28, with the notes in this Series.} \emph{made havoc of the Church}, and Menander\footnote{Menander is first mentioned by Justin M. (\emph{Apolog}. I. cap. 26): “Menander, also a Samaritan, of the town Capparetæa, a disciple of Simon, and inspired by devils, we know to have deceived many while he was in Antioch by his magical art. He persuaded those who adhered to him that they should never die.” Irenæus (I. xxiii. 5) adds that Menander announced himself as the Saviour sent by the Invisibles, and taught that the world was created by Angels. See also Tertullian (\emph{de Animâ}, cap. 50.)}, and Carpocrates\footnote{Carpocrates, a Platonic philosopher, who taught at Alexandria (125 \textsc{a.d.} \emph{circ}.), held that the world and all things in it were made by Angels far inferior to the unbegotten (unknown) Father (Iren. I. xxv. 1; Tertullian, \emph{Adv. Hær.} cap. 3).}, Ebionites\footnote{Irenæus, I. 26: “Those who are called Ebionites agree that the world was made by God; but their opinions with respect to the Lord are like those of Cerinthus and Carpocrates.”} also, and Marcion\footnote{On Marcion, see note 5, on Cat. iv. 4.}, that mouthpiece of ungodliness. For he who proclaimed different gods, one the Good, the other the Just, contradicts the Son when He says, \emph{O righteous Father}\footnote{John xvii. 25.}. And he who says again that the Father is one, and the maker of the world another, opposes the Son when He says, \emph{If then God so clothes the grass of the field which to-day is, and to-morrow is cast into the furnace of fire}\footnote{Luke xii. 28.}; and, \emph{Who maketh His sun to rise on the evil and on the good, and sendeth rain on the just and on the unjust}\footnote{Matt. v. 45.}. Here again is a second inventor of more mischief, this Marcion. For being confuted by the testimonies from the Old Testament which are quoted in the New, he was the first who dared to cut those testimonies out\footnote{Marcion accepted only St. Luke’s Gospel, and mutilated that (Tertullian, \emph{Adv. Marcion}. iv. 2). He thus got rid of the testimony of the Apostles and eye-witnesses, Matthew and John, and represented the Law and the Gospel as contradictory revelations of two different Gods. For this Cyril calls him ‘a second inventor of mischief,’ Simon Magus (§ 14) being the first.}, and leave the preaching of the word of faith without witness, thus effacing the true God: and sought to undermine the Church’s faith, as if there were no heralds of it.

17. He again was succeeded by another, Basilides, of evil name, and dangerous character, a preacher of impurities\footnote{Basilides was earlier than Marcion, being the founder of a Gnostic sect at Alexandria in the reign of Hadrian (\textsc{a.d.} 117–138). His doctrines are described by Irenæus (I. xxvii. 3–7), and very fully by Hippolytus (\emph{Refut. omn. Hær}. VII. 2–15). The charge of teaching licentiousness attaches rather to the later followers of Basilides than to himself or his son Isidorus (Clem. Alex. \emph{Stromat}. III. cap. 1). Basilides wrote a Commentary on the Gospel in 24 books (\emph{Exegetica}), of which the 23rd is quoted by Clement of Alexandria (\emph{Stromat}. IV. cap. 12), and against which Agrippa Castor wrote a refutation. Origen (\emph{Hom}. I. \emph{in Lucam}.) says that Basilides wrote a Gospel bearing his own name. See Routh, \emph{Rell. Sacr}. I. p. 85; V. p. 106: Westcott, \emph{History of Canon of N.T.} iv. § 3.}. The contest of wickedness was aided also by Valentinus\footnote{“The doctrines of Valentinus are described fully by Irenæus (I. cap. i.) from whom S. Cyril takes this account. Valentinus, and Basilides, and Bardesanes, and Harmonious, and those of their company admit Christ’s conception and birth of the Virgin, but say that God the Word received no addition from the Virgin, but made a sort of passage through her, as through a tube, and made use of a phantom in appearing to men.” (Theodoret, \emph{Epist}. 145.)}, a preacher of thirty gods. The Greeks tell of but few: and the man who was called—but more truly was not—a Christian extended the delusion to full thirty. He says, too, that Bythus the Abyss (for it became him as being an abyss of wickedness to begin his teaching from the Abyss) begat Silence, and of Silence begat the Word. This Bythus was worse than the Zeus of the Greeks, who was united to his sister: for Silence was said to be the child of Bythus. Dost thou see the absurdity invested with a show of Christianity? Wait a little, and thou wilt be shocked at his impiety; for he asserts that of this Bythus were begotten eight Æons; and of them, ten; and of them, other twelve, male and female. But whence is the proof of these things? See their silliness from their fabrications. Whence hast thou the proof of the thirty Æons? Because, saith he, it is written, that \emph{Jesus was baptized}, \emph{being thirty years old}\footnote{Luke iii. 23.}. But even if He was baptized when thirty years old, what sort of demonstration is this from the thirty years? Are there then five gods, because He brake five loaves among five thousand? Or because he had twelve Disciples, must there also be twelve gods?

18. And even this is still little compared with the impieties which follow. For the last of the deities being, as he dares to speak, both male and female, this, he says, is Wisdom\footnote{Irenæus I. ii. 2.}. What impiety! For \emph{the Wisdom of God}\footnote{1 Cor. i. 24.} is Christ His Only-begotten Son: and he by his doctrine degraded the Wisdom of God into a female element, and one of thirty, and the last fabrication. He also says that Wisdom attempted to behold the first God, and not bearing His brightness fell from heaven, and was cast out of her thirtieth place. Then she groaned, and of her groans begat the Devil\footnote{Irenæus, l. c., and Hippolytus, who gives an elaborate account of the doctrines of Valentinus (L. VI. capp. xvi.–xxxii.), both represent Sophia, “Wisdom,” as giving birth not to Satan, but to a shapeless abortion, which was the origin of matter. According to Irenæus (I. iv. 2), Achamoth, the enthymesis of Sophia, gave birth to the Demiurge, and “from her tears all that is of a liquid nature was formed.” \par In Tertullian’s Treatise \emph{against the Valentinians} chap. xxii., Achamoth is said as by Cyril to have given birth to Satan: but in chap. xxiii. Satan seems to be identified (or interchanged) with the Demiurge.}, and as she wept over her fall made of her tears the sea. Mark the impiety. For of Wisdom how is the Devil begotten, and of prudence wickedness, or of light darkness? He says too that the Devil begat others, some of whom created the world: and that the Christ came down in order to make mankind revolt from the Maker of the world.

19. But hear whom they say Christ Jesus to be, that thou mayest detest them yet more. For they say that after Wisdom had been cast down, in order that the number of the thirty might not be incomplete, the nine and twenty Æons contributed each a little part, and formed the Christ\footnote{The account in Irenæus (I. ii. 6) is rather different: “The whole Pleroma of the Æons, with one design and desire, and with the concurrence of the Christ and the Holy Spirit, their Father also setting the seal of His approval on their conduct, brought together whatever each one had in himself of the greatest beauty and preciousness; and uniting all these contributions so as skilfully to blend the whole, they produced, to the honour and glory of Bythus, a being of most perfect beauty, the very star of the Pleroma, and its perfect fruit, namely Jesus.” \par Tertullian, \emph{Against the Valentinians}, chap. 12, gives a sarcastic description of this strange doctrine, deriving his facts (chap. 5) from Justin, Miltiades, “Irenæus, that very exact inquirer into all doctrines,” and Proculus.}: and they say that He also is both male and female\footnote{This statement does not agree with Irenæus (I. vii. 1), who says that the Valentinians represented the Saviour, that is Jesus, as becoming the bridegroom of Achamoth or Sophia.}. Can anything be more impious than this? Anything more wretched? I am describing their delusion to thee, in order that thou mayest hate them the more. Shun, therefore, their impiety, and \emph{do not even give greeting to}\footnote{2 John 10, 11: “Neither bid him God speed” (\textsc{A.V}.): “give him no greeting” (\textsc{R.V.}).} a man of this kind, lest thou have \emph{fellowship with the unfruitful works of darkness}\footnote{Ephes. v. 11.}: neither make curious inquiries, nor be willing to enter into conversation with them.

20. Hate all heretics, but especially him who is rightly named after mania\footnote{Eusebius in his brief notice of the Manichean heresy (\emph{Hist. Eccles}. vii. 31) plays, like S. Cyril, upon the name Manes as well suited to a madman.}, who arose not long ago in the reign of Probus\footnote{Marcus Aurelius Probus, Emperor \textsc{a.d.} 276–282, from being an obscure Illyrian soldier came to be universally esteemed the best and noblest of the Roman Emperors.}. For the delusion began full seventy years ago\footnote{Routh (\emph{\textsc{R.S}}\textsc{. V.} p. 12) comes to the conclusion that the famous disputation between Manes and Archelaus took place between July and December, \textsc{a.d.} 277. Accordingly these Lectures, being “full 70 years” later, could not have been delivered before the Spring of \textsc{a.d.} 348.}, and there are men still living who saw him with their very eyes. But hate him not for this, that he lived a short time ago; but because of his impious doctrines hate thou the worker of wickedness, the receptacle of all filth, who gathered up the mire of every heresy\footnote{Leo the Great (\emph{Serm}. xv. cap. 4) speaks of the madness of the later Manichees as including all errors and impieties: “all profanity of Paganism, all blindness of the carnal Jews, the illicit secrets of the magic art, the sacrilege and blasphemy of all heresies, flowed together in that sect as into a sort of cess-pool of all filth.” Leo summoned those whom they called the “elect,” both men and women, before an assembly of Bishops and Presbyters, and obtained from these witnesses a full account of the execrable practices of the sect, in which, as he declares, “their law is lying, their religion the devil, their sacrifice obscenity.”}. For aspiring to become pre-eminent among wicked men, he took the doctrines of all, and having combined them into one heresy filled with blasphemies and all iniquity, he makes havoc of the Church, or rather of those outside the Church, roaming about like a lion and devouring. Heed not their fair speech, nor their supposed humility: for they are serpents, \emph{a generation of vipers}\footnote{Matt. iii. 7.}. Judas too \emph{said Hail! Master}\footnote{Ib. xxvi. 49.}, even while he was betraying Him. Heed not their kisses, but beware of their venom.

21. Now, lest I seem to accuse him without reason, let me make a digression to tell who this Manes is, and in part what he teaches: for all time would fail to describe adequately the whole of his foul teaching. But \emph{for help in time of need}\footnote{Heb. iv. 16.}, store up in thy memory what I have said to former hearers, and will repeat to those now present, that they who know not may learn, and they who know may be reminded. Manes is not of Christian origin, God forbid! nor was he like Simon cast out of the Church, neither himself nor the teachers who were before him. For he steals other men’s wickedness, and makes their wickedness his own: but how and in what manner thou must hear.

22. There was in Egypt one Scythianus\footnote{Cyril takes his account of Manes from the “\emph{Acta Archelai et Manetis Disputationis,}” of which Routh has edited the Latin translation together with the Fragments of the Greek preserved by Cyril in this Lecture and by Epiphanius. There is an English translation of the whole in Clark’s “Ante-Nicene Christian Library.”}, a Saracen\footnote{The Saracens are mentioned by both Pliny and Ptolemy. See \emph{Dict. of Greek and Roman Geography.}} by birth, having nothing in common either with Judaism or with Christianity. This man, who dwelt at Alexandria and imitated the life of Aristotle\footnote{There is no mention of Aristotle in the \emph{Acta Archelai}, but Scythianus is stated (cap. li.) to have founded the sect in the time of the Apostles, and to have derived his duality of Gods from Pythagoras, and to have learned the wisdom of the Egyptians.}, composed four books\footnote{These four books are stated by Archelaus (\emph{Acta}, cap. lii.), to have been written for Manes by his disciple Terebinthus.}, one called a Gospel which had not the acts of Christ, but the mere name only, and one other called the book of Chapters, and a third of Mysteries, and a fourth, which they circulate now, the Treasure\footnote{In allusion to this name the history of the Disputation is called (\emph{Acta}, cap. i.) “The true Treasure.”}. This man had a disciple, Terebinthus by name. But when Scythianus purposed to come into Judæa, and make havoc of the land, the Lord smote him with a deadly disease, and stayed the pestilence\footnote{The true reading of this sentence, \textgreek{προαιρούμενον τὸν Σκυθιανόν}, instead of \textgreek{τὸν πρόειρῃμένον Σκ}., has been restored by Cleopas from the \textsc{ms.} in the Archiepiscopal library at Jerusalem. This reading agrees with the statement in \emph{Acta Archel}. cap. li.: “Scythianus thought of making an excursion into Judæa, with the purpose of meeting all those who had a reputation there as teachers; but it came to pass that he suddenly departed this life, without having been able to make any progress.”}.

23. But Terebinthus, his disciple in this wicked error, inherited his money and books and heresy\footnote{This statement agrees with the reading of the Vatican \textsc{ms.} of the \emph{Acta Archelai}, “omnibus quæcunque ejus fuerunt congregratis.”}, and came to Palestine, and becoming known and condemned in Judæa\footnote{In the \emph{Acta} there is no mention of Palestine, but only that he “set out for Babylonia, a province which is now held by the Persians.”} he resolved to pass into Persia: but lest he should be recognised there also by his name he changed it and called himself Buddas\footnote{Clem. Alex. (\emph{Strom}. i. 15): “Some also of the Indians obey the precepts of Boutta, and honour him as a god for his extraordinary sanctity.”}. However, he found adversaries there also in the priests of Mithras\footnote{Cf. \emph{Acta Arch}. cap. lii.: “A certain Parcus, however, a prophet, and Labdacus, son of Mithras, charged him with falsehood.” On the name Parcus and Labdacus, see \emph{Dict. Chr. Biogr.}, “Barcabbas,” and on the Magian worship of the Sun-god Mithras, see Rawlinson (\emph{Herodot}. Vol. I. p. 426).}: and being confuted in the discussion of many arguments and controversies, and at last hard pressed, he took refuge with a certain widow. Then having gone up on the housetop, and summoned the dæmons of the air, whom the Manichees to this day invoke over their abominable ceremony of the fig\footnote{See below, § 33.}, he was smitten of God, and cast down from the housetop, and expired: and so the second beast was cut off.

24. The books, however, which were the records of his impiety, remained; and both these and his money the widow inherited. And having neither kinsman nor any other friend, she determined to buy with the money a boy named Cubricus\footnote{Cf. \emph{Acta Arch}. cap. liii. “A boy about seven years old, named Corbicius.”}: him she adopted and educated as a son in the learning of the Persians, and thus sharpened an evil weapon against mankind. So Cubricus, the vile slave, grew up in the midst of philosophers, and on the death of the widow inherited both the books and the money. Then, lest the name of slavery might be a reproach, instead of Cubricus he called himself Manes, which in the language of the Persians signifies discourse\footnote{See a different account in \emph{Dict. Chr. Biogr.}, “Manes.”}. For as he thought himself something of a disputant, he surnamed himself Manes, as it were an excellent master of discourse. But though he contrived for himself an honourable title according to the language of the Persians, yet the providence of God caused him to become a self-accuser even against his will, that through thinking to honour himself in Persia, he might proclaim himself among the Greeks by name a maniac.

25. He dared too to say that he was the Paraclete, though it is written, \emph{But whosoever shall blaspheme against the Holy Ghost, hath no forgiveness}\footnote{Mark iii. 29.}. He committed blasphemy therefore by saying that he was the Holy Ghost: let him that communicates with those heretics see with whom he is enrolling himself. The slave shook the world, since \emph{by three things the earth is shaken, and the fourth it cannot bear,—if a slave became a king}\footnote{Prov. xxx. 21, 22.}. Having come into public he now began to promise things above man’s power. The son of the King of the Persians was sick, and a multitude of physicians were in attendance: but Manes promised, as if he were a godly man, to cure him by prayer. With the departure of the physicians, the life of the child departed: and the man’s impiety was detected. So the would-be philosopher was a prisoner, being cast into prison not for reproving the king in the cause of truth, not for destroying the idols, but for promising to save and lying, or rather, if the truth must be told, for committing murder. For the child who might have been saved by medical treatment, was murdered by this man’s driving away the physicians, and killing him by want of treatment.

26. Now as there are very many wicked things which I tell thee of him, remember first his blasphemy, secondly his slavery (not that slavery is a disgrace, but that his pretending to be free-born, when he was a slave, was wicked), thirdly, the falsehood of his promise, fourthly, the murder of the child, and fifthly, the disgrace of the imprisonment. And there was not only the disgrace of the prison, but also the flight from prison. For he who called himself the Paraclete and champion of the truth, ran away: he was no successor of Jesus, who readily went to the Cross, but this man was the reverse, a runaway. Moreover, the King of the Persians ordered the keepers of the prison to be executed: so Manes was the cause of the child’s death through his vain boasting, and of the gaolers’ death through his flight. Ought then he, who shared the guilt of murder, to be worshipped? Ought he not to have followed the example of Jesus, and said, \emph{If ye seek Me, let these go their way}\footnote{John xviii. 8.}? Ought he not to have said, like Jonas, \emph{Take me, and cast me into the sea: for this storm is because of me}\footnote{Jonah i. 12.}?

27. He escapes from the prison, and comes into Mesopotamia: but there Bishop Archelaus, a shield of righteousness, encounters him\footnote{The account of the discussion in this and the two following chapters is not now found in the Latin Version of the “Disputation,” but is regarded by Dr. Routh as having been derived by Cyril from some different copies of the Greek. The last paragraph of § 29, “These mysteries, \&c.,” is evidently a caution addressed to the hearers by Cyril himself (Routh, \emph{Rell. Sac}. V. 199).}: and having accused him before philosophers as judges, and having assembled an audience of Gentiles, lest if Christians gave judgment, the judges might be thought to shew favour,—Tell us what thou preachest, said Archelaus to Manes. And he, whose \emph{mouth was as an open sepulchre}\footnote{Ps. v. 9.}, began first with blasphemy against the Maker of all things, saying, The God of the Old Testament is the author of evils, as He says of Himself, \emph{I am a consuming fire}\footnote{Deut. iv. 24.}. But the wise Archelaus undermined his blasphemous argument by saying, “If the God of the Old Testament, as thou sayest, calls Himself a fire, whose Son is He who saith, \emph{I came to send fire on the earth}\footnote{Luke xii. 49.}? If thou findest fault with Him who saith, The Lord killeth, and maketh alive\footnote{1 Sam. ii. 6.}, why dost thou honour Peter, who raised up Tabitha, but struck Sapphira dead? If again thou findest fault, because He prepared fire, wherefore dost thou not find fault with Him who saith, \emph{Depart from Me into everlasting fire}\footnote{Matt. xxv. 41.}? If thou findest fault with Him who saith, \emph{I am God that make peace, and create evil}\footnote{Is. xlv. 7.}, explain how Jesus saith, I came not to send peace but a sword\footnote{Matt. x. 34.}. Since both speak alike, of two things one, either both are good, because of their agreement, or if Jesus is blameless in so speaking. why blamest thou Him that saith the like in the Old Testament?”

28. Then Manes answers him: “And what sort of God causes blindness? For it is Paul who saith, \emph{In whom the God of this world hath blinded the minds of them that believe not, lest the light of the Gospel should shine unto them}\footnote{2 Cor. iv. 4, \textgreek{νοήματα}, “thoughts.”}.” But Archelaus made a good retort, saying, “Read a little before: \emph{But if our Gospel is veiled, it is veiled in them that are perishing}\footnote{2 Cor. iv. 3.}. Seest thou that in them that are perishing it is veiled? For it is not right \emph{to give the things which are holy unto the dogs}\footnote{Matt. vii. 6.}. Again, Is it only the God of the Old Testament that hath blinded the minds of them that believe not? Hath not Jesus Himself said, \emph{For this cause speak I unto them in parables, that seeing they may not see}\footnote{Matt. xiii. 13. Both A.V. and R.V. follow the better reading: “because seeing they see not, \&c.”}? Was it from hating them that He wished them not to see? Or because of their unworthiness, since \emph{their eyes they had closed}\footnote{Matt. xiii. 15.}. For where there is wilful wickedness, there is also a withholding of grace: \emph{for to him that hath shall be given; but from him that hath not shall be taken even that which he seemeth to have}\footnote{Ib. xxv. 29; Luke viii. 18.}.

29. “But if some are right in their interpretation, we must say as follows\footnote{Instead of the reading of the Benedictine and earlier editions, \textgreek{εἰ δὲ δεῖ καὶ ὥς τινες ἐξηγοῦνται τοῦτο εἰπεῖν}, the \textsc{mss.} Roe and Casaubon combine \textgreek{δει και ως} into the one word \textgreek{δικαιως}, which is probably the right reading. Something, however, is still wanted to complete the construction, and Petrus Siculus (\emph{circ.} \textsc{a.d.} 870) who quotes the passage in his \emph{History of the Manichees}, boldly conjectures \textgreek{ἔστι καὶ οὕτως εἰπεῖν}. A simpler emendation would be—\textgreek{εἰ δὲ δικαίως τινὲς ἐξηγοῦνται, δεῖ τουτο εἰπεῖν}—which both completes the construction and explains the reading \textgreek{δεῖ καὶ ὡς}.} (for it is no unworthy expression)—If indeed He blinded the thoughts of them that believe not he blinded them for a good purpose, that they might look with new sight on what is good. For he said not, He blinded their soul, but, \emph{the thoughts of them that believe not}\footnote{\textgreek{νοήματα}, 2 Cor. iv. 4.}. And the meaning is something of this kind: ‘Blind the lewd thoughts of the lewd, and the man is saved: blind the grasping and rapacious thought of the robber, and the man is saved.’ But wilt thou not understand it thus? Then there is yet another interpretation. The sun also blinds those whose sight is dim: and they whose eyes are diseased are hurt by the light and blinded. Not that the sun’s nature is to blind, but that the substance of the eyes is incapable of seeing. In like manner unbelievers being diseased in their heart cannot look upon the radiance of the Godhead. Nor hath he said, \emph{‘He hath blinded their thoughts}, that they should not hear the Gospel:’ but, \emph{that the light of the glory of the Gospel of our Lord Jesus Christ should not shine unto them}. For to hear the Gospel is permitted to all: but the glory of the Gospel is reserved for Christ’s true children only. Therefore the Lord spoke in parables to those who could not hear\footnote{Matt. xiii. 13.}: but to the Disciples he explained the parables in private\footnote{Mark iv. 34.}: for the brightness of the glory is for those who have been enlightened, the blinding for them that believe not.” These mysteries, which the Church now explains to thee who art passing out of the class of Catechumens, it is not the custom to explain to heathen. For to a heathen we do not explain the mysteries concerning Father, Son, and Holy Ghost, nor before Catechumens do we speak plainly of the mysteries: but many things we often speak in a veiled way, that the believers who know may understand, and they who know not may get no hurt\footnote{See the note at the end of Procatechesis.}.

30. By such and many other arguments the serpent was overthrown: thus did Archelaus wrestle with Manes and threw him. Again, he who had fled from prison flees from this place also: and having run away from his antagonist, he comes to a very poor village, like the serpent in Paradise when he left Adam and came to Eve. But the good shepherd Archelaus taking forethought for his sheep, when he heard of his flight, straightway hastened with all speed in search of the wolf. And when Manes suddenly saw his adversary, he rushed out and fled: it was however his last flight. For the officers of the King of Persia searched everywhere, and caught the fugitive: and the sentence, which he ought to have received in the presence of Archelaus, is passed upon him by the king’s officers. This Manes, whom his own disciples worship, is arrested and brought before the king. The king reproached him with his falsehood and his flight: poured scorn upon his slavish condition, avenged the murder of his child, and condemned him also for the murder of the gaolers: he commands him to be flayed after the Persian fashion. And while the rest of his body was given over for food of wild beasts, his skin, the receptacle of his vile mind, was hung up before the gates like a sack\footnote{Disput. § 55. Compare the account of Manes in Socrates, Eccles. Hist. I. 22, in this series.}. He that called himself the Paraclete and professed to know the future, knew not his own flight and capture.

31. This man has had three disciples, Thomas, and Baddas, and Hermas. Let none read the Gospel according to Thomas\footnote{The Gospel of Thomas, an account of the Childhood of Jesus, is extant in three forms, two in Greek and one in Latin: these are all translated in Clark’s Ante-Nicene Library. The work is wrongly attributed by Cyril to a disciple of Manes, being mentioned long before Hippolytus (\emph{Refutation of all Heresies}, V. 2) and by Origen (\emph{Hom. I. in Lucam}): “There is extant also the Gospel according to Thomas.”}: for it is the work not of one of the twelve Apostles, but of one of the three wicked disciples of Manes. Let none associate with the soul-destroying Manicheans, who by decoctions of chaff counterfeit the sad look of fasting, who speak evil of the Creator of meats, and greedily devour the daintiest, who teach that the man who plucks up this or that herb is changed into it. For if he who crops herbs or any vegetable is changed into the same, into how many will husbandmen and the tribe of gardeners be changed\footnote{In the Disputation, § 9, Turbo describes these transformations: “Reapers must be transformed into hay, or beans, or barley, or corn, or vegetables, that they may be reaped and cut. Again if any one eats bread, he must become bread, and be eaten. If one kills a chicken, he will be a chicken himself. If one kills a mouse, he also will be a mouse.”}? The gardener, as we see, has used his sickle against so many: into which then is he changed? Verily their doctrines are ridiculous, and fraught with their own condemnation and shame! The same man, being the shepherd of a flock, both sacrifices a sheep and kills a wolf. Into what then is he changed? Many men both net fishes and lime birds: into which then are they transformed?

32. Let those children of sloth, the Manicheans, make answer; who without labouring themselves eat up the labourers’ fruits: who welcome with smiling faces those who bring them their food, and return curses instead of blessings. For when a simple person brings them anything, “Stand outside a while,” saith he, “and I will bless thee.” Then having taken the bread into his hands (as those who have repented and left them have confessed), “I did not make thee,” says the Manichee to the bread: and sends up curses against the Most High; and curses him that made it, and so eats what was made\footnote{See Turbo’s confession, Disput. § 9: “And when they are going to eat bread, they first pray, speaking thus to the bread: ‘I neither reaped thee, nor ground thee, nor kneaded thee, nor cast thee into the oven: but another did these things and brought thee to me, and I am not to blame for eating thee.’ And when he has said this to himself, he says to the Catechumen, ‘I have prayed for thee,’ and so he goes away.”}. If thou hatest the food, why didst thou look with smiling countenance on him that brought it to thee? If thou art thankful to the bringer, why dost thou utter thy blasphemy to God, who created and made it? So again he says, “I sowed thee not: may he be sown who sowed thee! I reaped thee not with a sickle: may he be reaped who reaped thee! I baked thee not with fire: may he be baked who baked thee!” A fine return for the kindness!

33. These are great faults, but still small in comparison with the rest. Their Baptism I dare not describe before men and women\footnote{On the rites of Baptism and Eucharist employed by the Manichees, see Dict. Chr. Biogr., \emph{Manicheans.}}. I dare not say what they distribute to their wretched communicants\footnote{The original runs: \textgreek{Οὐ τολμῶ εἰπεῖν, τίνι ἐμβάπτοντες τὴν ἰσχάδα, διδόασι τοῖς ἀθλίοις. διὰ συσσήμων δὲ μόνον δηλούσθω. ἄνδρες γὰρ τὰ ἐν τοῖς ἐνυπνιασμοῖς ἐνθυμείσθωσιν, καὶ γυναῖκες τὰ ἐν ἀφέδροις. Μιαίνομεν ἀληθας τὸ στόμα κ.τ.λ}.}….Truly we pollute our mouth in speaking of these things. Are the heathen more detestable than these? Are the Samaritans more wretched? Are Jews more impious? Are fornicators more impure\footnote{\textgreek{῾Ο μὲν γὰρ πορνεύσας, πρὸς μίαν ὥραν δ ἐπιθυμίαν τελεῖ τὴν πρᾶξιν· καταγινώσκων δὲ τῆς πράξεως ὡς μιανθεὶς οἶδε λουτροῦ ἐπιδεόμενος, καὶ γινώσκει τῆς πρὰξεως τὸ μυσαρόν. ῾Ο δὲ Μανιχαῖος θυσιαστηρίου μέσον, οὗ νομίζει, τίθησι ταῦτα, καὶ μιαίνει καὶ τὸ στόμα καὶ τὴν γλῶτταν. παρὰ τοιούτου στόματος, ἄνθρωπε κ.τ.λ}.}? But the Manichee sets these offerings in the midst of the altar as he considers it\footnote{\textgreek{οὗ νομίζει}. The Manichees boasted of their superiority to the Pagans in not worshipping God with altars, temples, images, victims, or incense (August. \emph{contra Faustum} XX. cap. 15). Yet they used the names, as Augustine affirms (\emph{l. c.} cap. 18): “Nevertheless I wish you would tell me why you call all those things which you approve in your own case by these names, temple, altar, sacrifice.”}. And dost thou, O man, receive instruction from such a mouth? On meeting this man dost thou greet him at all with a kiss? To say nothing of his other impiety, dost thou not flee from the defilement, and from men worse than profligates, more detestable than any prostitute?

34. Of these things the Church admonishes and teaches thee, and touches mire, that thou mayest not be bemired: she tells of the wounds, that thou mayest not be wounded. But for thee it is enough merely to know them: abstain from learning by experience. God thunders, and we all tremble; and they blaspheme. God lightens, and we all bow down to the earth; and they have their blasphemous sayings about the heavens\footnote{\textgreek{Κᾀκεῖνοι περὶ οὐρανῶν τὰς δυσφήμους ἔχουσι γλώσσας. ᾽Ιησοῦς λέγει περὶ τοῦ πατρὸς αὐτοῦ, ῞Οστις τὸν ἥλιον αὐτοῦ ἀνατέλλει ἐπὶ δικαίους καὶ ἀδίκους, καὶ βρέχει ἐπὶ πονηροὺς καὶ ἀγαθούς. κᾀκεῖνοι λέγουσιν, ὅτι οἱ ὑετοὶ ἐξ ἐρωτικῆς μανίας γίνονται, καὶ τολμῶσι λέγειν, ὅτι ἐστί τις παρθένος ἐν οὐρανῷ εὐειδὴς μετὰ νεανίσκου εὐειδοῦς, καὶ κατὰ τὴν τῶν καμηλῶν ἢ λύκων καιρὸν, τοὺς τῆς αἰσχρᾶς ἐπιθυμίας καιροὺς ἔχειν, καὶ κατὰ τὴν τοῦ χειμῶνος καιρὸν, μανιωδῶς αὐτὸν ἐπιτρέχειν τῇ παρθένῳ, καὶ τὴν μὲν φεύγειν φασί, τὸν δὲ ἐπιτρέχειν, εἶτα ἐπιτρέχοντα ἱδροῦν, ἀπὸ δὲ τῶν ἱδρώτων αὐτοῦ εἶναι τὸν ὑετόν. Ταῦτα γέγραπται ἐν τοῖς τῶν Μανιχαίων βιβλίοις· ταῦτα ἡμεῖς ἀνέγνωμεν, κ.τ.λ}.}. These things are written in the books of the Manichees. These things we ourselves have read, because we could not believe those who told of them: yes, for the sake of your salvation we have closely inquired into their perdition.

35. But may the Lord deliver us from such delusion: and may there be given to you a hatred against the serpent, that as they lie in wait for the heel, so you may trample on their head. Remember ye what I say. What agreement can there be between our state and theirs? \emph{What communion hath light with darkness}\footnote{2 Cor. vi. 14.}? What hath the majesty of the Church to do with the abomination of the Manichees? Here is order, here is discipline\footnote{Gr. \textgreek{ἐπιστήμη}. See note on Introductory Lect. § 4.}, here is majesty, here is purity: here even \emph{to look upon a woman to lust after her}\footnote{Matt. v. 28.} is condemnation. Here is marriage with sanctity\footnote{\textgreek{σεμνότατος} is the reading of the chief \textsc{mss.} But the printed editions have \textgreek{σεμνότητος}, comparing it with such phrases as \textgreek{στόμα ἀθεότητος} (vi. 15), and \textgreek{μετάνοια τῆς σωτηρίας} (xiv. 17).}, here steadfast continence, here virginity in honour like unto the Angels: here partaking of food with thanksgiving, here gratitude to the Creator of the world. Here the Father of Christ is worshipped: here are taught fear and trembling before Him who sends the rain: here we ascribe glory to Him who makes the thunder and the lightning.

36. Make thou thy fold with the sheep: flee from the wolves: depart not from the Church. Hate those also who have ever been suspected in such matters: and unless in time thou perceive their repentance, do not rashly trust thyself among them. The truth of the Unity of God has been delivered to thee: learn to distinguish the pastures of doctrine. Be an approved banker\footnote{This saying is quoted three times in the Clementine Homilies as spoken by our Lord. See Hom. II. § 51; III. § 50; XVIII. § 20: “Every man who wishes to be saved must become, as the Teacher said, a judge of the books written to try us. For thus He spake: \emph{Become experienced bankers}. Now the need of bankers arises from the circumstance that the spurious is mixed up with the genuine.” \par On the same saying, quoted as Scripture in the Apostolic Constitutions (II. § 36), Cotelerius suggests that in oral tradition, or in some Apocryphal book, the proverb was said to come from the Old Testament, and was added by some transcriber as a gloss in the margin of Matt. xxv. 27, or Luke xix. 23. Dionysius of Alexandria, Epist. VII., speaks of “the Apostolic word, which thus urges all who are endowed with greater virtue, ‘Be ye skillful money-changers,’” referring apparently as here to 1 Thess. v. 21, 22, “try all things, \&c.” (See Euseb. \emph{E.H.} VII. ch. 6 in this series: Suicer. \emph{Thesaurus}, \textgreek{Τραπεζίτης}: and Resch. (\emph{Agrapha}, pp. 233–239.)}, \emph{holding fast that which is good, abstaining from every form of evil}\footnote{1 Thess. v. 21, 22.}. Or if thou hast ever been such as they, recognise and hate thy delusion. For there is a way of salvation, if thou reject the vomit, if thou from thy heart detest it, if thou depart from them, not with thy lips only, but with thy soul also: if thou worship the Father of Christ, the God of the Law and the Prophets, if thou acknowledge the Good and the Just to be one and the same God.\footnote{Compare § 13 of this Lecture, where Cyril seems to refer especially to the heresy of Manes, as described in the \emph{Disputatio Archelai}, cap. 6: “If you are desirous of being instructed in the faith of Manes, hear it briefly from me. That man worships two gods, unbegotten, self-originate, eternal, opposed one to the other. The one he represents as good, and the other as evil, naming the one Light, and the other Darkness.”} And may He preserve you all, guarding you from falling or stumbling, stablished in the Faith, in Christ Jesus our Lord, to Whom be glory for ever and ever. Amen.

\xchapter{Lecture VII. The Father.}

\textsc{Ephesians iii. 14, 15}

For this cause I bow my knees unto the Father,…of whom all fatherhood in heaven and earth is named, \&c.

1. \textsc{Of} God as the sole Principle we have said enough to you yesterday\footnote{See Lecture VI. 1, and 5.}: by “enough” I mean, not what is worthy of the subject, (for to reach that is utterly impossible to mortal nature), but as much as was granted to our infirmity. I traversed also the bye-paths of the manifold error of the godless heretics: but now let us shake off their foul and soul-poisoning doctrine, and remembering what relates to them, not to our own hurt, but to our greater detestation of them, let us come back to ourselves, and receive the saving doctrines of the true Faith, connecting the dignity of Fatherhood with that of the Unity, and believing \textsc{In One God the Father}: for we must not only believe in one God; but this also let us devoutly receive, that He is the Father of the Only-begotten, our Lord Jesus Christ.

2. For thus shall we raise our thoughts higher than the Jews\footnote{“In Athanasius, \emph{Quæstio} i. \emph{ad Antiochum}, tom. II. p. 331, Monarchia is opposed to Polytheism: ‘If we worship One God, it is manifest that we agree with the Jews in believing in a Monarchia: but if we worship three gods, it is evident that we follow the Greeks by introducing Polytheism, instead of piously worshipping One Only God.’” (Suicer, \emph{Thesaurus}, \textgreek{Μοναρχία}.)}, who admit indeed by their doctrines that there is One God, (for what if they often denied even this by their idolatries?); but that He is also the Father of our Lord Jesus Christ, they admit not; being of a contrary mind to their own Prophets, who in the Divine Scriptures affirm, \emph{The Lord said unto me, Thou art My Son, this day have I begotten thee}\footnote{Ps. ii. 7.}. And to this day they \emph{rage and gather themselves together against the Lord, and against His Anointed}\footnote{Ib. ii. 2.}, thinking that it is possible to be made friends of the Father apart from devotion towards the Son, being ignorant that \emph{no man cometh unto the Father but by}\footnote{John xiv. 6.} the Son, who saith, \emph{I am the Door, and I am the Way}\footnote{Ib. x. 9.}. He therefore that refuseth the Way which leadeth to the Father, and he that denieth the Door, how shall he be deemed worthy of entrance unto God? They contradict also what is written in the eighty-eighth Psalm, \emph{He shall call Me, Thou art my Father, my God, and the helper of my salvation. And I will make him my first-born, high among the kings of the earth}\footnote{Ps. lxxxix. 26, 27.}. For if they should insist that these things are said of David or Solomon or any of their successors, let them shew how \emph{the throne} of him, who is in their judgment described in the prophecy, is \emph{as the days of heaven, and as the sun before God, and as the moon established for ever}\footnote{\emph{vv}. 29, 36, 37.}. And how is it also that they are not abashed at that which is written, \emph{From the womb before the morning-star have I begotten thee}\footnote{Ps. cx. 3: “From the womb of the morning thou hast the dew of thy youth” (R.V.).}: also this, \emph{He shall endure with the sun, and before the moon, from generation to generation}\footnote{Ps. lxxii. 5.}. To refer these passages to a man is a proof of utter and extreme insensibility.

3. Let the Jews, however, since they so will, suffer their usual disorder of unbelief, both in these and the like statements. But let us adopt the godly doctrine of our Faith, worshipping one God the Father of the Christ, (for to deprive Him, who grants to all the gift of generation, of the like dignity would be impious): and let us \textsc{Believe in One God the Father}, in order that, before we touch upon our teaching concerning Christ, the faith concerning the Only-begotten may be implanted in the soul of the hearers, without being at all interrupted by the intervening doctrines concerning the Father.

4. For the name of the Father, with the very utterance of the title, suggests the thought of the Son: as in like manner one who names the Son thinks straightway of the Father also\footnote{Compare Athanasius (\emph{de Sententiâ Dionyssi}, § 17): “Each of the names I have mentioned is inseparable and indivisible from that next to it. I spoke of the Father, and before bringing in the Son, I designated Him also in the Father. I brought in the Son, and even if I had not previously mentioned the Father, in any wise He would have been presupposed in the Son.”}. For if a Father, He is certainly the Father of a Son; and if a Son, certainly the Son of a Father. Lest therefore from our speaking thus, \textsc{In One God, the Father Almighty, Maker of Heaven and Earth, and of All Things Visible and Invisible}, and from our then adding this also, \textsc{And in One Lord Jesus Christ}, any one should irreverently suppose that the Only-begotten is second in rank to heaven and earth,—for this reason before naming them we named \textsc{God the Father}, that in thinking of the Father we might at the same time think also of the Son: for between the Son and the Father no being whatever comes.

5. God then is in an improper sense\footnote{\textgreek{καταχρηστικῶς}. A technical term in Grammar, applied to the use of a word in a derived or metaphorical sense. See Aristotle’s description of the various kinds of metaphor, \emph{Poet}. § xxi. 7–16. The opposite to \textgreek{καταχρηστικῶς} is \textgreek{κυρίως}, as used in a parallel passage by Athanasius, \emph{Oratio} i. \emph{contra Arianos}, § 21 fin. “It belongs to the Godhead alone, that the Father is properly (\textgreek{κυρίως}) Father, and the Son properly Son.”} the Father of many, but by nature and in truth of One only, the Only-begotten Son, our Lord Jesus Christ; not having attained in course of time to being a Father, but being ever the Father of the Only-begotten\footnote{“And in Them, and Them only, does it hold, that the Father is ever Father, and the Son ever Son.” (Athan., \emph{as above}.)}. Not that being without a Son before, He has since by change of purpose become a Father: but before every substance and every intelligence, before times and all ages, God hath the dignity of Father, magnifying Himself in this more than in His other dignities; and having become a Father, not by passion\footnote{Compare vi. 6: \textgreek{ὁ γεννηθεὶς ἀπαθῶς.} The importance attached to the assertion of a “passionless generation” arose from the objections offered by Eusebius of Nicomedia and others to the word \textgreek{ὁμοούσιος} when proposed by Constantine at Nicæa. We learn from Eusebius of Cæsarea (\emph{Epist ad suæ parœciæ homines}, § 4) that the Emperor himself explained that the word was used “not in the sense of the affections (\textgreek{πάθη}) of bodies,” because “the immaterial, and intellectual, and incorporeal nature could not be the subject of any corporeal affection.” Again, in § 7, Eusebius admits that “there are grounds for saying that the Son is ‘one in essence’ with the Father, not in the way of bodies, nor like mortal beings, for He is not such by division of essence, or by severance, no, nor by any affection, or alteration, or changing of the Father’s essence and power.” (See the next note.)}, or union, not in ignorance, not by effluence\footnote{Athanasius (\emph{Expos. Fidei}, § 1): “Word not pronounced nor mental, nor an effluence of the Perfect, nor a dividing of the passionless nature.” Also (\emph{de Decretis}, § 11): “God being without parts is Father of the Son without partition or passion; for there is neither effluence of the Immaterial, nor influx from without, as among men.”}, not by diminution, not by alteration, \emph{for every good gift and every perfect gift is from above, coming down from the Father of lights, with whom can be no variation, neither shadow of turning}\footnote{James i. 17.}. Perfect Father, He begat a perfect Son, and delivered all things to Him who is begotten: (for \emph{all things}, He saith, \emph{are delivered unto Me of My Father}\footnote{Matt. xi. 27.}:) and is honoured by the Only-begotten: for, \emph{I honour My Father}\footnote{John viii. 49.}, saith the Son; and again, \emph{Even as I have kept My Father’s commandments, and abide in His love}\footnote{John xv. 10.}. Therefore we also say like the Apostle, \emph{Blessed be the God and Father of our Lord Jesus Christ, the Father of mercies, and God of all consolation}\footnote{2 Cor. i. 3.}: and, \emph{We bow our knees unto the Father from whom all fatherhood in heaven and on earth is named}\footnote{Eph. iii. 14, 15.}: glorifying Him with the Only-begotten: for \emph{he that denieth the Father, denieth the Son also}\footnote{1 John ii. 22: “This is the Antichrist, even he that denieth the Father and the Son” (\textsc{R.V.}).}: and again, \emph{He that confesseth the Son, hath the Father also}\footnote{\emph{v}. 23, bracketed in the A.V. as spurious, but rightly restored in R.V.}; knowing that Jesus Christ is Lord to the glory of God the Father\footnote{Phil. ii. 11.}.

6. We worship, therefore, as the Father of Christ, the Maker of heaven and earth, \emph{the God of Abraham, Isaac, and Jacob}\footnote{Ex. iii. 6.}; to whose honour the former temple also, over against us here, was built. For we shall not tolerate the heretics who sever the Old Testament from the New\footnote{Compare Lect. iv. 33.}, but shall believe Christ, who says concerning the temple, \emph{Wist ye not that I must be in My Father’s house}\footnote{Luke ii. 49.}? and again, \emph{Take these things hence, and make not my Father’s house a house of merchandise}\footnote{John ii. 16.}, whereby He most clearly confessed that the former temple in Jerusalem was His own Father’s house. But if any one from unbelief wishes to receive yet more proofs as to the Father of Christ being the same as the Maker of the world, let him hear Him say again, \emph{Are not two sparrows sold for a farthing, and not one of them shall fall on the ground without My Father which is in heaven}\footnote{Matt. x. 29. S. Cyril instead of “your Father” writes “my Father which is in heaven:” so Origen and Athanasius.}; this also, \emph{Behold the fowls of the heaven that they sow not, neither do they reap, nor gather into barns; and your heavenly Father feedeth them}\footnote{Matt. vi. 26.}; and this, \emph{My Father worketh hitherto, and I work}\footnote{John v. 17.}.

7. But lest any one from simplicity or perverse ingenuity should suppose that Christ is but equal in honour to righteous men, from His saying, \emph{I ascend to My Father, and your}\footnote{John xx. 17. On this text, quoted again in Cat. xi. 19, see the three Sermons of Bishop Andrewes \emph{On the Resurrection.}} \emph{Father}, it is well to make this distinction beforehand, that the name of the Father is one, but the power of His operation\footnote{\textgreek{ἐνεργεία}, meaning here, the operation of God, by nature in begetting His Son, by adoption in making many sons.} manifold. And Christ Himself knowing this has spoken unerringly, \emph{I go to My Father, and your Father:} not saying ‘to our Father,’ but distinguishing, and saying first what was proper to Himself, \emph{to My Father}, which was by nature; then adding, \emph{and your Father}, which was by adoption. For however high the privilege we have received of saying in our prayers, \emph{Our Father}, \emph{which art in heaven}, yet the gift is of loving-kindness. For we call Him Father, not as having been by nature begotten of Our Father which is in heaven; but having been transferred from servitude to sonship by the grace of the Father, through the Son and Holy Spirit, we are permitted so to speak by ineffable loving-kindness.

8. But if any one wishes to learn how we call God “Father,” let him hear Moses, the excellent schoolmaster, saying, \emph{Did not this thy Father Himself buy thee, and make thee, and create thee}\footnote{Deut. xxxii. 6.}? Also Esaias the Prophet, \emph{And now, O Lord. Thou art our Father: and we all are clay, the works of Thine hands}\footnote{Is. lxiv. 8.}. For most clearly has the prophetic gift declared that not according to nature, but according to God’s grace, and by adoption, we call Him Father.

9. And that thou mayest learn more exactly that in the Divine Scriptures it is not by any means the natural father only that is called father, hear what Paul says:—\emph{For though ye should have ten thousand tutors in Christ, yet have ye not many fathers: for in Christ Jesus I begat you through the Gospel}\footnote{1 Cor. iv. 15.}. For Paul was father of the Corinthians, not by having begotten them after the flesh, but by having taught and begotten them again after the Spirit. Hear Job also saying, \emph{I was a father of the needy}\footnote{Job xxix. 16.}: for he called himself a father, not as having begotten them all, but as caring for them. And God’s Only-begotten Son Himself, when nailed in His flesh to the tree at the time of crucifixion, on seeing Mary, His own Mother according to the flesh, and John, the most beloved of His disciples, said to him, \emph{Behold! thy mother}, and to her, \emph{Behold! thy Son}\footnote{John xix. 26, 27.}: teaching her the parental affection due to him\footnote{\textgreek{φιλοστοργία} might be applied to the mutual affection of mother and son, but the context shews that it refers here to parental love only; see Polybius, V. § 74, 5; Xenoph. \emph{Cyrop}. I. § 3, 2.}, and indirectly explaining that which is said in Luke, and \emph{His father and His mother marvelled at Him}\footnote{Luke ii. 33.}: words which the tribe of heretics snatch up, saying that He was begotten of a man and a woman. For like as Mary was called the mother of John, because of her parental affection, not from having given him birth, so Joseph also was called the father of Christ, not from having begotten Him (for \emph{he knew her not}, as the Gospel says, \emph{until she had brought forth her first-born Son}\footnote{Matt. i. 25.}), but because of the care bestowed on His nurture.

10. Thus much then at present, in the way of a digression, to put you in remembrance. Let me, however, add yet another testimony in proof that God is called the Father of men in an improper sense. For when in Esaias God is addressed thus, \emph{For Thou art our Father, though Abraham be ignorant of us}\footnote{Is. lxiii. 16.}, and \emph{Sarah travailed not with us}\footnote{Ib. li. 2.}, need we inquire further on this point? And if the Psalmist says, \emph{Let them be troubled from His countenance, the Father of the fatherless, and Judge of the widows}\footnote{Ps. lxviii. 5. Cyril quotes as usual from the Septuagint (Ps. lxvii. 6), where the clause \textgreek{ταραχθήσονται ἀπὸ προσώπου αὐτοῦ}, answering to nothing in the Hebrew, is evidently an interpolation, and may have crept in from a marginal quotation of Is. lxiv. 2.}, is it not manifest to all, that when God is called the Father of orphans who have lately lost their own fathers, He is so named not as begetting them of Himself, but as caring for them and shielding them. But whereas God, as we have said, is in an improper sense the Father of men, of Christ alone He is the Father by nature, not by adoption: and the Father of men in time, but of Christ before all time, as He saith, \emph{And now, O Father, glorify Thou Me with Thine own self, with the glory which I had with Thee before the world was}\footnote{John xvii. 5.}.

11. We believe then \textsc{In One God the Father} the Unsearchable and Ineffable, \emph{Whom no man hath seen}\footnote{1 Tim. ii. 16.}, \emph{but the Only-begotten alone hath declared Him}\footnote{John i. 18.}. \emph{For He which is of God, He hath seen God}\footnote{John vi. 46: \emph{He hath seen the Father}. The weight of authority is against the reading (\textgreek{τὸν θεόν}) which Cyril follows.}: whose face the Angels do alway behold in heaven\footnote{Matt. xviii. 10.}, behold, however, each according to the measure of his own rank. But the undimmed vision of the Father is reserved in its purity for the Son with the Holy Ghost.

12. Having reached this point of my discourse, and being reminded of the passages just before mentioned, in which God was addressed as the Father of men, I am greatly amazed at men’s insensibility. For God with unspeakable loving-kindness deigned to be called the Father of men,—He in heaven, they on earth,—and He the Maker of Eternity, they made in time,—He \emph{who holdeth the earth in the hollow of His hand}, they upon the earth as grasshoppers\footnote{Is. xl. 12 and 22.}. Yet man forsook his heavenly Father, and said to the stock, \emph{Thou art my father, and to the stone, Thou hast begotten me}\footnote{Jer. ii. 27.}. And for this reason, methinks, the Psalmist says to mankind, \emph{Forget also thine own people, and thy father’s house}\footnote{Ps. xlv. 10.}, whom thou hast chosen for a father, whom thou hast drawn upon thyself to thy destruction.

13. And not only stocks and stones, but even Satan himself, the destroyer of souls, have some ere now chosen for a father; to whom the Lord said as a rebuke, \emph{Ye do the deeds of your father}\footnote{John viii. 41.}, that is of the devil, he being the father of men not by nature, but by fraud. For like as Paul by his godly teaching came to be called the father of the Corinthians, so the devil is called the father of those who of their own will \emph{consent unto him}\footnote{Ps. l. 18.}.

For we shall not tolerate those who give a wrong meaning to that saying, \emph{Hereby know we the children of God, and the children of the devil}\footnote{1 John iii. 10.}, as if there were by nature some men to be saved, and some to be lost. Whereas we come into such holy sonship not of necessity but by choice: nor was the traitor Judas by nature a son of the devil and of perdition; for certainly he would never have cast out devils at all in the name of Christ: \emph{for Satan casteth not out Satan}\footnote{Mark iii. 23.}. Nor on the other hand would Paul have turned from persecuting to preaching. But the adoption is in our own power, as John saith, \emph{But as many as received Him, to them gave He power to become the children of God, even to them that believe in His name}\footnote{John i. 12.}. For not before their believing, but from their believing they were counted worthy to become of their own choice the children of God.

14. Knowing this, therefore, let us walk spiritually, that we may be counted worthy of God’s adoption. \emph{For as many as are led by the Spirit of God, they are the sons of God}\footnote{Rom. viii. 14.}. For it profiteth us nothing to have gained the title of Christians, unless the works also follow; lest to us also it be said, \emph{If ye were Abraham’s children, ye would do the works of Abraham}\footnote{John viii. 39.}. \emph{For if we call on Him as Father, who without respect of persons judgeth according to every man’s work, let us pass the time of our sojourning here in fear}\footnote{1 Pet. i. 17.}, \emph{loving not the world, neither the things that are in the world: for if any man love the world, the love of the Father is not in him}\footnote{1 John ii. 15.}. Wherefore, my beloved children, let us by our works offer glory to \emph{our Father which is in heaven, that they may see our good works, and glorify our Father which is in heaven}\footnote{Matt. v. 16.}. \emph{Let us cast all our care upon Him, for our Father knoweth what things we have need of}\footnote{1 Pet. v. 7; Matt. vi. 8.}.

15. But while honouring our heavenly Father let us honour also \emph{the fathers of our flesh}\footnote{Heb. xii. 9.}: since the Lord Himself hath evidently so appointed in the Law and the Prophets, saying, \emph{Honour thy father and thy mother, that it may be well with thee, and thy days shall be long in the land}\footnote{Deut. v. 16.}. And let this commandment be especially observed by those here present who have fathers and mothers. \emph{Children, obey your parents in all things: for this is well pleasing to the Lord}\footnote{Col. iii. 20.}. For the Lord said not, \emph{He that loveth father or mother is not worthy of Me}, lest thou from ignorance shouldest perversely mistake what was rightly written, but He added, \emph{more than Me}\footnote{Matt. x. 37.}. For when our fathers on earth are of a contrary mind to our Father in heaven, then we must obey Christ’s word. But when they put no obstacle to godliness in our way, if we are ever carried away by ingratitude, and, forgetting their benefits to us, hold them in contempt, then the oracle will have place which says, \emph{He that curseth father or mother, let him die the death}\footnote{Ex. xxi. 17; Lev. xx. 9; Matt. xv. 4.}.

16. The first virtue of godliness in Christians is to honour their parents, to requite the troubles of those who begat them\footnote{Compare for the thought Euripides, \emph{Medea}, 1029–1035.}, and with all their might to confer on them what tends to their comfort (for if we should repay them ever so much, yet we shall never be able to return their gift of life\footnote{\textgreek{ἀντιγεννῆσαι}. Jeremy Taylor (\emph{Ductor Dubitantium}, Book III. cap. ii. §17) mentions several stories in which a parent is nourished from a daughter’s breast, who thus ‘saves the life she cannot give.’}), that they also may enjoy the comfort provided by us, and may confirm us in those blessings which Jacob the supplanter shrewdly seized; and that our Father in heaven may accept\footnote{On the change of Moods, see Jelf, \emph{Greek Grammar}, § 809. The second verb (\textgreek{καταξίωσειεν}) expresses a wish and a consequence which might follow, if the first (\textgreek{στηρίξωσιν}) wish be realized, as it probably may be. Cf. Herod. ix. 51.} our good purpose, and judge us worthy \emph{to shine amid righteous as the sun in the kingdom of our Father\footnote{Matt. xiii. 43.}}: To whom be the glory, with the Only-begotten our Saviour Jesus Christ, and with the Holy and Life-giving Spirit, now and ever, to all eternity. Amen.

\xchapter{Lecture VIII. Almighty.}

\textsc{Jeremiah xxxix. 18, 19} \textsc{(Septuagint).}

\emph{The Great, the strong God, Lord of great Counsel, and mighty in His works, the Great God, the Lord Almighty and of great name}\footnote{The text is translated from the Septuagint, in which S. Cyril found the title \textsc{Almighty} (\textgreek{Παντοκράτωρ}), one of the usual equivalents in the Septuagint for \emph{Lord of Hosts} (\emph{Sabaoth}). In the English A.V. and R.V. the passage stands thus: Jer. xxxii. 18, 19: \emph{The Great, the Mighty God, the LORD of Hosts, is His name, Great in counsel, and mighty in work}.}.

1. \textsc{By} believing \textsc{In One God} we cut off all misbelief in many gods, using this as a shield against Greeks; and every opposing power of heretics; and by adding, \textsc{In One God the Father}, we contend against those of the circumcision, who deny the Only-begotten Son of God. For, as was said yesterday, even before explaining the truths concerning our Lord Jesus Christ, we made it manifest at once, by saying “The Father,” that He is the Father of a Son: that as we understand that God is, so we may understand that He has a Son. But to those titles we add that He is also “\textsc{Almighty};” and this we affirm because of Greeks and Jews\footnote{“For even the Jewish nation had wicked heresies: for of them were…the Pharisees, who ascribe the practice of sinners to fortune and fate; and the Basmotheans, who deny providence and say that the world is made by spontaneous motion” (\emph{Apost. Const}. VI. 6). Compare Euseb. (\emph{E.H}. IV. 22.)} together, and all heretics.

2. For of the Greeks some have said that God is the soul of the world\footnote{Cicero, \emph{De Natura Deorum}, Lib. I. 27: “Pythagoras thought that God was the soul pervading all nature.” The doctrine was accepted both by Stoics and Platonists, and became very general. Cf. Virg. \emph{Georg}. iv. 221: \par Deum namque ire per omnis \par Terrasque, tractusque maris, cælumque profundum. \par and \emph{Æn}. vi. 726: \par Spiritus intus alit, totamque infusa per artus \par Meus agitat molem, et magno se corpore miscet.}: and others that His power reaches only to heaven, and not to earth as well. Some also sharing their error and misusing the text which says, “\emph{And Thy truth unto the clouds}\footnote{Ps. xxxvi. 5. Cyril appears to have borrowed this statement from Clement of Alexandria, who states (\emph{Stromat}. V. xiv. § 91) that from this Psalm the thought occurred to Aristotle to let Providence come down as far as to the Moon.},” have dared to circumscribe God’s providence by the clouds and the heaven, and to alienate from God the things on earth; having forgotten the Psalm which says, \emph{If I go up into heaven, Thou art there, if I go down into hell, Thou art present}\footnote{Ps. cxxxix. 8.}. For if there is nothing higher than heaven, and if hell is deeper than the earth, He who rules the lower regions reaches the earth also.

3. But heretics again, as I have said before, know not One Almighty God. For He is Almighty who rules all things, who has power over all things. But they who say that one God is Lord of the soul, and some other of the body, make neither of them perfect, because either is wanting to the other\footnote{See note on Lect. IV. 4.}. For how is he almighty, who has power over the soul, but not over the body? And how is he almighty who has dominion over bodies, but no power over spirits? But these men the Lord confutes, saying on the contrary, \emph{Rather fear ye Him which is able to destroy both soul and body in hell}\footnote{Matt. x. 28.}. For unless the Father of our Lord Jesus Christ has the power over both, how does He subject both to punishment? For how shall He be able to take the body which is another’s and cast it into hell, \emph{except He first bind the strong man, and spoil his goods}\footnote{Ib. xii. 29.}?

4. But the Divine Scripture and the doctrines of the truth know but One God, who rules all things by His power, but endures many things of His will. For He rules even over the idolaters, but endures them of His forbearance: He rules also over the heretics who set Him at nought, but bears with them because of His long-suffering: He rules even over the devil, but bears with him of His long-suffering, not from want of power; as if defeated. For \emph{he is the beginning of the Lord’s creation, made to be mocked}\footnote{Job xl. 14, \textgreek{τοῦτ᾽ ἔστιν ἀρχὴ πλάσματος Κυρίου, πεποιῃμένον ἐγκαταπαίζεσθαι ὑπὸ τῶν ἀγγέλων αὐτοῦ}. In this description of Behemoth the Septuagint differs much from the Hebrew, which is thus rendered in our English Versions, xl. 19: \emph{He is the chief of the ways of God: he (only,} R.V.\emph{) that made him can make his sword to approach unto him}. Compare Job xli. 5: \emph{Wilt thou play with him as with a bird?} and Ps. civ. 26: \emph{There is that Leviathan whom thou hast formed to play therein} (Sept. \emph{to take thy pastime with him}). See Baruch iii. 17, with the note in the Speaker’s Commentary.}, not by Himself, for that were unworthy of Him, but \emph{by the Angels} whom He hath made. But He suffered him to live, for two purposes, that he might disgrace himself the more in his defeat, and that mankind might be crowned with victory. O all wise providence of God! which takes the wicked purpose for a groundwork of salvation for the faithful. For as He took the unbrotherly purpose of Joseph’s brethren for a groundwork of His own dispensation, and, by permitting them to sell their brother from hatred, took occasion to make him king whom He would; so he permitted the devil to wrestle, that the victors might be crowned; and that when victory was gained, he might be the more disgraced as being conquered by the weaker, and men be greatly honoured as having conquered him who was once an Archangel.

5. Nothing then is withdrawn from the power of God; for the Scripture says of Him, \emph{for all things are Thy servants}\footnote{Ps. cxix. 91.}. All things alike are His servants, but from all these One, His only Son, and One, His Holy Spirit, are excepted; and all the things which are His servants serve the Lord through the One Son and in the Holy Spirit. God then rules all, and of His long-suffering endures even murderers and robbers and fornicators, having appointed a set time for recompensing every one, that if they who have had long warning are still impenitent in heart, they may receive the greater condemnation. They are kings of men, who reign upon earth, but not without the power from above: and this Nebuchadnezzar once learned by experience, when he said; \emph{For His kingdom is an everlasting kingdom, and His power from generation to generation}\footnote{Dan. iv. 34.}.

6. Riches, and gold, and silver are not, as some think, the devil’s\footnote{On this doctrine of the Manicheans see Archelaus (\emph{Disputatio}, cap. 42), Epiphanius (\emph{Hæres}. lxvi. § 81). Compare Clement. Hom. xv. cap. 9: “To all of us possessions are sins.” Plato (\emph{Laws}, V. 743): “I can never agree with them that the rich man will be really happy, unless he is also good: but for one who is eminently good to be also extremely rich is impossible.”}: for \emph{the whole world of riches is for the faithful man, but for the faithless not even a penny}\footnote{Prov. xvii. 6, according to the Septuagint. See note on Cat. V. 2, where the same passage is quoted. Clement of Alexandria (\emph{Stromat.} II. 5) refers to it in connexion with the passage of Plato quoted in the preceding note. S. Augustine also quotes and explains it in \emph{Epist}. 153, § 26.}. Now nothing is more faithless than the devil; and God says plainly by the Prophet, \emph{The gold is Mine, and the silver is Mine, and to whomsoever I will I give it}\footnote{The former clause is from Haggai ii. 8; the latter, taken from the words of the Tempter in Luke iv. 6, is quoted both by Cyril and by other Fathers as if from Haggai. Chrysostom (\emph{Hom}. xxxiv. § 5, in 1 Cor. xiii.) treats the use which some made of the misquotation as ridiculous.}. Do thou but use it well, and there is no fault to be found with money: but whenever thou hast made a bad use of that which is good, then being unwilling to blame thine own management, thou impiously throwest back the blame upon the Creator. A man may even be justified by money: \emph{I was hungry, and ye gave Me meat}\footnote{Matt. xxv. 35, 36.}: that certainly was from money. \emph{I was naked, and ye clothed Me}: that certainly was by money. And wouldest thou learn that money may become a door of the kingdom of heaven? \emph{Sell}, saith He, \emph{that thou hast, and give to the poor, and thou shalt have treasure in heaven}\footnote{Ib. xix. 21.}.

7. Now I have made these remarks because of those heretics who count possessions, and money, and men’s bodies accursed\footnote{The connexion of \textgreek{σώματα} with money and possessions suggests the not uncommon meaning “slaves.” See Polyb. xviii. 18 § 6: \textgreek{καὶ τὴν ἐνδουχίαν ἀπέδοντο καὶ τὰ σώματα, καὶ σὺν τουτοις ἔτι τινὰς τῶν κτήσεων}, “household furniture, and slaves, and besides these some also of their lands.” See \emph{Dictionary of Christian Antiquities}, “Slavery,” where it is shewn that Christians generally and even Bishops still possessed slaves throughout the 4th Century. \par But here it is perhaps more probable that Cyril refers, as before, Cat. iv. § 23, to the Manichean doctrine of the body as the root of sin.}. For I neither wish thee to be a slave of money, nor to treat as enemies the things which God has given thee for use. Never say then that riches are the devil’s: for though he say, \emph{All these will I give thee, for they are delivered unto me}\footnote{Matt. iv. 9; Luke iv. 6.}, one may indeed even reject his assertion; for we need not believe the liar: and yet perhaps he spake the truth, being compelled by the power of His presence: for he said not, \emph{All these will I give thee}, for they are mine, but, \emph{for they are delivered unto me}. He grasped not the dominion of them, but confessed that he had been entrusted\footnote{For \textgreek{ἐγκεχειρῆσθαι}, the reading of all the printed Editions, which hardly yields a suitable sense, we should probably substitute \textgreek{ἐγκεχειρίσθαι}. A similar confusion of the two verbs occurs in Polybius (\emph{Hist}. VIII. xviii. 6); the proper use of the latter is seen in Joh. Damasc. (\emph{De Fide Orthod}. II. 4, quoted by Cleopas), who speaks of Satan as being “of these Angelic powers the chief of the earthly order, and entrusted by God with the guardianship of the earth” (\textgreek{τῆς γῆς τὴν φυλακὴν ἐγχειρισθεὶς παρὰ Θεοῦ}).} with them, and was for a time dispensing them. But at a proper time interpreters should inquire whether his statement is false or true\footnote{On this point compare Irenæus (\emph{Hær}. V. xxi.–xxiv.), and Gregory of Nyssa (\emph{Orat. Catech}. § 5).}.

8. God then is One, the Father, the Almighty, whom the brood of heretics have dared to blaspheme. Yea, they have dared to blaspheme the Lord of Sabaoth\footnote{The reference is to Manes, of whom his disciple Turbo says (\emph{Archelai Disput}. § 10), “the name Sabaoth, which is honourable and mighty with you, he declares to be the nature of man, and the parent of lust: for which reason the simple, he says, worship lust, and think it to be a god.”}, \emph{who sitteth above the Cherubim}\footnote{Ps. lxxx. 1.}: they have dared to blaspheme the Lord Adonai\footnote{\textgreek{᾽Αδωναΐ}, Heb. \texthebrew{י\texthebrew{נָדֹאַ}}, “the Lord,” an old form of the Plural of majesty, used of God only.}: they have dared to blaspheme Him who is in the Prophets the Almighty God\footnote{\textgreek{παντοκράτορα}, Heb. \texthebrew{י\texthebrew{דַּשׁ לא}”}, El-Shaddai, “God Almighty.”}. But worship thou One God the Almighty, the Father of our Lord Jesus Christ. Flee from the error of many gods, flee also from every heresy, and say like Job, \emph{But I will call upon the Almighty Lord, which doeth great things and unsearchable, glorious things and marvellous without number}\footnote{Job v. 8, 9. Cyril’s quotation agrees with the Codex Alexandrinus of the Septuagint, which has \textgreek{παντοκράτορα} , “Almighty,” while the Vatican and other \textsc{mss.} read \textgreek{τὸν πάντων δεσπότην}.}, and, \emph{For all these things there is honour from the Almighty}\footnote{Job xxxvii. 23: \emph{God hath upon Him terrible majesty} (R.V.). The Vatican and Alexandrine \textsc{mss.} of the Septuagint read \textgreek{ἐπὶ τούτοις μεγάλη ἡ δόξα καὶ τιμὴ παντοκράτορος}. (\emph{For these things great is the glory and honour of the Almighty}.) But Cyril’s text is the same as the Aldine and Complutensian.}: to Whom be the glory for ever and ever. Amen.

\xchapter{Lecture IX. On the Words, Maker of Heaven and Earth, and of All Things Visible and Invisible.}

\textsc{Job xxxviii. 2–3}

\emph{Who is this that hideth counsel from Me, and keepeth words in his heart, and thinketh to hide them from Me}\footnote{The Septuagint, from which Cyril quotes the text, differs much from the Hebrew, and from the English Versions: \emph{Who is this that darkeneth counsel by words without knowledge? Gird up now thy loins like a man: for I will demand of thee, and answer thou Me}.}\emph{?}

1. \textsc{To} look upon God with eyes of flesh is impossible: for the incorporeal cannot be subject to bodily sight: and the Only begotten Son of God Himself hath testified, saying, \emph{No man hath seen God at any time}\footnote{John i. 18.}. For if according to that which is written in Ezekiel any one should understand that Ezekiel saw Him, yet what saith the Scripture? \emph{He saw the likeness of the glory of the Lord}\footnote{Ezekiel i. 28.}; not the Lord Himself, but \emph{the likeness of His glory}, not the glory itself, as it really is. And when he saw merely \emph{the likeness of the glory}, and not the glory itself, he fell to the earth from fear. Now if the sight of the likeness of the glory brought fear and distress upon the prophets, any one who should attempt to behold God Himself would to a certainty lose his life, according to the saying, \emph{No man shall see My face and live}\footnote{Exod. xxxiii. 20.}. For this cause God of His great loving-kindness spread out the heaven as a veil of His proper Godhead, that we should not perish. The word is not mine, but the Prophet’s. \emph{If Thou shalt rend the heavens, trembling will take hold of the mountains at sight of Thee, and they will flow down}\footnote{Is. lxiv. 1, Septuagint. R.V. \emph{Oh that Thou wouldest rend the heavens, that Thou wouldest come down, that the mountains might flow down.}}. And why dost thou wonder that Ezekiel fell down on seeing \emph{the likeness of the glory?} when Daniel at the sight of Gabriel, though but a servant of God, straightway shuddered and fell on his face, and, prophet as he was, dared not answer him, until the Angel transformed himself into the likeness of a son of man\footnote{Dan. x. 9, 16, 18.}. Now if the appearing of Gabriel wrought trembling in the Prophets, had God Himself been seen as He is, would not all have perished?

2. The Divine Nature then it is impossible to see with eyes of flesh: but from the works, which are Divine, it is possible to attain to some conception of His power, according to Solomon, who says, \emph{For by the greatness and beauty of the creatures proportionably the Maker of them is seen}\footnote{Wisdom xiii. 5. Compare Theophilus of Antioch \emph{To Autolycus}, I. 5, 6: “God cannot indeed be seen by human eyes, but is beheld and perceived through His providence and works.…He is not visible to eyes of flesh, since He is incomprehensible.”}. He said not that from the creatures the Maker is seen, but added \emph{proportionably}. For God appears the greater to every man in proportion as he has grasped a larger survey of the creatures: and when his heart is uplifted by that larger survey, he gains withal a greater conception of God.

3. Wouldest thou learn that to comprehend the nature of God is impossible? The Three Children in the furnace of fire, as they hymn the praises of God, say \emph{Blessed art thou that beholdest the depths, and sittest upon the Cherubim}\footnote{Song of the Three Children, 32.}. Tell me what is the nature of the Cherubim, and then look upon Him who sitteth upon them. And yet Ezekiel the Prophet even made a description of them, as far as was possible, saying that \emph{every one has four faces}, one of a man, another of a lion, another of an eagle, and another of a calf; and that each one had six wings\footnote{In Ezekiel i. 6–11, the four living creatures have each \emph{four} wings, as also in x. 21 according to the Hebrew. But in the latter passage, according to the Vatican text of the Septuagint, each has \emph{eight} wings, as Codd. R. and Casaub. read here. Cyril seems to have confused the number in Ezekiel with that in Is. vi. 2: \emph{each one had six wings}. By “a wheel of four sides” Cyril explains Ez. i. 16: \emph{a wheel in the midst of a wheel}, as meaning two circles set at right angles to each other, like the equator and meridian on a globe.}, and they had eyes on all sides; and that under each one was a wheel of four sides. Nevertheless though the Prophet makes the explanation, we cannot yet understand it even as we read. But if we cannot understand the throne, which he has described, how shall we be able to comprehend Him who sitteth thereon, the Invisible and Ineffable God? To scrutinise then the nature of God is impossible: but it is in our power to send up praises of His glory for His works that are seen.

4. These things I say to you because of the following context of the Creed, and because we say, \textsc{We Believe in One God, the Father Almighty, Maker of Heaven and Earth, and of All Things Visible and Invisible;} in order that we may remember that the Father of our Lord Jesus Christ is the same as He that made the heaven and the earth\footnote{Compare Cat. iv. 4. Irenæus (I. x. 1): “The Church, though dispersed throughout the whole world, even to the ends of the earth, yet received from the Apostles and their disciples the Faith in One God the Father Almighty, Maker of heaven, and earth, and the sea and all that therein is.” Tertullian (\emph{de Præscriptione Hæret}. cap. xiii.) “The rule of faith is that whereby we believe that there is One God only, and none other than the Creator of the world, who brought forth all things out of nothing through His own Word first of all sent forth.”}, and that we may make ourselves safe against the wrong paths of the godless heretics, who have dared to speak evil of the All wise Artificer of all this world\footnote{Compare Cat. vi. 13, 27.}, men who see with eyes of flesh, but have the eyes of their understanding blinded.

5. For what fault have they to find with the vast creation of God?—they, who ought to have been struck with amazement on beholding the vaultings of the heavens: they, who ought to have worshipped Him who reared the sky as a dome, who out of the fluid nature of the waters formed the stable substance of the heaven. For God said, \emph{Let there be a firmament in the midst of the water}\footnote{Gen. i. 6.}. God spake once for all, and it stands fast, and falls not. The heaven is water, and the orbs therein, sun, moon, and stars are of fire: and how do the orbs of fire run their course in the water? But if any one disputes this because of the opposite natures of fire and water, let him remember the fire which in the time of Moses in Egypt flamed amid the hail, and observe the all-wise workmanship of God. For since there was need of water, because the earth was to be tilled, He made the heaven above of water that when the region of the earth should need watering by showers, the heaven might from its nature be ready for this purpose.

6. But what? Is there not cause to wonder when one looks at the constitution of the sun? For being to the sight as it were a small body he contains a mighty power; appearing from the East, and sending forth his light unto the West: whose rising at dawn the Psalmist described, saying: \emph{And he cometh forth out of his chamber as a bridegroom}\footnote{Ps. xix. 5.}. He was describing the brightness and moderation of his state on first becoming visible unto men: for when he rides at high noon, we often flee from his blaze: but at his rising he is welcome to all as a bridegroom to look on.

Observe also his arrangement (or rather not his, but the arrangement of Him who by an ordinance determined his course), how in summer he rises higher and makes the days longer, giving men good time for their works: but in winter contracts his course, that the period of cold may be increased, and that the nights becoming longer may contribute to men’s rest, and contribute also to the fruitfulness of the products of the earth\footnote{The common reading \textgreek{ἵνα μὴ τοῦ ψύχους πλείων γένηται ὁ χρόνος, ἀλλ᾽ ἵνα αἱ νύκτες, κ.τ.λ}. gives a meaning contrary to the facts. The translation follows the \textsc{mss.} Roe, Casaubon, which omit \textgreek{μή} and for \textgreek{ἀλλά} read \textgreek{καί}. Compare Whewell’s \emph{Astromony}, p. 22: “The length of the year is so determined as to be adapted to the constitution of most vegetables: or the construction of vegetables is so adjusted as to be suited to the length which the year really has, and unsuited to a duration longer or shorter by any considerable portion. The vegetable clock-work is so set as to go for a year.” \emph{Ibid}. p. 34: “The terrestrial day, and consequently the length of the cycle of light and darkness, being what it is, we find various parts of the constitution both of animals and vegetables, which have a periodical character in their functions, corresponding to the diurnal succession of external conditions, and we find that the length of the period, as it exists in their constitution, coincides with the length of the natural day.”}. See also how the days alternately respond each to other in due order, in summer increasing, and in winter diminishing; but in spring and autumn granting equal intervals one to another. And the nights again complete the like courses; so that the Psalmist also says of them, \emph{Day unto day uttereth speech, and night unto night proclaimeth knowledge}\footnote{Ps. xix. 2. Compare a beautiful passage of Theophilus of Antioch (\emph{To Autolycus}, vi.).}. For to the heretics who have no ears, they all but cry aloud, and by their good order say, that there is none other God save the Creator who hath set them their bounds, and laid out the order of the Universe\footnote{Lucretius, V. 1182: \par “They saw the skies in constant order run, \par The varied seasons and the circling sun, \par Apparent rule, with unapparent cause, \par And thus they sought in gods the source of laws.}.

7. But let no one tolerate any who say that one is the Creator of the light, and another of darkness\footnote{See note 3 on Cat. iii. 33.}: for let him remember how Isaiah says, \emph{I am the God who made the light, and created darkness}\footnote{Is. xlv. 7. Compare the Homily of Chrysostom on this text.}. Why, O man, art thou vexed thereat? Why art thou offended at the time that is given thee for rest\footnote{Whewell, \emph{Astromomy}. p. 38: “Animals also have a period in their functions and habits; as in the habits of waking, sleeping, eating, \&c., and their well-being appears to depend on the coincidence of this period with the length of the natural day.”}? A servant would have had no rest from his masters, had not the darkness necessarily brought a respite. And often after wearying ourselves in the day, how are we refreshed in the night, and he who was yesterday worn with toils, rises vigorous in the morning because of the night’s rest\footnote{Chrysostom, VI. p. 171: “As the day brings man out to his work, so the night succeeding releases him from his countless toils and thoughts, and lulling his weary eyes to sleep, and closing their lids, prepares him to welcome the sunbeam again with his force in full vigour.”}? And what more helpful to wisdom than the night\footnote{Clement of Alexandria (\emph{Stromat}. IV. 22, E. Tr.): “And in this way they seem to have called the night Euphrone, since then the soul released from the perceptions of sense turns in on itself, and has a truer hold of intelligence (\textgreek{φρόνησις}).”}? For herein oftentimes we set before our minds the things of God; and herein we read and contemplate the Divine Oracles. And when is our mind most attuned to Psalmody and Prayer? Is it not at night? And when have we often called our own sins to remembrance? Is not at night\footnote{Chrysostom (Tom. II. p. 793): “We usually take the reckoning of our money early in the morning, but of our actions, of all that we have said and done by day, let us demand of ourselves the account after supper, and even after nightfall, as we lie upon our bed, with none to trouble, none to disturb us. And if we see anything done amiss, let us chastise our conscience, let us rebuke our mind, let us so vehemently impugn our account, that we may no more dare to rise up and bring ourselves to the same pit of sin, being mindful of the scourging at night.”}? Let us not then admit the evil thought, that another is the maker of darkness: for experience shews that this also is good and useful.

8. They ought to have felt astonishment and admiration not only at the arrangement of sun and moon, but also at the well-ordered choirs of the stars, their unimpeded courses, and their risings in the seasons due to each: and how some are signs of summer, and others of winter; and how some mark the season for sowing, and others shew the commencement of navigation\footnote{Clem. Alex. (\emph{Stromat}. VI. 11): “The same is true also of Astronomy, for being engaged in the investigation of the heavenly bodies, as to the form of the universe, and the revolution of the heaven, and the motion of the stars, it brings the soul nearer to the Creative Power, and teaches it to be quick in perceiving the seasons of the year, the changes of the atmosphere, and the risings of the stars; since navigation also and husbandry are full of benefit from this science.” Compare Lactantius (\emph{De Irâ Dei}, cap. xiii.).}. And a man sitting in his ship, and sailing amid the boundless waves, steers his ship by looking at the stars. For of these matters the Scripture says well, \emph{And let them be for signs, and for seasons, and for years}\footnote{Gen. i. 14.}, not for fables of astrology and nativities. But observe how He has also graciously given us the light of day by gradual increase: for we do not see the sun at once arise; but just a little light runs on before, in order that the pupil of the eye may be enabled by previous trial to look upon his stronger beam: see also how He has relieved the darkness of the night by rays of moonlight.

9. \emph{Who is the father of the rain? And who hath begotten the drops of dew\footnote{Job xxxviii. 28.}}? Who condensed the air into clouds, and bade them carry the waters of the rain\footnote{Whewell, \emph{Astronomy}, p. 88: “\emph{Clouds} are produced by aqueous vapour when it returns to the state of water.” p. 89: “Clouds produce \emph{rain}. In the formation of a cloud the precipitation of moisture probably forms a fine watery \emph{powder}, which remains suspended in the air in consequence of the minuteness of its particles: but if from any cause the precipitation is collected in larger portions, and becomes \emph{drops}, these descend by their weight and produce a shower.” Compare Aristotle, \emph{Meterologica}, I. ix. 3; Ansted, \emph{Physical Geography}, p. 210.}, now \emph{bringing golden-tinted clouds from the north}\footnote{Job xxxvii. 22: “Out of the north cometh golden splendour” (\textsc{R.V.}).}, now changing these into one uniform appearance, and again transforming them into manifold circles and other shapes? \emph{Who can number the clouds in wisdom}\footnote{Job xxxviii. 37.}? Whereof in Job it saith, And He knoweth the separations of the clouds\footnote{Job xxxvii. 16: “Dost thou know the balancings of the clouds?” In the Septuagint \textgreek{διάκρισιν νεφῶν} may mean “the separate path of the clouds” (Vulg. “\emph{semitas nubium},”) or “the dissolving,” as in Aristotle (\emph{Meteorol.} I. vii. 10: \textgreek{διακρίνεσθαι καὶ διαλύεσθαι τὸ διάτμίζον ὑγρὸν ὑπὸ τοῦ πλήθους τῆς θερμῆς ἀναθυμιάσεως, ὥστε μὴ συνίστασθαι ῥαδίως εἰς ὕδωρ}. “The moist vapour is separated and dissolved by the great heat of the evaporation, so that it does not easily condense into water.” Cf. Plato, \emph{Sophistes} 243 B: \textgreek{διακρίσεις καὶ συγκρίσεις}.}, \emph{and hath bent down the heaven to the earth}\footnote{Job xxxviii. 37 (according to the Septuagint): “And who is he that numbereth the clouds by wisdom, and bent down the heaven to the earth?” \textsc{A.V., R.V.} “Or who can pour out the bottles of heaven?”}: and, \emph{He who numbereth the clouds in wisdom}: and, \emph{the cloud is not rent under Him}\footnote{Job xxvi. 8: “He bindeth up the waters in His thick clouds; and the cloud is not rent under them.”}. For so many measures of waters lie upon the clouds, yet they are not rent: but come down with all good order upon the earth. Who \emph{bringeth the winds out of their treasuries}\footnote{Ps. cxxxv. 7.}? \emph{And who}, as we said before, \emph{is he that hath begotten the drops of dew? And out of whose womb cometh the ice}\footnote{Job xxxviii. 28.}? For its substance is like water, and its strength like stone. And at one time the water becomes \emph{snow like wool}, at another it ministers to Him \emph{who scattereth the mist like ashes}\footnote{Ps. cxlvii. 16: “He scattereth the hoar frost like ashes.” The Hebrew \texthebrew{ד\texthebrew{וׂפכּ}} is rendered by \textgreek{πάχνη}, “hoar frost,” in Job xxxviii. 29, but here by \textgreek{ὀμιχλη}, “mist.”}, and at another it is changed into a stony substance; since \emph{He governs the waters as He will}\footnote{Job xxxvii. 10: “the breadth of the waters is straitened” (Marg. \textsc{R.V.} “congealed”). The word \textgreek{οἰακίζει} in the Septuagint means to “steer,” Lat. “gubernare” to “turn as by a helm.”}. Its nature is uniform, and its action manifold in force. Water becomes in vines \emph{wine that maketh glad the heart of man:} and in olives \emph{oil that maketh man’s face to shine:} and is transformed also into \emph{bread that strengtheneth man’s heart}\footnote{Ps. civ. 15.}, and into fruits of all kinds which He hath created\footnote{There is a similar passage on the various effects of water in Cat. xvi. 12. Chrysostom (\emph{de Statuis}, Hom. xii. 2), Epiphanius (\emph{Ancoratus}, p. 69), and other Fathers, appear to reproduce both the thoughts and words of Cyril.}.

10. What should have been the effect of these wonders? Should the Creator have been blasphemed? Or worshipped rather? And so far I have said noticing of the unseen works of His wisdom. Observe, I pray you, the spring, and the flowers of every kind in all their likeness still diverse one from another; the deepest crimson of the rose, and the purest whiteness of the lily: for these spring from the same rain and the same earth, and who makes them to differ? Who fashions them? Observe, pray, the exact care: from the one substance of the tree there is part for shelter, and part for divers fruits: and the Artificer is One. Of the same vine part is for burning\footnote{For \textgreek{καῦσιν}, “burning,” Morel and Milles, with Cod. Coisl., read \textgreek{καῦστιν}, a rare word explained by Hesychius as the “growth” or “foliage” of the vine: but this is fully expressed in what follows, and the reading \textgreek{καῦσιν} is confirmed by Virgil (\emph{Georg}. ii. 408): “Primus devecta cremato sarmenta” (Reischl).}, and part for shoots, and part for leaves, and part for tendrils, and part for clusters.

Admire also the great thickness of the knots which run round the reed, as the Artificer hath made them. From one and the same earth come forth creeping things, and wild beasts, and cattle, and trees, and food; and gold, and silver, and brass, and iron, and stone. The nature of the waters is but one, yet from it comes the substance of fishes and of birds; whereby\footnote{For the construction of \textgreek{ἵνα} with the Indicative \textgreek{ἵπτανται}, see Bernhardy, \emph{Syntax}, p. 401. Winer (\emph{Gram. N. T.} III. sect. xli. c).} as the former swim in the waters, so the birds fly in the air.

11. \emph{This great and wide sea, therein are things creeping innumerable}\footnote{Ps. civ. 25.}. Who can describe the beauty of the fishes that are therein? Who can describe the greatness of the whales, and the nature\footnote{Gr. \textgreek{ὑπόστασιν}, literally “substance.”} of its amphibious animals, how they live both on dry land and in the waters? Who can tell the depth and the breadth of the sea, or the force of its enormous waves? Yet it stays at its bounds, because of Him who said, \emph{Hitherto shalt thou come, and no further, but within thyself shall thy waves be broken}\footnote{Job xxxviii. 11.}. Which sea also clearly shews the word of the command imposed upon it, since after it has run up, it leaves upon the beach a visible line made by the waves, shewing, as it were, to those who see it, that it has not passed its appointed bounds.

12. Who can discern the nature of the birds of the air? How some carry with them a voice of melody, and others are variegated with all manner of painting on their wings, and others fly up into mid air and float motionless, as the hawk: for by the Divine command \emph{the hawk spreadeth out his wings and floateth motionless, looking towards the south}\footnote{Ib. xxxix. 26.}. What man can behold the eagle’s lofty flight? If then thou canst not discern the soaring of the most senseless of the birds, how wouldest thou understand the Maker of all?

13. Who among men knows even the names of all wild beasts? Or who can accurately discern the physiology of each? But if of the wild beasts we know not even the mere names, how shall we comprehend the Maker of them? God’s command was but one, which said, \emph{Let the earth bring forth wild beasts, and cattle, and creeping things, after their kinds}\footnote{Gen. i. 24.} and from one earth\footnote{Instead of \textgreek{φωνῆς} (Milles), or \textgreek{πηγῆς} (Bened. Roe, Casaub.) the recent Editors have restored \textgreek{τῆς γῆς} with the Jerusalem and Munich \textsc{mss.}, and Basil.}, by one command, have sprung diverse natures, the gentle sheep and the carnivorous lion, and various instincts\footnote{Gr. \textgreek{κινήσεις} “movements,” “impulses.” Aristotle (\emph{Historia Animalium}. IX. vii. 1) remarks that many imitations of man’s mode of life may be observed in the habits of other animals.} of irrational animals, bearing resemblance to the various characters of men; the fox to manifest the craft that is in men, and the snake the venomous treachery of friends, and the neighing horse the wantonness of young men\footnote{Jer. v. 8.}, and the laborious ant, to arouse the sluggish and the dull: for when a man passes his youth in idleness, then he is instructed by the irrational animals, being reproved by the divine Scripture saying, \emph{Go to the ant, thou sluggard, see and emulate her ways, and become wiser than she}\footnote{Prov. vi. 6. Instead of the epithet “laborious” (\textgreek{γεωργότατος} ) some \textsc{mss.} have “agile” or “restless” (\textgreek{γοργότατος}).}. For when thou seest her treasuring up her food in good season, imitate her, and treasure up for thyself fruits of good works for the world to come. And again, \emph{Go to the bee, and learn how industrious she is}\footnote{After the description of the ant, Prov. vi. 6–8, there follows in the Septuagint a similar reference to the bee: “Or go to the bee, and learn how industrious she is, and how comely she makes her work, and the produce of her labours kings and commons adopt for health, and she is desired and esteemed by all, and though feeble in strength has been exalted by her regard for wisdom.” The interpolation is supposed to be of Greek origin, as containing “idiomatic Greek expressions which would not occur to a translator from the Hebrew” (Delitzsch).}: how, hovering round all kinds of flowers, she collects her honey for thy benefit: that thou also, by ranging over the Holy Scriptures, mayest lay hold of salvation for thyself, and being filled with them mayest say, \emph{How sweet are thy words unto my throat, yea sweeter than honey and the honeycomb unto my mouth}\footnote{Ps. cxix. 103.}.

14. Is not then the Artificer worthy the rather to be glorified? For what? If thou knowest not the nature of all things, do the things that have been made forthwith become useless? Canst thou know the efficacy of all herbs? Or canst thou learn all the benefit which proceeds from every animal? Ere now even from venomous adders have come antidotes for the preservation of men\footnote{Compare Bacon (\emph{Natural Hist.} 965): “I would have trial made of two other kinds of bracelets, for comforting the heart and spirits: one of the trochisch of vipers, made into little pieces of beads; for since they do great good inwards (especially for pestilent agues), it is like they will be effectual outwards, where they may be applied in greater quantity. There would be trochisch likewise made of snakes; whose flesh dried is thought to have a very good opening and cordial virtue.” \emph{Ib}. 969: “The writers of natural magic commend the wearing of the spoil of a snake, for preserving of health.” Thomas Jackson (\emph{On the Creed}, VIII. 8, § 4): “The poisonous bitings of the scorpion are usually cured by the oil of scorpions.”}. But thou wilt say to me, “The snake is terrible.” Fear thou the Lord, and it shall not be able to hurt thee. “A scorpion stings.” Fear the Lord, and it shall not sting thee. “A lion is bloodthirsty.” Fear thou the Lord, and he shall lie down beside thee, as by Daniel. But truly wonderful also is the action of the animals: how some, as the scorpion, have the sharpness in a sting; and others have their power in their teeth; and others do battle with their claws; while the basilisk’s power is his gaze\footnote{Shakespeare (\emph{Richard III}. Act. i. Sc. ii.). \par Glo. “Thine eyes, sweet lady, have infected mine.” \par Anne. “Would they were basilisks to strike thee dead.” \par Compare Bacon (\emph{De Augmentis}, VII. cap. ii): “The fable goes of the basilisk, that if he see you first, you die for it, but if you see him first, he dies.” Bacon refers to Pliny (\emph{Nat. Hist}. viii. 33).}. So then from this varied workmanship understand the Creator’s power.

15. But these things perhaps thou knowest not: thou wouldest have nothing in common with the creatures which are without thee. Enter now into thyself, and from thine own nature consider its Artificer. What is there to find fault with in the framing of thy body? Be master of thyself, and nothing evil shall proceed from any of thy members. Adam was at first without clothing in Paradise with Eve, but it was not because of his members that he deserved to be cast out. The members then are not the cause of sin, but they who use their members amiss; and the Maker thereof is wise. Who prepared the recesses of the womb for child-bearing? Who gave life to the lifeless thing within it? \emph{Who knitted us with sinews and bones, and clothed us with skin and flesh}\footnote{Job x. 11.}, and, as soon as the child was born, brought streams of milk out of the breasts? How grows the babe into a boy, and the boy into a youth, and then into a man; and, still the same, passes again into an old man, while no one notices the exact change from day to day? Of the food, how is one part changed into blood, and another separated for excretion, and another part changed into flesh? Who gives to the heart its unceasing motion? Who wisely guarded the tenderness of the eyes with the fence of the eyelids\footnote{Xenophon (\emph{Memor. Socratis}. I. cap. iv): “And moreover does not this also seem to thee like a work of providence, that, whereas the sight is weak, the Creator furnished it with eyelids for doors, which are opened whenever there is need to use the sight, but are closed in sleep.”}? For as to the complicated and wonderful contrivance of the eyes, the voluminous books of the physicians hardly give us explanation. Who distributes the one breath to the whole body? Thou seest, O man, the Artificer, thou seest the wise Creator.

16. These points my discourse has now treated at large, having left out many, yea, ten thousand other things, and especially things incorporeal and invisible, that thou mayest abhor those who blaspheme the wise and good Artificer, and from what is spoken and read, and whatever thou canst thyself discover or conceive, \emph{from the greatness and beauty of the creatures mayest proportionably see the maker of them}\footnote{Wisdom xiii. 5.}, and bending the knee with godly reverence to the Maker of the worlds, the worlds, I mean, of sense and thought, both visible and invisible, thou mayest with a grateful and holy tongue, with unwearied lips and heart, praise God and say, \emph{How wonderful are Thy works, O Lord; in wisdom hast Thou made them all}\footnote{Ps. civ. 24.}. For to Thee belongeth honour, and glory, and majesty, both now and throughout all ages. Amen.

Appendix to Lecture IX.

\textsc{Note}.—In the manuscripts which contain this discourse under the name of “A Homily of S. Basil \emph{on God as Incomprehensible},” some portions are changed to suit that subject: but the conclusion especially is marked by great addition and variation, which it is well to reproduce here. Accordingly in place of the words in §15: \textgreek{τί μεμπτόν}, “What is there to find fault with?” and the following, the manuscripts before mentioned have it thus:

“What is there to find fault with in the framing of the body? Come forth into the midst and speak. Control thine own will, and nothing evil shall proceed from any of thy members. For every one of these has of necessity been made for our use. Chasten thy reasoning unto piety, submit to God’s commandments, and none of these members sin in working and serving in the uses for which they were made. If thou be not willing, the eye sees not amiss, the ear hears nothing which it ought not, the hand is not stretched out for wicked greed, the foot walketh not towards injustice, thou hast no strange loves, committest no fornication, covetest not thy neighbour’s wife. Drive out wicked thoughts from thine heart, be as God made thee, and thou wilt rather give thanks to thy Creator.

Adam at first was without clothing, faring daintily in Paradise: and after he had received the commandment, but failed to keep it, and wickedly stretched forth his hand (not because the hand wished this, but because his will stretched forth his hand to that which was forbidden), because of his disobedience he lost also the good things he had received. Thus the members are not the cause of sin to those who use them, but the wicked mind, as the Lord says, \emph{For out of the heart proceed evil thoughts, fornications, adulteries, envyings, and such like}. In what things thou choosest, therein thy limbs serve thee; they are excellently made for the service of the soul: they are provided as servants to thy reason. Guide them well by the motion of piety; bridle them by the fear of God; bring them into subjection to the desire of temperance and abstinence, and they will never rise up against thee to tyrannise over thee; but rather they will guard thee, and help thee more mightily in thy victory over the devil, while expecting also the incorruptible and everlasting crown of the victory. Who openeth the chambers of the womb? Who, \&c.”

At the end of the same section, after the words “Wise Creator,” this is found: “Glorify Him in His unsearchable works, and concerning Him whom thou art not capable of knowing, inquire not curiously what His essence is. It is better for thee to keep silence, and in faith adore, according to the divine Word, than daringly to search after things which neither thou canst reach, nor Holy Scripture hath delivered to thee. These points my discourse has now treated at large, that thou mayest abhor those who blaspheme the wise and good Artificer, and rather mayest thyself also say, \emph{How wonderful are Thy works O Lord; in wisdom hast Thou made them all}. To Thee be the glory, and power, and worship, with the Holy Spirit, now and ever, and throughout all ages. Amen.”

\xchapter{Lecture X. On the Clause, and in One Lord Jesus Christ, with a Reading from the First Epistle to the Corinthians.}

\emph{For though there be that are called gods, whether in heaven or on earth}\footnote{1 Cor. viii. 5, 6. Cyril omits the clause: \emph{as there be gods many and lords many}.}\emph{; yet to us there is One God, the Father, of whom are all things, and we in Him; and One Lord Jesus Christ, through whom are all things, and we through Him.}

1. \textsc{They} who have been taught to believe “\textsc{In One God the Father Almighty},” ought also to believe in His Only-begotten Son. For \emph{he that denieth the Son, the same hath not the Father}\footnote{1 John ii. 23.}. \emph{I am the Door}\footnote{Ib. x. 9.}, saith Jesus; \emph{no one cometh unto the Father but through Me}\footnote{Ib. xiv. 6.}. For if thou deny the Door, the knowledge concerning the Father is shut off from thee. \emph{No man knoweth the father, save the Son, and he to whomsoever the Son shall reveal Him}\footnote{Matt. xi. 27.}. For if thou deny Him who reveals, thou remainest in ignorance. There is a sentence in the Gospels, saying, \emph{He that believeth not on the Son, shall not see life; but the wrath of God abideth on him.}\footnote{John iii. 36.} For the Father hath indignation when the Only-begotten Son is set at nought. For it is grievous to a king that merely his soldier should be dishonoured; and when one of his nobler officers or friends is dishonoured, then his anger is greatly increased: but if any should do despite to the king’s only-begotten son himself, who shall appease the father’s indignation on behalf of his only-begotten son?

2. If, therefore, any one wishes to shew piety towards God, let him worship the Son, since otherwise the Father accepts not his service. The Father spake with a loud voice from heaven, saying, \emph{This is My beloved Son, in whom I am well pleased}\footnote{Matt. iii. 17.}. The Father was well pleased; unless thou also be well pleased in Him, thou hast not life. Be not thou carried away with the Jews when they craftily say, There is one God alone; but with the knowledge that God is One, know that there is also an Only-begotten Son of God. I am not the first to say this, but the Psalmist in the person of the Son saith, \emph{The Lord said unto Me, Thou art My Son}\footnote{Ps. ii. 7.}. Heed not therefore what the Jews say, but what the Prophets say. Dost thou wonder that they who stoned and slew the Prophets, set at nought the Prophets’ words?

3. Believe thou \textsc{In One Lord Jesus Christ, the Only-Begotten Son of God}. For we say “One Lord Jesus Christ,” that His Sonship may be “Only-begotten:” we say “One,” that thou mayest not suppose another: we say “One,” that thou mayest not profanely diffuse the many names\footnote{\textgreek{τὸ πολυώνυμον}, a word used by the Greek Poets of their gods, as by Homer (\emph{Hymn to Demeter}, 18, 32) of Zeus, \textgreek{Κρόνου πολυώνυμος υἱός}. Cf. Soph. \emph{Ant}. 1115; Æschyl. \emph{Prom}. V. 210.} of His action among many sons. For He is called a Door\footnote{John x. 7, 9. Cyril calls Christ a “spiritual,” or “rational” (\textgreek{λογική}) door, and applies the same term to His sheep, below. Origen (\emph{In Evang. Joh.} Tom. i. cap. 29): \textgreek{Θύρα ὁ Σωτηρ ἀναγέγραπται}, \emph{ibid}. \textgreek{φιλάνθρωπος δὲ ὢν…ποιμὴν γινεται}.}; but take not the name literally for a thing of wood, but a spiritual, a living Door, discriminating those who enter in. He is called a Way\footnote{John xiv. 6.}, not one trodden by feet, but leading to the Father in heaven; He is called a Sheep\footnote{Ib. i. 29; Is. liii. 7, 8; Acts viii. 32.}, not an irrational one, but the one which through its precious blood cleanses the world from its sins, which is led before the shearers, and knows when to be silent. This Sheep again is called a Shepherd, who says, \emph{I am the Good Shepherd}\footnote{John x. 11.}: a Sheep because of His manhood, a Shepherd because of the loving-kindness of His Godhead. And wouldst thou know that there are rational sheep? the Saviour says to the Apostles, \emph{Behold, I send you as sheep in the midst of wolves}\footnote{Matt. x. 10, 16.}. Again, He is called a Lion\footnote{Gen. xlix. 9; Apoc. v. 5.}, not as a devourer of men, but indicating as it were by the title His kingly, and stedfast, and confident nature: a Lion He is also called in opposition to the lion our adversary, who roars and devours those who have been deceived\footnote{1 Pet. v. 8.}. For the Saviour came, not as having changed the gentleness of His own nature, but as the strong \emph{Lion of the tribe of Judah}\footnote{Ps. cxviii. 22.}, saving them that believe, but treading down the adversary. He is called a Stone, not a lifeless stone, cut out by men’s hands, but a \emph{chief corner-stone}\footnote{Is. xxviii. 16; 1 Pet. ii. 4–6.}, on whom \emph{whosoever believeth shall not be put to shame}.

4. He is called \textsc{Christ}, not as having been anointed by men’s hands, but eternally anointed by the Father to His High-Priesthood on behalf of men\footnote{The reading of the earlier Editions \textgreek{ὑπὲρ ἀνθρώπων} is free from all difficulty, and so the more likely to have been substituted for what is at first sight more difficult \textgreek{ὑπὲρ ἄνθρωπον}, the reading of Cod. Coislin. adopted by the Benedictine and subsequent Editors. The idea of a super-human Priesthood to which the Son in His Divine nature was anointed by the Father from eternity is repeated by Cyril in § 14 of this Lecture, and in Cat. xi. 1, 14. See Index, “Priesthood,” and the reference there given to a fuller consideration of the subject in the Introduction.}. He is called Dead, not as having abode among the dead, as all in Hades, but as being alone \emph{free among the dead}\footnote{Ps. lxxxviii. 5.}. He is called Son of Man, not as having had His generation from earth, as each of us, but as coming upon the clouds \textsc{To Judge Both Quick and Dead\footnote{John v. 27. Comparing what Cyril says here with Cat. iv. 15, and xv. 10, we see that he means to explain why Christ is called the “Son of Man” when “He cometh again from heaven,” and “no more from earth.” The preceding clause refers to His first coming in the flesh, as differing in the manner of His conception and birth from other men.}}. He is called \textsc{Lord}, not improperly as those who are so called among men, but as having a natural and eternal Lordship\footnote{Cf. Athanas. (\emph{c. Arian.} II. xv. 14), “That very Word who was by nature Lord, and was then made man, hath by means of a servant’s form been made Lord of all and Christ.”}. He is called \textsc{Jesus} by a fitting name, as having the appellation from His salutary healing. He is called \textsc{Son}, not as advanced by adoption, but as naturally begotten. And many are the titles of our Saviour; lest, therefore, His manifold appellations should make thee think of many sons, and because of the errors of the heretics, who say that Christ is one, and Jesus another, and the Door another, and so on\footnote{Cf. Irenæus (III. xvi. 8): “All therefore are outside the Dispensation, who under pretence of knowledge understand that Jesus was one, and Christ another, and the Only-begotten another (from whom again is the Word), and the Saviour another.” The Cerinthians, Ebionites, Ophites, and Valentinians are mentioned by Irenæus as thus separating the Christ from Jesus.}, the Faith secures thee beforehand, saying well, \textsc{In One Lord Jesus Christ}: for though the titles are many, yet their subject is one.

5. But the Saviour comes in various forms to each man for his profit\footnote{Cf. Athanas. (\emph{Epist}. X.): “Since He is rich and manifold, He varies Himself according to the individual capacity of each soul.”}. For to those who have need of gladness He becomes a Vine; and to those who want to enter in He stands as a Door; and to those who need to offer up their prayers He stands a mediating High Priest. Again, to those who have sins He becomes a Sheep, that He may be sacrificed for them. \emph{He is made all things to all men}\footnote{1 Cor. ix. 22.}, remaining in His own nature what He is. For so remaining, and holding the dignity of His Sonship in reality unchangeable, He adapts Himself to our infirmities, just as some excellent physician or compassionate teacher; though He is Very Lord, and received not the Lordship by advancement\footnote{\textgreek{ἐκ προκοπῆς}. We learn from Athanasius (\emph{c. Arian}. i. 37, 38, 40), that from St. Paul’s language \emph{Philipp}. ii. 9: “Wherefore also God highly exalted Him, \&c.,” and from Ps. xlv. 7: “Thou hast loved righteousness and hated iniquity: therefore God, thy God, hath anointed thee with the oil of gladness above thy fellows,” the Arians argued that Christ first received Divine honour as Son and Lord as the reward of His obedience as Man. Athanasius replies (c. 40): “He was not from a lower state promoted; but rather, existing as God, He took the form of a servant, and in taking it was not promoted but humbled Himself. Where then is there here any reward of virtue, or what advancement (\textgreek{προκοπή}) and promotion in humiliation?” \par The same doctrine had been previously held by the disciples of Paul of Samosata, who said that Christ was not originally God, but after His Incarnation was by advance (\textgreek{ἐκ προκοπῆς}) made God, from being made by nature a mere man: see Athanas. (\emph{de Decretis}, § 24, \emph{c. Arian}. i. 38). S. Cyril refers to the error and uses the same word, in xi. 1, 7, 13, 15, 17, and xiv. 27.}, but has the dignity of His Lordship from nature, and is not called Lord improperly\footnote{\textgreek{καταχρηστικῶς}, i.e. in a secondary or metaphorical sense. Cf. vii. 5.}, as we are, but is so in verity, since by the Father’s bidding\footnote{\textgreek{νεύματι}, “command” or “bidding,” as expressed by nodding the head.} He is Lord of His own works. For our lordship is over men of equal rights and like passions, nay often over our elders, and often a young master rules over aged servants. But in the case of our Lord Jesus Christ the Lordship is not so: but He is first Maker, then Lord\footnote{Origen (\emph{De Principiis}, I. ii. 10) had argued that “even God cannot be called omnipotent, unless there exist those over whom He may exercise His power,” and therefore creation must have been eternal, or God could not have been eternally Omnipotent. In other passages Origen declares it an impiety to hold that matter is co-eternal with God (\emph{De Princip}. II. i. 4), and yet maintains that there were other worlds before this, and that there was never a time when there was no world existing. \par Methodius, in a fragment of his work \emph{On things Created}, preserved by Photius, and quoted by Bishop Bull (\emph{Def. Fid. Nic.} II. xiii. 9), argues against these theories of Origen, that in John i. 2 the words “The same was in the beginning with God” indicate the authority (\textgreek{τὸ ἐξουσιαστικόν}) of the Word which He had with the Father before the world came into existence; since from all eternity God the Father, together with His Word, possessed the Almighty power whereby whenever He would He could create worlds to rule over. \par Dean Church remarks that “On the other hand Tertullian, \emph{contra Hermog}. 3, considering the attributes in question to belong not to the Divine Nature, but Office, denies that God was Almighty (Lord?) from eternity; while the Greeks affirmed this (vid. Cyril Alex. \emph{in Joann}. xvii. 8, p. 963; Athan. \emph{Orat}. ii. 12–14), as understanding by the term the inherent but latent attribute of doing what He had not yet done, \textgreek{τὸ ἐξουσιαστικόν}.” \par Cleopas, the Jerusalem Editor, regards the passage as directed against Paul of Samosata, who asserted that Christ had become God, and received His kingdom and Lordship only after His Incarnation, and remarks:—“S. Cyril evidently regards the Lordship of Jesus Christ as twofold: one that which from eternity belonged to Him as God, which he calls natural, according to which ‘He was ever both Lord and King, as being by nature God’ (Cyril Alex. \emph{in Johann}. cap. xvii.); and the other the Lordship in time relative to the creatures, by which He exercises dominion over the works created by Him, as being their Maker.“}: first He made all things by the Father’s will, then, He is Lord of the things which were made by Him.

6. \emph{Christ the Lord} is He who was \emph{born in the city of David}\footnote{Luke ii. 11.}. And wouldest thou know that Christ is Lord with the Father even before His Incarnation\footnote{Among those who denied the Divine præ-existence of Christ Cleopas enumerates Ebion, Carpocrates, Theodotus, Artemon, Paul of Samosata, Marcellus, and Photinus.}, that thou mayest not only accept the statement by faith, but mayest also receive proof from the Old Testament? Go to the first book, Genesis: God saith, \emph{Let us make man}, not ‘in My image,’ but, \emph{in Our image}\footnote{Gen. i. 26.}. And after Adam was made, the sacred writer says, \emph{And God created man; in the image of God created He him}\footnote{Ib. i. 27.}. For he did not limit the dignity of the Godhead to the Father alone, but included the Son also: that it might be shewn that man is not only the work of God, but also of our Lord Jesus Christ, who is Himself also Very God. This Lord, who works together with the Father, wrought with Him also in the case of Sodom, according to the Scripture: \emph{And the Lord rained upon Sodom and Gomorrah fire and brimstone from the Lord out of heaven}\footnote{Ib. xix. 24.}. This Lord is He who afterwards was seen of Moses, as much as he was able to see. For the Lord is loving unto man, ever condescending to our infirmities.

7. Moreover, that you may be sure that this is He who was seen of Moses, hear Paul’s testimony, when he says, \emph{For they all drank of a spiritual rock that followed them; and the rock was Christ}\footnote{1 Cor. x. 4.}. And again: \emph{By faith Moses forsook Egypt}\footnote{Heb. xi. 27.}, and shortly after he says, \emph{accounting the reproach of Christ greater riches than the treasures in Egypt}\footnote{Heb. xi. 26. Quoting from memory Cyril mistakes the order of the two sentences.}. This Moses says to Him, \emph{Shew me Thyself}. Thou seest that the Prophets also in those times saw the Christ, that is, as far as each was able.\emph{ Shew me Thyself, that I may see Thee with understanding}\footnote{Ex. xxxiii. 13. Cyril means that even before His Incarnation Christ was seen as far as was possible by Prophets such as Moses. This view was held by many of the Fathers before Cyril. See Justin M. (\emph{Tryph}. § 56 ff.); Tertull. (\emph{adv. Praxean}, § 16); Euseb. (\emph{Demonstr. Evang}. V. 13–16).}. But He saith, \emph{There shall no man see My face, and live}\footnote{Ex. xxxiii. 20.}. For this reason then, because no man could see the face of the Godhead and live, He took on Him the face of human nature, that we might see this and live. And yet when He wished to shew even that with a little majesty, when \emph{His face did shine as the sun}\footnote{Matt. xvii. 2.}, the disciples fell down affrighted. If then His bodily countenance, shining not in the full power of Him that wrought, but according to the capacity of the Disciples, affrighted them, so that even thus they could not bear it, how could any man gaze upon the majesty of the Godhead? ‘A great thing,’ saith the Lord, ‘thou desirest, O Moses: and I approve thine insatiable desire, \emph{and I will do this thing}\footnote{Ex. xxxiii. 17. Gr. \textgreek{λόγον}, “word,” in imitation of the Hebrew idiom.} for thee, but according as thou art able. \emph{Behold, I will put thee in the clift of the rock}\footnote{Ex. xxxiii. 22.}: for as being little, thou shalt lodge in a little space.’

8. Now here I wish you to make safe what I am going to say, because of the Jews. For our object is to prove that the Lord Jesus Christ was with the Father. The \textsc{Lord} then says to Moses, \emph{I will pass by before thee with My glory, and will proclaim the name of the \textsc{Lord} before thee}\footnote{Ex. xxxiii. 19. Literally “will call in the name of the Lord (Jehovah):” compare Gen. iv. 26.}. Being Himself the \textsc{Lord}, what \textsc{Lord} doth He proclaim? Thou seest how He was covertly teaching the godly doctrine of the Father and the Son. And again, in what follows it is written word for word: \emph{And the \textsc{Lord} descended in the cloud, and stood with him there, and proclaimed the name of the \textsc{Lord}. And the \textsc{Lord} passed by before him, and proclaimed, The \textsc{Lord}, the \textsc{Lord} God, merciful and gracious, long-suffering, and abundant in goodness and truth, both keeping righteousness and shewing mercy unto thousands, taking away iniquities, and transgressions, and sins}\footnote{Ex. xxxiv. 5–7. For “keeping righteousness and shewing mercy,” the Hebrew has only “keeping mercy.”}. Then in what follows, \emph{Moses bowed his head and worshipped}\footnote{Ex. xxxiv. 8.} before the Lord who proclaimed the Father, and said: Go Thou then, O Lord, in the midst of us\footnote{Ib. xxxiv. 9.}.

9. This is the first proof: receive now a second plain one. \emph{The \textsc{Lord} said unto my Lord, sit Thou on My right hand}\footnote{Ps. cx. 1. Heb. “An oracle of Jehovah unto my lord.” Cyril’s argument is based upon the common mistake of supposing that \textgreek{Κύριος} represents the same Hebrew word in both parts of the sentence.}. The \textsc{Lord} says this to the Lord, not to a servant, but to the Lord of all, and His own Son, to whom He put all things in subjection. \emph{But when He saith that all things are put under Him, it is manifest that He is excepted, which did put all things under Him}, and what follows; \emph{that God may be all in all}\footnote{1 Cor. xv. 27, 28.}. The Only-begotten Son is Lord of all, but the obedient Son of the Father, for He grasped not the Lordship\footnote{Cyril evidently alludes to Philip. ii. 6, “Who being in the form of God thought it not a prize to be on an equality with God:” for the right interpretation of which passage, see Dean Gwynn’s notes in the \emph{Speaker’s Commentary}.}, but received it by nature of the Father’s own will. For neither did the Son grasp it, nor the Father grudge to impart it. He it is who saith, \emph{All things are delivered unto Me of My Father}\footnote{Matt. xi. 27; Luke x. 22. On this text Athanasius wrote a special treatise (\emph{In illud ‘Omnia,’ \&c}.), against the arguments of Arius, Eusebius, and their fellows, who said,—“If all things were delivered (meaning by ‘all’ the Lordship of Creation), there was once a time when He had them not. But it He had them not, He is not of the Father, for if He were, He would on that account have had them always.” \par Again (\emph{contr. Arian. Orat}. III. cap. xxvii. § 36), Athanasius argues: “Lest a man, perceiving that the Son has all that the Father hath, from the exact likeness and identity of what He hath, should wander into the impiety of Sabellius, considering Him to be the Father, therefore He has said, \emph{Was given unto Me}, and \emph{I received}, and \emph{Were delivered to Me}, only to shew that He is not the Father, but the Father’s Word, and the Eternal Son, who, because of His likeness to the Father, has eternally what He has from Him, and because He is the Son, has from the Father what eternally He hath.”}; “delivered unto Me, not as though I had them not before; and I keep them well, not robbing Him who hath given them.”

10. The Son of God then is Lord: He is Lord, who was born in Bethlehem of Judæa, according to the Angel who said to the shepherds, \emph{I bring you good tidings of great joy, that unto you is born this day in the city of David Christ the Lord}\footnote{Luke ii. 10, 11.}: of whom an Apostle says elsewhere, \emph{The word which God sent unto the children of Israel, preaching the gospel of peace by Jesus Christ: He is Lord of all}\footnote{Acts x. 36.}. But when he says, \emph{of all}, do thou except nothing from His Lordship: for whether Angels, or Archangels, or principalities, or powers, or any created thing named by the Apostles, all are under the Lordship of the Son. Of Angels He is Lord, as thou hast it in the Gospels, \emph{Then the Devil departed from Him, and the Angels came and ministered unto Him}\footnote{Matt. iv. 11.}; for the Scripture saith not, they succoured Him, but they \emph{ministered unto Him}, that is, like servants. When He was about to be born of a Virgin, Gabriel was then His servant, having received His service as a peculiar dignity. When He was about to go into Egypt, that He might overthrow the gods of Egypt made with hands\footnote{Isa. xix. 1. “Behold, the Lord rideth upon a swift cloud, and cometh unto Egypt: and the idols of Egypt shall be moved at His presence.” The prophecy was supposed by many of the Fathers to have been fulfilled by the flight into Egypt. Cf. Athanas. (\emph{Ep}. LXI. \emph{ad Maximum}, § 4): “As a child He came down to Egypt, and brought to nought its idols made with hands:” and (\emph{de Incarn}. § 36): “Which of the righteous men or kings went down into Egypt, so that at his coming the idols of Egypt fell?” On the passage of Isaiah see Delitzsch, and Kay (\emph{Speaker’s Commentary}).}, again \emph{an Angel appeareth to Joseph in a dream}\footnote{Matt. ii. 13.}. After He had been crucified, and had risen again, an Angel brought the good tidings, and as a trustworthy servant said to the women, \emph{Go, tell His disciples that He is risen, and goeth before you into Galilee; lo, I have told you}\footnote{Ib. xxviii. 7.}: almost as if he had said, “I have not neglected my command, I protest that I have told you; that if ye disregard it, the blame may not be on me, but on those who disregard it.” This then is the One Lord Jesus Christ, of whom the lesson just now read speaks: \emph{For though there be many that are called gods, whether in heaven or in earth}, and so on, \emph{yet to us there is One God, the Father, of whom are all things, and we in Him; and One Lord, Jesus Christ, through whom are all things, and we through Him}\footnote{1 Cor. viii. 5, 6.}.

11. And He is called by two names, Jesus Christ; Jesus, because He saves,—Christ, because He is a Priest\footnote{Compare Eusebius (\emph{Eccl. Hist}. I. cap. iii.), a passage which Cyril seems to have followed in his explanation of the names ‘Jesus’ and ‘Christ.’}. And knowing this the inspired Prophet Moses conferred these two titles on two men distinguished above all\footnote{For the common reading \textgreek{ἐγκρίτοις πάντων} Cod. Mon. I. has \textgreek{ἐκκρίτοις π.} which is required both by the construction and the sense. The change may have been caused by the occurrence of \textgreek{ἐγκρίτων} just below.}: his own successor in the government, Auses\footnote{Eusebius (\emph{u.s}): “His successor, therefore, who had not hitherto borne the name Jesus, but had been called by another name, Auses, which had been given him by his parents, he now called Jesus, bestowing the name upon him as a gift of honour far greater than any kingly diadem.” Auses is a common corruption of the name Oshea. See the note on the passage of Eusebius in this series.}, he renamed Jesus; and his own brother Aaron he surnamed Christ\footnote{Eusebius: “He consecrated a man high-priest of God, in so far as that was possible, and him he called Christ.” Cf. Lev. iv. 5, 16; vi. 22: \textgreek{ὁ ἱερεὺς ὁ Χριστός}}, that by two well-approved men he might represent at once both the High Priesthood, and the Kingship of the One Jesus Christ who was to come. For Christ is a High Priest like Aaron; since He \emph{glorified not Himself to be made a High Priest, but He that spake unto Him, Thou art a Priest for ever after the order of Melchizedek}\footnote{Heb. v. 4, 5, 6. Cyril omits from his quotation the reference to Ps. ii. 7: “Thou art My Son: this day have I begotten Thee.”}. And Jesus the son of Nave was in many things a type of Him. For when he began to rule over the people, he began from Jordan\footnote{Josh. iii. 1.}, whence Christ also, after He was baptized, began to preach the gospel. And the son of Nave appoints twelve to divide the inheritance\footnote{Ib. xiv. 1.}; and twelve Apostles Jesus sends forth, as heralds of the truth, into all the world. The typical Jesus saved Rahab the harlot when she believed: and the true Jesus says, \emph{Behold, the publicans and the harlots go before you into the kingdom of God}\footnote{Matt. xxi. 31.}. With only a shout the walls of Jericho fell down in the time of the type: and because Jesus said, \emph{There shall not be left here one stone upon another}\footnote{Matt. xxiv. 2.}, the Temple of the Jews opposite to us is fallen, the cause of its fall not being the denunciation but the sin of the transgressors.

12. There is One Lord Jesus Christ, a wondrous name, indirectly announced beforehand by the Prophets. For Esaias the Prophet says, \emph{Behold, thy Saviour cometh, having His own reward}\footnote{Isa. lxii. 11: “Behold, thy salvation cometh; behold, his reward is with him.”}. Now Jesus in Hebrew is by interpretation \emph{Saviour}. For the Prophetic gift, foreseeing the murderous spirit of the Jews against their Lord\footnote{\textgreek{τὸ κυριοκτόνον τῶν ᾽Ιουδαίων}.}, veiled His name, lest from knowing it plainly beforehand they might plot against Him readily. But He was openly called Jesus not by men, but by an Angel, who came not by his own authority, but was sent by the power of God, and said to Joseph, \emph{Fear not to take unto thee Mary thy wife; for that which is con}\emph{ceivedin her is of the Holy Ghost. And she shall bring forth a Son, and thou shalt call His name Jesus}\footnote{Matt. i. 20.}. And immediately he renders the reason of this name, saying, \emph{for He shall save His people from their sins}. Consider how He who was not yet born could have a \emph{people}, unless He was in being before He was born\footnote{The Anathema appended to the Creed of Nicæa condemns those who said \textgreek{πρὶν γεννηθῆναι οὐκ ἦν.} On this Eusebius of Cæsarea (\emph{Epist}. § 9) remarks: “Moreover to anathematize ‘Before His generation He was not,’ did not seem preposterous, in that it is confessed by all, that the Son of God was before the generation according to the flesh.”}. This also the Prophet says in His person, \emph{From the bowels of my mother hath He made mention of My name}\footnote{Isa. xlix. 1.}; because the Angel foretold that He should be called Jesus. And again concerning Herod’s plot again he says, \emph{And under the shadow of His hand hath He hid Me}\footnote{Ib. xlix. 2.}.

13. Jesus then means according to the Hebrew “Saviour,” but in the Greek tongue “The Healer;” since He is physician of souls and bodies, curer of spirits, curing the blind in body\footnote{\textgreek{τυφλῶν αἰσθητῶν}.}, and leading minds into light, healing the visibly lame, and guiding sinners’ steps to repentance, saying to the palsied, \emph{Sin no more}, and, \emph{Take up thy bed and walk}\footnote{John v. 14, 8.}. For since the body was palsied for the sin of the soul, He ministered first to the soul that He might extend the healing to the body. If, therefore, any one is suffering in soul from sins, there is the Physician for him: and if any one here is of little faith, let him say to Him, \emph{Help Thou mine unbelief}\footnote{Mark ix. 24.}. If any is encompassed also with bodily ailments, let him not be faithless, but let him draw nigh; for to such diseases also Jesus ministers\footnote{Compare the fragment of the Apology of Quadratus presented to Hadrian 127 \textsc{a.d.}, preserved by Eusebius (\emph{H.E}. IV. iii.): “But the works of our Saviour were always present, for they were genuine:—those that were healed, and those that arose from the dead, who were seen not only when they were healed and when they were raised, but were also always present; and not merely while the Saviour was on earth, but also after His death they were alive for a long while, so that some of them survived even to our times.” See the notes on the passage of Eusebius, in this series.}, and let him learn that Jesus is the Christ.

14. For that He is Jesus the Jews allow, but not further that He is Christ. Therefore saith the Apostle, \emph{Who is the liar, but he that denieth that Jesus is the Christ}\footnote{1 John ii. 22.}? But Christ is a High Priest, \emph{whose priesthood passes not to another}\footnote{Heb. vii. 24.}, neither having begun His Priesthood in time\footnote{On the opinion that Christ was from all eternity the true High Priest of the Creation, see Index, \emph{Priesthood}, and the reference there given to the Introduction. Cf. x. 4: xi. 1. Athan (c. Arian. \emph{Or}. ii. 12, \emph{J. H. N.}).}, nor having any successor in His High-Priesthood: as thou heardest on the Lord’s day, when we were discoursing in the congregation\footnote{The word ‘synaxis’ was used by the early Christians to distinguish their assemblies from the Jewish ‘synagogue,’ a word formed from the same root and more regularly. ‘Synaxis’ came to be used more especially of a celebration of the Eucharist. See Suicer, \emph{Thesaurus}, \textgreek{Σύναξις}.} on the phrase, \emph{After the Order of Melchizedek}. He received not the High-Priesthood from bodily succession, nor was He anointed with oil prepared by man\footnote{\textgreek{σκευαστῷ}, Ex. xxx. 22–25: “a perfume compounded (\textgreek{μυρεψικόν} ) after the art of the perfumer” (\textsc{R.V.}).}, but before all ages by the Father; and He so far excels the others as \emph{with an oath} He is made Priest: \emph{For they are priests without an oath, but He with an oath by Him that said, The Lord sware, and will not repent}\footnote{Heb. vii. 21.}. The mere purpose of the Father was sufficient for surety: but the mode of assurance is twofold, namely that with the purpose there follows the oath also, \emph{that by two immutable things, in which it was impossible for God to lie, we might have strong encouragement}\footnote{Ib. vi. 18.} for our faith, who receive Christ Jesus as the Son of God.

15. This Christ, when He was come, the Jews denied, but the devils confessed. But His forefather David was not ignorant of Him, when he said, \emph{I have ordained a lamp for mine Anointed}\footnote{Ps. cxxxii. 17. The “lamp for the Anointed” was commonly applied by the Fathers to John the Baptist. Compare John v. 35, and Bishop Westcott’s note there.}: which lamp some have interpreted to be the brightness of Prophecy\footnote{2 Pet. i. 19. The supposed reference in the Psalm to the lamp of prophecy is mentioned by Eusebius (\emph{Demonstr. Evang}. IV. cap. 16).}, others the flesh which He took upon Him from the Virgin, according to the Apostle’s word, \emph{But we have this treasure in earthen vessels}\footnote{2 Cor. iv. 7. The reference of the ‘lamp’ to Christ’s Incarnation is mentioned by Eusebius (\emph{u.s.}) and other Fathers.}. The Prophet was not ignorant of Him, when He said, \emph{and announceth unto men His Christ}\footnote{Amos. iv. 13: “and declareth unto man what is his thought.” For \texthebrew{ו\&x\#o@”\texthebrew{־המ}}, ‘what is his thought,’ the \textsc{LXX.} read \texthebrew{ו\texthebrew{ׂחישִֹמְ}}, ‘His Anointed,’ \textgreek{τὸν Χριστὸν αὐτοῦ}.}. Moses also knew Him, Isaiah knew Him, and Jeremiah; not one of the Prophets was ignorant of Him. Even devils recognised Him, for He rebuked them, and the Scripture says, \emph{because they knew that He was Christ}\footnote{Luke iv. 41.}. The Chief-priests knew Him not, and the devils confessed Him: the Chief Priests knew Him not, and a woman of Samaria proclaimed Him, saying, \emph{Come, see a man which told me all things that ever I did. Is not this the Christ}\footnote{John iv. 29.}?

16. This is Jesus Christ who came \emph{a High-Priest of the good things to come}\footnote{Heb. ix. 11.}; who for the bountifulness of His Godhead imparted His own title to us all. For kings among men have their royal style which others may not share: but Jesus Christ being the Son of God gave us the dignity of being called Christians. But some one will say, The name of “Christians” is new, and was not in use aforetime\footnote{\textgreek{οὐκ ἐπολιτεύετο}, “was not in citizenship,” “not naturalised.” Cf. Sueton. \emph{Nero}. cap. 16: “Christiani, genus hominum superstitionis novae et maleficae.”}: and new-fashioned phrases are often objected to on the score of strangeness\footnote{\textgreek{τὸ ξένον}.}. The prophet made this point safe beforehand, saying, \emph{But upon My servants shall a new name be called, which shall be blessed upon the earth}\footnote{Isa. lxv. 15, 16. The LXX. here depart from the meaning of the Hebrew: “\emph{He shall call His servants by another name: so that he who blesseth himself in the earth shall bless himself in the God of truth}” (R.V.).}. Let us question the Jews: Are ye servants of the Lord, or not? Shew then your new name. For ye were called Jews and Israelites in the time of Moses, and the other prophets, and after the return from Babylon, and up to the present time: where then is your new name? But we, since we are servants of the Lord, have that new name: \emph{new} indeed, but the \emph{new name, which shall be blessed upon the earth}. This name caught the world in its grasp: for Jews are only in a certain region, but Christians reach to the ends of the world: for it is the name of the Only-begotten Son of God that is proclaimed.

17. But wouldest thou know that the Apostles knew and preached the name of Christ, or rather had Christ Himself within them? Paul says to his hearers, \emph{Or seek ye a proof of Christ that speaketh in me}\footnote{2 Cor. xiii. 3.}? Paul proclaims Christ, saying, \emph{For we preach not ourselves, but Christ Jesus as Lord, and ourselves your servants for Jesus’ sake}\footnote{Ib. iv. 5.}. Who then is this? The former persecutor. O mighty wonder! The former persecutor himself preaches Christ. But wherefore? Was he bribed? Nay there was none to use this mode of persuasion. But was it that he saw Him present on earth, and was abashed? He had already been taken up into heaven. He went forth to persecute, and after three days the persecutor is a preacher in Damascus. By what power? Others call friends as witnesses for friends but I have presented to you as a witness the former enemy: and dost thou still doubt? The testimony of Peter and John, though weighty, was yet of a kind open to suspicion: for they were His friends. But of one who was formerly his enemy, and afterwards dies for His sake, who can any longer doubt the truth?

18. At this point of my discourse I am truly filled with wonder at the wise dispensation of the Holy Spirit; how He confined the Epistles of the rest to a small number, but to Paul the former persecutor gave the privilege of writing fourteen. For it was not because Peter or John was less that He restrained the gift; God forbid! But in order that the doctrine might be beyond question, He granted to the former enemy and persecutor the privilege of writing more, in order that we all might thus be made believers. For \emph{all were amazed} at Paul, and said, \emph{Is not this he that} was formerly a persecutor\footnote{Acts ix. 21.}? Did he not come hither, that he might lead us away bound to Jerusalem? Be not amazed, said Paul, I know that \emph{it is hard for me to kick against the pricks:} I know that \emph{I am not worthy to be called an Apostle, because I persecuted the Church of God}\footnote{1 Cor. xv. 9.}; but I did it \emph{in ignorance}\footnote{1 Tim. i. 13.}: for I thought that the preaching of Christ was destruction of the Law, and knew not that He came Himself to fulfil the Law and not to destroy it\footnote{Matt. v. 17.}. \emph{But the grace of God was exceeding abundant in me}\footnote{1 Tim. i. 14.}.

19. Many, my beloved, are the true testimonies concerning Christ. The Father bears witness from heaven of His Son: the Holy Ghost bears witness, descending bodily in likeness of a dove: the Archangel Gabriel bears witness, bringing good tidings to Mary: the Virgin Mother of God\footnote{\textgreek{ἡ θεοτόκος}—\emph{ Deipara}. Gibbon (Chap. xlvii. 34) says, “It is not easy to fix the invention of this word, which La Croze (\emph{Christianisme des Indes}, tom. i. p. 16) ascribes to Eusebius of Cæsarea and the Arians. The orthodox testimonies are produced by Cyril (of Alexandria) and Petavius (\emph{Dogmat. Theolog.} tom. v. L. v. cap. 15, p. 254, \&c.), but the veracity of the Saint is questionable, and the epithet of \textgreek{θεοτόκος} so easily slides from the margin to the text of a Catholic \textsc{ms.}” This passage is justly described as “Gibbon’s calumny” by Dr. Newman: see his notes on the title \textgreek{θεοτόκος} (\emph{Athan. c. Arian. Or.} ii. cap. 12, n.; \emph{Or}. iii. capp. 14, 29, 33). The word is certainly used by Origen (\emph{Deut}. xxii. 13, Lommatzch. Tom. x. p. 378): “She that is already betrothed is called a wife, as also in the case of Joseph and the Theotokos.” Cf. Archelaus (\emph{Disput. cum Mane}, cap. xxxiv. “qui de Maria Dei Genetrice natus est”); Eusebius (\emph{de Vita Constantini}, III. cap. 43: “The pious Empress adorned with rare memorials the place of the travail of the Theotokos”). For other examples see Suicer’s \emph{Thesaurus}, \textgreek{θεοτόκος}, Pearson, \emph{Creed}, Art. iii. notes l, m, n, o, and Routh, \emph{Reliq. Sacr.} ii. p. 332.} bears witness: the blessed place of the manger bears witness. Egypt bears witness, which received the Lord while yet young in the body\footnote{“Chrysostom describing the flourishing state of the Church in Egypt in those times, says: ‘Egypt welcomes and saves Him when a fugitive and plotted against, and receives a beginning as it were of its appropriation to Him, in order that when it shall hear Him proclaimed by the Apostles, it may in their day also be honoured as having been first to welcome Him’” (Cleopas).}: Symeon bears witness, who received Him in his arms, and said, \emph{Now, Lord, lettest Thou Thy servant depart in peace, according to Thy word; for mine eyes have seen Thy salvation, which Thou hast prepared before the face of all people}\footnote{Luke ii. 29, 30.}. Anna also, the prophetess, a most devout widow, of austere life, bears witness of Him. John the Baptist bears witness, the greatest among the Prophets, and leader of the New Covenant, who in a manner united both Covenants in Himself, the Old and the New. Jordan is His witness among rivers; the sea of Tiberias among seas: blind and lame bear witness, and dead men raised to life, and devils saying, \emph{What have we to do with Thee, Jesus? we know Thee, who Thou art, the Holy One of God}\footnote{Mark i. 24.}. Winds bear witness, silenced at His bidding: five loaves multiplied into five thousand bear Him witness. The holy wood of the Cross bears witness, seen among us to this day, and from this place now almost filling the whole world, by means of those who in faith take portions from it\footnote{See Cat. iv. 10, note 7.}. The palm-tree\footnote{The Bordeaux Pilgrim, who visited the Holy Places of Jerusalem, \textsc{a.d.} 333, \emph{c}. speaks of this palm-tree as still existing. The longevity of the palm was proverbial: cf. Aristot. (\emph{De Longitudine Vitæ}, c. iv. 2).} on the ravine bears witness, having supplied the palm-branches to the children who then hailed Him. Gethsemane\footnote{The same Pilgrim (as quoted by the Benedictine Editor) says, “There is also the rock where Judas Iscariot betrayed Christ.” Compare Cat. xiii. 38.} bears witness, still to the thoughtful almost shewing Judas. Golgotha\footnote{See Index, \emph{Golgotha}.}, the holy hill standing above us here, bears witness to our sight: the Holy Sepulchre bears witness, and the stone which lies there\footnote{See the passage of the Introduction referred to in Index, \emph{Sepulchre}.} to this day. The sun now shining is His witness, which then at the time of His saving Passion was eclipsed\footnote{See Cat. ii. 15, note 8, and xiii. 25, 34, 38. On the supernatural character of the darkness mentioned in the Gospels see Meyer, \emph{Commentary}, Matt. xxvii. 45. An eclipse of the sun was of course impossible, as the moon was full. Mr. J. R. Hind (\emph{Historical Eclipses}, “Times,” 19th July, 1872) states that the solar eclipse, mentioned by Phlegon the freedman of Hadrian, which occurred on Nov. 24, \textsc{a.d.} 29, and was partial at Jerusalem, is “the only solar eclipse that could have been visible at Jerusalem during the period usually fixed for the ministry of Christ.” He adds, “The Moon was eclipsed on the generally received date of the Crucifixion, 3 April, \textsc{a.d.} 33. I find she had emerged from the earth’s dark shadow a quarter of an hour before she rose at Jerusalem (6:36 p.m.), but the penumbra continued upon her disc for an hour afterwards.” Thus the “darkness from the sixth hour unto the ninth” cannot be explained as the natural effect of an eclipse either solar or lunar.}: the darkness is His witness, which was then from the sixth hour to the ninth: the light bears witness, which shone forth from the ninth hour until evening. The Mount of Olives bears witness, that holy mount from which He ascended to the Father: the rain-bearing clouds are His witnesses, having received their Lord: yea, and the gates of heaven bear witness [having received their Lord\footnote{This clause is omitted in Codd. Mon. 1, 2, Roe, Casaub., and is probably repeated from the preceding line: such repetitions, however, are not uncommon in Cyril’s style.}], concerning which the Psalmist said, \emph{Lift up your doors, O ye Princes, and be ye lift up ye everlasting doors; and the King of Glory shall come in}\footnote{Ps. xxiv. 7. The first clause is mistranslated by the \textsc{LXX.} from whom Cyril quotes.}. His former enemies bear witness, of whom the blessed Paul is one, having been a little while His enemy, but for a long time His servant: the Twelve Apostles are His witnesses, having preached the truth not only in words, but also by their own torments and deaths: \emph{the shadow of Peter}\footnote{Acts v. 15.} bears witness, having healed the sick in the name of Christ. The handkerchiefs and aprons bear witness, as in like manner by Christ’s power they wrought cures of old through Paul\footnote{Ib. xix. 12.}. Persians\footnote{The persecution of the Christians in Persia by Sapor II. is described at length by Sozomen (\emph{E.H.} II. cc. ix.–xv., in this Series). It commenced in \textsc{a.d.} 343, and was going on at the date of these Lectures and long after. “During fifty years the Cross lay prostrate in blood and ashes” (\emph{Dict. Bib.} ‘Sassanidæ’). Compare Neander, \emph{Church History}, Tom. III. p 148, Bohn.)} and Goths\footnote{The Goths here mentioned are the \emph{Gothi minores} dwelling on the north of the Danube, where Ulfilas, “the Apostle of the Goths” (311–381), converted many of his countrymen to Christianity. After suffering severe persecution, he was allowed by the Constantius to take refuge with his Arian converts in Mœsia and Thrace. This migration took place in 348 \textsc{a.d.}, the same year in which Cyril’s Lectures were delivered.}, and all the Gentile converts bear witness, by dying for His sake, whom they never saw with eyes of flesh: the devils, who to this day\footnote{See Index, \emph{Exorcism}.} are driven out by the faithful, bear witness to Him.

20. So many and diverse, yea and more than these, are His witnesses: is then the Christ thus witnessed any longer disbelieved? Nay rather if there is any one who formerly believed not, let him now believe: and if any was before a believer, let him receive a greater increase of faith, by believing in our Lord Jesus Christ, and let him understand whose name he bears. Thou art called a Christian: be tender of the name; let not our Lord Jesus Christ, the Son of God, be blasphemed through thee: but rather \emph{let your good works shine before men}\footnote{Matt. v. 16.} that they who see them may in Christ Jesus our Lord glorify the Father which is in heaven: To whom be the glory, both now and for ever and ever. Amen.

\xchapter{Lecture XI. On the Words, the Only-Begotten Son of God, Begotten of the Father Very God Before All Ages, by Whom All Things Were Made.}

\textsc{Hebrews i. 1}

God, who at sundry times and in divers manners spake in times past unto the Fathers by the Prophets, hath in these last days spoken unto us by His Son.

1. \textsc{That} we have hope in Jesus Christ has been sufficiently shewn, according to our ability, in what we delivered to you yesterday. But we must not simply believe in Christ Jesus nor receive Him as one of the many who are improperly called Christs\footnote{\textgreek{Compare x. 11, 15; xvi. 13: xxi. 1.}}. For they were figurative Christs, but He is the true Christ; not having risen by advancement\footnote{\textgreek{ἐκ προκοπῆς}. See x. 5. note 8.} from among men to the Priesthood, but ever having the dignity of the Priesthood from the Father\footnote{Compare x. 14, note 9.}. And for this cause the Faith guarding us beforehand lest we should suppose Him to be one of the ordinary Christs, adds to the profession of the Faith, that we believe \textsc{In One Lord Jesus Christ, the Only-Begotten Son of God}.

2. And again on hearing of a “Son,” think not of an adopted son but a Son by nature\footnote{\textgreek{θετόν}. Athanasius (\emph{de Sententiâ Dionysii}, § 23), represents Arius as saying that the Word “is not by nature (\textgreek{κατὰ φύσιν}) and in truth Son of God, but is called Son, He too, by adoption (\textgreek{κατὰ θέσιν}), as a creature.” Again (\emph{c. Arian. Orat}. iii. 19), he says, “This is the true God and the Life eternal, and we are made sons through Him by adoption and grace (\textgreek{θέσει καὶ χάριτι}).” Cf. vii. 10, and § 4, below.}, an Only-begotten Son, having no brother. For this is the reason why He is called “Only-begotten,” because in the dignity of the Godhead, and His generation from the Father, He has no brother. But we call Him the Son of God, not of ourselves, but because the Father Himself named Christ\footnote{The \textsc{mss.} all read \textgreek{αὐτὸν Χριστόν} which might mean “Christ and no other.” But \textgreek{Χριστόν} is probably a gloss introduced from the margin.} His Son\footnote{Compare the passages in which Cyril quotes Ps. ii. 7, as Cat. vii. 2; x. 2; xi. 5; xii. 18.}: and a true name is that which is set by fathers upon their children\footnote{“It was one of the especial rights of a father to choose the names for his children, and to alter them if he pleased” (\emph{Dict. Greek and Roman Antiq}. “Nomen. 1 Greek.”) The right to the name given by the father is the subject of one of the Private Orations of Demosthenes (\textgreek{Πρὸς Βοιωτὸν περὶ τοῦ ὀνόματος}).}.

3. Our Lord Jesus Christ erewhile became Man, but by the many He was unknown. Wishing, therefore, to teach that which was not known, He called together His disciples, and asked them, \emph{Whom do men say that I, the Son of Man, am}\footnote{Matt. xiii. 16.}? —not from vain-glory, but wishing to shew them the truth, lest dwelling with God, the Only-begotten of God\footnote{Compare iv. 7: “God of God begotten;” xiii. 3 and 13: “God the Son of God.” Here however, the \textsc{mss.} vary, and the reading of Cod. Coisl. \textgreek{Υἱῷ Θεοῦ μονογενεῖ} is approved by the Benedictine Editor, though not adopted. The confusion of \textgreek{Υἱῷ} and \textgreek{Θεῷ} is like that in John i. 18.}, they should think lightly of Him as if He were some mere man. And when they answered that some said Elias, and some Jeremias, He said to them, They may be excused for not knowing, but ye, My Apostles, who in My name cleanse lepers, and cast out devils, and raise the dead, ought not to be ignorant of Him, through whom ye do these wondrous works. And when they all became silent (for the matter was too high for man to learn), Peter, the foremost of the Apostles and chief herald\footnote{\textgreek{ὁ πρωτοστάτης τῶν ᾽Αποστόλων καὶ τῆς ᾽Εκκλησίας κορυφαῖος κήρυξ}. Cf. ii. 19.} of the Church, neither aided by cunning invention, nor persuaded by human reasoning, but enlightened in his mind from the Father, says to Him, \emph{Thou art the Christ}, not only so, but \emph{the Son of the living God}. And there follows a blessing upon his speech (for in truth it was above man), and as a seal upon what he had said, that it was the Father who had revealed it to him. For the Saviour says, \emph{Blessed art thou, Simon Barjona, for flesh and blood hath not revealed it to thee, but My Father which is in heaven}\footnote{Matt. xvi. 17.}. He therefore who acknowledges our Lord Jesus Christ the Son of God, partakes of this blessedness; but he who denies the Son of God is a poor and miserable man.

4. Again, I say, on hearing of a Son, understand it not merely in an improper sense, but as a Son in truth, a Son by nature, without beginning\footnote{Athanasius (\emph{de Synodis}, § 15) quotes a passage from the \emph{Thalia} of Arius, in which he says: “We praise Him as without beginning, because of Him who has a beginning: and adore Him as eternal, because of Him who in time has come to be. He who is without beginning made the Son a beginning of things created.” \par It is important, therefore, to notice the sense in which Cyril here calls the Son \textgreek{ἄναρχος}. The word has two meanings, which should be clearly distinguished, (i) \emph{unoriginate,} (ii) \emph{without beginning in time}. The former referring to origin, or cause, can properly be applied to the One true God, or to God the Father only, as it is used by Clement of Alexandria (\emph{Protrept}. cap. v. § 65: \textgreek{τὸν πάντων ποιητὴν…ἀγνοοῦντες, τὸν ἄναρχον Θεόν}. [\emph{Strom}. IV. cap. xxv. § 164: \textgreek{ὁ Θεὸς δὲ ἄναρχος ἀρχὴ τῶν ὅλων παντελὴς ἀρχῆς ποιητικός}]. [\emph{Stromat}. V. cap. xiv. § 142: \textgreek{ἐξ ἀρχῆς ἀνάρχου}]. Methodius (\emph{ob}. 312 \textsc{a.d.} \emph{circ.}) in a fragment (\emph{On things created}, § 8, English Trans. Clark’s Ante-Nic. Libr.) comments thus on John i. 1–3: “And so after \emph{the peculiar unbeginning beginning}, who is the Father, He (the Word) is the beginning of other things, ‘by whom all things are made.’” \par In this sense Cyril has said (iv. 4) that God alone is “unbegotten, unoriginate:” and in xi. 20 he explains this more fully,—“Suffer none to speak of a beginning of the Son in time (\textgreek{χῥονικὴν ἀρχήν}), but as a timeless beginning acknowledge the Father. For the Father is the beginning of the Son, timeless, incomprehensible, without beginning.” From a confusion of the two meanings the word came to be improperly applied in the sense of “unoriginate” to the Son, and to the Spirit; and this improper usage is condemned in the 49th \emph{Apostolic Canon}, which Hefele regards as amongst the most ancient Canons, and taken from the \emph{Apostolic Constitutions}, vi. 11: “If any Bishop or Presbyter shall baptize not according to our Lord’s ordinance into the Father, and Son, and Spirit, but into \emph{three Unoriginates}, or three Sons, or three Paracletes let him be deposed.” (ii.) Athanasius frequently calls the Son \textgreek{ἄναρχος} in the sense of ‘timeless,’ as being the co-eternal brightness (\textgreek{ἀπαύγασμα}) of the Eternal Light: see \emph{de Sent. Dionys.} §§ 15, 16, 22; “God is the Eternal Light, which never either began or shall cease: accordingly the Brightness is ever before Him, and co-exists with Him, without beginning and ever-begotten (\textgreek{ἄναρχον καὶ ἀειγενές}).”}; not as having come out of bondage into a higher state of adoption\footnote{\textgreek{εἰς προκοπὴν υἱοθεσίας.} Cf. § 2, note 4.}, but a Son eternally begotten by an inscrutable and incomprehensible generation. And in like manner on hearing of the First-born\footnote{\textgreek{Πρωτότοκον}. The word occurs in Heb. i. 6, which had been read in the Lesson before this Lecture. The exact dogmatic sense of the word is carefully explained by Athanasius (\emph{c. Arian.} Or. ii. 62): “The same cannot be both Only-begotten and Firstborn, except in different relations;—that is, Only-begotten, because of His generation from the father, as has been said; and First-born, because of His condescension to the creation, and His making the many His brethren.” See Mr. Robertson’s discussion of the word \textgreek{πρωτότοκος} (\emph{Athan.} p. 344, in this series), and Bp. Bull (\emph{Def. Fid. Nic}. iii. 5–8).}, think not that this is after the manner of men; for the first-born among men have other brothers also. And it is somewhere written, \emph{Israel is My son, My first-born}\footnote{Ex. iv. 22.}. But Israel is, as Reuben was, a first-born son rejected: for Reuben went up to his father’s couch; and Israel cast his Father’s Son out of the vineyard, and crucified Him.

To others also the Scripture says, \emph{Ye are the sons of the Lord your God}\footnote{Deut. xiv. 1.}: and in another place, \emph{I have said, Ye are gods, and ye are all sons of the Most High}\footnote{Ps. lxxxii. 6.}. \emph{I have said}, not, “I have begotten.” They, when God so \emph{said}, received the sonship, which before they had not: but He was not begotten to be other than He was before; but was begotten from the beginning Son of the Father, being above all beginning and all ages, Son of the Father, in all things like\footnote{\textgreek{ἐν πᾶσιν ὅμοιος}. See the note on iv. 7. That the phrase was not equivalent to \textgreek{ὁμοούσιος}, and did not adequately express the relation of the Son to the Father is clearly shewn by Athanasius (\emph{de Synodis}, cap. iii. § 53).} to Him who begat Him, eternal of a Father eternal, Life of Life begotten, and Light of Light, and Truth of Truth, and Wisdom of the Wise, and King of King, and God of God, and Power of Power\footnote{The additions which the Benedictine Editor has here made to the earlier text, as represented by Milles, may be conveniently shewn in brackets. \textgreek{ἀλλὰ Υἱὸς [τοῦ Πατρὸς]* ἐξ ἀρχῆς ἐγεννήθη, [ὑπεράνω πάσης ἀρχῆς καὶ αἰώνων τυγχάνων]*, Υιὸς τοῦ Πατρὸς [ἐν πᾶσιν]}† \textgreek{ὅμοιος τῷ γεγεννηκότι· [ἀΐδιος ἐξ ἀϊδίου Πατρός,]* ζωὴ ἐκ ζωῆς γεγεννημένος. …καὶ Θεὸς ἐκ Θεοῦ, [καὶ δύναμις ἐκ δυνάμεως}]‡. \par * Codd. Coisl. Ottob. Mon. 2. † Coisl. Ottob. Roe, Casaub. Mon. 1, 2. \par  ‡ Coisl. Ottob. Mon. 1, 2.}.

5. If then thou hear the Gospel saying, \emph{The book of the generation of Jesus Christ, the Son of David, the Son of Abraham}\footnote{Matt. i. 1.}, understand “according to the flesh.” For He is the Son of David \emph{at the end of the ages}\footnote{Heb. ix. 26.}, but the Son of God \textsc{Before All Ages}, without beginning\footnote{See § 4, note 3.}. The one, which before He had not, He received; but the other, which He hath, He hath eternally as begotten of the Father. Two fathers He hath: one, David, according to the flesh, and one, God, His Father in a Divine manner\footnote{\textgreek{Θεϊκῶς}.}. As the Son of David, He is subject to time, and to handling, and to genealogical descent: but as Son according to the Godhead\footnote{\textgreek{τὸ μὲν κατὰ τὸν Δαβίδ.…τὸ δὲ κατὰ τὴν Θεότητα}.}, He is subject neither to time nor to place, nor to genealogical descent: for \emph{His generation who shall declare}\footnote{Isa. liii. 8. Compare § 7, below.}? \emph{God is a Spirit}\footnote{John iv. 24.}; He who is a Spirit hath spiritually begotten, as being incorporeal, an inscrutable and incomprehensible generation. The Son Himself says of the Father, \emph{The Lord said unto Me, Thou art My Son, to-day have I begotten Thee}\footnote{Ps. ii. 7.}. Now this \emph{to-day} is not recent, but eternal: a timeless to-day, before all ages. \emph{From the womb, before the morning star, have I begotten Thee}\footnote{Ps. cx. 3. “From the womb of the morning thou hast the dew of thy youth” (\textsc{R.V.}). There is a remarkable various reading in Codd. Roe, Casaub. \textgreek{Τό εἶ σύ, ἄχρονον καὶ ἀΐδιον· τὸ δὲ σήμερον πρόσφατον, ἀλλ᾽ οὐκ ἀΐδιον, οἰκειουμένου τοῦ Πατρὸς καὶ τὴν κάτω γέννησιν. Καὶ πάλιν λέγει· ᾽Εκ γαστρὸς πρὸ ἑωσφόρου γεγέννηκά σε· τοῦτο μόνον τῆς Θεότητος· Πίστευσον, κ.τ.λ} The words “\emph{Thou art My Son},” are thus referred to the eternal generation, and “\emph{This day}” to the birth in time: whereas in the received text, followed in our translation, \textgreek{σήμερον} refers to the timeless and eternal generation of the Son. The former interpretation of Ps. ii. 7 is found in many Fathers, as for example in Tertullian (\emph{adv. Prax.} vii. xi.), and Methodius (\emph{Conviv. Virg.} VIII. cap. ix.): “He says ‘Thou art,’ and not ‘Thou hast become,’ shewing that He had not recently attained to the position of Son.…But the expression, ‘This day have I begotten Thee,’ signifies that He willed that existing already before the ages in heaven He should also be begotten for the world, that is that He who was before unknown should be made known.’ The same interpretation was held by many Fathers, some referring \textgreek{σήμερον} to the Nativity, as Cyprian (\emph{adv. Judæos Testim}. ii. 8), others to the Baptism (Justin M. \emph{Dialog}. cap. lxxxviii.; Tertullian. \emph{adv. Marcion}. iv. 22). Athanasius (\emph{c. Arian.} iv. § 27), has a long discussion on the question whether Ps. cx. 3, \textgreek{ἐκ γαστρὸς πρὸ ἑωσφόρου γεγέννηκά σε}, refers to the eternal generation of the Son, or to His Nativity.}.

6. Believe thou therefore on \textsc{Jesus Christ}, \textsc{Son} of the living \textsc{God}, and a Son \textsc{Only-Begotten}, according to the Gospel which says, \emph{For God so loved the world, that He gave His Only-begotten Son, that whosoever believeth on Him should not perish, but have everlasting life}\footnote{John iii. 16.}. And again,\emph{He that believeth on the Son is not judged, but hath passed out of death into life}\footnote{Ib. iii. 18; v. 24.}. \emph{But he that believeth not the Son shall not see life, but the wrath of God abideth on him}\footnote{Ib. iii. 36.}. And John testified concerning Him, saying, \emph{And we beheld His glory, glory as of the only-begotten from the father,—full of grace and truth}\footnote{Ib. i. 14.}: at whom the devils trembled and said, \emph{Ah! what have we to do with Thee, Jesus, Thou Son of the living God}\footnote{Luke iv. 34.}.

7. He is then the Son of God by nature and not by adoption\footnote{\textgreek{φύσει καὶ οὐ θέσει}. Cf. § 2, note 4.}, begotten of the Father. \emph{And he that loveth Him that begat, loveth Him also that is begotten of Him}\footnote{1 John v. 1.}; but he that despiseth Him that is begotten casts back the insult upon Him who begat. And whenever thou hear of God begetting, sink not down in thought to bodily things, nor think of a corruptible generation, lest thou be guilty of impiety. \emph{God is a Spirit}\footnote{John iv. 24. Cf. § 5.}, His generation is spiritual: for bodies beget bodies, and for the generation of bodies time needs must intervene; but time intervenes not in the generation of the Son from the Father. And in our case what is begotten is begotten imperfect: but the Son of God was begotten perfect; for what He is now, that is He also from the beginning\footnote{\textgreek{γεγεννημένος ἀνάρχως}. Cf. § 5, note 4.}, begotten without beginning. We are begotten so as to pass from infantile ignorance to a state of reason: thy generation, O man, is imperfect, for thy growth is progressive. But think not that it is thus in His case, nor impute infirmity to Him who hath begotten. For if that which He begot was imperfect, and acquired its perfection in time, thou art imputing infirmity to Him who hath begotten; if so be, the Father did not bestow from the beginning that which, as thou sayest, time bestowed afterwards\footnote{\textgreek{ὃ χρόνος}. Bened. c. Codd. Roe, Casaub. Coisl. \textgreek{ὃ χρόνοις} Ottob. Mon. 1, 2. A. With the latter reading, the meaning will be—“if He did not bestow from the beginning, as thou sayest, what He bestowed in after times.” Cyril does not here address his auditor, but an imaginary opponent,—“O man.” \par Compare Athan. (\emph{de Synodis}, § 26).}.

8. Think not therefore that this generation is human, nor as Abraham begat Isaac. For in begetting Isaac, Abraham begat not what he would, but what another granted. But in God the Father’s begetting there is neither ignorance nor intermediate deliberation\footnote{The Arians appear to have made use of a dilemma: If God begat with will and purpose, these preceded the begetting, and so \textgreek{ἦν ποτε ὅτε οὐκ ἦν}, there was a time when the Son was not: if without will and purpose, then He begat in ignorance and of necessity. The answer is fully given by Athanasius (\emph{c. Arian}. iii. 58–67, pp. 425–431 in this Series).}. For to say that He knew not what He was begetting is the greatest impiety; and it is no less impious to say, that after deliberation in time He then became a Father. For God was not previously without a Son, and afterwards in time became a Father; but hath the Son eternally, having begotten Him not as men beget men, but as Himself only knoweth, who begat Him before all ages \textsc{Very God}.

9. For the Father being Very God begat the Son like unto Himself, Very God\footnote{Athanasius (\emph{ad Episcopos Ægypti}, § 13), referring to 1 John v. 20, \emph{This is the true} (\textgreek{ἁληθινός}) \emph{God}, writes: “But these men (the Arians), as if in contradiction to this, allege that Christ is not the true God, but that He is only called God, as are other creatures, in regard of His participation in the Divine nature.” Again (\emph{c. Arian}. iii. 9), “He gave us to know that of the true Father He is the true Offspring (\textgreek{ἀληθινὸν γέννημα}).}; not as teachers beget disciples, not as Paul says to some, \emph{For in Christ Jesus I begat you through the Gospel}\footnote{1 Cor. iv. 15.}. For in this case he who was not a son by nature became a son by discipleship, but in the former case He was a Son by nature, a true Son. Not as ye, who are to be illuminated, are now becoming sons of God: for ye also become sons, but by adoption of grace, as it is written, \emph{But as many as received Him, to them gave He the right to become children of God, even to them that believe on His name: which were begotten not of blood, nor of the will of the flesh, nor of the will of man, but of God}\footnote{John i. 12, 13.}. And we indeed are begotten of water and of the Spirit, but not thus was Christ begotten of the Father. For at the time of His Baptism addressing Him, and saying, \emph{This is My Son}\footnote{Matt. iii. 17.}, He did not say, “This has now become My Son,” but, \emph{This is My Son;} that He might make manifest, that even before the operation of Baptism He was a Son.

10. The Father begat the Son, not as among men mind begets word. For the mind is substantially existent in us; but the word when spoken is dispersed into the air and comes to an end\footnote{Compare Athanasius (\emph{de Sententiâ Dionysii}, § 23): “the mind creates the word, being manifested in it, and the word shews the mind, having originated therein.” Tertullian (\emph{adv. Prax}. vii.): “You will say what is a word but a voice and sound of the mouth, and (as the Grammarians teach) air when struck against, intelligible to the ear, but for the rest a sort of void, empty, and incorporeal thing.” Cf. Athan. (\emph{de Synodis}, § 12): \textgreek{ἀνυπόστατον}.}. But we know Christ to have been begotten not as a word pronounced\footnote{\textgreek{προφορικόν}. See Cat. iv. 8, note 9.}, but as a Word substantially existing\footnote{\textgreek{ἐνυπόστατον}. \emph{ibid}. So the Spirit is described in Cat. xvii. 5 “not uttered or breathed by the mouth and lips of the Father and the Son, nor dispersed into the air, but personally subsisting (\textgreek{ἐνυπόστατον}).”} and living; not spoken by the lips, and dispersed, but begotten of the Father eternally and ineffably, in substance\footnote{\textgreek{ἐν ὑποστάσει}.}. For, \emph{In the beginning was the Word, and the Word was with God, and the Word was God}\footnote{John i. 1.}, sitting at God’s right hand;—the Word understanding the Father’s will, and creating all things at His bidding: the Word, which came down and went up; for the word of utterance when spoken comes not down, nor goes up; the Word speaking and saying, \emph{The things which I have seen with My Father, these I speak}\footnote{John viii. 38.}: the Word possessed of power, and reigning over all things: for \emph{the Father hath committed all things unto the Son}\footnote{Matt. xi. 27; John v. 22.}.

11. The Father then begat Him not in such wise as any man could understand, but as Himself only knoweth. For we profess not to tell in what manner He begat Him, but we insist that it was not in this manner. And not we only are ignorant of the generation of the Son from the Father, but so is every created nature. \emph{Speak to the earth, if perchance it may teach thee}\footnote{Job xii. 8.}: and though thou inquire of all things which are upon the earth, they shall not be able to tell thee. For the earth cannot tell the substance of Him who is its own potter and fashioner. Nor is the earth alone ignorant, but the sun also\footnote{In saying that the earth, the sun, and the heavens know not their Maker, Cyril is simply using figurative language like that of the passage of Job just quoted. There is no reason to suppose that he accepted Origen’s theory (\emph{de Principiis}, II. cap. 7), that the heavenly bodies are living and rational beings, capable of sin.}: for the sun was created on the fourth day, without knowing what had been made in the three days before him; and he who knows not the things made in the three days before him, cannot tell forth the Maker Himself. Heaven will not declare this: for at the Father’s bidding \emph{the heaven also was like smoke established}\footnote{Isa. li. 6: \emph{the heavens shall vanish away like smoke}.} by Christ. Nor shall \emph{the heaven of heavens} declare this, \emph{nor the waters which are above the heavens}\footnote{Ps. cxlviii. 4.}. Why then art thou cast down, O man, at being ignorant of that which even the heavens know not? Nay, not only are the heavens ignorant of this generation, but also every angelic nature. For if any one should ascend, were it possible, into the first heaven, and perceiving the ranks of the Angels there should approach and ask them how God begat His own Son, they would say perhaps, “We have above us beings greater and higher; ask them.” Go up to the second heaven and the third; attain, if thou canst, to Thrones, and Dominions, and Principalities, and Powers: and even if any one should reach them, which is impossible, they also would decline the explanation, for they know it not.

12. For my part, I have ever wondered at the curiosity of the bold men, who by their imagined reverence fall into impiety. For though they know nothing of Thrones, and Dominions, and Principalities, and Powers, the workmanship of Christ, they attempt to scrutinise their Creator Himself. Tell me first, O most daring man, wherein does Throne differ from Dominion, and then scrutinise what pertains to Christ. Tell me what is a Principality, and what a Power, and what a Virtue, and what an Angel: and then search out their Creator, \emph{for all things were made by Him}\footnote{John i. 3.}. But thou wilt not, or thou canst not ask Thrones or Dominions. What else is there that \emph{knoweth the deep things of God\footnote{1 Cor. ii. 10, 11.}}, save only the Holy Ghost, who spake the Divine Scriptures? But not even the Holy Ghost Himself has spoken in the Scriptures concerning the generation of the Son from the Father. Why then dost thou busy thyself about things which not even the Holy Ghost has written in the Scriptures? Thou that knowest not the things which are written, busiest thou thyself about the things which are not written? There are many questions in the Divine Scriptures; what is written we comprehend not, why do we busy ourselves about what is not written? It is sufficient for us to know that God hath begotten One Only Son.

13. Be not ashamed to confess thine ignorance, since thou sharest ignorance with Angels. Only He who begat knoweth Him who was begotten, and He who is begotten of Him knoweth Him who begat. He who begat knoweth what He begat: and the Scriptures also testify that He who was begotten is God\footnote{I have followed the reading of Codd. Coisl. Roe, Casaub. Mon. A., which is approved though not adopted by the Benedictine Editor. The common text is manifestly interpolated: “And the Holy Spirit of God testifies in the Scriptures, that He who was begotten without beginning is God. \emph{For what man knoweth, \&c}.” This insertion of 1 Cor. ii. 11 interrupts the argument, and is a useless repetition of the allusion to the same passage in § 12.}. For \emph{as the Father hath life in Himself, so also hath He given to the Son to have life in Himself}\footnote{John v. 26.}; and, \emph{that all men should honour the Son, even as they honour the Father}\footnote{Ib. v. 23.}; and, \emph{as the Father quickeneth} whom He will, even so the Son quickeneth whom He will\footnote{Ib. v. 21.}. Neither He who begat suffered any loss, nor is anything lacking to Him who was begotten (I know that I have said these things many times, but it is for your safety that they are said so often): neither has He who begat, a Father, nor He who was begotten, a brother. Neither was He who begat changed into the Son\footnote{See iv. 8, note 8, on the Sabellian doctrine, and Athanas. (\emph{de Synodis}, § 16, note 10 in this series).}, nor did He who was begotten become the Father\footnote{The doctrine of Sabellius might be expressed in two forms, either the Father became the Son, or the Son became the Father. Both forms are here denied. The Jerusalem Editor thinks there is an allusion to the Arian argument mentioned by Athanasius (\emph{c. Arian. Or}. I. cap. vi. 22): “If the Son is the Father’s offspring and Image, and is like in all things to the Father, then it necessarily holds that as He is begotten so He begets, and He too becomes father of a son.” But the close connexion of the two clauses is in favour of the reference to the Sabellian \textgreek{υἱοπατορία}.}. Of One Only Father there is One Only-begotten Son: neither two Unbegotten\footnote{\textgreek{ἀγέννητοι}. The context shews that this, not \textgreek{ἀγένητοι}, is here the right form. Athanasius seems to have used \textgreek{ἀγέννητος} in both senses “Un-begotten,” as here, and “unoriginate.” Thus (\emph{c. Arian. Or}. i. cap. ix. § 30) he says of the Arians: “Their further question ‘whether the Unoriginate be one or two,’ shews how false are their views.” Compare Bp. Lightfoot’s Excursus on Ignatius, \emph{Ephes}. § 7, and Mr. Robertson’s notes on Athanasius in this Series.}, nor two Only-begotten; but One Father, Unbegotten (for He is Unbegotten who hath no father); and One Son, eternally begotten of the Father; begotten not in time, but before all ages; not increased by advancement, but begotten that which He now is.

14. We believe then \textsc{In the Only-Begotten Son of God, Who Was Begotten of the Father Very God}. For the True God begetteth not a false god, as we have said, nor did He deliberate and afterwards beget\footnote{See above, § 8, note 3.}; but He begat eternally, and much more swiftly than our words or thoughts: for we speaking in time, consume time; but in the case of the Divine Power, the generation is timeless. And as I have often said, He did not bring forth the Son from non-existence into being, nor take the non-existent into sonship\footnote{Athan. (\emph{c. Arian}. I. ix. 31) “speaking against the Lord, ‘He is of nothing,’ and ‘He was not before His generation.’”}: but the Father, being Eternal, eternally and ineffably begat One Only Son, who has no brother. Nor are there two first principles; but the Father is \emph{the head of the Son}\footnote{1 Cor. xi. 3.}; the beginning is One. For the Father begot the Son \textsc{Very God}, called Emmanuel; and Emmanuel \emph{being interpreted is, God with us}\footnote{Matt. i. 23.}.

15. And wouldest thou know that He who was begotten of the Father, and afterwards became man, is God? Hear the Prophet saying, \emph{This is our God, none other shall be accounted of in comparison with Him. He hath found out every way of knowledge, and given it to Jacob His servant, and to Israel His beloved. Afterwards He was seen on earth, and conversed among men}\footnote{Baruch iii. 35–37. The last verse was understood by Cyril, as by many of the Greek and Latin Fathers, to be a prophecy of the Incarnation: but in reality it refers to “knowledge” (\textgreek{ἐπιστήμη}, \emph{v}. 36), and should be translated “she was seen upon earth.” See notes on the passage in the \emph{Speaker’s Commentary}.}. Seest thou herein God become man, after the giving of the law by Moses? Hear also a second testimony to Christ’s Deity, that which has just now been read, \emph{Thy throne, O God, is for ever and ever}\footnote{Heb. i. 8.}. For lest, because of His presence here in the flesh, He should be thought to have been advanced after this to the Godhead, the Scripture says plainly, \emph{Therefore God, even Thy God, hath anointed Thee with the oil of gladness above Thy fellows}\footnote{Ib. i. 9. See x. 14, note 9.}. Seest thou Christ as God anointed by God the Father?

16. Wouldest thou receive yet a third testimony to Christ’s Godhead? Hear Esaias saying, \emph{Egypt hath laboured, and the merchandise of Ethiopia:} and soon after, \emph{In Thee shall they make supplication, because God is in Thee, and there is no God save Thee. For Thou art God, and we knew it not, the God of Israel, the Saviour}\footnote{Isa. xlv. 14, 15: “They shall make supplication unto thee, saying, surely God is in thee.” The words are addressed to Jerusalem as the city of God. Cyril applies them to the Son, misled by the Septuagint.}. Thou seest that the Son is God, having in Himself God the Father: saying almost the very same which He has said in the Gospels: \emph{The Father is in Me, and I am in the Father}\footnote{John xiv. 11.}. He says not, I am the Father, but \emph{the Father is in Me, and I am in the Father}. And again He said not, I and the Father am\footnote{Athanasius (\emph{c. Arian. Or}. iv. § 9), arguing for the \textgreek{ὁμοούσιον} says: “These are two, because there is Father and Son, that is the Word; and one, because one God. For if this is not so, He would have said, I am the Father, or, I and the Father am.”} one, but, \emph{I and the Father am one}, that we should neither separate them, nor make a confusion of Son-Father\footnote{See iv. 8, notes 7 and 8.}. One they are because of the dignity pertaining to the Godhead, since God begat God. One in respect of their kingdom; for the Father reigns not over these, and the Son over those, lifting Himself up against His Father like Absalom: but the kingdom of the Father is likewise the kingdom of the Son. One they are, because there is no discord nor division between them: for what things the Father willeth, the Son willeth the same. One, because the creative works of Christ are no other than the Father’s; for the creation of all things is one, the Father having made them through the Son: \emph{For He spake, and they were made; He commanded, and they were created}, saith the Psalmist\footnote{Psa. xxxiii. 9; cxlviii. 5. S. Cyril explains the creative “Fiat” in Gen. i. as addressed by the Father to the Son.}. For He who speaks, speaks to one who hears: and He who commands, gives His commandment to one who is present with Him.

17. The Son then is \textsc{Very God}, having the Father in Himself, not changed into the Father; for the Father was not made man, but the Son. For let the truth be freely spoken\footnote{We learn from Socrates (\emph{Eccl. Hist}. I. 24), that after the Nicene Council “those who objected to the word \textgreek{ὁμοούσιος} conceived that those who approved it favoured the opinion of Sabellius.” Marcellus of Ancyra, who was deposed on a charge of Sabellianism, and who did not in fact make clear the distinct personality of the Son, had been warmly supported by the friends of Athanasius. Cyril apparently fears to incur their censure, if he too strongly condemned the Sabellian view.}. The Father suffered not for us, but the Father sent Him who suffered. Neither let us say, There was a time when the Son was not; nor let us admit a Son who is the Father\footnote{Cyril here rejects both the opposite errors, Arianism, “There was a time when the Son was not,” and Sabellianism, “a Son who is the Father.”}: but let us walk in the king’s highway; let us turn aside neither on the left hand nor on the right. Neither from thinking to honour the Son, let us call Him the Father; nor from thinking to honour the Father, imagine the Son to be some one of the creatures. But let One Father be worshipped through One Son, and let not their worship be separated. Let One Son be proclaimed, sitting at the right hand of the Father before all ages: sharing His throne not by advancement in time after His Passion, but by eternal possession.

18. \emph{He who hath seen the Son, hath seen the Father}\footnote{John xiv. 9.}: for in all things the Son is like to Him who begat Him\footnote{See above, § 4, note 9.}; begotten Life of Life and Light of Light, Power of Power, God of God; and the characteristics of the Godhead are unchangeable\footnote{\textgreek{ἀπαράλλακτοι}. The word was used by the Orthodox Bishops at Nicæa, who said that “the Word must be described as the True power and Image of the Father, in all things like the Father and Himself incapable of change.” See the notes of Dr. Newman and Mr. Robertson on Athanasius (\emph{de Decretis}, § 20).} in the Son; and he who is counted worthy to behold Godhead in the Son, attains to the fruition of the Father. This is not my word, but that of the Only-begotten: \emph{Have I been so long time with you, and hast thou not known Me, Philip? He that hath seen Me, hath seen the Father}\footnote{John xiv. 9.}. And to be brief, let us neither separate them, nor make a confusion\footnote{See iv. 8, note 8.}: neither say thou ever that the Son is foreign to the Father, nor admit those who say that the Father is at one time Father, and at another Son: for these are strange and impious statements, and not the doctrines of the Church. But the Father having begotten the Son, remained the Father and is not changed. He begat Wisdom, yet lost not wisdom Himself; and begat Power, yet became not weak: He begat God, but lost not His own Godhead: and neither did He lose anything Himself by diminution or change; nor has He who was begotten any thing wanting. Perfect is He who begat, Perfect that which was begotten: God was He who begat, God He who was begotten; God of all Himself, yet entitling the Father His own God. For He is not ashamed to say, I \emph{ascend unto My Father and your Father, and to My God and your God}\footnote{John xx. 17.}.

19. But lest thou shouldest think that He is in a like sense Father of the Son and of the creatures, Christ drew a distinction in what follows. For He said not, “I ascend to our Father,” lest the creatures should be made fellows of the Only-begotten; but He said, \emph{My Father and your Father}; in one way Mine, by nature; in another yours, by adoption. And again, \emph{to my God and your God}, in one way Mine, as His true and Only-begotten Son, and in another way yours, as His workmanship\footnote{Compare Cat. vii. 7. The Jerusalem Editor observes that the expression “My God” is understood by the Fathers generally as spoken by Christ in reference to His human nature, but Cyril applies this, as well as the other expression “My Father,” to the Divine nature. So Hilary (\emph{de Trinit}. iv. 53): “idcirco Deus ejus est, quia ex eo natus in Deum est.” Compare Epiphanius (\emph{Hær.} lxix. 55).}. The Son of God then is \textsc{Very God}, ineffably begotten before all ages (for I say the same things often to you, that it may be graven upon your mind). This also believe, that God has a Son: but about the manner be not curious, for by searching thou wilt not find. Exalt not thyself, lest thou fall: \emph{think upon those things only which have been commanded thee}\footnote{Ecclus. iii. 22.}. Tell me first what He is who begat, and then learn that which He begat; but if thou canst not conceive the nature of Him who hath begotten, search not curiously into the manner of that which is begotten.

20. For godliness it sufficeth thee to know, as we have said, that God hath One Only Son, One naturally begotten; who began not His being when He was born in Bethlehem, but \textsc{Before All Ages}. For hear the Prophet Micah saying, \emph{And thou, Bethlehem, house of Ephrata, art little to be among the thousands of Judah. Out of thee shall come forth unto Me a Ruler, who shall feed My people Israel: and His goings forth are from the beginning, from days of eternity}\footnote{Micah v. 2; on the various readings \textgreek{ὀλγιοστὸς εἶ, μὴ ὀλ, εἶ οὐκ ὀλ. εἶ}, found in the \textsc{mss.} of Cyril, see the Commentaries on the quotation of the passage in Matt. ii. 6.}. Think not then of Him who is now come forth out of Bethlehem\footnote{Codd. Roe, Casaub. have a different reading—“Think not then of His having now been born in Bethlehem, and (nor) suppose Him as the Son of Man to be altogether recent, but worship, \&c.” This is rightly regarded by the Benedictine and other Editors as an interpolation intended to avoid the apparent tendency of Cyril’s language in the received text to separate the Virgin’s Son from the Eternal Word. Had Cyril so written after the Nestorian controversy arose, he would have appeared to favour the Nestorian formula that “Mary did not give birth to the Deity.” Compare Swainson (\emph{Nicene Creed}, Ch. ix. § 7.) What Cyril really means is that we are not to think of Christ simply as man, but to worship Him as God.}, but worship Him who was eternally begotten of the Father. Suffer none to speak of a beginning of the Son in time, but as a timeless Beginning acknowledge the Father. For the Father is the Beginning of the Son, timeless, incomprehensible, without beginning\footnote{Compare § 4, note 3.}. The fountain of the river of righteousness, even of the Only-begotten, is the Father, who begat Him as Himself only knoweth. And wouldest thou know that our Lord Jesus Christ is King Eternal? Hear Him again saying, \emph{Your father Abraham rejoiced to see My day, and he saw it, and was glad}\footnote{John viii. 56.}. And then, when the Jews received this hardly, He says what to them was still harder, \emph{Before Abraham was, I am}\footnote{Ib. viii. 58.}. And again He saith to the Father, \emph{And now, Father, glorify Thou Me with Thine own self, with the glory which I had with Thee before the world was}\footnote{Ib. xvii. 5.}. He says plainly, “before the world was, I had the glory which is with Thee.” And again when He says, \emph{For Thou lovedst Me before the foundation of the world}\footnote{John xvii. 24.}, He plainly declares, “The glory which I have with thee is from eternity.”

21. We believe then \textsc{In One Lord Jesus Christ, the Only-Begotten Son of God, Begotten of His Father Very God Before All Worlds, by Whom All Things Were Made}. For \emph{whether they be thrones, or dominions, or principalities, or powers, all things were made through Him}\footnote{Col. i. 16.}, and of things created none is exempted from His authority. Silenced be every heresy which brings in different creators and makers of the world; silenced the tongue which blasphemes the Christ the Son of God; let them be silenced who say that the sun is the Christ, for He is the sun’s Creator, not the sun which we see\footnote{Compare Cat. vi. 13, and xv. 3: “Here let converts from the Manichees gain instruction, and no longer make those lights their gods; nor impiously think that this sun which shall be darkened is Christ.”}. Silenced be they who say that the world is the workmanship of Angels\footnote{The creation of the world was ascribed to Angels by the Gnostics generally, \emph{e.g.} by Simon Magus (Irenæus, \emph{adv. Hæres}. I. xxiii. § 2), Menander (\emph{ibid}. § 5), Saturninus (\emph{ibid}. xxiv. 1), Basilides (\emph{ibid}. § 3), Carpocrates (\emph{ibid}. xxv. 1).}, who wish to steal away the dignity of the Only-begotten. For whether visible or invisible, whether thrones or dominions, or anything that is named, all things were made by Christ. He reigns over the things which have been made by Him, not having seized another’s spoils, but reigning over His own workmanship, even as the Evangelist John has said, \emph{All things were made by Him, and without Him was not anything made}\footnote{John i. 3.}. All things were made by Him, the Father working by the Son.

22. I wish to give also a certain illustration of what I am saying, but I know that it is feeble; for of things visible what can be an exact illustration of the Divine Power? But nevertheless as feeble be it spoken by the feeble to the feeble. For just as any king, whose son was a king, if he wished to form a city, might suggest to his son, his partner in the kingdom, the form of the city, and he having received the pattern, brings the design to completion; so, when the Father wished to form all things, the Son created all things at the Father’s bidding, that the act of bidding might secure to the Father His absolute authority\footnote{On the doctrine of Creation by the Son as held by Cyril, see the reference to the Introduction in the Index, \emph{Creation}.}, and yet the Son in turn might have authority over His own workmanship, and neither the Father be separated from the lordship over His own works, nor the Son rule over things created by others, but by Himself. For, as I have said, Angels did not create the world, but the Only-begotten Son, begotten, as I have said, before all ages, \textsc{By Whom All Things Were Made}, nothing having been excepted from His creation. And let this suffice to have been spoken by us so far, by the grace of Christ.

23. But let us now recur to our profession of the Faith, and so for the present finish our discourse. Christ made all things, whether thou speak of Angels, or Archangels, of Dominions, or Thrones. Not that the Father wanted strength to create the works Himself, but because He willed that the Son should reign over His own workmanship, God Himself giving Him the design of the things to be made. For honouring His own Father the Only-begotten saith, \emph{The Son can do nothing of Himself, but what He seeth the Father do; for what things soever He doeth, these also doeth the Son likewise}\footnote{John v. 19.}. And again, \emph{My Father worketh hitherto, and I work}\footnote{Ib. v. 17.}, there being no opposition in those who work. \emph{For all Mine are Thine, and Thine are Mine}, saith the Lord in the Gospels\footnote{Ib. xvii. 10.}. And this we may certainly know from the Old and New Testaments. For He who said, \emph{Let us make man in our image and after our likeness}\footnote{Gen. i. 26.}, was certainly speaking to some one present. But clearest of all are the Psalmist’s words, \emph{He spake and they were made; He commanded, and they were created}\footnote{Ps. cxlviii. 5.}, as if the Father commanded and spake, and the Son made all things at the Father’s bidding. And this Job said mystically, \emph{Which alone spread out the heaven, and walketh upon the sea as on firm ground}\footnote{Job ix. 8.}; signifying to those who understand that He who when present here walked upon the sea is also He who aforetime made the heavens. And again the Lord saith, \emph{Or didst Thou take earth, and fashion clay into a living being}\footnote{Ib. xxxviii. 14.}? then afterwards, \emph{Are the gates of death opened to Thee through fear, and did the door-keepers of hell shudder at sight of Thee}\footnote{Ib. xxxviii. 17.}? thus signifying that He who through loving-kindness descended into hell, also in the beginning made man out of clay.

24. Christ then is the Only-begotten Son of God, and Maker of the world. For \emph{He was in the world, and the world was made by Him;} and \emph{He came unto His own}, as the Gospel teaches us\footnote{John i. 10, 11.}. And not only of the things which are seen, but also of the things which are not seen, is Christ the Maker at the Father’s bidding. For \emph{in Him}, according to the Apostle, \emph{were all things created that are in the heavens, and that are upon the earth, things visible and invisible, whether thrones, or dominions, or principalities, or powers; all things have been created by Him and for Him; and He is before all, and} \emph{in Him all things consist}\footnote{Col. i. 16, 17.}. Even if thou speak of the worlds, of these also Jesus Christ is the Maker by the Father’s bidding. For \emph{in these last days God spake unto us by His Son, whom He appointed heir of all things, by whom also He made the worlds\footnote{Heb. i. 2.}}. To whom be the glory, honour, might, now and ever, and world without end. Amen.

\xchapter{Lecture XII. On the words Incarnate, and Made Man.}

\textsc{Isaiah vii. 10–14}

\emph{“And the Lord spoke again unto Ahaz, saying, Ask thee a sign, \&c.:}” and “\emph{Behold! a virgin shall conceive, and bear a son, and shall call His name Emmanuel,} \&c.”

1. \textsc{Nurslings} of purity and disciples of chastity, raise we our hymn to the Virgin-born God\footnote{This passage supplies a complete answer to the suspicion of a quasi-Nestorian tendency referred to in note 6, on xi. 20. See x. 19, note 2, on the title \textgreek{Θεοτόκος}.} with lips full of purity. Deemed\footnote{The Present Participle (\textgreek{καταξιούμενοι}) means that the Candidates for Baptism were already on the way to be admitted to Holy Communion. Compare Cat. i. 1, where the same Candidates are addressed as “partakers of the mysteries of Christ, as yet by calling only, but ere long by grace also.”} worthy to partake of the flesh of the Spiritual Lamb\footnote{Aubertin remarks on this passage that “this spiritual Lamb, consisting of head and feet, can be received only by the spiritual mouth.” This explanation, however true in itself, cannot fairly be held to express fully the meaning of Cyril. See the section of the Introduction referred to in the Index, “Eucharist.”}, let us take the head together with the feet\footnote{Ex. xii. 9: \emph{the head with the feet}. The same figurative interpretation is given by Eusebius (\emph{Eccl. Hist}. I. ii. § 1): “In Christ there is a twofold nature; and the one—in so far as He is thought of as God—resembles the head of the body, while the other may be compared with the feet,—in so far as He, for the sake of our salvation, put on human nature with the same passions as our own.”}, the Deity being understood as the head, and the Manhood taken as the feet. Hearers of the Holy Gospels, let us listen to John the Divine\footnote{\textgreek{᾽Ιωάννῃ τῷ Θεολόγω}. The title is given to Moses by Philo Judæus (\emph{Vita Mos}. III. § 11), to Prophets by Eusebius (\emph{Demostr. Evang}. ii. 9), to Apostles by Athanasius (\emph{de Incarn}. § 10: \textgreek{τῶν αὐτοῦ τοῦ Σωτῆρος θεολόγων ἀνδρῶν}), and especially to St. John, because the chief purpose of his Gospel was to set forth the Deity of Christ. See note on Revel. i. 1, in \emph{Speaker’s Commentary}, and Suicer, \emph{Thesaurus}, \textgreek{Θεολόγος}.}. For he who said, \emph{In the beginning was the Word, and the Word was with God, and the Word was God}\footnote{John i. 1.}, went on to say, \emph{and the Word was made flesh}\footnote{Ib. i. 14.}. For neither is it holy to worship the mere man, nor religious to say that He is God only without the Manhood. For if Christ is God, as indeed He is, but took not human nature upon Him, we are strangers to salvation. Let us then worship Him as God, but believe that He also was made Man. For neither is there any profit in calling Him man without Godhead nor any salvation in refusing to confess the Manhood together with the Godhead. Let us confess the presence of Him who is both King and Physician. For Jesus the King when about to become our Physician, \emph{girded Himself with the linen} of humanity\footnote{Ib. xiii. 4.}, and healed that which was sick. The perfect Teacher of babes\footnote{Rom. ii. 20.} became a babe among babes, that He might give wisdom to the foolish. The Bread of heaven came down on earth\footnote{John vi. 32, 33, 50.} that He might feed the hungry.

2. But the sons of the Jews by setting at nought Him that came, and looking for him who cometh in wickedness, rejected the true Messiah, and wait for the deceiver, themselves deceived; herein also the Saviour being found true, who said, \emph{I am come in My Father’s name, and ye receive Me not: but if another shall come in his own name, him ye will receive}\footnote{Ib. v. 43. Cf. 2 John 7.}. It is well also to put a question to the Jews. Is the Prophet Esaias, who saith that Emmanuel shall be born of a virgin, true or false\footnote{Isa. vii. 14.}? For if they charge him with falsehood, no wonder: for their custom is not only to charge with falsehood, but also to stone the Prophets. But if the Prophet is true, point to the Emmanuel, and say, Whether is He who is to come, for whom ye are looking, to be born of a virgin or not? For if He is not to be born of a virgin, ye accuse the Prophet of falsehood: but if in Him that is to come ye expect this, why do ye reject that which has come to pass already?

3. Let the Jews, then, be led astray, since they so will: but let the Church of God be glorified. For we receive God the Word made Man in truth, not, as heretics say\footnote{Carpocrates, Cerinthus, the Ebionites, \&c. See Irenæus (\emph{Hær.} I. xxv. § 1; xxvi. §§ 1, 2).}, of the will of man and woman, but of \textsc{The Virgin and the Holy Ghost\footnote{Dr. Swainson (\emph{Creeds}, Chap. vii. § 7), speaking of the Creed of Cyril of Jerusalem, says that “the words \textgreek{σαρκωθέντα καὶ ἐνανθρωπήσαντα} are found in it, but no reference whatever is made to the birth from the Virgin.” The present passage, and that in Cat. iv. § 9, “begotten of the Holy Virgin and the Holy Ghost,” seems to shew that such a clause formed part of the Creed which Cyril was expounding. The genuineness of both passages is attested by all the \textsc{mss.} and Dr. Swainson was mistaken in charging the Editors of the Oxford Translation with having omitted to “mention that Touttée was himself doubtful as to the words within the brackets” [\textgreek{ἐκ Παρθένου καὶ Πνεύματος ῾Αγίου}]. The brackets are added by Dr. Swainson himself, and Touttée had no doubt of the genuineness of the words: on the contrary he believed them to be part of the Creed itself. His note is as follows: “The words \emph{of the Virgin and Holy Ghost} I have caused to be printed in larger letters as if taken from the Symbol: although they are wanting in the Title of this Lecture and in § 13, where the third Article of the Creed is referred to. But they are read in nearly all the Latin and Greek Symbols, and are referred to in Cat. iv. § 9.”}} according to the Gospel, \textsc{Made Man\footnote{\textgreek{ἐνανθρωπήσαντα}. The word occurs in the true Nicene formula, where, as Dr. Swainson thinks, it is “scarcely ambiguous,” “it is defective.” Both the Verb and the Substantive \textgreek{ἐνανθρώπησις} are constantly used by Athanasius to denote the Incarnation in a perfectly general way, without any indication of ambiguity or defect. In the Creed proposed by Eusebius of Cæsarea instead of \textgreek{ἐνανθρωπήσαντα} we find \textgreek{ἐν ἀνθρώποις πολιτευσάμενον}; and in the \emph{Expositio Fidei} ascribed to Athanasius, but of somewhat doubtful authenticity, the Incarnation is described thus \textgreek{ἐκ τῆς ἀχράντου παρθένου Μαρίας τὸν ἡμέτερον ἀνείληφεν ἄνθρωπον Χριστὸν ᾽Ιησοῦν}. In the Apollinarian controversy the attempt was made to interpret \textgreek{ἐνηνθρώπησεν} as meaning not that “He became Man,” but that “He assumed a man,” \emph{i.e.} that “the man was first formed and then assumed” (Gregory, \emph{Epist. ad Cledon}, quoted by Swainson, p. 83), or else merely that “He dwelt among men.” But the context of the passages in which Cyril uses the word (iv. 9; xii. 3) clearly shews that he employed it in the perfectly orthodox sense which it has in the Nicene Formula and in Athanasius.}}, not in seeming but in truth. And that He was truly Man made of the Virgin, wait for the proper time of instruction in this Lecture, and thou shalt receive the proofs\footnote{See below, § 21 ff. Cyril means that the direct proof cannot be given at once, because there are many errors to be set aside first. Compare the end of § 4.}: for the error of the heretics is manifold. And some have said that He has not been born at all of a virgin\footnote{See Cat. iv. 9, notes 3, 4.}: others that He has been born, not of a virgin, but of a wife dwelling with a husband. Others say that the Christ is not God made Man, but a man made God\footnote{Athanasius (\emph{contra Arian. Or.} I. § 9) quotes as from Arius, \emph{Thalia}, “Christ is not Very God, but He, as others, was made God (\textgreek{ἐθεοποιήθη} ) by participation.” The Eusebians in the Confession of Faith called Macrostichos (\textsc{a.d.} 344) condemned this view as being held by the disciples of Paul of Samosata, “who say that after the incarnation He was by advance made God, from being made by nature a mere man.” The orthodox use of the word \textgreek{Θεοποιεῖσθαι} is seen in Athan. \emph{de Incarnat.} § 54: \textgreek{αὐτὸς ἐνηνθρώπησεν, ἵνα ἡμεῖς θεοποιηθῶμεν}.}. For they dared to say that not He—the pre-existent Word—was made Man; but a certain man was by advancement crowned.

4. But remember thou what was said yesterday concerning His Godhead. Believe that He the Only-begotten Son of God—He Himself was again begotten of a Virgin. Believe the Evangelist John when he says, \emph{And the Word was made flesh, and dwelt among us}\footnote{John i. 14.}. For the Word is eternal, \textsc{Begotten of the Father Before All Worlds}: but the flesh He took on Him recently for our sake. Many contradict this, and say: “What cause was there so great, for God to come down into humanity? And, is it at all God’s nature to hold intercourse with men? And, is it possible for a virgin to bear, without man?” Since then there is much controversy, and the battle has many forms, come, let us by the grace of Christ, and the prayers of those who are present, resolve each question.

5. And first let us inquire for what cause Jesus came down. Now mind not my argumentations, for perhaps thou mayest be misled but unless thou receive testimony of the Prophets on each matter, believe not what I say: unless thou learn from the Holy Scriptures concerning the Virgin, and the place, the time, and the manner, \emph{receive not testimony from man}\footnote{John v. 34.}. For one who at present thus teaches may possibly be suspected: but what man of sense will suspect one that prophesied a thousand and more years beforehand? If then thou seekest the cause of Christ’s coming, go back to the first book of the Scriptures. In six days God made the world: but the world was for man. The sun however resplendent with bright beams, yet was made to give light to man, yea, and all living creatures were formed to serve us: herbs and trees were created for our enjoyment. All the works of creation were good, but none of these was an image of God, save man only. The sun was formed by a mere command, but man by God’s hands: \emph{Let us make man after our image, and after our likeness}\footnote{Gen. i. 26.}. A wooden image of an earthly king is held in honour; how much more a rational image of God?

But when this the greatest of the works of creation was disporting himself in Paradise, the envy of the Devil cast him out. The enemy was rejoicing over the fall of him whom he had envied: wouldest thou have had the enemy continue to rejoice? Not daring to accost the man because of his strength, he accosted as being weaker the woman, still a virgin: for it was after the expulsion from Paradise that \emph{Adam knew his wife}\footnote{Ib. iv. 1.}.

6. Cain and Abel succeeded in the second generation of mankind: and Cain was the first murderer. Afterwards a deluge was poured abroad because of the great wickedness of men: fire came down from heaven upon the people of Sodom because of their transgression. After a time God chose out Israel: but Israel also turned aside, and the chosen race was wounded. For while Moses stood before God in the mount, the people were worshipping a calf instead of God. In the lifetime of Moses, the law-giver who had said, \emph{Thou shalt not commit adultery}, a man dared to enter a place of harlotry and transgress\footnote{Numb. xxv. 6.}. After Moses, Prophets were sent to cure Israel: but in their healing office they lamented that they were not able to overcome the disease, so that one of them says, \emph{Woe is me! for the godly man is perished out of the earth, and there is none that doeth right among men}\footnote{Micah vii. 2.}: and again, \emph{They are all gone out of the way, they are together became unprofitable; there is none that doeth good, no, not one}\footnote{Ps. xiv. 3; Rom. iii. 12.}: and again, \emph{Cursing and stealing, and adultery, and murder are poured out upon the land}\footnote{Hosea iv. 2.}. \emph{Their sons and their daughters} \emph{they sacrificed unto devils}\footnote{Ps. cvi. 37.}. \emph{They used auguries, and enchantments, and divinations}\footnote{2 Chron. xxxiii. 6.}. \emph{And again, they fastened their garments with cords, and made hangings attached to the altar}\footnote{Amos ii. 8: \emph{they lay themselves down beside every altar upon clothes taken in pledge} (R.V.).}.

7. Very great was the wound of man’s nature; \emph{from the feet to the head there was no soundness in it;} none could apply \emph{mollifying ointment, neither oil, nor bandages}\footnote{Isa. i. 6.}. Then bewailing and wearying themselves, the Prophets said, \emph{Who shall give salvation out of Sion}\footnote{Ps. xiv. 7.}? And again, \emph{Let Thy hand be upon the man of Thy right hand, and upon the son of man whom Thou madest strong for Thyself: so will not we go back from Thee}\footnote{Ib. lxxx. 17, 18.}. And another of the Prophets entreated, saying, \emph{Bow the heavens, O Lord and come down}\footnote{Ps. cxliv. 5.}. The wounds of man’s nature pass our healing. \emph{They slew Thy Prophets, and cast down Thine altars}\footnote{1 Kings xix. 10.}. The evil is irretrievable by us, and needs thee to retrieve it.

8. The Lord heard the prayer of the Prophets. The Father disregarded not the perishing of our race; He sent forth His Son, the Lord from heaven, as healer: and one of the Prophets saith, \emph{The Lord whom ye seek, cometh, and shall suddenly come}\footnote{Mal. iii. 1.}. Whither? \emph{The Lord} shall come \emph{to His own temple}, where ye stoned Him. Then another of the Prophets, on hearing this, saith to him: In speaking of the salvation of God, speakest thou quietly? In preaching the good tidings of God’s coming for salvation, speakest thou in secret? \emph{O thou that bringest good tidings to Zion, get thee up into the high mountain. Speak to the cities of Judah}. What am I to speak? \emph{Behold our God! Behold! the Lord cometh with strength}\footnote{Isa. xl. 9, 10.}! Again the Lord Himself saith, \emph{Behold! I come, and I will dwell in the midst of thee, saith the Lord. And many nations shall flee unto the Lord}\footnote{Zech. ii. 10, 11.}. The Israelites rejected salvation through Me: \emph{I come to gather all nations and tongues}\footnote{Isa. lxvi. 18.}. For \emph{He came to His own and His own received Him not}\footnote{John i. 11.}. Thou comest and what dost Thou bestow on the nations? \emph{I come to gather all nations, and I will leave on them a sign}\footnote{Isa. lxvi. 19, a passage interpreted by the Fathers of the sign of the Cross. Eusebius (\emph{Demonstr. Evang.} vi. 25): “Who, on seeing that all who have believed in Christ use as a seal the symbol of salvation, would not reasonably be astonished at hearing the Lord’s saying of old time, \emph{And they shall come, and see My glory, and I will leave a sign upon them?}” Cf. Cat. iv. 14; xiii. 36.}. For from My conflict upon the Cross I give to each of My soldiers a royal seal to bear upon his forehead. Another also of the Prophets said, \emph{He bowed the heavens also, and came down; and darkness was under His feet}\footnote{Ps. xviii. 9. The “feet,” interpreted allegorically, mean the Humanity, and the “darkness” the mystery of the Incarnation. See Euseb. \emph{Demonstr. Evang}. vi. 1, § 2.}. For His coming down from heaven was not known by men.

9. Afterwards Solomon hearing his father David speak these things, built a wondrous house, and foreseeing Him who was to come into it, said in astonishment, \emph{Will God in very deed dwell with men on the earth}\footnote{1 Kings viii. 27; 2 Chron. vi. 18.}? Yea, saith David by anticipation in the Psalm inscribed \emph{For Solomon}, wherein is this, \emph{He shall come down like rain into a fleece}\footnote{Ps. lxxii. Title, and v. 6.}: \emph{rain}, because of His heavenly nature, and into a fleece, because of His humanity. For rain, coming down into a fleece, comes down noiselessly: so that the Magi, not knowing the mystery of the Nativity, say, \emph{Where is He that is born King of the Jews}\footnote{Matt. ii. 2.}? and Herod being troubled inquired concerning Him who was born, and said, \emph{Where is the Christ to be born}\footnote{Ib. ii. 4.}?

10. But who is this that cometh down? He says in what follows, \emph{And with the sun He endureth, and before the moon generations of generations}\footnote{Ps. lxxii. 5.}. And again another of the Prophets saith, \emph{Rejoice greatly, O daughter of Sion, shout, O daughter of Jerusalem. Behold! thy King cometh unto thee, just and having salvation}\footnote{Zech. ix. 9.}. Kings are many; of which speakest thou, O Prophet? Give us a sign which other Kings have not. If thou say, A king clad in purple, the dignity of the apparel has been anticipated. If thou say, Guarded by spear-men, and sitting in a golden chariot, this also has been anticipated by others. Give us a sign peculiar to the King whose coming thou announcest. And the Prophet maketh answer and saith, \emph{Behold! thy King cometh unto thee, just, and having salvation: He is meek, and riding upon an ass and a young foal}, not on a chariot. Thou hast a unique sign of the King who came. Jesus alone of kings sat upon an unyoked\footnote{\textgreek{ἀσαγῆ}, a rare word, formed from \textgreek{σάγη}, “harness.”} foal, entering into Jerusalem with acclamations as a king. And when this King is come, what doth He? \emph{Thou also by the blood of the covenant hast sent forth thy prisoners out of the pit wherein is no water}\footnote{Zech. ix. 11.}.

11. But He might perchance even sit upon a foal: give us rather a sign, where the King that entereth shall stand. And give the sign not far from the city, that it may not be unknown to us: and give us the sign plain before our eyes, that even when in the city we may behold the place. And the Prophet again makes answer, saying: \emph{And His feet shall stand in that day upon the Mount of Olives which is before Jerusalem on the east}\footnote{Zech. xiv. 4. “There is an excellent view from the city of the Mount of Olives which stands up over against it, especially from the height of Golgotha where Cyril was delivering his Lectures” (Cleopas).}. Does any one standing within the city fail to behold the place?

12. We have two signs, and we desire to learn a third. Tell us what the Lord doth when He is come. Another Prophet saith, \emph{Behold! our God}, and afterwards, \emph{He will come and save us. Then the eyes of the blind shall be opened, and the ears of the deaf shall hear: then shall the lame man leap as an hart, and the tongue of the stammerers shall be distinct}\footnote{Isa. xxxv. 4–6.}. But let yet another testimony be told us. Thou sayest, O Prophet, that the Lord cometh, and doeth signs such as never were: what other clear sign tellest thou? \emph{The Lord Himself entereth into judgment with the elders of His people, and with the princes thereof}\footnote{Ib. iii. 14.}. A notable sign! The Master judged by His servants, the elders, and submitting to it.

13. These things the Jews read, but hear not: for they have stopped the ears of their heart, that they may not hear. But let us believe in Jesus Christ, as having come in the flesh and \emph{been made Man}, because we could not receive Him otherwise. For since we could not look upon or enjoy Him as He was, He became what we are, that so we might be permitted to enjoy Him. For if we cannot look full on the sun, which was made on the fourth day, could we behold God its Creator\footnote{Cf. Epist. Barnab. § 13: “For had He not come in flesh, how could we men have been safe in beholding Him? For in beholding the Sun, which being the work of His hands shall cease to be, men have no strength to fix their eyes upon him.”}? The Lord came down in fire on Mount Sinai, and the people could not bear it, but said to Moses, \emph{Speak thou with us, and we will hear; and let not God speak to us, lest we die}\footnote{Exod. xx. 19.}: and again, \emph{For who is there of all flesh that hath heard the voice of the living God speaking out of the midst of the fire, and shall live}\footnote{Deut. v. 26.}? If to hear the voice of God speaking is a cause of death, how shall not the sight of God Himself bring death? And what wonder? Even Moses himself saith, \emph{I exceedingly fear and quake}\footnote{Heb. xii. 21.}.

14. What wouldest thou then? That He who came for our salvation should become a minister of destruction because men could not bear Him? or that He should suit His grace to our measure? Daniel could not bear the vision of an Angel, and wert thou capable of the sight of the Lord of Angels? Gabriel appeared, and Daniel fell down: and of what nature or in what guise was he that appeared? His countenance was like lightning\footnote{Dan. x. 6.}; not like the sun: \emph{and his eyes as lamps of fire}, not as a furnace of fire: \emph{and the voice of his words as the voice of a multitude}, not as the voice of twelve legions of angels; nevertheless the Prophet fell down. And the Angel cometh unto him, saying, \emph{Fear not, Daniel, stand upright: be of good courage, thy words are heard}\footnote{Dan. x. 12.}. And Daniel says, I stood up trembling\footnote{Ib. x. 11.}: and not even so did he make answer, until the likeness of a man’s hand touched him. And when he that appeared was changed into the appearance of a man, then Daniel spoke: and what saith he? \emph{O my Lord, at the vision of Thee my inward parts were turned within me, and no strength remaineth in me, neither is there breath left in me}\footnote{Ib. x. 16, 17.}. If an Angel appearing took away the Prophet’s voice and strength, would the appearance of God have allowed him to breathe? And until \emph{there touched me as it were a vision of a man}\footnote{Ib. x. 18.}, saith the Scripture, Daniel took not courage. So then after trial shewn of our weakness, the Lord assumed that which man required: for since man required to hear from one of like countenance, the Saviour took on Him the nature of like affections, that men might be the more easily instructed.

15. Learn also another cause. Christ came that He might be baptized, and might sanctify Baptism: He came that He might work wonders, walking upon the waters of the sea. Since then before His appearance in flesh, \emph{the sea saw Him and fled, and Jordan was turned back}\footnote{Ps. cxiv. 3.}, the Lord took to Himself His body, that the sea might endure the sight, and Jordan receive Him without fear. This then is one cause; but there is also a second. Through Eve yet virgin came death; through a virgin, or rather from a virgin, must the Life appear: that as the serpent beguiled the one, so to the other Gabriel might bring good tidings\footnote{Justin M. (\emph{Tryph}. § 100): “Eve, when she was a virgin and undefiled, having conceived the word of the serpent, brought forth disobedience and death: but the Virgin Mary received faith and joy, when the Angel Gabriel announced the good tidings to her.”}. Men forsook God, and made carved images of men. Since therefore an image of man was falsely worshipped as God, God became truly Man, that the falsehood might be done away. The Devil had used the flesh as an instrument against us; and Paul knowing this, saith, \emph{But I see another law in my members warring against the law of my mind, and bringing me into captivity}\footnote{Rom. vii. 23.}, and the rest. By the very same weapons, therefore, wherewith the Devil used to vanquish us, have we been saved. The Lord took on Him from us our likeness, that He might save man’s nature: He took our likeness, that He might give greater grace to that which lacked; that sinful humanity might become partaker of God. \emph{For where sin abounded, grace did much more abound}\footnote{Ib. v. 20.}. It behoved the Lord to suffer for us; but if the Devil had known Him, he would not have dared to approach Him. \emph{For had they known it, they would not have crucified the Lord of Glory}\footnote{1 Cor. ii. 8.}. His body therefore was made a bait to death that the dragon\footnote{Death is here called “the dragon,” as in xiv. 17 he is called “the invisible whale,” in allusion to the case of Jonah.}, hoping to devour it, might disgorge those also who had been already devoured\footnote{On Christ’s descent into Hades compare iv. 11; xiv. 19; and Eusebius (\emph{Dem. Evang}. x. 50), and Athanasius (\emph{c. Arian. Or}. iii. 56): “The Lord, at Whom the keepers of hell’s gates shuddered and set open hell. The Lord, Whom death as a dragon flees.”}. For \emph{Death prevailed and devoured}; and again, \emph{God wiped away every tear from off every face}\footnote{Isa. xxv. 8. The first clause, \emph{He hath swallowed up death for ever} (R.V.), is mistranslated in the Septuagint.}.

16. Was it without reason that Christ was made Man? Are our teachings ingenious phrases and human subtleties? Are not the Holy Scriptures our salvation? Are not the predictions of the Prophets? Keep then, I pray thee, this deposit\footnote{\textgreek{ταύτην τὴν παρακαταθηκην}. 1 Tim. vi. 20; 2 Tim. i. 14.} undisturbed, and let none remove thee: believe that God became Man. But though it has been proved possible for Him to be made Man, yet if the Jews still disbelieve, let us hold this forth to them: What strange thing do we announce in saying that God was made Man, when yourselves say that Abraham received the Lord as a guest\footnote{Gen. xviii. 1 ff.}? What strange thing do we announce, when Jacob says, \emph{For I have seen God face to face, and my life is preserved}\footnote{Ib. xxxii. 30.}? The Lord, who ate with Abraham, ate also with us. What strange thing then do we announce? Nay more, we produce two witnesses, those who stood before Lord on Mount Sinai: Moses was in a \emph{clift of the rock}\footnote{Ex. xxxiii. 22.}, and Elias was once in a clift of the rock\footnote{1 Kings xix. 8.}: they being present with Him at His Transfiguration on Mount Tabor, \emph{spoke} to the Disciples \emph{of His decease which fire should accomplish at Jerusalem}\footnote{Luke ix. 30, 31. On the tradition that Mt. Tabor was the place of the Transfiguration, accepted by S. Jerome and other Fathers, compare Lightfoot (\emph{Hor. Hebr}. in Marc. ix. 2).}. But, as I said before, it has been proved possible for Him to be made man: and the rest of the proofs may be left for the studious to collect.

17. My statement, however, promised to declare\footnote{Cat. xii. 5. For \textgreek{εὑρεῖν} the recent Editors with \textsc{mss. A.R.C.} and Grodecq. have \textgreek{ἐρεῖν}.} also the time of the Saviour’s and the place: and I must not go away convicted of falsehood, but rather send away the Church’s novices\footnote{\textgreek{νεήλυδας·}} well assured. Let us therefore inquire the time when our Lord came: because His coming is recent, and is disputed: and because \emph{Christ Jesus is the same yesterday, and to-day, and for ever}\footnote{Heb. xiii. 8. Cyril is supposed to refer to two objections to the Incarnation, one founded on the lateness of Christ’s coming, the other on the Divine immutability. But the meaning of the passage is not clear, and the construction of the second sentence is incomplete.}. Moses then, the prophet, saith, \emph{A Prophet shall the Lord your God raise up unto you of your brethren, like unto me}\footnote{Deut. xviii. 15; Acts vii. 37.}: but let that “\emph{like unto me}” be reserved awhile to be examined in its proper place\footnote{\textgreek{ἐξεταζόμενον} , a clear instance of the Gerundive, or quasi-Future, sense of the Present Participle, common in Cyril. “This intention is not fulfilled in the sequel of these Lectures” (\textsc{R.W.C.}).}. But when cometh this Prophet that is expected? Recur, he says, to what has been written by me: examine carefully Jacob’s prophecy addressed to Judah: \emph{Judah, thee may thy brethren praise}, and afterwards, not to quote the whole, \emph{A prince shall not fail out of Judah, nor a ruler from his loins, until He come, for whom it is reserved; and He is the expectation, not of the Jews but of the Gentiles}\footnote{Gen. xlix. 8, 10.}. He gave, therefore, as a sign of Christ’s advent the cessation of the Jewish rule. If they are not now under the Romans, the Christ is not yet come: if they still have a prince of the race of Judah and of David\footnote{According to Cyril (§ 19, below) and other Fathers, the continuance of Jewish rulers ceased on the accession of Herod an Idumean. Compare Justin M. (\emph{Tryphon} §§ 52, 120); Eusebius (\emph{Demonstr. Evang}. VIII. 1). On modern interpretations of the passage see Delitzsch (\emph{New Commentary on Genesis}), Briggs (\emph{Messianic Prophecy}, p. 93), Cheyne (\emph{Isaiah}, Vol. II. p. 189), Driver (\emph{Journal of Philology}, No. 27, 1885).}, he is not yet come that was expected. For I am ashamed to tell of their recent doings concerning those who are now called Patriarchs\footnote{A full and interesting account of the Jewish Patriarchs of the West established at Tiberias from the time of Antoninus Pius till the close of the 4th century is contained in Dean Milman’s \emph{History of the Jews}, Vol. III. Compare Epiphanius (\emph{Hæres}. xxx. § 3 ff.).} among them, and what their descent is, and who their mother: but I leave it to those who know. But He that cometh as \emph{the expectation of the Gentiles}, what further sign then hath He? He says next, \emph{Binding his foal unto the vine}\footnote{Gen. xlix. 11.}. Thou seest that foal which was clearly announced by Zachariah\footnote{Zech. ix. 9, quoted above, § 10.}.

18. But again thou askest yet another testimony of the time. \emph{The \textsc{Lord} said unto Me, Thou art My Son; this day have I begotten Thee}: and a few words further on, \emph{Thou shalt rule them with a rod of iron}\footnote{Ps. ii. 7, 9. The passage is interpreted by Cyril (xi. 5) of the eternal generation of the Son: here it refers to His Incarnation, or perhaps is meant only to identify the Son of God with him who “shall rule with a rod of iron.”}. I have said before that the kingdom of the Romans is clearly called a rod of iron; but what is wanting concerning this let us further call to mind out of Daniel. For in relating and interpreting to Nebuchadnezzar the image of the statue, he tells also his whole vision concerning it: and that a stone cut out of a mountain without hands, that is, not set up by human contrivance, should overpower the whole world: and he speaks most clearly thus; \emph{And in the days of} \emph{those kingdoms the God of heaven shall set up a kingdom, which shall never be destroyed, and His kingdom shall not be left to another people}\footnote{Dan. ii. 44.}.

19. But we seek still more clearly the proof of the time of His coming. For man being hard to persuade, unless he gets the very years for a clear calculation, does not believe what is stated. What then is the season, and what the manner of the time? It is when, on the failure of the kings descended from Judah, Herod a foreigner succeeds to the kingdom? The Angel, therefore, who converses with Daniel says, and do thou now mark the words, \emph{And thou shalt know and understand: From the going forth of the word for making answer}\footnote{Sep. \textgreek{τοῦ ἀποκριθῆναι}, a frequent meaning of the Hebrew \texthebrew{ב\texthebrew{ישִּׁהָלְ}}, by which the Greek Translators understood the answer of Darius to the Letter of Tatnai and his companions. Both A.V. and R.V. render the word “to restore.”}, \emph{and for the building of Jerusalem, until Messiah the Prince are seven weeks and three score and two weeks}\footnote{Dan. ix. 25.}. Now three score and nine weeks of years contain four hundred and eighty-three years. He said, therefore, that after the building of Jerusalem, four hundred and eighty-three years having passed, and the rulers having failed, then cometh a certain king of another race, in whose time the Christ is to be born. Now Darius the Mede\footnote{Darius the Mede (Dan. v. 31) succeeded Belshazzar as king in Babylon \textsc{b.c.} 538, the date assigned in Dan. ix. 1 to the prophecy of the 70 years. But “Darius the king” in whose 6th year (\textsc{b.c.} 516) the Temple was finished (Ezra vi. 15) was Darius Hystaspis, king of Persia, whom Cyril here confounds with “Darius the Mede.” He also fails to distinguish the rebuilding of the Temple, \textsc{b.c.} 516, from the rebuilding of the City by permission of Artaxerxes Longimanus, \textsc{b.c.} 444 (\emph{Nehemiah}, ii. 1).} built the city in the sixth year of his own reign, and first year of the 66th Olympiad according to the Greeks. Olympiad is the name among the Greeks of the games celebrated after four years, because of the day which in every four years of the sun’s courses is made up of the three\footnote{In speaking of three supernumeracy hours in the year instead of nearly six, Cyril seems to follow the division of the diurnal period into twelve parts, not twenty-four. The Jews had derived this division either from the Egyptians, or more probably from the Babylonians: see Herodotus, II. 109.} (supernumerary) hours in each year. And Herod is king in the 186th Olympiad, in the 4th year thereof. Now from the 66th to the 186th Olympiad there are 120 Olympiads intervening, and a little over. So then the 120 Olympiads make up 480 years: for the other three years remaining are perhaps taken up in the interval between the first and fourth years. And there thou hast the proof according to the Scripture which saith, \emph{From the going forth of the word that Jerusalem be restored and built until Messiah the Prince are seven weeks and sixty-two weeks}. Of the times, therefore, thou hast for the present this proof, although there are also other different interpretations concerning the aforesaid weeks of years in Daniel.

20. But now hear the place of the promise, as Micah says, \emph{And thou, Bethlehem, house of Ephrathah, art thou little to be among the thousands of Judah? For out of thee shall come forth unto Me a ruler, to be governor in Israel: and His goings forth are front the beginning, from the days of eternity}\footnote{Micah v. 2, quoted also in Cat. xi. 20, where see note.}. But assuredly as to the places, thou being an inhabitant of Jerusalem, knowest also beforehand what is written in the hundred and thirty-first psalm. \emph{Lo! we heard of it at Ephrathah, we found it in the plains of the wood\footnote{Ps. cxxxii. 6. The Psalmist refers to the recovery of the Ark, but Cyril interprets the passage mystically of Christ, and the place of His Nativity.}}. For a few years ago the place was woody\footnote{The Benedictine Editor thinks that in calling the place “woody” Cyril refers to a grove planted by Hadrian in honour of Adonis, which had been destroyed about sixteen years before, when Helena built the Church at Bethlehem: see Eusebius, \emph{Life of Constantine}, III. 43. But Cyril evidently means that the wood of which the Psalmist speaks had remained till a few years before. Ephrâthah is the ancient name of Bethlehem (Gen. xxxv. 19; xlviii. 7), and by “the fields of the wood” is probably meant Kirjath-Jearim, “the city of woods,” where the Ark was found by David (2 Sam vi. 2; 1 Chron. xiii. 6).}. Again thou hast heard Habakkuk say to the Lord, \emph{When the years draw nigh, thou shalt be made known, when the time is come, thou shalt be shewn}\footnote{Hab. iii. 2: (R.V.) \emph{O Lord, revive Thy work in the midst of the years, in the midst of the years make it known.} The Septuagint gives a different sense: \emph{In the midst of two lives} (or, \emph{living beings}) \emph{shalt Thou be known: when the years draw nigh Thou shalt be recognised: when the time is come, Thou shalt be shewn.} The two latter clauses seem to be different renderings of the same Hebrew words.}. And what is the sign, O Prophet, of the Lord’s coming? And presently he saith, \emph{In the midst of two lives shalt thou be known}\footnote{\textgreek{ἑξῆς}. This clause comes before the preceding quotation: Cyril misplaces them. In the Vatican and other \textsc{mss.} of the Sept. and in some Fathers \textgreek{ζώων} (“living creatures”) is found in place of \textgreek{ζωῶν} “lives;” but the latter reading is evidently required by the interpretation which follows in Cyril. Origen (\emph{de Principiis,} I. 4), who recognises both readings (“In medio vel duorum animalium, vel duarum vitarum, cognosceris,”) interprets the “two living beings” of the Son and the Spirit. Eusebius (\emph{Demonstr. Evang} VI. 15) observes that \textgreek{ζωων} is to be read as perispomenon from the Singular \textgreek{ζωή}, and interprets it of Christ’s life with God, and life on earth. Theodoret says, in commenting on the passage, “To me it seems that the Prophet means not “living beings” (\textgreek{ζῶα}) but “lives” (\textgreek{ζωάς}), the present life, and that which is to come, between which is the appearance of the Righteous Judge.”}, plainly saying this to the Lord, “Having come in the flesh thou livest and diest, and after rising from the dead thou livest again.” Further, from what part of the region round Jerusalem cometh He? From east, or west, or north, or south? Tell us exactly. And he makes answer most plainly and says, \emph{God shall come from Teman}\footnote{Hab. iii. 3. Cyril interprets the word \textgreek{Θαιμάν} (Heb. \texthebrew{ן\texthebrew{מָיתּ}”}) as a common Noun meaning “South,” and the Vulgate has here “ab Austro veniet.” The prophecy is thus referred to Bethlehem, as lying to the South of Jerusalem. Eusebius (\emph{Dem. Evang}. VI. 15) mentions this as the rendering of Theodotion in his Greek Version, about 180 \textsc{a.d.} As a proper name Teman denotes a district and town in the southern part of Idumea, so called from a grandson of Esau (Gen. xxxvi. 11, 15, 42; Jer. xlix. 7, 20; Ezek. xxv. 13; Amos i. 12; Obad. 9).} (now Teman is by interpretation ‘south’) \emph{and the Holy One from Mount Paran}\footnote{The following note is slightly abridged from the Edition of Alexandrides of Jerusalem. “Previous Editions read \textgreek{ἔξ ὄρους φαρὰν κατασκίου δασέος}. This reading is found in Cod. Vat. and other \textsc{mss.} of the Septuagint, but \textgreek{φαράν} is omitted in the Aldine and many other copies nor was it read in the \textsc{mss.} of the Sept. in Jerome’s time, as is clear from his comments on the passage. In the \textsc{mss.} of Cyril, Ottob. \textsc{R.C.V.} Monac. I. and II. it is wanting. Paran is the name of the desert towards the S. of Palestine lying between it and Egypt (Gen. xxi. 21; Num. i. 12). There was also a Mount Paran (Deut. xxxiii. 2). But since Cyril applies the prophecy to Bethlehem, and the “shady thickly-wooded mountain” of Habakkuk is identified with “the plains of the wood” of David, we may safely conclude that Cyril did not read \textgreek{φαράν} in his copies of the Septuagint, nor write it in his Lecture: but the reading crept in from the later copyists, accustomed to the reading \textgreek{φαράν} in the Septuagint.”}, \emph{shady, woody}: what the Psalmist spake in like words, \emph{We found it in the plains of the wood}.

21. We ask further, of whom cometh He and how? And this Esaias tells us: \emph{Behold! the virgin shall conceive in her womb, and shall bring forth a Son, and they shall call His name Emmanuel}\footnote{Isa. vii. 14. The objection of the Jews that the Hebrew word “Almah” means “a young woman,” whether married or not, is mentioned by Justin M. (\emph{Tryph}. 43, 67, 71), and by Eusebius (\emph{Dem. Evang}. VII. i. 315).}. This the Jews contradict, for of old it is their wont wickedly to oppose the truth: and they say that it is not written “the virgin,” but “the damsel.” But though I assent to what they say, even so I find the truth. For we must ask them, If a virgin be forced, when does she cry out and call for helpers, after or before the outrage? If, therefore, the Scripture elsewhere says, \emph{The betrothed damsel cried, and there was none to save her}\footnote{Deut. xxii. 27.}, doth it not speak of a virgin?

But that you may learn more plainly that even a virgin is called in Holy Scripture a “damsel,” hear the Book of the Kings, speaking of Abishag the Shunamite, \emph{And the damsel was very fair}\footnote{1 Kings i. 4. Cyril’s argument is fully justified by the actual usage of “Almah,” which certainly refers to unmarried women in Gen. xxiv. 43; Ex. ii. 8; Cant. i. 3. The same is probably the meaning in Ps. lxviii. 25: “in the midst were the damsels playing with the timbrels.” There is no passage in which the word can be shewn to mean a married woman.}: for that as a virgin she was chosen and brought to David is admitted.

22. But the Jews say again, This was said to Ahaz in reference to Hezekiah. Well, then, let us read the Scripture: \emph{Ask thee a sign of the Lord thy God, in the depth or in the height}\footnote{Isa. vii. 11.}. And the sign certainly must be something astonishing. For the water from the rock was a sign, the sea divided, the sun turning back, and the like. But in what I am going to mention there is still more manifest refutation of the Jews. (I know that I am speaking at much length, and that my hearers are wearied: but bear with the fulness of my statements, because it is for Christ’s sake these questions are moved, and they concern no ordinary matters.) Now as Isaiah spoke this in the reign of Ahaz, and Ahaz reigned only sixteen years, and the prophecy was spoken to him within these years, the objection of the Jews is refuted by the fact that the succeeding king, Hezekiah, son of Ahaz, was twenty-five years old when he began to reign: for as the prophecy is confined within sixteen years, he must have been begotten of Ahaz full nine years before the prophecy. What need then was there to utter the prophecy concerning one who had been already begotten even before the reign of father Ahaz\footnote{Compare Justin M. (\emph{Tryph.} § 77), Euseb. (\emph{Demonstr. Evang}. L. VII. c. i. 317).}? For he said not, \emph{hath conceived}, but “\emph{the virgin shall conceive},” speaking as with foreknowledge\footnote{In the Hebrew the word used is a Participle, and describes what Isaiah sees in a prophetic vision; “\emph{Behold, the damsel—with child}.”}.

23. We know then for certain that the Lord was to be born of a Virgin, but we have to shew of what family the Virgin was. \emph{The Lord sware in truth unto David, and will not set it aside. Of the fruit of body will I set upon thy throne}\footnote{Ps. cxxxii. 11.}: \emph{ and again, seed will I establish for ever, and his throne as the days of heaven}\footnote{Ib. lxxxix. 22.}. \emph{ And afterwards, Once have I sworn by My holiness that I will not lie unto David. His seed shall endure for ever, and his throne as the sun before Me, and as the moon established for ever}\footnote{\emph{vv}. 35–37.}. Thou seest that the discourse is of Christ, not of Solomon. For Solomon’s throne endured not as the sun. But if any deny this, because Christ sat not on David’s throne of wood, we will bring forward that saying, \emph{The Scribes and the Pharisees sit in Moses’ seat}\footnote{Matt. xxiii. 2.}: for it signifies not his wooden seat, but the authority of his teaching. In like manner then I would have you seek for David’s throne not the throne of wood, but the kingdom itself. Take, too, as my witnesses the children who cried aloud,\emph{Hosanna to the Son of David}\footnote{Ib. xxi. 9.}, \emph{blessed is the King of Israel}\footnote{John xii. 13.}. And the blind men also say, \emph{Son of David, have mercy on us}\footnote{Matt. xx. 30.}. Gabriel too testifies plainly to Mary, saying, \emph{And the Lord God shall give unto Him the throne of His father David}\footnote{Luke i. 32.}. Paul also saith, \emph{Remember Jesus Christ raised from the dead, of the seed of David, according to my Gospel}\footnote{2 Tim. ii. 8.}: and in the beginning of the Epistle to the Romans he saith, \emph{Which was made of the seed of David according to the flesh}\footnote{Rom. i. 3.}. Receive thou therefore Him that was born of David, believing the prophecy which saith, \emph{And in that day there shall be a root of Jesse, and He that shall rise to rule over the Gentiles: in Him shall the Gentiles trust}\footnote{Is. xi. 10; Rom. xv. 12.}.

24. But the Jews are much troubled at these things. This also Isaiah foreknew, saying, \emph{And they shall wish that they had been burnt with fire: for unto us a child is born} (not unto them), \emph{unto us a Son is given}\footnote{Isa. ix. 5.}. Mark thou that at first He was the Son of God, then was given to us. And a little after he says, \emph{And of His peace there is no bound}\footnote{\emph{v.} 7.}. The Romans have bounds: of the kingdom of the Son of God there is no bound. The Persians and the Medes have bounds, but the Son has no bound. Then next, \emph{upon the throne of David, and upon his kingdom to order it}. The Holy Virgin, therefore, is from David.

25. For it became Him who is most pure, and a teacher of purity, to have come forth from a pure bride-chamber. For if he who well fulfils the office of a priest of Jesus abstains from a wife, how should Jesus Himself be born of man and woman? \emph{For thou}, saith He in the Psalms, art \emph{He that took Me out of the womb}\footnote{Ps. xxii. 9.}. Mark that carefully, He that took Me out of the womb, signifying that He was begotten without man, being taken from a virgin’s womb and flesh. For the manner is different with those who are begotten according to the course of marriage.

26. And from such members He is not ashamed to assume flesh, who is the framer of those very members. But then who telleth us this? The Lord saith unto Jeremiah: \emph{Before I formed thee in the belly, I knew thee: and before thou camest forth out of the womb, I sanctified thee}\footnote{Jer. i. 5.}. If, then, in fashioning man He was not ashamed of the contact, was He ashamed in fashioning for His own sake the holy Flesh, the veil of His Godhead? It is God who even now creates the children in the womb, as it is written in Job, \emph{Hast thou not poured me out as milk, and curdled me like cheese? Thou hast clothed me with skin and flesh, and hast knit me together with bones and sinews}\footnote{Job x. 10, 11.}. There is nothing polluted in the human frame except a man defile this with fornication and adultery. He who formed Adam formed Eve also, and male and female were formed by God’s hands. None of the members of the body as formed from the beginning is polluted. Let the mouths of all heretics be stopped who slander their bodies, or rather Him who formed them. But let us remember Paul’s saying, \emph{Know ye not that your bodies are the temples of the Holy Ghost which is in you}\footnote{1 Cor. vi. 19.}? And again the Prophet hath spoken before in the person of Jesus, \emph{My flesh is from them}\footnote{Hos. ix. 12. R.V. \emph{Woe also to them, when I depart from them.} The Seventy mistook \texthebrew{י\texthebrew{רִוּשֹבְּ}}, “at my departure,” for \texthebrew{י\texthebrew{רְשָֹבִּ}}, “my flesh.”}: and in another place it is written, \emph{Therefore will He give them up, until the time that she bringeth forth}\footnote{Mic. v. 3.}. And what is the sign? He tells us in what follows, \emph{She shall bring forth, and the remnant of their brethren shall return}. And what are the nuptial pledges of the Virgin, the holy bride? \emph{And I will betroth thee unto Me in faithfulness}\footnote{Hos. ii. 20.}. And Elizabeth, talking with Mary, speaks in like manner: \emph{And blessed is she that believed; for there shall be a performance of those things which were told her from the Lord}\footnote{Luke i. 45.}.

27. But both Greeks and Jews harass us and say that it was impossible for the Christ to be born of a virgin. As for the Greeks we will stop their mouths from their own fables. For ye who say that stones being thrown were changed into men\footnote{See the story of Pyrrha and Deucalion in Pindar, \emph{Ol}. ix. 60: \textgreek{ἄτερ δ᾽ εὐνᾶς κτησάσθαν λίθινον γόνον}, and in Ovid. \emph{Metam} i. 260 ff.}, how say ye that it is impossible for a virgin to bring forth? Ye who fable that a daughter was born from the brain\footnote{Athena was said to have sprung armed from the head of Zeus: Pindar, \emph{Ol}. vii. 65: \textgreek{κορυφὰν κατ᾽ ἄκραν ἀνορούσαισ᾽ ἀλάλαξεν ὑπερμάκει βοᾷ}. Cf. Hes. \emph{Theog}. 924.}, how say ye that it is impossible for a son to have been born from a virgin’s womb? Ye who falsely say that Dionysus was born from the thigh of your Zeus\footnote{Eurip. \emph{Bacchae}. 295; Ovid. \emph{Metam}. iv. 11.}, how set ye at nought our truth? I know that I am speaking of things unworthy of the present audience: but in order that thou in due season mayest rebuke the Greeks, we have brought these things forward answering them from their own fables.

28. But those of the circumcision meet thou with this question: Whether is harder, for an aged woman, barren and past age, to bear, or for a virgin in the prime of youth to conceive? Sarah was barren, and though it had ceased to be with her after the manner of women, yet, contrary to nature, she bore a child. If, then, it is against nature for a barren woman to conceive, and also for a virgin, either, therefore, reject both, or accept both. For it is the same God\footnote{Codd. Mon. i, A: \textgreek{ὁ γὰρ αὐτὸς Θεός}. Bened. \textgreek{ὁ γὰρ Θεὸς αὐτός}.} who both wrought the one and appointed the other. For thou wilt not dare to say that it was possible for God in that former case, and impossible in this latter. And again: how is it natural for a man’s hand to be changed in a single hour into a different appearance and restored again? How then was the hand of Moses made white as snow, and at once restored again? But thou sayest that God’s will made the change. In that case God’s will has the power, and has it then no power in this case? That moreover was a sign concerning the Egyptians only, but this was a sign given to the whole world. But whether is the more difficult, O ye Jews? For a virgin to bear, or for a rod to be quickened into a living creature? Ye confess that in the case of Moses a perfectly straight rod became like a serpent, and was terrible to him who cast it down, and he who before held the rod fast, fled from it as from a serpent; for a serpent in truth it was: but he fled not because he feared that which he held, but because he dreaded Him that had changed it. A rod had teeth and eyes like a serpent: do then seeing eyes grow out of a rod, and cannot a child be born of a virgin’s womb, if God wills? For I say nothing of the fact that Aaron’s rod also produced in a single night what other trees produce in several years. For who knows not that a rod, after losing its bark, will never sprout, not even if it be planted in the midst of rivers? But since God is not dependent on the nature of trees, but is the Creator of their natures, the unfruitful, and dry, and barkless rod budded, and blossomed, and bare almonds. He, then, who for the sake of the typical high-priest gave fruit supernaturally to the rod, would He not for the sake of the true High-Priest grant to the Virgin to bear a child?

29. These are excellent suggestions of the narratives: but the Jews still contradict, and do not yield to the statements concerning the rod, unless they may be persuaded by similar strange and supernatural births. Question them, therefore, in this way: of whom in the beginning was Eve begotten? What mother conceived her the motherless? But the Scripture saith that she was born out of Adam’s side. Is Eve then born out of a man’s side without a mother, and is a child not to be born without a father, of a virgin’s womb? This debt of gratitude was due to men from womankind: for Eve was begotten of Adam, and not conceived of a mother, but as it were brought forth of man alone. Mary, therefore, paid the debt, of gratitude, when not by man but of herself alone in an immaculate way she conceived of the Holy Ghost by the power of God.

30. But let us take what is yet a greater wonder than this. For that of bodies bodies should be conceived, even if wonderful, is nevertheless possible: but that the dust of the earth should become a man, this is more wonderful. That clay moulded together should assume the coats and splendours of the eyes, this is more wonderful. That out of dust of uniform appearance should be produced both the firmness of bones, and the softness of lungs, and other different kinds of members, this is wonderful. That clay should be animated and travel round the world self moved, and should build houses, this is wonderful. That clay should teach, and talk, and act as carpenter, and as king, this is wonderful. Whence, then, O ye most ignorant Jews, was Adam made? Did not God take dust from the earth, and fashion this wonderful frame? Is then clay changed into an eye, and cannot a virgin bear a son. Does that which for men is more impossible take place, and is that which is possible never to occur?

31. Let us remember these things, brethren: let us use these weapons in our defence. Let us not endure those heretics who teach Christ’s coming as a phantom. Let us abhor those also who say that the Saviour’s birth was of husband and wife; who have dared to say that He was the child of Joseph and Mary, because it is written, \emph{And he took unto him his wife}\footnote{Matt. i. 24.}. For let us remember Jacob who before he received Rachel, said to Laban, \emph{Give me my wife}\footnote{Gen. xxix. 21.}. For as she before the wedded state, merely because there was a promise, was called the wife of Jacob, so also Mary, because she had been betrothed, was called the wife of Joseph. Mark also the accuracy of the Gospel, saying, \emph{And in the sixth month the Angel Gabriel was sent from God unto a city of Galilee, named Nazareth, to a virgin espoused to a man whose name was Joseph}\footnote{Luke i. 26, 27.}, and so forth. And again when the census took place, and Joseph went up to enrol himself, what saith the Scripture? \emph{And Joseph also went up from Galilee, to enrol himself with Mary who was espoused to him, being great with child}\footnote{Ib. ii. 4, 5.}. For though she was with child, yet it said not “with his wife,” but with her who was \emph{espoused to him}. For \emph{God sent forth His Son}, says Paul, not made of a man and a woman, but \emph{made of a woman}\footnote{Gal. iv. 4.} only, that is of a virgin. For that the virgin also is called a woman, we shewed before\footnote{See above, § 21.}. For He who makes souls virgin, was born of a Virgin.

32. But thou wonderest at the event: even she herself who bare him wondered at this. For she saith to Gabriel, \emph{How shall this be to me, since I know not a man?} But he says, \emph{The Holy Ghost shall came upon thee, and the power of the Highest shall overshadow thee: wherefore also the holy thing which is to be born shall be called the Son of God}\footnote{Luke i. 34, 35.}. Immaculate and undefiled was His generation: for where the Holy Spirit breathes, there all pollution is taken away: undefiled from the Virgin was the incarnate generation of the Only-begotten. And if the heretics gainsay the truth, the Holy Ghost shall convict them: that overshadowing power of the Highest shall wax wroth: Gabriel shall stand face to face against them in the day of judgment: the place of the manger, which received the Lord, shall put them to shame. The shepherds, who then received the good tidings, shall bear witness; and the host of the Angels who sang praises and hymns, and said, \emph{Glory to God in the highest, and on earth peace among men of His good pleasure}\footnote{Luke ii. 14.}: the Temple into which He was then carried up on the fortieth day: the pairs of turtle-doves, which were offered on His behalf\footnote{Ib. ii. 24. In Lev. xii. 8 one pair only of turtles is prescribed, to be offered for the mother, not for the child. But the reading \textgreek{τὰ ζεύγη} in Cyril is confirmed by that in St. Luke, \textgreek{τοῦ καθαρισμοῦ αὐτῶν}. See the authorities in Tischendorf.}: and Symeon who then took Him up in his arms, and Anna the prophetess who was present.

33. Since God then beareth witness, and the Holy Ghost joins in the witness, and Christ says, \emph{Why do ye seek to kill me, a man who has told you the truth}\footnote{John vii. 19; viii. 40.}? let the heretics be silenced who speak against His humanity, for they speak against Him, who saith, \emph{Handle me, and see; for a spirit hath not flesh and bones, as ye see me have}\footnote{Luke xxiv. 39.}. Adored be the Lord the Virgin-born, and let Virgins acknowledge the crown of their own state: let the order also of Solitaries acknowledge the glory of chastity for we men are not deprived of the dignity of chastity. In the Virgin’s womb the Saviour’s period of nine months was passed: but the Lord was for thirty and three years a man: so that if a virgin glories\footnote{\textgreek{σεμνύνεται}. Rivet, misled by a double error in the old Latin version, “veneratur,” accused Cyril of approving the worship of the Virgin Mary.} because of the nine months, much more we because of the many years.

34. But let us all by God’s grace run the race of chastity, \emph{young men and maidens, old men and children}\footnote{Ps. cxlviii. 12.}; not going after wantonness, but praising the name of Christ. Let us not be ignorant of the glory of chastity: for its crown is angelic, and its excellence above man. Let us be chary of our bodies which are to shine as the sun: let us not for short pleasure defile so great, so noble a body: for short and momentary is the sin, but the shame for many years and for ever. Angels walking upon earth are they who practise chastity: the Virgins have their portion with Mary the Virgin. Let all vain ornament be banished, and every hurtful glance, and all wanton gait, and every flowing robe, and perfume enticing to pleasure. But in all for perfume let there be the prayer of sweet odour, and the practice\footnote{\textgreek{ἡ τῶν ἀγαθῶν πρᾶξις}, Cod. A.} of good works, and the sanctification of our bodies: that the Virgin-born Lord may say even of us, both men who live in chastity and women who wear the crown, \emph{I will dwell in them; and walk in them, and I will be their God, and they shall be My people}\footnote{2 Cor. vi. 16.}. To whom be the glory for ever and ever. Amen.

\xchapter{Lecture XIII. On the words, Crucified and Buried.}

\textsc{Isaiah liii. 1, 7}

Who hath believed our report? and to whom is the arm of the Lord revealed?…He is brought as a lamb to the slaughter, \&c.

1. \textsc{Every} deed of Christ is a cause of glorying to the Catholic Church, but her greatest of all glorying is in the Cross; and knowing this, Paul says, \emph{But God forbid that I should glory, save in the Cross of Christ}\footnote{Gal. vi. 14.}. For wondrous indeed it was, that one who was blind from his birth should receive sight in Siloam\footnote{Cf. Athanas. (\emph{de Incarn}. § 18, 49).}; but what is this compared with the blind of the whole world? A great thing it was, and passing nature, for Lazarus to rise again on the fourth day; but the grace extended to him alone, and what was it compared with the dead in sins throughout the world? Marvellous it was, that five loaves should pour forth food for the five thousand; but what is that to those who are famishing in ignorance through all the world? It was marvellous that she should have been loosed who had been bound by Satan eighteen years: yet what is this to all of us, who were fast bound in the chains of our sins? But the glory of the Cross led those who were blind through ignorance into light, loosed all who were held fast by sin, and ransomed the whole world of mankind.

2. And wonder not that the whole world was ransomed; for it was no mere man, but the only-begotten Son of God, who died on its behalf. Moreover one man’s sin, even Adam’s, had power to bring death to the world; but \emph{if by the trespass of the one death reigned} over the world, how shall not life much rather reign \emph{by the righteousness of the One}\footnote{Rom. v. 17, 18.}? And if because of the tree of food they were then cast out of paradise, shall not believers now more easily enter into paradise because of the Tree of Jesus? If the first man formed out of the earth brought in universal death, shall not He who formed him out of the earth bring in eternal life, being Himself the Life? If Phinees, when he waxed zealous and slew the evil-doer, staved the wrath of God, shall not Jesus, who slew not another, but \emph{gave up Himself for a ransom}\footnote{1 Tim. ii. 6.}, put away the wrath which is against mankind?

3. Let us then not be ashamed of the Cross of our Saviour, but rather glory in it. \emph{For the word of the Cross is unto Jews a stumbling-block, and unto Gentiles foolishness}, but to us salvation: and \emph{to them that are perishing it is foolishness, but unto us which are being saved it is the power of God}\footnote{1 Cor. i. 18, 23.}. For it was not a mere man who died for us, as I said before, but the Son of God, God made man. Further; if the lamb under Moses drove the destroyer\footnote{Ex. xii. 23.} far away, did not much rather the \emph{Lamb of God, which taketh away the sin of the world}\footnote{John i. 29.}, deliver us from our sins? The blood of a silly sheep gave salvation; and shall not the Blood of the Only-begotten much rather save? If any disbelieve the power of the Crucified, let him ask the devils; if any believe not words, let him believe what he sees. Many have been crucified throughout the world, but by none of these are the devils scared; but when they see even the Sign of the Cross of Christ, who was crucified for us, they shudder\footnote{Cf. Cat. i. 3; xvii. 35, 36.}. For those men died for their own sins, but Christ for the sins of others; for He \emph{did no sin, neither was guile found in His mouth}\footnote{1 Pet. ii. 22, quoted from Isa. liii. 9.}. It is not Peter who says this, for then we might suspect that he was partial to his Teacher; but it is Esaias who says it, who was not indeed present with Him in the flesh, but in the Spirit foresaw His coming in the flesh. Yet why now bring the Prophet only as a witness? take for a witness Pilate himself, who gave sentence upon Him, saying, \emph{I find no fault in this Man}\footnote{Luke xxiii. 14.}: and when he gave Him up, and had washed his hands, he said, \emph{I am innocent of the blood of this just person}\footnote{Matt. xxvii. 24.}. There is yet another witness of the sinlessness of Jesus,—the robber, the first man admitted into Paradise; who rebuked his fellow, and said, “\emph{We receive the due reward of our deeds; but this man hath done nothing amiss}\footnote{Luke xxiii. 41. Cf. Cat. xiii. 30, 31. The Benedictine Editor remarks, “We know not whence Cyril took the notion that the two robbers were present at the trial of Jesus.” He may have inferred from the words \textgreek{ἐν τῷ αὐτῷ κρίματι} that the sentence of crucifixion was pronounced on them at the same time as on Jesus.}; for we were present, both thou and I, at His judgment.”

4. Jesus then really suffered for all men; for the Cross was no illusion\footnote{\textgreek{δόκησις}. Cf. Ignat. \emph{Smyrn}. § 2: “He suffered truly, as also He raised Himself truly: not as certain unbelievers say, that He suffered in semblance (\textgreek{τὸ δοκεῖν αὐτὸν πεπονθέναι}).” See § 37, below.}, otherwise our redemption is an illusion also. His death was not a mere show\footnote{\textgreek{φαντασιώδης}. Athanas. \emph{c. Apollinar}. § 3: “Supposing the exhibition and the endurance of the Passion to be a mere show (\textgreek{φαντασίαν}).”}, for then is our salvation also fabulous. If His death was but a show, they were true who said, \emph{We remember that that deceiver said, while He was yet alive, After three days I rise again}\footnote{Matt. xxvii. 63.}. His Passion then was real: for He was really crucified, and we are not ashamed thereat; He was crucified, and we deny it not, nay, I rather glory to speak of it. For though I should now deny it, here is Golgotha to confute me, near which we are now assembled; the wood of the Cross confutes me, which was afterwards distributed piecemeal from hence to all the world\footnote{Cf. iv. 10; x. 19.}. I confess the Cross, because I know of the Resurrection; for if, after being crucified, He had remained as He was, I had not perchance confessed it, for I might have concealed both it and my Master; but now that the Resurrection has followed the Cross, I am not ashamed to declare it.

5. Being then in the flesh like others, He was crucified, but not for the like sins. For He was not led to death for covetousness, since He was a Teacher of poverty; nor was He condemned for concupiscence, for He Himself says plainly, \emph{Whosoever shall look upon a woman to lust after her, hath committed adultery with her already}\footnote{Matt. v. 28.}; not for smiting or striking hastily, for He turned the other cheek also to the smiter; not for despising the Law, for He was the fulfiller of the Law; not for reviling a prophet, for it was Himself who was proclaimed by the Prophets; not for defrauding any of their hire, for He ministered without reward and freely; not for sinning in words, or deeds, or thoughts, He \emph{who did no sin, neither was guile found in His mouth; who when He was reviled, reviled not again; when He suffered, threatened not}\footnote{1 Pet. ii. 22, 23.}; \emph{who came to His passion, not unwillingly, but willing; yea, if any dissuading Him say even now, Be it far from Thee, Lord}, He will say again, \emph{Get thee behind Me, Satan}\footnote{Matt. xvi. 22, 23.}.

6. And wouldest thou be persuaded that He came to His passion willingly? others, who foreknow it not, die unwillingly; but He spoke before of His passion: \emph{Behold, the Son of man is betrayed to be crucified}\footnote{Ib. xxvi. 2.}. But knowest thou wherefore this Friend of man shunned not death? It was lest the whole world should perish in its sins. \emph{Behold, we go up to Jerusalem, and the Son of man shall be betrayed, and shall be crucified}\footnote{Ib. xx. 18.}\emph{; and again, He stedfastly set His face to go to Jerusalem}\footnote{Luke ix. 5.}. And wouldest thou know certainly, that the Cross is a glory to Jesus? Hear His own words, not mine. Judas had become ungrateful to the Master of the house, and was about to betray Him. Having but just now gone forth from the table, and drunk His cup of blessing, in return for that drought of salvation he sought to shed righteous blood. \emph{He who did eat of His bread, was lifting up his heel against Him}\footnote{Ps. xli. 9.}; his hands were but lately receiving the blessed gifts\footnote{“\textgreek{τὰς εὐλογίας}. The word has this meaning in Chrysostom and Cyril of Alexandria also; afterwards it came to signify consecrated bread, distinct from that of the Eucharist. Vid. Bingham, \emph{Antiq}. xv. 4, § 3.” (\textsc{R.W.C.}) \par The custom of sending the bread of the Eucharist was forbidden in the latter part of the 4th century by the Synod of Laodicea, Canon 14: “At Easter the Host shall no more be sent into foreign dioceses as \emph{eulogiae}.” Bp. Hefele (\emph{Councils} II. p. 308) says—“It was a custom in the ancient Church, not indeed to consecrate, but to bless those of the several breads of the same form laid on the altar which were not needed for the Communion, and to employ them partly for the maintenance of the Clergy, and partly for distributing them to those of the faithful who did not communicate at the Mass.” See Eusebius (\emph{Hist. Eccles}. V. 24), with the note thereon in this Series.}, and presently for the wages of betrayal he was plotting His death. And being reproved, and having heard that word, \emph{Thou hast said}\footnote{Matt. xxvi. 25.}, he again went out: then said Jesus, \emph{The hour is come, that the Son of man should be glorified}\footnote{John xii. 23.}. Seest thou how He knew the Cross to be His proper glory? What then, is Esaias not ashamed of being sawn asunder\footnote{See Cat. ii. 14, note 4.}, and shall Christ be ashamed of dying for the world? \emph{Now is the Son of man glorified}\footnote{John xiii. 31.}. Not that He was without glory before: for He was \emph{glorified with the glory} which was \emph{before the foundation of the world}\footnote{Ib. xvii. 5.}. He was ever glorified as God; but now He was to be glorified in wearing the Crown of His patience. He gave not up His life by compulsion, nor was He put to death by murderous violence, but of His own accord. Hear what He says: \emph{I have power to lay down My life, and I have power to take it again}\footnote{Ib. x. 18.}: I yield it of My own choice to My enemies; for unless I chose, this could not be. He came therefore of His own set purpose to His passion, rejoicing in His noble deed, smiling at the crown, cheered by the salvation of mankind; not ashamed of the Cross, for it was to save the world. For it was no common man who suffered, but God in man’s nature, striving for the prize of His patience.

7. But the Jews contradict this\footnote{There is so close a resemblance between the remainder of this Lecture and the explanation of the same Article of the Creed by Rufinus, that “I have no doubt,” says the Benedictine Editor, “that Rufinus drew from Cyril’s fountains.” Cf. Rufin. \emph{de Symbolo}, § 19, \emph{sqq}.}, ever ready, as they are, to cavil, and backward to believe; so that for this cause the Prophet just now read says, \emph{Lord, who hath believed our report}\footnote{Isa. lii. 15.}? Persians believe\footnote{Cf. Acts ii. 9: \emph{Parthians and Medes and Elamites}. These Jewish converts of the day of Pentecost would naturally be the first heralds of the Gospel in their respective countries. On the dispersion of the Apostles, “Parthia, according to tradition, was allotted to Thomas as his field of labour” (Euseb. \emph{Hist. Eccl}. III. 1; cf. I. 13). An earlier notice of the tradition is found in the \emph{Clementine Recognitions}, L. IX. c. 29, where the Pseudo-Clement professes to have received a letter from “Thomas, who is preaching the Gospel among them.”}, and Hebrews believe not; \emph{they shall see, to whom He was not spoken of, and they that have not heard shall understand}\footnote{Rom. xv. 21, quoted from Isaiah, \emph{u s}.}, while they who study these things shall set at nought what they study. They speak against us, and say, “Does the Lord then suffer? What? Had men’s hands power over His sovereignty?” Read the Lamentations; for in those Lamentations, Jeremias, lamenting you, wrote what is worthy of lamentations. He saw your destruction, he beheld your downfall, he bewailed Jerusalem which then was; for that \emph{which now is}\footnote{Gal. iv. 25.} shall not be bewailed; for that Jerusalem crucified the Christ, but that \emph{which now is} worships Him. Lamenting then he says, \emph{The breath of our countenance, Christ the Lord was taken in our corruptions}\footnote{Lam. iv. 20: \emph{The breath of our nostrils, the anointed of the Lord, was taken in their pits.}}. Am I then stating views of my own? Behold he testifies of the Lord Christ seized by men. And what is to follow from this? Tell me, O Prophet. He says, \emph{Of whom we said, Under His shadow we shall live among the nations}\footnote{Ibid.}. For he signifies that the grace of life is no longer to dwell in Israel, but among the Gentiles.

8. But since there has been much gainsaying by them, come, let me, with the help of your prayers, (as the shortness of the time may allow,) set forth by the grace of the Lord some few testimonies concerning the Passion. For the things concerning Christ are all put into writing, and nothing is doubtful, for nothing is without a text. All are inscribed on the monuments of the Prophets; clearly written, not on tablets of stone, but by the Holy Ghost. Since then thou hast heard the Gospel speaking concerning Judas, oughtest thou not to receive the testimony to it? Thou hast heard that He was pierced in the side by a spear; oughtest thou not to see whether this also is written? Thou hast heard that He was crucified in a garden; oughtest thou not to see whether this also is written? Thou hast heard that He was sold for thirty pieces of silver; oughtest thou not to learn what prophet spake this? Thou hast heard that He was given vinegar to drink; learn where this also is written. Thou hast heard that His body was laid in a rock, and that a stone was set over it; oughtest thou not to receive this testimony also from the prophet? Thou hast heard that He was crucified with robbers; oughtest thou not to see whether this also is written? Thou hast heard that He was buried; oughtest thou not to see whether the circumstances of His burial are anywhere accurately written? Thou hast heard that He rose again; oughtest thou not to see whether we mock thee in teaching these things? For \emph{our speech and our preaching is not in persuasive words of man’s wisdom}\footnote{1 Cor. ii. 4. The simple style of the New Testament is defended by Origen, \emph{c. Celsum}, iii. 68, and in many other passages.}. We stir now no sophistical contrivances; for these become exposed; we do not conquer words with words\footnote{Cyril alludes to the same proverb in the \emph{Homily on the Paralytic}, c. 14: “Word resists word, but a deed is irresistible.” The Jerusalem Editor refers to Gregory Nazianzen (Tom. II. p. 596): \textgreek{Δόγῳ παλαίει πᾶς λόγος}.}, for these come to an end; but \emph{we preach Christ Crucified}\footnote{1 Cor. i. 23.}, who has already been preached aforetime by the Prophets. But do thou, I pray, receive the testimonies, and seal them in thine heart. And, since they are many, and the rest of our time is narrowed into a short space, listen now to a few of the more important as time permits; and having received these beginnings, be diligent and seek out the remainder. Let not thine hand be only stretched out to receive, but let it be also ready to work\footnote{Ecclus. iv. 31: \emph{Let not thine hand be stretched out to receive, and shut when thou shouldest repay.} The passage is quoted in the \emph{Didaché}, c. iv., Barnab. \emph{Epist}. c. xix., and \emph{Constit. Apost}. VII. 11.}. God gives all things freely. \emph{For if any of you lack wisdom, let him ask of God who giveth}\footnote{James i. 5.}, and he shall receive. May He through your prayer grant utterance to us who speak, and faith to you who hear.

9. Let us then seek the testimonies to the Passion of Christ: for we are met together, not now to make a speculative exposition of the Scriptures, but rather to be certified of the things which we already believe. Now thou hast received from me first the testimonies concerning the coming of Jesus; and concerning His walking on the sea, for it is written, \emph{Thy way is in the sea}\footnote{Ps. lxxvii. 19. The Benedictine Editor, with no authority but the Latin version by Grodecq, inserts a quotation of Job ix. 8: \emph{Who walketh on the sea, as on a pavement.} Cf. xi. 23.}. Also concerning divers cures thou hast on another occasion received testimony. Now therefore I begin from whence the Passion began. Judas was the traitor, and he came against Him, and stood, speaking words of peace, but plotting war. Concerning him, therefore, the Psalmist says, \emph{My friends and My neighbours drew near against Me, and stood}\footnote{Ps. xxxviii. 11.}. And again, \emph{Their words were softer than oil, yet be they spears}\footnote{Ib. lv. 21.}. \emph{Hail, Master}\footnote{Matt. xxvi. 49.}; yet he was betraying his Master to death; he was not abashed at His warning, when He said, \emph{Judas, betrayest than the Son of Man with a kiss}\footnote{Luke xxii. 48.}? for what He said to him was just this, Recollect thine own name; Judas means \emph{confession}\footnote{Cf. Phil. Jud. \emph{de Plantatione Noë}, II § 33: “And his name was called Judah, which being interpreted is “confession to the Lord.” In Gen. xlix. 8 the name is differently interpreted: “Judah, thou art he whom thy brethren shall praise.” The root has both senses “to confess,” and “to praise,” which are closely allied since to “confess” is to “give God the glory” (Josh. vii. 19).}; thou hast covenanted, thou hast received the money, make confession quickly. \emph{O God, pass not over My praise in silence; for the mouth of the wicked, and the mouth of the deceitful, are opened against Me; they have spoken against Me with a treacherous tongue, they have compassed Me about also with words of hatred}\footnote{Ps. cix. 1–3.}. But that some of the chief-priests also were present, and that He was put in bonds before the gates of the city, thou hast heard before, if thou rememberest the exposition of the Psalm, which has told the time and the place; how \emph{they returned at evening, and hungered like dogs, and encompassed the city}\footnote{Ps. lix. 6. The exposition was probably given in a sermon preached to the whole congregation, not in these Lectures.}.

10. Listen also for the thirty pieces of silver. \emph{And I will say to them, If it be good in your sight, give me my price, or refuse}\footnote{Zech. xi. 12.}, and the rest. One price is owing to Me from you for My healing the blind and lame, and I receive another; for thanksgiving, dishonour, and for worship, insult. Seest thou how the Scripture foresaw these things? \emph{And they weighed for My price thirty pieces of silver}\footnote{Ib.}. How exact the prophecy! how great and unerring the wisdom of the Holy Ghost! For he said, not ten, nor twenty, but thirty, exactly as many as there were. Tell also what becomes of this price, O Prophet! Does he who received it keep it? or does he give it back? and after he has given it back, what becomes of it? The Prophet says then, \emph{And I took the thirty pieces of silver, and cast them into the house of the Lord, into the foundry}\footnote{Ib. xi. 13.}. Compare the Gospel with the Prophecy: \emph{Judas}, it says, \emph{repented himself, and cast down the pieces of silver in the temple, and departed}\footnote{Matt. xxvii. 3, 5.}.

11. But now I have to seek the exact solution of this seeming discrepancy. For they who make light of the prophets, allege that the Prophet says on the one hand, \emph{And I cast them into the house of the Lord, into the foundry}, but the Gospel on the other hand, \emph{And they gave them for the potter’s field}\footnote{Matt. xxvii. 3, 7.}. Hear then how they are both true. For those conscientious Jews forsooth, the high-priests of that time, when they saw that Judas repented and said, \emph{I have sinned, in that I have betrayed innocent blood}, reply, \emph{What is that to us, see thou to that}\footnote{Ib. v. 4.}. Is it then nothing to you, the crucifiers? but shall he who received and restored the price of murder see to it, and shall ye the murderers not see to it? Then they say among themselves, \emph{It is not lawful to cast them into the treasury, because it is the price of blood}\footnote{Ib. v. 6.}. Out of your own mouths is your condemnation; if the price is polluted, the deed is polluted also: but if thou art fulfilling righteousness in crucifying Christ, why receivest thou not the price of it? But the point of iniquity is this: how is there no disagreement, if the Gospel says, \emph{the potter’s field}, and the Prophet, \emph{the foundry}? Nay, but not only people who are goldsmiths, or brass-founders, have a foundry, but potters also have foundries for their clay. For they sift off the fine and rich and useful earth from the gravel, and separate from it the mass of the refuse matter, and temper the clay first with water, that they may work it with ease into the forms intended. Why then wonderest thou that the Gospel says plainly \emph{the potter’s field}, whereas the Prophet spoke his prophecy like an enigma, since prophecy is in many places enigmatical?

12. They bound Jesus, and brought Him into the hall of the High-priest. And wouldest thou learn and know that this also is written? Esaias says, \emph{Woe unto their soul, for they have taken evil counsel against themselves, saying, Let us bind the Just, for He is troublesome to us}\footnote{Isa. iii. 9: (R.V.) \emph{they have rewarded evil unto themselves. Say ye of the righteous, that it shall be well with him}/ In the Septuagint, from which Cyril quotes, there is an evident interpolation of Wisdom ii. 12: \emph{Let us lie in wait for the righteous; because he is not for our turn} (\textgreek{δύσχρηστος}, as in Cyril).}. And truly, \emph{Woe unto their soul! } Let us see how Esaias was sawn asunder, yet after this the people was restored. Jeremias was cast into the mire of the cistern, yet was the wound of the Jews healed; for the sin was less, since it was against man. But when the Jews sinned, not against man, but against God in man’s nature, \emph{Woe unto their soul}!—\emph{Let us bind the Just}; could He not then set Himself free, some one will say; He, who freed Lazarus from the bonds of death on the fourth day, and loosed Peter from the iron bands of a prison? Angels stood ready at hand, saying, \emph{Let us burst their bands in sunder}\footnote{Ps. ii. 3.}; but they hold back, because their Lord willed to undergo it. Again, He was led to the judgment-seat before the Elders; thou hast already the testimony to this, \emph{The Lord Himself will come into judgment with the ancients of His people, and with the princes thereof}\footnote{Isa. iii. 14.}.

13. But the High-priest having questioned Him, and heard the truth, is wroth; and the wicked officer of wicked men smites Him; and the countenance, which had shone as the sun, endured to be smitten by lawless hands. Others also come and spit on the face of Him, who by spittle had healed the man who was blind from his birth. \emph{Do ye thus requite the Lord? This people is foolish and unwise}\footnote{Deut. xxxii. 6.}. And the Prophet greatly wondering, says, \emph{Lord, who hath believed our report}\footnote{Isa. liii. 1.}? for the thing is incredible, that God, the Son of God, and \emph{the Arm of the Lord}\footnote{Ibid.}, should suffer such things. But that they who are being saved may not disbelieve, the Holy Ghost writes before, in the person of Christ, who says, (for He who then spake these things, was afterward Himself an actor in them,) \emph{I gave My back to the scourges}; (for Pilate, \emph{when he had scourged Him, delivered Him to be crucified}\footnote{Isa. l. 6.; Matt. xxvii. 26.};) \emph{and My cheeks to smitings; and My face I turned not away from the shame of spittings;} saying, as it were, “Though knowing before that they will smite Me, I did not even turn My cheek aside; for how should I have nerved My disciples against death for truth’s sake, had I Myself dreaded this?” I said. \emph{He that loveth his life shall lose it}\footnote{John xii. 25.}: if I had loved My life, how was I to teach without practising what I taught? First then, being Himself God, He endured to suffer these things at the hands of men; that after this, we men, when we suffer such things at the hands of men for His sake, might not be ashamed. Thou seest that of these things also the prophets have clearly written beforehand. Many, however, of the Scripture testimonies I pass by for want of time, as I said before; for if one should exactly search out all, not one of the things concerning Christ would be left without witness.

14. Having been bound, He came from Caiaphas to Pilate,—is this too written? yes; \emph{And having bound Him, they led Him away as a present to the king of Jarim}\footnote{Hosea x. 6: (R.V.) \emph{It also shall be carried unto Assyria for a present to king Jareb.} This passage is applied in the same manner to Luke xxiii. 7 by Justin M. (\emph{Tryph}. § 103), Tertullian (\emph{c.} \emph{Marcion}. iv. 42), and Rufinus (\emph{de Symbolo}, § 21), who adds,—“And rightly does the Prophet add the name ‘Jarim,’ which means ‘a wild vine,’ for Herod was…a wild vine, i.e. of an alien stock.” For the various interpretations of the name see the Commentaries on Hosea v. 13, and x. 6; Schrader, \emph{Cuneijorm Inscriptions}, II. § 439, Driver, \emph{Introduction to O. T. Literature}, p. 283.}. But here some sharp hearer will object, “Pilate was not a king,” (to leave for a while the main parts of the question,) “how then having bound Him, led they Him as a present to the king?” But read thou the Gospel; \emph{When Pilate heard that He was of Galilee, he sent Him to Herod}\footnote{Luke xxiii. 6, 7.}; for Herod was then king, and was present at Jerusalem. And now observe the exactness of the Prophet; for he says, that He was sent as a present; for \emph{the same day Pilate and Herod were made friends together, for before they were at enmity}\footnote{Ibid. xxiii. 12.}. For it became Him who was on the eve of making peace between earth and heaven, to make the very men who condemned Him the first to be at peace; for the Lord Himself was there present, \emph{who reconciles}\footnote{Job xii. 24: (R.V.) \emph{He taketh away the heart of the chiefs of the people of the earth.} The rendering “who reconciles” (\textgreek{ὁ διαλλάσσων,} Sept.) is forbidden by the context.} \emph{the hearts of the princes of the earth}. Mark the exactness of the Prophets, and their true testimony.

15. Look with awe then at the Lord who was judged. He suffered Himself to be led and carried by soldiers. Pilate sat in judgment, and He who sitteth on the right hand of the Father, stood and was judged\footnote{Some \textsc{mss.} have \textgreek{ἠνεσχετο} or \textgreek{ἠνείχετο}, “He submitted to stand.”}. The people whom He had redeemed from the land of Egypt, and oftimes from other places, shouted against Him, \emph{Away with Him, away with Him, crucify Him}\footnote{Josh. xix. 15.}. Wherefore, O ye Jews? because He healed your blind? or because He made your lame to walk, and bestowed His other benefits? So that the Prophet in amazement speaks of this too, \emph{Against whom have ye opened your mouth, and against whom have ye let loose your tongue}\footnote{Isa. lvii. 4.}? and the Lord Himself says in the Prophets, \emph{Mine heritage became unto Me as a lion in the forest; it gave its voice against Me; therefore have I hated it}\footnote{Jer. xii. 8.}. I have not refused them, but they have refused Me; in consequence thereof I say, \emph{I have forsaken My house}\footnote{Ibid. v. 7.}.

16. When He was judged, He held His peace; so that Pilate was moved for Him, and said, \emph{Hearest Thou not what these witness against Thee}\footnote{Matt. xxvii. 13.}? Not that He knew Him who was judged, but he feared his own wife’s dream which had been reported to him. And Jesus held His peace. The Psalmist says, \emph{And I became as a man that heareth not; and in whose mouth are no reproofs}\footnote{Ps. xxxviii. 14.}; and again, \emph{But I was as a deaf man and heard not; and as a dumb man that openeth not his mouth}\footnote{Ibid. v. 13.}. Thou hast before heard concerning this\footnote{“Perhaps in some Homily” (Ben. Ed.).}, if thou rememberest.

17. But the soldiers who crowd around mock Him, and their Lord becomes a sport to them, and upon their Master they make jests. \emph{When they looked on Me, they shaked their heads}\footnote{Ps. cix. 25.}. Yet the figure of kingly state appears; for though in mockery, yet they bend the knee. And the soldiers before they crucify Him, put on Him a purple robe, and set a crown on His head; for what though it be of thorns? Every king is proclaimed by soldiers; and Jesus also must in a figure be crowned by soldiers; so that for this cause the Scripture says in the Canticles, \emph{Go forth, O ye daughters of Jerusalem, and look upon King Solomon in the crown wherewith His mother crowned Him}\footnote{Cant. iii. 11.}. And the crown itself was a mystery; for it was a remission of sins, a release from the curse.

18. Adam received the sentence, \emph{Cursed is the ground in thy labours; thorns and thistles shall it bring forth to thee}\footnote{Gen. iii. 17, 18. By mistaking one letter in the Hebrew, the Seventy give the meaning “in thy labours” instead of “for thy sake.”}. For this cause Jesus assumes the thorns, that He may cancel the sentence; for this cause also was He buried in the earth, that the earth which had been cursed might receive the blessing instead of a curse. At the time of the sin, they clothed themselves with fig-leaves; for this cause Jesus also made the fig-tree the last of His signs. For when about to go to His passion, He curses the fig-tree, not every fig-tree, but that one alone, for the sake of the figure; saying, \emph{No more let any man eat fruit of thee}\footnote{Mark xi. 1.}; let the doom be cancelled. And because they aforetime clothed themselves with fig-leaves, He came at a season when food was not wont to be found on the fig-tree. Who knows not that in winter-time the fig-tree bears no fruit, but is clothed with leaves only? Was Jesus ignorant of this, which all knew? No, but though He knew, yet He came as if seeking; not ignorant that He should not find, but shewing that the emblematical curse extended to the leaves only.

19. And since we have touched on things connected with Paradise, I am truly astonished at the truth of the types. In Paradise was the Fall, and in a Garden was our Salvation. From the Tree came sin, and until the Tree sin lasted. \emph{In the evening, when the Lord walked in the Garden, they hid themselves}\footnote{Gen. iii. 8.}; and in the evening the robber is brought by the Lord into Paradise. But some one will say to me, “Thou art inventing subtleties; shew me from some prophet the Wood of the Cross; except thou give me a testimony from a prophet, I will not be persuaded. Hear then from Jeremias, and assure thyself; \emph{I was like a harmless lamb led to be slaughtered; did I not know it}\footnote{Jer. xi. 19: \emph{I was like a tame} (R.V. \emph{gentle}) \emph{lamb that is led to the slaughter; and I knew not that they had devised devices against me.} Cyril’s interrogative rendering is not admissible.}? (for in this manner read it as a question, as I have read it; for He who said, \emph{Ye know that after two days comes the passover, and the Son of Man is betrayed to be crucified}\footnote{Matt. xxvi. 2.}, did He not know?) \emph{I was like a harmless lamb led to be slaughtered; did I not know it?} (but what sort of lamb? let John the Baptist interpret it, when he says, Behold the Lamb of God, that taketh away the sin of the world\footnote{John i. 29.}.) \emph{They devised against Me a wicked device, saying}\footnote{Jer. xi. 19.}—(He who knows the devices, knew He not the result of them? And what said they?)—\emph{Come, and let us place a beam upon His bread}\footnote{Ibid. R.V. \emph{Let us destroy the tree with the fruit thereof.} The word rendered \emph{fruit} is literally \emph{bread}. The phrase is evidently proverbial. The Hebrew word which means “destroy” is misinterpreted by \textgreek{ἐμβάλωμεν} in the Greek. Hence arose the fanciful application of the passage to the cross laid on the body of Christ to be borne by Him. Justin M. (\emph{Tryph}. lxxii.) charges the Jews with having recently cut out the passage because of the supposed reference to Christ. Tertullian (\emph{adv. Judæos}, c. 10) writes: “Of course on His body that ‘wood’ was put; for so Christ has revealed, calling His body ‘bread.’” He gives the same interpretation elsewhere (\emph{adv. Marcion}. III. 19; IV. 40). Cf. Cyprian (\emph{Testimonia ad Quirinum}, Lib. II. 15); Athanas. (\emph{de Incarn}. § 33).}—(and if the Lord reckon thee worthy, thou shalt hereafter learn, that His body according to the Gospel bore the figure of bread;)—\emph{Come} then, \emph{and let us place a beam upon His bread, and cut Him off out of the land of the living;}—(life is not cut off, why labour ye for nought?)—\emph{And His name shall be remembered no more.} Vain is your counsel; for \emph{before the sun His Name}\footnote{Ps. lxxii. 17.} abideth in the Church. And that it was Life, which hung on the Cross, Moses says, weeping, \emph{And thy life shall be hanging before thine eyes; and thou shalt be afraid day and night, and thou shalt not trust thy life}\footnote{Deut. xxviii. 66.}. And so too, what was just now read as the text, \emph{Lord, who hath believed our report?}

20. This was the figure which Moses completed by fixing the serpent to a cross, that whoso had been bitten by the living serpent, and looked to the brasen serpent, might be saved by believing\footnote{Num. xxi. 9; John iii. 14. The Jerusalem Editor asks, “How did Moses complete the figure by fixing the serpent to a cross? First he set up the wood and fixed it in the earth as a post: then by putting the brazen serpent athwart (\textgreek{πλαγίως}, he formed a figure of the Cross.” Cf. Barnab. \emph{Epist}. c. xii.; Justin M. (\emph{Apol}. i. c. 60); Iren. (\emph{Hæres.} IV. c. 2); Tertull. \emph{adv. Judæos}, c. 10).}. Does then the brazen serpent save when crucified, and shall not the Son of God incarnate save when crucified also? On each occasion life comes by means of wood. For in the time of Noe the preservation of life was by an ark of wood. In the time of Moses the sea, on beholding the emblematical rod, was abashed at him who smote it; is then Moses’ rod mighty, and is the Cross of the Saviour powerless? But I pass by the greater part of the types, to keep within measure. The wood in Moses’ case sweetened the water; and from the side of Jesus the water flowed upon the wood.

21. The beginning of signs under Moses was blood and water; and the last of all Jesus’ signs was the same. First, Moses changed the river into blood; and Jesus at the last gave forth from His side water with blood. This was perhaps on account of the two speeches, his who judged Him, and theirs who cried out against Him; or because of the believers and the unbelievers. For Pilate said, \emph{I am innocent} and washed his hands in water; they who cried out against Him said, \emph{His blood be upon us}\footnote{Matt. xxvii. 24, 25.}: there came therefore these two out of His side; the water, perhaps, for him who judged Him; but for them that shouted against Him the blood. And again it is to be understood in another way; the blood for the Jews, and the water for the Christians: for upon them as plotters came the condemnation from the blood; but to thee who now believest, the salvation which is by water. For nothing has been done without a meaning. Our fathers who have written comments have given another reason of this matter. For since in the Gospels the power of salutary Baptism is twofold, one which is granted by means of water to the illuminated, and a second to holy martyrs, in persecutions, through their own blood, there came out of that saving Side blood and water\footnote{John xix. 34. Cf. Cat. iii. 10. Origen (\emph{In Lib. Judic}. Hom. vii. § 2): “It is the Baptism of blood alone that can render us purer than the Baptism of water has done.” Cf. Origen (\emph{in Ev. Matt}. Tom. xvi. 6): “If Baptism promises remission of sins, as we have received concerning Baptism in water and the Spirit, and if one who has endured the Baptism of Martyrdom receives remission of sins, then with good reason martyrdom may be called a Baptism.” For a summary of the “Patristic Interpretation” of the passage, see Bp. Westcott, \emph{Speaker’s Commentary}.)}, to confirm the grace of the confession made for Christ, whether in baptism, or on occasions of martyrdom. There is another reason also for mentioning the Side. The woman, who was formed from the side, led the way to sin; but Jesus who came to bestow the grace of pardon on men and women alike, was pierced in the side for women, that He might undo the sin.

22. And whoever will inquire, will find other reasons also; but what has been said is enough, because of the shortness of the time, and that the attention of my hearers may not become sated. And yet we never can be tired of hearing concerning the crowning of our Lord, and least of all in this most holy Golgotha. For others only hear, but we both see and handle. Let none be weary; take thine armour against the adversaries in the cause of the Cross itself; set up the faith of the Cross as a trophy against the gainsayers. For when thou art going to dispute with unbelievers concerning the Cross of Christ, first make with thy hand the sign of Christ’s Cross, and the gainsayer will be silenced. Be not ashamed to confess the Cross; for Angels glory in it, saying, \emph{We know whom ye seek, Jesus the Crucified}\footnote{Matt. xxviii. 5.}. Mightest thou not say, O Angel, “I know whom ye seek, my Master?” But, “I,” he says with boldness, “I know the Crucified.” For the Cross is a Crown, not a dishonour.

23. Now let us recur to the proof out of the Prophets which I spoke of. The Lord was crucified; thou hast received the testimonies. Thou seest this spot of Golgotha! Thou answerest with a shout of praise, as if assenting. See that thou recant not in time of persecution. Rejoice not in the Cross in time of peace only, but hold fast the same faith in time of persecution also; be not in time of peace a friend of Jesus, and His foe in time of wars. Thou receivest now remission of thy sins, and the gifts of the King’s spiritual bounty; when war shall come, strive thou nobly for thy King. Jesus, the Sinless, was crucified for thee; and wilt not thou be crucified for Him who was crucified for thee? Thou art not bestowing a favour, for thou hast first received; but thou art returning a favour, repaying thy debt to Him who was crucified for thee in Golgotha. Now Golgotha is interpreted, “the place of a skull.” Who were they then, who prophetically named this spot Golgotha, in which Christ the true Head endured the Cross? As the Apostle says, \emph{Who is the Image of the Invisible God;} and a little after, \emph{and He is the Head of the body, the Church}\footnote{Col. i. 15, 18.}. And again, \emph{The Head of every man is Christ}\footnote{1 Cor. xi. 3.}; and again, \emph{Who is the Head of all principality and power}\footnote{Col. ii. 10.}. The Head suffered in “the place of the skull.” O wondrous prophetic appellation! The very name also reminds thee, saying, “Think not of the Crucified as of a mere man; He is \emph{the Head of all principality and power}. That Head which was crucified is the Head of all power, and has for His Head the Father; \emph{for the Head of the man is Christ, and the Head of Christ is God}\footnote{1 Cor. xi. 3.}.”

24. Christ then was crucified for us, who was judged in the night, when it was cold, and therefore a \emph{fire of coals}\footnote{John xviii. 18.} was laid. He was crucified at the third hour; \emph{and from the sixth} \emph{hour there was darkness until the ninth hour}\footnote{Matt. xxvii. 45.}; but from the ninth hour there was light again. Are these things also written? Let us inquire. Now the Prophet Zacharias says, \emph{And it shall come to pass in that day, that there shall not be light, and there shall be cold and frost one day;} (the cold on account of which Peter warmed himself;) \emph{And that day shall be known unto the Lord}\footnote{Zech. xiv. 6, 7.}; (what, knew He not the other days? days are many, but \emph{this is the day} of the Lord’s patience, \emph{which the Lord made}\footnote{Ps. cxviii. 24.};)—\emph{And that day shall be known unto the Lord, not day, and not night: } what is this dark saying which the Prophet speaks? That day is neither day nor night? what then shall we name it? The Gospel interprets it, by relating the event. It was not day; for the sun shone not uniformly from his rising to his setting, but from the sixth hour till the ninth hour, there was darkness at mid-day. The darkness therefore was interposed; but \emph{God called the darkness night}\footnote{Gen. i. 5.}. Wherefore it was neither day nor night: for neither was it all light, that it should be called day; nor was it all darkness, that it should be called night; but after the ninth hour the sun shone forth. This also the Prophet foretels; for after saying, \emph{Not day, nor night,} he added, \emph{And at evening time it shall be light}\footnote{Zech. xiv. 7. Cf. Euseb. (\emph{Dem. Evang.} x. 7): “It was not day, because of the noon-tide darkness: and again it was not night, because of the day which followed upon it, which he represented by a sign in saying, \emph{at evening time there shall be light}.}. Seest thou the exactness of the prophets? Seest thou the truth of the things which were written aforetime?

25. But dost thou ask exactly at what hour the sun failed\footnote{\textgreek{ἐξέλιπεν}. See Cat. x. 19, note 2. \emph{Acta Pilati}. c. xi.}? was it the fifth hour, or the eighth, or the tenth? Tell, O Prophet, the exact time thereof to the Jews, who are unwilling to hear; when shall the sun go down? The Prophet Amos answers, \emph{And it shall come to pass in that day, saith the Lord God, that the sun shall go down at noon} (for there was darkness from the sixth hour;) \emph{and the light shall grow dark over the earth in the day}\footnote{Amos viii. 9. Cf. Euseb. (\emph{Dem. Ev}. x. 6).}.” What sort of season is this, O Prophet, and what sort of day? \emph{And I will turn your feasts into mourning}; for this was done in the days of unleavened bread, and at the feast of the Passover: then afterwards he says, \emph{And I will make Him as the mourning of an Only Son, and those with Him as a day of anguish}\footnote{Amos viii. 10.}; for in the day of unleavened bread, and at the feast, their women were wailing and weeping, and the Apostles had hidden themselves and were in anguish. Wonderful then is this prophecy.

26. But, some one will say, “Give me yet another sign; what other exact sign is there of that which has come to pass? Jesus was crucified; and He wore but one coat, and one cloak: now His cloak the soldiers shared among themselves, having rent it into four; but His coat was not rent, for when rent it would have been no longer of any use; so about this lots are cast by the soldiers; thus the one they divide, but for the other they cast lots. Is then this also written? They know, the diligent chanters\footnote{Synod of Laodicea, Can. xvi. 15: “Besides the appointed singers, who mount the ambo and sing from the book, others shall not sing in the Church.” Hefele thinks that this was not intended to forbid the laity to take any part in the Church music, but only to forbid those who were not cantors to take the lead. See Bingham, \emph{Antiquities}, III. c. 7; XIV. c. 1.} of the Church, who imitate the Angel hosts, and continually sing praises to God: who are thought worthy to chant Psalms in this Golgotha, and to say, \emph{They parted My garments among them, and upon My vesture they did cast lots}\footnote{Ps. xxii. 18, quoted in John xix. 24.}. The “lots” were what the soldiers cast\footnote{\textgreek{κλῆρος δὲ ἦν ὁ λαχμός}. Bishop Hall, \emph{Contemplations}, Book IV. 32, speaks of the soldiers’ “barbarous \emph{sortitions}.” The technical term is “sortilege.” Cf. \emph{Evang. Pet}. § 4; Justin M. \emph{Dial}. 97.}.

27. Again, when He had been judged before Pilate, He was clothed in red; for there they put on Him a purple robe. Is this also written? Esaias saith, \emph{Who is this that cometh from Edom? the redness of His garments is from Bosor}\footnote{Isa. lxiii. 1, 2.}; (who is this who in dishonor weareth purple? For Bosor has some such meaning in Hebrew\footnote{Bozrah means a “sheepfold,” and is the name of a city in Idumea. Cyril’s interpretation rests on a false derivation.}.) \emph{Why are Thy garments red, and Thy raiment as from a trodden wine-press? } But He answers and says, \emph{All day long have I stretched forth Mine hands unto a disobedient and gainsaying people}\footnote{Isa. lxv. 2. “It is a commonplace in patristic literature that the Crucifixion was prefigured by Isa. lxv. 2.” (Dr. C. Taylor, \emph{Hermas and the Four Gospels}, p. 49.) Cf. Barnab. \emph{Epist}. c. xii.; \emph{Didache} xvi.; Justin M. (\emph{Apolog}. I. c. 35; \emph{Tryph}. cc. 97, 114); Tertull. (\emph{contra Jud}. xii); Irenæ. IV. xxxiii. 12.}.

28. He stretched out His hands on the Cross, that He might embrace the ends of the world; for this Golgotha is the very centre of the earth. It is not my word, but it is a prophet who hath said, \emph{Thou hast wrought salvation in the midst of the earth}\footnote{Ps. lxxiv. 12. The passage does not refer to Palestine especially: “in the midst of the earth” is equivalent to “in the sight of all nations.” Cf. \emph{Orac. Sibyll}. viii. 302: “He shall spread out His hands, and span the whole world,” quoted by Dr. Taylor, “The Teaching,” p. 103.}. He stretched forth human hands, who by His spiritual hands had established the heaven; and they were fastened with nails, that His manhood, which here the sins of men, having been nailed to the tree, and having died, sin might die with it, and we might rise again in righteousness. \emph{For since by one man came death, by} One \emph{Man came also life}\footnote{Rom. v. 12, 17.}; by One Man, the Saviour, dying of His own accord: for remember what He said, \emph{I have power to} \emph{lay down My life, and I have power to take it again}\footnote{John x. 18.}.

29. But though He endured these things, having come for the salvation of all, yet the people returned Him an evil recompense. Jesus saith, \emph{I thirst}\footnote{Ib. ix. 28.},—He who had brought forth the waters for them out of the craggy rock; and He asked fruit of the Vine which He had planted. But what does the Vine? This Vine, which was by nature of the holy fathers, but of Sodom by purpose of heart;—(for \emph{their Vine is of Sodom, and their tendrils of Gomorrah}\footnote{Deut. xxxii. 32.};)—this Vine, when the Lord was athirst, having filled a sponge and put it on a reed, offers Him vinegar. \emph{They gave Me also gall for My meat, and in My thirst, they gave Me vinegar to drink}\footnote{Ps. lxix. 21.}. Thou seest the clearness of the Prophets’ description. But what sort of gall put they into My mouth? \emph{They gave Him,} it says, \emph{wine mingled with myrrh}\footnote{Mark xv. 23.}. Now myrrh is in taste like gall, and very bitter. Are these things what ye recompense unto the Lord? Are these thy offerings, O Vine, unto thy Master? Rightly did the Prophet Esaias aforetime bewail you, saying, \emph{My well-beloved had a vineyard in a hill in a fruitful place;} and (not to recite the whole) \emph{I waited}, he says, \emph{that it should bring forth grapes}; I thirsted that it should give wine; \emph{but it brought forth thorns}\footnote{Isa. v. 1, 2.}; for thou seest the crown, wherewith I am adorned. What then shall I now decree? I will command the clouds that they rain no rain upon it\footnote{Ib. \emph{v}. 6. Cf. Tertull. \emph{adv. Marcion.} III. c. 23; \emph{contra Jud}. c. 13: “The clouds being celestial benefits which were commanded not to be forthcoming to the house of Israel; for it ‘had borne \emph{thorns},’ whereof that house of Israel had wrought a crown for Christ.” \emph{Constitt. Apost}. VI. § 5: “He has taken away from them the Holy Spirit and the prophetic rain, and has replenished His Church with spiritual grace.”}. For the clouds which are the Prophets were removed from them, and are for the future in the Church; as Paul says, \emph{Let the Prophets speak two or three, and let the others judge}\footnote{1 Cor. xiv. 29.}; and again, \emph{God gave in the Church, some, Apostles, and some, Prophets}\footnote{Eph. iv. 11.}. Agabus, who bound his own feet and hands, was a prophet.

30. Concerning the robbers who were crucified with Him, it is written, \emph{And He was numbered with the transgressors}\footnote{Isa. liii. 12.}. Both of them were before this transgressors, but one was so no longer. For the one was a transgressor to the end, stubborn against salvation; who, though his hands were fastened, smote with blasphemy by his tongue. When the Jews passing by wagged their heads, mocking the Crucified, and fulfilling what was written, \emph{When they looked on Me, they shaked their heads}\footnote{Ps. cix. 25.}, he also reviled with them. But the other rebused the reviler; and it was to him the end of life and the beginning of restoration; the surrender of his soul a first share in salvation. And after rebuking the other, he says, \emph{Lord, remember me}\footnote{Luke xxiii. 40. ff.}; for with Thee is my account. Heed not this man, for the eyes of his understanding are blinded; but remember me. I say not, remember my works, for of these I am afraid. Every man has a feeling for his fellow-traveller; I am travelling with Thee towards death; remember me, Thy fellow-wayfarer. I say not, Remember me now, but, \emph{when Thou comest in Thy kingdom}.

31. What power, O robber, led thee to the light? Who taught thee to worship that despised Man, thy companion on the Cross? O Light Eternal, which gives light to them that are in darkness! Therefore also he justly heard the words, \emph{Be of good cheer}\footnote{\textgreek{θάρσει}. An addition to the text of Luke xxiii. 43 in Codex Bezae.}; not that thy deeds are worthy of good cheer; but that the King is here, dispensing favours. The request reached unto a distant time; but the grace was very speedy. \emph{Verily I say unto thee, This day shalt thou be with Me in Paradise;} because \emph{to-day} thou hast \emph{heard My voice,} and hast not \emph{hardened thine heart}\footnote{Ps. xcv. 7, 8.}. Very speedily I passed sentence upon Adam, very speedily I pardon thee. To him it was said, \emph{In the day wherein ye eat, ye shall surely die}\footnote{Gen. ii. 17.}; but thou to-day hast obeyed the faith, to-day is thy salvation. Adam by the Tree fell away; thou by the Tree art brought into Paradise. Fear not the serpent; he shall not cast thee out; for he is \emph{fallen from heaven}\footnote{Luke x. 18.}. And I say not unto thee, This day shalt thou depart, but, \emph{This day shalt thou be with Me}. Be of good courage: thou shalt not be cast out. Fear not the flaming sword; it shrinks from its Lord\footnote{Gen. iii. 24. S. Ambrose (\emph{Ps}. cxix. Serm. xx. § 12): “All who desire to return to Paradise must be tried by fire: for not in vain the Scripture saith that when Adam and Eve were driven out of their abode in Paradise, God placed at the gate of Eden a flaming sword which turned every way.”}. O mighty and ineffable grace! The faithful Abraham had not yet entered, but the robber enters\footnote{Cf. Iren. V. c. 5, § 1; Athan. (\emph{Expos. Fid}. c. i.): “He shewed us.…an entrance into Paradise from which Adam was cast out, and into which he entered again by means of the thief.” S. Leo (\emph{de Pass. Dom}. Serm. II. c. 1): “Excedit humanam conditionem ista promissio: nec tam de ligno Crucis, quam de throno editur potestatis.”}! Moses and the Prophets had not yet entered, and the robber enters though a breaker of the law. Paul also wondered at this before thee, saying, \emph{Where sin abounded, there grace did much more abound}\footnote{Rom. v. 20.}. They who had borne the heat of the day had not yet entered; and he of the eleventh hour entered. Let none murmur against the goodman of the house, for he says, \emph{Friend, I do thee no wrong; is it not} \emph{lawful for Me to do what I will with Mine own}\footnote{Matt. xx. 12 ff.}? The robber has a will to work righteousness, but death prevents him; I wait not exclusively for the work, but faith also I accept. I am come who \emph{feed My sheep among the lilies}\footnote{Cant. vi. 3.}, I am come to feed them in the gardens. I have \emph{found} a \emph{sheep that was lost}\footnote{Luke xv. 5, 6.}, but I lay it on My shoulders; for he believes, since he himself has said, \emph{I have gone astray like a lost sheep}\footnote{Ps. cxix. 176.}; \emph{Lord, remember me when Thou comest in Thy kingdom.}

32. Of this garden I sang of old to My spouse in the Canticles, and spoke to her thus. \emph{I am come into My garden, My sister, My spouse}\footnote{Cant. v. 1.}; (\emph{now in the place where He was crucified was a garden}\footnote{John xix. 41.};) and what takest Thou thence? \emph{I have gathered My myrrh;} having drunk wine mingled with myrrh, and vinegar, after receiving which, He said, It is finished\footnote{Ib. 30.}. For the mystery has been fulfilled; the things that are written have been accomplished; sins are forgiven. For \emph{Christ being come an High-Priest of the good things to come, by the greater and more perfect tabernacle, not made with hands, that is to say, not of this creation, nor yet by the blood of goats and calves, but by His own blood, entered in once for all into the holy place, having obtained eternal redemption; for if the blood of bulls and of goats, and the ashes of an heifer, sprinkling the defiled, sanctifieth to the purifying of the flesh, how much more the blood of Christ}\footnote{Heb. ix. 11.}? And again, \emph{Having therefore, brethren, boldness to enter into the holiest by the blood of Jesus, by a new and living way, which He hath consecrated for us, through the veil, that is to say, His flesh}\footnote{Ib. x. 19.}. And because His flesh, this veil, was dishonoured, therefore the typical veil of the temple was rent asunder, as it is written, \emph{And, behold, the veil of the temple was rent in twain from the top to the bottom}\footnote{Matt. xxvii. 51.}; for not a particle of it was left; for since the Master said, \emph{Behold, your house is left unto you desolate}\footnote{Ib. xxiii. 38.}, the house brake all in pieces.

33. These things the Saviour endured, \emph{and made peace through the Blood of His Cross, for things in heaven, and things in earth}\footnote{Col. i. 20.}. For we were enemies of God through sin, and God had appointed the sinner to die. There must needs therefore have happened one of two things; either that God, in His truth, should destroy all men, or that in His loving-kindness He should cancel the sentence. But behold the wisdom of God; He preserved both the truth of His sentence, and the exercise of His loving-kindness. Christ took our sins in \emph{His body on the tree, that we by His death might die to sin, and live unto righteousness}\footnote{1 Pet. ii. 24.}. Of no small account was He who died for us; He was not a literal sheep; He was not a mere man; He was more than an Angel; He was God made man. The transgression of sinners was not so great as the righteousness of Him who died for them; the sin which we committed was not so great as the righteousness which He wrought who laid down His life for us,—who laid it down when He pleased, and took it again when He pleased. And wouldest thou know that He laid not down His life by violence, nor yielded up the ghost against His will? He cried to the Father, saying, \emph{Father, into Thy hands I commend My spirit}\footnote{Luke xxiii. 46.}; I commend it, that I may take it again. And having said these things, \emph{He gave up the ghost}\footnote{Matt. xxvii. 50.}; but not for any long time, for He quickly rose again from the dead.

34. The Sun was darkened, because of \emph{the Sun of Righteousness}\footnote{Mal. iv. 2.}. Rocks were rent, because of the spiritual Rock. Tombs were opened, and the dead arose, because of Him who was \emph{free among the dead}\footnote{Ps. lxxxviii. 5.}; \emph{He sent forth His prisoners out of the pit wherein is no water}\footnote{Zech. ix. 11.}. Be not then ashamed of the Crucified, but be thou also bold to say, \emph{He beareth our sins, and endureth grief for us, and with His stripes we are healed}\footnote{Isa. liii. 4, 5.}. Let us not be unthankful to our Benefactor. And again; \emph{for the transgression of my people was He led to death; and I will give the wicked for His burial, and the rich for His death}\footnote{Ib. vv. 8, 9.}. Therefore Paul says plainly, \emph{that Christ died for our sins according to the Scriptures, and that He was buried, and that He hath risen again the third day according to the Scriptures}\footnote{1 Cor. xv. 3, 4.}.

35. But we seek to know clearly where He has been buried. Is His tomb made with hands? Is it, like the tombs of kings, raised above the ground? Is the Sepulchre made of stones joined together? And what is laid upon it? Tell us, O Prophets, the exact truth concerning His tomb also, where He is laid, and where we shall seek Him? And they say, \emph{Look into the solid rock which ye have hewn}\footnote{Isa. li. 1.}. \emph{Look in} and behold. Thou hast in the Gospels \emph{In a sepulchre hewn in stone, which was hewn out of a rock}\footnote{Matt. xxvii. 60; Mark xv. 46; Luke xxiii. 50.}. And what happens next? What kind of door has the sepulchre? Again another Prophet says, \emph{They cut off My life in a dungeon}\footnote{Lam. iii. 53: \textgreek{ἐν λάκκῳ}, “in a pit,” or “well.” Cf. Jer. xxxvii. 16.}, \emph{and cast a stone upon Me.} I, who am the \emph{Chief corner-stone, the elect, the precious}\footnote{1 Pet. ii. 6.}, lie for a little time within a stone—I who am a stone of stumbling to the Jews, and of salvation to them who believe. \emph{The Tree of life}\footnote{Gen. ii. 9; iii. 22. Methodius (\emph{Sympos}. ix. c. 3): “He that hath not believed in Christ, nor hath understood that He is the first principle and the Tree of Life, \&c.”}, therefore was planted in the earth, that the earth which had been cursed might enjoy the blessing, and that the dead might be released.

36. Let us not then be ashamed to confess the Crucified. Be the Cross our seal made with boldness by our fingers on our brow, and on everything; over the bread we eat, and the cups we drink; in our comings in, and goings out; before our sleep, when we lie down and when we rise up; when we are in the way, and when we are still\footnote{Cf. Cat. iv. 14, note 3; Euseb. (\emph{Dem. Ev}. ix. 14).}. Great is that preservative; it is without price, for the sake of the poor; without toil, for the sick; since also its grace is from God. It is the Sign of the faithful, and the dread of devils: for He \emph{triumphed over them in it, having made a shew of them openly}\footnote{Col. ii. 15.}; for when they see the Cross they are reminded of the Crucified; they are afraid of Him, who \emph{bruised the heads of the dragon}\footnote{Ps. lxxiv. 13.}. Despise not the Seal, because of the freeness of the gift; out for this the rather honour thy Benefactor.

37. And if thou ever fall into disputation and hast not the grounds of proof, yet let Faith remain firm in thee; or rather, become thou well learned, and then silence the Jews out of the prophets, and the Greeks out of their own fables. They themselves worship men who have been thunderstricken\footnote{See Cat. vi. 11, note 2.}: but the thunder when it comes from heaven, comes not at random. If they are not ashamed to worship men thunderstricken and abhorred of God, art thou ashamed to worship the beloved Son of God, who was crucified for thee? I am ashamed to tell the tales about their so-called Gods, and I leave them because of time; let those who know, speak. And let all heretics also be silenced. If any say that the Cross is an illusion, turn away from him. Abhor those who say that Christ was crucified to our fancy\footnote{\textgreek{κατὰ φαντασίαν}. Cf. Ignat. \emph{Trall}. 9, 10; Cat. iv. 9; xiii. 4.} only; for if so, and if salvation is from the Cross, then is salvation a fancy also. If the Cross is fancy, the Resurrection is fancy also; but \emph{if Christ be not risen, we are yet in our sins}\footnote{1 Cor. xv. 17.}. If the Cross is fancy, the Ascension also is fancy; and if the Ascension is fancy, then is the second coming also fancy, and everything is henceforth unsubstantial.

38. Take therefore first, as an indestructible foundation, the Cross, and build upon it the other articles of the faith. Deny not the Crucified; for, if thou deny Him, thou hast many to arraign thee. Judas the traitor will arraign thee first; for he who betrayed Him knows that He was condemned to death by the chief-priests and elders. The thirty pieces of silver bear witness; Gethsemane bears witness, where the betrayal occurred; I speak not yet of the Mount of Olives, on which they were with Him at night, praying. The moon in the night bears witness; the day bears witness, and the sun which was darkened; for it endured not to look on the crime of the conspirators. The fire will arraign thee, by which Peter stood and warmed himself; if thou deny the Cross, the eternal fire awaits thee. I speak hard words, that thou may not experience hard pains. Remember the swords that came against Him in Gethsemane, that thou feel not the eternal sword. The house of Caiaphas\footnote{The house of Caiaphas and Pilate’s Prætorium (§ 41), and Mount Zion itself (Cat. xvi. 18), on which they both stood are described by Cyril as being in his time ruined and desolate. Eusebius (\emph{Dem. Ev}. VIII. 406), referring to the prophecy of Micah (iii. 12), repeated by Jeremiah (xxvi. 18), that \emph{Zion shall be plowed as a field, and Jerusalem shall become heaps}, testifies that he had seen with his own eyes the place being ploughed and sown by strangers, and adds that in his own time the stones for both public and private buildings were taken from the ruins. The Bordeaux Pilgrim (333 \textsc{a.d.}) says, “It is evident where the house of Caiaphas the Priest was; and there is still the pillar at which Christ was scourged:” this pillar is described by Jerome (\emph{Ep}. 86) as supporting the portico of the Church which by his time had been built on the spot. Prudentius \emph{circ.} 400 \textsc{a.d.}):— \par “Impia blasphemi cecidit domus alta Caiphae.… \par Vinctus in his Dominus stetit ædibus atque columnae \par Annexus tergum dedit ut servile flagellis. \par Perstat adhuc, templumque gerit veneranda columna.” \par (Benedictine Editor.)} will arraign thee, shewing by its present desolation the power of Him who was erewhile judged there. Yea, Caiaphas himself will rise up against thee in the day of judgment, the very servant will rise up against thee, who smote Jesus with the palm of his hand; they also who bound Him, and they who led Him away. Even Herod shall rise up against thee; and Pilate; as if saying, Why deniest thou Him who was slandered before us by the Jews, and whom we knew to have done no wrong? For I Pilate then washed my hands. The false witnesses shall rise up against thee, and the soldiers who arrayed Him in the purple robe, and set on Him the crown of thorns, and crucified Him in Golgotha, and cast lots for His coat. Simon the Cyrenian will cry out upon thee, who bore the Cross after Jesus.

39. From among the stars there will cry out upon thee, the darkened Sun; among the things upon earth, the Wine mingled with myrrh; among reeds, the Reed; among herbs, the Hyssop; among the things of the sea, the Sponge; among trees, the Wood of the Cross;—the soldiers, too, as I have said, who nailed Him, and cast lots for His vesture; the soldier who pierced His side with the spear; the women who then were present; the veil of the temple then rent asunder; the hall of Pilate, now laid waste by the power of Him who was then crucified; this holy Golgotha, which stands high above us, and shews itself to this day, and displays even yet how because of Christ the rocks were then riven\footnote{Cf. Lucian. Antioch. ap. Rufin. \emph{Hist. Eccl.} ix. c. 6; “Golothana rupes sub patibuli onere disrupta.”}; the sepulchre nigh at hand where He was laid; and the stone which was laid on the door, which lies to this day by the tomb; the Angels who were then present; the women who worshipped Him after His resurrection; Peter and John, who ran to the sepulchre; and Thomas, who thrust his hand into His side, and his fingers into the prints of the nails. For it was for our sakes that he so carefully handled Him; and what thou, who wert not there present, wouldest have sought, he being present, by God’s Providence, did seek.

40. Thou hast Twelve Apostles, witnesses of the Cross; and the whole earth, and the world of men who believe on Him who hung thereon. Let thy very presence here now persuade thee of the power of the Crucified. For who now brought thee to this assembly? what soldiers? With what bonds wast thou constrained? What sentence held thee fast here now? Nay, it was the Trophy of salvation, the Cross of Jesus that brought you all together. It was this that enslaved the Persians, and tamed the Scythians; this that gave to the Egyptians, for cats and dogs and their manifold errors, the knowledge of God; this, that to this day heals diseases; that to this day drives away devils, and overthrows the juggleries of drugs and charms.

41. This shall appear again with Jesus from heaven\footnote{Cf. Cat. xv. 22.}; for the trophy shall precede the king: that seeing \emph{Him whom they pierced}\footnote{Zech. xii. 10.}, and knowing by the Cross Him who was dishonoured, the Jews may repent and mourn; (but \emph{they shall mourn tribe by tribe}\footnote{Ib. \emph{v}. 12.}, for they shall repent, when there shall be no more time for repentance;) and that we may glory, exulting in the Cross, worshipping the Lord who was sent, and crucified for us, and worshipping also God His Father who sent Him, with the Holy Ghost: To whom be glory for ever and ever. Amen.

\xchapter{Lecture XIV. On the Words, And Rose Again from the Dead on the Third Day, and Ascended into the Heavens, and Sat on the Right Hand of the Father.}

\textsc{1 Cor. xv. 1–4}

\emph{Now I make known unto you, brethren, the gospel which I preached unto you}….\emph{that He hath been raised on the third day according to the Scriptures, \&c.}

\emph{Rejoice, O Jerusalem, and keep high festival, all ye that love Jesus;} for He is risen. Rejoice, all ye that mourned before\footnote{Is. lxvi. 10.}, when ye heard of the daring and wicked deeds of the Jews: for He who was spitefully entreated of them in this place is risen again. And as the discourse concerning the Cross was a sorrowful one, so let the good tidings of the Resurrection bring joy to the hearers. Let mourning be turned into gladness, and lamentation to joy: and let our mouth be filled with joy and gladness, because of Him, who after His resurrection, said Rejoice\footnote{Matt. xxviii. 9, “All hail.” The usual greeting, \textgreek{Χαίρετε}, “Rejoice.”}. For I know the sorrow of Christ’s friends in these past days; because, as our discourse stopped short at the Death and the Burial, and did not tell the good tidings of the Resurrection, your mind was in suspense, to hear what you were longing for.

Now, therefore, the Dead is risen, He who was \emph{free among the dead}\footnote{Ps. lxxxviii. 5: \emph{Cast off among the dead} (R.V.); \emph{Cast away} (Margin).}, and the deliverer of the dead. He who in dishonour wore patiently the crown of thorns, even He arose, and crowned Himself with the diadem of His victory over death.

2. As then we set forth the testimonies concerning His Cross, so come let us now verify the proofs of His Resurrection also: since the Apostle before us\footnote{\textgreek{ὁ παρών}. i.e. in the text. 1 Cor. xv. 4.} affirms, \emph{He was buried, and has been raised on the third day according to the Scriptures.} As an Apostle, therefore, has sent us back to the testimonies of the Scriptures, it is good that we should get full knowledge of the hope of our salvation; and that we should learn first whether the divine Scriptures tell us the season of His resurrection, whether it comes in summer or in autumn, or after winter; and from what kind of place the Saviour has risen, and what has been announced in the admirable Prophets as the name of the place of the Resurrection, and whether the women, who sought and found Him not, afterwards rejoice at finding Him; in order that when the Gospels are read, the narratives of these holy Scriptures may not be thought fables nor rhapsodies.

3. That the Saviour then was buried, ye have heard distinctly in the preceding discourse, as Isaiah saith, His burial shall be in peace\footnote{Is. lvii. 2: \emph{He entereth into peace} (R.V.).}: for in His burial He made peace between heaven and earth, bringing sinners unto God: and, \emph{that the righteous is taken out of the way of unrighteousness}\footnote{Is. lvii. 1: \emph{that the righteous is taken away from the evil to come} (R.V.).}: and, \emph{His burial shall be in peace}: and, \emph{I will give the wicked for His burial}\footnote{Is. liii. 9: \emph{they made His grave with the wicked} (R.V.).}. There is also the prophecy of Jacob saying in the Scriptures, \emph{He lay down and couched as a lion, and as a lion’s whelp}: who shall rouse Him up\footnote{Gen. xlix. 9.}? And the similar passage in Numbers, He couched, He lay down as a lion, and as a lion’s whelp\footnote{Num. xxiv. 9.}. The Psalm also ye have often heard, which says, \emph{And Thou hast brought me down into the dust of death}\footnote{Ps. xxii. 15.}. Moreover we took note of the spot, when we quoted the words, \emph{Look unto the rock, which ye have hewn}\footnote{\textgreek{ἐπεσημειωσάμεθα}, “noted for ourselves;” Middle Voice. Is. li. 1: quoted in Cat. xiii. 35.}. But now let the testimonies concerning His resurrection itself go with us on our way.

4. First, then, in the 11th Psalm He says, \emph{For the misery of the poor, and the sighing of the needy, now will I arise, saith the Lord}\footnote{Ps. xii. 5.}. But this passage still remains doubtful with some: for He often rises up also in anger\footnote{Ib. vii. 6: “Arise, O Lord, in Thine anger.}, to take vengeance upon His enemies.

Come then to the 15th Psalm, which says distinctly: \emph{Preserve Me, \textsc{O Lord}, for in Thee} \emph{have I put my trust}\footnote{Ps. xvi. 1.}: and after this, \emph{their assemblies of blood will I not join, nor make mention of their names between my lips}\footnote{Ib. xvi. 4: “their drink-offerings of blood will I not offer.” The Psalmist abhors the bloody rites, and the very names of the false gods.}; since they have refused me, and chosen Cæsar as their king\footnote{John xix 15. Cyril applies to the Jews what the Psalmist says concerning those that hasten after another god.}: and also the next words, \emph{I foresaw the} \textsc{Lord} \emph{alway before Me, because He is at My right hand, that I may not be moved}\footnote{Ps. xvi. 8.}: and soon after \emph{Yea and even until night my reins chastened me}\footnote{Ib. 7. Quoting from memory, Cyril transposes these sentences.}. And after this He says most plainly, \emph{For Thou wilt not leave My soul in hell}\footnote{Ib. 10. R.V. \emph{in Sheol}, Sept. \emph{in Hades}.}\emph{; neither wilt Thou suffer Thine Holy One to see corruption}. He said not, neither wilt Thou suffer Thine Holy One to see death, since then He would not have died; but \emph{corruption}, saith He, I see not, and shall not abide in death. \emph{Thou hast made known to Me the ways of life}\footnote{Ib. 11.}. Behold here is plainly preached a life after death. Come also to the 29th Psalm, \emph{I will extol Thee, O} \textsc{Lord}\emph{, for Thou hast lifted Me up, and hast not made My foes to rejoice over Me}\footnote{Ib. xxx. 1.}. What is it that took place? Wert thou rescued from enemies, or wert thou released when about to be smitten? He says himself most plainly, \emph{O \textsc{Lord}, Thou hast brought up My soul from hell}\footnote{Ib. 3. R.V. \emph{from Sheol}, Sept. \emph{from Hades}.}. There he says, Thou wilt not leave, prophetically: and here he speaks of that which is to take place as having taken place, \emph{Thou hast brought up. Thou hast saved Me from them that go down into the pit}\footnote{Ib. 3.}. At what time shall the event occur? \emph{Weeping shall continue for the evening, and joy cometh in the morning}\footnote{Ib. 5.}: for in the evening was the sorrow of the disciplines, and in the morning the joy of the resurrection.

5. But wouldst thou know the place also? Again He saith in Canticles, \emph{I went down into the garden of nuts}\footnote{Cant. iv. 11.}; for it was a garden where He was crucified\footnote{John xix. 41. See Index, \emph{Golgotha}.}. For though it has now been most highly adorned with royal gifts, yet formerly it was a garden, and the signs and the remnants of this remain. \emph{A garden enclosed, a fountain sealed}\footnote{Cant. iv. 12.}, by the Jews who said, \emph{We remember that that deceiver said while He was yet alive, After three days, I will rise: command, therefore, that the sepulchre be made sure; and further on, So they went, and made the sepulchre sure, sealing the stone with the guard}\footnote{Matt. xxvii. 63, 65.}. And aiming well at these, one saith, \emph{and in rest Thou shalt judge them}\footnote{Job vii. 18:.…\emph{try him every moment}. Heb. \texthebrew{ע\texthebrew{גרֶֶ}} “a wink,” as in Job xxi. 13, misinterpreted in both passages by the \textsc{LXX.} as meaning “rest.”}. But who is the fountain that is sealed, or who is interpreted as being a \emph{well-spring of living water}\footnote{Cant. iv. 15.}? It is the Saviour Himself, concerning whom it is written, \emph{For with Thee is the fountain of life}\footnote{Ps. xxxvi. 9.}.

6. But what says Zephaniah in the person of Christ to the disciples? \emph{Prepare thyself, be rising at the dawn: all their gleaning is destroyed}\footnote{Zeph. iii. 7: \emph{they rose early and corrupted all their doings}. The passage is wholly is understood by the Seventy, whom S. Cyril follows.}: the gleaning, that is, of the Jews, with whom there is not a cluster, nay not even a gleaning of salvation left; for their vine is cut down. See how He says to the disciples, \emph{Prepare thyself, rise up at dawn}: at dawn expect the Resurrection.

And farther on in the same context of Scripture He says, \emph{Therefore wait thou for Me, saith the \textsc{Lord,} until the day of My Resurrection at the Testimony}\footnote{Zeph. iii. 8: \emph{until the day that I rise up to the prey}. For \texthebrew{ד\texthebrew{עלְ}}, \emph{to the prey}, the \textsc{LXX.} seem to have read \texthebrew{ד\texthebrew{ע}“\texthebrew{לְְ}}, \emph{to the testimony}. About ten years before these Lectures were delivered, Eusebius (\emph{Life of Constantine}, III. c. xxviii.), speaking of the discovery of the Holy Sepulchre, \textsc{a.d.} 326, calls it “a testimony to the Resurrection of the Saviour clearer than any voice could give.”}. Thou seest that the Prophet foresaw the place also of the Resurrection, which was to be surnamed “the Testimony.” For what is the reason that this spot of Golgotha and of the Resurrection is not called, like the rest of the Churches, a Church, but a Testimony? Why, perhaps, it was because of the Prophet, who had said, \emph{until the day of My Resurrection at the Testimony}.

7. And who then is this, and what is the sign of Him that rises? In the words of the Prophet that follow in the same context, He says plainly, \emph{For then will I turn to the peoples a language}\footnote{Zeph. iii. 9: \emph{a pure language}.}: since, after the Resurrection, when the Holy Ghost was sent forth the gift of tongues was granted, \emph{that they might serve the Lord under one yoke}\footnote{Ib. \emph{to serve him with one consent} (Marg. \emph{shoulder}).}. And what other token is set forth in the same Prophet, that they should serve the \textsc{Lord} \emph{under one yoke?} \emph{From beyond the rivers of Ethiopia they shall bring me offerings}\footnote{Ib. v. 10.}. Thou knowest what is written in the Acts, when the Ethiopian eunuch came from beyond the rivers of Ethiopia\footnote{Acts viii. 27.}. When therefore the Scriptures tell both the time and the peculiarity of the place, when they tell also the signs which followed the Resurrection, have thou henceforward a firm faith in the Resurrection, and let no one stir thee from confessing \emph{Christ risen from the} dead\footnote{2 Tim. ii. 8.}.

8. Now take also another testimony in the 87th Psalm, where Christ speaks in the Prophets, (for He who then spoke came afterwards among us): \emph{O \textsc{Lord}, God of My salvation, I have cried day and night before Thee,} and a little, farther on, \emph{I became as it were a man without help, free among the dead}\footnote{Ps. lxxxviii. 1, 4, 5.}. He said not, I became a man without help; but, \emph{as it were a man without help. } For indeed He was crucified not from weakness, but willingly and His Death was not from involuntary weakness. \emph{I was counted with them that go down into the pit}\footnote{Ib. v. 4.}. And what is the token? \emph{Thou hast put away Mine acquaintance far from Me}\footnote{Ib. v. 8.} (for the disciples have fled). \emph{Wilt Thou shew wonders to the dead}\footnote{Ib. v. 10.}? Then a little while afterwards: \emph{And unto Thee have I cried, O} \textsc{Lord}\emph{; and in the morning shall my prayer come before Thee}\footnote{Ib. v. 13.}. Seest thou how they shew the exact point of the Hour, and of the Passion and of the Resurrection?

9. And whence hath the Saviour risen? He says in the Song of Songs: \emph{Rise up, come, My neighbour}\footnote{Cant. ii. 10: \emph{Rise up, my love, my fair one, and come away.}}: and in what follows, \emph{in a cave of the rock}\footnote{v. 14: \emph{in the clefts of the rock}.}! A cave of the rock He called the cave which was erewhile before the door of the Saviour’s sepulchre, and had been hewn out of the rock itself, as is wont to be done here in front of the sepulchres. For now it is not to be seen, since the outer cave was cut away at that time for the sake of the present adornment. For before the decoration of the sepulchre by the royal munificence, there was a cave in the front of the rock\footnote{See Index, \emph{Sepulchre}.}. But where is the rock that had in it the cave? Does it lie near the middle of the city, or near the walls and the outskirts? And whether is it within the ancient walls, or within the outer walls which were built afterwards? He says then in the Canticles: \emph{in a cave of the rock, close to the outer wall}\footnote{Cant. ii. 14: \emph{in the clefts of the rock, in the secret places of the stairs.} The Revised Version reads, \emph{in the covert of the steep place.}}.

10. At what season does the Saviour rise? Is it the season of summer, or some other? In the same Canticles immediately before the words quoted He says, The winter is past, the rain is past and gone\footnote{Cant. ii. 11. In \textgreek{παρῆλθεν, ἐπορεύθη ἑαυτῷ} the \textsc{LXX.} have imitated the pleonastic use of \texthebrew{ן\texthebrew{ֹל}} after verbs of motion, corresponding to our idiom “Go away with you,” and to the \emph{Dativus Ethicus} in Greek and Latin. See Gesenius \emph{Lexicon} on this use of \texthebrew{ל\texthebrew{ְ}}, and Ewald, \emph{Introductory Grammar}, § 217, l. 2.}; the flowers appear on the earth; the time of the pruning is come\footnote{Cant. ii. 12: \emph{the singing of birds}. The Hebrew word (\texthebrew{ר\texthebrew{ימִזָ֯}} means either “cutting,” as in the \textsc{LXX}. \textgreek{τομῆς}, Symmachus \textgreek{κλαδεύσεως} , and \textsc{R.V.} Marg. “pruning,” or as in \textsc{A.V.} “singing.”}. Is not then the earth full of flowers now, and are they not pruning the vines? Thou seest how he said also that the winter is now past. For when this month Xanthicus\footnote{Xanthicus is the name of the sixth month in the Macedonian Calendar, corresponding nearly to the Jewish Nisan (Josephus, \emph{Antiq}. II. xiv. 6), and to the latter part of Lent and Easter. On the tradition that the Creation took place at this season, see S. Ambrose, \emph{Hexæmeron}, I. c. 4, § 13.} is come, it is already spring. And this is the season, the first month with the Hebrews, in which occurs the festival of the Passover, the typical formerly, but now the true. This is the season of the creation of the world: for then God said, Let the earth bring forth herbage of grass, yielding seed after his kind and after his likeness\footnote{Gen. i. 11: \emph{grass, the herb yielding seed}. \par The \textsc{LXX.} give an irregular construction, \par \textgreek{Βοτανὴν χόρτου σπεῖρον σπέρμα}.}. And now, as thou seest, already every herb is yielding seed. And as at that time God made the sun and moon and gave them courses of equal day (and night), so also a few days since was the season of the equinox.

At that time God said, let us make man after our image and after our likeness\footnote{Gen. i. 26. “The ancient Church very accurately distinguished between \textgreek{εἰκών} (\emph{image}) and \textgreek{ὁμοίωσις} (\emph{likeness}), and the Greek Church does the same in its Confession. The latter phrase expresses man’s destination, which is not to be regarded as carried out at the moment of creation. (Dorner, \emph{System of Christian Doctrine, E. Tr.} II. p. 78). The \emph{image} lies in the permanent capacities of man’s nature (Gen. ix. 6; 1 Cor. xi. 7; Jas. iii. 9), the \emph{likeness} in their realisation in moral conformity with God (\textgreek{ὁμοήθειαν Θεοῦ}, Ignatius, \emph{Magnes} vi). “The image of God is a comprehensive thing.…To this belongs man’s intellective power, his liberty of will, his dominion over the other creatures flowing from the two former. These make up the \textgreek{τὸ οὐσιῶδες}, that part of that divine image which is natural and essential to man, and consequently can never be wholly blotted out, defaced, or extinguished, but still remains even in man fallen. But beside these the Church of God hath ever acknowledged, in the first man, certain additional ornaments, and as it were complements of the divine image, such as immortality, grace, holiness, righteousness, whereby man approached more nearly to the similitude and likeness of God. These were (if I may so speak) the lively colours wherein the grace, the beauty, and lustre of the divine image principally consisted; these colours faded, yea, were defaced and blotted out by man’s transgression. (Bull, \emph{The State of Man before the Fall}, Vol. ii. p. 114, \emph{Ox}.). Cf. Iren. (V. vi. § 1; xvi. § 2); Tertullian (\emph{de Baptismo}, c. 5); Clem. Alex. (\emph{Exhort}. c. 12); Origen (\emph{c. Cels}. IV. 30).}. And the image he received, but the likeness through his disobedience he obscured. At the same season then in which he lost this the restoration also took place. At the same season as the created man through disobedience was cast out of Paradise, he who believed was through obedience brought in. Our Salvation then took place at the same season as the Fall: when the flowers appeared, and the pruning was come.

11. A garden was the place of His Burial, and a vine that which was planted there: and He hath said, \emph{I am the vine}\footnote{John xv. 1. The Benedictine Editor has a different punctuation: “and the vine which was planted there hath said, \emph{And I am the Vine}.”}! He was planted therefore in the earth in order that the curse which came because of Adam might be rooted out. The earth was condemned to thorns and thistles: the true Vine sprang up out of the earth, that the saying might be fulfilled, \emph{Truth sprang up out of the earth, and righteousness} \emph{looked down from heaven}\footnote{Ps. lxxxv. 11.}. And what will He that is buried in the garden say? \emph{I have gathered My myrrh with My spices}: and again, \emph{Myrrh and aloes, with all chief spices}\footnote{Cant. v. 1; iv. 14. Compare Cat. xiii. 32.}. Now these are the symbols of the burying; and in the Gospels it is said, \emph{The women came unto the sepulchre bringing the spices which they had prepared}\footnote{Luke xxiv. 1.}\emph{: Nicodemus also bringing a mixture of myrrh and aloes}\footnote{John xix. 39.}. And farther on it is written, \emph{I did eat My bread with My honey}\footnote{Cant. vi. 1: \emph{my honeycomb with my honey}.}: the bitter before the Passion, and the sweet after the Resurrection. Then after He had risen He entered through closed doors: but they believed not that it was He: for \emph{they supposed that they beheld a spirit}\footnote{Luke xxiv. 37.}. But He said, \emph{Handle Me and see}. Put your fingers into the print of the nails, as Thomas required. \emph{And while they yet believed not for joy, and wondered, He said unto them, Have ye here anything to eat? And they gave Him a piece of a broiled fish and honeycomb}\footnote{Ib. v. 41.}. Seest thou how that is fulfilled, \emph{I did eat My bread with My honey}.

12. But before He entered through the closed doors, the Bridegroom and Suitor\footnote{\textgreek{ὁ θεραπευτής.} In connexion with “Bridegroom,” and “Him whom my soul loveth” the meaning “Suitor” is more appropriate than “Physician.”} of souls was sought by those noble and brave women. They came, those blessed ones, to the sepulchre, and sought Him Who had been raised, and the tears were still dropping from their eyes, when they ought rather to have been dancing with joy for Him that had risen. Mary came seeking Him, according to the Gospel, and found Him not: and presently she heard from the Angels, and afterwards saw the Christ. Are then these things also written? He says in the Song of Songs, \emph{On my bed I sought Him whom my soul loved}. At what season? \emph{By night on my bed I sought Him Whom my soul loved: Mary}, it says, \emph{came while it was yet dark. On my bed I sought Him by night, I sought Him, and I found Him not}\footnote{Cant. iii. 1; John xx. 1.}. And in the Gospels Mary says, \emph{They have taken away my Lord, and I know not where they have laid Him}\footnote{John xx. 13.}. But the Angels being then present cure their want of knowledge; for they said, \emph{Why seek ye the living among the dead}\footnote{Luke xxiv. 5.}? He not only rose, but had also the dead with Him when He rose\footnote{Matt. xxvii. 52.}. But she knew not, and in her person the Song of Songs said to the Angels, \emph{Saw ye Him Whom my soul} loved? It was but a little that I passed from them (that is, from the two Angels), \emph{until I found Him Whom my soul loved. I held Him, and would not let Him go}\footnote{Cant. iii. 3, 4.}.

13. For after the vision of the Angels, Jesus came as His own Herald; and the Gospel says, \emph{And behold Jesus met them, saying, All hail! and they came and took hold of His feet}\footnote{Matt. xxviii. 9.}. They took hold of Him, that it might be fulfilled, \emph{I will hold Him, and will not let Him go}. Though the woman was weak in body, her spirit was manful. \emph{Many waters quench not love, neither do rivers drown it}\footnote{Cant. viii. 7.}; He was dead whom they sought, yet was not the hope of the Resurrection quenched. And the Angel says to them again, \emph{Fear not ye}; I say not to the soldiers, \emph{fear not}, but to you\footnote{Matt. xxviii. 5. The emphatic \textgreek{ὑμεῖς} is rightly interpreted by Cyril as distinguishing the women from the frightened sentinels.}; as for them, let them be afraid, that, taught by experience, they may bear witness and say, \emph{Truly this was the Son of God}\footnote{Matt. xxvii. 54.}; but you ought not to be afraid, \emph{for perfect love casteth out fear}\footnote{1 John iv. 18.}. \emph{Go, tell His disciples that He is risen}\footnote{Matt. xxviii. 7.}; and the rest. And they depart with joy, yet full of fear; is this also written? yes, the second Psalm, which relates the Passion of Christ, says, \emph{Serve the Lord with fear, and rejoice unto Him with trembling}\footnote{Ps. ii. 11.};—\emph{rejoice}, because of the risen Lord; but \emph{with trembling}, because of the earthquake, and the Angel who appeared as lightning.

14. Though, therefore, Chief Priests and Pharisees through Pilate’s means sealed the tomb; yet the women beheld Him who was risen. And Esaias knowing the feebleness of the Chief Priests, and the women’s strength of faith, says, \emph{Ye women, who come from beholding, come hither}\footnote{Isa. xxvii. 11: \emph{The women shall come, and set them on fire}.}\emph{; for the people hath no understanding};—the Chief Priests want understanding, while women are eye-witnesses. And when the soldiers came into the city to them, and told them all that had come to pass, they said to them, \emph{Say ye, His disciples came by night, and stole Him away while we slept}\footnote{Matt. xxviii. 13.}? Well therefore did Esaias foretell this also, as in their persons, \emph{But tell us, and relate to us another deceit}\footnote{Isa. xxx. 10.}. He who rose again, is up, and for a gift of money they persuade the soldiers; but they persuade not the kings of our time. The soldiers then surrendered the truth for silver; but the kings of this day have, in their piety, built this holy Church of the Resurrection of God our Saviour, inlaid with silver and wrought with gold, in which we are assembled\footnote{Cf. Euseb. (\emph{Life of Const}. III. 36.).}; and embellished it with the treasures of silver and gold and precious stones. \emph{And if this come to the governor’s ears}, they say, \emph{we will persuade him}\footnote{Matt. xxxviii. 14.}. Yea, though ye persuade the soldiers, yet ye will not persuade the world; for why, as Peter’s guards were condemned when he escaped out of the prison, were not they also who watched Jesus Christ condemned? Upon the former, sentence was pronounced by Herod, for they were ignorant and had nothing to say for themselves; while the latter, who had seen the truth, and concealed it for money, were protected by the Chief Priests. Nevertheless, though but a few of the Jews were persuaded at the time, the world became obedient. They who hid the truth were themselves hidden; but they who received it were made manifest by the power of the Saviour, who not only rose from the dead, but also raised the dead with Himself. And in the person of these the Prophet Osee says plainly, \emph{After two days will He revive us, and in the third day we shall rise again, and shall live in His sight}\footnote{Hos. vi. 2.}.

15. But since the disobedient Jews will not be persuaded by the Divine Scriptures, but forgetting all that is written gainsay the Resurrection of Jesus, it were good to answer them thus: On what ground, while you say that Eliseus and Elias raised the dead, do you gainsay the Resurrection of our Saviour? Is it that we have no living witnesses now out of that generation to what we say? Well, do you also bring forward witnesses of the history of that time. But that is written;—so is this also written: why then do ye receive the one, and reject the other? They were Hebrews who wrote that history; so were all the Apostles Hebrews: why then do ye disbelieve the Jews\footnote{Instead of \textgreek{τοῖς ᾽Ιουδαίοις} the Jerusalem Editor adopts from Cod. A. \textgreek{τοῖς ἰδίοις}, “Your own countrymen,” a better reading in this place, if it had more support from \textsc{mss.} The Latin in Milles has only “Cur igitur non creditis?”}? Matthew who wrote the Gospel wrote it in the Hebrew tongue\footnote{The statements of Papias, Irenæus, Origen, Eusebius, Epiphanius, and Jerome, concerning a Hebrew Gospel of S. Matthew are ably discussed by Dr. Salmon (\emph{Introduction to N.T.} Lect. X.), who comes to the conclusion that the Canonical Gospel was not translated from Hebrew (Aramaic), but originally written in Greek.}; and Paul the preacher was a Hebrew of the Hebrews; and the twelve Apostles were all of Hebrew race: then fifteen Bishops of Jerusalem were appointed in succession from among the Hebrews\footnote{This statement may have been derived either from Eusebius (\emph{Hist. Eccl.}. IV. c. 5), or from the “written records” (\textgreek{ἐγγράφων}), from which he had learned that “until the siege of the Jews which took place under Adrian (135 \textsc{a.d.}), there were fifteen bishops in succession there, all of whom are said to have been of Hebrew descent.” See the list of names, and the notes on the passage in this Series.}. What then is your reason for allowing your own accounts, and rejecting ours, though these also are written by Hebrews from among yourselves.

16. But it is impossible, some one will say, that the dead should rise; and yet Eliseus twice raised the dead,—when he was alive, and also when dead. Do we then believe, that when Eliseus was dead, a dead man who was cast upon him and touched him, arose and is Christ not risen? But in that case, the dead man who touched Eliseus, arose, yet he who raised him continued nevertheless dead: but in this case both the Dead of whom we speak Himself arose, and many dead were raised without having even touched Him. For \emph{many bodies of the Saints which slept arose, and they came out of the graves after His Resurrection, and went into the Holy City}\footnote{Matt. xxvii. 52, 53.}, (evidently this city, in which we now are\footnote{The Archdeacon of Jerusalem, Photius Alexandrides, observes that “by this parenthetic explanation Cyril perhaps wished to refute the opinion which some favoured that these saints which slept and were raised entered into the heavenly Jerusalem.” See Euseb. \emph{Dem. Evang}. IV. 12.},) \emph{and appeared unto many}. Eliseus then raised a dead man, but he conquered not the world; Elias raised a dead man, but devils are not driven away in the name of Elias. We are not speaking evil of the Prophets, but we are celebrating their Master more highly; for we do not exalt our own wonders by disparaging theirs; for theirs also are ours; but by what happened among them, we win credence for our own.

17. But again they say, “A corpse then lately dead was raised by the living; but shew us that one three days dead can possibly arise, and that a man should be buried, and rise after three days.” If we seek for Scripture testimony in proof of such facts, the Lord Jesus Christ Himself supplies it in the Gospels, saying, \emph{For as Jonas was three days and three nights in the whale’s belly; so shall the Son of man be three days and three nights in the heart of the earth}\footnote{Matt. xii. 40.}. And when we examine the story of Jonas, great is the force\footnote{\textgreek{“ἐνέργεια} [Forte \textgreek{ἐνάργεια}, Edit.].” This conjecture of the Benedictine Editor is recommended by the very appropriate sense “distinctness of the resemblance,” but seems to have no \textsc{ms.} authority.} of the resemblance. Jesus was sent to preach repentance; Jonas also was sent: but whereas the one fled, not knowing what should come to pass; the other came willingly, to give repentance unto salvation. Jonas was asleep in the ship, and snoring amidst the stormy sea; while Jesus also slept, the sea, according to God’s providence\footnote{\textgreek{κατ᾽ οἰκονομίαν}.}, began to rise, to shew in the sequel the might of Him who slept. To the one they said, \emph{Why art thou snoring? Arise, call upon thy God, that God may save us}\footnote{Jonah i. 6.}; but in the other case they say unto the Master, \emph{Lord, save us}\footnote{Matt. viii. 25, 26.}. Then they said, \emph{Call upon thy God}; here they say, \emph{save Thou}. But the one says, \emph{Take me, and cast me into the sea; so shall the sea be calm unto you}\footnote{Jonah i. 12.}; the other, Himself \emph{rebuked the winds and the sea, and there was a great calm}\footnote{Matt. viii. 25, 26.}. The one was cast into a whale’s belly: but the other of His own accord went down thither, where the invisible whale of death is. And He went down of His own accord, that death might cast up those whom he had devoured, according to that which is written, \emph{I will ransom them from the power of the grave; and from the hand of death I will redeem them}\footnote{Hosea xiii. 14.}.

18. At this point of our discourse, let us consider whether is harder, for a man after having been buried to rise again from the earth, or for a man in the belly of a whale, having come into the great heat of a living creature, to escape corruption. For what man knows not, that the heat of the belly is so great, that even bones which have been swallowed moulder away? How then did Jonas, who was three days and three nights in the whale’s belly, escape corruption? And, seeing that the nature of all men is such that we cannot live without breathing, as we do, in air, how did he live without a breath of this air for three days? But the Jews make answer and say, The power of God descended with Jonas when he was tossed about in hell. Does then the Lord grant life to His own servant, by sending His power with him, and can He not grant it to Himself as well? If that is credible, this is credible also; if this is incredible, that also is incredible. For to me both are alike worthy of credence. I believe that Jonas was preserved, for \emph{all things are possible with God}\footnote{Matt. xix. 26.}; I believe that Christ also was raised from the dead; for I have many testimonies of this, both from the Divine Scriptures, and from the operative power even at this day\footnote{Cf. Cat. iv. 13; xiii. 3.} of Him who arose,—who descended into hell alone, but ascended thence with a great company; for He went down to death, \emph{and many bodies of the saints which slept arose}\footnote{Matt. xxvii. 52.} through Him.

19. Death was struck with dismay on beholding a new visitant descend into Hades, not bound by the chains of that place. Wherefore, O porters of Hades, were ye scared at sight of Him? What was the unwonted fear that possessed you? Death fled, and his flight betrayed his cowardice. The holy prophets ran unto Him, and Moses the Lawgiver, and Abraham, and Isaac, and Jacob; David also, and Samuel, and Esaias, and John the Baptist, who bore witness when he asked, \emph{Art Thou He that should come, or look we for another}\footnote{Ib. xi. 3.}? All the Just were ransomed, whom death had swallowed; for it behoved the King whom they had proclaimed, to become the redeemer of His noble heralds. Then each of the Just said, \emph{O death, where is thy victory? O grave, where is thy sting}\footnote{1 Cor. xv. 55. On the opinion that the Patriarchs, Prophets, and Righteous men were redeemed by Christ in Hades, compare Irenæus (\emph{Hær}. I. xxvii. § 3; IV. xxvii. §2), Clem. Alex. (\emph{Stromat}. vi. c. 6), Origen (\emph{In Genes}. Hom. xv. § 5).}? For the Conqueror hath redeemed us.

20. Of this our Saviour the Prophet Jonas formed the type, when he prayed out of the belly of the whale, and said, \emph{I cried in my affliction}, and so on; \emph{out of the belly of hell}\footnote{Jonah ii. 2.}, and yet he was in the whale; but though in the whale, he says that he is in Hades; for he was a type of Christ, who was to descend into Hades. And after a few words, he says, in the person of Christ, prophesying most clearly, \emph{My head went down to the chasms of the mountains}\footnote{Ib. v. 6: (R.V.) \emph{I went down to the bottoms of the mountains: the earth with her bars closed upon me for ever.}}; and yet he was in the belly of the whale. What mountains then encompass thee? I know, he says, that I am a type of Him, who is to be laid in the Sepulchre hewn out of the rock. And though he was in the sea, Jonas says, \emph{I went down to the earth}, since he was a type of Christ, who went down into the heart of the earth. And foreseeing the deeds of the Jews who persuaded the soldiers to lie, and told them, \emph{Say that they stole Him away}, he says, \emph{By regarding lying vanities they forsook their own mercy}\footnote{v. 8.}. For He who had mercy on them came, and was crucified, and rose again, giving His own precious blood both for Jews and Gentiles; yet say they, \emph{Say that they stole Him away}, having regard to \emph{lying vanities}\footnote{By \emph{lying vanities} are meant in the original “vain idols.”}. But concerning His Resurrection, Esaias also says, \emph{He who brought up from the earth the great Shepherd of the sheep}\footnote{Isa. lxiii. 11; (R.V.), \emph{Where is He that brought them up out of the sea with the shepherds} (Marg. \emph{shepherd}) \emph{of His flock?} Cyril’s reading, \textgreek{ἐκ τῆς γῆς} instead of \textgreek{ἐκ τῆς θαλάσσης} is found in the Alexandrine \textsc{ms.} of the Septuagint. Athanasius (\emph{Ad Serapion,} Ep. i. 12) has the same reading and interpretation as Cyril. By “\emph{the shepherds}” are probably meant Moses and Aaron: cf. Ps. lxxvii. 20: \emph{Who leddest Thy people like sheep by the hand of Moses and Aaron}. \par Heb. xiii. 20: \emph{Now the God of peace, that brought again from the dead our Lord Jesus, that great Shepherd of the sheep,} \&c. The word “great” is added by the Author of the Epistle to the Hebrews not by Isaiah.}; he added the word, \emph{great}, lest He should be thought on a level with the shepherds who had gone before Him.

21. Since then we have the prophecies, let faith abide with us. Let them fall who fall through unbelief, since they so will; but thou hast taken thy stand on the rock of the faith in the Resurrection. Let no heretic ever persuade thee to speak evil of the Resurrection. For to this day the Manichees say, that, the resurrection of the Saviour was phantom-wise, and not real, not heeding Paul who says, \emph{Who was made of the seed of David according to the flesh}; and again, \emph{By the resurrection of Jesus Christ our Lord from the dead}\footnote{Rom. i. 3, 4. Cyril in his incomplete quotation of \emph{v}. 4 makes \textgreek{᾽Ιησοῦ Χριστοῦ τοῦ Κ. ἡμ.} depend on \textgreek{ἀναστάσεως}. The right order and construction is given in R.V. \emph{who was declared to be the Son of God}.\emph{…by the resurrection of the dead; even Jesus Christ our Lord.}}. And again he aims at them, and speaks thus, \emph{Say not in} \emph{thine heart, who shall ascend into heaven; or who shall descend into the deep? that is, to bring up Christ from the dead}\footnote{Rom. x. 6, 7.}; and in like manner warning as he has elsewhere written again, \emph{Remember Jesus Christ raised from the dead}\footnote{2 Tim. ii. 8.}; and again, \emph{And if Christ be not risen, then is our preaching vain, and your faith also vain. Yea, and we are found false witnesses of God; because we testified of God that He raised up Christ, whom He raised not up}\footnote{1 Cor. xv. 14, 15.}. But in what follows he says, \emph{But now is Christ risen from the dead, the first fruits of them that are asleep}\footnote{Ib. v. 20.};\emph{—And He was seen of Cephas, then of the twelve}; (for if thou believe not the one witness, thou hast twelve witnesses;) \emph{then He was seen of above five hundred brethren at once}\footnote{Ib. 5, 6.}; (if they disbelieve the twelve, let them admit the five hundred;) \emph{after that He was seen of James}\footnote{Ib. 7. This appearance of Christ to James is not mentioned in the Gospels. Jerome (\emph{Catalog. Script. Eccles.} p. 170 D) mentions a tradition that James had taken an oath that he would eat no bread from the hour in which he had drunk the Cup of the Lord, until he should see Him rising from the dead. Wherefore the Saviour immediately after He had risen appeared to James and commanded him to eat.}, His own brother, and first Bishop of this diocese. Seeing then that such a Bishop originally\footnote{For \textgreek{τοιούτου τοίνυν ἐπισκόπου πρωτοτύπως ἰδόντος} Codd. Roe, Casaub. have \textgreek{τοῦ τοίνυν πρωτοτύπου ἐπισκόπου ἰδόντος}, which gives the better sense—“since therefore the primary Bishop saw, \&c.” On the meaning of \textgreek{παροικία}, and the extent of a primitive Diocese, see Bingham. IX. c. 2.} saw Christ Jesus when risen, do not thou, his disciple, disbelieve him. But thou sayest that His brother James was a partial witness; \emph{afterwards He was seen also of me}\footnote{1 Cor. xv. 8.} Paul, His enemy; and what testimony is doubted, when an enemy proclaims it? “I, \emph{who was before a persecutor}\footnote{1 Tim. i. 13.}, now preach the glad tidings of the Resurrection.”

22. Many witnesses there are of the Saviour’s resurrection.—The night, and the light of the full moon; (for that night was the sixteenth\footnote{If the Crucifixion took place on the 14th of Nisan, the following night would begin the 15th, and the next night the 16th.};) the rock of the sepulchre which received Him; the stone also shall rise up against the face of the Jews, for it saw the Lord; even the stone which was then rolled away\footnote{Cf. Cat. xiii. 39.}, itself bears witness to the Resurrection, lying there to this day. Angels of God who were present testified of the Resurrection of the Only-begotten: Peter and John, and Thomas, and all the rest of the Apostles; some of whom ran to the sepulchre, and saw the burial-clothes, in which He was wrapped before, lying there after the Resurrection; and others handled His hands and His feet, and beheld the prints of the nails; and all enjoyed together that Breath of the Saviour, and were counted worthy to forgive sins in the power of the Holy Ghost. Women too were witnesses, who took hold of His feet, and who beheld the mighty earthquake, and the radiance of the Angel who stood by: the linen clothes also which were wrapped about Him, and which He left when He rose;—the soldiers, and the money given to them; the spot itself also, yet to be seen;—and this house of the holy Church, which out of the loving affection to Christ of the Emperor Constantine of blessed memory, was both built and beautified as thou seest.

23. A witness to the resurrection of Jesus is Tabitha also, who was in His name raised from the dead\footnote{Acts ix. 41.}; for how shall we disbelieve that Christ is risen, when even His Name raised the dead? The sea also bears witness to the resurrection of Jesus, as thou hast heard before\footnote{See § 17, above.}. The drought of fishes also testifies, and the fire of coals there, and the fish laid thereon. Peter also bears witness, who had erst denied Him thrice, and who then thrice confessed Him; and was commanded to feed His spiritual\footnote{\textgreek{νοητά}.} sheep. To this day stands Mount Olivet, still to the eyes of the faithful all but displaying Him Who ascended on a cloud, and the heavenly gate of His ascension. For from heaven He descended to Bethlehem, but to heaven He ascended from the Mount of Olives\footnote{St. Luke (xxiv. 50) describes the Ascension as taking place at Bethany, but the tradition, which Cyril follows, had long since fixed the scene on the summit of the Mount of Olives, a mile nearer to Jerusalem; and here the Empress Helena had built the Church of the Ascension (Eusebius, \emph{Life of Constantine}, III. 43; \emph{Demonstr. Evang.} VI. xviii. 26). There is nothing in Cyril’s language to warrant the Benedictine Editor’s suggestion that he alludes to the legend, according to which the marks of Christ’s feet were indelibly impressed on the spot from which He ascended. In the next generation St. Augustine seems to countenance the miraculous story (\emph{In Joh. Evang}. Tract xlvii.): “There are His footsteps, now adored, where last He stood, and whence He ascended into heaven.” The supposed trace of one foot is still shewn on Mount Olivet; “the other having been removed by the Turks is now to be found in the Chapel of S. Thecla, which is in the Patriarch’s Palace” (Jerusalem Ed.). Compare Stanley, \emph{Sinai and Palestine}, c. xiv.; Dictionary of Bible, \emph{Olives Mount of}.}; at the former place beginning His conflicts among men, but in the latter, crowned after them. Thou hast therefore many witnesses; thou hast this very place of the Resurrection; thou hast also the place of the Ascension towards the east; thou hast also for witnesses the Angels which there bore testimony; and the cloud on which He went up, and the disciples who came down from that place.

24. The course of instruction in the Faith would lead me to speak of the Ascension also; but the grace of God so ordered\footnote{\textgreek{ᾠκονόμησε}. In this word, as also in the phrase below, \textgreek{κατ᾽ οἰκονομίαν τῆς Θείας χάριτος}, Cyril refers to the order of reading the Scriptures as part of a dispensation established by Divine grace.} it, that thou heardest most fully concerning it, as far as our weakness allowed, yesterday, on the Lord’s day; since, by the providence of divine grace, the course of the Lessons\footnote{\textgreek{ἀναγνωσμάτων} a term including the portions of Scripture (\textgreek{περικοπαί}) appointed for the Epistle and Gospel as well as the daily lessons from the Old and New Testaments.} in Church included the account of our Saviour’s going up into the heavens\footnote{The section Luke xxiv. 36–53, which in the Eastern Church is the Gospel for Ascension Day, is also one of the “eleven morning Gospels of the Resurrection (\textgreek{εὐαγγέλια ἀναστασιμὰ ἑωθινά}), which were read in turn, one every Sunday at Matins.” \emph{Dictionary of Chr. Antiq.} “Lectionary.” This Lecture being delivered on Monday, the Section in question had been read on the preceding day.}; and what was then said was spoken principally for the sake of all, and for the assembled body of the faithful, yet especially for thy sake\footnote{\textgreek{μάλιστα μὲν…ἐξαιρέτως δέ}.}. But the question is, didst thou attend to what was said? For thou knowest that the words which come next in the Creed teach thee to believe in Him “Who \textsc{rose again the third day, and ascended into Heaven, and sat down on the right hand of the Father}.” I suppose then certainly that thou rememberest the exposition; yet I will now again cursorily put thee in mind of what was then said. Remember what is distinctly written in the Psalms, \emph{God is gone up with a shout}\footnote{Ps. xlvii. 5.}; remember that the divine powers also said to one another, \emph{Lift up your gates, ye Princes}\footnote{Ps. xxiv. 7: \emph{Lift up, O gates, your heads}. The order of the Hebrew words misled the Greek Translators.}, and the rest; remember also the Psalm which says, \emph{He ascended on high, He led captivity captive}\footnote{Ps. lxviii. 18. On the reading \textgreek{ἀνέβη}, found in a few \textsc{mss.} of the Septuagint, see Tischendorf’s note on Eph. iv. 8.}; remember the Prophet who said, \emph{Who buildeth His ascension unto heaven}\footnote{Amos ix. 6: (R.V.) \emph{It is He that buildeth His chambers in the heaven.} (A.V.) \emph{His stories}. Marg. \emph{ascensions}, or \emph{spheres}. Sept. \textgreek{τὴν ἀνάβασιν αὐτοὐ}.}; and all the other particulars mentioned yesterday because of the gainsaying of the Jews.

25. For when they speak against the ascension of the Saviour, as being impossible, remember the account of the carrying away of Habakkuk: for if Habakkuk was transported by an Angel, being carried by the hair of his head\footnote{Bel and the Dragon, \emph{v}. 33: Compare Ezek. viii. 3.}, much rather was the Lord of both Prophets and Angels, able by His own power to make His ascent into the Heavens on a cloud from the Mount of Olives. Wonders like this thou mayest call to mind, but reserve the preeminence for the Lord, the Worker of wonders; for the others were borne up, but He bears up all things. Remember that Enoch was translated\footnote{Heb. xi. 5.}; but Jesus ascended: remember what was said yesterday concerning Elias, that Elias was taken up in a chariot of fire\footnote{2 Kings ii. 11.}; but that \emph{the chariots of} Christ \emph{are ten thousand-fold even thousands upon thousands}\footnote{Ps. lxviii. 17: \textgreek{χιλιάδες εὐθηνούντων}. The Hebrew means literally “thousands of repetition,” i.e. many thousands: \textgreek{εὐθηνεῖν}, “to abound.”}: and that Elias was taken up, towards the east of Jordan; but that Christ ascended at the east of the brook Cedron: and that Elias went \emph{as into heaven}\footnote{Sept. \textgreek{ὡς εἰς τὸν οὐρανόν}. In 1 Macc. ii. 58 the \textsc{mss.} vary between \textgreek{ἕως} and \textgreek{ὡς}, but the latter (says Fritzsche) “is an alteration made to agree with 2 Kings ii. 11. But there the reference is to the \emph{intended} exaltation of Elijah into heaven, and therefore \textgreek{ὡς} is rightly used (Kühner, \emph{Gramm}. § 604, note; Jelf, § 626, Obs. 1), while here the thing is referred to as an \emph{accomplished historical fact}.” The distinction here drawn by Cyril is therefore hypercritical, as is seen below in § 26, where he writes, \textgreek{᾽Ηλίας μὲν γὰρ ἀνελήφθη εἰς οὐρανόν}.}; but Jesus, into heaven: and that Elias said that a double portion in the Holy Spirit should be given to his holy disciple; but that Christ granted to His own disciples so great enjoyment of the grace of the Holy Ghost, as not only to have It in themselves, but also, by the laying on of their hands, to impart the fellowship of It to them who believed.

26. And when thou hast thus wrestled against the Jews,—when thou hast worsted them by parallel instances, then come further to the pre-eminence of the Saviour’s glory; namely, that they were the servants, but He the Son of God. And thus thou wilt be reminded of His pre-eminence, by the thought that a servant of Christ was caught up to the third heaven. For if Elias attained as far as the first heaven, but Paul as far as the third, the latter, therefore, has obtained a more honourable dignity. Be not ashamed of thine Apostles; they are not inferior to Moses, nor second to the Prophets; but they are noble among the noble, yea, nobler still. For Elias truly was taken up into heaven; but Peter has the keys of the kingdom of heaven, having received the words, \emph{Whatsoever thou shalt loose on earth shall be loosed in heaven}\footnote{Matt. xvi. 19.}. Elias was taken up only to heaven; but Paul both into \emph{heaven}, and into \emph{paradise}\footnote{2 Cor. xii. 2, 4.} (for it behoved the disciples of Jesus to receive more manifold grace), and \emph{heard unspeakable words, which it is not lawful for man to utter}. But Paul came down again from above, not because he was unworthy to abide in the third heaven, but in order that after having enjoyed things above man’s reach, and descended in honour, and having preached Christ, and died for His sake, he might receive also the crown of martyrdom. But I pass over the other parts of this argument, of which I spoke yesterday in the Lord’s-day congregation; for with understanding hearers, a mere reminder is sufficient for instruction.

27. But remember also what I have often said\footnote{See Cat. iv. 7; xi. 17. The clause, \textgreek{καὶ καθίσαντα ἐκ δεξιῶν τοῦ Πατρός}, does not occur in the original form of the Nicene Creed, but is found in the Confession of Faith contained in \emph{Const. Apost}. c. 41, in the four Eusebian Confessions of Antioch (341, 2 \textsc{a.d.}), and in the Macrostichos (344 \textsc{a.d.}). An equivalent clause is found in the brief Confession of Hippolytus (circ. 220 \textsc{a.d.}) \emph{Contra Hæres. Noeti}, c. 1: “\textgreek{καὶ ὄντα ἐν δεξίᾳ τοῦ Πατρός},” and in Tertullian, \emph{De Virgin. Veland}. c. 1: “Regula quidem Fidei una omnino est, sola immobilis et irreformabilis,.…sedentem nunc ad dextram Patris:” \emph{de Præscriptione}, c 13: “Regula est autem fidei.…sedisse ad dexteram Patris:” \emph{adversus Praxean}, c. 2: “sedere ad dexteram Patris.”} concerning the Son’s sitting at the right hand of the Father; because of the next sentence in the Creed, which says, “\textsc{and ascended into Heaven, and sat down at the right hand of the Father}.” Let us not curiously pry into what is properly meant by the throne; for it is incomprehensible: but neither let us endure those who falsely say, that it was after His Cross and Resurrection and Ascension into heaven, that the Son began to sit on the right hand of the Father. For the Son gained not His throne by advancement\footnote{\textgreek{ἐκ προκοπῆς}. Cf. Cat. x. 5, note 8.}; but throughout His being (and His being is by an eternal generation\footnote{\textgreek{ἀφ᾽ οὗπερ ἔστιν, (ἔστι δὲ ἀεὶ γεννηθείς}). In both clauses \textgreek{ἔστιν} is emphatic.}) He also sitteth together with the Father. And this throne the Prophet Esaias having beheld before the incarnate coming of the Saviour, says, \emph{I saw the Lord sitting on a throne, high and lifted up\footnote{Is. vi. 1.}}, and the rest. For the Father \emph{no man hath seen at any time\footnote{John i. 18.}}, and He who then appeared to the Prophet was the Son. The Psalmist also says, \emph{Thy throne is prepared of old; Thou art from everlasting\footnote{Ps. xciii. 2.}}. Though then the testimonies on this point are many, yet because of the lateness of the time, we will content ourselves even with these.

28. But now I must remind you of a few things out of many which are spoken concerning the Son’s sitting at the right hand of the Father. For the hundred and ninth Psalm says plainly, \emph{The \textsc{Lord} said unto my Lord, Sit Thou on My right hand, until I make Thine enemies Thy footstool}\footnote{Ps. cx. 1.}. And the Saviour, confirming this saying in the Gospels, says that David spoke not these things of himself, but from the inspiration of the Holy Ghost, saying, \emph{How then doth David in the Spirit call Him Lord, saying, The Lord said unto my Lord, Sit Thou on My right hand}\footnote{Matt. xxii. 43.}? and the rest. And in the Acts of the Apostles, Peter on the day of Pentecost standing with the Eleven\footnote{Acts ii. 34.}, and discoursing to the Israelites, has in very words cited this testimony from the hundred and ninth Psalm.

29. But I must remind you also of a few other testimonies in like manner concerning the Son’s sitting at the right hand of the Father. For in the Gospel according to Matthew it is written, \emph{Nevertheless, I say unto you, Henceforth ye shall see the Son of Man sitting on the right hand of power}\footnote{Matt. xxvi. 64.}, and the rest: in accordance with which the Apostle Peter also writes, \emph{By the Resurrection of Jesus Christ, who is on the right hand of God, having gone into heaven}\footnote{1 Pet. iii. 22.}. And the Apostle Paul, writing to the Romans, says, \emph{It is Christ that died, yea rather, that is risen again, who is even at the right hand of God}\footnote{Rom. viii. 34.}. And charging the Ephesians, he thus speaks, \emph{According to the working of His mighty power, which He wrought in Christ when He raised Him from the dead, and set Him at His own right hand}\footnote{Eph. i. 19, 20.}; and the rest. And the Colossians he taught thus, \emph{If ye then be risen with Christ, seek the things above, where Christ is seated at the right hand of God}\footnote{Col. iii. 1.}. And in the Epistle to the Hebrews he says, \emph{When He had made purification of our sins, He sat down on the right hand of the Majesty on high}\footnote{Heb. i. 3.}. And again, \emph{But unto which of the Angels hath He said at any time, Sit thou at My right hand, until I make thine enemies thy footstool}\footnote{Ib. v. 13.}? And again, \emph{But He, when He had offered one sacrifice for all men, for ever sat down on the right hand of God; from henceforth expecting till His enemies be made His footstool}\footnote{Ib. x. 12.}. And again, \emph{Looking unto Jesus, the author and perfecter of our faith; Who for the joy that was set before Him endured the Cross, despising shame, and is set down on the right hand of the throne of God}\footnote{Ib. xii. 2. On Cyril’s omission of Mark xvi. 19. see Westcott and Hort.}.

30. And though there are many other texts concerning the session of the Only-begotten on the right hand of God, yet these may suffice us at present; with a repetition of my remark, that it was not after His coming in the flesh\footnote{\textgreek{τὴν ἔνσαρκον παρουσίαν}. Cf. § 27.} that He obtained the dignity of this seat; no, for even before all ages, the Only-begotten Son of God, our Lord Jesus Christ, ever possesses the throne on the right hand of the Father. Now may He Himself, the God of all, who is Father of the Christ, and our Lord Jesus Christ, who came down, and ascended, and sitteth together with the Father, watch over your souls; keep unshaken and unchanged your hope in Him who rose again; raise you together with Him from your dead sins unto His heavenly gift; count you worthy to be \emph{caught up in the clouds, to meet the Lord in the air}\footnote{1 Thess. iv. 17.}, in His fitting time; and, until that time arrive of His glorious second advent, write all your names in the Book of the living, and having written them, never blot them out (for the names of many, who fall away, are blotted out); and may He grant to all of you to believe on Him who rose again, and to look for Him who is gone up, and is to come again, (to come, but not from the earth; for be on your guard, O man, because of the deceivers who are to come;) Who sitteth on high, and is here present together with us, \emph{beholding the} \emph{order of each, and the steadfastness of his faith}\footnote{Col. ii. 5.}. For think not that because He is now absent in the flesh, He is therefore absent also in the Spirit. He is here present in the midst of us, listening to what is said of Him, and beholding thine inward thoughts, and \emph{trying the reins and hearts}\footnote{Ps. vii. 9.};—who also is now ready to present those who are coming to baptism, and all of you, in the Holy Ghost to the Father, and to say, \emph{Behold, I and the children whom God hath given Me}\footnote{Isa. viii. 18; Heb. ii. 13.}:—To whom be glory for ever. Amen.

\xchapter{Lecture XV. On the Clause, And Shall Come in Glory to Judge the Quick and the Dead; Of Whose Kingdom There Shall Be No End.}

\textsc{Daniel vii. 9–14}

\emph{I beheld till thrones were placed, and one that was ancient of days did sit,} and then, \emph{I saw in a vision of the night, and behold one like unto the Son of Man came with the clouds of heaven, \&c.}

1. \textsc{We} preach not one advent only of Christ, but a second also, far more glorious than the former. For the former gave a view of His patience; but the latter brings with it the crown of a divine kingdom. For all things, for the most part, are twofold in our Lord Jesus Christ: a twofold generation; one, of God, before the ages; and one, of a Virgin, at the close of the ages: His descents twofold; one, the unobserved, \emph{like rain on a fleece}\footnote{Ps. lxxii. 6. See xii. 9; and § 10, below.}; and a second His open coming, which is to be. In His former advent, He was wrapped in swaddling clothes in the manger; in His second, He \emph{covereth Himself with light as with a garment}\footnote{Ps. civ. 2.}. In His first coming, \emph{He endured the Cross, despising shame}\footnote{Heb. xii. 2.}; in His second, He comes attended by a host of Angels, receiving glory\footnote{Cyril’s contrast of the two Advents seems to be partly borrowed from Justin M. (\emph{Apol}. i. 52; \emph{Tryph}. 110). See also Tertullian (\emph{Adv. Judæos}, c. 14); Hippolytus (\emph{De Antichristo}, 44).}. We rest not then upon His first advent only, but look also for His second. And as at His first coming we said, \emph{Blessed is He that cometh in the Name of the Lord}\footnote{Matt. xxi. 9; xxiii. 39.}, so will we repeat the same at His second coming; that when with Angels we meet our Master, we may worship Him and say, \emph{Blessed is He that cometh in the Name of the Lord}. The Saviour comes, not to be judged again, but to judge them who judged Him; He who before held His peace when judged\footnote{Ib. xxvi. 63.}, shall remind the transgressors who did those daring deeds at the Cross, and shall say, \emph{These things hast thou done, and I kept silence}\footnote{Ps. l. 21.}. Then, He came because of a divine dispensation, teaching men with persuasion; but this time they will of necessity have Him for their King, even though they wish it not.

2. And concerning these two comings, Malachi the Prophet says, \emph{And the Lord whom ye seek shall suddenly come to His temple}\footnote{Mal. iii. 1–3.}; behold one coming. And again of the second coming he says, \emph{And the Messenger of the covenant whom ye delight in. Behold, He cometh, saith}\footnote{The Benedictine Editor by omitting \textgreek{λέγει}, obtains the sense, \emph{He cometh, even the Lord Almighty}. But \textgreek{λέγει} is supported by the \textsc{mss.} of Cyril, as well as by the Septuagint and Hebrew.} \emph{the Lord Almighty. But who shall abide the day of His coming? or who shall stand when He appeareth? Because He cometh in like a refiner’s fire, and like fullers’ herb; and He shall sit as a refiner and purifier}. And immediately after the Saviour Himself says, \emph{And I will draw near to you in judgment; and I will be a swift witness against the sorcerers, and against the adulteresses, and against those who swear falsely in My Name}\footnote{Mal. iii. 5.}, and the rest. For this cause Paul warning us beforehand says, \emph{If any man buildeth on the foundation gold, and silver, and precious stones, wood, hay, stubble; every man’s work shall be made manifest; for the day shall declare it, because it shall be revealed in fire}\footnote{1 Cor. iii. 12.}. Paul also knew these two comings, when writing to Titus and saying, \emph{The grace of God hath appeared which bringeth salvation unto all men, instructing us that, denying ungodliness and worldly lusts, we should live soberly, and godly, and righteously in this present world; looking for the blessed hope, and appearing of the glory of the great God and our Saviour Jesus Christ}\footnote{Titus ii. 11. The Benedictine Editor adopts \textgreek{τοῦ Σωτῆρος} instead of \textgreek{ἡ σωτήριος} against the authority of the best \textsc{mss.} of Cyril.}. Thou seest how he spoke of a first, for which he gives thanks; and of a second, to which we look forward. Therefore the words also of the Faith which we are announcing were just now delivered thus\footnote{\textgreek{νῦν παρεδόθη}. Cyril means that at the beginning of this present Lecture he had delivered to the Catechumens those articles of the Creed which he was going to explain. Compare Cat. xviii. 21, where we see that Cyril first announces (\textgreek{ἐπαγγέλλω}) the words which the learners repeat after him (\textgreek{ἀπαγγέλλω}). \par The clause, \textsc{Whose Kingdom shall have no end}, was not contained in the original form of the Creed of Nicæa, \textsc{a.d.} 325, but its substance is found in many earlier writings. Compare Justin M. (\emph{Tryph}. § 46: \textgreek{καὶ αὐτοῦ ἐστιν ἡ αἰώνιος βασιλεία}); \emph{Const. Apost.} vii. 41; the Eusebian Confessions 1st and 4th Antioch, and the Macrostich, \textsc{a.d.} 341, 342, 344. Bp. Bull asserts that the Creed of Jerusalem, containing this clause, was no other than the ancient Eastern Creed, first directed against the Gnostics of the Sub-Apostolic age (\emph{Judicium Eccl. Cathol}. vi. 16).}; that we believe in Him, who also \textsc{ascended into the heavens, and sat} \textsc{down on the right hand of the Father, and shall come in glory to judge quick and dead; whose kingdom shall have no end}.

3. Our Lord Jesus Christ, then, comes from heaven; and He comes with glory at the end of this world, in the last day. For of this world there is to be an end, and this created world is to be re-made anew\footnote{The Benedictine Editor suggests that Cyril “is refuting those who said that the Universe was to perish utterly, an opinion which seems to be somehow imputed to Origen by Methodius, or Proclus, in Epiphanius (\emph{Hæres}. lxiv. 31, 32).” On Origen’s much controverted opinions concerning the beginning and end of the world, see Huet. \emph{Origeniana}, II. 4–6: and Bp. Westcott, \emph{Dictionary of Christian Biography,} “Origen,” pp. 137, 138.}. For since corruption, \emph{and theft, and adultery}, and every sort of sins \emph{have been poured forth over the earth, and blood has been mingled with blood}\footnote{Hos. iv. 2.} in the world, therefore, that this wondrous dwelling-place may not remain filled with iniquity, this world passeth away, that the fairer world may be made manifest. And wouldest thou receive the proof of this out of the words of Scripture? Listen to Esaias, saying, \emph{And the heaven shall be rolled together as a scroll; and all the stars shall fall, as leaves from a vine, and as leaves fall from a fig-tree\footnote{Is. xxxiv. 4.}}. The Gospel also says, \emph{The sun shall be darkened, and the moon shall not give her light, and the stars shall fall from heaven\footnote{Matt. xxiv. 29.}}. Let us not sorrow, as if we alone died; the stars also shall die; but perhaps rise again. And the Lord rolleth up the heavens, not that He may destroy them, but that He may raise them up again more beautiful. Hear David the Prophet saying, \emph{Thou, Lord, in the beginning didst lay the foundations of the earth, and the heavens are the work of Thy hands; they shall perish, but Thou remainest}\footnote{Ps. cii. 25, 26; Heb. i. 10–12.}. But some one will say, Behold, he says plainly that \emph{they shall perish}. Hear in what sense he says, \emph{they shall perish}; it is plain from what follows; \emph{And they all shall wax old as doth a garment; and as a vesture shalt Thou fold them up, and they shall be changed}. For as a man is said to “perish,” according to that which is written, \emph{Behold, how the righteous perisheth, and no man layeth it to heart}\footnote{Is. lvii. 1.}, and this, though the resurrection is looked for; so we look for a resurrection, as it were, of the heavens also. \emph{The sun shall be turned into darkness, and the moon into blood}\footnote{Joel ii. 31.}. Here let converts from the Manichees gain instruction, and no longer make those lights their gods; nor impiously think, that this sun which shall be darkened is Christ\footnote{Cat. vi. 13; xi. 21.}. And again hear the Lord saying, \emph{Heaven and earth shall pass away, but My words shall not pass away}\footnote{Matt. xxiv. 35.}; for the creatures are not as precious as the Master’s words.

4. The things then which are seen shall pass away, and there shall come the things which are looked for, things fairer than the present; but as to the time let no one be curious. For \emph{it is not for you}, He says, \emph{to know times or seasons, which the Father hath put in His own power}\footnote{Acts i. 7.}. And venture not thou to declare when these things shall be, nor on the other hand supinely slumber. For he saith, \emph{Watch, for in such an hour as ye expect not the Son of Man cometh}\footnote{Matt. xxiv. 42, 44; Ib. v. 3.}. But since it was needful for us to know the signs of the end, and since we are looking for Christ, therefore, that we may not die deceived and be led astray by that false Antichrist, the Apostles, moved by the divine will, address themselves by a providential arrangement to the True Teacher, and say, \emph{Tell us, when shall these things be, and what shall be the sign of Thy coming, and of the end of the world}\footnote{Ib. vv. 3 and 4.}? We look for Thee to come again, but \emph{Satan transforms himself into an Angel of light}; put us therefore on our guard, that we may not worship another instead of Thee. And He, opening His divine and blessed mouth, says, \emph{Take heed that no man mislead you}. Do you also, my hearers, as seeing Him now with the eyes of your mind, hear Him saying the same things to you; \emph{Take heed that no man mislead you}. And this word exhorts you all to give heed to what is spoken; for it is not a history of things gone by, but a prophecy of things future, and which will surely come. Not that we prophesy, for we are unworthy; but that the things which are written will be set before you, and the signs declared. Observe thou, which of them have already come to pass, and which yet remain; and make thyself safe.

5. \emph{Take heed that no man mislead you: for many shall come in My name, saying, I am Christ, and shall mislead many}. This has happened in part: for already Simon Magus has said this, and Menander\footnote{Cat. vi. 14, 16.}, and some others of the godless leaders of heresy; and others will say it in our days, or after us.

6. A second sign. \emph{And ye shall hear of wars and rumours of wars}\footnote{Matt. xxiv. 6. The war with Sapor II., King of Persia, which broke out immediately on the death of Constantine, and continued throughout the reign of Constantius, was raging fiercely at the date of these Lectures, the great battle of Singara being fought in the year 348 \textsc{a.d.}}. Is there then at this time war between Persians and Romans for Mesopotamia, or no? Does nation rise up against nation and kingdom against kingdom, or no? \emph{And there shall be famines and pesti}\emph{lences and earthquakes in divers places}. These things have already come to pass; and again, \emph{And fearful sights from heaven, and mighty storms}\footnote{Luke xxi. 11. Jerome in the \emph{Chronicon} mentions a great earthquake in 346 \textsc{a.d.}, by which Dyrrachium was destroyed, and Rome and other cities of Italy greatly injured (Ben. Ed.). \par Cyril substitutes \textgreek{χειμῶνες} for \textgreek{σημεῖα}, the better reading in Luke xxi. 11.}. \emph{Watch therefore}, He says; \emph{for ye know not at what hour your Lord doth come}\footnote{Matt. xxiv. 42.}.

7. But we seek our own sign of His coming; we Churchmen seek a sign proper to the Church\footnote{\textgreek{ἑκκλησιαστικός}, when applied to persons, means either, as here, an orthodox member of the Church in contrast to a heretic, pagan, or Jew (Origen, \emph{in Job} xx. 6), or more particularly a Cleric as opposed to a layman (Cat. xvii. 10).}. And the Saviour says, \emph{And then shall many be offended, and shall betray one another, and shall hate one another}\footnote{Matt. xxiv. 10.}. If thou hear that bishops advance against bishops, and clergy against clergy, and laity against laity even unto blood, be not troubled\footnote{“S. Cyril here describes the state of the Church, when orthodoxy was for a while trodden under foot, its maintainers persecuted, and the varieties of Arianism, which took its place, were quarreling for the ascendancy. Gibbon quotes two passages, one from a pagan historian of the day, another from a Father of the Church, which fully bear out S. Cyril’s words. What made the state of things still more deplorable, was the defection of some of the orthodox party, as Marcellus, into opposite errors: while the subsequent secessions of Apollinaris and Lucifer show what lurking disorders there were within it at the time when S. Cyril wrote. (Vid. \emph{infr}. 9.) The passages referred to are as follows: ‘The Christian Religion,’ says Ammianus, ‘in itself plain and simple, he (Constantius) confounded by the dotage of superstition. Instead of reconciling the parties by the weight of his authority, he cherished and propagated, by vain disputes, the differences which his vain curiosity had excited. The highways were covered with troops of Bishops, galloping from every side to the assemblies, which they called synods; and while they laboured to reduce the whole sect to their own particular opinions, the public establishment of the posts was almost ruined by their hasty and repeated journeys.’ \emph{Hist}. xxi. 16. S. Hilary of Poictiers thus speaks of Asia Minor, the chief seat of the Arian troubles: ‘It is a thing equally deplorable and dangerous, that there are as many creeds as opinions among men, as many doctrines as inclinations, and as many sources of blasphemy as there are faults among us; because we make creeds arbitrarily, and explain them as arbitrarily. The Homoousion is rejected and received and explained away by successive synods. The partial or total resemblance of the Father and of the Son is a subject of dispute for these unhappy divines. Every year, nay, every moon, we make new creeds to describe invisible mysteries. We repent of what we have done, we defend those who repent, we anathematize those whom we defended. We condemn either the doctrine of others in ourselves, or our own in that of others; and reciprocally tearing one another to pieces, we have been the cause of each other’s ruin,’ \emph{ad Constant}. ii. 4, 5. Gibbon’s translations are used, which, though diffuse, are faithful in their matter. What a contrast do these descriptions present to Athanasius’ uniform declaration, that the whole question was really settled at Nicæa, and no other synod or debate was necessary!”—(\textsc{R.W.C.}). Compare, for example, the account of the seditions in Antioch and in Constantinople, in Socrates, \emph{Eccles. Hist}. i. 24; i., 12–14, and Athanas. \emph{Hist. Arianorum}, passim.}; for it has been written before. Heed not the things now happening, but the things which are written; and even though I who teach thee perish, thou shalt not also perish with me; nay, even a hearer may become better than his teacher, and he who came last may be first, since even those about the eleventh hour the Master receives. If among Apostles there was found treason, dost thou wonder that hatred of brethren is found among bishops? But the sign concerns not only rulers, but the people also; for He says, \emph{And because iniquity shall abound, the love of the many shall wax cold}\footnote{Matt. xxiv. 12.}. Will any then among those present boast that he entertains friendship unfeigned towards his neighbour? Do not the lips often kiss, and the countenance smile, and the eyes brighten forsooth, while the heart is planning guile, and the man is plotting mischief with words of peace?

8. Thou hast this sign also: \emph{And this Gospel of the kingdom shall be preached in all the world for a witness unto all nations, and then shall the end come}\footnote{Matt. xxiv. 14.}. And as we see, nearly the whole world is now filled with the doctrine of Christ.

9. And what comes to pass after this? He says next, \emph{When therefore ye see the abomination of desolation, which was spoken of by Daniel the Prophet, standing in the Holy Place, let him that readeth understand}\footnote{Ib. v. 15.}. And again, \emph{Then if any man shall say unto you, Lo, here is the Christ, or, Lo, there; believe it not}\footnote{Ib. v. 23.}. Hatred of the brethren makes room next for Antichrist; for the devil prepares beforehand the divisions among the people, that he who is to come may be acceptable to them. But God forbid that any of Christ’s servants here, or elsewhere, should run over to the enemy! Writing concerning this matter, the Apostle Paul gave a manifest sign, saying, \emph{For that day shall not come, except there came first the falling away, and the man of sin be revealed, the son of perdition, who opposeth and exalteth himself against all that is called God, or that is worshipped; so that he sitteth in the temple of God, shewing himself that he is God. Remember ye not that when I was yet with you, I told you these things? And now ye know that which restraineth, to the end that he may be revealed in his own season. For the mystery of iniquity doth already work, only there is one that restraineth now, until he be taken out of the way. And then shall the lawless one be revealed, whom the Lord Jesus shall slay with the breath of His mouth, and shall destroy with the brightness of His coming. Even him, whose coming is after the working of Satan, with all power and signs and lying wonders, and with all deceit of unrighteousness for them that are perishing}\footnote{2 Thess. ii. 3–10.}. Thus wrote Paul, and now is the falling away. For men have fallen away from the right faith\footnote{The prediction was supposed by earlier Fathers to refer to a personal Antichrist whom they expected to come speedily. See Justin M. (\emph{Tryph}. § 110: \textgreek{ὁ τῆς ἀποστασίας ἄνθρωπος}; \emph{ib}. § 32: “He who is to speak blasphemous and daring things against the Most High is already at the doors.” Iren. \emph{Hær}. V. 25. Cyril in this passage regards the heresies of his time as the apostasy in general, but looks also for a personal Antichrist: (§§ 11, 12).}; and some preach the identity of the Son with the Father\footnote{\textgreek{υἱοπατορία} . On this contemptuous name for Sabellianism, see Cat. iv. 8; xi. 16. The Third (Eusebian) Confession, or Third of Antioch, \textsc{a.d.} 341, anathematizes any who hold the doctrines of Marcellus of Ancyra or Sabellius, or Paul of Samosata (Athan. \emph{de Synodis}, § 24 note 10, p. 462, in this Series, and Mr. Robertson’s \emph{Prolegomena}, p. xliv.). In the \emph{Ecthesis}, or \emph{Statement of Faith}, § 2, Athanasius writes: “Neither do we hold a Son-Father, as do the Sabellians, calling Him of one but (\emph{a sole and?}) not the same essence, and thus destroying the existence of the Son.” As to Marcellus, see Athanasius, \emph{Hist. Arian}. § 6 (p. 271), and the letter of Julius in the \emph{Apologia c. Arian.} § 32 (p. 116): also notes 3, 4 on § 27 below.}, and others dare to say that Christ was brought into being out of nothing\footnote{See Athanasius, \emph{De Synod}. § 15: “Arius and those with him thought and professed thus: ‘God made the Son out of nothing, and called Him His Son:’” and \emph{Expos. Fidei}, § 2: “We do not regard as a creature, or thing made, or as made out of nothing, God the Creator of all, the Son of God, the true Being from the true Being, the Alone from the Alone, inasmuch as the like glory and power was eternally and conjointly begotten of the Father.” The 4th (Eusebian) Confession, or 4th of Antioch, \textsc{a.d.} 342, ends thus: “Those who say that the Son was from nothing,.…the Catholic Church regards as aliens.”}. And formerly the heretics were manifest; but now the Church is filled with heretics in disguise\footnote{Athan. \emph{Adversus Arianos, Or.} i. 1: “One heresy and that the last which has now risen as forerunner of Antichrist, the Arian as it is called, considering that other heresies, her elder sisters, have been openly proscribed, in her craft and cunning affects to array herself in Scripture language, like her father the devil, and is forcing her way back into the Church’s paradise, \&c.” The supposed date of this Oration is 8 or 10 years later than that of Cyril’s Lectures.}. For men have fallen away from the truth, and \emph{have itching ears}\footnote{2 Tim. iv. 3.}. Is it a plausible discourse? all listen to it gladly. Is it a word of correction? all turn away from it. Most have departed from right words, and rather choose the evil, than desire the good\footnote{A reading supported by the best \textsc{mss.} and approved by the Benedictine Editor gives a different sense, “and rather choose to seem than resolve to be,” inverting the proverb “esse quam videri.”}. This therefore is \emph{the falling away}, and the enemy is soon to be looked for: and meanwhile he has in part begun to send forth his own forerunners\footnote{In the passage quoted above in note 5 the Arian heresy is called a “forerunner” (\textgreek{πρόδρομος}) of Antichrist.}, that he may then come prepared upon the prey. Look therefore to thyself, O man, and make safe thy soul. The Church now charges thee before the Living God; she declares to thee the things concerning Antichrist before they arrive. Whether they will happen in thy time we know not, or whether they will happen after thee we know not; but it is well that, knowing these things, thou shouldest make thyself secure beforehand.

10. The true Christ, the Only-begotten Son of God, comes no more from the earth. If any come making false shows\footnote{\textgreek{φαντασιοκοπῶν}, a rare word, rendered “frantic” in Ecclus. iv. 30: its more precise meaning seems to be “making a false show,” which is here applied to a false Christ, and again in § 14 to the father of lies who makes a vain show of false miracles.} in the wilderness, go not forth; if they say, \emph{Lo, here is the Christ, Lo, there, believe it not}\footnote{Matt. xxiv. 23.}. Look no longer downwards and to the earth; for the Lord descends from heaven; not alone as before, but with many, escorted by tens of thousands of Angels; nor secretly as the dew on the fleece\footnote{Ps. lxxii. 6. Cf. § 1, note 1.}; but shining forth openly as the lightning. For He hath said Himself, \emph{As the lightning cometh out of the east, and shineth even unto the west, so shall also the coming of the Son of Man be}\footnote{Matt. xxiv. 27.}; and again, \emph{And they shall see the Son of Man coming upon the clouds with power and great glory, and He shall send forth His Angels with a great trumpet}\footnote{Matt. xxiv. v. 30.}; and the rest.

11. But as, when formerly He was to take man’s nature, and God was expected to be born of a Virgin, the devil created prejudice against this, by craftily preparing among idol-worshippers\footnote{\textgreek{ἐν εἰδωλολατρείᾳ} may mean either “in idol-worship,” or “among idolaters,” the abstract being used for the concrete, as in Rom. iii. 30: \textgreek{δικαιώσει περιτομήν}.} fables of false gods, begetting and begotten of women, that, the falsehood having come first, the truth, as he supposed, might be disbelieved; so now, since the true Christ is to come a second time, the adversary, taking occasion by\footnote{\textgreek{ἐφόδιον}, “provision for a journey,” is here equivalent in meaning to \textgreek{ἀφορμή}, “a starting point,” or “favourable occasion.”} the expectation of the simple, and especially of them of the circumcision, brings in a certain man who is a magician\footnote{Antichrist is described by Hippolytus (\emph{De Christo et Antichristo}, § 57, as “a son of the devil, and a vessel of Satan,” who will rule and govern “after the manner of the law of Augustus, by whom the Roman empire was established, sanctioning everything thereby.” Cf. Iren. \emph{Hær}. V. 30, § 3; Dictionary of Christian Biography, \emph{Antichrist}: “The sharp precision with which St. Paul had pointed to ‘\emph{the} man of sin,’ ‘\emph{the} lawless one,’ ‘\emph{the} adversary,’ ‘\emph{the} son of perdition,’ led men to dwell on that thought rather than on the many \textgreek{ψευδόχριστοι} of whom Christ Himself had spoken.”}, and most expert in sorceries and enchantments of beguiling craftiness; who shall seize for himself the power of the Roman empire, and shall falsely style himself Christ; by this name of Christ deceiving the Jews, who are looking for the Anointed\footnote{\textgreek{τὸν ᾽Ηλειμυένον}, Aquila’s rendering of \texthebrew{ח\texthebrew{ישמ}}, adopted by the Jews in preference to \textgreek{τὸν Χριστόν}, from hatred of the name Christ or Christian. Hippolytus, \emph{ubi supra}, § 6: “The Saviour came into the world in the Circumcision, and he (Antichrist) will come in the same manner:” ib. § 14: “As Christ springs from the tribe of Judah, so Antichrist is to spring from the tribe of Dan.” This expectation was grounded by Hippolytus on Gen. xlix. 17.}, and seducing those of the Gentiles by his magical illusions.

12. But this aforesaid Antichrist is to come when the times of the Roman empire shall have been fulfilled, and the end of the world is now drawing near\footnote{The fourth kingdom in the prophecy of Daniel (vii. 7, 23) was generally understood by early Christian writers to be the Roman Empire; and its dissolution was to be speedily followed by the end of the world. See § 13 below; Irenæus, V. 26; and Hippolytus, \emph{ubi supra}, §§ 19, 28.}. There shall rise up together ten kings of the Romans, reigning in different parts perhaps, but all about the same time; and after these an eleventh, the Antichrist, who by his magical craft shall seize upon the Roman power; and of the kings who reigned before him, \emph{three he shall humble}\footnote{Dan. vii. 24: \emph{and he shall put down three kings} (R.V.).}, and the remaining seven he shall keep in subjection to himself. At first indeed he will put on a show of mildness (as though he were a learned and discreet person), and of soberness and benevolence\footnote{The Jerusalem Editor quotes as from Hippolytus a similar description of Antichrist (§ 23): ”In his first steps he will be gentle, loveable, quiet, pious, pacific, hating injustice, detesting gifts, not allowing idolatry, \&c.” But the treatise is a forgery of unknown date, apparently much later than Cyril.}: and by the lying signs and wonders of his magical deceit\footnote{Iren. V. 28, § 2: “Since the demons and apostate spirits are at his service, he through their means performs wonders, by which he leads the inhabitants of the earth astray.”} having beguiled the Jews, as though he were the expected Christ, he shall afterwards be characterized by all kinds of crimes of inhumanity and lawlessness, so as to outdo all unrighteous and ungodly men who have gone before him; displaying against all men, but especially against us Christians, a spirit murderous and most cruel, merciless and crafty\footnote{Iren. V. 25, § 4: “He shall remove his kingdom into that city (Jerusalem), and shall sit in the Temple of God, leading astray those who worship him as if he were Christ.” \par According to the genuine treatise of Hippolytus Antichrist was to restore the kingdom of the Jews (\emph{De Antichristo}, § 25), to collect the Jews out of every country of the Dispersion, making them his own, as though they were his own children, and promising to restore their country, and establish again their kingdom and nation, in order that he may be worshipped by them as God (§ 54), and he will lead them on to persecute the saints, i.e. the Christians (§ 56). Compare the elaborate description of Antichrist and his cruelty in Lactantius, \emph{Div. Inst}. vii. 17; \emph{Epit}. § 71.}. And after perpetrating such things for three years and six months only, he shall be destroyed by the glorious second advent from heaven of the only-begotten Son of God, our Lord and Saviour Jesus, the true Christ, who shall slay Antichrist \emph{with the breath of His mouth}\footnote{2 Thess ii. 8. Cf. Iren. V. 25, § 3: Hippol. § 64.}, and shall deliver him over to the fire of hell.

13. Now these things we teach, not of our own invention, but having learned them out of the divine Scriptures used in the Church\footnote{\textgreek{ἐκκλησιαζομένων}. Cf. Cat. iv. 35, 36, where Cyril distinguishes the Scriptures \textgreek{ἃς καὶ ἐν ᾽Εκκλησίᾳ μετὰ παρρησίας ἀναγινώσκομεν} from \textgreek{ὅσα ἐν ᾽Εκκλησίαις μὴ ἀναγινώσκεται}.}, and chiefly from the prophecy of Daniel just now read; as Gabriel also the Archangel interpreted it, speaking thus: \emph{The fourth beast shall be a fourth kingdom upon earth, which shall surpass all kingdoms}\footnote{Dan. vii. 23: (R.V.) \emph{shall be diverse from all the kingdoms}.}. And that this kingdom is that of the Romans, has been the tradition of the Church’s interpreters. For as the first kingdom which became renowned was that of the Assyrians, and the second, that of the Medes and Persians together, and after these, that of the Macedonians was the third, so the fourth kingdom now is that of the Romans\footnote{Irenæus (V. 26) identifies the fourth kingdom with “the empire which now rules.” Hippolytus, \emph{de Antichristo}, § 25: “\emph{A fourth beast dreadful and terrible: it had iron teeth and claws of brass}. And who are these but the Romans?”}. Then Gabriel goes on to interpret, saying, \emph{His ten horns are ten kings that shall arise; and another king shall rise up after them, who shall surpass in wickedness all who were before him}\footnote{Dan. vii. 24.}; (he says, not only the ten, but also all who have been before him;) \emph{and he shall subdue three kings}; manifestly out of the ten former kings: but it is plain that by subduing three of these ten, he will become the eighth king; \emph{and he shall speak words against the Most High}\footnote{Dan. v. 25. Dean Church compares Rev. xvii. 11: \emph{And the beast that was, and is not, even he is the eighth, and is of the seven, and goeth into perdition}. See also Iren. V. 26, § 1.}. A blasphemer the man is and lawless, not having received the kingdom from his fathers, but having usurped the power by means of sorcery.

14. And who is this, and from what sort of working? Interpret to us, O Paul. \emph{Whose coming}, he says, \emph{is after the working of Satan, with all power and signs and lying wonders}\footnote{2 Thess. ii. 9. Lactantius (\textsc{a.d.} 300 \emph{circ.}), \emph{Div. Inst}. vii. 17: “that king.…will also be a prophet of lies; and he will constitute and call himself God, and will order himself to be worshipped as the Son of God; and power will be given him to do signs and wonders, by the sight of which he may entice men to adore him.” Cf. \emph{Epitome}, lxxi.}; implying, that Satan has used him as an instrument, working in his own person through him; for knowing that his judgment shall now no longer have respite, he wages war no more by his ministers, as is his wont, but henceforth by himself more openly\footnote{“Vid. Iren. \emph{Hær} V. 26, 2” (R.W.C.). The passage is quoted by Eusebius (\emph{Eccl. Hist}. iv. 18), from a lost work of Justin M. \emph{Against Marcion}: “Justin well said that before the coming of the Lord Satan never dared to blaspheme God, as not yet knowing his own condemnation, because it was started by the prophets in parables and allegories. But after our Lord’s advent having learnt plainly from His words and those of the Apostles that everlasting fire is prepared for him,.…he by means of such men as these blasphemes the Lord who brings the judgment upon him, as being already condemned.” \par S. Cyril seems to expect that Antichrist will be an incarnation of Satan, as did Hippolytus (\emph{de Antichr}. § 6): “The Saviour appeared in the form of man, and he too will come in the form of a man.”}. And \emph{with all signs and lying wonders}; for the father of falsehood will make a show\footnote{\textgreek{φαντασιοκοπεῖ}. See above, § 10, note 9, and the equivalent phrase in § 17: \textgreek{σημείων καὶ τεράτων φαντασίας ἐδείκνυον}.} of the works of falsehood, that the multitudes may think that they see a dead man raised, who is not raised, and lame men walking, and blind men seeing, when the cure has not been wrought.

15. And again he says, \emph{Who opposeth and exalteth himself against all that is called God, or that is worshipped}; (\emph{against every God}; Antichrist forsooth will abhor the idols,) \emph{so that he seateth himself in the temple of God}\footnote{2 Thess. ii. 4.}. What temple then? He means, the Temple of the Jews which has been destroyed. For God forbid that it should be the one in which we are! Why say we this? That we may not be supposed to favour ourselves. For if he comes to the Jews as Christ, and desires to be worshipped by the Jews, he will make great account of the Temple, that he may more completely beguile them; making it supposed that he is the man of the race of David, who shall build up the Temple which was erected by Solomon\footnote{See § 12, notes 3, 4, and Hippolytus, \emph{ubi supra}: “The Saviour raised up and shewed His holy flesh like a temple; and he will raise a temple of stone in Jerusalem.” “Cyril wrote this before Julian’s attempt to rebuild the Jewish Temple” (\textsc{R.W.C}.).}. And Antichrist will come at the time when there shall not be left one stone upon another in the Temple of the Jews, according to the doom pronounced by our Saviour\footnote{Matt. xxiv. 2. Cyril refers the whole prophecy to the time of Christ’s second coming at the end of the world, not regarding the destruction of Jerusalem and its Temple by Titus as fulfilling any part of the prediction.}; for when, either decay of time, or demolition ensuing on pretence of new buildings, or from any other causes, shall have overthrown all the stones, I mean not merely of the outer circuit, but of the inner shrine also, where the Cherubim were, then shall he come \emph{with all signs and lying wonders}, exalting himself against all idols; at first indeed making a pretence of benevolence, but afterwards displaying his relentless temper, and that chiefly against the Saints of God. For he says, \emph{I beheld, and the same horn made war with the saints}\footnote{Dan. vii. 21. Here again Cyril follows Hippolytus, § 25: “And under this (horn) was signified none other than Antichrist.}; and again elsewhere, \emph{there shall be a time of trouble, such as never was since there was a nation upon earth, even to that same time}\footnote{Ib. xii. 1.}. Dreadful is that beast, a mighty dragon, unconquerable by man, ready to devour; concerning whom though we have more things to speak out of the divine Scriptures, yet we will content ourselves at present with thus much, in order to keep within compass.

16. For this cause the Lord knowing the greatness of the adversary grants indulgence to the godly, saying, \emph{Then let them which be in Judæa flee to the mountains}\footnote{Matt. xxiv. 16.}. But if any man is conscious that he is very stout-hearted, to encounter Satan, let him stand (for I do not despair of the Church’s nerves), and let him say, \emph{Who shall separate us from the love of Christ and the rest}\footnote{Rom. viii. 35.}? But, let those of us who are fearful provide for our own safety; and those who are of a good courage, stand fast: \emph{for then shall be great tribulation, such as hath not been from the beginning of the world until now, no, nor ever shall be}\footnote{Matt. xxiv. 21.}. But thanks be to God who hath confined the greatness of that tribulation to a few days; for He says, \emph{But for the elect’s sake those days shall be shortened}\footnote{Ib. v. 22.}; and Antichrist shall reign for three years and a half only. We speak not from apocryphal books, but from Daniel; for he says, \emph{And they shall be given into his hand until a time and times and half a time}\footnote{Dan. vii. 25. By “apocryphal” books Cyril probably means all such as were not allowed to be read in the public service of the Church: see Cat. iv. 33, note 3; and Bp. Westcott’s note on the various meanings of the word \textgreek{ἀπόκρυφος}, \emph{Hist. of the Canon}, P. III. c. 1. That the Apocalypse of St. John is included under this term by Cyril, appears probable from the following reasons suggested by the Benedictine Editor. (1) It is not mentioned in the list of the Canonical Scriptures in iv. 36. (2) The earlier writers whom Cyril follows in this Lecture, Irenæus, \emph{Hær.} V., 26, § 1, and Hippolytus, \emph{De Antichristo}, § 34, combine the testimony of the Apocalypse with that of Daniel. The omission in Cyril therefore cannot have been accidental.}\emph{. A time} is the one year in which his coming shall for a while have increase; and \emph{the times} are the remaining two years of iniquity, making up the sum of the three years; and \emph{the half a time} is the six months. And again in another place Daniel says the same thing, \emph{And he sware by Him that liveth for ever that it shall be for a time, and times, and half a time}\footnote{Dan. xii. 7.}. And some peradventure have referred what follows also to this; namely, \emph{a thousand two hundred and ninety days}\footnote{Ib. v. 11.}; and this, \emph{Blessed is he that endureth and cometh to the thousand three hundred and five and thirty days}\footnote{Ib. v. 12.}. For this cause we must hide ourselves and flee; for perhaps \emph{we shall not have gone over the cities of Israel, till the Son of Man be come}\footnote{Matt. x. 23.}.

17. Who then is the blessed man, that shall at that time devoutly witness for Christ? For I say that the Martyrs of that time excel all martyrs. For the Martyrs hitherto have wrestled with men only; but in the time of Antichrist they shall do battle with Satan in his own person\footnote{\textgreek{αὐτοπροσώπως} . See above, § 14, note 2. Some \textsc{mss.} read \textgreek{ἀντιπροσώπως}, “face to face,” as in xii. 32, \textgreek{ἀντιπρόσωπος}.}. And former persecuting kings only put to death; they did not pretend to raise the dead, nor did they make false shows\footnote{See above, § 14, note 3.} of signs and wonders. But in his time there shall be the evil inducement both of fear and of deceit, \emph{so that if it be possible the very elect shall be deceived}\footnote{Matt. xxiv. 24.}. Let it never enter into the heart of any then alive to ask, “What did Christ more? For by what power does this man work these things? Were it not God’s will, He would not have allowed them.” The Apostle warns thee, and says beforehand, \emph{And for this cause God shall send them a working of error}; (\emph{send}, that is, \emph{shall allow to happen};) not that they might make excuse, but \emph{that they might be condemned}\footnote{2 Thess. ii. 11, 12: (R\textsc{.V.}) \emph{That they all might be judged}. Cyril has \textgreek{κατακριθῶσι}}. Wherefore? \emph{They}, he says, \emph{who believed not the truth}, that is, the true Christ, \emph{but had pleasure in unrighteousness}, that is, in Antichrist. But as in the persecutions which happen from time to time, so also then God will permit these things, not because He wants power to hinder them, but because according to His wont He will through patience crown His own champions like as He did His Prophets and Apostles; to the end that having toiled for a little while they may inherit the eternal kingdom of heaven, according to that which Daniel says, \emph{And at that time thy people shall be delivered, every one that shall be found written in the book} (manifestly, the book of life); \emph{and many of them that sleep in the dust of the earth shall awake, some to everlasting life, and same to shame and everlasting contempt; and they that be wise shall shine as the brightness of the firmament; and of the many righteous}\footnote{Dan. xii. 1, 2: (R.V.) \emph{they that turn many to righteousness}. Cyril follows the rendering of the Septuagint, \textgreek{ἀπὸ τῶν δικαίων τῶν πολλῶν}, which gives no proper construction.}, \emph{as the stars for ever and ever}.

18. Guard thyself then, O man; thou hast the signs of Antichrist; and remember them not only thyself, but impart them also freely to all. If thou hast a child according to the flesh, admonish him of this now; if thou hast begotten one through catechizing\footnote{Compare 1 Cor. iv. 15: \emph{I begat you through the gospel}. Clem. Alex. \emph{Strom}. iii. c. 15: \textgreek{τῷ διὰ τῆς ἀληθοῦς κατηχήσεως γεννήσαντι κεῖταί τις μισθός}.}, put him also on his guard, lest he receive the false one as the True. For the \emph{mystery of iniquity doth already work}\footnote{2 Thess. ii. 7.}. I fear these wars of the nations\footnote{See above, §§ 6, 7.}; I fear the schisms of the Churches; I fear the mutual hatred of the brethren. But enough on this subject; only God forbid that it should be fulfilled in our days; nevertheless, let us be on our guard. And thus much concerning Antichrist.

19. But let us wait and look for the Lord’s coming upon the clouds from heaven. Then shall Angelic trumpets sound; \emph{the dead in Christ shall rise first}\footnote{1 Thes. iv. 16.},—the godly persons who are alive shall be caught up in the clouds, receiving as the reward of their labours more than human honour, inasmuch as theirs was a more than human strife; according as the Apostle Paul writes, saying, \emph{For the Lord Himself shall descend from heaven with a shout, with the voice of the Archangel, and with the trump of God: and the dead in Christ shall rise first. Then we which are alive and remain shall be caught up together with them in the clouds, to meet the Lord in the air; and so shall we ever be with the Lord}\footnote{Ib. vv. 16, 17.}.

20. This coming of the Lord, and the end of the world, were known to the Preacher; who says, \emph{Rejoice, O young man, in thy youth}, and the rest\footnote{Eccles. xi. 9. The Preacher’s description of old age and death is interpreted by Cyril of the end of the world, as it had been a century before by Gregory Thaumaturgus, in his paraphrase of the book.}; \emph{Therefore remove anger\footnote{Ib. \emph{v}. 10: (\textsc{R.V.}) \emph{sorrow}. Marg. Or, \emph{vexation}, Or, \emph{provocation}.} from thy heart, and put away evil from thy flesh};…\emph{and remember thy Creator}…\emph{or ever the evil days come\footnote{Ib. xii. 1.}},.…\emph{or ever the sun, and the light, and the moon, and the stars be darkened}\footnote{Ib. v. 2.},.…\emph{and they that look out of the windows be darkened}\footnote{Ib. v. 3.}; (signifying the faculty of sight;) \emph{or ever the silver cord be loosed}; (meaning the assemblage of the stars, for their appearance is like silver;) \emph{and the flower of gold be broken}\footnote{Ib. v. 6. According to the usual interpretation death is here represented by the breaking of a chain and the lamp which hangs from it. Cf. Delitzsch, and \emph{Speaker’s Commentary}, in loc. for other interpretations. \par \textgreek{τὸ ἀνθέμιον τοῦ χρυσιόυ} (Sept.), by which Cyril understood camomile (\textgreek{ἀνθεμίς}), more probably meant a pattern of flowers embossed on the vessel of gold: \emph{vid}. Xenoph. \emph{Anab}. V. 4, § 32: \textgreek{ἐστιγμένους ἀνθέμια}, “damasked with \emph{flowers}.”}; (thus veiling the mention of the golden sun; for the camomile is a well-known plant, having many ray-like leaves shooting out round it;) \emph{and they shall rise up at the voice of the sparrow, yea, they shall look away from the height, and terrors shall be in the way}\footnote{Eccles. xii. 5. Cyril means rightly that the aged shrink from a giddy height, and from imaginary dangers of the road. For \emph{the voice of the sparrow}, see below, § 21, note 4.}. What shall they see? \emph{Then shall they see the Son of man coming on the clouds of heaven; and they shall mourn tribe by tribe}\footnote{Matt. xxiv. 30; Zech. xii. 12.}. And what shall come to pass when the Lord is come? \emph{The almond tree shall blossom, and the grasshopper shall grow heavy, and the caper-berry shall be scattered abroad}\footnote{Eccles. xii. 5.}. And as the interpreters say, the blossoming almond signifies the departure of winter; and our bodies shall then after the winter blossom with a heavenly flower\footnote{“Dr. Thomson (\emph{The Land and the Book}, p. 319) says of the almond tree, “It is the type of old age, whose hair is white” (Speaker’s Commentary).}. \emph{And the grasshopper shall grow in substance} (that means the winged soul clothing itself with the body\footnote{The step, once as active as a \emph{grasshopper}, or locust, shall grow heavy and slow. For other interpretations see Delitzsch.},) \emph{and the caper-berry shall be scattered abroad} (that is, the transgressors who are like thorns shall be scattered\footnote{\emph{The caper-berry} (\textgreek{κάππαρις}) \emph{shall fail}, i.e. no longer stimulate appetite. But \textgreek{διασχεδασθήσεται} (Sept. Cyril) means that the old man shall be like a caper-berry which when fully ripe bursts it husks and scatters its seeds: so \textsc{R.V.} (Margin); \emph{The caper-berry shall burst}. Greg. Thaumat. \emph{Metaphr. Eccles}. “The transgressors are cast out of the way, like a black and despicable caper-plant.”}).

21. Thou seest how they all foretell the coming of the Lord. Thou seest how they know \emph{the voice of the sparrow}. Let us know what sort of voice this is. \emph{For the Lord Himself shall descend from heaven with a shout, with the voice of the Archangel, and with the trump of God}\footnote{1 Thes. ii. 16.}. The Archangel shall make proclamation and say to all, \emph{Arise to meet the Lord}\footnote{Compare the spurious \emph{Apocalypse of John}: “And at the voice of the bird every plant shall arise; that is, At the voice of the Archangel all the human race shall arise” (English Trs. \emph{Ante-Nic. Libr}. p. 496). According to the Talmud the meaning is, “Even a bird awakes him” (Delitzsch).}. And fearful will be that descent of our Master. David says, \emph{God shall manifestly come, even our God, and shall not keep silence; a fire shall burn before Him, and a fierce tempest round about Him}, and the rest\footnote{Ps. l. 3.}. The Son of Man shall come to the Father, according to the Scripture which was just now read, \emph{on the clouds of heaven}, drawn by \emph{a stream of fire}\footnote{Dan. vii. 13, 10.}, which is to make trial of men. Then if any man’s works are of gold, he shall be made brighter; if any man’s course of life be like stubble, and unsubstantial, it shall be burnt up by the fire\footnote{1 Cor. iii. 12, 13. On \textgreek{ἀνυπόστατον}, see Index. On \textgreek{δοκιμαστικόν}, compare \emph{The Teaching of the Apostles}, § 16: “Then all created mankind shall come to the fire of testing (\textgreek{δοκιμασίας}), and many shall be offended and perish.”}. And the Father \emph{shall sit}, having \emph{His garment white as snow, and the hair of His head like pure wool}\footnote{Dan. vii. 9.}. But this is spoken after the manner of men; wherefore? Because He is the King of those who have not been defiled with sins; \emph{for}, He says, \emph{I will make your sins white as snow, and as wool}\footnote{Is. i. 18.}, which is an emblem of forgiveness of sins, or of sinlessness itself. But the Lord who shall come from heaven on the clouds, is He who ascended on the clouds; for He Himself hath said, \emph{And they shall see the Son of Man coming on the clouds of heaven, with power and great glory}\footnote{Matt. xxiv. 30.}.

22. But what is the sign of His coming? lest a hostile power dare to counterfeit it. \emph{And then shall appear}, He says, \emph{the sign of the Son of Man in heaven}\footnote{Ib.}. Now Christ’s own true sign is the Cross; a sign of a luminous Cross shall go before the King\footnote{Cat. xiii. 4. In the letter to Constantius, three or four years later than this Lecture, Cyril treats the appearance at that time of a luminous Cross in the sky as a fulfilment of Matt. xxiv. 30: but he there adds (\emph{Ep. ad Constantium}, § 6) that our Lord’s prediction “was both fulfilled at that present time, and shall again be fulfilled more largely.” On the opinion that “the sign of the Son of Man in heaven” should be the Cross, see Suicer, \emph{Thesaurus}, \textgreek{Σταυρός}. It is not improbable that the earliest trace of this interpretation is found in \emph{The Teaching of the Apostles}, § 16: “Then shall appear the signs of the Truth: the first the sign of a (cross) spreading out (\textgreek{ἐκπετάσεως}) in heaven.”}, plainly declaring Him who was formerly crucified: that the Jews who before \emph{pierced Him} and plotted against Him, when they see it, may \emph{mourn tribe by tribe}\footnote{Zech. xii. 12.}, saying, “This is He who was buffeted, this is He whose face they spat on, this is He whom they bound with chains, this is He whom of old they crucified, and set at nought\footnote{Cf. Barnab. \emph{Epist}. c. vii.: “For they shall see Him in that day wearing the long scarlet robes about His flesh, and shall say, Is not this He, whom once we crucified, and set at nought, and spat upon (\emph{al}. and pierced, and mocked)?”}. Whither, they will say, shall we flee from the face of Thy wrath?” But the Angel hosts shall encompass them, so that they shall not be able to flee anywhere. The sign of the Cross shall be a terror to His foes; but joy to His friends who have believed in Him, or preached Him, or suffered for His sake. Who then is the happy man, who shall then be found a friend of Christ? That King, so great and glorious, attended by the Angel-guards, the partner of the Father’s throne, will not despise His own servants. For that His Elect may not be confused with His foes, \emph{He shall send forth His Angels with a great trumpet, and they shall gather together His elect from the four winds}\footnote{Matt. xxiv. 31.}. He despised not Lot, who was but one; how then shall He despise many righteous? \emph{Come, ye blessed of My Father}\footnote{Ib. xxv. 34.}, will He say to them who shall then ride on chariots of clouds, and be assembled by Angels.

23. But some one present will say, “I am a poor man,” or again, “I shall perhaps be found at that time sick in bed;” or, “I am but a woman, and I shall be taken at the mill: shall we then be despised?” Be of good courage, O man; the Judge is no respecter of persons; \emph{He will not judge according to a man’s appearance, nor reprove according to his speech}\footnote{Is. xi. 3: (R.V.) \emph{He shall not judge after the sight of his eyes, nor reprove after the hearing of his ears.}}. He honours not the learned before the simple, nor the rich before the needy. Though thou be in the field, the Angels shall take thee; think not that He will take the landowners, and leave thee the husbandman. Though thou be a slave, though thou be poor, be not any whit distressed; He who \emph{took the form of a servant}\footnote{Phil. ii. 7.} despises not servants. Though thou be lying sick in bed, yet it is written, \emph{Then shall two be in one bed; the one shall be taken, and the other left}\footnote{Luke xvii. 34.}. Though thou be of compulsion put to grind, whether thou be man or woman\footnote{Ib. v. 35.}; though thou be in fetters\footnote{The Jerusalem \textsc{ms.} (A) alone has the true reading \textgreek{πέδας}, which is confirmed by \textgreek{πεπεδημένους} in the quotation following, instead of \textgreek{παῖδας}, which is quite inappropriate, and evidently an itacism.}, and sit beside the mill, yet He \emph{who by His might bringeth out them that are bound}\footnote{Ex. xi. 5.}, will not overlook thee. He who brought forth Joseph out of slavery and prison to a kingdom, shall redeem thee also from thy afflictions into the kingdom of heaven. Only be of good cheer, only work, only strive earnestly; for nothing shall be lost. Every prayer of thine, every Psalm thou singest is recorded; every alms-deed, every fast is recorded; every marriage duly observed is recorded; continence\footnote{\textgreek{᾽Εγκράτεια}. “Id est viduitas” (Ben. Ed.). This special reference of the word to widowhood is to some extent confirmed by 1 Cor. vii. 9: \textgreek{εἰ δὲ οὐκ ἐγκρατεύονται}, and is rendered highly probable by Cyril’s separate mention of marriage and virginity.} kept for God’s sake is recorded; but the first crowns in the records are those of virginity and purity; and thou shalt shine as an Angel. But as thou hast gladly listened to the good things, so listen again without shrinking to the contrary. Every covetous deed of thine is recorded; thine every act of fornication is recorded, thine every false oath is recorded, every blasphemy, and sorcery, and theft, and murder. All these things are henceforth to be recorded, if thou do the same now after having been baptized; for thy former deeds are blotted out.

24. \emph{When the Son of Man}, He says, \emph{shall come in His glory, and all the Angels with Him}\footnote{Matt. xxv. 31.}. Behold, O man, before what multitudes thou shalt come to judgment. Every race of mankind will then be present. Reckon, therefore, how many are the Roman nation; reckon how many the barbarian tribes now living, and how many have died within the last hundred years; reckon how many nations have been buried during the last thousand years; reckon all from Adam to this day. Great indeed is the multitude; but yet it is little, for the Angels are many more. They are \emph{the ninety and nine sheep}, but mankind is the single \emph{one}\footnote{Matt. xviii. 12; Luke xv. 4. Ambrose, \emph{Expos. in Luc}. VII. 210: “Rich is that shepherd of whose flock we are but the one hundredth part. Of Angels and Archangels, of Dominions, Powers, Thrones, and others He hath countless flocks, whom He hath left upon the mountains.” Cf. Gregor, Nyss. \emph{Contra Eunom. Or}. xii.}. For according to the extent of universal space, must we reckon the number of its inhabitants. The whole earth is but as a point in the midst of the one heaven, and yet contains so great a multitude; what a multitude must the heaven which encircles it contain? And must not the heaven of heavens contain unimaginable numbers\footnote{There is much variation in the reading and punctuation of this passage. I have followed the text adopted by the Jerusalem Editor with Codd. A. Roe. Casaub. and Grodecq, in preference to the Benedictine text, with which the Editor himself is dissatisfied.}? And it is written, \emph{Thousand thousands ministered unto Him, and ten thousand times ten thousand stood before Him}\footnote{Dan. vii. 10.}; not that the multitude is only so great, but because the Prophet could not express more than these. So there will be present at the judgment in that day, God, the Father of all, Jesus Christ being seated with Him, and the Holy Ghost present with Them; and an angel’s trumpet shall summon us all to bring our deeds with us. Ought we not then from this time forth to be sore troubled? Think it not a slight doom, O man, even apart from punishment, to be condemned in the presence of so many. Shall we not choose rather to die many deaths, than be condemned by friends?

25. Let us dread then, brethren, lest God condemn us; who needs not examination or proofs, to condemn. Say not, In the night I committed fornication, or wrought sorcery, or did any other thing, and there was no man by. Out of thine own conscience shalt thou be judged, thy \emph{thoughts the meanwhile accusing or else excusing, in the day when God shall judge the secrets of men\footnote{Rom. ii. 15, 16.}}. The terrible countenance of the Judge will force thee to speak the truth; or rather, even though thou speak not, it will convict thee. For thou shalt rise clothed with thine own sins, or else with thy righteous deeds. And this has the Judge Himself declared—for it is Christ who judges—\emph{for neither doth the Father judge any man, but he hath given all judgment unto the Son}\footnote{John v. 22.}, not divesting Himself of His power, but judging through the Son; the Son therefore judgeth by the will\footnote{\textgreek{νεύματι}. Cat. xi. 22.} of the Father; for the wills of the Father and of the Son are not different, but one and the same. What then says the Judge, as to whether thou shalt bear thy works, or no? \emph{And before Him shall they gather all nations}\footnote{Matt. xxv. 32.}: (for in the presence of Christ \emph{every knee must bow, of things in heaven, and things in earth, and things under the earth}\footnote{Phil. ii. 10.}:) \emph{and He shall separate them one from another, as the shepherd divideth his sheep from the goats}. How does the shepherd make the separation? Does he examine out of a book which is a sheep and which a goat? or does he distinguish by their evident marks? Does not the wool show the sheep, and the hairy and rough skin the goat? In like manner, if thou hast been just now cleansed from thy sins, thy deeds shall be henceforth as pure wool; and thy robe shall remain unstained, and thou shalt ever say, \emph{I have put off my coat, how shall I put it on}\footnote{Cant. V. 3. Compare Cat. iii. 7; xx. (Mystag. ii.) 2.}? By thy vesture shalt thou be known for a sheep. But if thou be found hairy, like Esau, who was rough with hair, and wicked in mind, who for food lost his birthright and sold his privilege, thou shalt be one of those on the left hand. But God forbid that any here present should be cast out from grace, or for evil deeds be found among the ranks of the sinners on the left hand!

26. Terrible in good truth is the judgment, and terrible the things announced. The kingdom of heaven is set before us, and everlasting fire is prepared. How then, some one will say, are we to escape the fire? And how to enter into the kingdom? \emph{I was an hungered}, He says, \emph{and ye gave Me meat}. Learn hence the way; there is here no need of allegory, but to fulfil what is said. \emph{I was an hungered, and ye gave Me meat; I was thirsty, and ye gave Me drink; I was a stranger, and ye took Me in; naked, and ye clothed Me; I was sick, and ye visited Me; I was in prison, and ye came unto Me}\footnote{Matt. xxv. 35.}. These things if thou do, thou shalt reign together with Him; but if thou do them not, thou shalt be condemned. At once then begin to do these works, and abide in the faith; lest, like the foolish virgins, tarrying to buy oil, thou be shut out. Be not confident because thou merely possessest the lamp, but constantly keep it burning. Let the light of thy good works shine before men\footnote{Matt. v. 16.}, and let not Christ be blasphemed on thy account. Wear thou a garment of incorruption\footnote{The prayer for the Catechumens in the \emph{Apostolic Constitutions}, viii. 6, contains a petition that God would “vouchsafe to them the laver of regeneration, and the garment of incorruption, which is the true life.”}, resplendent in good works; and whatever matter thou receivest from God to administer as a steward, administer profitably. Hast thou been put in trust with riches? Dispense them well. Hast thou been entrusted with the word of teaching? Be a good steward thereof. Canst thou attach the souls of the hearers\footnote{\textgreek{προσθεῖναι}. Cf. Acts ii. 41: \textgreek{προσετέθησαν} . According to some \textsc{mss.} the sentence would run thus: “Hast thou been entrusted with the word of teaching? Be a good steward of thy hearers’ souls. Hast thou power to rule (\textgreek{προστῆναι})? Do this diligently.”}? Do this diligently. There are many doors of good stewardship. Only let none of us be condemned and cast out; that we may with boldness meet Christ the Everlasting King, who reigns for ever. For He doth reign for ever, who shall be judge of quick and dead, because for quick and dead He died. And as Paul says, \emph{For to this end Christ both died and lived again, that He might be Lord both of the dead and living}\footnote{Rom. xiv. 9.}.

27. And shouldest thou ever hear any say that the kingdom of Christ shall have an end, abhor the heresy; it is another head of the dragon, lately sprung up in . A certain one has dared to affirm, that after the end of the world Christ shall reign no longer\footnote{Marcellus, Bishop of Ancyra, and his pupil Photinus, are anathematized in the Creed called \textgreek{Μακρόστιχος} as holding that Christ first became “Son of God when He took our flesh from the Virgin.…For they will have it that then Christ began His Kingdom, and that it will have an end after the consummation of all and the judgment. Such are the disciples of Marcellus and Scotinus of Galatian Ancyra, \&c.” See Newman on Athanasius, \emph{de Synodis}, § 26, (5), notes \emph{a} and \emph{b}. Compare the description of Marcellus in the Letter of the Oriental Bishops who had withdrawn from the Council of Sardica to Philippopolis (\textsc{a.d.} 344). “There has arisen in our days a certain Marcellus of Galatia, the most execrable pest of all heretics, who with sacrilegious mind, and impious mouth, and wicked argument seeks to set bounds to the perpetual, eternal, and timeless kingdom of our Lord Christ, saying that He began to reign four hundred years since, and shall end at the dissolution of the present world” (Hilar. Pictav. \emph{Ex Opere Hist.} Fragm. iii.).}; he has also dared to say, that the Word having come forth from the Father shall be again absorbed into the Father, and shall be no more\footnote{“The person meant by Cyril, though he withholds the name, is Marcellus of Ancyra; who having written a book against the Arian Sophist Asterius to explain the Apostle’s statement concerning the subjection of the Son to the Father, was thought to be renewing the heresy of Paul of Samosata. On this account he was reproved by the Bishops at the Council of Jerusalem, \textsc{a.d.} 335, for holding false opinions, and being ordered to recant his opinions promised to burn his book. Afterwards he applied to Constantine, by whom he was remitted to the Council of Constantinople, \textsc{a.d.} 336, and deposed by the Bishops. As however he was acquitted by the Councils of Rome, \textsc{a.d.} 342, and of Sardica, \textsc{a.d.} 347, it became a matter of dispute whether he was really heretical.…From the fragments of his books transcribed by Eusebius, you may possibly acquit him of the Sabellian heresy and the confusion of the Father and the Son, but certainly not of the heresy concerning the end of Christ’s kingdom, and the abandonment by the Word of the human nature which He assumed for our sake; so express are his words recorded by Eusebius in the beginning of the 2nd Book \emph{Contra Marcellum}, pp. 50, 51.” (Ben. Ed.) Cf. \emph{Dict. Chr. Biogr.} “Eusebius of Cæsarea,” p. 341; and note 3 on § 9 above.}; uttering such blasphemies to his own perdition. For he has not listened to the Lord, saying, The Son abideth for ever\footnote{John viii. 25.}. He has not listened to Gabriel, saying, And He shall reign over the house of Jacob for ever, and of His kingdom there shall be no end\footnote{Luke i. 33.}. Consider this text. Heretics of this day teach in disparagement of Christ, while Gabriel the Archangel taught the eternal abiding of the Saviour; whom then wilt thou rather believe? wilt thou not rather give credence to Gabriel? Listen to the testimony of Daniel in the text\footnote{\textgreek{τὴν παροῦσαν}.}; I saw in a vision of the night, and behold, one like the Son of Man came with the clouds of heaven, and came to the Ancient of days.….And to Him was given the honour, and the dominion, and the kingdom: and all peoples, tribes, and languages shall serve Him; His dominion is an everlasting dominion, which shall not pass away, and His kingdom shall not be destroyed\footnote{Dan. vii. 13, 14.}. These things rather hold fast, these things believe, and cast away from thee the words of heresy; for thou hast heard most plainly of the endless kingdom of Christ.

28. The like doctrine thou has also in the interpretation of the \emph{Stone, which was cut out of a mountain without hands}, which is \emph{Christ according to the flesh}\footnote{Ib. ii. 45; Rom. ix. 5.}; \emph{And His kingdom shall not be left to another people}. David also says in one place, \emph{Thy throne, O God, is for ever and ever}\footnote{Ps. xlv. 6.}; and in another place, \emph{Thou, Lord, in the beginning hast laid the foundations of the earth, \&c., they shall perish, but Thou remainest, \&c.; but Thou art the same, and Thy years shall not fail}\footnote{Ib. cii. 25–27.}: words which Paul has interpreted of the Son\footnote{Heb. i. 10–12.}.

29. And wouldest thou know how they who teach the contrary ran into such madness? They read wrongly that good word of the Apostle, \emph{For He must reign, till He hath put all enemies under His feet}\footnote{1 Cor. xv. 25.}; and they say, when His enemies shall have been put under His feet, He shall cease to reign, wrongly and foolishly alleging this. For He who is king before He has subdued His enemies, how shall He not the rather be king, after He has gotten the mastery over them.

30. They have also dared to say that the Scripture, \emph{When all things shall be subjected unto Him, then shall the Son also Himself be subjected unto Him that subjected all things unto Him}\footnote{1 Cor. xv. 28. Theodoret \emph{Comment. in Epist. i. ad Cor.} xv. 28: “This passage the followers of Arius and Eunomius carry continually on their tongue, thinking in this way to disparage the dignity of the Only-begotten.”},—that this Scripture shews that the Son also shall be absorbed into the Father. Shall ye then, O most impious of all men, ye the creatures of Christ, continue? and shall Christ perish, by whom both you and all things were made? Such a word is blasphemous. But further, how shall all things be made subject unto Him? By perishing, or by abiding? Shall then the other things, when subject to the Son abide, and shall the Son, when subject to the Father, not abide? For He shall be subjected, not because He shall then begin to do the Father’s will (for from eternity He \emph{doth} always \emph{those things that please Him\footnote{John viii. 29.}}), but because, then as before, He obeys the Father, yielding, not a forced obedience, but a self-chosen accordance; for He is not a servant, that He should be subjected by force, but a Son, that He should comply of His free choice and natural love.

31. But let us examine them; what is the meaning of “until” or “as long as?” For with the very phrase will I close with them, and try to overthrow their error. Since they have dared to say that the words, \emph{till He hath put His enemies under His feet}, shew that He Himself shall have an end, and have presumed to set bounds to the eternal kingdom of Christ, and to bring to an end, as far as words go, His never-ending sovereignty, come then, let us read the like expressions in the Apostle: \emph{Nevertheless, death reigned from Adam till Moses}\footnote{Rom. v. 14. “\textgreek{ἄχρι} from \textgreek{ἄκρος}, as \textgreek{μέχρι} from \textgreek{μῆκος, μακρός}” (L. and Sc.). It is not always possible to mark this distinction in translation: cf. Lobeck, \emph{Phrynichus}, p. 14; Viger, \emph{De Idiot. Gr.} p. 419.}. Did men then die up to that time, and did none die any more after Moses, or after the Law has there been no more death among men? Well then, thou seest that the word “unto” is not to limit time; but that Paul rather signified this,—“And yet, though Moses was a righteous and wonderful man, nevertheless the doom of death, which was uttered against Adam, reached even unto him, and them that came after him; and this, though they had not committed the like sins as Adam, by his disobedience in eating of the tree.”

32. Take again another similar text. \emph{For until this day…when Moses is read, a vail lieth upon their heart\footnote{2 Cor. iii. 14, 15.}}. Does \emph{until this day} mean only “until Paul?” Is it not \emph{until this day} present, and even to the end? And if Paul say to the Corinthians, \emph{For we came even as far as unto you in preaching the Gospel of Christ, having hope when your faith increases to preach the Gospel in the regions beyond you}\footnote{Ib. x. 14, 15, 16.}, thou seest manifestly that \emph{as far as} implies not the end, but has something following it. In what \emph{sense then shouldest thou remember that Scripture, till} He hath put all enemies under His feet\footnote{1 Cor. xv. 25.}? According as Paul says in another place, \emph{And exhort each other daily, while it is called to-day}\footnote{Heb. iii. 13.}; meaning, “continually.” For as we may not speak of the “beginning of the days” of Christ, so neither suffer thou that any should ever speak of the end of His kingdom. For it is written, \emph{His kingdom is an everlasting kingdom}\footnote{Dan. vii. 14, 27.}.

33. And though I have many more testimonies out of the divine Scriptures, concerning the kingdom of Christ which has no end for ever, I will be content at present with those above mentioned, because the day is far spent. But thou, O hearer, worship only Him as thy King, and flee all heretical error. And if the grace of God permit us, the remaining Articles also of the Faith shall be in good time declared to you. And may the God of the whole world keep you all in safety, bearing in mind the signs of the end, and remaining unsubdued by Antichrist. Thou hast received the tokens of the Deceiver who is to come; thou hast received the proofs of the true Christ, who shall openly come down from heaven. Flee therefore the one, the False one; and look for the other, the True. Thou hast learnt the way, how in the judgment thou mayest be found among those on the right hand; guard \emph{that which is committed to thee}\footnote{1 Tim. vi. 20.} concerning Christ, and be conspicuous in good works, that thou mayest stand with a good confidence before the Judge, and inherit the kingdom of heaven:—Through whom, and with whom, be glory to God with the Holy Ghost, for ever and ever. Amen.

\xchapter{Lecture XVI. On the Article, And in One Holy Ghost, the Comforter, Which Spake in the Prophets.}

\textsc{1 Corinthians xii. 1, 4}

Now concerning spiritual gifts, brethren, I would not have you ignorant.…Now there are diversities of gifts, but the same Spirit, \&c.

1. \textsc{Spiritual} in truth is the grace we need, in order to discourse concerning the Holy Spirit; not that we may speak what is worthy of Him, for this is impossible, but that by speaking the words of the divine Scriptures, we may run our course without danger. For a truly fearful thing is written in the Gospels, where Christ has plainly said, \emph{Whosoever shall speak a word against the Holy Ghost, it shall not be forgiven him, neither in this world, nor in that which is to come}\footnote{Matt. xii. 32.}. And there is often fear, lest a man should receive this condemnation, through speaking what he ought not concerning Him, either from ignorance, or from supposed reverence. The Judge of quick and dead, Jesus Christ, declared that he hath no forgiveness; if therefore any man offend, what hope has he?

2. It must therefore belong to Jesus Christ’s grace itself to grant both to us to speak without deficiency, and to you to hear with discretion; for discretion is needful not to them only who speak, but also to them that hear, lest they hear one thing, and misconceive another in their mind. Let us then speak concerning the Holy Ghost nothing but what is written; and whatsoever is not written, let us not busy ourselves about it. The Holy Ghost Himself spoke the Scriptures; He has also spoken concerning Himself as much as He pleased, or as much as we could receive. Let us therefore speak those things which He has said; for whatsoever He has not said, we dare not say.

3. There is One Only Holy Ghost, the Comforter; and as there is One God the Father, and no second Father;—and as there is One Only-begotten Son and Word of God, who hath no brother;—so is there One Only Holy Ghost, and no second spirit equal in-honour to Him. Now the Holy Ghost is a Power most mighty, a Being divine and unsearchable; for He is living and intelligent, a sanctifying principle of all things made by God through Christ. He it is who illuminates the souls of the just; He was in the Prophets, He was also in the Apostles in the New Testament. Abhorred be they who dare to separate the operation of the Holy Ghost! There is One God, the Father, Lord of the Old and of the New Testament: and One Lord, Jesus Christ, who was prophesied of in the Old Testament, and came in the New; and One Holy Ghost, who through the Prophets preached of Christ, and when Christ was come, descended, and manifested Him\footnote{At the end of this section there follows in the Coislin \textsc{ms.} a long interpolation consisting of two parts. The former is an extract taken word for word from Gregory of Nyssa, \emph{Oratio Catechetica}, ii. c, which may be read in this series: \textgreek{᾽Αλλ᾽ ὡς Θεοῦ Λόγον ἀκούσαντες .…σύνδρομον ἔχουσαν τῇ βουλήσει τὴν δύναμιν}. Of the second passage the Benedictine Editor says: “I have not been able to discover who is the author. No one can assign it to our Cyril, although the doctrine it contains is in full agreement with his: but he explains all the same points more at large in his two Lectures (xvi. xvii.). The passage is very ancient and undoubtedly older than the eleventh century, which is the date of the Cod. Coislin. Therefore in the controversy of the Latins against the Greeks concerning the Procession of the Holy Ghost it is important to notice what is taught in this passage, and also brought forward as a testimony by S. Thomas (Aquinas), that “The Holy Ghost is of the Godhead of the Father and the Son (ex Patris et Filii divinitate existere).” To me indeed these words seem to savour altogether not of the later but of the more ancient theology of the Greeks, and to be earlier than the controversies of the Greeks against the Latins.” \par This second passage is as follows:— \par “For the Spirit of God is good. \emph{And Thy good Spirit}, says David, \emph{shall lead me in the land of righteousness}. This then is the Spirit of God in which we believe: the blessed Spirit, the eternal, immutable, unchangeable, ineffable: which rules and reigns over all productive being, both visible and invisible natures: which is Lord both of Angels and Archangels, Powers, Principalities, Dominions, Thrones: the Creator of all being, enthroned with the glory of the Father and the Son, reigning without beginning and without end with the Father and the Son, before the created substances: Who sanctifies the \emph{ministering spirits sent forth for the sake of those who are to inherit salvation}: Who came down upon the holy and blessed Virgin Mary, of whom was born Christ according to the flesh; came down also upon the Lord Himself in bodily form of a dove in the river Jordan: Who came upon the Apostles on the day of Pentecost in form of fiery tongues; Who gives and supplies all spiritual gifts in the Church, \textsc{Who Proceedeth from the Father}: Who is of the Godhead of the Father and the Son; Who is of one substance with the Father and the Son, inseparable and indivisible.”}.

4. Let no one therefore separate the Old from the New Testament\footnote{Cf. Cat. iv. 33; vii. 6. Irenæus, \emph{Hæres}. III. xxi. 4; IV. ix. 1. In Eusebius, \emph{E.H.} V. 13, Rhodon says that Apelles attributed the prophecies to an adverse spirit and rejected them as false and self-contradictory. Similar blasphemies against the holy Prophets are imputed to Manes by Epiphanius (\emph{Hæres}. lxvi. 30).}; let no one say that the Spirit in the former is one, and in the latter another; since thus he offends against the Holy Ghost Himself, who with the Father and the Son together is honoured, and at the time of Holy Baptism is included with them in the Holy Trinity. For the Only-begotten Son of God said plainly to the Apostles, \emph{Go ye, and make disciples of all the nations, baptizing them into the name of the Father, and of the Son, and of the Holy Ghost}\footnote{Matt. xxviii. 19. The same text is used with much force by S. Basil (\emph{De Spir. S.} cap. xxiv.).}. Our hope is in Father, and Son, and Holy Ghost. We preach not three Gods\footnote{Cat. xi. 4, note 3. See Newman’s notes on Athanasius, \emph{Contra Arian. Or.} I. viii. 1; Ib. \emph{Or.} III. xxv. 9; Ib. xxvii. 3. Marcion’s doctrine of three first principles (\textgreek{τριῶν ἀρχῶν λόγος}) is discussed by Epiphanius (\emph{Hæres}. xlii. 6, 7). See also Tertull. \emph{Contra Marcion}. I. 15; Euseb. \emph{Hist. Eccles}. V. 13.}; let the Marcionites be silenced; but with the Holy Ghost through One Son, we preach One God. The Faith is indivisible; the worship inseparable. We neither separate the Holy Trinity, like some; nor do we as Sabellius work confusion.\footnote{\textgreek{συναλοιφήν}, iv. 8; xi. 16; xv. 9.} But we know according to godliness One Father, who sent His Son to be our Saviour; we know One Son, who promised that He would send the Comforter from the Father; we know the Holy Ghost, who spake in the Prophets, and who on the day of Pentecost descended on the Apostles in the form of fiery tongues, here, in Jerusalem, in the Upper Church of the Apostles\footnote{Cat. xvii. 13. Epiphanius (\emph{De Mensuris et Ponder. c.} 14): “And he (Hadrian) found the city all levelled to the ground, except a few houses, and the Church of God which was small: where the Disciples, on their return after the Saviour was taken up from the Mount of Olives, went up into the upper chamber: for there it had been built, that is on Sion.” Cf. Stanley, \emph{Sinai and Palestine}, c. xiv. 3: “Within the precincts of that Mosque (of the Tomb of David) is a vaulted Gothic chamber, which contains within its four walls a greater confluence of traditions than any other place of like dimensions in Palestine. It is startling to hear that this is the scene of the Last Supper, of the meeting after the Resurrection, of the miracle of Pentecost, of the residence and death of the Virgin, of the burial of Stephen.”}; for in all things the choicest privileges are with us. Here Christ came down from heaven; here the Holy Ghost came down from heaven. And in truth it were most fitting, that as we discourse concerning Christ and Golgotha here in Golgotha, so also we should speak concerning the Holy Ghost in the Upper Church; yet since He who descended there jointly partakes of the glory of Him who was crucified here, we here speak concerning Him also who descended there: for their worship is indivisible.

5. We would now say somewhat concerning the Holy Ghost; not to declare His substance with exactness, for this were impossible; but to speak of the diverse mistakes of some concerning him, lest from ignorance we should fall into them; and to block up the paths of error, that we may journey on the King’s one highway. And if we now for caution’s sake repeat any statement of the heretics, let it recoil on their heads, and may we be guiltless, both we who speak, and ye who hear.

6. For the heretics, who are most profane in all things, have \emph{sharpened their tongue}\footnote{Ps. cxl. 3.} against the Holy Ghost also, and have dared to utter impious things; as Irenæus the interpreter has written in his injunctions against heresies\footnote{Irenæus is called “the interpreter” in the same general sense as other ecclesiastical authors (Cat. xiii. 21; xv. 20), on account of his frequent comments upon the Scriptures. The full title of his work was \emph{A Refutation and Subversion of Knowledge falsely so called} (Euseb. \emph{Hist. Eccles}. V. c. 7). Cyril’s expression (\textgreek{ἐν τοῖς προστάγμασι}) is sufficiently appropriate to the hortatory purpose professed by Irenæus in his preface. But the Benedictine Editor thinks that the word \textgreek{προστάγμασι} may be an interpolation arising from the following words \textgreek{πρὸς τὰς}.…The meaning would then be “in his writings \emph{Against Heresies},” the usual short title of the work.}. For some of them have dared to say that they were themselves the Holy Ghost;—of whom the first was Simon\footnote{Cat. vi. 14, note 10.}, the sorcerer spoken of in the Acts of the Apostles; for when he was cast out, he presumed to teach such doctrines: and they who are called Gnostics, impious men, have spoken other things against the Spirit\footnote{Irenæus (I. xxix § 4; xxx. § 1).}, and the wicked Valentinians\footnote{Ib. I. ii. §§ 5, 6.} again something else; and the profane Manes dared to call himself the Paraclete sent by Christ\footnote{Cat. vi. 25.}. Others again have taught that the Spirit is different in the Prophets and in the New Testament.\footnote{Cat. iv. 33. See § 3, note 3, above.} Yea, and great is their error, or rather their blasphemy. Such therefore abhor, and flee from them who blaspheme the Holy Ghost, and have no forgiveness. For what fellowship hast thou with the desperate, thou, who art now to be baptized, into the Holy Ghost also\footnote{i.e. as well as into the Father and the Son.}? If he who attaches himself to a thief, and consenteth with him, is subject to punishment, what hope shall he have, who offends against the Holy Ghost?

7. Let the Marcionists also be abhorred, who tear away from the New Testament the sayings of the Old\footnote{See \emph{Dict. Christ. Biography}, Marcion, p. 283; and Tertullian (\emph{Adv. Marcion}. IV. 6): “His whole aim centres in this that he may establish a diversity between the Old and New Testaments, so that his own Christ may be separate from the Creator, as belonging to the rival god, and as alien from the Law and the Prophets.}. For Marcion first, that most impious of men, who first asserted three Gods\footnote{Cf. § 4, note 5, above.}, knowing that in the New Testament are contained testimonies of the Prophets concerning Christ, cut out the testimonies taken from the Old Testament, that the King might be left without witness. Abhor those above-mentioned Gnostics, men of knowledge by name, but fraught with ignorance; who have dared to say such things of the Holy Ghost as I dare not repeat.

8. Let the Cataphrygians\footnote{Phrygians, or Cataphrygians (\textgreek{οἱ κατὰ φρύγας}) was the name given to the followers of the Phrygian Montanus. See the account of Montanism in Eusebius, \emph{Hist. Eccl}. V. xvi., and the note there in this Series.} also be thy abhorrence, and Montanus, their ringleader in evil, and his two so-called prophetesses, Maximilla and Priscilla. For this Montanus, who was out of his mind and really mad (for he would not have said such things, had he not been mad), dared to say that he was himself the Holy Ghost,—he, miserable man, and filled with all uncleanness and lasciviousness; for it suffices but to hint at this, out of respect for the women who are present. And having taken possession of Pepuza, a very small hamlet of Phrygia, he falsely named it Jerusalem; and cutting the throats of wretched little children, and chopping them up into unholy food, for the purpose of their so-called mysteries\footnote{The charges of lust and cruelty brought against the Montanists by Cyril and Epiphanius (\emph{Hær} 48) seem to rest on no trustworthy evidence, and are not mentioned by Eusebius, a bitter foe to the sect.},—(wherefore till but lately in the time of persecution we were suspected of doing this, because these Montanists were called, falsely indeed, by the common name of Christians;)—yet he dared to call himself the Holy Ghost, filled as he was with all impiety and inhuman cruelty, and condemned by an irrevocable sentence.

9. And he was seconded, as was said before, by that most impious Manes also, who combined what was bad in every heresy\footnote{On Manes, see Cat. vi. 20. ff.}; who being the very lowest pit of destruction, collected the doctrines of all the heretics, and wrought out and taught a yet more novel error, and dared to say that he himself was the Comforter, whom Christ promised to send. But the Saviour when He promised Him, said to the Apostles, \emph{But tarry ye in the city of Jerusalem, until ye be endued with power from on high}\footnote{Luke xxiv. 49.}. What then? did the Apostles who had been dead two hundred years, wait for Manes, \emph{until they should be endued with the power}; and will any dare to say, that they were not forthwith full of the Holy Ghost? Moreover it is written, \emph{Then they laid their hands on and they received the Holy Ghost}\footnote{Acts viii. 17.}; was not this before Manes, yea, many years before, when the Holy Ghost descended on the day of Pentecost?

10. Wherefore was Simon the sorcerer condemned? Was it not that he came to the Apostles, and said, \emph{Give me also this power, that on whomsoever I lay hands, he may receive the Holy Ghost}? For he said not, “Give me also the fellowship of the Holy Ghost,” but “Give me the power;” that he might sell to others that which could not be sold, and which he did not himself possess. He offered money also to them who had no possessions\footnote{Acts viii. 19. \textgreek{ἀκτήμοσι}. Cf. § 19: \textgreek{ἀκτημονοῦσι}, and § 22: \textgreek{ἀκτημοσύνην}.}; and this, though he saw men bringing the prices of the things sold, and laying them at the Apostles’ feet. And he considered not that they who trod under foot the wealth which was brought for the maintenance of the poor, were not likely to give the power of the Holy Ghost for a bribe. But what say they to Simon? \emph{Thy money perish with thee, because thou hast thought to purchase the gift of God with money}\footnote{Ib. v. 20.}; for thou art a second Judas, for expecting to buy the grace of the Spirit with money. If then Simon, for wishing to get this power for a price, is to \emph{perish}, how great is the impiety of Manes, who said that he was the Holy Ghost? Let us hate them who are worthy of hatred; let us turn away from them from whom God turns away; let us also ourselves say unto God with all boldness concerning all heretics, \emph{Do not I hate them, O Lord, that hate Thee, and am not I grieved with Thine enemies}\footnote{Ps. cxxxix. 21.}? For there is also an enmity which is right, according as it is written, \emph{I will put enmity between thee and her seed}\footnote{Gen. iii. 15.}; for friendship with the serpent works enmity with God, and death.

11. Let then thus much suffice concerning those outcasts; and now let us return to the divine Scriptures, and let us \emph{drink waters out of our own cisterns} [that is, the holy Fathers\footnote{The words \textgreek{ἁγίων πατέρων} are not found in the \textsc{mss.} Mon. 1. Mon. 2. Vind. Roe. Casaub. nor in Grodecq. Whether meant to refer, as the Benedictine Editor thinks, to the writers of the Old Testament, or to Christian authors, they are an evident gloss.}], \emph{and out of our own springing wells}\footnote{Prov. v. 15.}. Drink we of \emph{living water, springing up into everlasting life}\footnote{John iv. 14, quoted more fully at the end of the section.}; \emph{but this spake} the Saviour \emph{of the Spirit, which they that believe on Him should receive}\footnote{Ib. vii. 38, 39.}. For observe what He says, \emph{He that believeth on Me} (not simply this, but), \emph{as the Scripture hath said} (thus He hath sent thee back to the Old Testament), \emph{out of his belly shall flow rivers of living water}, not rivers perceived by sense, and merely watering the earth with its thorns and trees, but bringing souls to the light. And in another place He says, \emph{But the water that I shall give him, shall be in him a well of living water springing up} \emph{into everlasting life},—a new kind of water living and springing up, springing up unto them who are worthy.

12. And why did He call the grace of the Spirit water? Because by water all things subsist; because water brings forth grass and living things; because the water of the showers comes down from heaven; because it comes down one in form, but works in many forms. For one fountain watereth the whole of Paradise, and one and the same rain comes down upon all the world, yet it becomes white in the lily, and red in the rose, and purple in violets and hyacinths, and different and varied in each several kind: so it is one in the palm-tree, and another in the vine, and all in all things; and yet is one in nature, not diverse from itself; for the rain does not change itself, and come down first as one thing, then as another, but adapting itself to the constitution of each thing which receives it, it becomes to each what is suitable\footnote{Compare a similar passage on rain in Cat. ix. 9, 10.}. Thus also the Holy Ghost, being one, and of one nature, and indivisible, divides to each His grace, \emph{according as He will}\footnote{1 Cor. xii. 11.}: and as the dry tree, after partaking of water, puts forth shoots, so also the soul in sin, when it has been through repentance made worthy of the Holy Ghost, brings forth clusters of righteousness. And though He is One in nature, yet many are the virtues which by the will of God and in the Name of Christ He works. For He employs the tongue of one man for wisdom; the soul of another He enlightens by Prophecy; to another He gives power to drive away devils; to another He gives to interpret the divine Scriptures. He strengthens one man’s self-command; He teaches another the way to give alms; another He teaches to fast and discipline himself; another He teaches to despise the things of the body; another He trains for martyrdom: diverse in different men, yet not diverse from Himself, as it is written, \emph{But the manifestation of the Spirit is given to every man to profit withal. For to one is given through the Spirit the word of wisdom; and to another the word of knowledge according to the same Spirit; to another faith, in the same Spirit; and to another gifts of healing, in the same Spirit; and to another workings of miracles; and to another prophecy; and to another discernings of spirits; and to another divers kinds of tongues; and to another the interpretation of tongues: but all these worketh that one and the same Spirit, dividing to every man severally as He will}\footnote{Ib. vv. 7–11.}.

13. But since concerning spirit in general many diverse things are written in the divine Scriptures, and there is fear lest some out of ignorance fall into confusion, not knowing to what sort of spirit the writing refers; it will be well now to certify you, of what kind the Scripture declares the Holy Spirit to be. For as Aaron is called Christ, and David and Saul and others are called Christs\footnote{See Cat. x. 11; xi. 1.}, but there is only one true Christ, so likewise since the name of spirit is given to different things, it is right to see what is that which is distinctively called the Holy Spirit. For many things are called spirits. Thus an Angel is called spirit, our soul is called spirit, and this wind which is blowing is called spirit; great virtue also is spoken of as spirit; and impure practice is called spirit; and a devil our adversary is called spirit. Beware therefore when thou hearest these things, lest from their having a common name thou mistake one for another. For concerning our soul the Scripture says, \emph{His spirit shall go forth, and he shall return to his earth}\footnote{Ps. cxlvi. 4.}: and of the same soul it says again, \emph{Which formeth the spirit of man within him}\footnote{Zech. xii. 1.}. And of the Angels it is said in the Psalms, \emph{Who maketh His Angels spirits, and His ministers a flame of fire}\footnote{Ps. civ. 4.}. And of the wind it saith, \emph{Thou shalt break the ships of Tarshish with a violent spirit}\footnote{Ps. xlviii. 7.}; and, \emph{As the tree in the wood is shaken by the spirit}\footnote{Is. vii. 2.}; and, \emph{Fire, hail, snow, ice, spirit of storm}\footnote{Ps. cxlviii. 8.}. And of good doctrine the Lord Himself says, \emph{The words that I have spoken unto you, they are spirit}\footnote{John vi. 63.}, \emph{and they are life}; instead of, “are spiritual.” But the Holy Spirit is not pronounced by the tongue; but He is a Living Spirit, who gives wisdom of speech, Himself speaking and discoursing.

14. And wouldest thou know that He discourses and speaks? Philip by revelation of an Angel went down to the way which leads to Gaza, when the Eunuch was coming; and the Spirit said to Philip, \emph{Go near, and join thyself to this chariot}\footnote{Acts viii. 29.}. Seest thou the Spirit talking to one who hears Him? Ezekiel also speaks thus, \emph{The Spirit of the Lord came upon me, and said unto me, Thus saith the Lord}\footnote{Ezek. xi. 5.}. And again, \emph{The Holy Ghost said\footnote{Acts xiii. 2.}}, unto the Apostles who were in Antioch, \emph{Separate me now Barnabas and Saul for the work whereunto I have called them}. Beholdest thou the Spirit living, separating, calling, and with authority sending forth? Paul also said, \emph{Save that the Holy Ghost witnesseth in every city, saying that bonds and afflictions await me}\footnote{Ib. xx. 23.}. For this good Sanctifier of the Church, and her Helper, and Teacher, the Holy Ghost, the Comforter, of whom the Saviour said, \emph{He shall} \emph{teach you all things} (and He said not only, \emph{He shall teach}, but also, \emph{He shall bring to your remembrance whatever I have said unto you}\footnote{John xiv. 26.}; for the teachings of Christ and of the Holy Ghost are not different, but the same)—He, I say, testified before to Paul what things should befall him, that he might be the more stout-hearted, from knowing them beforehand. Now I have spoken these things unto you because of the text, \emph{The words which I have spoken unto you, they are spirit}; that thou mayest understand this, not of the utterance of the lips\footnote{Ib. vi. 63. The Holy Spirit is more than words pronounced by the tongue, even than our Lord’s own words, which he called \emph{spirit}.}, but of the good doctrine in this passage.

15. But sin also is called spirit, as I have already said; only in another and opposite sense, as when it is said, \emph{The spirit of whoredom caused them to err}\footnote{Hosea iv. 12.}. The name “spirit” is given also to the \emph{unclean spirit}, the devil; but with the addition of, “the unclean;” for to each is joined its distinguishing name, to mark its proper nature. If the Scripture speak of the soul of man, it says \emph{the spirit} with the addition, \emph{of the man}; if it mean the wind, it says, \emph{spirit of storm}; if sin, it says, \emph{spirit of whoredom}; if the devil, it says, \emph{an unclean spirit}: that we may know which particular thing is spoken of, and thou mayest not suppose that it means the Holy Ghost; God forbid! For this name of spirit is common to many things; and every thing which has not a solid body is in a general way called spirit\footnote{Origen, \emph{de Principiis}, i. § 2: “It is the custom of Holy Scripture, when it would designate anything contrary to this more dense and solid body, to call it spirit.”}. Since, therefore, the devils have not such bodies, they are called spirits: but there is a great difference; for the unclean devil, when he comes upon a man’s soul (may the Lord deliver from him every soul of those who hear me, and of those who are not present), he comes like a wolf upon a sheep, ravening for blood, and ready to devour. His coming is most fierce; the sense of it most oppressive; the mind becomes darkened; his attack is an injustice also, and so is his usurpation of another’s possession. For he makes forcible use of another’s body, and another’s instruments, as if they were his own; he throws down him who stands upright (for he is akin to him who \emph{fell from heaven}\footnote{Luke x. 18.}); he twists the tongue and distorts the lips; foam comes instead of words; the man is filled with darkness; his eye is open, yet the soul sees not through it; and the miserable man gasps convulsively at the point of death. The devils are verily foes of men, using them foully and mercilessly.

16. Such is not the Holy Ghost; God forbid! For His doings tend the contrary way, towards what is good and salutary. First, His coming is gentle; the perception of Him is fragrant; His burden most light; beams of light and knowledge gleam forth before His coming\footnote{In this contrast between the evil spirit and the Spirit of God Cyril’s language rises to true eloquence, far surpassing a somewhat similar description, which may have been known to him, in Euseb. \emph{Dem. Evang.} V. 132.}. He comes with the bowels of a true guardian: for He comes to save, and to heal, to teach, to admonish, to strengthen, to exhort, to enlighten the mind, first of him who receives Him, and afterwards of others also, through him. And as a man, who being previously in darkness then suddenly beholds the sun, is enlightened in his bodily sight, and sees plainly things which he saw not, so likewise he to whom the Holy Ghost is vouchsafed, is enlightened in his soul, and sees things beyond man’s sight, which he knew not; his body is on earth, yet his soul mirrors forth the heavens. He sees, like Esaias, \emph{the Lord sitting upon a throne high and lifted up}\footnote{Is. vi. 1.}; he sees, like Ezekiel, \emph{Him who is above the Cherubim}\footnote{Ezek. x. 1.}; he sees like Daniel, \emph{ten thousand times ten thousand, and thousands of thousands}\footnote{Dan. vii. 10.}; and the man, who is so little, beholds the beginning of the world, and knows the end of the world, and the times intervening, and the successions of kings,—things which he never learned: for the True Enlightener is present with him. The man is within the walls of a house; yet the power of his knowledge reaches far and wide, and he sees even what other men are doing.

17. Peter was not with Ananias and Sapphira when they sold their possessions, but he was present by the Spirit; \emph{Why}, he says, \emph{hath Satan filled thine heart to lie to the Holy Ghost}\footnote{Acts v. 3.}? There was no accuser; there was no witness; whence knew he what had happened? \emph{Whiles it remained was it not thine own? and after it was sold, was it not in thine own power? why hast thou conceived this thing in thine heart}\footnote{Ib. v. 4.}? The \emph{unlettered\footnote{Ib. iv. 13.}} Peter, through the grace of the Spirit, learnt what not even the wise men of the Greeks had known. Thou hast the like in the case also of Elisseus. For when he had freely healed the leprosy of Naaman, Gehazi received the reward, the reward of another’s achievement; and he took the money from Naaman, and bestowed it in a dark place. But the \emph{darkness is not} hidden from the Saints\footnote{Ps. cxxxix. 12.}. And when he came, Elisseus asked him; and like Peter, when he said, \emph{Tell me whether ye} \emph{sold the land for so much}\footnote{Acts v. 8.}? he also enquires, \emph{Whence comest thou, Gehazi}\footnote{2 Kings v. 25.}? Not in ignorance, but in sorrow ask I \emph{whence comest thou}? From darkness art thou come, and to darkness shalt thou go; thou hast sold the cure of the leper, and the leprosy is thy heritage. I, he says, have fulfilled the bidding of Him who said to me, \emph{Freely ye have received, freely give}\footnote{Matt. x. 8.}; but thou hast sold this grace; receive now the condition of the sale. But what says Elisseus to him? \emph{Went not mine heart with thee}? I was here shut in by the body, but the spirit which has been given me of God saw even the things afar off, and shewed me plainly what was doing elsewhere. Seest thou how the Holy Ghost not only rids of ignorance, but invests with knowledge? Seest thou how He enlightens men’s souls?

18. Esaias lived nearly a thousand years ago; and he beheld Zion \emph{as a booth}. The city was still standing, and beautified with public places, and robed in majesty; yet he says, \emph{Zion shall be ploughed as a field}\footnote{Micah iii. 12; ascribed by Cyril to Isaiah.}, foretelling what is now fulfilled in our days\footnote{Cf. Euseb. \emph{Dem. Evang}. vi. 13: “In our own time we have seen with our eyes the Sion of old renown being ploughed by Romans with yokes of oxen, and Jerusalem in a state of utter desolation as the oracle itself says, like a lodge in a garden of cucumbers. As Cyril at that time saw the Prophet’s prediction fulfilled, so we also to the present day see most plainly the fulfilment of the divine oracle, and Sion ploughed before our eyes: for except the Church of the Apostles, with the houses lying around it, and the house of Caiaphas and the cemeteries, all the remaining space of this hill, lying without the city, is under plough.” (Jerusalem Editor).}. And observe the exactness of the prophecy; for he said, \emph{And the daughter of Zion shall be left as a booth in a vineyard, as a lodge in a garden of cucumbers}\footnote{Isa. i. 8. \textgreek{ὀπωροφυλάκιον} is the hut of the watchman who guarded the crop when ripening for harvest. \textgreek{Σικυήλατον} is explained by Basil in his comment on the passage of Isaiah as “A place that produces quick-growing and perishable fruits.” This agrees with the etymological sense of the word as “a forcing-bed for cucumbers” (Hippocrates apud Fritzsche, “\emph{Der Brief des Jeremia},” \emph{v}. 70). On the form \textgreek{σικυηράτῳ}, see the notes on the Epistle of Jeremy in the Speaker’s Commentary.}. And now the place is filled with gardens of cucumbers. Seest thou how the Holy Spirit enlightens the saints? Be not therefore carried away to other things, by the force of a common term, but keep fast the exact meaning.

19. And if ever, while thou hast been sitting here, a thought concerning chastity or virginity has come into thy mind, it has been His teaching. Has not often a maiden, already at the bridal threshold\footnote{\textgreek{παστάδας}. On the meaning of \textgreek{παστάς} see the notes on Herodotus, II. 148, 169 in Bähr, and Rawlinson. Here it appears to mean the cloister or colonnade which gave access to the bridal chamber, \textgreek{θάλαμος}.}, fled away, He teaching her the doctrine of virginity? Has not often a man distinguished at court\footnote{\textgreek{ἐν παλατίοις}.}, scorned wealth and rank, under the teaching of the Holy Ghost? Has not often a young man, at the sight of beauty, closed his eyes, and fled from the sight, and escaped the defilement? Askest thou whence this has come to pass? The Holy Ghost taught the soul of the young man. Many ways of covetousness are there in the world; yet Christians refuse possessions: wherefore? because of the teaching of the Holy Ghost. Worthy of honour is in truth that Spirit, holy and good; and fittingly are we baptized into Father, Son, and Holy Ghost. A man, still clothed with a body, wrestles with many fiercest demons; and often the demon, whom many men could not master with iron bands, has been mastered by the man himself with words of prayer, through the power which is in him of the Holy Ghost; and the mere breathing of the Exorcist\footnote{Compare Procat. § 9; Cat. xx. 3.} becomes as fire to that unseen foe. A mighty ally and protector, therefore, have we from God; a great Teacher of the Church, a mighty Champion on our behalf. Let us not be afraid of the demons, nor of the devil; for mightier is He who fighteth for us. Only let us open to Him our doors; \emph{for He goeth about seeking such as are worthy}\footnote{Wisdom vi. 16. Compare the saying in Clem. Alex. \emph{Quis dives salvetur?} § 31: \textgreek{αὐτὸν ζητεῖν τοὺς εὖ πεισομένους ἀξίους τε ὄντας τοῦ Σωτῆρος μαθητάς}. The Jerusalem Editor quotes from Origen (\emph{Prolog. in Cantic.)} a passage which may have been known to Cyril: “This Comforter therefore goeth about seeking if He may discover any worthy and receptive souls to whom He may reveal the greatness of the love which is in God.”} and searching on whom He may confer His gifts.

20. And He is called the Comforter, because He comforts and encourages us, and \emph{helpeth our infirmities; for we know not what we should pray for as we ought; but the Spirit Himself maketh intercession for us, with groanings which cannot be uttered}\footnote{Rom. viii. 26.}, that is, makes intercession to God. Oftentimes a man for Christ’s sake has been outraged and dishonoured unjustly; martyrdom is at hand; tortures on every side, and fire, and sword, and savage beasts, and the pit. But the Holy Ghost softly whispers to him, “\emph{Wait thou on the Lord}\footnote{Ps. xxviii. 14; xxxvii. 34.}, O man; what is now befalling thee is a small matter, the reward will be great. Suffer a little while, and thou shalt be with Angels for ever. \emph{The sufferings of this present time art not worthy to be compared with the glory which shall be revealed in us}\footnote{Rom. viii. 18.}.” He portrays to the man the kingdom of heaven; He gives him a glimpse of the paradise of delight; and the martyrs, whose bodily countenances are of necessity turned to their judges, but who in spirit are already in Paradise, despise those hardships which are seen.

21. And wouldest thou be sure that by the power of the Holy Ghost the Martyrs bear their witness? The Saviour says to His disciples, \emph{And when they bring you unto} \emph{the synagogues, and the magistrates, and authorities, be not anxious how ye shall answer, or what ye shall say; for the Holy Ghost shall teach you in that very hour, what ye ought to say}\footnote{Luke xii. 11, 12.}. For it is impossible to testify as a martyr for Christ’s sake, except a man testify by the Holy Ghost; for if \emph{no man can say that Jesus Christ is the Lord, but by the Holy Ghost}\footnote{1 Cor. xii. 3. \textgreek{Μαρτυρῆσαι}, “to bear witness by death.”}, how shall any man give his own life for Jesus’ sake, but by the Holy Ghost?

22. Great indeed, and all-powerful in gifts, and wonderful, is the Holy Ghost. Consider, how many of you are now sitting here, how many souls of us are present. He is working suitably for each, and being present in the midst, beholds the temper of each, beholds also his reasoning and his conscience, and what we say, and think, and believe\footnote{Codd. Monac. Vind. Roe. Casaub. add \textgreek{καὶ τί πιστεύομεν}.}. Great indeed is what I have now said, and yet is it small. For consider, I pray, with mind enlightened by Him, how many Christians there are in all this diocese, and how many in the whole province\footnote{The terms \textgreek{παροικία}, the See of a Bishop, and \textgreek{ἐπαρχία}, the Province of a Metropolitan, were both adopted from the corresponding divisions of the Roman Empire. See Bingham, \emph{Antt}. Book IX. i. §§ 2–6.} of Palestine, and carry forward thy mind from this province, to the whole Roman Empire; and after this, consider the whole world; races of Persians, and nations of Indians, Garbs and Sarmatians, Gauls and Spaniards, and Moors, Libyans and Ethiopians, and the rest for whom we have no names; for of many of the nations not even the names have reached us. Consider, I pray, of each nation, Bishops, Presbyters, Deacons, Solitaries, Virgins, and laity besides; and then behold their great Protector, and the Dispenser of their gifts;—how throughout the world He gives to one chastity, to another perpetual virginity, to another almsgiving, to another voluntary poverty, to another power of repelling hostile spirits. And as the light, with one touch of its radiance sheds brightness on all things, so also the Holy Ghost enlightens those who have eyes; for if any from blindness is not vouchsafed His grace, let him not blame the Spirit, but his own unbelief.

23. Thou hast seen His power, which is in all the world; tarry now no longer upon earth, but ascend on high. Ascend, I say, in imagination even unto the first heaven, and behold there so many countless myriads of Angels. Mount up in thy thoughts, if thou canst, yet higher; consider, I pray thee, the Archangels, consider also the Spirits; consider the Virtues, consider the Principalities, consider the Powers, consider the Thrones, consider the Dominions\footnote{S. Basil (\emph{De Spiritu S.} c. xvi. § 38), after quoting the same passage, Col. i. 16, proceeds—\textgreek{εἴτε κυριότητες, καὶ εἴ τινές εἰσιν ἕτεραι λογικαὶ φύσεις ἁκατονόμαστοι}. The last word shews that Basil had in mind this passage of Cyril, who after the names of nations in § 22, adds \textgreek{καὶ τοὺς λοίπους ἀκατονομάστους ἡμῖν}.};—of all these the Comforter is the Ruler from God, and the Teacher, and the Sanctifier. Of Him Elias has need, and Elisseus, and Esaias, among men; of Him Michael and Gabriel have need among Angels. Naught of things created is equal in honour to Him: for the families of the Angels, and all their hosts assembled together, have no equality with the Holy Ghost. All these the all-excellent power of the Comforter overshadows. And they indeed are sent forth to minister\footnote{Heb. i. 14.}, but He searches even the deep things of God, according as the Apostle says, For the Spirit searcheth all things, yea, the deep things of God. For what man knoweth the thing of a man, save the spirit of the man which is in him? even so the things of God knoweth no man, but the Spirit of God\footnote{1 Cor. ii. 10, 11.}.

24. He preached concerning Christ in the Prophets; He wrought in the Apostles; He to this day seals the souls in Baptism. And the Father indeed gives to the Son; and the Son shares with the Holy Ghost. For it is Jesus Himself, not I, who says, \emph{All things are delivered unto Me of My Father}\footnote{Matt. xi. 27.}; and of the Holy Ghost He says, \emph{When He, the Spirit of Truth, shall come}, and the rest.…\emph{He shall glorify Me; for He shall receive of Mine, and shall shew it unto you}\footnote{John xvi. 13, 14.}. The Father through the Son, with the Holy Ghost, is the giver of all grace; the gifts of the Father are none other than those of the Son, and those of the Holy Ghost; for there is one Salvation, one Power, one Faith; One God, the Father; One Lord, His only-begotten Son; One Holy Ghost, the Comforter. And it is enough for us to know these things; but inquire not curiously into His nature or substance\footnote{In regard to the caution with which St. Cyril here speaks, we must remember that the heresy of Macedonius had not yet given occasion to the formal discussion and determination of the “nature and substance” of the Holy Ghost.}: for had it been written, we would have spoken of it; what is not written, let us not venture on; it is sufficient for our salvation to know, that there is Father, and Son, and Holy Ghost.

25. This Spirit descended upon the seventy Elders in the days of Moses. (Now let not the length of the discourse, beloved, produce weariness in you: but may He the very subject of our discourse grant strength to every one, both to us who speak, and to you who listen!) This Spirit, as I was saying, came down upon the seventy Elders in the time of Moses; and this I say to thee, that I may now prove, that He knoweth all things, and worketh \emph{as He will}\footnote{1 Cor. xii. 11.}. The seventy Elders were chosen; \emph{And the Lord came down in a cloud, and took of the Spirit that was upon Moses, and put it upon the seventy Elders}\footnote{Num. xi. 24, 25. “Modad” is the form of the name in the \textsc{LXX.}}; not that the Spirit was divided, but that His grace was distributed in proportion to the vessels, and the capacity of the recipients. Now there were present sixty and eight, and they prophesied; but Eldad and Modad were not present: therefore that it might be shewn that it was not Moses who bestowed the gift, but the Spirit who wrought, Eldad and Modad, who though called, had not as yet presented themselves, did also prophesy\footnote{The apocryphal book of Eldad and Modad is mentioned by Hermas, \emph{Shepherd}, Vis. ii. § 3. S. Basil, \emph{Liber de Spir. S.} cap. 61, referring to Num. xi. 26, says that the Spirit rested permanently only upon Eldad and Modad.}.

26. Jesus the Son of Nun, the successor of Moses, was amazed; and came to him and said, “Hast thou heard that Eldad and Modad are prophesying? They were called, and they came not; \emph{my lord Moses, forbid them}\footnote{Num. xi. 28.}.” “I cannot forbid them,” he says, “for this grace is from Heaven; nay, so far am I from forbidding them, that I myself am thankful for it. I think not, however, that thou hast said this in envy; \emph{art} thou \emph{jealous for my sake}, because that they prophesy, and thou prophesiest not yet? Wait for the proper season; \emph{and oh that all the Lord’s people may be prophets, whenever the Lord shall give His Spirit upon them}\footnote{Num. xi. 29.}!” saying this also prophetically, \emph{whenever the Lord shall give}; “For as yet then He has not given it; so thou hast it not yet.”—Had not then Abraham this, and Isaac, and Jacob, and Joseph? And they of old, had they it not? Nay, but the words, “\emph{whenever the Lord shall give}” evidently mean “give it upon all; as yet indeed the grace is partial, then it shall be given lavishly.” And he secretly alluded to what was to happen among us on the day of Pentecost; for He Himself came down among us. He had however also come down upon many before. For it is written, \emph{And Jesus the son of Nun was filled with a spirit of wisdom; for Moses had laid his hands upon him}\footnote{Deut. xxxiv. 9.}. Thou seest the figure everywhere the same in the Old and New Testament;—in the days of Moses, the Spirit was given by laying on of hands; and by laying on of hands Peter\footnote{Acts viii. 18. On this passage of Cyril, see the section on “\emph{Chrism}” in the Introduction.} also gives the Spirit. And on thee also, who art about to be baptized, shall His grace come; yet in what manner I say not, for I will not anticipate the proper season.

27. He also came down upon all righteous men and Prophets; Enos, I mean, and Enoch, and Noah, and the rest; upon Abraham, Isaac, and Jacob; for as regards Joseph, even Pharaoh perceived that he had \emph{the Spirit of God within him}\footnote{Gen. xli. 38.}. As to Moses, and the wonderful works wrought by the Spirit in his days, thou hast heard often: This Spirit Job also had, that most enduring man, and all the saints, though we repeat not all their names. He also was sent forth when the Tabernacle was in making, and filled with wisdom the wise-hearted men who were with Bezaleel\footnote{Ex. xxxi. 1–6; xxxvi. 1.}.

28. In the might of this Spirit, as we have it in the Book of Judges, Othniel judged\footnote{Judges iii. 10.}; Gideon\footnote{Ib. vi. 34.} waxed strong; Jephtha conquered\footnote{Ib. xi. 29.}; Deborah, a woman, waged war; and Samson, so long as he did righteously, and grieved Him not, wrought deeds above man’s power. And as for Samuel and David, we have it plainly in the Books of the Kingdoms, how by the Holy Ghost they prophesied themselves, and were rulers of the prophets;—and Samuel was called \emph{the Seer}\footnote{1 Sam. ix. 9.}; and David says distinctly, \emph{The Spirit of the Lord spake by me}\footnote{2 Sam. xxiii. 2.}, and in the Psalms, \emph{And take not thy Holy Spirit from me}\footnote{Ps. li. 11.}, and again, \emph{Thy good Spirit shall lead me in the land of righteousness}\footnote{Ps. cxliii. 10.}. And as we have it in Chronicles, Azariah\footnote{2 Chron. xv. 1.}, in the time of King Asa, and Jahaziel\footnote{Ib. xx. 14.} in the time of King Jehoshaphat, partook of the Holy Ghost; and again, another Azariah, he who was stoned\footnote{Ib. xxiv. 20, 21.}. And Ezra says, \emph{Thou gavest also Thy good Spirit to instruct them}\footnote{Neh. ix. 20. Ezra and Nehemiah form one book “Ezra” in the Hebrew Canon.}. But as touching Elias who was taken up, and Elisseus, those inspired\footnote{\textgreek{πνευματοφόρων}, used only twice in the Sept. (Hosea ix. 7; Zeph. iii. 4), and in an unfavourable sense. With Cyril’s use of it compare Theophilus, \emph{Ad Autolyc}. ii. 9: \textgreek{Θεοῦ ἀνθρώπους πνευματοφόρους Πνεύματος ἁγίου}.} and wonder-working men, it is manifest, without our saying so, that they were full of the Holy Ghost.

29. And if further a man peruse all the books of the Prophets, both of the Twelve, and of the others, he will find many testimonies concerning. the Holy Ghost; as when Micah says in the person of God, \emph{surely I will perfect power by the Spirit the Lord}\footnote{Mic. iii. 8.}; and Joel cries, \emph{And it shall come to pass afterwards}, saith God, \emph{that I will pour out My Spirit upon all flesh}\footnote{Joel ii. 28.}, and the rest; and Haggai, \emph{Because I am with you, saith the Lord of Hosts}\footnote{Haggai ii. 4.}; and \emph{My Spirit remaineth in the midst of you}\footnote{Ib. v. 5.}; and in like manner Zechariah, \emph{But receive My words and My statutes which I command by My Spirit, to My servants the Prophets}\footnote{Zech. i. 6.}; and other passages.

30. Esaias too, with his majestic voice, says, \emph{And the Spirit of God shall rest upon Him,} \emph{the spirit of wisdom and understanding, the spirit of counsel and might, the spirit of knowledge and godliness; and the Spirit of the fear of God shall fill Him}\footnote{Is. xi. 2.}; signifying that the Spirit is one and undivided, but His operations various. So again, \emph{Jacob My servant},…..\emph{I have put My Spirit upon Him}\footnote{Ib. xliv. 1; xlii. 1.}. And again, \emph{I will pour My Spirit upon thy seed}\footnote{Ib. xliv. 3.}; and again, \emph{And now the Lord Almighty and His Spirit hath sent Me}\footnote{Ib. xlviii. 16.}; and again, \emph{This is My covenant with them, saith the Lord, My Spirit which is upon thee}\footnote{Is. lix. 21.}; and again, \emph{The Spirit of the Lord is upon Me, because He hath anointed Me}\footnote{Is. lxi. 1.}, and the rest; and again in his charge against the Jews, \emph{But they rebelled and vexed His Holy Spirit}\footnote{Ib. lxiii. 10.}, and \emph{Where is He that put His Holy Spirit within them}\footnote{v. 11.}? Also thou hast in Ezekiel (if thou be not now weary of listening), what has already been quoted, \emph{And the Spirit fell upon me, and said unto me, Speak; Thus saith the Lord}\footnote{Ezek. xi. 5.}. But the words, \emph{fell upon me} we must understand in a good sense, that is “lovingly;” and as Jacob, when he had found Joseph, fell upon his neck; as also in the Gospels, the loving father, on seeing his son who had returned from his wandering, \emph{had compassion, and ran and fell on his neck, and kissed him}\footnote{Gen. xlvi. 29; Luke xv. 20.}. And again in Ezekiel, \emph{And he brought me in a vision by the Spirit of God into Chaldæa, to them of the captivity}\footnote{Ezek. xi. 24.}. And other texts thou heardest before, in what was said about Baptism; \emph{Then will I sprinkle clean water upon you}\footnote{Ib. xxxvi. 25; Cat. iii. 16.}, and the rest; \emph{a new heart also will I give you, and a new spirit will I put within you}\footnote{Ib. v. 26.}; and then immediately, \emph{And I will put My Spirit within you}\footnote{Ib. v. 27.}. And again, \emph{The hand of the Lord was upon me, and carried me out in the Spirit of the Lord}\footnote{Ezek. xxxvii. 1.}.

31. He endued with wisdom the soul of Daniel, that young as he was he should become a judge of Elders. The chaste Susanna was condemned as a wanton; there was none to plead her cause; for who was to deliver her from the rulers? She was led away to death, she was now in the hands of the executioners. But her Helper was at hand, the Comforter, the Spirit who sanctifies every rational nature. Come hither to me, He says to Daniel; young though thou be, convict old men infected with the sins of youth; for it is written, \emph{God raised up the Holy Spirit upon a young stripling}\footnote{Susanna, \emph{v}. 45.}; and nevertheless, (to pass on quickly,) by the sentence of Daniel that chaste lady was saved. We bring this forward as a testimony; for this is not the season for expounding. Nebuchadnezzar also knew that the Holy Spirit was in Daniel; for he says to him, \emph{O Belteshazzar, master of the magicians, of whom I know, that the Holy Spirit of God is in thee}\footnote{Dan. iv. 9.}. One thing he said truly, and one falsely; for that he had the Holy Spirit was true, but he was not the \emph{master of the magicians}, for he was no magician, but was wise through the Holy Ghost. And before this also, he interpreted to him the vision of the Image, which he who had seen it himself knew not; for he says, Tell me the vision, which I who saw it know not\footnote{Ib. ii. 26, 31.}. Thou seest the power of the Holy Ghost; that which they who saw it, know not, they who saw it not, know and interpret.

32. And indeed it were easy to collect very many texts out of the Old Testament, and to discourse more largely concerning the Holy Ghost. But the time is short; and we must be careful of the proper length of the lecture. Wherefore, being for the present content awhile with passages from the Old Testament, we will, if it be God’s pleasure, proceed in the next Lecture to the remaining texts out of the New Testament. And may the God of peace, through our Lord Jesus Christ, and through the love of the Spirit, count all of you worthy of His spiritual and heavenly gifts:—To whom be glory and power for ever and ever. Amen.

\xchapter{Lecture XVII. Continuation of the Discourse on the Holy Ghost.}

\textsc{1 Corinthians xii. 8}

For to one is given by the Spirit the word of wisdom, \&c.

1. \textsc{In} the preceding Lecture, according to our ability we set before you, our beloved hearers\footnote{\textgreek{ταῖς τῆς ὑμετέρας ἀγάπης ἀκοαῖς}. Compare § 30, below: \textgreek{συγγώμην αἰτῶ παρὰ τῆς ὑμετέρας ἀγάπης}. Ignat. \emph{Philadelph}. c. iv. (Long recension): \textgreek{θαρρῶν γράφω τῇ ἀξιοθέῳ ἀγάπη ὑμῶν}. “Caritas” is constantly used in the same manner.}, some small portion of the testimonies concerning the Holy Ghost; and on the present occasion, we will, if it be God’s pleasure, proceed to treat, as far as may be, of those which remain out of the New Testament: and as then to keep within due limit of your attention we restrained our eagerness (for there is no satiety in discoursing concerning the Holy Ghost), so now again we must say but a small part of what remains. For now, as well as then, we candidly own that our weakness is overwhelmed by the multitude of things written. Neither to-day will we use the subtleties of men, for that is unprofitable; but merely call to mind what comes from the divine Scriptures; for this is the safest course, according to the blessed Apostle Paul, who says, \emph{Which things also we speak, not in words which man’s wisdom teacheth, but which the Holy Ghost teacheth, comparing spiritual things with spiritual}\footnote{1 Cor. ii. 13.}. Thus we act like travellers or voyagers, who having one goal to a very long journey, though hastening on with eagerness, yet by reason of human weakness are wont to touch in their way at divers cities or harbours.

2. Therefore though our discourses concerning the Holy Ghost are divided, yet He Himself is undivided, being one and the same. For as in speaking concerning the Father, at one time we taught how He is the one only Cause\footnote{Cat. vi.}; and at another, how He is called Father\footnote{Ib. vii.}, or Almighty\footnote{Ib. viii.}; and at another, how He is the Creator\footnote{Ib. ix.} of the universe; and yet the division of the Lectures made no division of the Faith, in that He, the Object of devotion, both was and is One;—and again, as in discoursing concerning the Only-begotten Son of God we taught at one time concerning His Godhead\footnote{Cat. x. xi.}, and at another concerning His Manhood\footnote{Ib. xii. xv.}, dividing into many discourses the doctrines concerning our Lord Jesus Christ, yet preaching undivided faith towards Him;—so now also though the Lectures concerning the Holy Spirit are divided, yet we preach faith undivided towards Him. For it is one and the Self-same Spirit who \emph{divides} His gifts \emph{to every man severally as He will}\footnote{1 Cor. xii. 11.}, Himself the while remaining undivided. For the Comforter is not different from the Holy Ghost, but one and the self-same, called by various names; who lives and subsists, and speaks, and works; and of all rational natures made by God through Christ, both of Angels and of men, He is the Sanctifier\footnote{Compare Basil. \emph{de Sp. Sancto}, c. 38: “By the Father’s will the ministering spirits subsist, and by the operation of the Son they are brought into existence, and by the presence of the Holy Ghost are perfected: and the perfection of Angels is sanctification and continuance therein.”}.

3. But lest any from lack of learning, should suppose from the different titles of the Holy Ghost that these are divers spirits, and not one and the self-same, which alone there is, therefore the Catholic Church guarding thee beforehand hath delivered to thee in the profession of the faith, that thou “\textsc{believe in one Holy Ghost the Comforter, who spake by the Prophets};” that thou mightest know, that though His names be many, the Holy Spirit is but one;—of which names, we will now rehearse to you a few out of many.

4. He is called the Spirit, according to the Scripture just now read, \emph{For to one is given by the Spirit the word of wisdom}\footnote{1 Cor. xii. 8.}. He is called the Spirit of Truth, as the Saviour says, \emph{When He, the Spirit of Truth, is come}\footnote{John xvi. 13.}. He is called also the Comforter, as He said, \emph{For if I go not away, the Comforter will not come unto you}\footnote{Ib. v. 7.}. But that He is one and the same, though called by different titles, is shewn plainly from the following. For that the Holy Spirit and the Comforter are the same, is declared in those words, \emph{But the Comforter, which is the Holy Ghost}\footnote{John xiv. 26.}; and that the Comforter is the same as the Spirit of Truth, is declared, when it is said, \emph{And I will give you another Comforter, that He may abide with you for ever, even the Spirit of Truth}\footnote{Ib. vv. 16, 17.}; and again, \emph{But when the Comforter is come whom I will send unto you from the Father, even the Spirit of Truth}\footnote{Ib. xv. 26.}. And He is called the Spirit of God, according as it is written, \emph{And I saw the Spirit of God descending}\footnote{John i. 32.}; and again, \emph{For as many as are led by the Spirit of God, they are the sons of God}\footnote{Rom. viii. 14.}. He is called also the Spirit of the Father, as the Saviour says, \emph{For it is not ye that speak, but the Spirit of your Father which speaketh in you}\footnote{Matt. x. 20.}; and again Paul saith, \emph{For this cause I bow my knees unto the Father}, and the rest;…\emph{that He would grant you to be strengthened by His Spirit}\footnote{Eph. iii. 14–16.}. He is also called the Spirit of the Lord, according to that which Peter spake, \emph{Why is it that ye have agreed together to tempt the Spirit of the Lord}\footnote{Acts v. 9.}? He is called also the Spirit of God and Christ, as Paul writes, \emph{But ye are not in the flesh, but in the Spirit, if so be that the Spirit of God dwell in you. But if any man have not the Spirit of Christ, he is none of His}\footnote{Rom. viii. 9.}. He is called also the Spirit of the Son of God, as it is said, \emph{And because ye are sons, God hath sent forth the Spirit of His Son}\footnote{Gal. iv. 6.}. He is called also the Spirit of Christ, as it is written, \emph{Searching what or what manner of time the Spirit of Christ which was in them did signify}\footnote{1 Pet. i. 11.}; and again, \emph{Through your prayer, and the supply of the Spirit of Jesus Christ}\footnote{Phil. i. 19.}.

5. Thou wilt find many other titles of the Holy Ghost besides. Thus He is called the Spirit of Holiness, as it is written, \emph{According to the Spirit of Holiness}\footnote{Rom. i. 4.}. He is also called the Spirit of adoption, as Paul saith, \emph{For ye received not the spirit of bondage again unto fear, but ye received the Spirit of adoption, whereby we cry, Abba, Father}\footnote{Ib. viii. 15.}. He is also called the Spirit of revelation, as it is written, \emph{May give you the Spirit of wisdom and revelation in the knowledge of Him}\footnote{Eph. i. 17.}. He is also called the Spirit of promise, as the same Paul says, \emph{In whom ye also after that ye believed, were sealed with the Holy Spirit of promise}\footnote{Ib. v. 13.}. He is also called the Spirit of grace, as when he says again, \emph{And hath done despite to the Spirit of grace}\footnote{Heb. x. 29.}. And by many other such-like titles is He named. And thou heardest plainly in the foregoing Lecture, that in the Psalms He is called at one time \emph{the good Spirit}\footnote{Cat. xvi. 28; Ps. cxliii. 10.}, and at another the \emph{princely Spirit}\footnote{\textgreek{ἡγεμονικῷ}, Sept. Ps. li. 12: \textsc{R.V.} \emph{Uphold me with a free spirit}.}; and in Esaias He was styled \emph{the Spirit of wisdom and understanding, of counsel, and might, of knowledge, and of godliness, and of the fear of God}\footnote{Is. xi. 2.}. By all which Scriptures both those before and those now alleged, it is established, that though the titles of the Holy Ghost be different, He is one and the same; living and subsisting, and always present together with the Father and the Son\footnote{Origen, in the Catena on St. John iii. 8: “This also shews that the Spirit is a Being (\textgreek{οὐσίαν}): for He is not, as some suppose, an energy of God, having according to them no individuality of subsistence. And the Apostle also, after enumerating the gifts of the Spirit, immediately added, \emph{But all these worketh the one and the same Spirit, dividing to each one severally as He will.} Now if He willeth and worketh and divideth, He is surely an energizing Being, but not an energy” (Suicer, \emph{Thesaurus}, \textgreek{Πνευμα}).}; not uttered or breathed from the mouth and lips of the Father or the Son, nor dispersed into the air, but having a real substance\footnote{\textgreek{ἐνυπόστατον}. Cf. Cat. xi. 10; xvi. 13, note 5.}, Himself speaking, and working, and dispensing, and sanctifying; even as the Economy of salvation which is to usward from the Father and the Son and the Holy Ghost, is inseparable and harmonious and one, as we have also said before. For I wish you to keep in mind those things which were lately spoken, and to know clearly that there is not one Spirit in the Law and the Prophets, and another in the Gospels and Apostles; but that it is One and the Self-same Holy Spirit, which both in the Old and in the New Testament, spoke the divine Scriptures\footnote{Cat. iv. 16; xvi. 4.}.

6. This is the Holy Ghost, who came upon the Holy Virgin Mary; for since He who was conceived was Christ the Only-begotten, the \emph{power of the Highest overshadowed her}, and the \emph{Holy Ghost came upon her}\footnote{Luke i. 35.}, and sanctified her, that she might be able to receive Him, \emph{by whom all things were made}\footnote{John i. 3.}. But I have no need of many words to teach thee that generation was without defilement or taint, for thou hast learned it. It is Gabriel who says to her, I am the herald of what shall be done, but have no part in the work. Though an Archangel, I know my place; and though I joyfully bid thee All hail, yet how thou shalt bring forth, is not of any grace of mine. \emph{The Holy Ghost shall come upon thee, and the power of the Highest shall overshadow thee; therefore also that Holy Thing which shall be born of thee shall be called the Son of God}\footnote{Luke i. 35.}.

7. This Holy Spirit wrought in Elisabeth; for He recognises not virgins only, but matrons also, so that their marriage be lawful. \emph{And} \emph{Elisabeth was filled with the Holy Ghost}\footnote{Luke i. 41.}, and prophesied; and that noble hand-maiden says of her own Lord, \emph{And whence is this to me, that the Mother of my Lord should come to me}\footnote{Ib. v. 43.}? For Elisabeth counted herself blessed. Filled with this Holy Spirit, Zacharias also, the father of John, prophesied\footnote{Ib. v. 67.}, telling how many good things the Only-begotten should procure, and that John should be His harbinger\footnote{Ib. v. 76.} through baptism. By this Holy Ghost also it was revealed to just Symeon, \emph{that he should not see death, till he had seen the Lord’s Christ}\footnote{Luke ii. 26–35.}; and he received Him in his arms, and bore clear testimony in the Temple concerning Him.

8. And John also, who had been filled with the Holy Ghost from his mother’s womb\footnote{Cat. iii. 6.}, was for this cause sanctified, that he might baptize the Lord; not giving the Spirit himself, but preaching glad tidings of Him who gives the Spirit. For he says, \emph{I indeed baptize you with water unto repentance, but He that cometh after me}, and the rest; \emph{He shall baptize you with the Holy Ghost and with fire}\footnote{Matt. iii. 11.}. But wherefore with fire? Because the descent of the Holy Ghost was in fiery tongues; concerning which the Lord says joyfully, \emph{I am come to send fire on the earth; and what will I, if it be already kindled}\footnote{Luke xii. 49.}?

9. This Holy Ghost came down when the Lord was baptized, that the dignity of Him who was baptized might not be hidden; as John says, \emph{But He which sent me to baptize with water, the same said unto me, Upon whomsoever thou shalt see the Spirit descending and remaining upon Him, the same is He which baptizeth with the Holy Ghost}\footnote{John i. 33.}. But see what saith the Gospel; \emph{the heavens were opened}; they were opened because of the dignity of Him who descended; for, \emph{lo}, he says, \emph{the heavens were opened, and he saw the Spirit of God descending as a dove, and lighting upon Him}\footnote{Matt. iii. 16.}: that is, with voluntary motion in His descent. For it was fit, as some have interpreted, that the primacy and first-fruits\footnote{\textgreek{τὰς ἀπαρχὰς καὶ τὰ πρωτεῖα}. The order is inverted in the translation. Cf. Hermas, \emph{Sim}. viii. 7 \textgreek{ἔχοντες ζῆλόν τινα ἐν ἀλλήλοις περὶ πρωτείων}.} of the Holy Spirit promised to the baptized should be conferred upon the manhood of the Saviour, who is the giver of such grace. But perhaps He came down in the form of a dove, as some say, to exhibit a figure of that dove who is pure and innocent and undefiled, and also helps the prayers for the children she has begotten, and for forgiveness of sins\footnote{The Benedictine Editor adds the two last words \textgreek{τύπον παραδηλοῦν} from \textsc{mss.} Roe. Casaub. as necessary to the construction, and adds the following note. “The text thus emended is capable of two senses. The first, that the Holy Spirit came down in the form of a dove, a pure and harmless bird, to shew that He is Himself as it were a mystic dove in His simplicity and love of children, for whose new birth and remission of sins at Baptism He unites His prayers with Christ’s, as Cyril teaches in Cat. xvi. 20: and that Christ was for the like cause mystically foreshown in Canticles as having eyes like a dove’s. The other sense is, that the Spirit descended in the form of a dove on Christ’s Humanity in order to shew this to be as it were a dove in innocence, holiness, love of children, and concurrence with the Holy Spirit in their regeneration.…Either sense is admissible, and maintained by many of the Fathers: but I prefer the former.” This interpretation is confirmed by Tertullian (\emph{de Baptismo}, c. viii.), who says that the Holy Spirit glided down on the Lord “in the shape of a dove” in order that the nature of the Holy Spirit might be declared by means of a creature of simplicity and innocence.”}; even as it was emblematically foretold that Christ should be thus manifested in the appearance of His eyes; for in the Canticles she cries concerning the Bridegroom, and says, \emph{Thine eyes are as doves by the rivers of water}\footnote{Cant. v. 12. \textgreek{ἐπὶ πληρώματα ὑδάτων} (Sept.). The usual meaning of \textgreek{ὀφθαλμοφανῶς} is “manifestly to the eyes,” Esther viii. 13.}.

10. Of this dove, the dove of Noe, according to some, was in part a figure\footnote{Tertullian, \emph{ibid.} “Just as after the waters of the deluge, by which the old iniquity was purged—after the baptism, so to say, of the world—a dove was the herald which announced to the earth the assuagement of celestial wrath,.…so to our flesh, as it emerges from the font after its old sins, flies the dove of the Holy Spirit, bringing us the peace of God, sent out from heaven where the Church is, the typified ark.” Compare also Hippolytus, \emph{The Holy Theophany}, §§ 8, 9, a treatise with which Cyril has much in common.}. For as in his time by means of wood and of water there came salvation to themselves, and the beginning of a new generation, and the dove returned to him towards evening with an olive branch; thus, say they, the Holy Ghost also descended upon the true Noe, the Author of the second birth, who draws together into one the wills of all nations, of whom the various dispositions of the animals in the ark were a figure:—Him at whose coming the spiritual wolves feed with the lambs, in whose Church the calf, and the lion, and the ox, feed in the same pasture, as we behold to this day the rulers of the world guided and taught by Churchmen. The spiritual dove therefore, as some interpret, came down at the season of His baptism, that He might shew that it is He who by the wood of the Cross saves them who believe, He who at eventide should grant salvation through His death.

11. And these things perhaps should be otherwise explained; but now again we must hear the words of the Saviour Himself concerning the Holy Ghost. For He says, \emph{Except a man be born of water and of the Spirit, he cannot enter into the kingdom of God}\footnote{John iii. 5.}. And that this grace is from the Father, He thus states, \emph{How much more shall your heavenly Father give the Holy Spirit to them that ask him}\footnote{Luke xi. 13.}. And that we ought to worship God in the Spirit, He shews thus, \emph{But the hour cometh and now is, when the true worshippers shall worship the Father in Spirit and in truth; for the Father also seeketh such to worship Him. God is a Spirit; and they that worship Him must worship Him in spirit and in truth}\footnote{John iv. 23.}. And again, \emph{But if I by the} \emph{Spirit of God cast out devils}\footnote{Matt. xii. 28.}; and immediately afterwards, \emph{Therefore I say unto you, All manner of sin and blasphemy shall be forgiven unto men; but the blasphemy against the Holy Ghost shall not be forgiven. And whosoever shall speak a word against the Son of man, it shall be forgiven him; but whosoever shall speak a word against the Holy Ghost, it shall not be forgiven him, neither in this world, neither in the world to come}\footnote{Ib. v. 31.}. And again He says, \emph{And I will pray the Father, and He shall give you another Comforter, that He may be with you for ever, the Spirit of Truth; whom the world cannot receive, because it seeth Him not, neither knoweth Him; but ye know Him, for He abideth with you, and shall be in you}\footnote{John xiv. 16.}. And again He says, \emph{These things have I spoken unto you being yet present with you. But the Comforter, which is the Holy Ghost, whom the Father will send in My name, He shall teach you all things, and bring to your remembrance all things that I said unto you}\footnote{Ib. v. 25.}. And again He says, \emph{But when the Comforter is come, whom I will send unto you from the Father, even the Spirit of Truth, which proceedeth from the Father, He shall testify of Me}\footnote{Ib. xv. 26.}. And again the Saviour says, \emph{For if I go not away, the Comforter will not come unto you}\footnote{Ib. xvi. 7.}…..\emph{And when He is come, He will convince the world of sin, of righteousness, and of judgment}\footnote{Ib. v. 8.}; and afterwards again, \emph{I have yet many things to say unto you, but ye cannot bear them now. Howbeit, when He the Spirit of Truth is come, He will declare unto you all the truth; for He shall not speak from Himself; but whatsoever He shall hear that shall He speak, and He shall announce unto you the things to come. He shall glorify Me, for He shall take of Mine, and shall announce it unto you. All things that the Father hath are mine; therefore said I, That He shall take of Mine, and shall announce it unto you}\footnote{Ib. v. 12–15.}. I have read to thee now the utterances of the Only-begotten Himself, that thou mayest not give heed to men’s words.

12. The fellowship of this Holy Spirit He bestowed on the Apostles; for it is written, \emph{And when He had said this, He breathed on them, and saith unto them, Receive ye the Holy Ghost: whose soever sins ye remit, they are remitted unto them; and whose soever sins ye retain, they are retained}\footnote{John xx. 22.}. This was the second time He breathed on man (His first breath\footnote{Gen. ii. 7: \emph{and breathed into his nostrils the breath of life.} Compare Cat. xiv. 10.} having been stifled through wilful sins); that the Scripture might be fulfilled, \emph{He went up breathing upon thy face, and delivering thee from affliction}\footnote{Nahum ii. 1. The Septuagint, followed by Cyril, differs widely from the Hebrew: (\textsc{R.V.}) \emph{He that dasheth in pieces is come up before thy face.}}. But whence went He up? From Hades; for thus the Gospel relates, that then after His resurrection He breathed on them. But though He bestowed His grace then, He was to lavish it yet more bountifully; and He says to them, “I am ready to give it even now, but the vessel cannot yet hold it; for a while therefore receive ye as much grace as ye can bear; and look forward for yet more; \emph{but tarry ye in the city of Jerusalem, until ye be clothed with power from on high}\footnote{Luke xxiv. 39.}. \emph{Receive} it in part now; then, ye shall wear it in its fulness. For he who receives, often possesses the gift but in part; but he who is clothed, is completely enfolded by his robe. “Fear not,” He says, “the weapons and darts of the devil; for ye shall bear with you the power of the Holy Ghost.” But remember what was lately said, that the Holy Ghost is not divided, but only the grace which is given by Him.

13. Jesus therefore went up into heaven, and fulfilled the promise. For He said to them, \emph{I will pray the Father, and He shall give you another Comforter}\footnote{John xiv. 16.}. So they were sitting, looking for the coming of the Holy Ghost; \emph{and when the day of Pentecost was fully come}, here, in this city of Jerusalem,—(for this honour also belongs to us\footnote{Cat. iii. 7; xvi. 5. Bp. Pearson (\emph{Lectiones in Acta Apost}. I. § 18): “Rightly said Cyril, Bishop of Jerusalem, ‘All prerogatives are with us.’ And the Emperor Justin called her ‘Mother of the Christian name.’ Jerome also (\emph{Ep}. 17, 3), said: ‘The whole mystery of our Faith is native of that province and city.‘”}; and we speak not of the good things which have happened among others, but of those which have been vouchsafed among ourselves,)—on the day of Pentecost, I say, they were sitting, and the Comforter came down from heaven, the Guardian and Sanctifier of the Church, the Ruler of souls, the Pilot of the tempest-tossed, who leads the wanderers to the light, and presides over the combatants, and crowns the victors.

14. But He came down to clothe the Apostles with power, and to baptize them; for the Lord says, \emph{ye shall be baptized with the Holy Ghost not many days hence}\footnote{Acts i. 5.}. This grace was not in part, but His power was in full perfection; for as he who plunges into the waters and is baptized is encompassed on all sides by the waters, so were they also baptized completely by the Holy Ghost. The water however flows round the outside only, but the Spirit baptizes also the soul within, and that completely. And wherefore wonderest thou? Take an example from matter; poor indeed and common, yet useful for the simpler sort. If the fire passing in through the mass of the iron makes the whole of it fire, so that what was cold becomes burning and what was black is made bright,—if fire which is a body thus penetrates and works without hindrance in iron which is also a body, why wonder that the Holy Ghost enters into the very inmost recesses of the soul?

15. And lest men should be ignorant of the greatness of the mighty gift coming down to them, there sounded as it were a heavenly trumpet, For \emph{suddenly there came from heaven a sound as of the rushing of a mighty wind}\footnote{Acts ii. 2.}, signifying the presence of Him who was to grant power unto men to seize with violence the kingdom of God; that both their eyes might see the fiery tongues, and their ears hear the sound. \emph{And it filled all the house where they were sitting}; for the house became the vessel of the spiritual water; as the disciples sat within, the whole house was filled. Thus they were entirely baptized according to the promise, and invested soul and body with a divine garment of salvation. \emph{And there appeared unto them cloven tongues like as of fire, and it sat upon each of them; and they were all filled with the Holy Ghost}. They partook of fire, not of burning but of saving fire; of fire which consumes the thorns of sins, but gives lustre to the soul. This is now coming upon you also, and that to strip away and consume your sins which are like thorns, and to brighten yet more that precious possession of your souls, and to give you grace; for He gave it then to the Apostles. And He sat upon them in the form of fiery tongues, that they might crown themselves with new and spiritual diadems by fiery tongues upon their heads. A fiery sword barred of old the gates of Paradise; a fiery tongue which brought salvation restored the gift.

16. \emph{And they began to speak with other tongues as the Spirit gave them utterance}\footnote{Ib. v. 4.}. The Galilean Peter or Andrew spoke Persian or Median. John and the rest of the Apostles spoke every tongue to those of Gentile extraction; for not in our time have multitudes of strangers first begun to assemble here from all quarters, but they have done so since that time. What teacher can be found so great as to teach men all at once things which they have not learned? So many years are they in learning by grammar and other arts to speak only Greek well; nor yet do all speak this equally well; the Rhetorician perhaps succeeds in speaking well, and the Grammarian sometimes not well, and the skilful Grammarian is ignorant of the subjects of philosophy. But the Holy Spirit taught them many languages at once, languages which in all their life they never knew. This is in truth vast wisdom, this is power divine. What a contrast of their long ignorance in time past to their sudden, complete and varied and unaccustomed exercise of these languages!

17. The multitude of the hearers was confounded;—it was a second confusion, in the room of that first evil one at Babylon. For in that confusion of tongues there was division of purpose, because their thought was at enmity with God; but here minds were restored and united, because the object of interest was godly. The means of falling were the means of recovery. Wherefore they marvelled, saying\footnote{Acts ii. 8.}, \emph{How hear we them speaking}? No marvel if ye be ignorant; for even Nicodemus was ignorant of the coming of the Spirit, and to him it was said, \emph{The Spirit breatheth where it listeth, and thou hearest the voice thereof, but canst not tell whence it cometh, and whither it goeth}\footnote{John iii. 8: (R.V.) \emph{The wind bloweth}: (Marg.) Or, \emph{The Spirit breatheth}. It is impossible to preserve the double meaning in English.}; but if, even though I hear His voice, I know not whence he cometh, how can I explain, what He is Himself in substance?

18. \emph{But others mocking said, They are full of new wine}\footnote{Acts ii. 13.}, and they spoke truly though in mockery. For in truth the wine was new, even the grace of the New Testament; but this new wine was from a spiritual Vine, which had oftentimes ere this borne fruit in Prophets, and had budded in the New Testament. For as in things sensible, the vine ever remains the same, but bears new fruits in its seasons, so also the self-same Spirit continuing what He is, as He had often wrought in Prophets, now manifested a new and marvellous work. For though His grace had come before to the Fathers also, yet here it came exuberantly; for formerly men only partook of the Holy Ghost, but now they were baptized completely.

19. But Peter who had the Holy Ghost, and who knew what he possessed, says, “\emph{Men of Israel}, ye who preach Joel, but know not the things which are written, \emph{these men are not drunken as ye suppose}\footnote{Ib. v. 15.}. Drunken they are, not however as ye suppose, but according to that which is written, \emph{They shall be drunken with the fatness of thy house; and thou shalt make them drink of the torrents of thy pleasure}\footnote{Ps. xxxvi. 8.}. They are drunken, with a sober drunkenness, deadly to sin and life-giving to the heart, a drunkenness contrary to that of the body; for this last causes forgetfulness even of what was known, but that bestows the knowledge even of what was not known. They are drunken, for they have drunk the wine of the spiritual vine, which says, \emph{I am the vine and ye} \emph{are the branches}\footnote{John xv. 5.}. But if ye are not persuaded by me, understand what I tell you from the very time of the day; for \emph{it is the third hour of the day}\footnote{Acts ii. 25, and 15.}. For He who, as Mark relates, was crucified at the third hour, now at the third hour sent down His grace. For His grace is not other than the Spirit’s grace, but He who was then crucified, who also gave the promise, made good that which He promised. And if ye would receive a testimony also, \emph{Listen}, he says: “\emph{But this is that which was spoken by the prophet Joel; And it shall come to pass after this, saith God, I will pour forth of My Spirit}\footnote{Joel ii. 28.}—(and this word, \emph{I will pour forth}, implied a rich gift; \emph{for God giveth not the Spirit by measure, for the Father loveth the Son, and hath given all things into His hand}\footnote{John iii. 34, 35.}; and He has given Him the power also of bestowing the grace of the All-holy Spirit on whomsoever He will);—\emph{I will pour forth of My Spirit unto all flesh, and your sons and your daughters shall prophesy}; and afterwards, \emph{Yea, and on My servants and on My handmaidens I will pour out in those days of My Spirit, and they shall prophesy}\footnote{Joel ii. 29.}.” The Holy Ghost is no respecter of persons; for He seeks not dignities, but piety of soul. Let neither the rich be puffed up, nor the poor dejected, but only let each prepare himself for reception of the Heavenly gift.

20. We have said much to-day, and perchance you are weary of listening; yet more still remains. And in truth for the doctrine of the Holy Ghost there were need of a third lecture; and of many besides. But we must have your indulgence on both points. For as the Holy Festival of Easter is now at hand we have this day lengthened our discourse and yet we had not room to bring before you all the testimonies from the New Testament which we ought. For many passages are still to come from the Acts of the Apostles in which the grace of the Holy Ghost wrought mightily in Peter and in all the Apostles together; many also from the Catholic Epistles, and the fourteen Epistles of Paul; out of all which we will now endeavour to gather a few, like flowers from a large meadow, merely by way of remembrance.

21. For in the power of the Holy Ghost, by the will of Father and Son, Peter stood with the Eleven, and lifting up his voice, (according to the text, \emph{Lift up thy voice with strength, thou that bringest good tidings to Jerusalem}\footnote{Is. xl. 9.}), captured in the spiritual net of his words, \emph{about three thousand souls}. So great was the grace which wrought in all the Apostles together, that, out of the Jews, those crucifiers of Christ, this great number believed, and were baptized in the Name of Christ, and \emph{continued steadfastly in the Apostles’ doctrine and in the prayers}\footnote{Acts ii. 42.}. And again in the same power of the Holy Ghost, \emph{Peter and John went up into the Temple at the hour of prayer, which was the ninth hour}\footnote{Ib. iii. 1.}, and in the Name of Jesus healed the man at the Beautiful gate, who had been lame from his mother’s womb for forty years; that it might be fulfilled which was spoken, \emph{Then shall the lame man leap as an hart}\footnote{Is. xxxv. 6.}. And thus, as they captured in the spiritual net of their doctrine five thousand believers at once, so they confuted the misguided rulers of the people and chief priests, and that, not through their own wisdom, for \emph{they were unlearned and ignorant men}\footnote{Acts iv. 13.}, but through the mighty power of the Holy Ghost; for it is written, \emph{Then Peter filled with the Holy Ghost said to them}\footnote{Ib. v. 8.}. So great also was the grace of the Holy Ghost, which wrought by means of the Twelve Apostles in them who believed, that \emph{they were of one heart and of one soul}\footnote{Ib. v. 32.}, and their enjoyment of their goods was common, the possessors piously offering the prices of their possessions, and no one among them wanting aught; while Ananias and Sapphira, who attempted to lie to the Holy Ghost, underwent their befitting punishment.

22. \emph{And by the hands of the Apostles were many signs and wonders wrought among the people}\footnote{Acts v. 12.}. And so great was the spiritual grace shed around the Apostles, that gentle as they were, they were the objects of dread; for \emph{of the rest durst no man join himself to them; but the people magnified them; and multitudes were added of those who believed on the Lord, both of men and women}; and the streets were filled with the sick on their beds and couches, \emph{that as Peter passed by, at least his shadow might overshadow some of them}. And \emph{the multitude also of the cities round about came} unto this holy Jerusalem, \emph{bringing sick folk, and them that were vexed with unclean spirits, and they were healed every one} in this power of the Holy Ghost\footnote{Ib. vv. 13–16.}.

23. Again, after the Twelve Apostles had been cast into prison by the chief priests for preaching Christ, and had been marvellously delivered from it at night by an Angel, and were brought before them in the judgment hall from the Temple, they fearlessly rebuked them in their discourse to them concerning Christ, and added this, that \emph{God hath also given His Holy Spirit to them that obey Him}\footnote{Ib. v. 32.}. And when they had been scourged, they went their way rejoicing, and ceased not to \emph{teach and preach Jesus as the Christ}\footnote{Acts v. 42.}.

24. And it was not in the Twelve Apostles only that the grace of the Holy Spirit wrought, but also in the first-born children of this once barren Church, I mean the seven Deacons; for these also were chosen, as it is written, being \emph{full of the Holy Ghost and of wisdom}\footnote{Ib. vi. 3.}. Of whom Stephen, rightly so named\footnote{Ib. v. 8. \textgreek{Στέφανος}, “a crown.”}, the first fruits of the Martyrs, a man \emph{full of faith and of the Holy Ghost, wrought great wonders and miracles among the people}, and vanquished those who disputed with him; \emph{for they were not able to resist the wisdom and the Spirit by which he spake}\footnote{Ib. v. 10.}. But when he was maliciously accused and brought to the judgment hall, he was radiant with angelic brightness; for \emph{all they who sat in the council, looking steadfastly on him, saw his face, as it had been the face of an Angel}\footnote{Ib. v. 15.}. And having by his wise defence confuted the Jews, those \emph{stiffnecked men, uncircumcised in heart and ears, ever resisting the Holy Ghost}\footnote{Ib. vii. 51.}, he beheld \emph{the heavens opened}, and saw \emph{the Son of Man standing on the right hand of God}. He saw Him, not by his own power, but, as the Divine Scripture says, \emph{being full of the Holy Ghost, he looked up steadfastly into heaven, and saw the glory of God, and Jesus standing on the right hand of God}\footnote{Ib. v. 55.}.

25. In this power of the Holy Ghost, Philip also in the Name of Christ at one time in the city of Samaria drove away the unclean spirits, \emph{crying out with a loud voice}; and healed the palsied and the lame, and brought to Christ great multitudes of them that believe. To whom Peter and John came down, and with prayer, and the laying on of hands, imparted the fellowship of the Holy Ghost, from which Simon Magus alone was declared an alien, and that justly. And at another time Philip was called by the Angel of the Lord in the way, for the sake of that most godly Ethiopian, the Eunuch, and heard distinctly the Spirit Himself saying, \emph{Go near, and join thyself to this chariot}\footnote{Ib. viii. 5.}. He instructed the Eunuch, and baptized him, and so having sent into Ethiopia a herald of Christ, according as it is written, \emph{Ethiopia shall soon stretch out her hand unto God}\footnote{Ps. lxviii. 31.}, he was caught away by the Angel, and preached the Gospel in the cities in succession.

26. With this Holy Spirit Paul also had been filled after his calling by our Lord Jesus Christ. Let godly Ananias come as a witness to what we say, he who in Damascus said to him, \emph{The Lord, even Jesus who appeared to thee in the way which thou camest, hath sent me, that thou mayest receive thy sight, and be filled with the Holy Ghost}\footnote{Acts ix. 17.}. And straightway the Spirit’s mighty working changed the blindness of Paul’s eyes into newness of sight; and having vouchsafed His seal unto his soul, made him \emph{a chosen vessel} to \emph{bear the Name} of the Lord who had appeared to him, \emph{before kings and the children of Israel}, and rendered the former persecutor an ambassador and good servant,—one, who \emph{from Jerusalem, and even unto Illyricum, fully preached the Gospel}\footnote{Rom. xv. 19.}, and instructed even imperial Rome, and carried the earnestness of his preaching as far as Spain, undergoing conflicts innumerable, and performing signs and wonders. Of him for the present enough.

27. In the power of the same Holy Spirit Peter also, the chief of the Apostles and the bearer of the keys\footnote{\textgreek{κλειδοῦχος}. Cf. Matt. xvi. 19; Cat. ii. 19; xi. 3.} of the kingdom of heaven, healed Æneas the paralytic in the Name of Christ at Lydda, which is now Diospolis, and at Joppa raised from the dead Tabitha rich in good works. And being on the housetop in a trance, he saw heaven opened, and by means of the vessel let down as it were a sheet full of beasts of every shape and sort, he learnt plainly to call no man common or unclean, though he should be of the Greeks\footnote{Acts x. 11–16.}. And when he was sent for by Cornelius, he heard clearly the Holy Ghost Himself saying, \emph{Behold, men seek thee; but arise and get thee down, and go with them, nothing doubting; for I have sent them}\footnote{Ib. v. 19.}. And that it might be plainly shewn that those of the Gentiles also who believe are made partakers of the grace of the Holy Ghost, when Peter was come to Cesarea, and was teaching the things concerning Christ, the Scripture says concerning Cornelius and them who were with him; \emph{While Peter yet spake these words, the Holy Ghost fell on all them which heard the word; so that they of the circumcision also which came with Peter were astonished, and when they understood it said that on the Gentiles also was poured out the gift of the Holy Ghost}\footnote{Ib. v. 44.}.

28. And in Antioch also, a most renowned city of Syria, when the preaching of Christ took effect, Barnabas was sent hence as far as Antioch to help on the good work, being a \emph{good man, and full of the Holy Ghost, and of faith}\footnote{Ib. xi. 24.}; who seeing a great harvest of believers in Christ, brought Paul from Tarsus to Antioch, as his fellow-combatant. And when crowds had been instructed by them and assembled in the Church, \emph{it came to pass that the disciples were called Christians first in Antioch}\footnote{Ib. v. 26. Cf. Is. lxv. 15.}; the Holy Ghost, methinks, bestowing on the believers that new Name, which had been promised before by the Lord. And the grace of the Spirit being shed forth by God more abundantly in Antioch, there were there prophets and teachers of whom Agabus was one\footnote{Acts xi. 28.}. And \emph{as they ministered to the Lord and fasted, the Holy Ghost said, Separate Me Barnabas and Saul for the work whereunto I have called them}. And after hands had been laid on them, \emph{they were sent forth by the Holy Ghost}\footnote{Ib. xiii. 2–4.}. Now it is manifest, that the Spirit which speaks and sends, is a living Spirit, subsisting, and operating, as we have said.

29. This Holy Spirit, who in unison with Father and Son has established the New Covenant in the Church Catholic, has set us free from the burdens of the law grievous to be borne,—those I mean, concerning things common and unclean, and meats, and sabbaths, and new moons, and circumcision, and sprinklings, and sacrifices; which were given for a season, and \emph{had a shadow of the good things to come}\footnote{Heb. x. 1.}, but which, when the truth had come, were rightly withdrawn. For when Paul and Barnabas were sent to the Apostles, because of the question moved at Antioch by them who said that it was necessary to be circumcised and to keep the customs of Moses, the Apostles who were here at Jerusalem by a written injunction set free the whole world from all the legal and typical observances; yet they attributed not to themselves the full authority in so great a matter, but send an injunction in writing, and acknowledge this: \emph{For it hath seemed good unto the Holy Ghost and to us, to lay upon you no greater burden than these necessary things; that ye abstain from things sacrificed to idols, and from blood, and from things strangled, and from fornication}\footnote{Acts xv. 28, 29. \textgreek{ἐπιστολή} means a \emph{message} or \emph{injunction} whether verbal or written.}; shewing evidently by what they wrote, that though the writing was by the hands of human Apostles, yet the decree is universal from the Holy Ghost: which decree Paul and Barnabas took and confirmed unto all the world.

30. And now, having proceeded thus far in my discourse, I ask indulgence from your love\footnote{See note 1 on § 1, above.}, or rather from the Spirit who dwelt in Paul, if I should not be able to rehearse everything, by reason of my own weakness, and your weariness who listen. For when shall I in terms worthy of Himself declare the marvellous deeds wrought by the operation of the Holy Ghost in the Name of Christ? Those wrought in Cyprus upon Elymas the sorcerer, and in Lystra at the healing of the cripple, and in Cilicia and Phrygia and Galatia and Mysia and Macedonia? or those at Philippi (the preaching, I mean, and the driving out of the spirit of divination in the Name of Christ; and the salvation by baptism of the jailer with his whole house at night after the earthquake); or the events at Thessalonica; and the address at Areopagus in the midst of the Athenians; or the instructions at Corinth, and in all Achaia? How shall I worthily recount the mighty deeds which were wrought at Ephesus through Paul, by the Holy Ghost\footnote{Acts xix. 1–6.}? Whom they of that City knew not before, but came to know Him by the doctrine of Paul; and when Paul had laid his hands on them, and the Holy Ghost had come upon them, \emph{they spake with tongues, and prophesied}. And so great spiritual grace was upon him, that not only his touch wrought cures, but even the \emph{handkerchiefs and napkins}\footnote{Ib. v. 12.}, brought from his body, healed diseases, and scared away the evil spirits; and at last \emph{they also who practised curious arts brought their books together, and burned them before all men}\footnote{Ib. v. 19.}.

31. I pass by the work wrought at Troas on Eutychus, who \emph{being borne down by his sleep fell down from the third loft, and was taken up dead}; yet was saved alive by Paul\footnote{Ib. xx. 9–12.}. I also pass by the prophecies addressed to the Elders of Ephesus whom he called to him in Miletus, to whom he openly said, \emph{That the Holy Ghost witnesseth in every city, saying}\footnote{Ib. v. 23.}—and the rest; for by saying, \emph{in every city}, Paul made manifest that the marvellous works done by him in each city, were from the operative power of the Holy Ghost, by the will of God, and in the Name of Christ who spoke in him. By the power of this Holy Ghost, the same Paul was hastening to this holy city Jerusalem, and this, though Agabus by the Spirit foretold what should befall him; and yet he spoke to the people with confidence, declaring the things concerning Christ. And when brought to Cesarea, and set amid tribunals of justice, at one time before Felix, and at another before Festus the governor and King Agrippa, Paul obtained of the Holy Ghost grace so great, and triumphant in wisdom, that at last Agrippa himself the king of the Jews said, \emph{Almost thou persuadest me to be a Christian}\footnote{Ib. xxvi. 28. Cyril evidently understood \textgreek{ἐν ὀλίγῳ} to mean “\emph{almost}” (\textsc{A.V.}): but the more correct rendering is, “In brief thou wouldest persuade me to become a Christian.”}. This Holy Spirit granted to Paul, when he was in the island of Melita also, to receive no harm when bitten by the viper, and to effect divers cures on the diseased. This Holy Spirit guided him, the persecutor of old, as a herald of Christ, even as far as imperial Rome, and there he persuaded many of the Jews to believe in Christ, and to them who gainsaid he said plainly, \emph{Well spake the Holy Ghost by Esaias the Prophet, saying unto your fathers, and the rest}\footnote{Ib. xxviii. 25.}.

32. And that Paul was full of the Holy Ghost, and all his fellow Apostles, and they who after them believed in Father, Son, and Holy Ghost, hear from himself as he writes plainly in his Epistles; \emph{And my speech}, he says, \emph{and my preaching was not in persuasive words of man’s wisdom, but in demonstration of the Spirit and of power}\footnote{1 Cor. ii. 4.}. And again, \emph{But He who sealed us for this very purpose is God, who gave us the earnest of the Spirit}\footnote{2 Cor. i. 22.}. And again, \emph{He that raised up Jesus from the dead shall also quicken your mortal bodies by His Spirit which dwelleth in you}\footnote{Rom. viii. 11.}. And again, writing to Timothy, \emph{That good thing which was committed to thee guard through the Holy Ghost which was given to us}\footnote{2 Tim. i. 14: (R.V.) \emph{by the Holy Ghost which dwelleth in us}.}.

33. And that the Holy Ghost subsists, and lives, and speaks, and foretells, I have often said in what goes before, and Paul writes it plainly to Timothy: \emph{Now the Spirit speaketh expressly, that in later times some shall depart from the faith}\footnote{1 Tim. iv. 1.},—which we see in the divisions not only of former times but also of our own; so motley and diversified are the errors of the heretics. And again the same Paul says, \emph{Which in other generations was not made known unto the sons of men, as it hath now been revealed unto His Holy Apostles and Prophets in the Spirit}\footnote{Eph. iii. 5.}. And again, \emph{Wherefore, as saith the Holy Ghost}\footnote{Heb. iii. 7.}; and again, \emph{The Holy Ghost also witnesseth to us}\footnote{Ib. x. 15.}. And again he calls unto the soldiers of righteousness, saying, \emph{And take the helmet of salvation, and the sword of the Spirit, which is the Word of God, with all prayer and supplication}\footnote{Eph. vi. 17.}. And again, \emph{Be not drunk with wine, wherein is excess; but be filled with the Spirit, speaking to yourselves in psalms, and hymns, and spiritual songs}\footnote{Ib. v. 18, 19.}. And again, \emph{The grace of the Lord Jesus, and the love of God, and the communion of the Holy Ghost be with you all}\footnote{2 Cor. xiii. 14.}.

34. By all these proofs, and by more which have been passed over, is the personal, and sanctifying, and effectual power of the Holy Ghost established for those who can understand; for the time would fail me in my discourse if I wished to quote what yet remains concerning the Holy Ghost from the fourteen Epistles of Paul, wherein he has taught with such variety, completeness, and reverence. And to the power of the Holy Ghost Himself it must belong, to grant to us forgiveness for what we have omitted because the days are few, and upon you the hearers to impress more perfectly the knowledge of what yet remains; while from the frequent reading of the sacred Scriptures those of you who are diligent come to understand these things, and by this time, both from these present Lectures, and from what has before been told you, hold more steadfastly the Faith in “\textsc{One God the Father Almighty; and in our Lord Jesus Christ, His Only-Begotten Son; and in the Holy Ghost the Comforter}.” Though the word itself and title of Spirit is applied to Them in common in the sacred Scriptures,—for it is said of the Father, \emph{God is a Spirit}\footnote{John iv. 24.}, as it is written in the Gospel according to John; and of the Son, \emph{A Spirit before our face, Christ the Lord}\footnote{Lam. iv. 20. \emph{The breath of our nostrils, the anointed of the Lord}: referring to the captive king.}, as Jeremias the prophet says; and of the Holy Ghost, \emph{the Comforter, the Holy Ghost}\footnote{John xiv. 25.}, as was said;—yet the arrangement of articles in the Faith, if religiously understood, disproves the error of Sabellius also\footnote{The distinct mention in the Creed of three Persons excludes the error of Sabellius in confusing them. Cf. Cat. iv. 8; xvi. 14.}. Return we therefore in our discourse to the point which now presses and is profitable to you.

35. Beware lest ever like Simon thou come to the dispensers of Baptism in hypocrisy, thy heart the while not seeking the truth. It is ours to protest, but it is thine to secure thyself. If \emph{thou standest in faith}\footnote{Rom. xi. 20.}, blessed art thou; if thou hast fallen in unbelief, from this day forward cast away thine unbelief, and receive full assurance. For, at the season of baptism, when thou art come before the Bishops, or Presbyters, or Deacons\footnote{Cf. Bingham, \emph{Antiquities}, II. xx. 9. “When Cyril directs his Catechumens how they should behave themselves at the time of Baptism, when they came either before a bishop, or presbyter, or deacon, in city or village,—this may be presumed a fair intimation that then deacons were ordinarily allowed to minister Baptism in country places.” See further ‘Of the power granted anciently to deacons to baptize,’ Bingham, \emph{Lay Baptism}, I. i. 5.},—(for its grace is everywhere, in villages and in cities, on them of low as on them of high degree, on bondsmen and on freemen, for this grace is not of men, but the gift is from God through men,)—approach the Minister of Baptism, but approaching, think not of the face of him thou seest, but remember this Holy Ghost of whom we are now speaking. For He is present in readiness to seal thy soul, and He shall give thee that Seal at which evil spirits tremble, a heavenly and sacred seal, as also it is written, \emph{In whom also ye believed, and were sealed with the Holy Spirit of promise}\footnote{Eph. i. 13. Cf. Cat. i. 2, 3.}.

36. Yet He tries the soul. He casts not His pearls before swine; if thou play the hypocrite, though men baptize thee now, the Holy Spirit will not baptize thee\footnote{Cf. Procat. § 4: “The water will receive, but the Spirit will not accept thee.”}. But if thou approach with faith, though men minister in what is seen, the Holy Ghost bestows that which is unseen. Thou art coming to a great trial, to a great muster\footnote{\textgreek{στρατολογία}. Cf. Cat. iii. 3, \textgreek{μέλλετε στρατολογεῖσθαι}.}, in that one hour, which if thou throw away, thy disaster is irretrievable; but if thou be counted worthy of the grace, thy soul will be enlightened, thou wilt receive a power which thou hadst not, thou wilt receive weapons terrible to the evil spirits; and if thou cast not away thine arms, but keep the Seal upon thy soul, no evil spirit will approach thee; for he will be cowed; for verily by the Spirit of God are the evil spirits cast out.

37. If thou believe, thou shalt not only receive remission of sins, but also do things which pass man’s power\footnote{The same twofold grace is ascribed to Baptism in Cat. xiii. 23: “Thou receivest now remission of thy sins, and the gifts of the King’s spiritual bounty.”}. And mayest thou be worthy of the gift of prophecy also! For thou shalt receive grace according to the measure of thy capacity and not of my words; for I may possibly speak of but small things, yet thou mayest receive greater; since faith is a large affair\footnote{\textgreek{πραγματεία}. Cf. 2 Tim. ii. 4; and Luke xix. 13: \emph{Trade} (\textgreek{πραγματεύεοθε}) \emph{till I come}.}. All thy life long will thy guardian the Comforter abide with thee; He will care for thee, as for his own soldier; for thy goings out, and thy comings in, and thy plotting foes. And He will give thee gifts of grace of every kind, if thou grieve Him not by sin; for it is written, \emph{And grieve not the Holy Spirit of God, whereby ye were sealed unto the day of redemption}\footnote{Eph. iv. 30.}. What then, beloved, is it to preserve grace? Be ye ready to receive grace, and when ye have received it, cast it not away.

38. And may the very God of All, who spake by the Holy Ghost through the prophets, who sent Him forth upon the Apostles on the day of Pentecost in this place, Himself send Him forth at this time also upon you; and by Him keep us also, imparting His benefit in common to us all, that we may ever render up the fruits of the Holy Ghost, \emph{love, joy, peace, long-suffering, gentleness, goodness, faith, meekness, temperance}\footnote{Gal. v. 22, 23.}, in Christ Jesus our Lord:—By whom and with whom, together with the Holy Ghost, be glory to the Father, both now, and ever, and for ever and ever. Amen.

\xchapter{Lecture XVIII. On the Words, And in One Holy Catholic Church, and in the Resurrection of the Flesh, and the Life Everlasting.}

\textsc{Ezekiel xxxvii. 1}

The hand of the Lord was upon me, and carried me out in the Spirit of the Lord, and set me down in the midst of the valley which was full of bones.

1. \textsc{The} root of all good works is the hope of the Resurrection; for the expectation of the recompense nerves the soul to good works. For every labourer is ready to endure the toils, if he sees their reward in prospect; but when men weary themselves for nought, their heart soon sinks as well as their body. A soldier who expects a prize is ready for war, but no one is forward to die for a king who is indifferent about those who serve under him, and bestows no honours on their toils. In like manner every soul believing in a Resurrection is naturally careful of itself; but, disbelieving it, abandons itself to perdition. He who believes that his body shall remain to rise again, is careful of his robe, and defiles it not with fornication; but he who disbelieves the Resurrection, gives himself to fornication, and misuses his own body, as though it were not his own. Faith therefore in the Resurrection of the dead, is a great commandment and doctrine of the Holy Catholic Church; great and most necessary, though gainsaid by many, yet surely warranted by the truth. Greeks contradict it\footnote{Acts xvii. 32; xxvi. 24.}, Samaritans\footnote{Cf. § 12, below.} disbelieve it, heretics\footnote{Tertull. \emph{De Resurr. carnis}, cap. 2: “They acknowledge a half-resurrection, to wit of the soul only.” Compare Iren. I. xxiii. 5, on Menander’s assertion that his disciples attain to the resurrection by being baptized into him, and can die no more, but retain immortal youth: \emph{ib}. xxiv. 5. Basilides taught that “salvation belongs to the soul alone.” On the other forms of heresy concerning the Resurrection, see Suicer, \emph{Thesaurus}, \textgreek{᾽Ανάστασις}.} mutilate it; the contradiction is manifold, but the truth is uniform.

2. Now Greeks and Samaritans together argue against us thus. The dead man has fallen, and mouldered away, and is all turned into worms; and the worms have died also; such is the decay and destruction which has overtaken the body; how then is it to be raised? The shipwrecked have been devoured by fishes, which are themselves devoured. Of them who fight with wild beasts the very bones are ground to powder, and consumed by bears and lions. Vultures and ravens feed on the flesh of the unburied dead, and then fly away over all the world; whence then is the body to be collected? For of the fowls who have devoured it some may chance to die in India, some in Persia, some in the land of the Goths. Other men again are consumed by fire, and their very ashes scattered by rain or wind; whence is the body to be brought together again\footnote{The objections noticed in § 2 are discussed by Athenagoras, \emph{De Resurr}. capp. ii., iv.—viii.; Tatian. \emph{Or. ad Græcos}, cap. vi., Tertull. \emph{De Resurr. Carn}. cap. 63.}?

3. To thee, poor little feeble man, India is far from the land of the Goths, and Spain from Persia; but to God, who holds the whole \emph{earth in the hollow of His hand}\footnote{Is. xl. 12.}, all things are near at hand. Impute not then weakness to God, from a comparison of thy feebleness, but rather dwell on His power\footnote{On the argument from God’s power compare Athenagoras, \emph{De Resurr.} c. ix; Justin. M. \emph{De Resurr.} c. v; Theophil. \emph{ad Autolyc}. c. xiii.; Iren. V. iii. 2.}. Does then the sun, a small work of God, by one glance of his beams give warmth to the whole world; does the atmosphere, which God has made, encompass all things in the world; and is God, who is the Creator both of the sun, and of the atmosphere, far off from the world? Imagine a mixture of seeds of different plants (for as thou art weak concerning the faith, the examples which I allege are weak also), and that these different seeds are contained in thy single hand; is it then to thee, who art a man, a difficult or an easy matter to separate what is in thine hand, and to collect each seed according to its nature, and restore it to its own kind? Canst thou then separate the things in thine hand, and cannot God separate the things contained in His hand, and restore them to their proper place? Consider what I say, whether it is not impious to deny it?

4. But further, attend, I pray, to the very principle of justice, and come to thine own case. Thou hast different sorts of servants: and some are good and some bad; thou honourest therefore the good, and smitest the bad. And if thou art a judge, to the good thou awardest praise, and to the transgressors, punishment. Is then justice observed by thee a mortal man; and with God, the ever changeless King of all, is there no retributive justice\footnote{The argument from God’s justice is treated by Athenagor. \emph{De Resurr.} c. x. and xx.–xxiii.; Justin M. \emph{De Resurr.} c. viii.}? Nay, to deny it is impious. For consider what I say. Many murderers have died in their beds unpunished; where then is the righteousness of God? Yea, ofttimes a murderer guilty of fifty murders is beheaded once; where then shall he suffer punishment for the forty and nine? Unless there is a judgment and a retribution after this world, thou chargest God with unrighteousness. Marvel not, however, because of the delay of the judgment; no combatant is crowned or disgraced, till the contest is over; and no president of the games ever crowns men while yet striving, but he waits till all the combatants are finished, that then deciding between them he may dispense the prizes and the chaplets\footnote{\textgreek{τὴν στεφανηφορίαν}. Roe. Cas. A. Cf. Pind. \emph{Ol}. viii. 13; Eurip. \emph{Electr}. 862.}. Even thus God also, so long as the strife in this world lasts, succours the just but partially, but afterwards He renders to them their rewards fully.

5. But if according to thee there is no resurrection of the dead, wherefore condemnest thou the robbers of graves? For if the body perishes, and there is no resurrection to be hoped for, why does the violator of the tomb undergo punishment? Thou seest that though thou deny it with thy lips, there yet abides with thee an indestructible instinct of the resurrection.

6. Further, does a tree after it has been cut down blossom again, and shall man after being cut down blossom no more? And does the corn sown and reaped remain for the threshing floor, and shall man when reaped from this world not remain for the threshing? And do shoots of vine or other trees, when clean cut off and transplanted, come to life and bear fruit; and shall man, for whose sake all these exist, fall into the earth and not rise again? Comparing efforts, which is greater, to mould from the beginning a statue which did not exist, or to recast in the same shape that which had fallen? Is God then, who created us out of nothing, unable to raise again those who exist and are fallen\footnote{Athenag. \emph{De Resurr.} c. iii.: “If, when they did not exist, He made at their first formation the bodies of men, and their original elements, He will, when they are dissolved, in whatever manner that may take place, raise them again with equal ease.” Lactant. \emph{Institt}. VII. 23 fin.: \emph{Apost. Const}. V. 7.}? But thou believest not what is written of the resurrection, being a Greek: then from the analogy of nature consider these matters, and understand them from what is seen to this day. Wheat, it may be, or some other kind of grain, is sown; and when the seed has fallen, it dies and rots, and is henceforth useless for food. But that which has rotted, springs up in verdure; and though small when sown, springs up most beautiful. Now wheat was made for us; for wheat and all seeds were created not for themselves, but for our use; are then the things which were made for us quickened when they die, and do we for whom they were made, not rise again after our death\footnote{An eloquent statement of the argument for the resurrection from the analogies of nature occurs in Tertull. \emph{De Resurr.} c. xii. That it was not unknown to Cyril, seems probable from the concluding sentence: “And surely if all things rise again for man, for whom they have been provided—but not for man unless for his flesh also—how can the flesh itself perish utterly, for the sake and service of which nothing is allowed to perish.” Tertullian himself was probably indebted, as Bp. Lightfoot suggests, to Clemens. Rom. \emph{Epist. ad Corinth.} xxiv. Cf. Lactant. \emph{Div. Inst.} vii. 4.}?

7. The season is winter\footnote{Cf. Cat. iv. 30. These passages shew that the Lectures were delivered in a year when Easter fell early, as was the case in 348 \textsc{a.d.}}, as thou seest; the trees now stand as if they were dead: for where are the leaves of the fig-tree? where are the clusters of the vine? These in winter time are dead, but green in spring; and when the season is come, there is restored to them a quickening as it were from a state of death. For God, knowing thine unbelief, works a resurrection year by year in these visible things; that, beholding what happens to things inanimate, thou mayest believe concerning things animate and rational. Further, flies and bees are often drowned in water, yet after a while revive\footnote{In such cases there is, of course, no actual death.}; and species of dormice\footnote{The \textgreek{μυοξός} is supposed by the Benedictine Editor to be the toad (“Inventusque cavis bufo,” Virg. \emph{Georg}. i. 185), by others the marmot (mus Alpinus). More probably it is the dormouse (myoxis glis), which stores up provisions for the winter, though it sleeps through much of that season.}, after remaining motionless during the winter, are restored in the summer (for to thy slight thoughts like examples are offered); and shall He who to irrational and despised creatures grants life supernaturally, not bestow it upon us, for whose sake He made them?

8. But the Greeks ask for a resurrection of the dead still manifest; and say that, even if these creatures are raised, yet they had not utterly mouldered away; and they require to see distinctly some creature rise again after complete decay. God knew men’s unbelief, and provided for this purpose a bird, called a Phoenix\footnote{The story of the Phœnix as told by Herodotus, II. 73, is as follows: “They have also another sacred bird called the Phœnix, which I myself have never seen, except in pictures. Indeed it is a great rarity even in Egypt, only coming there (according to the accounts of the people of Heliopolis) once in five hundred years, when the old phœnix dies.…They tell a story of what this bird does, which does not seem to me to be credible; that he comes all the way from Arabia, and brings the parent bird, all plastered over with myrrh, to the temple of the Sun, and there buries the body.” \par The many variations and fabulous accretions of the story are detailed by Suicer, \emph{Thesaurus}, \textgreek{Φοῖνιξ}, and by Bp. Lightfoot in a long and interesting note on Clemens Rom. \emph{Epist. ad Cor.} xxv. Cyril borrows the story from Clement almost verbally, yet not without some variations, which will be noticed below. The legend with all its miraculous features is told by Ovid, \emph{Metamorph}. xv. 392, by Claudian, \emph{Phœnix}, and by the Pseudo-Lactantius in an Elegiac poem, \emph{Phœnix}, included in Weber’s \emph{Corpus Poetarum Latinorum}, and literally translated in Clark’s \emph{Ante-Nicene Library}. See also Tertull. \emph{De Resurr. Carn}. c. xiii.}. This bird, as Clement writes, and as many more relate, being the only one of its kind\footnote{\textgreek{μονογενὲς ὕπαρχον}, Clem. Rom. \emph{ubi supra}. Cf. Origen, \emph{contra Celsum}, iv. 98: \emph{Apost. Const}. V. 7: “a bird single in its kind, which they say is without a mate, and the only one in the creation.” Pseudo-Lactant. \emph{v}. 30. \par “Hoc nemus hos lucos avis incolit unica, phœnix, \par Unica, sed vivit morte refecta suâ”}, arrives in the land of the Egyptians at periods of five hundred years, shewing forth the resurrection, not in desert places, lest the occurrence of the mystery should remain unknown, but appearing in a notable city\footnote{“By day, in the sight of all” (Clem. R.) The city was Heliopolis, according to Herodotus and the other ancient authors. But Milton, \emph{Paradise Lost}, V. 272— \par ‘A phœnix gaz’d by all, as that \emph{sole} bird, \par When to enshrine his reliques in the Sun’s \par Bright temple to Ægyptian Thebes he flies.’ \par Why does Milton despatch his bird to Thebes rather than Heliopolis?” (Lightfoot).}, that men might even handle what would otherwise be disbelieved. For it makes itself a coffin\footnote{Ovid, \emph{Met}. xv. 405: “Fertque pius cunasque suas patriumque sepulcrum.” See the Commentaries on Job xxix. 18: \emph{I shall die in my nest, and I shall multiply my days as the sand. } Margin \textsc{R.V.} Or, \emph{the phœnix}.} of frankincense and myrrh and other spices, and entering into this when its years are fulfilled, it evidently dies and moulders away. Then from the decayed flesh of the dead bird a worm is engendered, and this worm when grown large is transformed into a bird;—and do not disbelieve this, for thou seest the offspring of bees also fashioned thus out of worms\footnote{The mode of reproduction in bees was regarded by Aristotle as mysterious, having in it something supernatural (\textgreek{θεῖον}): \emph{De Generatione Animal}. III. 10. 1, 27. In the story of the phœnix Herodotus makes no mention of the “worm.”}, and from eggs which are quite fluid thou hast seen wings and bones and sinews of birds issue. Afterwards the aforesaid Phoenix, becoming fledged and a full-grown Phoenix, like the former one, soars up into the air such as it had died, shewing forth to men a most evident resurrection of the dead. The Phoenix indeed is a wondrous bird, yet it is irrational, nor ever sang praise to God; it flies abroad through the sky, but it knows not who is the Only-begotten Son of God. Has then a resurrection from the dead been given to this irrational creature which knows not its Maker, and to us who ascribe glory to God and keep His commandments, shall there no resurrection be granted?

9. But since the sign of the Phoenix is remote and uncommon, and men still disbelieve our resurrection, take again the proof of this from what thou seest every day. A hundred or two hundred years ago, we all, speakers and hearers, where were we? Know we not the groundwork of the substance of our bodies? Knowest thou not how from weak and shapeless and simple\footnote{\textgreek{μονοειδής}.} elements we are engendered, and out of what is simple and weak a living man is formed? and how that weak element being made flesh is changed into strong sinews, and bright eyes, and sensitive nose, and hearing ears, and speaking tongue, and beating heart, and busy hands, and swift feet, and into members of all kinds\footnote{For a similar argument, see Lactant. \emph{De Resurr.} c xvii.}? and how that once weak element becomes a shipwright, and a builder, and an architect, and a craftsman of various arts, and a soldier, and a ruler, and a lawgiver, and a king? Cannot God then, who has made us out of imperfect materials, raise us up when we have fallen into decay? He who thus flames a body out of what is vile, cannot He raise the fallen body again? And He who fashions that which is not, shall He not raise up that which is and is fallen?

10. Take further a manifest proof of the resurrection of the dead, witnessed month by month in the sky and its luminaries\footnote{Clem. Rom. \emph{Epist. ad Cor}. xxiv: “Day and night shew unto us the resurrection. The night falleth asleep, and day ariseth; the day departeth, and night cometh on.”}. The body of the moon vanishes completely, so that no part of it is any more seen, yet it fills again, and is restored to its former state\footnote{Tertull. \emph{de Resurr. Carnis}, xii.: “Readorned also are the mirrors of the moon, which her monthly course had worn away.”…“The whole of this revolving order of things bears witness to the resurrection of the dead.”}; and for the perfect demonstration of the matter, the moon at certain revolutions of years suffering eclipse and becoming manifestly changed into blood, yet recovers its luminous body: God having provided this, that thou also, the man who art formed of blood, mightest not refuse credence to the resurrection of the dead, but mightest believe concerning thyself also what thou seest in respect of the moon. These therefore use thou as arguments against the Greeks; for with them who receive not what is written fight thou with unwritten weapons, by reasonings only and demonstrations; for these men know not who Moses is, nor Esaias, nor the Gospels, nor Paul.

11. Turn now to the Samaritans, who, receiving the Law only, allow not the Prophets. To them the text just now read from Ezekiel appears of no force, for, as I said, they admit no Prophets; whence then shall we persuade the Samaritans also? Let us go to the writings of the Law. Now God says to Moses, \emph{I am the God of Abraham, and of Isaac, and of Jacob}\footnote{Ex. iii. 6. Cf. Matt. xxii. 32: “\emph{He is not the God of the dead, but of the living}.”}; this must mean of those who have being and subsistence. For if Abraham has come to an end, and Isaac and Jacob, then He is the God of those who have no being. When did a king ever say, I am the king of soldiers, whom he had not? When did any display wealth which he possessed not? Therefore Abraham and Isaac and Jacob must subsist, that God may be the God of those who have being; for He said not, “I was their God,” but \emph{I am}. And that there is a judgment, Abraham shews in saying to the Lord, \emph{He who judgeth all the earth, shall He not execute judgment}\footnote{Gen. xviii. 25.}?

12. But to this the foolish Samaritans object again, and say that the souls possibly of Abraham and Isaac and Jacob continue, but that their bodies cannot possibly rise again. Was it then possible that the rod of righteous Moses should become a serpent, and is it impossible that the bodies of the righteous should live and rise again? And was that done contrary to nature, and shall they not be restored according to nature? Again, the rod of Aaron, though cut off and dead, budded, \emph{without the scent of waters}\footnote{Job xiv. 9.}, and though under a roof, sprouted forth into blossoms as in the fields; and though set in dry places, yielded in one night the flowers and fruit of plants watered for many years. Did Aaron’s rod rise, as it were, from the dead, and shall not Aaron himself be raised? And did God work wonders in wood, to secure to him the high-priesthood, and will He not vouchsafe a resurrection to Aaron himself? A woman also was made salt contrary to nature; and flesh was turned into salt; and shall not flesh be restored to flesh? Was Lot’s wife made a pillar of salt, and shall not Abraham’s wife be raised again? By what power was Moses’ hand changed, which even within one hour became as snow, and was restored again? Certainly by God’s command. Was then His command of force then, and has it no force now?

13. And whence in the beginning came man into being at all, O ye Samaritans, most senseless of all men? Go to the first book of the Scripture, which even you receive; \emph{And God formed man of the dust of the ground}\footnote{Gen. ii. 7.}. Is dust transformed into flesh, and shall not flesh be again restored to flesh? You must be asked too, whence the heavens had their being, and earth, and seas? Whence sun, and moon, and stars? How from the waters were made the things which fly and swim? And how from earth all its living things? Were so many myriads brought from nothing into being, and shall we men, who bear God’s image, not be raised up? Truly this course is full of unbelief, and the unbelievers are much to be condemned; when Abraham addresses the Lord as \emph{the Judge of all the earth}, and the learners of the Law disbelieve; when it is written that man is of the earth, and the readers disbelieve it\footnote{The anomalous construction \textgreek{ὅταν γέγραπται .…καὶ ἀπιστῶσιν} may be explained by the consideration, that the uncertainty expressed in \textgreek{ὅταν} attaches only to the latter Verb. See Winer’s \emph{Grammar of N.T. Greek}, P. III. sect. xlii. 5.}.

14. These questions, therefore, are for them, the unbelievers: but the words of the Prophets are for us who believe. But since some who have also used the Prophets believe not what is written, and allege against us that passage, \emph{The ungodly shall not rise up in judgment}\footnote{Ps. i. 5: \emph{The wicked shall not stand in the judgment} (\textsc{R.V.}).}, and, \emph{For if man go down to the grave he shall come up no more}\footnote{Job vii. 9.}, and, \emph{The dead shall not praise Thee, O Lord}\footnote{Ps. cxv. 17.},—for of what is well written, they have made ill use—it will be well in a cursory manner, and as far as is now possible, to meet them. For if it is said, that \emph{the ungodly shall not rise up in judgment}, this shews that they shall rise, not in judgment, but in condemnation; for God needs not long scrutiny, but close on the resurrection of the ungodly follows also their punishment. And if it is said, \emph{The dead shall not praise Thee, O Lord}, this shews, that since in this life only is the appointed time for repentance and pardon, for which they who enjoy it shall \emph{praise the Lord}, it remains not after death for them who have died in sins to give praise as the receivers of a blessing, but to bewail themselves; for praise belongs to them who give thanks, but to them who are under the scourge, lamentation. Therefore the just then offer praise; but they who have died in sins have no further season for confession\footnote{As to the bearing of this passage on the doctrine of Purgatory and prayer for the dead see note on xxiii. 10.}.

15. And respecting that passage, \emph{If a man go down to the grave, he shall come up no more}, observe what follows, for it is written, \emph{He shall come up no more, neither shall he return to his own house}. For since the whole world shall pass away, and every house shall be destroyed, how shall he return to his own house, there being henceforth a new and different earth? But they ought to have heard Job, saying, \emph{For there is hope of a tree; for if it be cut down, it will sprout again, and the tender branch thereof will not cease. For though the root thereof wax old in the earth, and the stock thereof die in the rocky ground; yet from the scent of water it will bud, and bring forth a crop like a new plant. But man when he dies, is gone; and when mortal man falls, is he no more}\footnote{Job xiv. 7–10.}? As it were remonstrating and reproving (for thus ought we to read the words \emph{is no more} with an interrogation\footnote{There is no indication of a question in the Septuagint version of the passage, which means in the Hebrew, \emph{and where is he?} (\textsc{A.V.} and \textsc{R.V.}): Vulg. \emph{ubi, quæso, est}?}); he says since a tree falls and revives, shall not man, for whom all trees were made, himself revive? And that thou mayest not suppose that I am forcing the words, read what follows; for after saying by way of question, \emph{When mortal man falls, is he no more?} he says, \emph{For if a man die, he shall live again\footnote{Job xiv. 14: \emph{For if a man die, shall he live again}? (\textsc{A.V.} and \textsc{R.V.}). By omitting the interrogation here, and inserting it above in \emph{v}. 10, Cyril exactly inverts the meaning.}}; and immediately he adds, \emph{I will wait till I be made again}\footnote{Ib. v. 14: (A.V.) \emph{All the days of my appointed time} (\textsc{R.V.} \emph{of my warfare}) \emph{will I wait, till my change} (\textsc{R.V.} \emph{release}) \emph{come.}}; and again elsewhere, \emph{Who shall raise up on the earth my skin, which endures these things}\footnote{Job xix. 26: (R.V.) \emph{and that he shall stand up at the last upon the earth: and after my skin hath been thus destroyed, \&c.} Cyril, as usual, follows the Septuagint.}. And Esaias the Prophet says, \emph{The dead men shall rise again, and they that are in the tombs shall awake}\footnote{Is. xxvi. 19.}. And the Prophet Ezekiel now before us, says most plainly, \emph{Behold I will open your graves, and bring you up out of your graves}\footnote{Ezek. xxxvii. 12.}. And Daniel says, \emph{Many of them that sleep in the dust of the earth shall arise, some to everlasting life, and some to everlasting shame}\footnote{Dan. xii. 2.}.

16. And many Scriptures there are which testify of the Resurrection of the dead; for there are many other sayings on this matter. But now, by way of remembrance only, we will make a passing mention of the raising of Lazarus on the fourth day; and just allude, because of the shortness of the time, to the widow’s son also who was raised, and merely for the sake of reminding you, let me mention the ruler of the synagogue’s daughter, and the rending of the rocks, and how \emph{there arose many bodies of the saints which slept}\footnote{Matt. xxvii. 52.}, their graves having been opened. But specially be it remembered that \emph{Christ has been raised from the dead}\footnote{1 Cor. xv. 20.}. I speak but in passing of Elias, and the widow’s son whom he raised; of Elisseus also, who raised the dead twice; once in his lifetime, and once after his death. For when alive he wrought the resurrection by means of his own soul\footnote{2 Kings iv. 34.}; but that not the souls only of the just might be honoured, but that it might be believed that in the bodies also of the just there lies a power, the corpse which was cast into the sepulchre of Elisseus, when it touched the dead body of the prophet, was quickened, and the dead body of the prophet did the work of the soul, and that which was dead and buried gave life to the dead, and though it gave life, yet continued itself among the dead. Wherefore? Lest if Elisseus should rise again, the work should be ascribed to his soul alone; and to shew, that even though the soul is not present, a virtue resides in the body of the saints, because of the righteous soul which has for so many years dwelt in it, and used it as its minister\footnote{“The worship of relics, and the belief in them as remedies and a protection against evil, originated in the 4th century. They first (?) appear in writings, none of which are earlier than the year 370: but they prevailed rapidly when they had once taken root” (Scudamore, \emph{Dict. Chr. Antiq.} “Relics,” p. 1770). Bingham (\emph{Ant}. xxiii. 4, § 7) quotes a law of Theodosius, “that no one should remove any dead body that was buried, from one place to another; that no one should sell or buy the relics of Martyrs: but if any one was minded to build over the grave where a martyr was buried, a church to be called a \emph{martyrium}, in respect to him, he should have liberty to do it.” The law wholly failed to suppress a superstition which was sanctioned by such men as Cyril, Basil, Chrysostom, Ambrose, and Augustine.}. And let us not foolishly disbelieve, as though this thing had not happened: for if handkerchiefs and aprons, which are from without, touching the bodies of the diseased, raised up the sick, how much more should the very body of the Prophet raise the dead?

17. And with respect to these instances we might say much, rehearsing in detail the marvellous circumstances of each event: but as you have been already wearied both by the superposed fast of the Preparation\footnote{\textgreek{ἐκ τῆς ὑπερθέσεως τῆς νηστείας τῆς παρασκευῆς}, Ed. Bened. “The ecclesiastical term \textgreek{τῆς ὑπερθέσεως} we have rendered, according to the interpretation received among the Latins, by the word ‘superpositio.’ The ancients meant by it a fast continued for two or three days without food. Moreover, since the great week was observed with severer fastings, there were many who passed either the whole week or four, three, or two days, namely the Preparation and the Holy Sabbath (Easter Eve), entirely fasting as is testified by S. Irenæus (Euseb. \emph{Hist}. V. 24) and others. The continuance of the fast throughout the Friday and Saturday was highly approved, as may be seen from the \emph{Apostolical Constitutions}, V. 18.” The passage referred to is as follows: “Do you therefore fast on the days of the Passover, beginning from the second day of the week until the Preparation and the Sabbath, six days, making use only of bread, and salt, and herbs, and water for your drink: but abstain on these days from wine and flesh, for they are days of lamentation and not of fasting. Do ye who are able fast throughout the Preparation and the Sabbath entirely, tasting nothing till the cockcrowing at night; but if any one is not able to combine them both, let the Sabbath at least be observed.”}, and by the watchings\footnote{The fast of the Great Sabbath was to be continued through the night, as prescribed in the \emph{Apost. Const}. V. 19: “Continue until cock-crowing and break off your fast at dawn of the first day of the week, which is the Lord’s day, keeping awake from evening until cock-crowing: and assembling together in the Church, watch and pray and beseech God, in your night-long vigil, reading the Law, the Prophets, and the Psalms, until the crowing of the cocks: and after baptizing your Catechumens, and reading the Gospel in fear and trembling, and speaking to the people the things pertaining to salvation, so cease from your mourning.” A chief reason for the watching was that Christ was expected to return at the same hour in which He rose. On the meaning of “superposition” see Routh’s note on the Synodical Epistle of Irenæus to Victor of Rome (\emph{Rell. Sac.} ii. p. 45, ss.), and the passage of Dionysius of Alexandria there quoted.}, let what has been cursorily spoken concerning them suffice for a while; these words having been as it were sown thinly, that you, receiving the seed like richest ground, may in bearing fruit increase them. But be it remembered, that the Apostles also raised the dead; Peter raised Tabitha in Joppa, and Paul raised Eutychus in Troas; and thus did all the other Apostles, even though the wonders wrought by each have not all been written. Further, remember all the sayings in the first Epistle to the Corinthians, which Paul wrote against them who said, \emph{How are the dead raised, and with} \emph{what manner of body do they come}\footnote{1 Cor. xv. 35.}? And how he says, \emph{For if the dead rise not, then is not Christ raised}\footnote{Ib. v. 16.}; and how he called them \emph{fools}\footnote{Ib. v. 36.}, who believed not; and remember the whole of his teaching there concerning the resurrection of the dead, and how he wrote to the Thessalonians, \emph{But we would not have you to be ignorant, brethren, concerning them which are asleep, that ye sorrow not, even as the rest which have no hope}\footnote{1 Thess. iv. 13.}, and all that follows: but chiefly that, \emph{And the dead in Christ shall rise first}\footnote{Ib. v. 16.}.

18. But especially mark this, how very pointedly\footnote{\textgreek{μονονουχὶ δακτυλοδεικτῶν}.} Paul says, \emph{For this corruptible must put on incorruption, and this mortal must put on immortality}\footnote{1 Cor. xv. 53.}. For this body shall be raised not remaining weak as now; but raised the very same body, though by putting on incorruption it shall be fashioned anew\footnote{\textgreek{μεταποιεῖται}. The meaning of this word as applied to the Eucharistic elements is fully discussed, and illustrated from its use by Cyril and other Fathers, by Dr. Pusey (\emph{Real Presence}, p. 189).},—as iron blending with fire becomes fire, or rather as He knows how, the Lord who raises us. This body therefore shall be raised, but it shall abide not such as it now is, but an eternal body; no longer needing for its life such nourishment as now, nor stairs for its ascent, for it shall be made \emph{spiritual}, a marvellous thing, such as we cannot worthily speak of. \emph{Then}, it is said, \emph{shall the righteous shine forth as the sun}\footnote{Matt. xiii. 43.}, and the moon, \emph{and as the brightness of the firmament}\footnote{Dan. xii. 3.}. And God, fore-knowing men’s unbelief, has given to little worms in the summer to dart beams of light from their body\footnote{Cyril refers to the glow-worm (\textgreek{πυγολαμπίς}, Aristot. \emph{Hist. Animal.} V. 19, 14), or some other species of Lampyris (Arist. \emph{de Partilus Animal}. I. 3. 3).}, that from what is seen, that which is looked for might be believed; for He who gives in part is able to give the whole also, and He who made the worm radiant with light, will much more illuminate a righteous man.

19. We shall be raised therefore, all with our bodies eternal, but not all with bodies alike: for if a man is righteous, he will receive a heavenly body, that he may be able worthily to hold converse with Angels; but if a man is a sinner, he shall receive an eternal body, fitted to endure the penalties of sins, that he may burn eternally in fire, nor ever be consumed\footnote{Cf. Cat. iv. 31.}. And righteously will God assign this portion to either company; for we do nothing without the body. We blaspheme with the mouth, and with the mouth we pray. With the body we commit fornication, and with the body we keep chastity. With the hand we rob, and by the hand we bestow alms; and the rest in like manner. Since then the body has been our minister in all things, it shall also share with us in the future the fruits of the past\footnote{\textgreek{τῶν γενομένων}. With the reading \textgreek{γινομένων} (Codd. Monn. Vind.), the meaning will be—“share with us in the future what shall happen to us then.” On the argument of this section compare the passages quoted on § 4, note 7.}.

20. Therefore, brethren, let us be careful of our bodies, nor misuse them as though not our own. Let us not say like the heretics, that this vesture of the body belongs not to us, but let us be careful of it as our own; for we must give account to the Lord of all things done through the body. Say not, none seeth me; think not, that there is no witness of the deed. Human witness oftentimes there is not; but He who fashioned us, an unerring \emph{witness}, abides \emph{faithful in heaven}\footnote{Ps. lxxxix. 37.}, and beholds what thou doest. And the stains of sin also remain in the body; for as when a wound has gone deep into the body, even if there has been a healing, the scar remains, so sin wounds soul and body, and the marks of its scars remain in all; and they are removed only from those who receive the washing of Baptism. The past wounds therefore of soul and body God heals by Baptism; against future ones let us one and all jointly guard ourselves, that we may keep this vestment of the body pure, and may not for practising fornication and sensual indulgence or any other sin for a short season, lose the salvation of heaven, but may inherit the eternal kingdom of God; of which may God, of His own grace, deem all of you worthy.

21. Thus much in proof of the Resurrection of the dead; and now, let me again recite to you the profession of the faith, and do you with all diligence pronounce it while I speak\footnote{Cat. V. 12, notes 7 and 4. Cf. Plat. Theaet. 204 C: \textgreek{ἐφ᾽ ἑκάστης λέξεως}, “each time we speak.”}, and remember it.

\_\_\_\_\_\_\_\_\_\_\_\_\_\_\_\_\_\_\_\_\_\_

22. The Faith which we rehearse contains in order the following, “\textsc{And in one Baptism of repentance for the remission of sins; and in one Holy Catholic Church; and in the resurrection of the flesh; and in eternal life}.” Now of Baptism and repentance I have spoken in the earliest Lectures; and my present remarks concerning the resurrection of the dead have been made with reference to the Article “In the resurrection of the flesh.” Now then let me finish what still remains to be said for the Article, “In one Holy Catholic Church,” on which, though one might say many things, we will speak but briefly.

23. It is called Catholic then because it extends over all the world, from one end of the earth to the other; and because it teaches universally and completely one and all the doctrines which ought to come to men’s knowledge, concerning things both visible and invisible, heavenly and earthly\footnote{Bishop Lightfoot (Ignatius\emph{, ad Smyrnæos}, viii.) traces the original and later senses of the word “Catholic” very fully. “In its earliest usages, therefore, as a fluctuating epithet of \textgreek{ἐκκλησία}, ‘catholic’ means ‘universal,’ as opposed to ‘individual,’ ‘particular.’ In its later sense, as a fixed attribute, it implies orthodoxy as opposed to heresy, conformity as opposed to dissent.” Commenting on this passage of Cyril, the Bishop adds that “these two latter reasons, that it (the Church) is comprehensive in doctrine, and that it is universal in application, can only be regarded as secondary glosses.”}; and because it brings into subjection to godliness the whole race of mankind, governors and governed, learned and unlearned; and because it universally treats and heals the whole class of sins, which are committed by soul or body, and possesses in itself every form of virtue which is named, both in deeds and words, and in every kind of spiritual gifts.

24. And it is rightly named (Ecclesia) because it calls forth\footnote{\textgreek{ἐκκαλεῖσθαι}. Cf. Heb. xii. 23.} and assembles together all men; according as the Lord says in Leviticus, \emph{And make an assembly for all the congregation at the door of the tabernacle of witness}\footnote{Lev. viii. 3: \textgreek{ἐκκλησίασον}.}. And it is to be noted, that the word \emph{assemble}, is used for the first time in the Scriptures here, at the time when the Lord puts Aaron into the High-priesthood. And in Deuteronomy also the Lord says to Moses, \emph{Assemble the people unto Me, and let them hear My words, that they may learn to fear Me}\footnote{Deut. iv. 10.}. And he again mentions the name of the Church, when he says concerning the Tables, \emph{And on them were written all the words which the Lord spake with you in the mount out of the midst of the fire in the day of the Assembly}\footnote{Ib. ix. 10: \textgreek{ἐκκλησίας}.}; as if he had said more plainly, in the day in which ye were called and gathered together by God. The Psalmist also says, \emph{I will give thanks unto Thee, O Lord, in the great Congregation; I will praise Thee among much people}\footnote{Ps. xxxv. 18; Heb. ii. 12.}.

25. Of old the Psalmist sang, \emph{Bless ye God in the congregations, even the Lord, (ye that are) from the fountains of Israel}\footnote{Ps. lxviii. 26: \textgreek{ἐν ἐκκλησίαις}.}. But after the Jews for the plots which they made against the Saviour were cast away from His grace, the Saviour built out of the Gentiles a second Holy Church, the Church of us Christians, concerning which he said to Peter, \emph{And upon this rock I will build My Church, and the gates of hell shall not prevail against it}\footnote{Matt. xvi. 18.}. And David prophesying of both these, said plainly of the first which was rejected, \emph{I have hated the Congregation of evil doers}\footnote{Ps. xxvi. 5.}; but of the second which is built up he says in the same Psalm, \emph{Lord, I have loved the beauty of Thine house}\footnote{Ps. xxvi. 8: Sept. \textgreek{εὐπρέπειαν} . \textsc{R.V.} and \textsc{A.V.} “habitation.”}; and immediately afterwards, \emph{In the Congregations will I bless thee, O Lord}\footnote{Ib. v. 12.}. For now that the one Church in Judæa is cast off, the Churches of Christ are increased over all the world; and of them it is said in the Psalms, \emph{Sing unto the Lord a new song, His praise in the Congregation of Saints}\footnote{Ps. cxlix. 1.}. Agreeably to which the prophet also said to the Jews, \emph{I have no pleasure in you, saith the Lord Almighty}\footnote{Mal. i. 10.}; and immediately afterwards, \emph{For from the rising of the sun even unto the going down of the same, My name is glorified among the Gentiles}\footnote{Ib. v. 11.}. Concerning this Holy Catholic Church Paul writes to Timothy, \emph{That thou mayest know how thou oughtest to behave thyself in the House of God, which is the Church of the Living God, the pillar and ground of the truth}\footnote{1 Tim. iii. 15.}.

26. But since the word Ecclesia is applied to different things (as also it is written of the multitude in the theatre of the Ephesians, \emph{And when he had thus spoken, he dismissed the Assembly}\footnote{Acts xix. 14.}), and since one might properly and truly say that there is a \emph{Church of evil doers}, I mean the meetings of the heretics, the Marcionists and Manichees, and the rest, for this cause the Faith has securely delivered to thee now the Article, “And in one Holy Catholic Church;” that thou mayest avoid their wretched meetings, and ever abide with the Holy Church Catholic in which thou wast regenerated. And if ever thou art sojourning in cities, inquire not simply where the Lord’s House is (for the other sects of the profane also attempt to call their own dens houses of the Lord), nor merely where the Church is, but where is the Catholic Church. For this is the peculiar name of this Holy Church, the mother of us all, which is the spouse of our Lord Jesus Christ, the Only-begotten Son of God (for it is written, \emph{As Christ also loved the Church and gave Himself for it}\footnote{Eph. v. 25.}, and all the rest,) and is a figure and copy of \emph{Jerusalem which is above, which is free, and the mother of us all}\footnote{Gal. iv. 26.}; which before was barren, but now has many children.

27. For when the first Church was cast off, in the second, which is the Catholic Church, God \emph{hath set,} as Paul says, \emph{first Apostles, secondly Prophets, thirdly teachers, then miracles, then gifts of healings, helps, governments, divers kinds of tongues}\footnote{1 Cor. xii. 28.}, and every sort of virtue, I mean wisdom and understanding, temperance and justice, mercy and loving-kindness, and patience unconquerable in persecutions. She, \emph{by the armour of righteousness on the right hand and on the left, by honour and dishonour}\footnote{2 Cor. vi. 7, 8.}, in former days amid persecutions and tribulations crowned the holy martyrs with the varied and blooming chaplets of patience, and now in times of peace by God’s grace receives her due honours from \emph{kings and those who are in high place}\footnote{1 Tim. ii. 2.}, and from every sort and kindred of men. And while the kings of particular nations have bounds set to their authority, the Holy Church Catholic alone extends her power without limit over the whole world; \emph{for God}, as it is written, \emph{hath made her border peace}\footnote{Ps. cxlvii. 14.}. But I should need many more hours for my discourse, if I wished to speak of all things which concern her.

\_\_\_\_\_\_\_\_\_\_\_\_\_\_\_\_\_\_\_\_\_\_

28. In this Holy Catholic Church receiving instruction and behaving ourselves virtuously, we shall attain the kingdom of heaven, and inherit \textsc{eternal life}; for which also we endure all toils, that we may be made partakers thereof from the Lord. For ours is no trifling aim, but our endeavour is for eternal life. Wherefore in the profession of the Faith, after the words, “\textsc{And in the resurrection of the flesh},” that is, of the dead (of which we have discoursed), we are taught to believe also “\textsc{in the life eternal},” for which as Christians we are striving.

29. The real and true life then is the Father, who through the Son in the Holy Spirit pours forth as from a fountain His heavenly gifts to all; and through His love to man, the blessings of the life eternal are promised without fail to us men also. We must not disbelieve the possibility of this, but having an eye not to our own weakness but to His power, we must believe; \emph{for with God all things are possible}. And that this is possible, and that we may look for eternal life, Daniel declares, \emph{And of the many righteous shall they shine as the stars for ever and ever}\footnote{Dan. xii. 3, Sept.}. And Paul says, \emph{And so shall we be ever with the Lord}\footnote{1 Thess. iv. 17.}: for the \emph{being for ever with the Lord} implies the life eternal. But most plainly of all the Saviour Himself says in the Gospel, \emph{And these shall go away into eternal punishment, but the righteous into life eternal}\footnote{Matt. xxv. 46.}.

30. And many are the proofs concerning the life eternal. And when we desire to gain this eternal life, the sacred Scriptures suggest to us the ways of gaining it; of which, because of the length of our discourse, the texts we now set before you shall be but few, the rest being left to the search of the diligent. They declare at one time that it is by faith; for it is written, \emph{He that believeth on the Son hath eternal life}\footnote{John iii. 36.}, and what follows; and again He says Himself, \emph{Verily, verily, I say unto you, He that heareth My word, and believeth Him that sent Me, hath eternal life}\footnote{Ib. v. 24.}, and the rest. At another time, it is by the preaching of the Gospel; for He says, that \emph{He that reapeth receiveth wages, and gathereth fruit unto life eternal}\footnote{Ib. iv. 36.}. At another time, by martyrdom and confession in Christ’s name; for He says, \emph{And he that hateth his life in this world, shall keep it unto life eternal}\footnote{Ib. xii. 25.}. And again, by preferring Christ to riches or kindred; \emph{And every one that hath forsaken brethren, or sisters}\footnote{Matt. xix. 29.}, and the rest, \emph{shall inherit eternal life}. Moreover it is by keeping the commandments, \emph{Thou shalt not commit adultery, Thou shalt not kill}\footnote{Ib. vv. 16–18.}, and the rest which follow; as He answered to him that came to Him, and said, \emph{Good Master, what shall I do that I may have eternal life}\footnote{Mark. x. 17.}? But further, it is by departing from evil works, and henceforth serving God; for Paul says, \emph{But now being made free from sin, and become servants to God, ye have your fruit unto sanctification, and the end eternal life}\footnote{Rom. vi. 22.}.

31. And the ways of finding eternal life are many, though I have passed over them by reason of their number. For the Lord in His loving-kindness has opened, not one or two only, but many doors, by which to enter into the life eternal, that, as far as lay in Him, all might enjoy it without hindrance. Thus much have we for the present spoken within compass concerning \textsc{the life eternal}, which is the last doctrine of those professed in the Faith, and its termination; which life may we all, both teachers and hearers, by God’s grace enjoy!

\_\_\_\_\_\_\_\_\_\_\_\_\_\_\_\_\_\_\_\_\_\_

32. And now, brethren beloved, the word of instruction exhorts you all, to prepare your souls for the reception of the heavenly gifts. As regards the Holy and Apostolic Faith delivered to you to profess, we have spoken through the grace of the Lord as many Lectures, as was possible, in these past days of Lent; not that this is all we ought to have said, for many are the points omitted; and these perchance are thought out better by more excellent teachers. But now the holy day of the Passover is at hand, and ye, beloved\footnote{\textgreek{τῆς ὑμετέρας ἐν Χριστῷ ἀγάπης}. Cf. Cat. xvii. 1, note 1. Athan. \emph{Epist. ad Epict}. § 2: \textgreek{παρὰ τῇ σῇ θεοσεβείά}. \emph{ad Serap}. iv. 1: \textgreek{παρὰ τῆς σῆς εὐλαβείας}.} in Christ, are to be enlightened \emph{by the Laver of regeneration}. Ye shall therefore again be taught what is requisite, if God so will; with how great devotion and order you must enter in when summoned, for what purpose each of the holy mysteries of Baptism is performed, and with what reverence and order you must go from Baptism to the Holy Altar of God, and enjoy its spiritual and heavenly mysteries; that your souls being previously enlightened by the word of doctrine, ye may discover in each particular the greatness of the gifts bestowed on you by God.

33. And after Easter’s Holy Day of salvation, ye shall come on each successive day, beginning from the second day of the week, after the assembly into the Holy Place of the Resurrection\footnote{The place meant is not the Church of the Resurrection in which the Service had been held, but the Anastasis or actual cave of the Resurrection, which Constantine had so enlarged by additional works that a discourse to the people could be held there: for Jerome (\emph{Epist}. 61) relates that Epiphanius had preached in that place in front of the Lord’s sepulchre to clergy and people in the hearing of John the Bishop (Ben. Ed.).}, and there, if God permit, ye shall hear other Lectures; in which ye shall again be taught the reasons of every thing which has been done, and shall receive the proofs thereof from the Old and New Testaments,—first, of the things done just before Baptism,—next, how ye were cleansed from your sins by the Lord, \emph{by the washing of water with the word}\footnote{Eph. v. 26.},—and how like Priests ye have become partakers of the Name of Christ,—and how the Seal of the fellowship of the Holy Ghost was given to you,—and concerning the mysteries at the Altar of the New Testament, which have taken their beginning from this place, both what the Divine Scriptures have delivered to us, and what is the power of these mysteries, and how ye must approach them, and when and how receive them;—and at the end of all, how for the time to come ye must behave yourselves worthily of this grace both in words and deeds, that you may all be enabled to enjoy the life everlasting. And these things shall be spoken, if it be God’s pleasure.

34. \emph{Finally, my brethren, rejoice in the Lord alway; again I will say, Rejoice: for your redemption hath drawn nigh}\footnote{Phil. iii. 1; and iv. 4; Luke xxi. 28.}, and the heavenly host of the Angels is waiting for your salvation. And there is now \emph{the voice of one crying in the wilderness, Prepare ye the way of the Lord}\footnote{Is. xl. 3.}; and the Prophet cries, \emph{Ho, ye that thirst, come ye to the water}\footnote{Ib. lv. 1.}; and immediately afterwards, \emph{Hearken unto me, and ye shall eat that which is good, and your soul shall delight itself in good things}\footnote{Ib. v. 2.}. And within a little while ye shall hear that excellent lesson which says, \emph{Shine, shine, O thou new Jerusalem; for thy light is come}\footnote{Is. lx. 1.}. Of this Jerusalem the prophet hath said, \emph{And afterwards thou shalt be called the city of righteousness, Zion, the faithful mother of cities}\footnote{Ib. i. 26.}; \emph{because of the law which went forth out of Zion, and the word of the Lord from Jerusalem}\footnote{Ib. ii. 3.}, which word has from hence been showered forth on the whole world. To her the Prophet also says concerning you, \emph{Lift up thine eyes round about, and behold thy children gathered together}\footnote{Ib. xlix. 18.}; and she answers, saying, \emph{Who are these that fly as a cloud, and as doves with their young ones to me}\footnote{Ib. lx. 8.}? (\emph{clouds} because of their spiritual nature, and \emph{doves}, from their purity). And again, she says, \emph{Who knoweth such things? or who hath seen it thus? did ever a land bring forth in one day? or was ever a nation born all at once? for as soon as Zion travailed, she brought forth her children}\footnote{Ib. lxvi. 8.}. And all things shall be filled with joy unspeakable because of the Lord who said, \emph{Behold, I create Jerusalem a rejoicing, and her people a joy}\footnote{Ib. lxv. 18.}.

35. And may these words be spoken now again over you also, \emph{Sing, O heavens, and be joyful, O earth}; and then; \emph{for the Lord hath had mercy on His people, and comforted the lowly of His people}\footnote{Ib. xlix. 13.}. And this shall come to pass through the loving-kindness of God, who says to you, \emph{Behold, I will blot out as a cloud thy transgressions, and as a thick cloud thy sins}\footnote{Is. xliv. 22.}. But ye who have been counted worthy of the name of Faithful (of whom it is written, \emph{Upon My servants shall be called a new name which shall be blessed on the earth}\footnote{Ib. lxv. 15.},) ye shall say with gladness, \emph{Blessed be the God and Father of our Lord Jesus Christ, who hath blessed us with every spiritual blessing in the heavenly places in Christ}\footnote{Eph. i. 3.}\emph{: in whom we have our redemption through His blood, the forgiveness of our sins, according to the riches of His grace, wherein He abounded towards us}\footnote{Ib. v. 7.}, and what follows; and again, \emph{But God being rich in mercy, for His great love wherewith He loved us, when we were dead through our trespasses, quickened us together with Christ}\footnote{Ib. ii. 4.}, and the rest. And again in like manner praise ye the Lord of all good things, saying, \emph{But when the kindness of God our Saviour, and His love towards man appeared, not by works of righteousness which we had done, but according to His mercy He saved us, through the washing of regeneration, and the renewing of the Holy Ghost, which He shed on us abundantly through Jesus Christ our Saviour, that being} \emph{justified by His grace, we might be made heirs, according to hope, of eternal life}\footnote{Tit. iii. 4.}. And may God Himself the Father of our Lord Jesus Christ, the Father of glory, \emph{give unto you a spirit of wisdom and revelation in the knowledge of Himself, the eyes of your understanding being enlightened}\footnote{Eph. i. 17, 18.}, and may He ever keep you in good works, and words, and thoughts; to Whom be glory, honour, and power, through our Lord Jesus Christ, with the Holy Ghost, now and ever, and unto all the endless ages of eternity. Amen\footnote{“At the end of this Lecture in the older of the Munich \textsc{mss}. there is the following addition: Many other Lectures were delivered year by year, both before Baptism and after the neophytes had been baptized. But these alone were taken down when spoken and written by some of the earnest students in the year 352 of the advent of our Lord and Saviour Jesus Christ. And in these you will find partly discussions of all the necessary doctrines of the Faith which ought to be known to men, and answers to the Greeks, and to those of the Circumcision, and to the Heresies, and the moral precepts of Christians of all kinds, by the grace of God. The year 352 according to the computation of the Greeks is the year 360 of the Christian era” (Rupp). \par The date at which the Lectures were delivered cannot possibly be so late as is here stated. See the section of the Introduction on the “Date.”}.

\xchapter{Lecture XIX. First Lecture on the Mysteries.}

Five Catechetical Lectures of the Same Author, to the Newly Baptized\footnote{This general title of the five following Lectures is omitted in many \textsc{mss.} “In Cod. Ottob. at the end of the special title of this first Mystagogic Lecture, after the words “to the end of the Epistle,” there follows the statement “Of the same author Cyril, and of John the Bishop” (Bened. Ed.). See Index, \emph{Authenticity}.}.

\textsc{With a Lesson from the} First General Epistle of Peter\textsc{, beginning at} \emph{Be sober, be vigilant}\textsc{, to the end of the Epistle.}

1. \textsc{I have} long been wishing, O true-born and dearly beloved children of the Church, to discourse to you concerning these spiritual and heavenly Mysteries; but since I well knew that seeing is far more persuasive than hearing, I waited for the present season; that finding you more open to the influence of my words from your present experience, I might lead you by the hand into the brighter and more fragrant meadow of the Paradise before us; especially as ye have been made fit to receive the more sacred Mysteries, after having been found worthy of divine and life-giving Baptism\footnote{This Lecture was delivered on the Monday after Easter in the Holy Sepulchre: see Cat. xviii. 33.}. Since therefore it remains to set before you a table of the more perfect instructions, let us now teach you these things exactly, that ye may know the effect\footnote{\textgreek{τὴν ἔμφασιν τὴν.…γεγενημένην} , is found in all the \textsc{mss.} “Nevertheless it would seem that we ought to read \textgreek{τῶν.…γεγενημένων}, which Grodecq either read or substituted” (Ben. Ed.). With the proposed reading the meaning would be—“the significance of the things done to you,” which agrees better with the meaning of \textgreek{ἔμφασις}.} wrought upon you on that evening of your baptism.

2. First ye entered into the vestibule\footnote{\textgreek{τὁν προαύλιον}, called below in § 11 “the outer chamber.” Cf. Procat. § 1, note 3. It appears from Tertullian, \emph{De Corona,} § 3, that the renunciation was made first in the Church, and afterwards in the Baptistery: “When we are going to enter the water, at that moment as well as just before in the Church under the hand of the President, we solemnly profess that we disown the devil, and his pomp, and his angels.”} of the Baptistery, and there facing towards the West ye listened to the command to stretch forth your hand, and as in the presence of Satan ye renounced him. Now ye must know that this figure is found in ancient history. For when Pharaoh, that most bitter and cruel tyrant, was oppressing the free and high-born people of the Hebrews, God sent Moses to bring them out of the evil bondage of the Egyptians. Then the door posts were anointed with the blood of a lamb, that the destroyer might flee from the houses which had the sign of the blood; and the Hebrew people was marvellously delivered. The enemy, however, after their rescue, \emph{pursued after them}\footnote{Ex. xiv. 9, 23.}, and saw the sea wondrously parted for them; nevertheless he went on, following close in their footsteps, and was all at once overwhelmed and engulphed in the Red Sea.

3. Now turn from the old to the new, from the figure to the reality. There we have Moses sent from God to Egypt; here, Christ, sent forth from His Father into the world: there, that Moses might lead forth an afflicted people out of Egypt; here, that Christ might rescue those who are oppressed in the world under sin: there, the blood of a lamb was the spell against\footnote{\textgreek{ἀποτρόπαιον}} the destroyer; here, the blood of the Lamb without blemish Jesus Christ is made the charm to scare\footnote{\textgreek{φυγαδευτήριον}, the word commonly used in the Septuagint for “a city of refuge.” But the Verb \textgreek{φυγαδεύω} is Transitive in 2 Macc. ix. 4, as well as in Xenophon and Demosthenes. The application of the blood of Christ in Baptism is represented by marking the sign of the Cross on the forehead. Compare the lines of Prudentius quoted by the Benedictine Editor: \par “Passio quæ nostram defendit sanguine frontem, \par Corporeamque domum signato collinit ore.”} evil spirits: there, the tyrant was pursuing that ancient people even to the sea; and here the daring and shameless spirit, the author of evil, was following thee even to the very streams of salvation. The tyrant of old was drowned in the sea; and this present one disappears in the water of salvation.

4. But nevertheless thou art bidden to say, with arm outstretched towards him as though he were present, “I renounce thee, Satan.” I wish also to say wherefore ye stand facing to the West; for it is necessary. Since the West is the region of sensible darkness, and he being darkness has his dominion also in darkness, therefore, looking with a symbolical meaning towards the West, ye renounce that dark and gloomy potentate. What then did each of you stand up and say? “I renounce thee, Satan,”—thou wicked and most cruel tyrant! meaning, “I fear thy might no longer; for that Christ hath overthrown, having partaken with me of flesh and blood, that through these He \emph{might by death destroy death}\footnote{Heb. ii. 14, 15.}, that I might not be made \emph{subject to bondage} for ever.” “I renounce thee,”—thou crafty and most subtle serpent. “I renounce thee,”—plotter as thou art, who under the guise of friendship didst contrive all disobedience, and work apostasy in our first parents. “I renounce thee, Satan,”—the artificer and abettor of all wickedness.

5. Then in a second sentence thou art taught to say, “and all thy works.” Now the works of Satan are all sin, which also thou must renounce;—just as one who has escaped a tyrant has surely escaped his weapons also. All sin therefore, of every kind, is included in the works of the devil. Only know this; that all that thou sayest, especially at that most thrilling hour, is written in God’s books; when therefore thou doest any thing contrary to these promises, thou shalt be judged as a \emph{transgressor}\footnote{Gal. ii. 18.}. Thou renouncest therefore the works of Satan; I mean, all deeds and thoughts which are contrary to reason.

6. Then thou sayest, “And all his pomp\footnote{Herod. II. 58: “The Egyptians were the first to introduce solemn assemblies (\textgreek{πανηγύρις}) and processions (\textgreek{πομπάς}).” At Rome the term “pompa” was applied especially to the procession with which the Ludi Circenses were opened and also to any grand ceremony or pageant.}.” Now the pomp of the devil is the madness of theatres\footnote{\textgreek{θεατρομανίαι}. Cf. Tertull. \emph{Apologet}. 38; “We renounce all your spectacles.…Among us nothing is ever said, or seen, or heard, which has anything in common with the madness of the Circus, the immodesty of the theatre, the atrocities of the arena, the useless exercises of the wrestling-ground.” He calls the theatre “that citadel of all impurities,” \emph{De Spectaculis}, c. 10, “immodesty’s peculiar abode,” c. 17, and gives a vivid description of the rage and fury of the Circus in c. 16.}, and horse-races, and hunting, and all such vanity: from which that holy man praying to be delivered says unto God, \emph{Turn away mine eyes from beholding vanity}\footnote{Ps. cxix. 37.}. Be not interested in the madness of the theatre, where thou wilt behold the wanton gestures of the players\footnote{\textgreek{μίμων}, the name either of a species of low comedy, “consisting more of gestures and mimicry than of spoken dialogue,” or of the persons who acted in them. Cyril’s description of the coarse and indecent character of the mimes is more than justified by the impartial testimony of Ovid, \emph{Trist}. ii. 497: \par “Quid si scripsissem mimos obscœna jocantes, \par Qui semper vetiti crimen amoris habent; \par In quibus assidue cultus procedit adulter, \par Verbaque dat stulto callida nupta viro. \par Nubilis hos virgo, matronaque, virque, puerque \par Spectat, et e magna parte Senatus adest. \par Nec satis incestis temerari vocibus aures; \par Assuescunt oculi multa pudenda pati.” \par A theatre is mentioned as one of the buildings erected by Hadrian in his new City Aelia Capitolina built on the site of Jerusalem; and that theatrical performances were continued in the time of Cyril we know from the accusation that in a time of famine he had sold one of the Church vestments, which was afterwards used upon the stage.}, carried on with mockeries and all unseemliness, and the frantic dancing of effeminate men\footnote{Lactantius, \emph{Epitome}, § 63: “Histrionici etiam impudici gestus, quibus infames fœminas imitantur, libidines, quæ saltando exprimunt, docent.”};—nor in the madness of them who in hunts\footnote{\textgreek{κυνηγεσίαις}, the so-called “venationes” of the Circus in which the “bestiarii” fought with wild beasts.} expose themselves to wild beasts, that they may pamper their miserable appetite; who, to serve their belly with meats, become themselves in reality meat for the belly of untamed beasts; and to speak justly, for the sake of their own god, their belly, they cast away their life headlong in single combats\footnote{The “bestiarii” were feasted in public on the day before their encounter with the beasts. See Tertull. \emph{Apologet}. § 42: “I do not recline in public at the feast of Bacchus, after the manner of the beast-fighters at their last banquet.” Ib. § 9: “Those also who dine on the flesh of wild beasts from the arena, who have keen appetites for bear and stag.” These latter, however, were chiefly the poor, to whom flesh was a rarity: Apuleius \emph{Metam}. iv. 14, quoted by Oehler.}. Shun also horse-races, that frantic and soul-subverting spectacle\footnote{\textgreek{ψυχὰς ἐκτραχήλιζον}, an allusion to the risk of a broken neck in the chariot-race. Tertull. \emph{de Spectaculis}, § 9: “Equestrianism was formerly practised in a simple way on horseback, and certainly its ordinary use was innocent: but when it was dragged into the games, it passed from a gift of God into the service of demons.” The presiding deity of the chariot-race was Poseidon (Hom. \emph{Il}. xxiii, 307; Pind. \emph{Ol}. i. 63; \emph{Pyth}. vi. 50; Soph. Œdip. Col. 712), and both this and the other shows of the Circus, and of the theatre, were connected with the worship of the gods of Greece and Rome, and therefore forbidden as idolatrous: “What high religious rites, what sacrifices precede, intervene, and follow, how many guilds, how many priesthoods, how many services are set astir” (Tert. \emph{de Spect}. § 7).}. For all these are the pomp of the devil.

7. Moreover, the things which are hung up at idol festivals\footnote{\textgreek{πανηγύρεσι}. The Panegyris was strictly a religious festival, but was commonly accompanied by a great fair or market, in which were sold not only such things as the worshippers might need for their offerings, e.g. frankincense, but also the flesh of the animals which had been sacrificed. Cf. \emph{Dictionary of Greek and Rom. Antiq}. “Panegyris.” Tertull. \emph{Apolog}. § 42: “We do not go to your spectacles: yet the articles that are sold there, if I need them, I shall obtain more readily at their proper places. We certainly buy no frankincense.”}, either meat or bread, or other such things polluted by the invocation of the unclean spirits, are reckoned in the pomp of the devil. For as the Bread and Wine of the Eucharist before the invocation of the Holy and Adorable Trinity were simple bread and wine, while after the invocation the Bread becomes the Body of Christ, and the Wine the Blood of Christ\footnote{Compare St. Paul’s argument against meats offered to idols, 1 Cor. x. 14–21: and on Cyril’s Eucharistic doctrine, see notes on Cat. xxii.}, so in like manner such meats belonging to the pomp of Satan, though in their own nature simple, become profane by the invocation of the evil spirit.

8. After this thou sayest, “and all thy service\footnote{The form of renunciation before Baptism is given in the \emph{Apostolic Constitutions}, VII. 41: “I renounce Satan, and his works, and his pomps, and his services, and his angels, and his inventions, and all things that are under him.” Cf. Tertull. \emph{De Spectaculis}, § 4: “When on entering the water, we make profession of the Christian faith in the words of its rule, we bear public testimony that we have renounced the devil, his pomp, and his angels.”}.” Now the service of the devil is prayer in idol temples; things done in honour of lifeless idols; the lighting of lamps\footnote{Herod. ii. 62: “At Sais, when the assembly takes place for the sacrifices (to Minerva, or Neith), there is one night on which the inhabitants all burn a multitude of lights in the open air round their houses.…These burn the whole night, and give to the festival the name of the Feast of Lamps (\textgreek{Λυχνοκαΐη}).”}, or burning of incense by fountains or rivers\footnote{Fountains and rivers had each its own deity or nymph, to whom sacrifices were offered, and incense burned.}, as some persons cheated by dreams or by evil spirits do [resort to this\footnote{\textgreek{ἐς τοῦτο διέβησαν}. These words are omitted in many \textsc{mss.}, and regarded by the Benedictine Editor as a spurious addition made to complete the construction. The words \textgreek{ἢ τοιαῦτα} at the end of the sentence are better omitted, as in several good \textsc{mss.}}], thinking to find a cure even for their bodily ailments. Go not after such things. The watching of birds, divination, omens, or amulets, or charms written on leaves, sorceries, or other evil arts\footnote{Cat. iv. 37; \emph{Apost. Const}. vi.: “Be not a diviner, for that leads to idolatry.…Thou shalt not use enchantments or purgations for thy child. Thou shalt not be a soothsayer nor a diviner by great or little birds. Nor shalt thou learn wicked arts; for all these things has the Law forbidden.” Deut. xviii. 10, 11.}, and all such things, are services of the devil; therefore shun them. For if after renouncing Satan and associating thyself with Christ\footnote{\emph{Apost. Const.} vii. 41: “And after his renunciation let him in his association (\textgreek{συντασσόμενος}) say, I associate myself with Christ.”}, thou fall under their influence, thou shalt find\footnote{\textgreek{πειραθήσῃ} (Cod. Mon. 1) is a better reading than \textgreek{πειρασθήσῃ}. Cf. Plat. \emph{Laches}, 188 E: \textgreek{τῶν ἔργων ἐπειράθην}.} the tyrant more bitter; perchance, because he treated thee of old as his own, and relieved thee from his hard bondage, but has now been greatly exasperated by thee; so thou wilt be bereaved of Christ, and have experience of the other. Hast thou not heard the old history which tells us of Lot and his daughters? Was not he himself saved with his daughters, when he had gained the mountain, while his wife became a pillar of salt, set up as a monument for ever, in remembrance of her depraved will and her turning back. Take heed therefore to thyself, and turn not again to \emph{what is behind}\footnote{Phil. iii. 13. On the pillar of salt, see \emph{Wisd}. x. 7: “Of whose wickedness even to this day the waste land that smoketh is a testimony,…and a standing pillar of salt is a monument of an unbelieving soul.” Joseph. \emph{Ant}. I. xi. 4: “Moreover I have seen it, for it remains even unto this day.” Bp. Lightfoot, \emph{Clem. Rom. Ep. ad Cor}. xi. remarks that the region abounds in pillars of salt, and “Mediæval and even modern travellers have delighted to identify one or other of these with Lot’s wife.”}, having put thine hand to the plough, and then turning back to the salt savour of this life’s doings; but escape to the mountain, to Jesus Christ, that \emph{stone hewn without hands}\footnote{Dan. ii. 35, 45.}, which has filled the world.

9. When therefore thou renouncest Satan, utterly breaking all thy covenant with him, that ancient \emph{league with hell}\footnote{Is. xxviii. 15.}, there is opened to thee the paradise of God, which He planted towards the East, whence for his transgression our first father was banished; and a symbol of this was thy turning from West to East, the place of light\footnote{Cf. S. Ambros. \emph{De Mysteriis}, c. ii. 7: “Ad orientem converteris; qui enim renunciat diabolo ad Christum convertitur:” “Where he plainly intimates.…that turning to the East was a symbol of their aversion from Satan and conversion unto Christ, that is, from darkness to light, from serving idols, to serve Him, who is the Sun of Righteousness and Fountain of Light” (Bingh. \emph{Ant}. xi. vii. 7).}. Then thou wert told to say, “I believe in the Father, and in the Son, and in the Holy Ghost, and in one Baptism of repentance\footnote{Cf. Didaché, vii. 1; Justin M. \emph{Apolog}. I. c. 61 A; Swainson, \emph{Creeds}, c. iii. on the short Baptismal Professions. “The writings of S. Cyprian distinctly tell us, that in his day the form of interrogation at Baptism was fixed and definite. He speaks of the “usitata et legitima verba interrogationis,”—and we know as distinctly that the interrogation included the words, “Dost thou believe in God the Father, in His Son Christ, in the Holy Spirit? Dost thou believe in remission of sins and eternal life through the Church?”}.” Of which things we spoke to thee at length in the former Lectures, as God’s grace allowed us.

10. Guarded therefore by these discourses, \emph{be sober}. \emph{For} our \emph{adversary the devil}, as was just now read, \emph{as a roaring lion, walketh about, seeking whom he may devour}\footnote{1 Pet. v. 9.}. But though in former times death was mighty and devoured, at the holy Laver of regeneration God has \emph{wiped away every tear from off all faces}\footnote{Is. xxv. 8; Rev. vii. 17.}. For thou shalt no more mourn, now that thou hast put off the old man; but thou shalt keep holy-day\footnote{\textgreek{πανηγυρίσεις}}, \emph{clothed in the garment of salvation}\footnote{Is. lxi. 10.}, even Jesus Christ.

11. And these things were done in the outer chamber. But if God will, when in the succeeding lectures on the Mysteries we have entered into the Holy of Holies\footnote{These words seem to imply that the Lectures on the Eucharist were to be delivered in the Holy Sepulchre, though the Mysteries themselves may be called metaphorically “the Holy of Holies.”}, we shall there know the symbolical meaning of the things which are there performed. Now to God the Father, with the Son and the Holy Ghost, be glory, and power, and majesty, forever and ever. Amen.

\xchapter{Lecture XX. On the Mysteries. II: Of Baptism.}

\textsc{Romans vi. 3–14}

Know ye not, that so many of us as were baptized into Jesus Christ, were baptized into His death? \&c..…for ye are not under the Law, but under grace.

1. \textsc{These} daily introductions into the Mysteries\footnote{\textgreek{μυσταγωγίαι}.}, and new instructions, which are the announcements of new truths, are profitable to us; and most of all to you, who have been renewed from an old state to a new. Therefore, I shall necessarily lay before you the sequel of yesterday’s Lecture, that ye may learn of what those things, which were done by you in the inner chamber\footnote{The renunciation and the profession of faith were made in the outer chamber or vestibule of the Baptistery.}, were symbolical.

2. As soon, then, as ye entered, ye put off your tunic; and this was an image of \emph{putting off the old man with his deeds}\footnote{Col. iii. 9.}. Having stripped yourselves, ye were naked; in this also imitating Christ, who was stripped naked on the Cross, and by His nakedness \emph{put off from Himself the principalities and powers, and openly triumphed over them on the tree}\footnote{Ib. ii. 15. Cyril’s use of this passage agrees best with the interpretation that Christ, having been clothed with the likeness of sinful flesh during His life on earth, submitted therein to the assaults of the powers of evil, but on the Cross threw off from Himself both it and them.}. For since the adverse powers made their lair in your members, ye may no longer wear that old garment; I do not at all mean this visible one, but the \emph{old man, which waxeth corrupt in the lusts of deceit}\footnote{Eph. iv. 22.}. May the soul which has once put him off, never again put him on, but say with the Spouse of Christ in the Song of Songs, \emph{I have put off my garment, how shall I put it on}\footnote{Cant. v. 3.}? O wondrous thing! ye were naked in the sight of all, and were not ashamed\footnote{See \emph{Dict. Christ. Antiq}. “Baptism,” § 48: \emph{The Unclothing of the Catechumens}: Bingh. \emph{Ant}. XI. xi. 1: All “persons were baptized naked, either in imitation of Adam in Paradise, or our Saviour upon the Cross, or to signify their putting off the body of sin, and the old man with his deeds.”}; for truly ye bore the likeness of the first-formed Adam, who was naked in the garden, and was not ashamed.

3. Then, when ye were stripped, ye were anointed with exorcised oil\footnote{Apost. Const. vii. 22: “But thou shalt beforehand anoint the person with holy oil (\textgreek{ἐλαίῳ}), and afterward baptize him with water, and in the conclusion shalt seal him with the ointment (\textgreek{μύρῳ}), that the anointing (\textgreek{χρῖσμα}) may be a participation of the Holy Spirit, and the water a symbol of the death, and the ointment the seal of the Covenants. But if there be neither oil nor ointment, water suffices both for anointing, and for a seal, and for a confession of Him who died, or indeed is dying with us.” The previous anointing “with oil sanctified by prayer” is mentioned in the \emph{Clementine Recognitions}, III. c. 67, and in the Pseudo-Justin, \emph{Quæstiones ad Orthodoxos}, Qu. 137. It was not however universal, and seems to have been unknown in Africa, not being mentioned by Clement of Alexandria (\emph{Pæd}. II. c. viii. \emph{On the use of ointments}), nor Tertullian, nor Augustine.}, from the very hairs of your head to your feet, and were made partakers of the good olive-tree, Jesus Christ. For ye were cut off from the wild olive-tree\footnote{On the significance of the wild olive-tree, see Irenæus, V. 10.}, and grafted into the good one, and were made to share the fatness of the true olive-tree. The exorcised oil therefore was a symbol of the participation of the fatness of Christ, being a charm to drive away every trace of hostile influence. For as the breathing of the saints, and the invocation of the Name of God, like fiercest flame, scorch and drive out evil spirits\footnote{See Index, “Exorcism.”}, so also this exorcised oil receives such virtue by the invocation of God and by prayer, as not only to burn and cleanse away the traces of sins, but also to chase away all the invisible powers of the evil one.

4. After these things, ye were led to the holy pool\footnote{\textgreek{κολυμβήθραν}. The pool or piscina was deep enough for total immersion, and large enough for many to be baptized at once. Cf. Bingh. \emph{Ant}. VIII. vii. 2; XI. xi. 2, 3. For engravings of the very ancient Baptisteries at Aquileia and Ravenna, shewing the form of the font or piscina, see \emph{Dict. Christian Ant}. “Baptistery.”} of Divine Baptism, as Christ was carried from the Cross to the Sepulchre which is before our eyes. And each of you was asked, whether he believed in the name of the Father, and of the Son, and of the Holy Ghost, and ye made that saving confession, and descended three times into the water, and ascended again; here also hinting by a symbol at the three days burial of Christ\footnote{The same significance is attributed to the trine immersion by many Fathers, but a different explanation is given by Tertullian (\emph{Adv. Praxean}, c. xxvi.): “Not once only, but three times, we are immersed into the several Persons at the mention of their several names.” Gregory of Nyssa (\emph{On the Baptism of Christ}, p. 520 in this Series) joins both reasons together: “By doing this thrice we represent for ourselves that grace of the Resurrection which was wrought in three days: and this we do, not receiving the Sacrament in silence, but while there are spoken over us the Names of the Three Sacred Persons on whom we believed, \&c.” Compare p. 529. Cf. \emph{Apost. Const} VIII. § 47, Can. 50: “If any Bishop or Presbyter does not perform the three immersions of one initiation, but one immersion made into the death of Christ, let him be deprived.” \par Milles in his note on this passage mentions that “this form of Baptism is still used in the Greek Church. See Eucholog. p. 355. Ed. Jac. Goar. and his notes p. 365.”}. For as our Saviour passed three days and three nights in the heart of the earth, so you also in your first ascent out of the water, represented the first day of Christ in the earth, and by your descent, the night; for as he who is in the night, no longer sees, but he who is in the day, remains in the light, so in the descent, as in the night, ye saw nothing, but in ascending again ye were as in the day. And at the self-same moment ye were both dying and being born; and that Water of salvation was at once your grave and your mother. And what Solomon spoke of others will suit you also; for he said, in that case, \emph{There is a time to bear and a time to die}\footnote{Eccles. iii. 2.}; but to you, in the reverse order, there was a time to die and a time to be born; and one and the same time effected both of these, and your birth went hand in hand with your death.

5. O strange and inconceivable thing! we did not really die, we were not really buried, we were not really crucified and raised again; but our imitation was in a figure, and our salvation in reality. Christ was actually crucified, and actually buried, and truly rose again; and all these things He has freely bestowed upon us, that we, sharing His sufferings by imitation, might gain salvation in reality. O surpassing loving-kindness! Christ received nails in His undefiled hands and feet, and suffered anguish; while on me without pain or toil by the fellowship of His suffering He freely bestows salvation.

6. Let no one then suppose that Baptism is merely the grace of remission of sins, or further, that of adoption; as John’s was a baptism\footnote{Tertullian (\emph{De Baptismo}, c. 10) denies that John’s Baptism availed for the remission of sins: “If repentance is a thing human, its baptism must necessarily be of the same nature: else if it had been celestial, it would have given both the Holy Spirit and the remission of sins.” Cyril’s doctrine is more in accordance with the language of the Fathers generally, and of St. Mark i. 4; Luke iii. 3.} conferring only remission of sins: whereas we know full well, that as it purges our sins, and ministers\footnote{\textgreek{πρόξενον}.} to us the gift of the Holy Ghost, so also it is the counterpart\footnote{\textgreek{ἀντίτυπον}. The “Antitype” is here the sign or memorial of that which is past, and no longer actually present: See note 6 on xxi. 1. Cf. Heb. ix. 24.} of the sufferings of Christ. For this cause Paul just now cried aloud and said, \emph{Or are ye ignorant that all we who were baptized into Christ Jesus, were baptized into His death? We were buried therefore with Him by baptism into His death}\footnote{Rom. vi. 3. In the following sentence several \textsc{mss.} have a different reading: “These things perhaps he said to some who were disposed to think that Baptism ministers remission of sins only, and not adoption, and that further it has not the fellowship, \&c.” Against this reading, approve by Milles, the Benedictine Editor argues that in Rom. vi. 3, 4, there is no reference to adoption, but only to the fellowship of Christ’s Passion, and that Cyril quotes the passage only to prove the latter, the gift of adoption being generally admitted, and therefore not in question.}. These words he spoke to some who were disposed to think that Baptism ministers to us the remission of sins, and adoption, but has not further the fellowship also, by representation, of Christ’s true sufferings.

7. In order therefore that we might learn, that whatsoever things Christ endured, \textsc{for us and for our salvation\footnote{This clause is contained in the Nicene Creed, and in that which was offered to the Council by Eusebius as the ancient Creed of Cæsarea. It probably formed part of the Creed of Jerusalem, though it is not found in the titles of the Lectures, nor specially explained.}} He suffered them in reality and not in appearance, and that we also are made partakers of His sufferings, Paul cried with all exactness of truth, \emph{For if we have been planted together with the likeness of His death, we shall be also with the likeness of His resurrection}. Well has he said, \emph{planted together}\footnote{Ib. vi. 5. Cyril gives the phrase “\emph{planted together}” a special application to those who had been baptized in the same place where Christ had been buried.}. For since the true Vine was planted in this place, we also by partaking in the Baptism of death have been \emph{planted together} with Him. And fix thy mind with much attention on the words of the Apostle. He said not, “For if we have been planted together with His death,” but, \emph{with the likeness of His death}. For in Christ’s case there was death in reality, for His soul was really separated from His body, and real burial, for His holy body was wrapt in pure linen; and everything happened really to Him; but in your case there was only a likeness of death and sufferings, whereas of salvation there was not a likeness but a reality.

8. Having been sufficiently instructed in these things, keep them, I beseech you, in your remembrance; that I also, unworthy though I be, may say of you, \emph{Now I love you}\footnote{1 Cor. xi. 2: \emph{Now I praise you, \&c}.}, \emph{because ye always remember me, and hold fast the traditions, which I delivered unto you}. And God, who has presented you \emph{as it were alive from the dead}\footnote{Rom. vi. 13.}, is able to grant unto you \emph{to walk in newness of life}\footnote{Ib. v. 4.}: because His is the glory and the power, now and for ever. Amen.

\xchapter{Lecture XXI. On the Mysteries. III: On Chrism.}

\textsc{1 John ii. 20–28}

But ye have an unction from the Holy One, \&c..…that, when He shall appear, we may have confidence, and not be ashamed before Him at His coming.

1. \textsc{Having} been \emph{baptized into Christ}, and \emph{put on Christ\footnote{Gal. iii. 27.}}, ye have been made conformable to the Son of God; for God having \emph{foreordained us unto adoption as sons\footnote{Eph. i. 5.}}, made us \emph{to be conformed to the body of Christ’s glory}\footnote{Phil. iii. 21.}. Having therefore become \emph{partakers of Christ}\footnote{Heb. iii. 14.}, ye are properly called Christs, and of you God said, \emph{Touch not My Christs}\footnote{Ps. cv. 15.}, or anointed. Now ye have been made Christs, by receiving the antitype\footnote{\textgreek{ἀντίτυπον}. Cat. xx. 6; xxiii. 20. Twice in this section as in Heb. ix. 24 (\textgreek{ἀντίτυπα τῶν ἀληθινῶν}), \textgreek{ἀντίτυπον} is the copy or figure representing the original pattern (\textgreek{τύπος}, cf. Acts vii. 44). Otherwise (as in Cat. x. 11; xiii. 19; xxii. 3) \textgreek{τύπος} is the figure to be subsequently realised in the antitype.} of the Holy Ghost; and all things have been wrought in you by imitation\footnote{\textgreek{εἰκονικῶς .…εἰκόνες τοῦ Χριστοῦ}.}, because ye are images of Christ. He washed in the river Jordan, and having imparted of the fragrance\footnote{\textgreek{χρώτων}, literally “tinctures.” The Ben. Ed. writes: “For \textgreek{φώτων} we have written \textgreek{χρώτων} with Codd. Coisl. Ottob. Roe, Casaub., \&c…But we must write \textgreek{χρώτων} from \textgreek{χρῶτα}, not \textgreek{χρῶτων} from \textgreek{χρῶτες}. Authors use the word \textgreek{χρῶτα} to signify the effluence of an odour. So Gregory of Nyssa takes it in his 3rd \emph{Homily on the Song of Songs}, p. 512; and S. Maximus in \emph{Question} 37 \emph{on Scripture}: ‘\textgreek{χρῶτα} we say is the godliness (\textgreek{εὐσέβειαν}) whereby S. Paul was \emph{to the one a savour of life unto life.}’…In the \emph{Procatechesis}, § 15, Cyril calls the waters of Baptism \textgreek{ὑδάτων χριστοφόρων ἐχόντων εὐωδίαν}. If however any one prefers the reading \textgreek{φώτων}, he may defend himself by the authority of Epiphanius, who in the \emph{Exposition of the Faith}, c. 15, says that Christ descending into the water gave rather than received,.…illuminating them, and empowering them for a type of what was to be accomplished in Him.” According to the Ebionite Gospel of St. Matthew in Epiphanius (\emph{Hær}. xxx. \emph{Ebionitæ}. c. 13), when Jesus came up out of the water a great light shone around the place: a tradition to which the Benedictine Editor thinks the reading \textgreek{φώτων} may refer. Justin M. (\emph{Dialog}. c. lxxxviii.): “When Jesus had stepped into the water, a fire was kindled in the Jordan.” Otto quotes the legend, as found in \emph{Orac. Sibyll}. vii. 81–83:— \par \textgreek{῞Ος σε Λόγον γέννησε Πατήρ Πνεῦμ᾽ ὄρνιν ἄφηκεν,} \par \textgreek{᾽Οξὺν ἀπαγγελτῆρα λόγων, Λόγον ὕδασιν ἁγνοῖς} \par \textgreek{῾Ραίνων, σὸν Βάπτισμα δι᾽ οὗ πυρὸς ἐξεφαάνθης .}} of His Godhead to the waters, He came up from them; and the Holy Ghost in the fulness of His being\footnote{\textgreek{οὐσιώδης ἐπιφοίησις ἐγένετο}. The Benedictine Editor understands this phrase as an allusion to the descent of the Holy Ghost on Jesus in a substantial bodily form. So Gregory Nazianzen (\emph{Orat} xliv. 17), says that the Holy Ghost descended on the Apostles \textgreek{οὐσιωδῶς καὶ σωματικῶς}. But Anastasius Sinaita interprets \textgreek{οὐσιωδῶς} in this latter passage as meaning “in the essence and reality of His (Divine ) Person:” and this latter sense agreeing with the frequent use of \textgreek{οὐσιωδής} by Athanasius is well rendered by Canon Mason (\emph{The Relation of Confirmation to Baptism}, p. 343, “in the fulness of His being.”} lighted on Him, like resting upon like\footnote{Cf. Greg. Naz. \emph{Orat}. xxxix: “The Sprit also bears witness to His Godhead, for he comes to that which is like Himself.”}. And to you in like manner, after you had come up from the pool of the sacred streams, there was given an Unction\footnote{Cf. Tertullian, \emph{De Baptismo}, c. 7: “Exinde egressi de lavacro perungimur benedictâ unctione.” It is clear that the Unction mentioned in these passages was conferred at the same time and place as Baptism. Whether it formed part of that Sacrament, or was regarded by Cyril as a separate and independent rite, has been made a matter of controversy. See Index, “Chrism.”}, the anti-type of that wherewith Christ was anointed; and this is the Holy Ghost; of whom also the blessed Esaias, in his prophecy respecting Him, said in the person of the Lord, \emph{The Spirit of the Lord is upon Me, because He hath anointed Me: He hath sent Me to preach glad tidings to the poor}\footnote{Is. lxi. 1.}.

2. For Christ was not anointed by men with oil or material ointment, but the Father having before appointed Him to be the Saviour of the whole world, anointed Him with the Holy Ghost, as Peter says, \emph{Jesus of Nazareth, whom God anointed with the Holy Ghost}\footnote{Acts x. 38.}. David also the Prophet cried, saying, \emph{Thy throne, O God, is for ever and ever; a sceptre of righteousness is the sceptre of Thy kingdom; Thou hast loved righteousness and hated iniquity; therefore God even Thy God hath anointed Thee with the oil of gladness above Thy fellows}\footnote{Ps. xlv. 6, 7.}. And as Christ was in reality crucified, and buried, and raised, and you are in Baptism accounted worthy of being crucified, buried, and raised together with Him in a likeness, so is it with the unction also. As He was anointed with an ideal\footnote{\textgreek{νοητῷ} cannot here be translated “spiritual” because of \textgreek{πνευματίκῆς} immediately following. Cf. i. 4, note.} oil of gladness, that is, with the Holy Ghost, called oil of gladness, because He is the author of spiritual gladness, so ye were anointed with ointment, having been made partakers and \emph{fellows of Christ}.

3. But beware of supposing this to be plain ointment. For as the Bread of the Eucharist, after the invocation of the Holy Ghost, is mere bread no longer\footnote{Compare xix. 7; xxiii. 7, 19; and the section on “\emph{Eucharist}” in the Introduction.}, but the Body of Christ, so also this holy ointment is no more simple ointment, nor (so to say) common, after invocation, but it is Christ’s gift of grace, and, by the advent of the Holy Ghost, is made fit to impart His Divine Nature\footnote{\textgreek{Χριστοῦ χάρισμα καὶ Πνεύματος ἁγίου παρουσίᾳ τῆς αὐτοῦ Θεότητος ἐνεργητικὸν γινόμενον}. The meaning of this passage seems to have been obscured by divergent views of the order and construction of the words. In the Oxford translation, followed by Dr. Pusey (\emph{Real Presence}, p. 357), the Chrism is “the gift of Christ, and by the presence of His godhead it causes in us the Holy Ghost.” The order of the operations proper to the two Divine Persons seems thus to be inverted. \par According to the Benedictine Editor, and Canon Mason (\emph{Relation of Confirmation to Baptism}, p. 344), it is “Christ’s gracious gift, and is made effectual to convey the Holy Ghost by the presence of His own Godhead,”—i.e. apparently, the Godhead of the Holy Ghost conveys the Holy Ghost. \par But according to the context “the presence” must be that of the Divine Person who has been invoked, namely the Holy Ghost: and this is clearly expressed in the order of the words \textgreek{Πνεύματος ἁγίου παρουσίᾳ τῆς αὐτοῦ θεότητος ἐνεργητικόν}. The connexion of the words \textgreek{Πν. ἁγ. παρουσίᾳ} is put beyond doubt by the Invocation in the Liturgy of S. James quoted in \emph{Myst}. V. 7, note 8. The true meaning thus seems to be that the Chrism is Christ’s gift of grace, and imparts His Divine nature by the presence of the Holy Ghost after the Invocation. This meaning is confirmed by the formula given in \emph{Apost. Const}. vii. 44, for the consecration of the Chrism: “Grant also now that this ointment may be made effectual in the baptized, that the sweet savour of Thy Christ may remain firm and stable in him, and that, having died with Him, he may rise again and live with Him.” The Chrism is thus regarded as “the Seal” which confirms the proper benefits of Baptism.}. Which ointment is symbolically applied to thy forehead and thy other senses\footnote{\textgreek{ἐπὶ μετώπου καὶ τῶν ἄλλων σου αἰσθητηρίων}. The forehead may be regarded as representing the sense of touch; or we may translate, according to the idiomatic use of \textgreek{ἄλλος}, “thy forehead and thine organs of sense besides.” See Winer, \emph{Grammar of N.T. Greek}, P. III. Sect. lix. 7: Riddell, \emph{Digest of Platonic Idioms}, § 46.}; and while thy body is anointed with the visible ointment, thy soul is sanctified by the Holy and life-giving Spirit.

4. And ye were first anointed on the forehead, that ye might be delivered from the shame, which the first man who transgressed bore about with him everywhere; and that \emph{with unveiled face ye might reflect as a mirror the glory of the Lord}\footnote{2 Cor. iii. 18.}. Then on your ears; that ye might receive the ears which are quick to hear the Divine Mysteries, of which Esaias said, \emph{The Lord gave me also an ear to hear}\footnote{Is. l. 4.}; and the Lord Jesus in the Gospel, \emph{He that hath ears to hear let him hear}\footnote{Matt. xi. 15.}. Then on the nostrils; that receiving the sacred ointment ye may say, \emph{We are to God a sweet savour of Christ, in them that are saved}\footnote{2 Cor. ii. 15.}. Afterwards on your breast; that having put on the \emph{breast-plate of righteousness}, ye may \emph{stand against the wiles of the devil}\footnote{Eph. vi. 14, and 11.}. For as Christ after His Baptism, and the visitation of the Holy Ghost, went forth and vanquished the adversary, so likewise ye, after Holy Baptism and the Mystical Chrism, having put on the whole armour of the Holy Ghost, are to stand against the power of the adversary, and vanquish it, saying, \emph{I can do all things through Christ which strengtheneth me}\footnote{Phil. iv. 13.}.

5. Having been counted worthy of this Holy Chrism, ye are called Christians, verifying the name also by your new birth. For before you were deemed worthy of this grace, ye had properly no right to this title, but were advancing on your way towards being Christians.

6. Moreover, you should know that in the old Scripture there lies the symbol of this Chrism. For what time Moses imparted to his brother the command of God, and made him High-priest, after bathing in water, he anointed him; and Aaron was called Christ or Anointed, evidently from the typical Chrism. So also the High-priest, in advancing Solomon to the kingdom, anointed him after he had bathed in Gihon\footnote{1 Kings i. 39.}. To them however these things happened in a figure, but to you not in a figure, but in truth; because ye were truly anointed by the Holy Ghost. Christ is the beginning of your salvation; for He is truly the First-fruit, and ye the mass\footnote{Rom. xi. 16.}; but if the First-fruit be holy, it is manifest that Its holiness will pass to the mass also.

7. Keep This unspotted: for it shall teach you all things, if it abide in you, as you have just heard declared by the blessed John, discoursing much concerning this Unction\footnote{1 John ii. 20: \emph{But ye have an unction} (\textgreek{χρῖσμα}) \emph{from the Holy One}.}. For this holy thing is a spiritual safeguard of the body, and salvation of the soul. Of this the blessed Esaias prophesying of old time said, \emph{And on this mountain},—(now he calls the Church a mountain elsewhere also, as when he says, \emph{In the last days the mountain of the Lord’s house shall be manifest}\footnote{Is. ii. 2.};)—\emph{on this mountain shall the Lord make unto all nations a feast; they shall drink wine, they shall drink gladness, they shall anoint themselves with ointment}\footnote{Ib. xxv. 6. The Septuagint differs much from the Hebrew, both here and in the following verse. \textsc{R.V.} “And in this mountain shall the Lord of host make unto all people a feast of fat things, a feast of wines on the lees, of fat things full of marrow, of wines on the lees well refined.”}. And that he may make thee sure, hear what he says of this ointment as being mystical; \emph{Deliver all these things to the nations, for the counsel of the Lord is unto all nations}\footnote{Ib. v. 7. R.V. “And He will destroy in this mountain the face of the covering that is cast over all peoples, and the veil that is spread over all nations.”}. Having been anointed, therefore, with this holy ointment, keep it unspotted and unblemished in you, pressing forward by good works, and being made well-pleasing to the Captain of your salvation, Christ Jesus, to whom be glory for ever and ever. Amen.

\xchapter{Lecture XXII. On the Mysteries. IV: On the Body and Blood of Christ.}

\textsc{1 Cor. xi. 23}

I received of the Lord that which also I delivered unto you, how that the Lord Jesus, in the night in which He was betrayed, took bread, \&c.

1. \textsc{Even} of itself\footnote{\textgreek{αὐτή} found in all \textsc{mss.} is changed for the worse into \textgreek{αὕτη} by the Benedictine Editor.} the teaching of the Blessed Paul is sufficient to give you a full assurance concerning those Divine Mysteries, of which having been deemed worthy, ye are become of \emph{the same body}\footnote{Introduction, “\emph{Eucharist}.” The word \textgreek{σύσσωμοι} has a different sense in Eph. iii. 6, where it is applied to the Gentiles as having been made members of Christ’s body the Church.} and blood with Christ. For you have just heard him say distinctly, \emph{That our Lord Jesus Christ in the night in which He was betrayed, took bread, and when He had given thanks He brake it, and gave to His disciples, saying, Take, eat, this is My Body: and having taken the cup and given thanks, He said, Take, drink, this is My Blood}\footnote{1 Cor. xi. 23. The clause “and gave to His disciples” is an addition taken from Matt. xxvi. 26. The part relating to the cup does not correspond exactly either with St. Paul’s language or with the Evangelists’.}. Since then He Himself declared and said of the Bread, \emph{This is My Body}, who shall dare to doubt any longer? And since He has Himself affirmed and said, \emph{This is My Blood}, who shall ever hesitate, saying, that it is not His blood?

2. He once in Cana of Galilee, turned the water into wine, akin to blood\footnote{\textgreek{οἰκεῖον αἵματι}. Cod. Scirlet. (Grodecq), Mesm. (Morel), Vindob.; Ben. Ed. \textgreek{οἰκείῳ νεύματι}, Codd. Monac. 1, 2, Genovef. Vatt. (Prevot.). Rupp. The whole passage is omitted in Codd. Coisl. R. Casaub. owing to the repetition of \textgreek{αἷμα} \par The reading \textgreek{οἰκείῳ νεύματι}, “by His own will,” introduces a superfluous thought, and destroys the very point of Cyril’s argument, in which the previous change of water into an element so different as wine is regarded as giving an \emph{a fortiori} probability to the change of that which is already “akin to blood” into blood itself. \par If Cyril thus seems to teach a physical change of the wine, it must be remembered that we are not bound to accept his view, but only to state it accurately. See however the section of the Introduction on his Eucharistic doctrine.}, and is it incredible that He should have turned wine into blood? When called to a bodily marriage, He miraculously wrought\footnote{\textgreek{ἐθαυματούργησε τὴν παραδοξοποιίαν}. Cf. Chrysost. \emph{Epist.} I. \emph{ad Olympiad. de Deo}, § 1, c.: \textgreek{τότε θαυματουργεῖ καὶ παραδοξοποιεῖ}.} that wonderful work; and \emph{on the children of the bride-chamber}\footnote{Matt. ix. 15.}, shall He not much rather be acknowledged to have bestowed the fruition of His Body and Blood\footnote{Ben. Ed.: “That the force of Cyril’s argument may be the better understood, we must observe that in Baptism is celebrated the marriage of Christ with the Christian soul; and that the consummation of this marriage is perfected through the union of bodies in the mystery of the Eucharist. Read Chrysostom’s \emph{Hom}. xx. \emph{in Ephes}.” Chrysostom’s words are: “In like manner therefore we become one flesh with Christ by participation (\textgreek{μετουσίας}).” But the participation expressed by \textgreek{μετουσία} does not necessarily refer to the Eucharist. From the use of the word in Cat. xxiii. 11, and in Athanasius (\emph{Contra Arianos, Or}. i.; \emph{de Synodis}. 19, 22, 25) the meaning rather seems to be that we are one flesh with Christ not by nature but by His gift.}?

3. Wherefore with full assurance let us partake as of the Body and Blood of Christ: for in the figure\footnote{See Index, \textgreek{Τύπος}, and the references there, and Waterland, \emph{On the Eucharist}, c. vii.} of Bread is given to thee His Body, and in the figure of Wine His Blood; that thou by partaking of the Body and Blood of Christ, mayest be made of the same body and the same blood with Him. For thus we come to bear Christ\footnote{\textgreek{Χριστοφόροι γινόμεθα}. Procat. 15.} in us, because His Body and Blood are distributed\footnote{Ben. Ed.: “\textgreek{᾽Αναδιδομένου}. The Codices Coisl. Roe, Casaub. Scirlet. Ottob. 2. Genovef. have \textgreek{ἀναδεδεγμένοι}, which does not agree well with the Genitives \textgreek{τοῦ σώματος} and \textgreek{τοῦ αἵματος}. It is evident that it was an ill-contrived emendation of \textgreek{ἀναδιδομένου}, the transcribers being offended at the distribution of Christ’s Body among our members. But Cyril uses even the same word in Cat. xxiii. 9: \textgreek{Οὗτος ὁ ἄρτος.…εἰς πᾶσάν σου τὴν σύστασιν ἀναδίδοται, εἰς ὠφέλειαν σώματος καὶ ψυχῆς}, ‘This Bread is distributed into thy whole system, to the benefit of body and soul.’” \textgreek{᾽Αναδιδομένου} is the reading of Milles and Rupp. For similar language see Justin M. \emph{Apol}. i. 66; Iren. V. ii. 2.} through our members; thus it is that, according to the blessed Peter, \emph{we become partakers of the divine nature}\footnote{2 Pet. i. 4.}.

4. Christ on a certain occasion discoursing with the Jews said, \emph{Except ye eat My flesh and drink My blood, ye have no life in you}\footnote{John vi. 53.}. They not having heard His saying in a spiritual sense were offended, and went back, supposing that He was inviting them to eat flesh.

5. In the Old Testament also there was shew-bread; but this, as it belonged to the Old Testament, has come to an end; but in the New Testament there is Bread of heaven, and a Cup of salvation, sanctifying soul and body; for as the Bread corresponds to our body, so is the Word\footnote{Ben. Ed.: “Here we are to understand (by \textgreek{ὁ Λόγος}) the Divine Word, not the bare discourse of God, but the second Person of the Holy Trinity, Christ Himself, the Bread of Heaven, as He testifies of Himself, John vi. 51: Him Cyril contrasts with the earthly shew-bread in the O.T.; otherwise he could not rightly from this sentence infer, by the particle \textgreek{οὖν}, “therefore,” that the Eucharist is truly the Body and Blood of Christ. And since he says, in Cat. xxiii. 15, that the Eucharistic food is “appointed for the substance of the soul,” for its benefit, that cannot be said of Christ’s body or of His soul, but only of the Word which is conjoined with both. Moreover that the Divine Word is the food of Angels and of the soul, is a common mode of speaking with all the Fathers. They often play on the ambiguity of this word (\textgreek{λόγος}), saying sometimes that the Divine Word, sometimes the word and oracles of God, are the food of our souls: both statements are true. For the whole life-giving power of the Eucharist is derived from the Word of God united to the flesh which He assumed: and the whole benefit of Eucharistic eating consists in the union of our soul with the Word, in meditation on His mysteries and sayings, and conformity thereto.”} appropriate to our soul.

6. Consider therefore the Bread and the Wine not as bare elements, for they are, according to the Lord’s declaration, the Body and Blood of Christ; for even though sense suggests this to thee, yet let faith establish thee. Judge not the matter from the taste, but from faith be fully assured without misgiving, that the Body and Blood of Christ have been vouchsafed to thee.

7. Also the blessed David shall advise thee the meaning of this, saying, \emph{Thou hast prepared a table before me in the presence of them that afflict me}\footnote{Ps. xxiii. 5.}. What he says, is to this effect: Before Thy coming, the evil spirits prepared a table for men\footnote{\textgreek{ἠλισγημένην}, a good restoration by Milles, with Codd. Roe, Casaub. Coislin. The earlier printed texts had \textgreek{ἠλυγισμένην}, “overshadowed.” Cf. Mal. i. 7: \textgreek{ἄρτους ἠλισγημένους ,}.…\textgreek{Τράπεζα Κυρίου ἠλισγημένη ἐστίν}}, polluted and defiled and full of devilish influence\footnote{Cyril refers to the idolatrous feasts, which St. Paul calls “the table of devils,” 1 Cor. x. 21.}; but since Thy coming. O Lord, \emph{Thou hast prepared a table before me}. When the man says to God, \emph{Thou hast prepared before me a table}, what other does he indicate but that mystical and spiritual Table, which God hath prepared for us \emph{over against}, that is, contrary and in opposition to the evil spirits? And very truly; for that had communion with devils, but this, with God. \emph{Thou hast anointed my head with oil}\footnote{Ps. xxiii. 5.}. With oil He anointed thine head upon thy forehead, for the seal which thou hast of God; that thou mayest be made \emph{the engraving of the signet, Holiness unto God}\footnote{Ex. xxviii. 36; Ecclus. xlv. 12. The plate of pure gold on the forefront of Aaron’s mitre was engraved with the motto, \emph{Holy unto the Lord}. This symbolism Cyril transfers to the sacramental Chrism, in which the forehead is signed with ointment, and the soul with the seal of God.}. And \emph{thy cup intoxicateth me, as very strong}\footnote{Ps. xxiii. 5: \emph{My cup runneth over}. Eusebius (\emph{Dem. Evang}. I. c. 10, § 28) applies the Psalm, as Cyril does, to the Eucharist.}. Thou seest that cup here spoken of, which Jesus took in His hands, and gave thanks, and said, \emph{This is My blood, which is shed for many for the remission of sins}\footnote{Matt. xxvi. 28.}.

8. Therefore Solomon also, hinting at this grace, says in Ecclesiastes, \emph{Come hither, eat thy bread with joy} (that is, the spiritual bread; \emph{Come hither}, he calls with the call to salvation and blessing), \emph{and drink thy wine with a merry heart} (that is, the spiritual wine); \emph{and let oil be poured out upon thy head} (thou seest he alludes even to the mystic Chrism); \emph{and let thy garments be always white, for the Lord is well pleased with thy works}\footnote{Eccles. ix. 7, 8.}; for before thou camest to Baptism, thy works were \emph{vanity of vanities}\footnote{For \textgreek{προσέλθῃς} (Bened.) we must read \textgreek{προσῆλθες}, or, with Monac. 1, \textgreek{προσελθεῖν}.}. But now, having put off thy old garments, and put on those which are spiritually white, thou must be continually robed in white: of course we mean not this, that thou art always to wear white raiment; but thou must be clad in the garments that are truly white and shining and spiritual, that thou mayest say with the blessed Esaias, \emph{My soul shall be joyful in my God; for He hath clothed me with a garment of salvation, and put a robe of gladness around me}\footnote{Is. lxi. 10.}.

9. Having learnt these things, and been fully assured that the seeming bread is not bread, though sensible to taste, but the Body of Christ; and that the seeming wine is not wine, though the taste will have it so, but the Blood of Christ\footnote{On this passage see the section of the Introduction referred to in the Index, “\emph{Eucharist}.”}; and that of this David sung of old, saying, \emph{And bread strengtheneth man’s heart, to make his face to shine with oil}\footnote{Ps. civ. 15.}, “strengthen thou thine heart,” by partaking thereof as spiritual, and “make the face of thy soul to shine.” And so having it unveiled with a pure conscience, mayest thou \emph{reflect as a mirror the glory of the Lord}\footnote{2 Cor. iii. 18.}, and proceed from \emph{glory to glory}, in Christ Jesus our Lord:—To whom be honour, and might, and glory, for ever and ever. Amen.

\xchapter{Lecture XXIII. On the Mysteries. V: On the Sacred Liturgy and Communion.}

On the Sacred Liturgy and Communion\footnote{This title is added by the Benedictine Editor. There is nothing corresponding to it in the Greek.}.

\textsc{1 Pet. ii. 1}

\emph{Wherefore putting away all filthiness, and all guile, and evil speaking}\footnote{The text is made up from memory of James i. 21: \textgreek{διὸ ἀποθέμενοι πᾶσαν ῥυπαρίαν}, and 1 Pet. ii. 1: \textgreek{ἀποθέμενοι οὖν πᾶσαν κακίαν καὶ πάντα δόλον καὶ ὑποκρίσεις καὶ πάσας καταλαλίας}.}\emph{, \&c.}

1. \textsc{By} the loving-kindness of God ye have heard sufficiently at our former meetings concerning Baptism, and Chrism, and partaking of the Body and Blood of Christ; and now it is necessary to pass on to what is next in order, meaning to-day to set the crown on the spiritual building of your edification.

2. Ye have seen then the Deacon who gives to the Priest water to wash\footnote{In the \emph{Apostolic Constitutions}, VIII. xi, this duty is assigned to a sub-deacon: “Let one of the sub-deacons bring water to wash the hands of the priests, which is a symbol of the purity of those souls that are devoted to God.” See \emph{Dictionary of Christian Antiquities,} “Lavabo.” The Priest who celebrates the Eucharist is here distinguished by the title \textgreek{ἱερεύς} from the other Presbyters who stood round the altar.}, and to the Presbyters who stand round God’s altar. He gave it not at all because of bodily defilement; it is not that; for we did not enter the Church at first\footnote{Cyril evidently refers to the custom of placing vessels of water outside the entrance of the Church. Bingham, \emph{Antiquities}, VIII. iii. 6. Chrysost. \emph{In Johannem Hom}. lxxiii. 3: “Do we then wash our hands when going into Church, and shall we not wash our hearts also?” That the same custom was observed in heathen Temples appears from Herod. I. 51: \textgreek{περιῤῥαντήρια δύο ἀνέθηκε} (See Bähr’s note). Compare also Joseph. \emph{Ant. Jud}. III. vi. 2.} with defiled bodies. But the washing of hands is a symbol that ye ought to be pure from all sinful and unlawful deeds; for since the hands are a symbol of action, by washing\footnote{\textgreek{[τῷ] νίψασθαι}. Rupp: “\textgreek{Τῷ} ex conjectura addidi.” Possibly the original reading was \textgreek{νιψάμενοι}, which would easily become altered through the presence of \textgreek{νιψασθαι} in the preceding line. This washing is not mentioned in the Liturgy of St. James.} them, it is evident, we represent the purity and blamelessness of our conduct. Didst thou not hear the blessed David opening this very mystery, and saying, \emph{I will wash my hands in innocency, and so will compass Thine Altar, O Lord}\footnote{Ps. xxvi. 6. In the Liturgy of Constantinople this Psalm was chanted by the Priest and Deacon while washing their hands at the Prothesis or Credence.}? The washing therefore of hands is a symbol of immunity\footnote{\textgreek{ἀνυπεύθυνος}.} from sin.

3. Then the Deacon cries aloud, “Receive ye one another; and let us kiss one another\footnote{These two directions by the Deacon are separated in the Liturgy of St. James: after the dismissal of the Catechumens, the Deacon says, “Take note one of another;” and after the Incense, Cherubic hymn, Oblation, Creed, and a short prayer “that we may be united one to another in the bond of peace and charity,” the Deacon says, “Let us salute (\textgreek{ἀγαπῶμεν}) one another with a holy kiss.” In the \emph{Apostolic Constitutions}, VIII. 11, there is but one such direction, and this comes before the washing of hands and the dismissal of the Catechumens, “Salute (\textgreek{ἀσπάσασθε}) ye one another with a holy kiss.”}.” Think not that this kiss is of the same character with those given in public by common friends. It is not such: but this kiss blends souls one with another, and courts entire forgiveness for them. The kiss therefore is the sign that our souls are mingled together, and banish all remembrance of wrongs. For this cause Christ said, \emph{If thou art offering thy gift at the altar, and there rememberest that thy brother hath aught against time, leave there thy gift upon the altar, and go thy way; first be reconciled to thy brother, and then come and offer thy gift}\footnote{Matt. v. 23. From Cyril’s reference to this passage “it may be inferred that the kiss of peace had been given before the gifts were brought to the altar, according to ancient custom attested by Justin M. \emph{Apolog}. i. c. 65: ‘Having ended the prayers’ (for the newly baptized) ‘we salute one another with a kiss. Then there is brought to the President of the brethren bread, and a cup of wine mixed with water’” (Ben. Ed.). There is the same order in the \emph{Apost. Const.} VIII. 12, and in the 19th Canon of the Synod of Laodicea; but in the Liturgy of S. James the gifts are offered before the kiss of peace.}. The kiss therefore is reconciliation, and for this reason holy: as the blessed Paul somewhere cried, saying, \emph{Greet ye one another with a holy kiss}\footnote{1 Cor. xvi. 20.}; and Peter, \emph{with a kiss of charity}\footnote{1 Pet. iii. 15.}.

4. After this the Priest cries aloud, “Lift up your hearts\footnote{The words are slightly varied in the Liturgies: thus in the Liturgy of St. James, “Let us lift up our mind and hearts;” in the \emph{Apost. Const}. viii. 12, “Lift up your mind.”}.” For truly ought we in that most awful hour to have our heart on high with God, and not below, thinking of earth and earthly things. In effect therefore the Priest bids all in that hour to dismiss all cares of this life, or household anxieties, and to have their heart in heaven with the merciful God. Then ye answer, “We lift them up unto the Lord:” assenting to it, by your avowal. But let no one come here, who could say with his mouth, “We lift up our hearts unto the Lord,” but in his thoughts have his mind concerned with the cares of this life. At all times, rather, God should be in our memory but if this is impossible by reason of human infirmity, in that hour above all this should be our earnest endeavour.

5. Then the Priest says, “Let us give thanks unto the Lord.” For verily we are bound to give thanks, that He called us, unworthy as we were, to so great grace; that He reconciled us when we were His foes; that He vouchsafed to us the Spirit of adoption. Then ye say, “It is meet and right:” for in giving thanks we do a meet thing and a right; but He did not right, but more than right, in doing us good, and counting us meet for such great benefits.

6. After this, we make mention of heaven, and earth, and sea\footnote{Compare the noble Eucharistic Preface in the Liturgy of St. James: “It is verily meet, right, becoming, and our bounden duty to praise Thee, to sing of Thee, to bless Thee, to worship Thee, to glorify Thee, to give thanks to Thee the Maker of every creature, visible and invisible, the Treasure of eternal blessings; the Fount of life and immortality, the God and Lord of all, whom the heavens of heavens do praise, and all the powers thereof, sun and moon and all the choir of the stars, earth, sea, and all that in them is, Jerusalem the heavenly assembly, Church of the firstborn that are written in the heavens, spirits of righteous men and prophets, souls of martyrs and Apostles. Angels, Archangels, Thrones, Dominions, Principalities, Authorities, and Powers dread, also the many-eyed Cherubim, and the six-winged Seraphim, which with twain of their wings cover their faces, and with twain their feet, and with twain do fly, crying one to another with unresting lips, in unceasing praises, singing with loud voice the triumphant hymn of Thy majestic glory, shouting, and glorifying, and crying aloud, and saying,—Holy, Holy, Holy, O Lord of Hosts, heaven and earth are full of Thy glory. Hosanna in the highest; blessed is He that cometh in the name of the Lord; Hosanna in the highest.”}; of sun and moon; of stars and all the creation, rational and irrational, visible and invisible; of Angels, Archangels, Virtues, Dominions, Principalities, Powers, Thrones; of the Cherubim with many faces: in effect repeating that call of David’s \emph{Magnify the Lord with me}\footnote{Ps. xxxiv. 3.}. We make mention also of the Seraphim, whom Esaias in the Holy Spirit saw standing around the throne of God, and with two of their wings veiling their face, and with twain their feet, while with twain they did fly, crying \emph{Holy, Holy, Holy, is the Lord of Sabaoth}\footnote{Is. vi. 2, 3.}. For the reason of our reciting this confession of God\footnote{\textgreek{θεολογίαν}, “the doctrine of the Godhead,” either of the Son in particular, or, as here, of the whole Trinity: cf. Athanas. \emph{Contra Arianos}, Or. i. § 18: \textgreek{νῦν ἐν τρίαδι ἡ θεολογία τελεία ἐστίν}.}, delivered down to us from the Seraphim, is this, that so we may be partakers with the hosts of the world above in their Hymn of praise.

7. Then having sanctified ourselves by these spiritual Hymns, we beseech the merciful God to send forth His Holy Spirit upon the gifts lying before Him; that He may make the Bread the Body of Christ, and the Wine the Blood of Christ\footnote{In the Liturgy of St. James the Triumphal Hymn is followed by the ‘Recital of the work of Redemption,’ and of ‘the Institution,’ by the ‘Great Olbation,’ and then by the ‘Invocation,’ as follows: “Have mercy upon us, O God, after Thy great mercy, and send forth on us, and on these gifts here set before Thee, Thine all-holy Spirit,.…that He may come, and by His holy, good, and glorious advent (\textgreek{παρουσίᾳ}) may sanctify this Bread and make it the holy Body of Thy Christ (\emph{Amen}), and this Cup the precious Blood of Thy Christ” (\emph{Amen}). In Cat. xix. 7, Cyril calls this prayer “the holy Invocation of the Adorable Trinity,” and in xxi. 3, “the Invocation of the Holy Ghost.”}; for whatsoever the Holy Ghost has touched, is surely sanctified and changed.

8. Then, after the spiritual sacrifice, the bloodless service, is completed, over that sacrifice of propitiation\footnote{See Index, “Sacrifice,” and the reference there to the Introduction. Compare Athenagoras (\emph{Apol}. c. xiii.): “What have I to do with burnt-offerings, of which God has no need? Though indeed it behoves us to bring a bloodless sacrifice, and the \emph{reasonable service}.”} we entreat God for the common peace of the Churches, for the welfare of the world\footnote{Cyril here gives a brief summary of the “Great Intercession,” in which, according to the common text of the Liturgy of St. James, there is a suffrage “for the peace and welfare (\textgreek{εὐστάθεια}) of the whole world, and of the holy Churches of God.” Mr. Hammond thinks that it has been taken from the Deacon’s Litany, and repeated by mistake in the Great Intercession. But from Chrysostom’s language (\emph{In Ep. ad Phil}. Hom. iii. p. 218; Guame, T. xi. p. 251), we must infer that the prayer \textgreek{ὑπὲρ εἰρήνης καὶ εὐσταθείας τοῦ κόσμου} formed part of the ‘Great Intercession’ in his Liturgy, as it does in the Clementine (\emph{Apost. Constit}. VIII. § 10).}; for kings; for soldiers and allies; for the sick; for the afflicted; and, in a word, for all who stand in need of succour we all pray and offer this sacrifice.

9. Then we commemorate also those who have fallen asleep before us, first Patriarchs, Prophets, Apostles, Martyrs, that at their prayers and intercessions God would receive our petition\footnote{In the Liturgies of St. James and St. Mark, and in the Clementine, there are similar commemorations of departed saints, especially “patriarchs, prophets, apostles, martyrs,” but nothing corresponding to the words, “that at their prayers and intercessions God would receive our petition.” See Index, \emph{Prayer} and \emph{Intercession}.}. Then on behalf also of the Holy Fathers and Bishops who have fallen asleep before us, and in a word of all who in past years have fallen asleep among us, believing that it will be a very great benefit to the souls\footnote{So Chrysostom (\emph{In} 1 \emph{Cor.} Hom. 41, p. 457 A): “Not in vain was this rule ordained by the Apostles, that in the dread Mysteries remembrance should be made of the departed: for they knew that it is a great gain to them, and a great benefit.”}, for whom the supplication is put up, while that holy and most awful sacrifice is set forth.

10. And I wish to persuade you by an illustration. For I know that many say, what is a soul profited, which departs from this world either with sins, or without sins, if it be commemorated in the prayer? For if a king were to banish certain who had given him offence, and then those who belong to them\footnote{\textgreek{οἱ τούτοις διαφέροντες}. “Hesychius, \textgreek{Διαφέρει, ἀνήκει}. Ubi Kusterus ait, \textgreek{ἀνήκει}, id est. “\emph{pertinet},” vel “\emph{attinet}” Routh, \emph{Scriptor. Eccles. Opuscula}, p. 441). Dr. Routh’s note refers to \emph{Nicæni Conc}. Can. xvi.: \textgreek{ὑφαρπάσαι τὸν τῷ ἑτέρῳ διαφέροντα}. Cf. \emph{Synodi Nic. ad Alexandrinos Epist.}: \textgreek{διαφέροντα τῇ Αἰγύπτῳ καὶ τῂ ἁγιωτάτῃ ᾽Αλεξανδρέων ἐκκλησία}.} should weave a crown and offer it to him on behalf of those under punishment, would he not grant a remission of their penalties? In the same way we, when we offer to Him our supplications for those who have fallen asleep, though they be sinners, weave no crown, but offer up Christ sacrificed for our sins\footnote{According to the Ben. Ed. the meaning is not “We offer Christ, who was sacrificed for our sins,” but “We offer for our sins Christ sacrificed,” i.e. “Christ lying on the altar as a victim sacrificed,” in allusion to Apoc. V. 6, 12. See Index, “Sacrifice.”}, propitiating our merciful God for them as well as for ourselves.

11. Then, after these things, we say that Prayer which the Saviour delivered to His own disciples, with a pure conscience entitling God our Father, and saying, \emph{Our Father, which art in heaven}. O most surpassing loving-kindness of God! On them who revolted from Him and were in the very extreme of misery has He bestowed such a complete forgiveness of evil deeds, and so great participation of grace, as that they should even call Him Father. \emph{Our Father, which art in heaven}; and they also are a heaven who \emph{bear the image of the heavenly}\footnote{1 Cor. xv. 49.}, in whom is God, \emph{dwelling and walking in them}\footnote{2 Cor. vi. 16.}.

12. \emph{Hallowed be Thy Name}. The Name of God is in its nature holy, whether we say so or not; but since it is sometimes profaned among sinners, according to the words, \emph{Through you My Name is continually blasphemed among the Gentiles}\footnote{Is. lii. 5; Rom. ii. 24.}, we pray that in us God’s Name may be hallowed; not that it comes to be holy from not being holy, but because it becomes holy in us, when we are made holy, and do things worthy of holiness.

13. \emph{Thy kingdom come}. A pure soul can say with boldness, \emph{Thy kingdom come}; for he who has heard Paul saying, \emph{Let not therefore sin reign in your mortal body}\footnote{Rom. vi. 12.}, and has cleansed himself in deed, and thought, and word, will say to God, \emph{Thy kingdom come}.

14. \emph{Thy will be done as in heaven so on earth}. God’s divine and blessed Angels do the will of God, as David said in the Psalm, \emph{Bless the Lord, all ye Angels of His, mighty in strength, that do His pleasure}\footnote{Ps. ciii. 20.}. So then in effect thou meanest this by thy prayer, “as in the Angels Thy will is done, so likewise be it done on earth in me, O Lord.”

15. \emph{Give us this day our substantial bread}. This common bread is not substantial bread, but this Holy Bread is substantial, that is, appointed for the substance of the soul\footnote{“It is manifest that the author derives the word \textgreek{ἐπιούσιος} from the two words \textgreek{ἐπί} and \textgreek{οὐσία}, as do many others: although the explanation which derives it from \textgreek{ἐπιούσῃ ἡμέρᾳ} is more probable. We render it “substantial” in accordance with Cyril’s meaning, with which the word “super-substantial does not agree” (Ben. Ed.).}. For this Bread \emph{goeth} not \emph{into the belly and is cast out into the draught}\footnote{Matt. xv. 17.}, but is distributed into thy whole system for the benefit of body and soul\footnote{Cat. xxii. § 3, note 1. Ben. Ed. “We are not to think that Cyril supposed the Body of Christ to be distributed and digested into our body; but in the usual way of speaking he attributes to the Holy Body that which belongs only to the species under which It is hidden. Nor does he deny that those species pass into the draught, but only the Body of Christ.” Cf. Iren. V. ii. 2, 3, and “Eucharistic Doctrine” in the Introduction.}. But by \emph{this day}, he means, “each day,” as also Paul said, \emph{While it is called to-day}\footnote{Heb. iii. 15.}.

16. \emph{And forgive us our debts as we also forgive our debtors}. For we have many sins. For we offend both in word and in thought, and very many things we do worthy of condemnation; and \emph{if we say that we have no sin}, we lie, as John says\footnote{1 John i. 8. \emph{We deceive ourselves}.}. And we make a covenant with God, entreating Him to forgive us our sins, as we also forgive our neighbours their debts. Considering then what we receive and in return for what, let us not put off nor delay to forgive one another. The offences committed against us are slight and trivial, and easily settled; but those which we have committed against God are great, and need such mercy as His only is. Take heed therefore, lest for the slight and trivial sins against thee thou shut out for thyself forgiveness from God for thy very grievous sins.

17. \emph{And lead us not into temptation, O Lord}. Is this then what the Lord teaches us to pray, that we may not be tempted at all? How then is it said elsewhere, “a man untempted, is a man unproved\footnote{Tertull. \emph{De Bapt}. c. 20: “For the word had gone before ‘that no one untempted should attain to the celestial kingdoms.’” \emph{Apost. Const}. II. viii.: “The Scripture says, ‘A man that is a reprobate (\textgreek{ἀδόκιμος}) is not tried (\textgreek{ἀπείραστος}) by God.’” Resch, \emph{Agrapha}, Logion 26, p. 188, quotes allusions to the saying in Jas. i. 12, 13; 2 Cor. xiii. 5, 6, 7, and concludes that it was recorded as a saying of our Lord in one of the un-canonical gospels (Luke i. 1), where it occurred in the context of the incident narrated in Matt. xxvi. 41, Mark xiv. 38.};” and again, \emph{My brethren, count it all joy when ye fall into divers temptations}\footnote{Jas. i. 2.}? But does perchance the entering into temptation mean the being overwhelmed by the temptation? For temptation is, as it were, like a winter torrent difficult to cross. Those therefore who are not overwhelmed in temptations, pass through, shewing themselves excellent swimmers, and not being swept away by them at all; while those who are not such, enter into them and are overwhelmed. As for example, Judas having entered into the temptation of the love of money, swam not through it, but was overwhelmed and was strangled\footnote{\textgreek{ἀπεπνίγη}. Matt. xxvii. 5: \textgreek{ἀπήγξατο}.} both in body and spirit. Peter entered into the temptation of the denial; but having entered, he was not overwhelmed by it, but manfully swam through it, and was delivered from the temptation\footnote{Compare the description of Peter’s repentance in Cat. ii. 19.}. Listen again, in another place, to a company of unscathed saints, giving thanks for deliverance from temptation, \emph{Thou, O God hast proved us; Thou hast tried us by fire like as silver is tried. Thou broughtest us into the net; Thou layedst afflictions upon our loins. Thou hast caused men to ride over our heads; we went through fire and water; and thou broughtest us out into a place of rest}\footnote{Ps. lxvi. 10–12.}. Thou seest them speaking boldly in regard to their having passed through and not been pierced\footnote{For \textgreek{ἐμπαρῆναι} the Ben. Ed. conjectures \textgreek{ἐμπαγῆναι} “to have been stuck fast.”}. \emph{But Thou broughtest us out into a place of rest}; now their coming into a place of rest is their being delivered from temptation.

18. \emph{But deliver us from the evil}. If \emph{Lead us not into temptation} implied the not being tempted at all, He would not have said, \emph{But deliver us from the evil}. Now evil is our adversary the devil, from whom we pray to be delivered\footnote{Cyril is here a clear witness for the reference of \textgreek{τοῦ πονηροῦ} to “the wicked one.”}. Then after completing the prayer thou sayest, \emph{Amen}\footnote{From § 14, \textgreek{εὐχόμενος τοῦτο λέγεις}, it seems probable that the whole Prayer was said by the people as well as by the Priest. See Introduction, “Eucharistic Rites.”}; by this \emph{Amen}, which means “So be it,” setting thy seal to the petitions of the divinely-taught prayer.

19. After this the Priest says, “Holy things to holy men.” Holy are the gifts presented, having received the visitation of the Holy Ghost; holy are ye also, having been deemed worthy of the Holy Ghost; the holy things therefore correspond to the holy persons\footnote{Compare Waterland on this passage, c. X. p 688.}. Then ye say, “One is Holy, One is the Lord, Jesus Christ\footnote{\emph{Apost. Const}. VIII. c. xiii: “Let the Bishop speak thus to the people: Holy things for holy persons. And let the people answer: There is One that is holy; there is one Lord, one Jesus Christ, blessed for ever, to the glory of God the Father.” The Liturgies of St. James and of Constantinople have nearly the same words: in the Liturgy of St. Mark the answer of the people is: One Father holy, one Son holy, one Spirit holy, in the unity of the Holy Spirit.}.” For One is truly holy, by nature holy; we too are holy, but not by nature, only by participation, and discipline, and prayer.

20. After this ye hear the chanter inviting you with a sacred melody to the communion of the Holy Mysteries, and saying, \emph{O taste and see that the Lord is good}\footnote{Ps. xxxiv. 9. In the \emph{Apostolic Constitutions} the “Sancta Sanctis” and its response are immediately followed by the “Gloria in excelsis,” and the “Hosanna.” Then the Clergy partake, and there follows a direction that this Psalm xxxiv. is to be said while all the rest are partaking. In the Liturgy of Constantinople there is the direction: “The Choir sings the communion antiphon (\textgreek{τὸ κοινωνικόν}) of the day or the saint.”}. Trust not the judgment to thy bodily palate\footnote{For \textgreek{μὴ ἐπιτρέπητε}, probably an itacism, we should read \textgreek{μὴ ἐπιτρέπεται}, as a question, the propriety of the change being indicated by the answer \textgreek{οὐχί}. “Is the judgment of this entrusted to the bodily palate? No, but, \&c.”} no, but to faith unfaltering; for they who taste are bidden to taste, not bread and wine, but the anti-typical\footnote{\textgreek{ἀντιτύπου σώματος}, “the antitypical Body,” not “the antitype of the Body,” which would require \textgreek{τοῦ σώματος}. Cf. Cat. xxi. § 1, note 6.} Body and Blood of Christ.

21. In approaching\footnote{Cat. xviii. 32: “with what reverence and order you must go from Baptism to the Holy Altar of God.”} therefore, come not with thy wrists extended, or thy fingers spread; but make thy left hand a throne for the right, as for that which is to receive a King\footnote{Cyril appears to be the earliest authority for thus placing the hands in the form of a Cross. A similar direction is given in the 101st Canon of the Trullan Council (692), and by Joh. Damasc. (\emph{De Fid. Orthod}. iv. 14). \emph{Dict. Chr. Ant.} “\emph{Communion}.” That the communicant was to receive the Bread in his own hands is clear from the language of Cyril and other Fathers. Cf. Clem. Alex. \emph{Strom}. I. c. i. § 5: “Some after dividing the Eucharist according to custom allow each of the laity himself to take his part.” See the passage of Origen quoted in the next note, and Tertull. \emph{Cor. Mil}. c. iii. “The Sacrament of the Eucharist, which the Lord commanded both (to be taken) at meal-times and by all, we take even in assemblies before dawn, and from the hand of none but the presidents.”}. And having hollowed thy palm, receive the Body of Christ, saying over it, \emph{Amen}. So then after having carefully hallowed thine eyes by the touch of the Holy Body, partake of it; giving heed lest thou lose any portion thereof\footnote{Origen. \emph{Hom. xiii. in Exod.} § 3: “I wish to admonish you by examples from your own religion: ye, who have been accustomed to attend the Sacred Mysteries, know how, when you receive the Body of the Lord, you guard it with all care and reverence, that no little part of it fall down, no portion of the consecrated gift slip away. For you believe yourselves guilty, and rightly so believe, if any part thereof fall through carelessness.”}; for whatever thou losest, is evidently a loss to thee as it were from one of thine own members. For tell me, if any one gave thee grains of gold, wouldest thou not hold them with all carefulness, being on thy guard against losing any of them, and suffering loss? Wilt thou not then much more carefully keep watch, that not a crumb fall from thee of what is more precious than gold and precious stones?

22. Then after thou hast partaken of the Body of Christ, draw near also to the Cup of His Blood; not stretching forth thine hands, but bending\footnote{\textgreek{κύπτων}, not kneeling, but standing in a bowing posture. Cf. Bingham, XV. c. 5, § 3.}, and saying with an air of worship and reverence, \emph{Amen}\footnote{\emph{Apost. Const}. VIII. c. 13: “Let the Bishop give the Oblation (\textgreek{προσφοράν}) saying, \emph{The Body of Christ}. And let him that receiveth say, \emph{Amen}. And let the Deacon hold the Cup, and when he delivers it say, \emph{The Blood of Christ, the Cup of Life}. And let him that drinketh say, \emph{Amen}.”}, hallow thyself by partaking also of the Blood of Christ. And while the moisture is still upon thy lips, touch it with thine hands, and hallow thine eyes and brow and the other organs of sense\footnote{Cat. xxi. 3, note 8.}. Then wait for the prayer, and give thanks unto God, who hath accounted thee worthy of so great mysteries\footnote{In the Liturgy of St. James, after all have communicated, “\emph{The Deacons and the People say}: Fill our mouths with Thy praise, O Lord, and fill our lips with joy, that we may sing of Thy glory, of Thy greatness, all the day. \emph{And again}: We render thanks to Thee, Christ our God, that Thou hast accounted us worthy to partake of Thy Body and Blood, \&c.”}.

23. Hold fast these traditions undefiled and, keep yourselves free from offence. Sever not yourselves from the Communion; deprive not yourselves, through the pollution of sins, of these Holy and Spiritual Mysteries. \emph{And the God of peace sanctify you wholly; and may your spirit, and soul, and body be preserved entire without blame at the coming of our Lord Jesus Christ}\footnote{1 Thess. v. 23.}:—To whom be glory and honour and might, with the Father and the Holy Spirit, now and ever, and world without end. Amen.
